% Merged research report. Standalone: no external graphics, no BibTeX run.
\documentclass[11pt]{article}
\usepackage[T1]{fontenc}
\usepackage[utf8]{inputenc}
\usepackage{lmodern}
\usepackage[margin=1in,headheight=15pt]{geometry}
\usepackage{amsmath,amssymb,amsthm,mathtools}
\usepackage{microtype}
\usepackage{booktabs,array,longtable,tabularx}
\usepackage{xcolor}
\usepackage{enumitem}
\usepackage{fancyhdr}
\usepackage{needspace}
\usepackage{xurl}
\usepackage{hyperref}
\usepackage{bookmark}
\definecolor{ink}{HTML}{233A45}
\definecolor{olive}{HTML}{596547}
\hypersetup{colorlinks=true,linkcolor=ink,citecolor=olive,urlcolor=ink,
  breaklinks=true,
  pdftitle={Holonomic Rigidity for Entire Hahn Functions},
  pdfauthor={Merged research report prepared for Vladimir Reshetnikov},
  pdfsubject={Arbitrary-rank Hahn fields, strong summability, linear
              differential and q-difference equations, algebraic
              differential equations, coefficient fields}}
\urlstyle{same}
\pagestyle{fancy}
\fancyhf{}
\fancyhead[L]{\small Holonomic rigidity for entire Hahn functions}
\fancyhead[R]{\small A sharp order-unit criterion}
\fancyfoot[C]{\thepage}
\renewcommand{\headrulewidth}{0.3pt}
\setlength{\parskip}{3.5pt}
\setlength{\parindent}{1.15em}
\setlength{\emergencystretch}{2em}
\widowpenalty=10000\clubpenalty=10000
\setcounter{tocdepth}{1}
\numberwithin{equation}{section}
\allowdisplaybreaks[2]
\newtheorem{maintheorem}{Theorem}
\renewcommand{\themaintheorem}{\Alph{maintheorem}}
\newtheorem{theorem}{Theorem}[section]
\newtheorem{lemma}[theorem]{Lemma}
\newtheorem{proposition}[theorem]{Proposition}
\newtheorem{corollary}[theorem]{Corollary}
\theoremstyle{definition}
\newtheorem{definition}[theorem]{Definition}
\newtheorem{convention}[theorem]{Convention}
\newtheorem{example}[theorem]{Example}
\newtheorem{question}[theorem]{Question}
\theoremstyle{remark}
\newtheorem{remark}[theorem]{Remark}
\newcommand{\C}{\mathbb C}
\newcommand{\R}{\mathbb R}
\newcommand{\Q}{\mathbb Q}
\newcommand{\N}{\mathbb N}
\newcommand{\Z}{\mathbb Z}
\newcommand{\No}{\mathbf{No}}
\newcommand{\K}{K_\Gamma}
\newcommand{\E}{\mathcal E_\Gamma}
\newcommand{\supp}{\operatorname{supp}}
\newcommand{\lc}{\operatorname{lc}}
\newcommand{\res}{\operatorname{res}}
\newcommand{\cf}{\operatorname{cf}}
\newcommand{\Dom}{\operatorname{Dom}_{\mathrm{str}}}
\newcommand{\fall}[2]{(#1)_{\underline{#2}}}
\newcommand{\abs}[1]{\lvert#1\rvert}
\newcommand{\repo}{\texttt{VladimirReshetnikov/Surreal}}
% Nonlinear part (source 10): any characteristic-zero coefficient field k.
\newcommand{\KK}{\mathbb K}
\newcommand{\EK}{\mathcal E_\Gamma(k)}
\newcommand{\Sfam}{\mathcal S}
\newcommand{\wt}{\operatorname{wt}}
\newcommand{\inpoly}{\operatorname{in}}
% Coefficient-field part (sources 11 and 12).
\newcommand{\FE}{\operatorname{Frac}\EK}
\newcommand{\trdeg}{\operatorname{trdeg}}
\newcommand{\dtrdeg}{\operatorname{dtrdeg}}
\newcommand{\rrk}{\operatorname{rrk}}
\newcommand{\cvr}{\operatorname{cvr}}
\newcommand{\red}{\operatorname{red}}
\newcommand{\code}{\operatorname{code}}
\newcommand{\ordz}{\operatorname{ord}_z}
\newcommand{\coeffz}[2]{[z^{#1}]\,#2}
\newcommand{\jet}[2]{\mathcal J_{#1,#2}}
% Theta-descent part (source 13).
\newcommand{\Oz}{\mathbf{Oz}}
\setlist[itemize]{leftmargin=1.5em,itemsep=2pt,topsep=4pt}
\setlist[enumerate]{leftmargin=1.8em,itemsep=3pt,topsep=4pt}

\begin{document}
\hypersetup{pageanchor=false}
\begin{titlepage}
\centering
\vspace*{0in}
{\color{olive}\large\sffamily SURREAL AND SURCOMPLEX RESEARCH}\par
\vspace{0.08in}
{\color{ink}\huge\bfseries Holonomic Rigidity\\[6pt] for Entire Hahn Functions}\par
\vspace{0.1in}
{\Large A sharp order-unit criterion for dilations,\\[4pt]
with uniform exclusion bounds,\\[4pt]
and the exact boundary of nonlinear differential rigidity}\par
\vspace{0.1in}
{\large Merged research report\quad\textperiodcentered\quad 22--23 September 2026}\par
\vspace{0.06in}
\begin{minipage}{0.94\textwidth}
\begin{abstract}
Fix a nonzero set-sized ordered abelian group $\Gamma$ and the complex Hahn
field $\K=\C((t^\Gamma))$. \emph{Entire} means strong Hahn summability of the
evaluation family at every point of this one field. We prove that every
strongly entire D-finite power series is a polynomial, in one or in finitely
many variables, and that a fixed differential operator already excludes every
nonterminating formal solution on an explicit exterior valuation region,
with a common degree bound.

For dilations there is an exact alternative. If $q\in\K^\times$ has infinite
multiplicative order, a nonpolynomial strongly entire solution of a nonzero
linear $q$-difference equation with polynomial coefficients exists if and only
if $v(q)\ne0$ and $\abs{v(q)}$ is an order unit of $\Gamma$. A nonzero
valuation is not enough: its integer multiples must be cofinal among all
scales. The proof covers units whose residues are roots of unity, and
coefficients and parameters with arbitrary well-ordered Hahn supports. A
partial theta series proves sharpness, and its exact strong evaluation domain
is computed for an arbitrary Hahn parameter. Quantitatively, the recurrence
supplies an exclusion bound that is attained on its own boundary, and a
coarsened exclusion bound valid for every noncofinal dilation. A mixed
differential--dilation extension is proved for nonzero noncofinal dilation
valuations.

The common mechanism is an elementary finite-step escape chain: bounded
valuation costs in a coefficient recurrence force a subsequence of evaluated
leading exponents that never increases. This contradicts strong summability
without replacing an arbitrary-rank valuation by a real-valued norm. In an
explicit surreal workspace with countable cofinality and no order unit, the
series $\sum_{n\ge0}\omega^{-\omega^n}z^n$ is strongly entire yet satisfies no
nonzero linear differential equation and no nonzero linear dilation equation
for any nontorsion parameter of that workspace.

Over $k((t^\Gamma))$, $k$ any field of characteristic zero, later sources
show: first-order algebraic differential equations have only polynomial
strongly entire solutions; without an order unit every nonpolynomial strongly
entire series is differentially transcendental, jointly with its dilates by
parameters whose quotients are not roots of unity; with one, a
Chebyshev-descended theta series has minimal differential order three; for
every $\Gamma$ no such series satisfies an equation of order $\le2$, or of
order three and jet degree $\le3$; linear mixed rigidity holds at unit
dilations; a dominant argument $f(P(z))$, $\deg P\ge2$, forces
polynomiality for every $\Gamma$; and unit-root recurrences have finitely
generated coefficient algebras. Divisibility of $\Gamma$ is
\emph{not} assumed, except in three threshold constructions that state it
(Section~\ref{hol:sub:divisibility}).
\end{abstract}
\end{minipage}
\vfill
\begin{minipage}{0.94\textwidth}
\small
\textbf{Status.} An AI-assisted research report; priority and independent
refereeing are not certified. Lean covers support and escape lemmas,
exponential and partial-theta domains, polynomial-orbit valuations,
constant-point avoidance, inward stability and torsion covariance.
The full classifications and the theorem packages of the ten later parts
remain pending. The
variable derivative is not the
Berarducci--Mantova derivation on surreal coefficients.
\end{minipage}
\end{titlepage}
\hypersetup{pageanchor=true}

\tableofcontents
\clearpage

\section{The question and the principal results}\label{hol:sec:intro}

\subsection{Why the question is scale sensitive}

A familiar way to specify a special function is to give a finite linear
differential equation. Another is to give a finite relation between its values
at $z,qz,q^2z,\ldots$. Over the complex numbers both descriptions produce
nonpolynomial entire functions. Over a surreal Hahn workspace, ``entire'' asks
for evaluation not only at ordinary complex constants but at every infinite
and infinitesimal scale available in that workspace. In a Hahn field an
infinite expression is legal only when its supports satisfy an order condition
and a finiteness condition, and an equation can force either condition to fail
at a sufficiently large argument. The question of this report is what a finite
linear equation can accommodate under that stronger global requirement.
Sections~\ref{hol:nl:sec:setting}--\ref{hol:nl:sec:consequences} ask the same
of a finite \emph{nonlinear} algebraic differential equation, and answer it in
order one; Sections~\ref{hol:cf:sec:setting}--\ref{hol:cf:sec:families}
answer it in every order when the value group has no order unit; and
Sections~\ref{hol:td:sec:setting}--\ref{hol:td:sec:boundary} show that with an
order unit the answer changes: an explicit nonpolynomial strongly entire
series satisfies an algebraic differential equation of order three.
Sections~\ref{hol:ot:sec:setting}--\ref{hol:ot:sec:theta} show that three is
the least such order: for every value group, no nonpolynomial strongly entire
series satisfies an equation of order at most two.
Sections~\ref{hol:ud:sec:setting}--\ref{hol:ud:sec:mixed} return to linear
equations that mix derivatives and dilations, and settle them by the relative
scale of the dilations; Sections~\ref{hol:nr:sec:setting}--%
\ref{hol:nr:sec:boundary} exclude equations of order three and total degree
at most three in the derivatives.
Sections~\ref{hol:th:sec:setting}--\ref{hol:th:sec:scope} show, for
$\Gamma=\R$, that minimal orders are unbounded and that the differentially
algebraic strongly entire series generate a field of transcendence degree
$2^{\aleph_0}$; Sections~\ref{hol:fh:sec:setting}--%
\ref{hol:fh:sec:realizations} give, at a single order-unit scale, explicit
series with coefficients in $\Q(t^\mu)$ that satisfy no algebraic differential
equation at all, with their whole relation ideals.
Sections~\ref{hol:pc:sec:setting}--\ref{hol:pc:sec:arithmetic} replace the
dilation by a polynomial argument of degree at least two, and there rigidity
holds for every value group, with explicit degree bounds.
Sections~\ref{hol:rc:sec:setting}--\ref{hol:rc:sec:surreal} turn from
functions to the ordinary linear recurrences with Hahn or surreal values that
underlie all these proofs, and read their individual coefficients: tame, in
one finitely generated algebra, when the characteristic roots are units, and
arbitrary for a first-order omnific recurrence.

The answer depends decisively on which operator is used, and for dilation
equations it depends further on whether the valuation of the particular
parameter reaches every scale of the value group.

\subsection{Conventions}\label{hol:sub:conventions}

We work with a fixed ordered abelian group $\Gamma$ and
\[
   \K=\C((t^\Gamma)).
\]
An element of $\K$ is a formal sum $a=\sum_{\gamma}a_\gamma t^\gamma$ with
$a_\gamma\in\C$ whose support $\supp(a)=\{\gamma:a_\gamma\ne0\}$ is a
well-ordered set. For $a\ne0$ put
\[
 v(a)=\min\supp(a),\qquad \lc(a)=a_{v(a)}\in\C^\times,
 \qquad \res(a)=\lc(a)\text{ when }v(a)=0,
\]
and put $v(0)=+\infty$ only where convenient; this symbol is never subtracted.
Products and sums are the usual Hahn operations, with
\begin{equation}\label{hol:eq:valuationlaws}
 v(ab)=v(a)+v(b),\qquad
 v(a+b)\ge\min\{v(a),v(b)\},
\end{equation}
with equality in the second law when the two valuations differ. The
coefficient field $\C$ is trivially valued, so every nonzero ordinary integer
and every factorial has valuation zero. Throughout, $\N=\{0,1,2,\ldots\}$, and
$\abs\lambda$ means $\max(\lambda,-\lambda)$ in the ordered group.

\begin{convention}[Standing hypothesis on the value group]
\label{hol:conv:group}
$\Gamma$ is a nonzero set-sized ordered abelian group. Divisibility is
\emph{not} part of this standing hypothesis. Every statement below that needs
divisibility says so in its own statement; those statements are collected in
Section~\ref{hol:sub:divisibility}.
\end{convention}

\begin{definition}[Strong summability, evaluation domain, entireness]
\label{hol:def:entire}
A set-indexed family $(b_i)_{i\in I}$ in $\K$ is \emph{strongly Hahn summable}
if $\bigcup_i\supp(b_i)$ is well ordered and each $\gamma\in\Gamma$ lies in the
supports of only finitely many members. Its sum is then defined
coefficientwise, by a finite sum in $\C$ at each exponent. For
$f=\sum_{n\ge0}a_nz^n\in\K[[z]]$ let $\Dom(f)$ be the set of $x\in\K$ for
which $(a_nx^n)_{n\ge0}$ is strongly summable, and let
\[
 \E=\{f:\Dom(f)=\K\}.
\]
In $d$ variables the family to be tested is $(a_\alpha x^\alpha)_{\alpha\in\N^d}$.
\end{definition}

The word \emph{entire} always means Definition~\ref{hol:def:entire} relative to
one fixed workspace. It does not mean that a complex coefficient function is
entire, it does not mean convergence in the fine topology inherited from the
full surreal class, and it is never a claim about all of $\No[i]$. The
definition also concerns the family \emph{before} cancellation between
different indices: one cannot first cancel a nonsummable family and then
declare it summable, and convergence of a numerical coefficient series does
not license infinitely many contributions at one Hahn exponent.

\begin{definition}[Order unit]\label{hol:def:orderunit}
A positive $\mu\in\Gamma$ is an \emph{order unit} if for every $\gamma\in\Gamma$
there is $N\in\N$ with $\abs\gamma\le N\mu$; equivalently, if the positive
integer multiples of $\mu$ are cofinal in $\Gamma$. Zero is not an order unit.
\end{definition}

\begin{definition}[The two operator classes]\label{hol:def:operators}
Let $D_z$ be the formal derivative in $z$, so $D_z(a)=0$ for every $a\in\K$,
and for $q\in\K^\times$ let $\sigma_qf(z)=f(qz)$. A series $f\in\K[[z]]$ is
\emph{D-finite over $\K(z)$} if
\begin{equation}\label{hol:eq:dlinear}
 \sum_{r=0}^R p_r(z)D_z^rf(z)=0,\qquad p_r\in\K[z],\qquad (p_0,\ldots,p_R)\ne0,
\end{equation}
and \emph{linearly $q$-difference finite} if
\begin{equation}\label{hol:eq:qlinear}
 \sum_{\ell=0}^r p_\ell(z)f(q^\ell z)=0,\qquad p_\ell\in\K[z],
 \qquad (p_0,\ldots,p_r)\ne0.
\end{equation}
These are formal identities in $\K[[z]]$, not assertions about limits of
partial sums. Rational coefficients give the same notions after denominators
are cleared.
\end{definition}

We say ``linearly $q$-difference finite'' rather than silently identifying
several meanings of ``$q$-holonomic''. In the sequence literature $q$ is often
an indeterminate with coefficients in $\Q(q)$; here $q$ is a \emph{fixed} Hahn
element, possibly with infinitely many terms. The differential terminology is
the usual characteristic-zero one~\cite{stanley}. \emph{Nontorsion} and
``of infinite multiplicative order'' are used interchangeably.

\subsection{The principal results}

\begin{maintheorem}[Differential rigidity, with a uniform exterior obstruction]
\label{hol:main:ode}
Let $\Gamma$ be as in Convention~\ref{hol:conv:group} and let
$0\ne L=\sum_{r,k}c_{r,k}z^kD_z^r$ be a nonzero linear differential operator
with coefficients in $\K[z]$. There are $B_L\in\Gamma$ with $B_L\ge0$ and
$N_L\in\N$, both determined by the finitely many coefficients of $L$, such
that every formal solution of $Lf=0$ that is strongly evaluable at even one
$x\ne0$ with $v(x)\le-B_L$ has degree less than $N_L$. Consequently
\[
 \{f\in\E:f\text{ is D-finite over }\K(z)\}=\K[z].
\]
\end{maintheorem}

\begin{maintheorem}[Sharp dilation criterion]\label{hol:main:q}
Let $\Gamma$ be as in Convention~\ref{hol:conv:group} and let $q\in\K^\times$
have infinite multiplicative order. The following are equivalent.
\begin{enumerate}[label=\textup{(\roman*)}]
\item There is a nonpolynomial $f\in\E$ that is linearly $q$-difference finite
over $\K(z)$.
\item $v(q)\ne0$ and $\abs{v(q)}$ is an order unit of $\Gamma$.
\end{enumerate}
When \textup{(ii)} fails, every fixed equation \eqref{hol:eq:qlinear} has a
uniform exterior obstruction and a uniform polynomial degree bound, exactly as
in Theorem~\ref{hol:main:ode}. When \textup{(ii)} holds, an explicit partial
theta series with parameter $q$ or $q^{-1}$ is a nonpolynomial entire
solution.
\end{maintheorem}

\begin{maintheorem}[Several-variable rigidity]\label{hol:main:multi}
Let $\Gamma$ be as in Convention~\ref{hol:conv:group} and let
$F\in\K[[z_1,\ldots,z_d]]$ be strongly Hahn summable at every point of $\K^d$.
If the mixed partial derivatives of $F$ span a finite-dimensional vector space
over $\K(z_1,\ldots,z_d)$, then $F$ is a polynomial.
\end{maintheorem}

The fourth principal theorem comes from the third source manuscript and is
proved in Sections~\ref{hol:nl:sec:setting}--\ref{hol:nl:sec:first}. It
allows an arbitrary coefficient field of characteristic zero; see
Convention~\ref{hol:nl:conv:field}, which also explains why that generality
is not transferred to Theorems~\ref{hol:main:ode}--\ref{hol:main:multi}.

\begin{maintheorem}[First-order nonlinear rigidity]\label{hol:nl:main:first}
Let $\Gamma$ be as in Convention~\ref{hol:conv:group}, let $k$ be any field of
characteristic zero, and let $\KK=k((t^\Gamma))$, with strong evaluation and
entireness defined as in Definition~\ref{hol:def:entire} with $k$ in place of
$\C$. For every nonzero $P\in\KK[z,Y_0,Y_1]$, every strongly entire formal
solution $f\in\KK[[z]]$ of
\[
 P(z,f,D_zf)=0
\]
is a polynomial. More strongly, if a nonpolynomial formal solution exists,
there is $\delta_0\in\Gamma_{>0}$, depending on the equation and on finitely
many coefficients of that solution, such that $t^{-\delta}\notin\Dom(f)$ for
every $\delta\ge\delta_0$. Consequently, for every nonpolynomial strongly
entire $f$, the series $f$ and $D_zf$ are algebraically independent over
$\KK(z)$.
\end{maintheorem}

Theorem~\ref{hol:nl:main:first} has no separant, irreducibility or
solved-form hypothesis on $P$. Its exclusion bound depends on the solution,
and Theorem~\ref{hol:nl:thm:nouniform} shows that this cannot be avoided,
in contrast with the operator-only bound of Theorem~\ref{hol:main:ode}. By
itself it does not say that a nonpolynomial entire function is differentially
transcendental in all orders: independence of $f$ and $D_zf$ does not exclude
a relation involving $D_z^2f$. Theorem~\ref{hol:cf:main:entire} below says
this when $\Gamma$ has no order unit, and Theorem~\ref{hol:td:main:boundary}
shows that the hypothesis cannot be dropped: with an order unit a
nonpolynomial strongly entire series can satisfy a relation of order three.
Theorem~\ref{hol:ot:main:second} below extends Theorem~\ref{hol:nl:main:first}
to order two for every $\Gamma$, so that three is the least order;
Question~\ref{hol:q:nonlinear}, once re-scoped to order exactly two, is now
re-scoped to the classification of the series of minimal order three.

The last three principal theorems come from sources~11 and~12 and are proved
in Sections~\ref{hol:cf:sec:setting}--\ref{hol:cf:sec:families}, over the
field $\KK=k((t^\Gamma))$ of Convention~\ref{hol:nl:conv:field} for an
arbitrary field $k$ of characteristic zero; $\EK$ is its set of strongly
entire series, $\FE$ its fraction field inside $\KK((z))$, and
$\jet\lambda jf=(D_z^jf)(\lambda z)$ (Convention~\ref{hol:cf:conv:field}).

\begin{maintheorem}[All-order rigidity without an order unit]
\label{hol:cf:main:entire}
Let $\Gamma$ be as in Convention~\ref{hol:conv:group}, and suppose that
$\Gamma$ has no order unit.
\begin{enumerate}[label=\textup{(\alph*)}]
\item \textup{(Sources~11 and~12.)} Every nonpolynomial $f\in\EK$ is
differentially transcendental over $\KK(z)$: the series $f,D_zf,D_z^2f,\ldots$
are algebraically independent over $\KK(z)$. Equivalently, every strongly
entire solution of a nonzero algebraic differential equation
$P(z,f,D_zf,\ldots,D_z^rf)=0$ over $\KK(z)$, of any finite order and with any
residue corner, is a polynomial.
\item \textup{(Source~12.)} Let $\Lambda\subseteq\KK^\times$ be such that
$\lambda/\lambda'$ is not a root of unity for distinct
$\lambda,\lambda'\in\Lambda$. For every nonpolynomial $f\in\EK$ the family
$\{\jet\lambda jf:\lambda\in\Lambda,\ j\ge0\}$ is algebraically independent
over $\KK(z)$. In particular, for every $q\in\KK^\times$ of infinite
multiplicative order, including $v(q)=0$, the series $(D_z^jf)(q^iz)$,
$i\in\Z$, $j\ge0$, are jointly algebraically independent over $\KK(z)$.
\item \textup{(Source~11.)} The elements of $\FE$ that are differentially
algebraic over $\KK(z)$ are exactly the elements of $\KK(z)$.
\end{enumerate}
\end{maintheorem}

For the coefficient-value rank $\cvr(f)$, the dimension of the rational span
of the values of the nonzero Taylor coefficients \eqref{hol:cf:eq:cvr}, there
is a test for a single series that needs neither entireness nor a hypothesis
on $\Gamma$.

\begin{maintheorem}[Coefficient-value-rank criterion; source~12]
\label{hol:cf:main:rank}
Let $\Gamma$ be as in Convention~\ref{hol:conv:group} and $f\in\KK[[z]]$. If
$\cvr(f)$ is infinite, the family of Theorem~\ref{hol:cf:main:entire}\,(b) is
algebraically independent over $\KK(z)$ for every $\Lambda$ as there.
\end{maintheorem}

\begin{maintheorem}[Order-unit dichotomy for mixed jets; source~12]
\label{hol:cf:main:sharp}
For $\Gamma$ as in Convention~\ref{hol:conv:group}, the following are
equivalent.
\begin{enumerate}[label=\textup{(\roman*)}]
\item $\Gamma$ has an order unit.
\item Some nonpolynomial $f\in\EK$ has all its coefficients in a subfield
$L\supseteq k$ of $\KK$ with $\trdeg_kL<\infty$.
\item For some nonpolynomial $f\in\EK$ and some $q\in\KK^\times$ of infinite
order, the family $\{(D_z^jf)(q^iz):i\in\Z,\ j\ge0\}$ is algebraically
dependent over $\KK(z)$.
\item \textup{(Added by the merge of source~13.)} Some nonpolynomial
$f\in\EK$ is differentially algebraic over $\KK(z)$.
\end{enumerate}
If $\mu$ is an order unit, the partial theta series $\Theta_p$ of
\eqref{hol:eq:theta} with $p=t^\mu$ is a witness for \textup{(ii)} and
\textup{(iii)}: its coefficients lie in $k(p)$, and
$\Theta_p(z)=1+z\Theta_p(pz)$. The series $\mathcal T_p$ of
Theorem~\ref{hol:td:main:boundary} is a witness for
\textup{(ii)}--\textup{(iv)}, and for \textup{(iii)} with every $q$ of infinite
order.
\end{maintheorem}

As proved by source~12, Theorem~\ref{hol:cf:main:sharp} concerned the
\emph{mixed} conclusion: its witness $\Theta_p$ satisfies a dilation equation,
and conditions (i)--(iii) do not say that an order unit produces a
nonpolynomial strongly entire \emph{differentially} algebraic series.
Clause~(iv) says this; it is added by the merge from
Theorem~\ref{hol:td:main:boundary} below, and its proof is in
Section~\ref{hol:td:sub:proofs}. Theorem~\ref{hol:cf:main:entire} answers
Question~\ref{hol:q:nonlinear}, as it then stood,
in every order for groups without an order unit, including the equations with
vanishing residue corner, and its part~(b) answers the same case of
Question~\ref{hol:q:mixedunit}. Both sources prove part~(a) by one argument,
printed once: a finite relation, linearized along its solution, puts all
Taylor coefficients into one finitely generated field
(Theorem~\ref{hol:cf:thm:coefficients}); such a field has a value group of
finite rational rank, which lies in a proper convex subgroup when there is no
order unit, while a nonpolynomial entire series needs coefficient values at
cofinally many scales. Part~(b) adds dilations through the classical
Skolem--Mahler--Lech theorem. Source~11 also proves a coarsening descent for
arbitrary coefficient fields (Theorem~\ref{hol:cf:thm:descent}), corollaries
for rational systems and for Painlev\'e~I in $\FE$, a joint-independence
criterion, and, over $\C((t^{\Gamma_\infty}))$, a continuum-sized family of
strongly entire series with jointly independent jets, so that the fraction
field has differential transcendence degree $2^{\aleph_0}$
(Theorem~\ref{hol:cf:thm:continuum}). The series
$\sum_{n\ge0}\omega^{-\omega^n}z^n$ satisfies no algebraic differential
equation of any order (Theorem~\ref{hol:cf:thm:Fomega}).

The eighth principal theorem comes from source~13 and is proved in
Sections~\ref{hol:td:sec:setting}--\ref{hol:td:sec:boundary}, over
$\KK=k((t^\Gamma))$ for an arbitrary field $k$ of characteristic zero
(Convention~\ref{hol:td:conv:field}). For $\mu>0$ in $\Gamma$ put $p=t^\mu$,
let $C_n$ be the monic Chebyshev polynomials, $C_0=2$, $C_1=z$,
$C_{n+1}=zC_n-C_{n-1}$, so that $C_n(u+u^{-1})=u^n+u^{-n}$, and put
\begin{equation}\label{hol:td:eq:series}
 \mathcal T_p(z)=1+\sum_{n\ge1}p^{n^2}C_n(z).
\end{equation}

\begin{maintheorem}[Exact boundary of differential algebraicity; source~13]
\label{hol:td:main:boundary}
For $\Gamma$ as in Convention~\ref{hol:conv:group} and every field $k$ of
characteristic zero, the following are equivalent.
\begin{enumerate}[label=\textup{(\roman*)}]
\item $\Gamma$ has an order unit.
\item Some nonpolynomial $f\in\EK$ is differentially algebraic over $\KK(z)$.
\item Some nonpolynomial $f\in\EK$ satisfies a nonzero algebraic differential
equation of order three over $\KK(z)$, while $f,D_zf,D_z^2f$ are algebraically
independent over $\KK(z)$ and $D_z^3f$ has degree two over
$\KK(z)(f,D_zf,D_z^2f)$.
\end{enumerate}
If $\mu$ is an order unit, $\mathcal T_p$ is a witness for \textup{(iii)}; it
is not D-finite over $\KK(z)$. For every $\mu>0$ the strong domain of
$\mathcal T_p$ is $\{0\}\cup\{x:v(x)\ge-N\mu\text{ for some }N\in\N\}$, so
$\mathcal T_p\in\EK$ exactly when $\mu$ is an order unit.
\end{maintheorem}

Theorem~\ref{hol:td:main:boundary} answers Question~\ref{hol:q:nonlinear} as it
stood after the merge of sources~11 and~12: with an order unit a nonpolynomial
entire series can satisfy an algebraic differential equation of order at
least two, a differentially algebraic element of $\FE$ need not be rational,
and the absence of an order unit is necessary for differential algebraic
rigidity. With Theorem~\ref{hol:nl:main:first}, the least order of such an
equation over all examples was then two or three, and the question was
re-scoped to order two; Theorem~\ref{hol:ot:main:second} below (source~14)
shows that it is three (Section~\ref{hol:sec:status}). The equation of
$\mathcal T_p$ is
explicit (Theorem~\ref{hol:td:thm:ode}); its residue corner polynomial vanishes
identically while its full corner polynomial is a nonzero constant
(Proposition~\ref{hol:td:prop:corner}). The direction (ii)$\Rightarrow$(i) is
Theorem~\ref{hol:cf:main:entire}\,(a), which source~13 proves again; so
Theorem~\ref{hol:cf:main:entire}\,(a) and~(c) hold exactly when $\Gamma$ has no
order unit (Corollary~\ref{hol:td:cor:sharp}). The series $\mathcal T_p$ is a
Chebyshev descent of the bilateral theta series $\sum_{n\in\Z}p^{n^2}u^n$ along
$z=u+u^{-1}$; the classical triple product gives all its zeros, and in the
surreal specialization $p=\omega^{-1}$ none of them is an omnific integer
(Corollary~\ref{hol:td:cor:omnific}).

The ninth principal theorem comes from source~14 and is proved in
Sections~\ref{hol:ot:sec:setting}--\ref{hol:ot:sec:theta}, again over
$\KK=k((t^\Gamma))$ for an arbitrary field $k$ of characteristic zero
(Convention~\ref{hol:ot:conv:field}).

\begin{maintheorem}[Second-order rigidity and the order-three threshold;
source~14]\label{hol:ot:main:second}
Let $\Gamma$ be as in Convention~\ref{hol:conv:group} and $k$ any field of
characteristic zero.
\begin{enumerate}[label=\textup{(\alph*)}]
\item Every $f\in\EK$ that satisfies a nonzero algebraic differential equation
$P(z,f,D_zf,D_z^2f)=0$ with $P\in\KK(z)[Y_0,Y_1,Y_2]$ is a polynomial.
Equivalently, for every nonpolynomial $f\in\EK$ the series $f$, $D_zf$ and
$D_z^2f$ are algebraically independent over $\KK(z)$.
\item If $\Gamma$ has an order unit, the least order of a nonzero algebraic
differential equation over $\KK(z)$ satisfied by a nonpolynomial element of
$\EK$ is exactly three. If $\Gamma$ has no order unit, no nonpolynomial
element of $\EK$ satisfies one of any order.
\end{enumerate}
\end{maintheorem}

Part~(a) needs no divisibility, rank, algebraic closedness or order-unit
hypothesis, and no condition on the residue corner. Its polynomiality
conclusion contains those of Theorem~\ref{hol:nl:main:first} and of
Theorem~\ref{hol:nl:thm:corner} in orders at most two, but not their explicit
exclusion bounds, and it holds in particular for equations with $\chi_P=0$,
such as \eqref{hol:nl:eq:degenerate}. The new ingredient is the
lower bound three; the witness of order three and the no-order-unit direction
in part~(b) are Theorems~\ref{hol:td:main:boundary}
and~\ref{hol:cf:main:entire}\,(a). The proof keeps the whole initial
polynomial at the infinitely many scales where it is not a monomial, reduces
the equation there to one fixed homogeneous polynomial over $k$, and shows
with logarithmic Euler derivatives that both endpoints of every such initial
polynomial lie in a finite set, while a nonpolynomial entire series has such
scales with arbitrarily large endpoints (Section~\ref{hol:ot:sub:endpoint}).
At order three an extra integer, the gap next to an endpoint, survives
(Theorem~\ref{hol:ot:thm:gaps}); for the equation of $\mathcal T_p$ it is
always one, the initial polynomials of every nonpolynomial entire solution
are eventually $bX^m(1+aX)$, and the coefficient values of such a solution
are eventually strictly convex (Theorems~\ref{hol:ot:thm:classification}
and~\ref{hol:ot:thm:convexity}). Consequences include rational two-state
systems and quotients with a polynomial denominator
(Corollaries~\ref{hol:ot:cor:systems} and~\ref{hol:ot:cor:meromorphic}).

The tenth principal theorem comes from source~18 and is proved in
Sections~\ref{hol:ud:sec:setting}--\ref{hol:ud:sec:mixed}, over
$\KK=k((t^\Gamma))$ for an arbitrary field $k$ of characteristic zero
(Convention~\ref{hol:ud:conv:field}). Its operators are the finite sums
\[
 L=\sum_{\ell=1}^m\sum_{r=0}^Rp_{\ell,r}(z)\,D_z^r\sigma_{\lambda_\ell},\qquad
 p_{\ell,r}\in\KK[z]\text{ not all zero},
\]
for multipliers $\lambda_1,\ldots,\lambda_m\in\KK^\times$.

\begin{maintheorem}[Mixed equations: a relative-scale criterion;
source~18]\label{hol:ud:main:mixed}
Let $\Gamma$ be as in Convention~\ref{hol:conv:group}, $k$ any field of
characteristic zero, and $\lambda_1,\ldots,\lambda_m\in\KK^\times$ with no
quotient $\lambda_i/\lambda_j$, $i\ne j$, a root of unity.
\begin{enumerate}[label=\textup{(\alph*)}]
\item If no $\abs{v(\lambda_i/\lambda_j)}$ is an order unit of $\Gamma$, then
for every such $L$ and every $\bar e\in\N$ there are $N\in\N$ and
$B\in\Gamma_{\ge0}$ such that every formal solution of $Lf=\mathsf h$, with
$\mathsf h\in\KK[z]$ of degree at most $\bar e$, that is strongly evaluable at
one $x\ne0$ with $v(x)<-B$ has degree less than $N$. In particular every
strongly entire solution is a polynomial.
\item Some nonpolynomial $f\in\EK$ satisfies $Lf=0$ for some such $L$ if and
only if $\abs{v(\lambda_i/\lambda_j)}$ is an order unit for some $i,j$; then
two of the multipliers and derivatives of order at most one suffice.
\end{enumerate}
In particular, for every $\Gamma$ and every $q\in\KK^\times$ of infinite
multiplicative order with $v(q)=0$, every strongly entire solution of a
nonzero equation \eqref{hol:eq:mixed} with coefficients in $\KK[z]$ is a
polynomial.
\end{maintheorem}

The last assertion answers Question~\ref{hol:q:mixedunit} positively, for
every $\Gamma$ (Corollary~\ref{hol:ud:cor:unitrigidity}). Part~(a) contains
Theorem~\ref{hol:main:ode}, the negative direction of
Theorem~\ref{hol:main:q} and Theorem~\ref{hol:thm:mixed}, now over every
characteristic-zero $k$ (Remark~\ref{hol:ud:rem:contains}); part~(b) answers
the existence half of Question~\ref{hol:q:severaldilations} for a prescribed
finite set of multipliers, but not for prescribed operators or systems. The
new ingredient is a valuation form of the Skolem--Mahler--Lech theorem: an
exponential polynomial $\sum_jR_j(n)\lambda_j^n$ with arbitrary Hahn
coefficients has eventually affine valuations on finitely many arithmetic
progressions, and eventually periodic ones at unit valuation
(Theorems~\ref{hol:ud:thm:unitvalues} and~\ref{hol:ud:thm:affine}); for
$v(P(n,q^n))$ this is the estimate that Question~\ref{hol:q:mixedunit} named
(Corollary~\ref{hol:ud:cor:bivariate}).

The eleventh principal theorem comes from source~19 and is proved in
Sections~\ref{hol:nr:sec:setting}--\ref{hol:nr:sec:boundary}, over
$\KK=k((t^\Gamma))$ for an arbitrary field $k$ of characteristic zero
(Convention~\ref{hol:nr:conv:field}); source~19 also re-derives
Theorem~\ref{hol:ot:main:second}\,(a) independently, and that proof is printed
once, in Section~\ref{hol:ot:sec:setting}.

\begin{maintheorem}[Order three through cubic jet degree; source~19]
\label{hol:nr:main:third}
Let $\Gamma$ be as in Convention~\ref{hol:conv:group} and $k$ any field of
characteristic zero. Every $f\in\EK$ satisfying
$P(z,f,D_zf,D_z^2f,D_z^3f)=0$ for a nonzero $P\in\KK(z)[Y_0,Y_1,Y_2,Y_3]$ of
total degree at most three in $Y_0,\ldots,Y_3$ is a polynomial. More
generally, the same holds for every equation of order at most three whose
exterior form, or each of whose irreducible factors over $k$, either takes a
nonzero value on the monomial jet curve $(1,T,T^2,T^3)$ or has a nonzero first
polar along it.
\end{maintheorem}

The degree in $z$ is unrestricted. At order three a nonmonomial initial
polynomial at a far break must satisfy the two endpoint conditions of
Theorem~\ref{hol:ot:thm:gaps}; when the first polar is nonzero these say that
the exponent next to each end lies on a fixed side of it, which fails at one of
the two ends far out (Proposition~\ref{hol:nr:prop:firstpolar}), and a secant
argument shows that a nonzero cubic form vanishing on the twisted cubic has a
nonzero first polar (Lemma~\ref{hol:nr:lem:secant}). Beyond degree three the
test can fail: the quartic $\mathsf H_3$ of order three and the quadratic
$\mathsf Q_4$ of order four have exactly the monomials and the doubling
binomials $az^n+bz^{2n}$ as formal solutions
(Theorems~\ref{hol:nr:thm:H3} and~\ref{hol:nr:thm:Q4}), and a sparse strongly
entire series satisfies both exterior equations at every scale without
satisfying either equation (Example~\ref{hol:nr:ex:sparse}). With an order
unit, $\mathcal T_p$ therefore satisfies no equation of order three and jet
degree below four (Corollary~\ref{hol:nr:cor:theta}). Consequences include
external surcomplex and Gaussian-omnific constants and algebraic two-state
systems (Corollary~\ref{hol:nr:cor:external},
Theorem~\ref{hol:nr:thm:systems}).

The twelfth principal theorem comes from source~16 and is proved in
Sections~\ref{hol:th:sec:setting}--\ref{hol:th:sec:scope}, over
$\KK=k((t^\R))$ with $k=\R$ or $\C$ (Convention~\ref{hol:th:conv:field}). For a
real $\mu>0$, $\mathcal T_{t^\mu}$ is \eqref{hol:td:eq:series} with $p=t^\mu$, and
$\operatorname{dord}(f)$ is the least $r$ with $f,D_zf,\ldots,D_z^rf$
algebraically dependent over $\KK(z)$.

\begin{maintheorem}[Independent theta scales and unbounded differential order;
source~16]\label{hol:th:main:hierarchy}
Let $\Gamma=\R$ and $k=\R$ or $\C$.
\begin{enumerate}[label=\textup{(\alph*)}]
\item If $\mu_1,\ldots,\mu_N>0$ and $\mu_1^{-1},\ldots,\mu_N^{-1}$ are linearly
independent over $\Q$, then $\mathcal T_{t^{\mu_1}},\ldots,\mathcal T_{t^{\mu_N}}$
are algebraically independent over $\KK(z)$; there is such a family of
cardinality $2^{\aleph_0}$, each member of minimal order three.
\item For every $N\ge1$ and $\mu_i=1/\sqrt{\mathrm p_i}$ with distinct primes
$\mathrm p_i$, some $g=\sum_{i=1}^N(\sum_{j<N}\mathsf c_{ij}z^j)\mathcal
T_{t^{\mu_i}}$ with $\mathsf c_{ij}\in\{0,1,\ldots,N\}$, one of at most
$(N+1)^{N^2}$ candidates, is strongly entire, differentially algebraic over
$\KK(z)$, and satisfies $N\le\operatorname{dord}(g)\le3N$.
\item The differential subfield $\mathcal D_{\mathrm{DA}}$ of $\FE$ generated
over $\KK(z)$ by all differentially algebraic strongly entire series has
$\trdeg_{\KK(z)}\mathcal D_{\mathrm{DA}}=2^{\aleph_0}$ and
$\dtrdeg_{\KK(z)}\mathcal D_{\mathrm{DA}}=0$; it is not finitely generated as a
differential field.
\end{enumerate}
\end{maintheorem}

Part~(a) rests on the exact Gauss profile
$\gamma_\delta(\mathcal T_{t^\mu})=-\delta^2/(4\mu)+O(1)$
(Theorem~\ref{hol:th:thm:profile}): distinct monomials of a relation have
distinct quadratic growth. Part~(b) combines it with a lemma of pure
differential algebra, bounded polynomial jet compression
(Theorem~\ref{hol:th:thm:compression}). With Theorem~\ref{hol:ot:main:second}
the minimal orders of nonpolynomial differentially algebraic strongly entire
series over $\KK$ form an unbounded set of integers at least three containing
three (Remark~\ref{hol:th:rem:ordertwo}). Source~16 also gives a valuation
proof of the minimal order three of $\mathcal T_p$
(Section~\ref{hol:th:sub:third}), a third route beside those of sources~13
and~14.

The thirteenth principal theorem comes from source~17 and is proved in
Sections~\ref{hol:fh:sec:setting}--\ref{hol:fh:sec:realizations}
(Convention~\ref{hol:fh:conv:field}). Let $(\mathsf b_n)_{n\ge2}$ be positive
integers with $\mathsf b_n/\mathsf b_{n-1}\to\infty$, for instance $n!$, and
for $A\subseteq\N_{\ge2}$ put $\mathcal G_A(p,z)=\sum_{n\in A}p^{\mathsf b_n}z^n$.

\begin{maintheorem}[Differential freedom at a single scale; source~17]
\label{hol:fh:main:single}
Let $k$ be a field of characteristic zero.
\begin{enumerate}[label=\textup{(\alph*)}]
\item For every finite family $A_1,\ldots,A_m$, after subtracting explicit
polynomials that account for the finitely many indices with patterns
$(\mathbf 1_{A_i}(n))_i$ occurring finitely often, the differential relation
ideal of $(\mathcal G_{A_i})_i$ over $k(p,z)$, for $\vartheta=zD_z$ and
$\vartheta_p=p\,\mathrm d/\mathrm dp$, is generated by the rational linear
relations among the patterns that occur infinitely often; if every $A_i$ has
infinitely many indices belonging to it alone, all mixed jets are
algebraically independent.
\item If $\Gamma$ has an order unit $\mu$, $\KK=k((t^\Gamma))$ and $p=t^\mu$,
every $\mathcal G_A$ with $A$ infinite is a nonpolynomial element of $\EK$ that
satisfies no nonzero algebraic differential equation over $\KK(z)$, and the
$D_z$-relations of a finite family over $\KK(z)$ are exactly those of~(a).
\item For $\KK=\R((t^\Q))$ or $\C((t^\Q))$, $\dtrdeg_{\KK(z)}\FE=2^{\aleph_0}$,
with a witnessing family whose Taylor coefficients all lie in $\Q(t)$.
\end{enumerate}
\end{maintheorem}

The coefficient values $\mathsf b_n\mu$ span $\Q\mu$, so neither
Theorem~\ref{hol:cf:main:entire} (which needs no order unit) nor
Theorem~\ref{hol:cf:main:rank} (which needs infinite coefficient-value rank)
applies: the freedom comes from the support, through a height-separation band
and polarization of the top-degree part of a relation
(Section~\ref{hol:fh:sub:finite}). With Theorems~\ref{hol:td:main:boundary}
and~\ref{hol:th:main:hierarchy}, an order-unit field thus carries strongly
entire series of order three, of arbitrarily large order (for $\Gamma=\R$), and
of no finite order. Source~17 also proves the independence of the
Berarducci--Mantova jets of the numbers $\sum_{n\in A}\omega^{-\mathsf b_n}$ over
a bounded-support field, the exact common strong domain
$\{x:v(x)\ge-m\text{ for some }m\}$ of these series in $\No[i]$, and
independent omnific codes with factorial floor profiles
(Section~\ref{hol:fh:sec:realizations}). Differential transcendence of one such
lacunary series is classical (Gr\"onwall, via Ostrowski~\cite{ostrowski}).

The fourteenth principal theorem comes from source~15 and is proved in
Sections~\ref{hol:pc:sec:setting}--\ref{hol:pc:sec:arithmetic}
(Convention~\ref{hol:pc:conv:field}). For monomials $Y^{\mathbf a}$ in
variables attached to arguments of degrees $e_1,\ldots,e_{\bar m}$, the
\emph{composition weight} is $\operatorname{cw}(\mathbf a)=\sum_\nu a_\nu e_\nu$.

\begin{maintheorem}[Polynomial-composition rigidity; source~15]
\label{hol:pc:main:composition}
Let $k$ be a field of characteristic zero and $\KK=k((t^\Gamma))$. Let
$P\in\KK[z]$ have degree $d\ge2$, let $P_1,\ldots,P_{\bar m}\in\KK[z]$ be
nonconstant of degrees $e_\nu<d$, let $L_1,\ldots,L_{\bar s}$ \textup{(}$\bar
s\ge1$\textup{)} and $L'_1,\ldots,L'_{\bar m}$ be nonzero linear differential
operators with coefficients in $\KK[z]$, let $0\ne c_*\in\KK[z]$, and let
$\mathsf G\in\KK[z][Y_1,\ldots,Y_{\bar m}]$ have only monomials of composition
weight $\operatorname{cw}(\mathbf a)<\bar sd$. Then every $f\in\EK$ with
\[
 c_*(z)\prod_{\ell=1}^{\bar s}L_\ell\bigl(f(P(z))\bigr)
 =\mathsf G\bigl(z,L'_1(f(P_1(z))),\ldots,L'_{\bar m}(f(P_{\bar m}(z)))\bigr)
\]
is a polynomial, of degree at most an explicit integer $B_{\mathrm{cmp}}$
computed from the degrees, the operator shifts and the integer roots of the
exact indicial polynomials $I_{L_\ell}(dT)$ \textup{(}Theorem~\ref{hol:pc:thm:bound}\textup{)}.
In particular, for every $\Gamma$, every strongly entire solution of a linear
equation $\sum_{j=0}^mL_j(f(P_j(z)))=g(z)$ with $L_m\ne0$ and a unique
argument $P_m$ of largest degree $d_m\ge2$ is a polynomial.
\end{maintheorem}

The lower arguments may be dilations $\lambda z$. So, unlike
Theorems~\ref{hol:main:q} and~\ref{hol:ud:main:mixed}, no order unit and no
condition on scales enters: a strictly dominant argument of degree at least
two defeats the partial theta mechanism (Remark~\ref{hol:pc:rem:dichotomy}).
The proof reads how a polynomial argument rescales the Gauss data of
Section~\ref{hol:nl:sub:scaled} and compares scales by an integer convexity
inequality. Source~15 also proves a matrix version, describes the solution
modules of linear equations with integer data and omnific coefficients, and
shows that every affine hypersurface over $\Q$ is the entire-solution locus of
one equation of this class, so omnific solvability is undecidable inside it
(Theorem~\ref{hol:pc:thm:universal}, Corollary~\ref{hol:pc:cor:undecidable}).

The fifteenth principal theorem comes from source~20 and is proved in
Sections~\ref{hol:rc:sec:setting}--\ref{hol:rc:sec:surreal}
(Convention~\ref{hol:rc:conv:field}). A \emph{coefficient-algebraic condition}
on a vector of Hahn series is a polynomial equation over $k$ in finitely many
of their coefficients (Definition~\ref{hol:rc:def:conditions}).

\begin{maintheorem}[Coefficient observables of recurrences; source~20]
\label{hol:rc:main:observables}
Let $k$ be a field of characteristic zero and $\KK=k((t^\Gamma))$.
\begin{enumerate}[label=\textup{(\alph*)}]
\item Let $\mathbf y_n$ be a finite vector of exponential polynomials
$\sum_jR_j(n)\lambda_j^n$, $R_j\in\KK[U]$, whose bases $\lambda_j$ have
valuation zero and residues $\zeta_1,\ldots,\zeta_m$. Every coefficient
sequence $n\mapsto[t^\gamma]$ of every coordinate lies in the one finitely
generated $k$-algebra $k[n,\zeta_1^n,\ldots,\zeta_m^n]$ of sequences.
\item For every set-sized family of coefficient-algebraic conditions, a finite
subfamily is equivalent to the whole family at every point of the orbit
$(\mathbf y_n)_{n\ge0}$, and the set of $n$ at which $\mathbf y_n$ satisfies the
family is a finite union of full residue classes modulo
$\mathsf M_{\mathrm{tor}}$, the exponent of the torsion subgroup of
$\langle\zeta_1,\ldots,\zeta_m\rangle$, and a finite set; one
$\mathsf M_{\mathrm{tor}}$ serves for every family.
\item Without units this fails completely: for every real sequence
$(\mathbf b_m)_{m\ge0}$ there is a purely infinite omnific integer
$\mathfrak w$ with $[\omega^\omega](\omega^n\mathfrak w)=\mathbf b_n$ for every
ordinary $n$, so the first-order omnific recurrence $y_{n+1}=\omega y_n$
realizes every subset of $\N$ as the nonzero set of one fixed Conway
coefficient.
\end{enumerate}
\end{maintheorem}

Part~(a) rests on the support envelope of Lemma~\ref{hol:ud:lem:envelope};
part~(b) adds Hilbert's basis theorem and a Skolem--Mahler--Lech period valid
for all monomials in the residues (Lemma~\ref{hol:rc:lem:torperiod}), for which
the period of source~18's valuation theorem is too small
(Remark~\ref{hol:rc:rem:periods}). The element in part~(c) is a stream already
in the omnific-notations report (Remark~\ref{hol:rc:rem:stream}); new is its
reading as a recurrence observable, which excludes a coefficientwise
Skolem--Mahler--Lech theorem for omnific recurrences
(Corollary~\ref{hol:rc:cor:nosml}), while the leading exponent of every
recurrence over $\No[i]$, $\Oz$ or $\Oz[i]$ stays eventually affine on finitely
many progressions (Theorem~\ref{hol:rc:thm:surreal}). Source~20 also
re-derives, independently of source~18, the valuation theorem for Hahn
exponential polynomials and the unit-dilation case of
Theorem~\ref{hol:ud:main:mixed}, printed once; its exact finite-group
criterion is a special case of that theorem
(Section~\ref{hol:rc:sub:mixed}).

For the linear theorems, the hypotheses separate two obstructions.
Differential equations have bounded valuation costs after passing to their
coefficient recurrence. Dilation
equations can have costs increasing linearly with the index, but only along
the single scale $v(q)$; a scale that is not an order unit cannot reach all of
$\Gamma$, even if its multiples are unbounded within their own Archimedean
component. Strong entireness detects that failure.

Beyond the first three, the report proves: a mixed differential--dilation rigidity
theorem at every nonzero noncofinal dilation scale
(Theorem~\ref{hol:thm:mixed}, extended by Theorem~\ref{hol:ud:main:mixed} to
unit dilations and to finite sets of multipliers); a recurrence exclusion bound that is
\emph{attained}, boundary included, by an explicit formal exponential
(Corollary~\ref{hol:cor:periodic} and
Proposition~\ref{hol:prop:expdomain}); a coarsened exclusion bound valid at
every nonzero noncofinal dilation valuation (Theorem~\ref{hol:thm:coarse}); and the exact
strong evaluation domain of the partial theta series for an arbitrary Hahn
parameter, including arbitrary tails (Theorem~\ref{hol:thm:thetadomain}).

\subsection{Where divisibility of \texorpdfstring{$\Gamma$}{Gamma} is used}
\label{hol:sub:divisibility}

The core arguments --- the escape mechanism, the orbit valuation lemmas, the
two conversions, the theta domain theorem, and all of
Theorems~\ref{hol:main:ode}--\ref{hol:main:multi} --- use only
Convention~\ref{hol:conv:group}. Divisibility of $\Gamma$ enters in exactly
three threshold constructions, with the hypothesis stated where it is used.
\begin{enumerate}[label=\textup{(\arabic*)}]
\item The refined recurrence threshold $B_{\mathrm{rec}}$ of
\eqref{hol:eq:radiusbound}, whose defining maximum divides a valuation
difference by the jump length $j$. Corollary~\ref{hol:cor:periodic} states it
under divisibility; Corollary~\ref{hol:cor:boundedcost} is the
divisibility-free predecessor, with a coarser threshold and the same
qualitative conclusion.
\item The refined coarsened threshold $\overline B_{\mathrm{rec}}$ of
\eqref{hol:eq:coarsebound}, for the same reason, in the quotient group.
Theorem~\ref{hol:thm:coarse} states the divisibility-free form and the
divisible refinement separately.
\item The worked sparse threshold in Example~\ref{hol:ex:sparse}, an instance
of (1).
\end{enumerate}
The torsion eigenspace equivalence in Proposition~\ref{hol:prop:torsion}
needs no divisibility: inward stability of strong evaluation replaces the
root-extraction argument of the source manuscripts.

The nonlinear part,
Sections~\ref{hol:nl:sec:setting}--\ref{hol:nl:sec:consequences}, uses no
divisibility of $\Gamma$ anywhere: the degree cutoff
\eqref{hol:nl:eq:kappa} divides ordinary integers in $\Q$, and its exclusion bounds are
chosen from finitely many strict inequalities with positive integer
multipliers.

The coefficient-field part,
Sections~\ref{hol:cf:sec:setting}--\ref{hol:cf:sec:families}, uses no
divisibility either. It works in $\Gamma\otimes_\Z\Q$ only to measure rational
rank, and the principal convex bound of Lemma~\ref{hol:cf:lem:principal} uses
integer multiples only.

The theta-descent part uses no divisibility. The order-three-threshold part,
Sections~\ref{hol:ot:sec:setting}--\ref{hol:ot:sec:theta}, states its theorems
for $\Gamma$ as in Convention~\ref{hol:conv:group}; to find scales at which the
initial polynomial is not a monomial it divides in the divisible hull
$\Gamma\otimes_\Z\Q$, after extending the field by
Lemma~\ref{hol:ot:lem:extension}, and transfers the polynomiality and
coefficient conclusions back to $\Gamma$, in which the hull is cofinal
(Lemmas~\ref{hol:ot:lem:extension} and~\ref{hol:ot:lem:breaks}).

The unit-dilation part, Sections~\ref{hol:ud:sec:setting}--%
\ref{hol:ud:sec:mixed}, uses no divisibility: its cost bounds are formed from
finitely many elements and a noncofinality witness, and its jump-sensitive
comparison multiplies by the jump instead of dividing by it. The
Newton-rigidity part, Sections~\ref{hol:nr:sec:setting}--%
\ref{hol:nr:sec:boundary}, passes to $\Gamma\otimes_\Z\Q$ exactly as the
order-three-threshold part does, through Lemma~\ref{hol:ot:lem:extension}, and
transfers its conclusions back.

The theta-hierarchy part works in $\Gamma=\R$, which is divisible, and states
it; the single-scale part uses no divisibility, and neither does the
polynomial-composition part, whose degree bounds are floors of ratios of
ordinary integers and whose convexity inequality multiplies by positive
integers. The recurrence part uses no divisibility either: its periods are
exponents of finite groups of residues, and its cost bounds are those of the
unit-dilation part.

No theorem in this report states a conclusion under
Convention~\ref{hol:conv:group} while resting on a proof step that divides by
an integer in $\Gamma$; the divisions of the order-three-threshold and
Newton-rigidity parts take place in $\Gamma\otimes_\Z\Q$. The division-free predecessors in (1) and (2) are
weaker in the constant only, not in the conclusion.

\subsection{Position relative to the collection and the literature}
\label{hol:sub:position}

\paragraph{Provenance.}
This report is built from thirteen manuscripts, numbered as they arrived in
the collection. Sources~08 and~09 were written independently on 22 September 2026
and both prove the differential rigidity theorem and the order-unit dilation
criterion (Theorems~\ref{hol:main:ode} and~\ref{hol:main:q}); source~08 is
the base, because it assumes only $\Gamma\ne0$ while source~09 assumed
divisibility, and every proof taken from~09 was re-read for a hidden division
(Section~\ref{hol:sub:divisibility}). Together they supply
Sections~\ref{hol:sec:intro}--\ref{hol:sec:certificates}. Source~10, delivered
later the same day, supplies the nonlinear part,
Sections~\ref{hol:nl:sec:setting}--\ref{hol:nl:sec:consequences} and
Theorem~\ref{hol:nl:main:first}. The repository snapshot consulted by
sources~08 and~09 is
\begin{center}\small\texttt{a3124af79f66b8b9c196d76b4cbc5ac3938907c4},\end{center}
and the one consulted by source~10 is
\begin{center}\small\texttt{048b72cf7cbfc8ab246e4f73788c10460cb3f6e0}.\end{center}
At the second pin this report consisted of sources~08 and~09 alone and its text
was identical to the text into which source~10 was merged, so source~10's
references to ``Question~15.1'' and ``Theorems~A--C'' of this report were
accurate there; that question is now Question~\ref{hol:q:nonlinear}, re-scoped
in Section~\ref{hol:sec:status}. The choices the merge made (renamed symbols,
results printed once, the placement of the nonlinear part after the linear
theory) are recorded in Section~\ref{hol:nl:sub:conventions}.

Sources~11 and~12, delivered together in a later batch and written
independently, are titled \emph{Coarsening Descent and Nonlinear Differential
Rigidity for Entire Hahn--Surcomplex Functions} and \emph{Coefficient-Field
Rigidity at Unbounded Surreal Scales}. Both answer the no-order-unit case of
Question~\ref{hol:q:nonlinear} by the same argument, and together they form
the coefficient-field part,
Sections~\ref{hol:cf:sec:setting}--\ref{hol:cf:sec:families}, with
Theorems~\ref{hol:cf:main:entire}--\ref{hol:cf:main:sharp}. Source~12 is the
base of that part: it adds dilations, answers the no-order-unit case of
Question~\ref{hol:q:mixedunit}, and gives the order-unit dichotomy. Source~11
adds coarsening descent, rationality in the fraction field, the rational-system
and Painlev\'e corollaries, the joint-independence criterion and the continuum
family. Both consulted the snapshot
\begin{center}\small\texttt{465a54b479a1ee842cbf7db1689a7d2f6bfe25e1},\end{center}
at which this report was exactly
Sections~\ref{hol:sec:intro}--\ref{hol:sec:status} and the conclusion, so
their references to Questions~19.1 and~19.4 and to the README were accurate.
The coefficient-field part is placed after Section~\ref{hol:sec:status} so
that the numbering of Sections~1--19, of Questions~19.1--19.6 and of
Theorems~A--D, which the shipped audit files cite, is unchanged; the new
questions are Questions~\ref{hol:cf:q:descent}--\ref{hol:cf:q:torsion}. The
merge's choices for this part are recorded in
Section~\ref{hol:cf:sub:conventions}.

Source~13, \emph{Theta Descent and the Exact Boundary of Differential
Rigidity}, delivered in a later batch, answers the order-unit case of
Question~\ref{hol:q:nonlinear} by an explicit series of minimal differential
order three, and so shows that the order unit is the exact boundary. It forms
the theta-descent part,
Sections~\ref{hol:td:sec:setting}--\ref{hol:td:sec:boundary}, with
Theorem~\ref{hol:td:main:boundary}, and adds clause~(iv) to
Theorem~\ref{hol:cf:main:sharp}. It consulted the snapshot
\begin{center}\small\texttt{4cdeaec43ff1692ed2ad9f5bdf761a41872975a5},\end{center}
at which this report was exactly the five-source text, so its references to
Question~19.1 and Theorem~E were accurate. The new part is placed after
Section~\ref{hol:cf:sec:families}, so that the numbering of Sections~1--22, of
Questions~19.1--19.9 and of Theorems~A--G is unchanged; its one new question is
Question~\ref{hol:td:q:domains}. The merge's choices for this part, including
two steps it writes out, are recorded in
Section~\ref{hol:td:sub:conventions}.

Source~14, \emph{The Order-Three Threshold for Entire Hahn--Surcomplex
Functions}, delivered in a later batch, answers the order-two question that
source~13 left open: no nonpolynomial strongly entire series satisfies an
algebraic differential equation of order at most two, for every value group.
It forms the order-three-threshold part,
Sections~\ref{hol:ot:sec:setting}--\ref{hol:ot:sec:theta}, with
Theorem~\ref{hol:ot:main:second}. It consulted the snapshot
\begin{center}\small\texttt{b895e8672990e8b5a56f97dd7246f9dd5f86c771},\end{center}
at which this report's text was still the five-source text and source~13 was
known to it only as a manuscript; its description of this report was accurate
there. The new part is placed after
Section~\ref{hol:td:sec:boundary}, so that the numbering of Sections~1--25, of
Questions~19.1--19.10 and of Theorems~A--H is unchanged; the conclusion moves
to the end. Its new questions are
Questions~\ref{hol:ot:q:lift}--\ref{hol:ot:q:omnific}. The merge's choices for
this part are recorded in Section~\ref{hol:ot:sub:conventions}.

Sources~18 and~19 arrived together in a later batch (batch~31, manuscripts~05
and~08), written independently of each other. Source~18, \emph{Unit-Dilation
Rigidity at Surreal Scales}, answers Question~\ref{hol:q:mixedunit}
positively for every value group and proves the relative-scale criterion of
Theorem~\ref{hol:ud:main:mixed}; it forms the unit-dilation part,
Sections~\ref{hol:ud:sec:setting}--\ref{hol:ud:sec:mixed}. It consulted the
snapshot
\begin{center}\small\texttt{bcac55ae6fd3f2b354e568ba1d2e94d496a2cc59}.\end{center}
Source~19, \emph{Low-Order Nonlinear Rigidity at Surreal Scales}, re-derives
Theorem~\ref{hol:ot:main:second}\,(a) independently, printed once, and proves
Theorem~\ref{hol:nr:main:third}; it forms the Newton-rigidity part,
Sections~\ref{hol:nr:sec:setting}--\ref{hol:nr:sec:boundary}. It consulted the
snapshot
\begin{center}\small\texttt{3d40856953f4bb0f5d069e45a3b4bab6eda33040}.\end{center}
At both snapshots this report was still the five-source text, so their
descriptions of it were accurate there; what has since become stale is
recorded in Sections~\ref{hol:ud:sub:conventions}
and~\ref{hol:nr:sub:conventions}. The two parts are placed after
Section~\ref{hol:ot:sec:theta}, so that the numbering of Sections~1--27, of
Questions~19.1--19.18 and of Theorems~A--I is unchanged; the conclusion moves
to the end. Their principal theorems are Theorems~\ref{hol:ud:main:mixed}
and~\ref{hol:nr:main:third}, and their new questions are
Questions~\ref{hol:ud:q:prescribed}--\ref{hol:ud:q:formal}
and~\ref{hol:nr:q:layers}--\ref{hol:nr:q:formal}.

Sources~16 and~17 (batch~31, manuscripts~02 and~03), written independently,
were merged next. Source~16, \emph{Independent Theta Families and Unbounded
Differential Order in Hahn--Surcomplex Analysis}, forms the theta-hierarchy
part, Sections~\ref{hol:th:sec:setting}--\ref{hol:th:sec:scope}, with
Theorem~\ref{hol:th:main:hierarchy}; it records the snapshot
\begin{center}\small\texttt{0fffc26c670c15c7db451ed899f98d4384f9df1e},\end{center}
at which this report's text was the five-source one and source~13's audits
were placed. Source~17, \emph{Differential Freedom at a Single Surreal Scale},
forms the single-scale part,
Sections~\ref{hol:fh:sec:setting}--\ref{hol:fh:sec:realizations}, with
Theorem~\ref{hol:fh:main:single}; it consulted \texttt{bcac55a}, the pin of
source~18. What has since become stale in them is recorded in
Sections~\ref{hol:th:sub:conventions} and~\ref{hol:fh:sub:conventions}. The two
parts follow Section~\ref{hol:nr:sec:boundary}, so that the numbering of
Sections~1--31, of Questions~19.1--19.35 and of Theorems~A--K is unchanged; the
conclusion moves to the end. Their new questions are
Questions~\ref{hol:th:q:spectrum}--\ref{hol:th:q:omnific}
and~\ref{hol:fh:q:heights}--\ref{hol:fh:q:arithmetic}.

Source~15 (batch~31, manuscript~01), \emph{Polynomial-Composition Rigidity at
Surreal Scales}, was merged last of its batch. It forms the
polynomial-composition part,
Sections~\ref{hol:pc:sec:setting}--\ref{hol:pc:sec:arithmetic}, with
Theorem~\ref{hol:pc:main:composition}; it consulted \texttt{bcac55a}, the pin
of source~18, at which this report was the five-source text. What has since
become stale in it, and one statement that was already incorrect there, are
recorded in Section~\ref{hol:pc:sub:conventions}. The part follows
Section~\ref{hol:fh:sec:realizations}, so that the numbering of
Sections~1--35, of Questions~19.1--19.54 and of Theorems~A--M is unchanged;
the conclusion moves to the end. Its new questions are
Questions~\ref{hol:pc:q:critical}--\ref{hol:pc:q:formal}.

Source~20 (batch~33, manuscript~01), \emph{Valuation Rigidity and
Coefficient Universality in Surreal Recurrences}, forms the recurrence part,
Sections~\ref{hol:rc:sec:setting}--\ref{hol:rc:sec:surreal}, with
Theorem~\ref{hol:rc:main:observables}. It consulted the snapshot
\begin{center}\small\texttt{efc5446229dbf61a97bdd5edae91dba20af9d131},\end{center}
at which this report was the seven-source text (sources~08--14) and
sources~15--19 were placed but not yet merged; its statement that the mixed
unit-valuation case was left open was accurate there and is stale now, and its
re-derivations of source~18's results are printed once
(Section~\ref{hol:rc:sub:conventions}). The part follows
Section~\ref{hol:pc:sec:arithmetic}, so that the numbering of
Sections~1--37, of Questions~19.1--19.62 and of Theorems~A--N is unchanged;
the conclusion moves to the end. Its new questions are
Questions~\ref{hol:rc:q:nonunit}--\ref{hol:rc:q:substitutions}.

\paragraph{The companion entire-function report.}
The collection already contains a merged report on cofinality, factorization
and scalar extension for entire Hahn functions at arbitrary
rank~\cite{repoentire}. It fixes the same class $\E$ over the same kind of
workspace, classifies the ring structure of $\E$, and computes the exact
strong evaluation domain after an enlargement of the value group. None of its
ring-theoretic conclusions is used to prove anything here: the present proofs
start from Hahn support arithmetic. Two points of contact matter.

First, we do not reprove the order-unit coarsening machinery. That report's
Lemma~9.1, labelled \texttt{ent:lem:coarsening}, states that if $h>0$ is an
order unit of $\Gamma$ there is an order-preserving homomorphism
$\pi:\Gamma\to\R$ with $\pi(h)=1$ whose kernel is the convex subgroup
$H=\{\gamma:n\abs\gamma<h\text{ for all }n\ge1\}$, so that
$\abs a_\pi=\exp(-\pi(v(a)))$ is a nontrivial real-valued non-Archimedean
absolute value inducing the intrinsic valuation topology. That is exactly the
sense in which sources~08 and~09 of this report use the phrase
``coarsening'', and we cite it rather than repeat it; see
Remark~\ref{hol:rem:coarsening}.

Second, and this is easy to get wrong:
\emph{the two order-unit dichotomies are not the same theorem} --- the
companion report's order unit decides Hermite interpolation and the B\'ezout
property of the ring $\E$ (its Theorem~9.3), whereas the order unit here
decides which annihilating operators a nonpolynomial member of $\E$ can
satisfy. Neither dichotomy implies the other, and the coincidence of the
invariant is not a coincidence of statement. In particular the existence of
nonpolynomial entire functions at countable cofinality, and the fact that
$\Gamma_\infty=\bigoplus_{j\ge0}\Q\omega^j$ has countable cofinality and no
order unit, are background from that report~\cite{repoentire}, not targets of
this one.

The collection's differential-equations report uses a derivation on surreal
\emph{scalars}~\cite{repodifferential}; here $D_z$ kills every scalar and
differentiates only the formal variable, and no theorem about that scalar
derivation is restated as ordinary D-finite rigidity. The rank-one analytic
report~\cite{reporankone} already contains the quadratic-exponent example
$\sum t^{n^2}z^n$; our partial theta witness is explicitly classical and is
not counted as a new function.

\paragraph{External literature.}
Stanley's D-finite/P-recursive correspondence is the standard algebraic
starting point~\cite{stanley}. We make the coefficient conversion explicit
because the valuations of the resulting recurrence, rather than ordinary
coefficient growth, drive the proof. Ramis treats growth of ordinary complex
entire solutions of linear $q$-difference equations and also studies
simultaneous differential and difference constraints~\cite{ramis}; the
arbitrary-rank strong summability condition here is a different global
condition, and none of its order-unit conclusions is inferred from an ordinary
complex growth estimate. Garoufalidis proves eventual quadratic
quasi-polynomial degree behaviour for $q$-holonomic sequences in a
rational-function and Puiseux setting~\cite{garoufalidis}; that is an
important antecedent to the idea that dilation equations constrain coefficient
scales, we do not claim to originate it, and we do not strengthen its
quadratic quasi-polynomial conclusion. Di Vizio's ultrametric
Maillet--Malgrange theorem treats Gevrey behaviour of formal solutions of
nonlinear ultrametric $q$-difference equations, including the valuation-zero
case with Diophantine conditions~\cite{divizio}; we make no nonlinear
generalization of it. The partial theta series belongs to the classical
partial-theta literature~\cite{andrewswarnaar}, and Higman's finite-word
theorem is a standard external ingredient in the positive-support monoid
lemma~\cite{higman}.

Terminology needs care in two further directions. Hahn meromorphic functions
in the analytic setting of M\"uller and Strohmaier involve a convergence
framework different from the formal strong evaluation used
here~\cite{mueller}. The Berarducci--Mantova derivation differentiates surreal
coefficients themselves and is not $D_z$~\cite{bm}. Work on operation-closed
surreal subfields gives context for set-sized workspaces~\cite{bg}, and the
Conway normal-form conventions used in Section~\ref{hol:sub:embedding} follow
current surreal field work~\cite{kks}; closure under a scalar derivation is
assumed in no theorem below.

The inspections of sources~08 and~09 were targeted: repository documentation,
the report catalogue, the entire-function report, and the primary sources
named above. Neither was a line-by-line audit of every repository manuscript
or an exhaustive bibliography search, and a repository content search for
\texttt{holonomic} returning no matches is not proof of absence. Source~10's
repository and literature comparison, equally targeted, is recorded in
Section~\ref{hol:nl:sub:predecessors}, that of sources~11 and~12 in
Section~\ref{hol:cf:sub:predecessors}, that of source~13 in
Section~\ref{hol:td:sub:predecessors}, that of source~14 in
Section~\ref{hol:ot:sub:predecessors}, those of sources~18 and~19 in
Sections~\ref{hol:ud:sub:predecessors} and~\ref{hol:nr:sub:predecessors}, and
those of sources~16 and~17 in Sections~\ref{hol:th:sub:predecessors}
and~\ref{hol:fh:sub:predecessors}, that of source~15 in
Section~\ref{hol:pc:sub:predecessors}, and that of source~20 in
Section~\ref{hol:rc:sub:predecessors}. Section~\ref{hol:sec:status} and
Sections~\ref{hol:cf:sub:nonclaims}, \ref{hol:td:sub:nonclaims},
\ref{hol:ot:sub:nonclaims}, \ref{hol:ud:sub:nonclaims},
\ref{hol:nr:sub:nonclaims}, \ref{hol:th:sub:nonclaims},
\ref{hol:fh:sub:nonclaims}, \ref{hol:pc:sub:nonclaims}
and~\ref{hol:rc:sub:nonclaims} record exactly what is and is not claimed.

\section{Hahn support and what evaluation means}\label{hol:sec:support}

\subsection{Support calculus}

\begin{lemma}[Hahn--Neumann support calculus and the positive monoid]
\label{hol:lem:support}
Let $A,B$ be well-ordered subsets of an ordered abelian group.
\begin{enumerate}[label=\textup{(\arabic*)}]
\item $A+B$ is well ordered, and each $\gamma$ has only finitely many
representations $\gamma=a+b$ with $(a,b)\in A\times B$.
\item If $S$ is a well-ordered subset of $\Gamma_{>0}$, the additive monoid
\[
 M(S)=\{s_1+\cdots+s_k:k\ge0,\ s_i\in S\}
\]
is well ordered, and each of its elements has only finitely many
representations as a finite word in letters from $S$.
\item Multiplication by a fixed Hahn element preserves strong summability, and
multiplication by a nonzero one preserves and reflects it.
\end{enumerate}
\end{lemma}

\begin{proof}
An infinite sequence in a well-ordered set has an infinite nondecreasing
subsequence. Given a strictly decreasing sequence $a_n+b_n$, pass first to a
subsequence with $a_n$ nondecreasing and then to one with $b_n$ nondecreasing,
a contradiction. If infinitely many distinct pairs sum to a fixed $\gamma$,
the distinct first coordinates contain an increasing sequence, forcing a
decreasing sequence of second coordinates. This is part~(1).

For part~(2), order finite words over $S$ by subsequence embedding with
coordinatewise increase. Higman's theorem says that words over a
well-quasi-ordered alphabet are well quasi ordered~\cite{higman}. An embedding
cannot decrease the sum, and increases it strictly unless the two words are
identical, since inserted letters are positive and a strict coordinate
increase is positive. A descending sequence of sums therefore gives a bad
sequence of words, and so do infinitely many distinct words with one fixed
sum. Both are impossible.

For part~(3), let $A$ be the union of the supports of a strongly summable
family and $B$ the support of the multiplier. The new supports lie in $A+B$.
Part~(1) leaves finitely many admissible pairs at a specified exponent, and
for each first coordinate only finitely many family members contribute. For a
nonzero multiplier apply the same to its inverse.
\end{proof}

In part~(2) words carry their order; a fixed finite multiset has only finitely
many permutations, so the commutative monoid version follows. The argument
does not assume that $S$ is bounded below by an order unit, and remains valid
at arbitrary valuation rank.

\begin{lemma}[The nonincreasing leading-exponent obstruction]
\label{hol:lem:nonincreasing}
A family of nonzero Hahn elements possessing an infinite subsequence whose
leading exponents are \emph{nonincreasing} is not strongly summable.
\end{lemma}

\begin{proof}
If that sequence of leading exponents takes infinitely many distinct values it
has an infinite strictly decreasing subsequence, so the union of the supports
is not well ordered. Otherwise it is eventually constant, and one exponent
lies in the supports of infinitely many members, so local finiteness fails.
\end{proof}

The second alternative is essential and is what makes the exclusion boundary
below a \emph{closed} region. Merely excluding a descending support does not
establish strong summability: the family $(1\cdot z^n)_{n\ge0}$ at $z=1$ has
support union $\{0\}$, yet every summand contributes at exponent zero.

\begin{lemma}[Two sufficient certificates]\label{hol:lem:certificate}
Let $(b_n)_{n\ge0}$ be a family in $\K$.
\begin{enumerate}[label=\textup{(\alph*)}]
\item If $v(b_n)$ tends cofinally to $+\infty$, meaning that for each
$\beta\in\Gamma$ all but finitely many nonzero $b_n$ satisfy $v(b_n)>\beta$,
then $(b_n)$ is strongly summable.
\item If $\supp(b_n)\subseteq h_n+M(S)$ for a well-ordered
$S\subseteq\Gamma_{>0}$, where the set of values of $h_n$ is well ordered and
each value is taken by only finitely many indices, then $(b_n)$ is strongly
summable.
\end{enumerate}
Moreover, grouping a strongly summable family into finite blocks preserves
strong summability and its sum.
\end{lemma}

\begin{proof}
For (a) fix $\beta$. Only finitely many supports meet $(-\infty,\beta]$, and a
finite union of well-ordered subsets of a linear order is well ordered. A
descending sequence in the full union would lie below its first member, hence
inside such a finite union. The same finiteness bounds the multiplicity of an
exponent. For (b) use parts~(1) and~(2) of Lemma~\ref{hol:lem:support}: a
fixed target exponent admits finitely many pairs $(h,m)\in\{h_n\}\times M(S)$
with $h+m$ equal to it, and finitely many indices carry each such $h$; actual
supports can only be smaller than the indicated sets. For the last assertion,
each block support is contained in the union of the old supports and the
blocks contributing at any exponent form a finite set, while coefficientwise
associativity is finite.
\end{proof}

Criterion~(a) and criterion~(b) are genuinely different tools and both are
used below. Criterion~(a) is the shorter route whenever the term valuations
really do escape every bound. Criterion~(b) is what survives at the boundary,
where the term valuations increase strictly but remain trapped below a fixed
element of $\Gamma$; Theorem~\ref{hol:thm:thetadomain} is exactly such a case,
and no leading-valuation argument alone can decide it.

\begin{lemma}[Inward stability of strong evaluation]\label{hol:lem:inward}
Let $f=\sum_na_nz^n\in\K[[z]]$, and let $x,y\in\K^\times$ with
$v(y)\ge v(x)$. If $x\in\Dom(f)$, then $y\in\Dom(f)$.
Thus membership of a nonzero argument in $\Dom(f)$ depends only on its
valuation, and entireness can be tested on monomial arguments.
\end{lemma}

\begin{proof}
Put $b_n=a_nx^n$ and $u=y/x$, so $(b_n)$ is strongly summable and
$v(u)=\eta\ge0$. Write $u=c\,t^\eta(1+w)$, where $c\in\C^\times$
and $w$ is zero or has positive valuation. Let $S$ be $\supp(w)$,
with $\eta$ adjoined when $\eta>0$, and put $M=M(S)$.
This is a well-ordered positive alphabet, so Lemma~\ref{hol:lem:support}
makes $M$ well ordered. Every finite power $u^n$ has support in $M$.

Let $A=\bigcup_n\supp(b_n)$. The supports of
$a_ny^n=b_nu^n$ lie in $A+M$, a well-ordered set.
For a fixed exponent $\gamma$ there are only finitely many pairs
$(\alpha,\mu)\in A\times M$ with $\alpha+\mu=\gamma$.
For each such $\alpha$, strong summability of $(b_n)$ leaves only
finitely many indices $n$. Hence $\gamma$ occurs in only finitely many
supports of $b_nu^n$, proving strong summability. Equal valuations allow
the argument in both directions; compare an arbitrary nonzero $y$ with
$t^{v(y)}$ to obtain the monomial test. Evaluation at zero is always finite.
\end{proof}

The statement fixes the full coefficient series $a_n$. It does not say
that their leading valuations alone determine the domain: higher
coefficient tails can still matter at a boundary scale.

\subsection{Positive valuation is not cofinal decay}

For $\beta>0$ write
\begin{equation}\label{hol:eq:delta}
 \Delta_\beta=\{\gamma\in\Gamma:\abs\gamma\le N\beta\text{ for some }N\in\N\},
\end{equation}
the smallest convex subgroup containing $\beta$. It is all of $\Gamma$ exactly
when $\beta$ is an order unit. If it is proper, there is $\gamma>0$ outside
it, and then
\begin{equation}\label{hol:eq:dominates}
 \gamma>N\beta\qquad\text{for every ordinary }N\in\N,
\end{equation}
since a positive element bounded by some $N\beta$ would lie in $\Delta_\beta$.

\begin{remark}[Sufficiency is not necessity]\label{hol:rem:positive}
Let $0<\epsilon<H$ in $\Gamma$ with $n\epsilon<H$ for every ordinary integer
$n$; for instance $\Gamma=\Q H\oplus\Q\epsilon$ ordered by the highest nonzero
coordinate. The family $(t^{n\epsilon})_{n\ge0}$ is strongly summable, since
its support is an increasing sequence and each exponent occurs once, and its
sum is the Hahn element $(1-t^\epsilon)^{-1}$. Yet its valuations do not tend
cofinally to $+\infty$. This is the most elementary warning against
identifying positive valuation with topological nilpotence, and it is why
criterion (b) of Lemma~\ref{hol:lem:certificate} is needed at all. No appeal
to convergence of partial sums at the $H$ scale is involved.
\end{remark}

\begin{remark}[The real-valued coarsening is imported, not reproved]
\label{hol:rem:coarsening}
When $\Gamma$ does have an order unit $h$, the companion
report~\cite[Lemma~9.1, \texttt{ent:lem:coarsening}]{repoentire} supplies an
order-preserving $\pi:\Gamma\to\R$ with $\pi(h)=1$ and kernel the convex
subgroup $H=\{\gamma:n\abs\gamma<h\text{ for all }n\ge1\}$, hence a nontrivial
real-valued non-Archimedean absolute value $\abs a_\pi=\exp(-\pi(v(a)))$
inducing the intrinsic topology. We cite that lemma and do not reprove it. No
proof in this report uses it: every argument below is carried out in $\Gamma$
itself, which is what makes the arguments valid when no order unit exists.
The coarsening is mentioned twice, as an optional rank-one picture of the
order-unit case (Remark~\ref{hol:rem:thetacoarse}) and for the warning of
Section~\ref{hol:sub:position} that the interpolation/B\'ezout dichotomy of
that report and the annihilating-operator dichotomy of this one share an
invariant but are different theorems.
\end{remark}

\subsection{The entire class and its closure properties}

Polynomials are entire, because all the evaluation families are finite. An
ordinary complex entire Taylor series need not be entire in the sense of
Definition~\ref{hol:def:entire}: at a monomial argument $x=t^\gamma$ the
family
\begin{equation}\label{hol:eq:monomialexp}
 \left(\frac{x^n}{n!}\right)_{n\ge0}
\end{equation}
is strongly summable exactly when $\gamma>0$, since $\gamma=0$ puts infinitely
many contributions at exponent zero and $\gamma<0$ gives a descending
sequence. Proposition~\ref{hol:prop:expdomain} computes the exact domain of
the corresponding family with an arbitrary Hahn argument.

\begin{proposition}[Closure]\label{hol:prop:closure}
$\E$ is a $\K$-algebra, $\sigma_q(\E)\subseteq\E$ for every $q\in\K^\times$,
and $D_z(\E)\subseteq\E$.
\end{proposition}

\begin{proof}
For addition, the evaluation supports lie in the union of two well-ordered
sets with finite exponent multiplicities. For multiplication, the doubly
indexed product family is strongly summable by
Lemma~\ref{hol:lem:support}\,(1) and (3), and grouping by ordinary total
degree preserves strong summability by the last clause of
Lemma~\ref{hol:lem:certificate}. Dilation is evaluation at $qx$. For $x\ne0$
the derivative family is
\[
 (n+1)a_{n+1}x^n=x^{-1}\cdot(n+1)a_{n+1}x^{n+1}.
\]
Discarding one term and multiplying each term by a nonzero ordinary integer
enlarges neither supports nor multiplicities, and multiplication by $x^{-1}$
preserves summability by Lemma~\ref{hol:lem:support}\,(3). At $x=0$ only the
constant term contributes. Iterate.
\end{proof}

\begin{proposition}[Faithfulness of evaluation]\label{hol:prop:faithful}
If $f\in\E$ vanishes at infinitely many distinct ordinary complex arguments,
then $f=0$ as a formal series. Consequently a linear equation with
rational-function coefficients that holds pointwise wherever its coefficients
are defined is also a formal equation in the sense of
Definition~\ref{hol:def:operators}.
\end{proposition}

\begin{proof}
Strong evaluation at $1$ says that the coefficient family $(a_n)$ has a
well-ordered union of supports and is locally finite at every exponent. Hence
for each $\gamma$
\[
 P_\gamma(T)=\sum_n(a_n)_\gamma T^n\in\C[T]
\]
is a polynomial, and its value at an ordinary complex $c$ is the coefficient
of $t^\gamma$ in $f(c)$. If $f$ vanishes at infinitely many such $c$, every
$P_\gamma$ has infinitely many roots and is zero, so every $a_n$ is zero. For
the last assertion, clear denominators; the left-hand side is in $\E$ by
Proposition~\ref{hol:prop:closure}, and nonzero polynomial denominators
exclude at most finitely many ordinary complex arguments.
\end{proof}

The nonexistence proofs may therefore be run on formal coefficients without
weakening their reading as functional obstructions, and no interchange of
infinite limits is needed anywhere.

The theorems below use only the valuation laws
\eqref{hol:eq:valuationlaws}, the existence of every monomial $t^\gamma$,
triviality of the valuation on $\C^\times$, and the support facts of this
section. They do not require algebraic closedness, spherical completeness,
Weierstrass preparation, or an embedding of $\Gamma$ into $\R$.

\subsection{Embedding the workspace into the surcomplex numbers}
\label{hol:sub:embedding}

If $\Gamma$ is an additive ordered subgroup of $\No$, the Conway normal-form
map
\begin{equation}\label{hol:eq:embedding}
 \sum_{\gamma\in S}c_\gamma t^\gamma
 \longmapsto\sum_{\gamma\in S}c_\gamma\omega^{-\gamma}
\end{equation}
identifies the real-coefficient Hahn field with a subfield of $\No$ and its
complexification with a subfield of $\No[i]$; a well-ordered $S$ becomes a
reverse-well-ordered set of Conway exponents. This viewpoint and its
set-sized support discipline are standard~\cite{bg,kks}. Here $t^\gamma$ is a
formal monomial and $\omega^{-\gamma}$ a Conway monomial; the notation does
not silently identify every formal power with a choice of surreal
exponentiation.

The results may therefore be read twice. Abstractly they concern Hahn fields
and need no surreal construction. Concretely they restrict the exact linear
descriptions available for surcomplex-valued power series on specified surreal
scales. Formula~\eqref{hol:eq:embedding} is the only bridge used, and no
global surcomplex exponential is invoked.

\section{The finite-step escape mechanism}\label{hol:sec:escape}

The following isolates the mechanism used throughout. It needs no formula for
the recurrence coefficients, no asymptotic slope for the valuations, and no
assumption that cancellations are rare.

\begin{lemma}[Escape chain]\label{hol:lem:escape}
Suppose a sequence $(a_n)$ in $\K$ satisfies, for every $n\ge N$,
\begin{equation}\label{hol:eq:generalrecurrence}
 c_0(n)a_n+\sum_{j=1}^sc_j(n)a_{n+j}=0,\qquad c_0(n)\ne0,
\end{equation}
and suppose $a_{n_0}\ne0$ for some $n_0\ge N$. Then $s\ge1$ and there is an
infinite sequence $n_0<n_1<n_2<\cdots$ with $j_r:=n_{r+1}-n_r\in\{1,\ldots,s\}$,
with $a_{n_r}\ne0$ and $c_{j_r}(n_r)\ne0$ for every $r$, and with
\begin{equation}\label{hol:eq:escapeineq}
 v(a_{n_{r+1}})-v(a_{n_r})\le v(c_0(n_r))-v(c_{j_r}(n_r)).
\end{equation}
\end{lemma}

\begin{proof}
Let $n\ge N$ with $a_n\ne0$. The leftmost term of
\eqref{hol:eq:generalrecurrence} is nonzero, so $s\ge1$ and at least one later
term is nonzero. If every nonzero later term had valuation strictly larger
than $v(c_0(n)a_n)$, the leading monomial of the leftmost term could not
cancel, contradicting \eqref{hol:eq:valuationlaws}. So some
$j\in\{1,\ldots,s\}$ has $c_j(n)a_{n+j}\ne0$ and
\[
 v(c_j(n)a_{n+j})\le v(c_0(n)a_n),
\]
which is \eqref{hol:eq:escapeineq}. Taking the least eligible $j$ makes the
recursion definite. The successor coefficient is nonzero and its index is
again at least $N$, so the construction never stops.
\end{proof}

\begin{remark}[The direction of the inequality matters]
\label{hol:rem:direction}
Solving the recurrence for its \emph{last} coefficient typically yields a
\emph{lower} bound on a valuation, which obstructs nothing. The comparison
above is made from the first nonzero coefficient forward and produces a later
coefficient whose valuation is controlled from above. Cancellation among the
other terms can raise the valuation of their sum, but cannot make every term
on the right have valuation strictly larger than the single nonzero term on
the left. Zeros and arbitrarily large cancellations elsewhere therefore do not
invalidate the chain; the chain records only that finite inequality,
repeatedly.
\end{remark}

\begin{theorem}[Evaluation exclusion criterion]\label{hol:thm:escapeexclusion}
In the situation of Lemma~\ref{hol:lem:escape}, let $x\in\K^\times$ and
suppose that for every $n\ge N$ and every $j\in\{1,\ldots,s\}$ with
$c_j(n)\ne0$,
\begin{equation}\label{hol:eq:boundincrements}
 v(c_0(n))-v(c_j(n))+jv(x)\le0.
\end{equation}
If $a_{n_0}\ne0$ for some $n_0\ge N$, then $(a_nx^n)_{n\ge0}$ is not strongly
summable. Equivalently: if $(a_nx^n)_{n\ge0}$ is strongly summable, then
$a_n=0$ for every $n\ge N$.
\end{theorem}

\begin{proof}
Build the chain of Lemma~\ref{hol:lem:escape} from $n_0$ and add
$n_{r+1}v(x)$ to each side of \eqref{hol:eq:escapeineq}, using
$n_{r+1}=n_r+j_r$:
\[
 v(a_{n_{r+1}}x^{n_{r+1}})
 \le v(a_{n_r}x^{n_r})+v(c_0(n_r))-v(c_{j_r}(n_r))+j_rv(x)
 \le v(a_{n_r}x^{n_r}).
\]
The evaluated leading exponents along the chain are therefore nonincreasing,
and Lemma~\ref{hol:lem:nonincreasing} applies.
\end{proof}

Note that \eqref{hol:eq:boundincrements} is a weak inequality. This is the
difference between an open and a closed excluded region, and it is available
only because Lemma~\ref{hol:lem:nonincreasing} also rules out a constant
leading exponent. Section~\ref{hol:sub:expdomain} exhibits an equation whose
excluded region is exactly the closed one, and
Section~\ref{hol:sub:coarseboundary} exhibits a coarsened comparison where the
boundary genuinely must be left in.

\begin{corollary}[Uniform bounded cost; no divisibility]
\label{hol:cor:boundedcost}
Suppose in \eqref{hol:eq:generalrecurrence} there is a single $B\in\Gamma$ with
$v(c_0(n))-v(c_j(n))\le B$ for all $n\ge N$ and all $j\ge1$ with $c_j(n)\ne0$.
Put $B^{+}=\max(B,0)\in\Gamma$. Then for every $x\ne0$ with $v(x)\le-B^{+}$,
strong summability of $(a_nx^n)_n$ forces $a_n=0$ for all $n\ge N$. Such $x$
exist: $x=t^{-B^{+}}$, and $x=t^{-\delta}$ for every $\delta\ge B^{+}$.
\end{corollary}

\begin{proof}
Since $B^{+}\ge0$ we have $v(x)\le0$, so $jv(x)\le v(x)\le-B^{+}\le-B$ for
$j\ge1$. Hence
$v(c_0(n))-v(c_j(n))+jv(x)\le B-B=0$, and
Theorem~\ref{hol:thm:escapeexclusion} applies. Every ordered abelian group
admits a maximum of finitely many elements and the monomial $t^{-B^{+}}$
exists, so no cofinal cyclic subgroup is needed.
\end{proof}

The absence of a cofinality requirement here is exactly why the mechanism
proves rigidity in groups of very large rank. It is also why the refined
bound below, which divides by $j$, is a refinement of the constant and never
of the conclusion.

\begin{corollary}[Refined periodic threshold; requires $\Gamma$ divisible]
\label{hol:cor:periodic}
Assume in addition that $\Gamma$ is divisible. Assume each slot of
\eqref{hol:eq:generalrecurrence} is either identically zero for $n\ge N$ or
nonzero for all $n\ge N$ with periodic valuation there of common period
$m\ge1$; enlarge $N$ past any initial nonperiodic range first.
Let $\alpha_{j,r}$ be the valuation of slot $j$ on the residue class
$r\bmod m$, the identically zero slots being omitted. Put
\begin{equation}\label{hol:eq:radiusbound}
 B_{\mathrm{rec}}
 =\max_{\substack{0\le r<m\\ 1\le j\le s,\ c_j\not\equiv0}}
   \frac{\alpha_{0,r}-\alpha_{j,r}}{j}\ \in\Gamma.
\end{equation}
If there is no active forward slot $j\ge1$, set $B_{\mathrm{rec}}=0$;
the recurrence then forces $a_n=0$ for all $n\ge N$ directly.
Then for every $x\ne0$ with $v(x)\le-B_{\mathrm{rec}}$, strong summability of
$(a_nx^n)_n$ forces $a_n=0$ for all $n\ge N$. In particular every solution
with $a_n\ne0$ for some $n\ge N$ --- so in particular every nonpolynomial
solution --- satisfies
\begin{equation}\label{hol:eq:excludedball}
 \Dom(f)\subseteq\{0\}\cup\{x\in\K^\times:v(x)>-B_{\mathrm{rec}}\},
 \qquad f=\sum_na_nz^n.
\end{equation}
\end{corollary}

\begin{proof}
The case with no forward slot was handled in the statement. Otherwise
divisibility makes each quotient $(\alpha_{0,r}-\alpha_{j,r})/j$ an element of
$\Gamma$, so $B_{\mathrm{rec}}$ is defined. If $v(x)\le-B_{\mathrm{rec}}$ then
$jv(x)\le-j\,(\alpha_{0,r}-\alpha_{j,r})/j=-(\alpha_{0,r}-\alpha_{j,r})$,
because multiplying a weak inequality by the positive integer $j$ preserves
it. Hence \eqref{hol:eq:boundincrements} holds and
Theorem~\ref{hol:thm:escapeexclusion} applies.
\end{proof}

The inequality in \eqref{hol:eq:excludedball} is strict on the right: at
$v(x)=-B_{\mathrm{rec}}$ the chain may keep a constant leading exponent rather
than descend, which is precisely why both halves of
Definition~\ref{hol:def:entire} are needed. The bound is an \emph{exclusion}
bound. It does not assert that every point on the other side lies in the
domain.

\section{Valuations on integer and geometric orbits}\label{hol:sec:orbits}

\subsection{Ordinary integers cannot create unbounded valuation loss}

\begin{lemma}[Integer evaluation]\label{hol:lem:integer}
Let $0\ne P(X)=\sum_kp_kX^k\in\K[X]$ and put $\beta(P)=\min_{p_k\ne0}v(p_k)$.
For all but finitely many nonnegative ordinary integers $n$,
\begin{equation}\label{hol:eq:integerstable}
 P(n)\ne0\qquad\text{and}\qquad v(P(n))=\beta(P).
\end{equation}
\end{lemma}

\begin{proof}
The polynomial $t^{-\beta(P)}P$ has coefficients of nonnegative valuation and
at least one of valuation zero, so its coefficientwise reduction
\[
 \overline P(X)=\sum_{v(p_k)=\beta(P)}\lc(p_k)X^k\in\C[X]
\]
is nonzero and has finitely many roots. The ordinary integers have distinct
images in $\C$. At every integer $n$ with $\overline P(n)\ne0$, the Hahn
coefficient of $t^{-\beta(P)}P(n)$ at exponent zero is nonzero, and no smaller
exponent occurs.
\end{proof}

Residue characteristic zero does real work here. The proof does not apply to
$p$-adic integer inputs, where positive integers can have arbitrarily large
valuation. Conversely no Diophantine approximation over $\C$ is relevant to
$v$: every nonzero complex number, however small in ordinary modulus, has
valuation zero.

\subsection{A geometric orbit of nonzero valuation}

\begin{lemma}[Eventually affine valuation]\label{hol:lem:affine}
Let $q\in\K^\times$ with $\lambda=v(q)\ne0$ and let
$0\ne P(X)=\sum_kp_kX^k\in\K[X]$. There are an index $k_P$ with $p_{k_P}\ne0$,
the value $\beta_P=v(p_{k_P})$, and $N_P\in\N$, such that
\begin{equation}\label{hol:eq:affine}
 v(P(q^n))=\beta_P+k_Pn\lambda\qquad(n\ge N_P).
\end{equation}
In particular $P(q^n)$ is eventually nonzero.
\end{lemma}

\begin{proof}
The nonzero terms have valuations $v(p_k)+kn\lambda$. For two distinct indices
the difference $v(p_k)-v(p_l)+(k-l)n\lambda$ is strictly monotone in $n$, so
its sign is eventually constant and it vanishes for at most one integer,
whether or not its values are cofinal in any convex subgroup. Finitely many
pairs give one common threshold beyond which every pairwise comparison is
strict and fixed, hence a unique minimizing index $k_P$. Its term cannot
cancel against terms of strictly larger valuation.
\end{proof}

\begin{remark}[A rank-one shortcut that fails]\label{hol:rem:shortcut}
The minimizing index need not be the smallest when $v(q)>0$. In
$\Gamma=\Q H\oplus\Q\epsilon$ with $H$ dominating every $n\epsilon$, take
$q=t^\epsilon$ and $P(X)=X+t^{-H}$: the constant term has smaller valuation
than $q^n$ for every $n$, however large. Conversely with $P(X)=1+t^{-H}X$ the
nonconstant term is the minimum for every $n$. The same phenomenon in the
other sign convention: in $\Gamma=\Q e_0\oplus\Q e_1$ with $e_1\gg e_0>0$,
$q=t^{e_0}$ and $P(X)=t^{e_1}+X$ give $v(P(q^n))=ne_0$ for every $n\ge0$, so
the nonzero constant coefficient never becomes the least-valued term. The
proof above uses the eventual ordering of finitely many affine functions, not
the Archimedean assertion that $n\epsilon$ eventually exceeds $H$, and not the
rank-one slogan that $q^n$ eventually beats every coefficient ratio.
\end{remark}

\subsection{Valuation-zero dilations, including residues that are roots of unity}

Distinct powers of a valuation-zero $q$ all have the same valuation, so the
previous lemma says nothing. Polynomial evaluation nevertheless has only a
finite periodic valuation pattern.

\begin{theorem}[Unit-orbit valuations]\label{hol:thm:unitorbit}
Let $q\in\K^\times$ have $v(q)=0$ and infinite multiplicative order, and let
$0\ne P\in\K[X]$. Then $v(P(q^n))$ is eventually finite and periodic. More
precisely, with $\zeta=\res(q)\in\C^\times$:
\begin{enumerate}[label=\textup{(\alph*)}]
\item if $\zeta$ is not a root of unity, $v(P(q^n))$ is eventually constant,
equal to the minimum valuation of the coefficients of $P$;
\item if $\zeta$ has exact order $h$, then $v(P(q^n))$ is eventually constant
on each residue class modulo $h$, so its eventual period divides $h$.
\end{enumerate}
\end{theorem}

\begin{proof}
For (a) write $P(X)=\sum_kp_kX^k$ and $\beta=\min_kv(p_k)$. The coefficient of
$P(q^n)$ at exponent $\beta$ is $\overline P(\zeta^n)$ with
$\overline P(X)=\sum_{v(p_k)=\beta}\lc(p_k)X^k$, a nonzero complex polynomial.
The points $\zeta^n$ are distinct, so this vanishes only finitely often, and
no exponent below $\beta$ occurs.

For (b) write
\[
 q=\zeta(1+u),\qquad \eta=v(u)>0,\qquad d=\lc(u).
\]
Here $u=q/\zeta-1$ has positive valuation because $q/\zeta$ has valuation zero
and residue $1$, and $u\ne0$ because otherwise $q=\zeta$ would have finite
order. For every positive ordinary integer $n$ the finite binomial expansion
gives
\begin{equation}\label{hol:eq:yn}
 y_n:=(1+u)^n-1,\qquad v(y_n)=\eta,\qquad \lc(y_n)=nd,
\end{equation}
the leading coefficient being nonzero in characteristic zero. There is no
convergence assertion: for each fixed $n$ only finitely many powers occur.

Fix $r\in\{0,\ldots,h-1\}$. For $n\equiv r\pmod h$ we have $\zeta^n=\zeta^r$,
hence $q^n=\zeta^r(1+y_n)$. Expand the nonzero polynomial
\begin{equation}\label{hol:eq:unitexpand}
 P(\zeta^r(1+Y))=\sum_{k=0}^{\deg P}b_{r,k}Y^k
\end{equation}
and set
\begin{equation}\label{hol:eq:betar}
 \beta_r=\min_{b_{r,k}\ne0}\bigl(v(b_{r,k})+k\eta\bigr),
 \qquad J_r=\{k:v(b_{r,k})+k\eta=\beta_r\}.
\end{equation}
Then $P(q^n)=\sum_kb_{r,k}y_n^k$, and its coefficient at exponent $\beta_r$ is
the value at $n$ of
\begin{equation}\label{hol:eq:residuepoly}
 R_r(T)=\sum_{k\in J_r}\lc(b_{r,k})d^kT^k\in\C[T].
\end{equation}
This polynomial is nonzero: its terms have distinct degrees and all the
indicated coefficients are nonzero. Outside its finitely many integer roots
the coefficient at $\beta_r$ survives, so $v(P(q^n))=\beta_r$ eventually on
that class. Finitely many classes give one finite exceptional range.
\end{proof}

\begin{remark}[A second normalization of the same proof]
\label{hol:rem:unitalt}
Sources~08 and~09 both prove Theorem~\ref{hol:thm:unitorbit}\,(b), by two
slightly different normalizations, and we record the other one because it is
occasionally more convenient. Instead of factoring out the complex root of
unity, put $q^h=1+u$ directly, with $\eta=v(u)>0$ and $u\ne0$ for the same
reason. Write $n=r+hk$ and expand the polynomial $Q_r(Y)=P(q^rY)$ as
$Q_r(1+V)=\sum_kb_{r,k}V^k$. The finite binomial expansion gives
$V_k=(1+u)^k-1$ with $v(V_k)=\eta$ and $\lc(V_k)=k\lc(u)$, and
$P(q^n)=Q_r(1+V_k)$. The minimum construction \eqref{hol:eq:betar} and the corresponding
nonzero residue polynomial \eqref{hol:eq:residuepoly}, now in the variable
$k$ rather than $n$, finish the argument. Neither version sums a binomial
series to a noninteger exponent; every expansion is finite, and arbitrary
supports in $q$ and in the coefficients of $P$ enter only through the leading
valuation and leading coefficient of finitely many Hahn elements.
\end{remark}

\begin{example}[A genuine two-phase pattern]\label{hol:ex:unitperiod}
Let $q=-(1+t^\eta)$ with $\eta>0$. Then $q$ is nontorsion, $\res(q)=-1$ has
order two, and for $P(X)=X-1$ and $n\ge1$
\[
 v(P(q^n))=v(q^n-1)=
 \begin{cases}
 0,&n\text{ odd},\\
 \eta,&n\text{ even}.
 \end{cases}
\]
For odd $n$ the residue is $-2$; for even positive $n$ the leading term is
$nt^\eta$. The valuation sequence is periodic, not the dilation itself, and a
claim of eventual \emph{constancy} without splitting residue classes would be
false. The escape argument uses the bounded valuation costs; it does not
replace $q$ by its residue.
\end{example}

\begin{example}[Finite exceptional indices really occur]
\label{hol:ex:unitexception}
For $q=1+t^\eta$ and $P(X)=X-1-7t^\eta$, the leading coefficient of $P(q^n)$
at exponent $\eta$ is $n-7$, so $v(P(q^n))=\eta$ for positive $n\ne7$ while
$v(P(q^7))=2\eta$. With $P(X)=X-1-3t^\eta$ the exceptional index is $n=3$,
where $P(q^3)=3t^{2\eta}+t^{3\eta}$. The exceptions are harmless for the
escape argument, since it needs only an eventual range, but they cannot be
ignored in a statement asserted for all indices. This is why
Lemma~\ref{hol:lem:integer} and Theorem~\ref{hol:thm:unitorbit} both allow
finitely many exceptions even when $\res(q)=1$.
\end{example}

For an ordinary complex $q$ with $\abs q_\C\ne1$, the complex modulus does not
alter these conclusions: as an element of the Hahn coefficient field, $q$
still has valuation zero. The ordinary modulus and the scale valuation answer
different questions.

\section{Differential rigidity}\label{hol:sec:ode}

\subsection{From an operator to a forward recurrence}

A formal series is D-finite exactly when its derivatives span a
finite-dimensional space over $\K(z)$. The D-finite/P-recursive equivalence is
classical~\cite[Theorem~1.5]{stanley}; the following records the exact finite
conversion, because the valuations of the resulting recurrence, not ordinary
coefficient growth, are what the proof consumes.

\begin{lemma}[Differential coefficient recurrence]\label{hol:lem:drec}
Write
\begin{equation}\label{hol:eq:operator}
 L=\sum_{r=0}^R\sum_{k=0}^mc_{r,k}z^kD_z^r,\qquad c_{r,k}\in\K,
\end{equation}
with at least one $c_{r,k}\ne0$, and let $f=\sum_{n\ge0}a_nz^n$. For $n\ge m$
the coefficient of $z^n$ in $Lf$ is
\begin{equation}\label{hol:eq:odecoeff}
 \sum_{s=-m}^{R}Q_s(n)a_{n+s},\qquad
 Q_s(X)=\sum_{r-k=s}c_{r,k}\fall{X+s}{r},
\end{equation}
where $\fall Yr=Y(Y-1)\cdots(Y-r+1)$ and $\fall Y0=1$. At least one $Q_s$ is
nonzero. Writing $s_0$ for the least active shift and $u=n+s_0$, every formal
solution of $Lf=0$ satisfies, for all sufficiently large $u$,
\begin{equation}\label{hol:eq:precursive}
 \sum_{j=0}^{s}P_j(u)a_{u+j}=0,\qquad P_j(X)=Q_{s_0+j}(X-s_0),
 \qquad P_0\ne0,
\end{equation}
with $s=R-s_0$. If $s=0$, every formal solution is a polynomial.
\end{lemma}

\begin{proof}
The operator $z^kD_z^r$ carries the term of index $n+r-k$ to degree $n$ and
multiplies it by $\fall{n+s}{r}$ with $s=r-k$, which gives
\eqref{hol:eq:odecoeff}. For a fixed shift $s$ the distinct derivative orders
$r$ contribute polynomials of distinct degrees in $X$, so a nonempty group
containing a nonzero coefficient cannot cancel identically and some $Q_s$ is
nonzero. Reindexing gives \eqref{hol:eq:precursive}. If $s=0$ the recurrence
reads $P_0(u)a_u=0$ with $P_0\ne0$, so by Lemma~\ref{hol:lem:integer} all but
finitely many $a_u$ vanish.
\end{proof}

\begin{proposition}[Uniform obstruction for P-recursive coefficients]
\label{hol:prop:precursive}
Suppose $(a_n)$ satisfies \eqref{hol:eq:precursive} for all $n\ge N_1$, with
$P_j\in\K[X]$ and $P_0\ne0$. Let $\beta_j=\beta(P_j)$ for the nonzero $P_j$
and put
\[
 B=\max\bigl(\{0\}\cup\{\beta_0-\beta_j:j\ge1,\ P_j\ne0\}\bigr)\ \ge0.
\]
There is $N\ge N_1$, beyond the finitely many exceptional integers of
Lemma~\ref{hol:lem:integer} for the finitely many nonzero $P_j$, such that for
every $x\ne0$ with $v(x)\le-B$, strong summability of $(a_nx^n)_n$ forces
$a_n=0$ for all $n\ge N$. If $\Gamma$ is divisible one may replace $B$ by the
smaller threshold
\begin{equation}\label{hol:eq:refinedthreshold}
 B_{\mathrm{rec}}
 =\max\Bigl(\Bigl\{\tfrac{\beta_0-\beta_j}{j}:j\ge1,\ P_j\ne0\Bigr\}\Bigr),
\end{equation}
which is \eqref{hol:eq:radiusbound} with period $m=1$.
If the displayed set is empty, take $B_{\mathrm{rec}}=0$; then every
solution terminates beyond $N$ without an evaluation hypothesis.
\end{proposition}

\begin{proof}
Discard the identically zero $P_j$. By Lemma~\ref{hol:lem:integer} there is
$N$ beyond which $v(P_j(n))=\beta_j$ for all the finitely many nonzero $P_j$
simultaneously. Apply Corollary~\ref{hol:cor:boundedcost} with this $B$, or
Corollary~\ref{hol:cor:periodic} with $m=1$ when $\Gamma$ is divisible. If no
$j\ge1$ remains, the recurrence forces the same vanishing directly.
\end{proof}

These are certified exterior obstructions, not a computation of the exact
evaluation domain of each solution. Individual solutions can have stricter
restrictions or can terminate much earlier.

\begin{proof}[Proof of Theorem~\ref{hol:main:ode}]
Lemma~\ref{hol:lem:drec} converts $L$ into \eqref{hol:eq:precursive}. Take
$B_L=B$ from Proposition~\ref{hol:prop:precursive} and take $N_L$ to be an
positive integer beyond the recurrence's starting index and beyond every nonnegative
integer root of the finitely many residue polynomials $\overline{P_j}$. If a
formal solution is strongly evaluable at one $x\ne0$ with $v(x)\le-B_L$, then
$a_n=0$ for all $n\ge N_L$, so its degree is less than $N_L$. An entire
solution is in particular strongly evaluable at the available monomial
$t^{-B_L}$, hence is a polynomial. Conversely a polynomial is entire and is
annihilated by a sufficiently high derivative, which gives the displayed
equality.
\end{proof}

The closed region $v(x)\le-B_L$, rather than an open one, comes from the weak
inequality in Theorem~\ref{hol:thm:escapeexclusion}, hence from
Lemma~\ref{hol:lem:nonincreasing}. A polynomial solution of degree at least
$N_L$ would in any case contradict the recurrence at its last nonzero
coefficient, since $P_0(n)$ no longer vanishes there.

Theorem~\ref{hol:main:ode} is stronger than merely excluding a globally
extended exponential. The operator coefficients may themselves be arbitrary
Hahn series carrying infinite surreal scales. There are nevertheless only
finitely many of them, so their leading valuations give a finite obstruction
that one additional exterior scale defeats.

\subsection{An exact boundary, not only a nonexistence result}
\label{hol:sub:expdomain}

\begin{proposition}[Formal exponential domain]\label{hol:prop:expdomain}
For $\eta\in\Gamma$ let
\[
 E_\eta(z)=\sum_{n\ge0}\frac{t^{n\eta}}{n!}z^n.
\]
Then $D_zE_\eta=t^\eta E_\eta$ and
\begin{equation}\label{hol:eq:expdomain}
 \Dom(E_\eta)=\{0\}\cup\{x\in\K^\times:v(x)>-\eta\}.
\end{equation}
\end{proposition}

\begin{proof}
The differential identity is a formal coefficient comparison. Put $y=t^\eta x$
for $x\ne0$, so the $n$th evaluation term is $y^n/n!$.

If $v(y)>0$, write $y=c\,t^{v(y)}(1+u)$ with $c\in\C^\times$ and $u$ zero or
of positive valuation, and put $S=\supp(u)\subseteq\Gamma_{>0}$. Since
$(1+u)^n$ is a finite integer power, $\supp(y^n)\subseteq h_n+M(S)$ with
$h_n=nv(y)$, and $(h_n)$ is strictly increasing, so its values are well
ordered and each is taken once. Lemma~\ref{hol:lem:certificate}\,(b) makes the
family strongly summable; the factor $1/n!$ is a nonzero complex scalar and
changes no support.

If $v(y)<0$ the leading exponents $nv(y)$ strictly decrease. If $v(y)=0$, then
$y^n/n!$ has the nonzero coefficient $\lc(y)^n/n!$ at exponent zero for every
$n$, so local finiteness fails. These cases exhaust $x\ne0$, and $v(y)>0$ is
$v(x)>-\eta$.
\end{proof}

The recurrence of $E_\eta$ is $(n+1)a_{n+1}-t^\eta a_n=0$, with one active
shift $j=1$, so $\alpha_0=\eta$, $\alpha_1=0$, and
\eqref{hol:eq:radiusbound} gives $B_{\mathrm{rec}}=\eta$ with no division by
an integer at all. Comparing with \eqref{hol:eq:expdomain}: the excluded
region $v(x)\le-\eta$ of Corollary~\ref{hol:cor:periodic} is here
\emph{exactly} the complement of the domain, boundary included. The general
exclusion bound can therefore be attained, and the closed form of the boundary
is not an artefact of a lossy estimate. This example is the reason for keeping
Lemma~\ref{hol:lem:nonincreasing} rather than the strictly descending version
alone.

There is no conflict with any scalar exponential on a surreal field: $E_\eta$
is an ordinary power series in $z$, not a total exponential evaluated by an
independent construction. Taking $\eta=0$ recovers the elementary observation
\eqref{hol:eq:monomialexp}.

\subsection{Algebraic and system consequences}

\begin{corollary}[Algebraic entire series]\label{hol:cor:algebraic}
An element of $\E$ algebraic over $\K(z)$ belongs to $\K[z]$.
\end{corollary}

\begin{proof}
Let $F=\K(z,f)$ be the finite extension generated by the algebraic series. In
characteristic zero its minimal polynomial is separable, and differentiating
that polynomial expresses $D_zf$ inside $F$, because the derivative with
respect to the algebraic variable is nonzero at $f$ and hence invertible in
$F$. So $D_z$ preserves $F$, all derivatives of $f$ lie in the
finite-dimensional $\K(z)$-space $F$, and $f$ is D-finite. Apply
Theorem~\ref{hol:main:ode}.
\end{proof}

This is not offered as a new algebraic rigidity theorem; it records that the
linear result subsumes the algebraic case without any factorization theorem
for $\E$.

\begin{corollary}[Rational differential systems]\label{hol:cor:systems}
If a finite vector $Y$ of strongly entire formal series satisfies
$D_zY=A(z)Y$ with $A\in M_r(\K(z))$, every component of $Y$ is a polynomial.
The same holds for an inhomogeneous system with a rational forcing vector,
provided all its solution components are strongly entire.
\end{corollary}

\begin{proof}
The $\K(z)$-span of the components is stable under $D_z$, so the successive
derivatives of each component span a finite-dimensional space. For $Y'=AY+b$
append the constant component $1$ and use
$\left(\begin{smallmatrix}A&b\\0&0\end{smallmatrix}\right)$.
\end{proof}

\begin{proposition}[Rational inhomogeneous equations]
\label{hol:prop:inhomogeneous}
The polynomiality conclusions of this report remain valid when the right-hand
side of a linear equation is a rational function of $z$, provided the
corresponding homogeneous hypothesis on $L$ or on $q$ holds.
\end{proposition}

\begin{proof}
Clear denominators in the coefficients and in the right-hand side. The result
has a polynomial right-hand side, so the coefficient recurrence is homogeneous
for all sufficiently large indices, and the escape statements were formulated
for eventual recurrences. The left-hand operator is still nonzero, having been
multiplied by a nonzero polynomial.
\end{proof}

This applies, for instance, to the inhomogeneous first-order theta identity
\eqref{hol:eq:thetafirst} when the parameter valuation is noncofinal: it is
not necessary to homogenize it in order to see the obstruction.

\begin{example}[A sparse recurrence does not escape the obstruction]
\label{hol:ex:airy}
The formal Airy-type equation $f''-zf=0$ has the recurrence
\[
 a_{n+3}=\frac{a_n}{(n+3)(n+2)}\quad(n\ge0),\qquad a_2=0.
\]
With $a_0=1$ and $a_1=0$ there are infinitely many nonzero coefficients, at
indices $3n$, all of valuation zero. At $x=t^{-\delta}$ their leading
exponents are $-3n\delta$, strictly decreasing for every $\delta>0$. Gaps
between nonzero coefficients do not help: the escape chain follows only the
surviving congruence class.
\end{example}

\begin{remark}[What transcendence has actually been proved]
\label{hol:rem:transcendence}
Every nonpolynomial member of $\E$ is transcendental over $\K(z)$ and
satisfies no nonzero \emph{linear} differential equation over $\K(z)$, and no
finite family $f,D_zf,\ldots,D_z^Rf$ is linearly dependent over $\K(z)$. This
does not imply differential transcendence in the nonlinear sense. A relation
\[
 P\bigl(z,f,D_zf,\ldots,D_z^Rf\bigr)=0
\]
involving products of $f$ and its derivatives lies outside the present proofs
unless it can separately be reduced to a finite linear system. The
distinction is not cosmetic: the escape mechanism exploits a finite
\emph{linear} cancellation in which one nonzero coefficient forces a later
nonzero coefficient, whereas nonlinear coefficient equations involve
convolutions of many earlier coefficients and supply no such forward step
without further work.

That further work is the nonlinear part of this report
(Sections~\ref{hol:nl:sec:setting}--\ref{hol:nl:sec:consequences}, from
source~10). It does not repair the escape chain; it replaces the recurrence by
a finite initial polynomial at one large scale. Every relation of order
$R\le1$ is thereby excluded (Theorem~\ref{hol:nl:main:first}), and so is every
relation of higher order whose residue corner polynomial is nonzero
(Theorem~\ref{hol:nl:thm:corner}). Relations of order at least two with a
vanishing residue corner remain outside the proofs of
Sections~\ref{hol:nl:sec:setting}--\ref{hol:nl:sec:consequences}. When
$\Gamma$ has no order unit they are excluded in every order by a third,
different argument, which bounds the field generated by the Taylor
coefficients (Theorem~\ref{hol:cf:main:entire}, sources~11 and~12). With an
order unit they are not always excluded: the series $\mathcal T_p$ of
Theorem~\ref{hol:td:main:boundary} (source~13) is a nonpolynomial strongly
entire solution of a relation of order three with vanishing residue corner.
Relations of order two are excluded for every $\Gamma$ by a fourth argument,
through initial polynomials at infinitely many scales
(Theorem~\ref{hol:ot:main:second}, source~14), so three is the least order.
\end{remark}

\section{Generic lines and several-variable rigidity}\label{hol:sec:multi}

A single differential equation in several variables is not enough to force
polynomiality: it could ignore most of the variables. The correct hypothesis
is finite-dimensionality of the span of \emph{all} mixed partial derivatives,
which is what Theorem~\ref{hol:main:multi} assumes.

\begin{lemma}[Avoiding countably many polynomial conditions]
\label{hol:lem:generic}
For any countable family of nonzero polynomials in $\K[z_1,\ldots,z_d]$ there
is $c\in\C^d$ at which none of them vanishes.
\end{lemma}

\begin{proof}
For one polynomial $H$, take the least Hahn exponent among its finitely many
nonzero coefficients; the coefficient at that exponent is a nonzero
$\overline H\in\C[z_1,\ldots,z_d]$, and $\overline H(c)\ne0$ implies
$H(c)\ne0$. The coefficients of the countably many polynomials $\overline H$
form a countable subset of $\C$, and the field they generate over $\Q$ is
countable. Choose $c_1,\ldots,c_d$ algebraically independent over that field,
successively outside the algebraic closure of a countable field, which remains
countable while $\C$ is uncountable.
\end{proof}

The same countable-field device is used in the companion report held in
\path{docs/surcomplex/entire-functions-at-arbitrary-rank/}
(\texttt{ent:mv:thm:exceptional}) to show that the set of polynomial
directions of a nonpolynomial series is not all of projective space.

\begin{proof}[Proof of Theorem~\ref{hol:main:multi}]
Suppose $F$ is not a polynomial and write it by total degree,
\begin{equation}\label{hol:eq:homogeneous}
 F(z)=\sum_{n\ge0}F_n(z),\qquad F_n\text{ homogeneous of degree }n,
\end{equation}
so infinitely many $F_n$ are nonzero. Let $g_1,\ldots,g_\sigma$ be a basis of
the finite-dimensional $\K(z_1,\ldots,z_d)$-span of all mixed derivatives of
$F$, taken to be mixed derivatives themselves with $g_1=F$. The span is stable
under differentiation, so there are rational matrices $A_i$ with
\begin{equation}\label{hol:eq:pfaffian}
 \partial_ig=A_i(z)g\qquad(1\le i\le d).
\end{equation}
These matrices have finitely many nonzero denominator polynomials; choose one
nonzero homogeneous component of each. Apply Lemma~\ref{hol:lem:generic} to
every nonzero $F_n$, to those finitely many chosen components, and to the
coordinate polynomials $z_i$, obtaining $c\in\C^d$ with $F_n(c)\ne0$ whenever
$F_n\ne0$ and with no denominator becoming the zero polynomial under
$z=wc$. Then
\[
 h(w)=F(wc)=\sum_{n\ge0}F_n(c)w^n
\]
is nonpolynomial.

For each $x\in\K$ the family
$(a_\alpha c^\alpha x^{\abs\alpha})_{\alpha\in\N^d}$ is strongly summable, by
entireness of $F$ at $(c_1x,\ldots,c_dx)$. Group it into the finite blocks of
fixed total degree; by the last clause of Lemma~\ref{hol:lem:certificate},
$\sum_nF_n(c)x^n$ is strongly summable, so $h\in\E$.

The formal chain rule, applied after clearing denominators in
\eqref{hol:eq:pfaffian}, gives
\begin{equation}\label{hol:eq:restricted}
 \frac{d}{dw}g(wc)=\Bigl(\sum_{i=1}^dc_iA_i(wc)\Bigr)g(wc).
\end{equation}
Every denominator here is a nonzero element of $\K[w]$. It may vanish at
$w=0$, which is harmless: read the identity in the Laurent series field
$\K((w))$, with matrices in $M_\sigma(\K(w))$. The finite $\K(w)$-span of the
components of $g(wc)$ is preserved by $d/dw$, so its first component $h$ is
D-finite over $\K(w)$, and Theorem~\ref{hol:main:ode} forces $h$ to be a
polynomial. This contradicts its construction.
\end{proof}

\begin{corollary}[Multivariate algebraic rigidity]\label{hol:cor:multialg}
If a strongly entire $F\in\K[[z_1,\ldots,z_d]]$ is algebraic over
$\K(z_1,\ldots,z_d)$, it is a polynomial.
\end{corollary}

\begin{proof}
Each partial derivation extends to the finite separable extension generated by
$F$, by differentiation of its minimal polynomial, so all mixed partial
derivatives lie in one finite-dimensional space. Apply
Theorem~\ref{hol:main:multi}.
\end{proof}

The generic line is an existence device, not a promised effective choice of
ordinary complex constants. The proof does not say that one prescribed line
detects every nonpolynomial multivariate series; the line is chosen only after
the countably many nonzero homogeneous parts and the finitely many rational
connection matrices are known.

\section{Dilation equations and the necessity of a cofinal scale}
\label{hol:sec:qnecessity}

\subsection{Exact coefficient conversion}

\begin{lemma}[Dilation coefficient recurrence]\label{hol:lem:qrec}
Let $q\in\K^\times$ have infinite multiplicative order and let $f=\sum a_nz^n$
satisfy \eqref{hol:eq:qlinear}. Write $p_\ell(z)=\sum_kp_{\ell,k}z^k$ and
$Q_k(X)=\sum_\ell p_{\ell,k}X^\ell$, let $d$ and $e$ be the largest and
smallest $k$ with $Q_k\ne0$, and put $s=d-e$. Then for every $n\ge0$
\begin{equation}\label{hol:eq:qrec}
 \sum_{h=0}^sR_h(q^n)a_{n+h}=0,\qquad
 R_h(X)=Q_{d-h}(q^hX),
\end{equation}
with $R_0=Q_d\ne0$ and $R_s=Q_e(q^s\,\cdot\,)\ne0$. Each nonzero $R_h$
satisfies $R_h(q^n)\ne0$ for all but finitely many $n$. If $s=0$, every formal
solution is a polynomial.
\end{lemma}

\begin{proof}
The coefficient of $z^{n+d}$ in $\sum_\ell p_\ell(z)f(q^\ell z)$ is
$\sum_{k=e}^{d}Q_k(q^{n+d-k})a_{n+d-k}$; substituting $h=d-k$ and using
$Q_{d-h}(q^{n+h})=Q_{d-h}(q^h\cdot q^n)$ gives \eqref{hol:eq:qrec}. Since $q$
has infinite order the points $q^n$ are distinct, and a nonzero univariate
polynomial vanishes at only finitely many of them; this already gives eventual
nonvanishing, independently of the finer valuation results. If $s=0$ the
relation reads $R_0(q^n)a_n=0$ for every $n$, so all but finitely many $a_n$
vanish.
\end{proof}

The powers $q^h$ in \eqref{hol:eq:qrec} matter: dropping them changes the
recurrence and can change its valuation constants. Equivalently, taking
$m=\max_\ell\deg p_\ell$ and reading off the coefficient of $z^{n+m}$ gives
$\sum_jP_j(q^n)a_{n+j}=0$ with
$P_j(X)=\sum_\ell p_{\ell,m-j}q^{\ell j}X^\ell$ and $P_0=Q_m\ne0$; both
conversions are checked on finite exact examples by the programs of
Section~\ref{hol:sec:verification}. If $m=0$ the equation already reads
$P_0(q^n)a_n=0$ and all solutions are polynomials with no entireness
assumption at all, so we may assume $s\ge1$ below.

\subsection{The valuation-zero case}

\begin{corollary}[Every unit dilation is globally rigid]
\label{hol:cor:unitrigid}
Let $q\in\K^\times$ have $v(q)=0$ and infinite multiplicative order. Every
linearly $q$-difference finite member of $\E$ is a polynomial, and each fixed
equation \eqref{hol:eq:qlinear} has a uniform exterior obstruction and a
uniform degree bound. If $\Gamma$ is in addition divisible, every
nonpolynomial formal solution has its domain bounded as in
\eqref{hol:eq:excludedball} by the constant \eqref{hol:eq:radiusbound} formed
from the eventual residue-class valuations of \eqref{hol:eq:qrec}.
\end{corollary}

\begin{proof}
Apply Theorem~\ref{hol:thm:unitorbit} simultaneously to the finitely many
nonzero $R_h$ of Lemma~\ref{hol:lem:qrec}. Their valuations take only finitely
many values after a common starting index, with common period one or the exact
order $h_0$ of $\res(q)$. All differences $v(R_0(q^n))-v(R_h(q^n))$ are
therefore bounded above by a single element of $\Gamma$, and
Corollary~\ref{hol:cor:boundedcost} gives the exterior obstruction and, via
the residue polynomials, the degree bound. When $\Gamma$ is divisible,
Corollary~\ref{hol:cor:periodic} applies with $m=1$ if the residue has
infinite order, and $m=h_0$ otherwise. In both cases take the starting
index beyond the exceptional range.
\end{proof}

This includes parameters arbitrarily close to a root of unity without being
one. The finite exceptional cancellations of Theorem~\ref{hol:thm:unitorbit}
can move the starting index but do not remove the bound.

\subsection{Nonzero but noncofinal valuation}

\begin{proposition}[Noncofinal dilation scales force polynomiality]
\label{hol:prop:noncofinal}
Let $q\in\K^\times$ with $\beta=v(q)\ne0$ and suppose $\abs\beta$ is not an
order unit of $\Gamma$. Every linearly $q$-difference finite member of $\E$ is
a polynomial, with a uniform exterior obstruction and a uniform degree bound
for each fixed equation. Divisibility of $\Gamma$ is not used.
\end{proposition}

\begin{proof}
By Lemma~\ref{hol:lem:affine} applied to the finitely many nonzero $R_h$ of
\eqref{hol:eq:qrec}, there is a common index beyond which
\[
 v(R_h(q^n))=\alpha_h+k_hn\beta,\qquad k_h\in\N.
\]
Put $V=\abs\beta>0$, choose $B\in\Gamma$ with $B\ge0$ and
$B\ge\abs{\alpha_0-\alpha_h}$ for all relevant $h$, and put
$M=\max_h\abs{k_0-k_h}\in\N$. Then for $h\ge1$
\begin{equation}\label{hol:eq:linearbound}
 v(R_0(q^n))-v(R_h(q^n))=\alpha_0-\alpha_h+(k_0-k_h)n\beta\le B+MnV.
\end{equation}
Because $\abs\beta$ is not an order unit, \eqref{hol:eq:dominates} supplies
$\gamma>0$ with $\gamma>NV$ for every $N\in\N$; in particular $MnV<\gamma$ for
all $n$, trivially so if $M=0$. Hence the right side of
\eqref{hol:eq:linearbound} is bounded above by the single element $B+\gamma$.
Corollary~\ref{hol:cor:boundedcost} with this bound excludes every $x\ne0$
with $v(x)\le-(B+\gamma)$, and the explicit witness $x=t^{-(B+\gamma)}$ is
available. Since an entire solution is evaluable there, it satisfies $a_n=0$
beyond the common starting index, hence is a polynomial of degree less than
that index. Nothing in this argument divides an element of $\Gamma$ by an
integer.
\end{proof}

The proof covers two different higher-rank situations. The group may have no
order unit at all, so every nonzero dilation valuation is noncofinal;
alternatively the group may have an order unit while the chosen $q$ lives on a
smaller scale. Global $q$-difference finiteness depends on the valuation of
that particular dilation, not only on the group. The argument also requires no
asymptotic formula for $v(a_n)$ and no eventually quadratic coefficient
valuation: a single infinite escape chain suffices, which is why cancellations
and sparse zero coefficients do not force a separate Skolem--Mahler--Lech
analysis.

\subsection{An intrinsic coarsened exclusion bound}

The previous proof produces more than one bad point. Let $\beta=v(q)\ne0$, let
$\Delta=\Delta_{\abs\beta}$ be the convex subgroup \eqref{hol:eq:delta}, and
assume $\Delta\ne\Gamma$. Let $\pi:\Gamma\to\Gamma/\Delta$ be the ordered
quotient and put $w=\pi\circ v$. With the eventual constants $\alpha_h$ of the
previous proof, put
\begin{equation}\label{hol:eq:coarsebound0}
 \overline B^{\flat}_{\mathrm{rec}}
 =\max\bigl(\{0\}\cup\{\pi(\alpha_0)-\pi(\alpha_h):1\le h\le s,\ R_h\ne0\}\bigr),
\end{equation}
and, when $\Gamma$ is divisible so that $\Gamma/\Delta$ is a divisible ordered
group,
\begin{equation}\label{hol:eq:coarsebound}
 \overline B_{\mathrm{rec}}
 =\max_{1\le h\le s,\ R_h\ne0}\frac{\pi(\alpha_0)-\pi(\alpha_h)}{h}.
\end{equation}

\begin{theorem}[Coarsened exclusion]\label{hol:thm:coarse}
Let $f=\sum a_nz^n$ be a nonpolynomial formal solution of
\eqref{hol:eq:qlinear} with $v(q)\ne0$ and $\Delta=\Delta_{\abs{v(q)}}$
proper.
\begin{enumerate}[label=\textup{(\alph*)}]
\item Under Convention~\ref{hol:conv:group} alone, every $x\ne0$ with
$-w(x)>\overline B^{\flat}_{\mathrm{rec}}$ lies outside $\Dom(f)$.
\item If $\Gamma$ is in addition divisible, every $x\ne0$ with
$-w(x)>\overline B_{\mathrm{rec}}$ lies outside $\Dom(f)$.
\end{enumerate}
\end{theorem}

\begin{proof}
Coarsening by $\Delta$ kills $n\beta$, since $n\beta\in\Delta$ for every $n$,
so the projected valuations of the recurrence coefficients are the constants
$\pi(\alpha_h)$. In case~(b), $-w(x)>\overline B_{\mathrm{rec}}$ gives
$hw(x)<-(\pi(\alpha_0)-\pi(\alpha_h))$ for every relevant $h\ge1$. In case~(a)
we have $w(x)<-\overline B^{\flat}_{\mathrm{rec}}\le0$, hence
$hw(x)\le w(x)<-\overline B^{\flat}_{\mathrm{rec}}
 \le-(\pi(\alpha_0)-\pi(\alpha_h))$. Either way the projected leading
valuation along an escape chain strictly decreases at every step. A strict
decrease in the ordered quotient is a strict decrease in $\Gamma$, so the
chain produces an infinite strictly descending sequence in the union of the
supports.
\end{proof}

The boundary is deliberately \emph{not} excluded here. Equal leading values in
the quotient can correspond to a strictly increasing sequence in $\Gamma$, and
hence to a strongly summable family; Section~\ref{hol:sub:coarseboundary}
exhibits exactly that. Unlike Corollary~\ref{hol:cor:periodic}, a coarsened
comparison alone cannot decide local finiteness at an exponent of $\Gamma$.

\section{Partial theta series: sharp witness and exact domain}
\label{hol:sec:theta}

\subsection{The classical identities}

For $p\in\K^\times$ with $\beta=v(p)>0$ set
\begin{equation}\label{hol:eq:theta}
 \Theta_p(z)=\sum_{n\ge0}p^{\binom n2}z^n
 =\sum_{n\ge0}p^{n(n-1)/2}z^n.
\end{equation}
This is a standard partial theta series up to a harmless sign
convention~\cite{andrewswarnaar}. None of its coefficients vanishes, so it is
not a polynomial. The coefficient identity
$p^{n(n-1)/2}=p^{\,n-1}p^{(n-1)(n-2)/2}$ for $n\ge1$ gives
\begin{equation}\label{hol:eq:thetafirst}
 \Theta_p(z)=1+z\Theta_p(pz),
\end{equation}
and applying $\sigma_p-1$ to $(1-z\sigma_p)\Theta_p=1$, with
$\sigma_pz=pz\sigma_p$, yields the homogeneous equation
\begin{equation}\label{hol:eq:thetahom}
 \Theta_p(z)-(1+z)\Theta_p(pz)+pz\Theta_p(p^2z)=0,
\end{equation}
which is also obtained by substituting $pz$ for $z$ in
\eqref{hol:eq:thetafirst} and subtracting the two inhomogeneous identities.
The existence of these identities is not the new part. What matters is
precisely where the series is strongly evaluable, for an arbitrary-rank value
group and an arbitrary Hahn parameter.

\subsection{Exact domain with arbitrary coefficient tails}

\begin{theorem}[Exact partial theta domain]\label{hol:thm:thetadomain}
Let $\Gamma$ be as in Convention~\ref{hol:conv:group} and let $p\in\K^\times$
with $\beta=v(p)>0$, and let $\Delta_\beta$ be as in \eqref{hol:eq:delta}.
Then
\begin{align}
 \Dom(\Theta_p)
 &=\{0\}\cup\{x\in\K^\times:v(x)\ge-N\beta\text{ for some }N\in\N\}
 \label{hol:eq:thetadomain}\\
 &=\{0\}\cup\{x\in\K^\times:(v\bmod\Delta_\beta)(x)\ge0\}.\nonumber
\end{align}
In particular $\Theta_p\in\E$ if and only if $\beta$ is an order unit of
$\Gamma$. No monomial hypothesis on $p$ or on $x$ is needed, and $\Gamma$ need
not be divisible.
\end{theorem}

\begin{proof}
Evaluation at $0$ is finite. For $x\ne0$ put $\delta=v(x)$ and write
\[
 p=c\,t^\beta(1+u),\qquad x=d\,t^\delta(1+w),\qquad c,d\in\C^\times,
\]
where each of $u,w$ is zero or has strictly positive valuation. Let
$S=\supp(u)\cup\supp(w)\subseteq\Gamma_{>0}$ and $M=M(S)$. The $n$th
evaluation term is
\[
 p^{\binom n2}x^n
 =c^{\binom n2}d^{\,n}\,t^{h_n}\,(1+u)^{\binom n2}(1+w)^n,
 \qquad h_n=\binom n2\beta+n\delta,
\]
and the last two factors are finite integer powers, so
\begin{equation}\label{hol:eq:thetasupport}
 \supp\bigl(p^{\binom n2}x^n\bigr)\subseteq h_n+M,
 \qquad v\bigl(p^{\binom n2}x^n\bigr)=h_n.
\end{equation}

Suppose first that $\delta\ge-N\beta$ for some $N\in\N$. Then
\[
 h_{n+1}-h_n=n\beta+\delta\ge(n-N)\beta>0\qquad(n>N),
\]
so $(h_n)$ is eventually strictly increasing. Its set of values is therefore
well ordered, including the finite initial segment, and each value is taken by
finitely many indices. Lemma~\ref{hol:lem:certificate}\,(b) applied to
\eqref{hol:eq:thetasupport} proves strong summability.

If no such $N$ exists, then $\delta<-N\beta$ for every $N$, so
$h_{n+1}-h_n=n\beta+\delta<0$ for every $n$. By
\eqref{hol:eq:thetasupport} the leading exponents are exactly $h_n$, and they
form a strictly decreasing sequence, so strong summability fails by
Lemma~\ref{hol:lem:nonincreasing}.

Finally, $\delta\ge-N\beta$ for some $N$ says exactly that the image of
$\delta$ in $\Gamma/\Delta_\beta$ is nonnegative: a nonnegative coset either
is positive, whence $\delta>0$, or is zero, whence $\delta\in\Delta_\beta$ and
$\delta\ge-N\beta$ for some $N$; a negative coset lies below every element of
$\Delta_\beta$. The last assertion is the case $\Delta_\beta=\Gamma$.
\end{proof}

This proof controls the whole support, not only the leading valuation. Leading
valuation formulas by themselves do not in general determine a Hahn evaluation
domain, because hidden tails can matter; here the common positive monoid $M$
is the extra certificate that makes the formula exact. Note also that the
entire calculation is carried out in $\Gamma$: nothing is divided by an
integer, and no real-valued surrogate for $v$ is introduced.

\begin{remark}[A second, shorter proof of the order-unit half]
\label{hol:rem:thetasecond}
When $\mu=\beta$ is an order unit, entireness also follows from the
cofinal-valuation criterion Lemma~\ref{hol:lem:certificate}\,(a) alone, and
the routes of sources~08 and~09 are both kept here because they are genuinely
different. For $x\ne0$ with $\alpha=v(x)$, the $n$th evaluated term has
valuation
\begin{equation}\label{hol:eq:thetaval}
 \frac{n(n-1)}2\mu+n\alpha.
\end{equation}
Given a target $\beta'\in\Gamma$, choose ordinary nonnegative integers $M',B'$
with $\alpha\ge-M'\mu$ and $\beta'\le B'\mu$, which is possible exactly
because $\mu$ is an order unit. Then \eqref{hol:eq:thetaval} is at least
$\bigl(\tfrac{n(n-1)}2-M'n\bigr)\mu$, which exceeds $B'\mu$ for all
sufficiently large ordinary $n$. So the term valuations tend cofinally to
$+\infty$ and Lemma~\ref{hol:lem:certificate}\,(a) applies. Only leading term
valuations are used, so unrestricted higher Hahn tails in $p$ do not
invalidate the conclusion. This shorter route is available only in the
order-unit case: at the boundary $v(x)=-N\beta$ with $\beta$ noncofinal, the
term valuations increase strictly but stay trapped below a fixed element of
$\Gamma$, so criterion (a) says nothing and the monoid certificate is
essential.
\end{remark}

\begin{remark}[The order-unit case has an optional rank-one picture]
\label{hol:rem:thetacoarse}
When $\beta$ is an order unit, the real-valued coarsening
$\pi:\Gamma\to\R$ of~\cite[Lemma~9.1]{repoentire}, normalized by
$\pi(\beta)=1$, gives the same field $\K$ a rank-one absolute value.
For $x\ne0$, its theta terms have coarsened valuations
$\binom n2+n\pi(v(x))$, tending to $+\infty$.
This is an optional normed-field picture of the order-unit case,
not a map that replaces the Hahn exponents by their images in $\R$.

Such a coefficientwise replacement need not define a Hahn series.
In $\Q H\oplus\Q\epsilon$ with $H$ dominating all $n\epsilon$,
the coarsening normalized by $H$ kills $\epsilon$.
The valid series $\sum_{n\ge0}t^{n\epsilon}$ would then place infinitely
many coefficients at exponent zero. Coarsening a valuation retains
the field and changes the value map; it is not an embedding into a
complex-coefficient Hahn field with smaller exponent group.
No proof here requires this optional coarsening.
\end{remark}

\subsection{Completion of the sharp criterion}

\begin{proof}[Proof of Theorem~\ref{hol:main:q}]
Assume (i) and suppose (ii) fails. If $v(q)=0$, Corollary
\ref{hol:cor:unitrigid} makes every entire solution a polynomial. If
$v(q)\ne0$ and $\abs{v(q)}$ is not an order unit,
Proposition~\ref{hol:prop:noncofinal} does the same. Either way (i) fails,
which proves (i) $\Rightarrow$ (ii) and supplies the stated uniform
obstruction and degree bound.

Assume (ii). If $v(q)>0$ is an order unit take $p=q$: by
Theorem~\ref{hol:thm:thetadomain}, $\Theta_q\in\E\setminus\K[z]$, and
\eqref{hol:eq:thetahom} is a nonzero linear $q$-difference equation with
polynomial coefficients. If $v(q)<0$ and $-v(q)$ is an order unit take
$p=q^{-1}$, so $v(p)=-v(q)>0$ is an order unit and $\Theta_{q^{-1}}\in\E$ by
the same theorem. Substituting $q^2z$ for $z$ in \eqref{hol:eq:thetahom},
with $p=q^{-1}$ so that $pz\mapsto qz$ and $p^2z\mapsto z$, gives
\begin{equation}\label{hol:eq:inverseq}
 \Theta_{q^{-1}}(q^2z)-(1+q^2z)\Theta_{q^{-1}}(qz)
 +qz\,\Theta_{q^{-1}}(z)=0,
\end{equation}
which uses only nonnegative powers of $\sigma_q$ and has polynomial
coefficients that are not all zero. Both signs of $v(q)$ are covered.
\end{proof}

The partial theta function is used as an established type of construction,
not advertised as a newly discovered special function. Its role is precise: it
realizes the positive side of a necessary and sufficient criterion, including
nonmonomial $q$ and arbitrary valuation rank.

\subsection{Two scales in the same rank-two field}

\begin{example}[One field, opposite answers]\label{hol:ex:ranktwo}
Let $\Gamma=\Q e_0\oplus\Q e_1$ with $e_1\gg e_0>0$, ordered by the highest
nonzero coordinate. Then $e_1$ is an order unit and $e_0$ is not.

For $q=t^{e_0}$, no nonpolynomial strongly entire series satisfies a nonzero
linear $q$-difference equation. The series $\Theta_q$ itself is not entire:
by Theorem~\ref{hol:thm:thetadomain} its exact domain is the valuation ring
obtained by forgetting the $e_0$ coordinate, which contains every $x$ with
$v(x)=-Ne_0$ and excludes $x=t^{-e_1}$.

For $q=t^{e_1}$, $\Theta_q$ is strongly entire, though the value group still
has rank two. The two conclusions occur in the \emph{same} Hahn field: the
distinction is the dilation scale, not the rank of the field.
\end{example}

Under $e_0=1$, $e_1=\omega$ and $t^\gamma=\omega^{-\gamma}$, the first dilation
is multiplication by $\omega^{-1}$ and the second by $\omega^{-\omega}$.
Quadratic multiples of the first valuation never reach the second scale;
quadratic multiples of the second dominate every scale in this workspace. A
fuller table for the same group, including unit dilations and the torsion
exception, appears in Section~\ref{hol:sub:tables}.

\subsection{The coarse boundary really can be included}
\label{hol:sub:coarseboundary}

For the theta series the forward recurrence is $a_{n+1}=p^na_n$, read off from
\eqref{hol:eq:thetafirst}, whose inhomogeneity affects only the index $n=0$
and is removed by Proposition~\ref{hol:prop:inhomogeneous}. Its two slots have
valuations $n\beta$ and $0$, both of which project to $0$ modulo
$\Delta_\beta$, so the projected gap is zero and
$\overline B_{\mathrm{rec}}=\overline B^{\flat}_{\mathrm{rec}}=0$.
Theorem~\ref{hol:thm:coarse} therefore excludes the points of negative
coarsened valuation, while Theorem~\ref{hol:thm:thetadomain} includes
\emph{all} points of zero coarsened valuation. The strict inequality in the
coarsened exclusion theorem cannot be replaced by a weak one in general. This
is the exact mirror image of Section~\ref{hol:sub:expdomain}, where the
uncoarsened bound was attained with its boundary excluded from the domain. The
two boundaries behave differently because one comparison is made in $\Gamma$
and the other in a quotient that has forgotten the scale doing the work.

\section{Torsion parameters and why they are excluded}\label{hol:sec:torsion}

If $q^{h}=1$ then $\sigma_q^{h}-1$ annihilates \emph{every} formal power
series, so the unrestricted definition of linear $q$-difference finiteness is
vacuous for such a parameter. This is not a technical inconvenience removable
by a better proof; it is a genuine counterexample to dropping the nontorsion
hypothesis in Theorem~\ref{hol:main:q}.

The difficulty persists even after excluding operators with zero action. For
$h>1$ and
\[
 f(z)=z^rg(z^{h}),
\]
one has $f(qz)=q^rf(z)$, and the operator $\sigma_q-q^r$ is not zero on all
formal series when $h>1$. A nonpolynomial entire $g$ produces a nonpolynomial
entire $f$, since evaluation of $f$ at $x$ is evaluation of $g$ at $x^{h}$
followed by multiplication by $x^r$. Finite-order dilations therefore have a
real periodic-support exception and are not a limiting case of
Theorem~\ref{hol:main:q}. In particular ``the residue of $q$ is a root of
unity'' and ``$q$ is a root of unity'' are fundamentally different
hypotheses: Example~\ref{hol:ex:unitperiod} is covered by the theorem and this
section is not.

There is nevertheless a meaningful eigenspace statement.

\begin{proposition}[Torsion covariance]\label{hol:prop:torsion}
Let $q\in\C^\times$ have exact finite order $m$ and let $0\le r<m$.
\begin{enumerate}[label=\textup{(\alph*)}]
\item A formal series satisfies $f(qz)=q^rf(z)$ if and only if
$f(z)=z^rG(z^m)$ for a unique $G\in\K[[z]]$.
\item If $G\in\E$ then $f=z^rG(z^m)\in\E$.
\item The converse of \textup{(b)} also holds, so
$f\in\E$ if and only if $G\in\E$, without divisibility of $\Gamma$.
\end{enumerate}
Moreover, every root of unity of $\K$ lies in $\C$, so the hypothesis
$q\in\C^\times$ loses no torsion element of the Hahn field. All parts use only Convention~\ref{hol:conv:group}.
\end{proposition}

\begin{proof}
For (a), coefficient comparison gives $(q^n-q^r)a_n=0$, so only degrees
congruent to $r$ modulo $m$ occur, and $G$ is read off uniquely.

For (b), evaluation of $z^rG(z^m)$ at $x$ is $x^r$ times the family for $G$ at
$x^m$, reindexed; multiplication by the fixed element $x^r$ preserves strong
summability by Lemma~\ref{hol:lem:support}\,(3), and at $x=0$ the evaluation
is finite.

For (c), suppose $f\in\E$ and let $y\in\K^\times$.
Set $\delta=v(y)$, $\gamma=\min(\delta,0)$, and $x=t^\gamma$.
Since $m\ge1$, we have $v(x^m)=m\gamma\le\delta$.
Strong evaluation of $f$ at $x$, restricted to the degrees $r+mn$
and multiplied by $x^{-r}$, gives strong evaluation of $G$ at $x^m$.
Lemma~\ref{hol:lem:inward} then gives strong evaluation at $y$.
At $y=0$ evaluation is finite. No root of $y$ is extracted.

Finally, let $\zeta\in\K$ be a root of unity. Then $v(\zeta)=0$; dividing by
$\res(\zeta)$ gives $1+u$ with $u$ zero or of positive valuation, and
$(1+u)^{m'}=1$ for some $m'\ge1$. If $u\ne0$ then $(1+u)^{m'}-1$ has valuation
$v(u)$ and leading coefficient $m'\lc(u)\ne0$ in characteristic zero, a
contradiction. So $u=0$ and $\zeta=\res(\zeta)\in\C$.
\end{proof}

The converse uses only that the values $m\gamma$ are coinitial in $\Gamma$,
as the explicit choice $\gamma=\min(\delta,0)$ shows. Actual $m$th roots
need not exist in the original Hahn field; this stronger algebraic
requirement is unnecessary for the entireness equivalence.

\section{An extension to mixed differential--dilation equations}
\label{hol:sec:mixed}

For a nonzero noncofinal dilation valuation, derivatives do not repair the
missing-scale problem. The extension is useful when a symbolic system allows
both differentiation and rescaling of the argument.

\begin{lemma}[Polynomial dependence on the index as well as the dilation]
\label{hol:lem:bivariate}
Let $q\in\K^\times$ with $\lambda=v(q)\ne0$ and let
$0\ne P(U,X)\in\K[U,X]$. There are $\beta\in\Gamma$, $k\in\N$ and $N\in\N$
with
\[
 v(P(n,q^n))=\beta+kn\lambda\qquad(n\ge N).
\]
\end{lemma}

\begin{proof}
Write $P(U,X)=\sum_kC_k(U)X^k$. By Lemma~\ref{hol:lem:integer} each nonzero
$C_k(n)$ has the fixed valuation $\beta_k=\beta(C_k)$ for all sufficiently
large $n$. The finitely many summands then have valuations
$\beta_k+kn\lambda$, and the pairwise-order argument of
Lemma~\ref{hol:lem:affine} produces a unique eventual minimizing index, whose
term cannot cancel.
\end{proof}

\begin{theorem}[Mixed rigidity at a noncofinal nonzero scale]
\label{hol:thm:mixed}
Let $\Gamma$ be as in Convention~\ref{hol:conv:group}, let $q\in\K^\times$
with $v(q)\ne0$, and suppose $\abs{v(q)}$ is not an order unit of $\Gamma$. If
a strongly entire $f$ satisfies a nonzero equation
\begin{equation}\label{hol:eq:mixed}
 \sum_{r,\ell}p_{r,\ell}(z)D_z^r\sigma_q^\ell f(z)=0,\qquad p_{r,\ell}\in\K[z],
\end{equation}
where the sum is finite and the coefficients are not all zero, then $f$ is a
polynomial. Each fixed equation has a uniform exterior obstruction and a
uniform degree bound.
\end{theorem}

\begin{proof}
Expand $p_{r,\ell}(z)=\sum_kc_{r,\ell,k}z^k$. With $s=r-k$, the term of
coefficient index $n+s$ contributes to degree $n$ the quantity
$c_{r,\ell,k}\fall{n+s}{r}q^{\ell(n+s)}a_{n+s}$, so the coefficient recurrence
has shift polynomials
\begin{equation}\label{hol:eq:mixedshift}
 P_s(U,X)=\sum_{r-k=s}\sum_\ell c_{r,\ell,k}\,q^{\ell s}\fall{U+s}{r}X^\ell,
\end{equation}
evaluated at $(U,X)=(n,q^n)$. For fixed $s$ and $\ell$, distinct $r$ give
distinct degrees in $U$, so a nonempty group of nonzero coefficients cannot
cancel identically and some $P_s$ is nonzero. Reindex from the smallest active
shift; replacing $U$ by $U-s_0$ and $X$ by $q^{-s_0}X$ preserves polynomiality
and nonzeroness, so the leftmost coefficient is $A_0(n)=P'_0(n,q^n)$ for a
nonzero bivariate $P'_0$, and Lemma~\ref{hol:lem:bivariate} makes it
eventually nonzero.

Lemma~\ref{hol:lem:bivariate} gives eventual affine valuations for all the
nonzero coefficients, so their differences obey a bound $B+Mn\abs{v(q)}$
exactly as in \eqref{hol:eq:linearbound}. Since $\abs{v(q)}$ is not an order
unit, \eqref{hol:eq:dominates} bounds this by one element of $\Gamma$, and
Corollary~\ref{hol:cor:boundedcost} finishes the proof.
\end{proof}

The hypothesis $v(q)\ne0$ is essential to \emph{this} proof. For $v(q)=0$ the
separate pure differential and pure dilation conclusions have been proved
above, but they are not a proof of the mixed conclusion. When $\Gamma$ has no
order unit, the mixed conclusion at unit valuation, for nonlinear relations
as well, is Theorem~\ref{hol:cf:main:entire}\,(b) of source~12, proved by a
different method. For every $\Gamma$, with or without an order unit, the
linear mixed conclusion at unit valuation is
Corollary~\ref{hol:ud:cor:unitrigidity} (source~18), which answers
Question~\ref{hol:q:mixedunit}: Lemma~\ref{hol:lem:bivariate} is replaced there
by the eventually periodic valuations of $P(n,q^n)$
(Corollary~\ref{hol:ud:cor:bivariate}), and the theorem above is a special
case of Theorem~\ref{hol:ud:thm:rigidity}.

\section{Surreal workspaces that separate the theories}
\label{hol:sec:surreal}

The results so far are stated in set-sized Hahn fields, which keeps the
foundations transparent and allows a literal surreal reading without
quantifying over a proper class as though it were a set.

\subsection{The workspace without an order unit}

Take
\begin{equation}\label{hol:eq:gammainfty}
 \Gamma_\infty=\bigoplus_{j\ge0}\Q e_j,
 \qquad 0<e_0\ll e_1\ll e_2\ll\cdots,
\end{equation}
ordered so that the highest index with a nonzero coordinate determines the
sign, equivalently $e_{j+1}>Ne_j>0$ for all $j\ge0$ and $N\ge1$. Every element
has only finitely many nonzero coordinates. One may identify $e_j$ with the
surreal $\omega^j$.

The group has no order unit: if $\gamma>0$ has largest active index $j$, then
$e_{j+1}>n\gamma$ for every ordinary $n$. The sequence $(e_n)$ is nevertheless
cofinal in the whole group, so this is a countable cofinality that cannot be
generated by repeatedly adding one positive element. This group and this
distinction appear in the companion entire-function
report~\cite{repoentire} and are background here, not a new observation.

\begin{corollary}[Universal nontorsion dilation rigidity]
\label{hol:cor:universal}
Suppose $\Gamma$ has no order unit. Then every nonpolynomial $f\in\E$
satisfies, simultaneously for all nontorsion $q\in\K^\times$ and all $r\ge0$,
\[
 f(z),f(qz),\ldots,f(q^rz)\quad\text{are linearly independent over }\K(z).
\]
Such fields can contain nonpolynomial strongly entire series, as
\eqref{hol:eq:gammainfty} and Theorem~\ref{hol:thm:explicit} show.
\end{corollary}

\begin{proof}
A linear dependence over $\K(z)$ becomes, after clearing denominators, a
nonzero equation \eqref{hol:eq:qlinear}. Since $\Gamma$ has no order unit,
no $\abs{v(q)}$ is an order unit, so Theorem~\ref{hol:main:q} forces $f$ to be
a polynomial.
\end{proof}

\begin{theorem}[An explicit universally nonholonomic entire series]
\label{hol:thm:explicit}
In $K_{\Gamma_\infty}$ the series
\begin{equation}\label{hol:eq:explicitF}
 F(z)=\sum_{n\ge0}t^{e_n}z^n
     =\sum_{n\ge0}\omega^{-\omega^n}z^n
\end{equation}
is strongly entire and nonpolynomial. It satisfies no nonzero linear
differential equation over $K_{\Gamma_\infty}(z)$ and, for every nontorsion
$q\in K_{\Gamma_\infty}^\times$, no nonzero polynomial-coefficient linear
$q$-difference equation. For every $q$ of nonzero valuation it satisfies no
nonzero mixed equation \eqref{hol:eq:mixed} either.
\end{theorem}

\begin{proof}
Fix $x\ne0$ and write $x=c\,t^\delta(1+w)$ with $c\in\C^\times$ and $w$ zero
or of positive valuation. Choose $M$ with $\delta$ involving only the
coordinates $e_0,\ldots,e_M$. The $n$th evaluation term has leading exponent
$h_n=e_n+n\delta$ and support contained in $h_n+M(\supp w)$. For $n>M$ the
highest coordinate of $h_n$ is $e_n$ with coefficient $1$, so $h_n>0$, and the
highest coordinate of
\[
 h_{n+1}-h_n=e_{n+1}-e_n+\delta
\]
is $e_{n+1}$ with coefficient $1$, so $(h_n)$ is eventually strictly
increasing. Lemma~\ref{hol:lem:certificate}\,(b) proves strong summability; at
$x=0$ only the constant term remains. All coefficients are nonzero, so $F$ is
nonpolynomial.

Alternatively, for any $\beta\in\Gamma_\infty$ choose $j$ larger than every
active coordinate index of $\delta$ and of $\beta$; for $n>j$ the largest
active coordinate of $h_n-\beta$ is $n$ with coefficient $1$, so $h_n>\beta$,
and Lemma~\ref{hol:lem:certificate}\,(a) applies for arbitrary $w$.
Cofinal leading valuations control the whole support, since every exponent
of a Hahn term is at least its leading exponent.

The differential conclusion is Theorem~\ref{hol:main:ode}. Since
$\Gamma_\infty$ has no order unit, Theorem~\ref{hol:main:q} --- equivalently
Corollary~\ref{hol:cor:universal} --- excludes every nontorsion dilation
equation, and Theorem~\ref{hol:thm:mixed} excludes the mixed equations at
every $q$ with $v(q)\ne0$. Rational coefficient equations reduce to the
polynomial-coefficient case by clearing denominators.
\end{proof}

The universal quantifier over $q$ matters. The parameter may be chosen after
$F$ is fixed and may itself have infinite support; each individual $v(q)$
still has finitely many coordinates in $\Gamma_\infty$, hence is noncofinal.
So the theorem excludes a linear relation for \emph{each} such parameter, not
merely for an advance list of simple dilations. The series has a simple
explicit coefficient rule, and this is a statement about finite linear
annihilating equations, not a claim that the rule is uncomputable or lacks a
finite description.

\subsection{Why the same formula is not entire on all surcomplex numbers}
\label{hol:sub:notallsurcomplex}

Let $e_\ast=\omega^\omega$, which dominates every integer multiple of every
positive element of $\Gamma_\infty$, and enlarge the group to
$\Gamma_\infty+\Q e_\ast$. At the surcomplex point $x=\omega^{e_\ast}
=t^{-e_\ast}$ the evaluated leading exponents are $e_n-ne_\ast$, whose
successive differences $e_{n+1}-e_n-e_\ast$ are negative. The defining family
is therefore not strongly summable, by
Lemma~\ref{hol:lem:nonincreasing}.

This explicit obstruction prevents the fixed-workspace theorem from being
misread as an assertion of a nonpolynomial power series entire on all of
$\No[i]$. The example illustrates two independent restrictions: a workspace
can be large enough that no single dilation scale is cofinal, yet small enough
to admit a countable sequence of coefficient scales escaping all its
arguments, and enlarging it further can destroy the very entireness on which
the theorem operates. The companion report studies this kind of boundary in
general~\cite{repoentire}; the calculation here is included only to state the
domain of the present example unambiguously, and no new universal
scalar-extension theorem is claimed.

\subsection{Two classification tables}\label{hol:sub:tables}

Let $\Gamma=\Q H\oplus\Q\epsilon$ with $H$ dominating every $n\epsilon$, so
that $H$ is an order unit and $\epsilon$ is not.

\begin{center}
\begin{tabularx}{\textwidth}{@{}l l X@{}}
\toprule
Dilation $q$ & Its valuation & Nonpolynomial entire solution of a nonzero
linear $q$-difference equation?\\
\midrule
$t^H$ & $H$ & Yes: $\Theta_{t^H}$, by \eqref{hol:eq:thetahom}.\\
$t^{-H}$ & $-H$ & Yes: $\Theta_{t^{H}}$, by \eqref{hol:eq:inverseq}.\\
$t^\epsilon$ & $\epsilon$ & No: positive, but noncofinal.\\
$1+t^H$ & $0$ & No: an infinite-order unit.\\
$2$ & $0$ & No: complex size does not change the Hahn valuation.\\
$-1$ & $0$ & Outside the theorem: a torsion parameter, with
periodic-support exceptions.\\
\bottomrule
\end{tabularx}
\end{center}

\begin{center}
\small
\begin{tabularx}{\textwidth}{@{}l l l X@{}}
\toprule
Workspace group & Parameter & Solution exists? & Reason\\
\midrule
$\Q$ & $v(q)\ne0$ & Yes & Every positive nonzero scale is an order unit.\\
$\Q$ & $v(q)=0$, nontorsion & No & Unit-orbit valuations are eventually
periodic.\\
$\Q e_0\oplus\Q e_1$, $e_1\gg e_0$ & $v(q)=e_0$ & No & A lower scale is not
cofinal.\\
the same group & $v(q)=e_1$ & Yes & The top scale is an order unit.\\
$\bigoplus_{n\ge0}\Q e_n$ & any nontorsion $q$ & No & No scale is an order
unit, although $F$ is entire.\\
\bottomrule
\end{tabularx}
\end{center}

In every row, nonpolynomial strongly entire D-finite series are impossible.
Nonpolynomial strongly entire differentially algebraic series exist exactly in
the rows whose group has an order unit, that is, in all rows except the last
row of the second table (Theorem~\ref{hol:td:main:boundary}), and their least
differential order is then three (Theorem~\ref{hol:ot:main:second}); this does
not depend on the parameter. Allowing derivatives beside the powers of
$\sigma_q$ changes no answer in either table: a nonzero mixed equation in
$D_z$ and powers of $\sigma_q$ has a nonpolynomial strongly entire solution
exactly in the rows answered ``yes'' (Corollary~\ref{hol:ud:cor:cyclic},
source~18); the torsion row stays outside that corollary, and there
$\sigma_{-1}-1$ already has nonpolynomial entire solutions
(Section~\ref{hol:sec:torsion}).
Torsion parameters are deliberately absent from the second table: their
annihilating operator may be an identity on every series. Finite valuation
rank alone is therefore not the criterion, and the distinction persists for
nonmonomial dilations, because the theorem uses the valuation rather than a
requirement that $q$ be a monomial.

\begin{corollary}[Detecting order units by functional equations]
\label{hol:cor:detect}
For a nonzero set-sized ordered abelian group $\Gamma$ the following are
equivalent:
\begin{enumerate}[label=\textup{(\roman*)}]
\item $\Gamma$ has an order unit;
\item some nontorsion $q\in\K^\times$ admits a nonpolynomial strongly entire
solution of a nonzero polynomial-coefficient linear $q$-difference equation.
\end{enumerate}
\end{corollary}

\begin{proof}
Theorem~\ref{hol:main:q} gives (ii) $\Rightarrow$ (i). For the converse take
$q=t^\mu$ for an order unit $\mu$; then $v(q)=\mu$ is an order unit and $q$ is
nontorsion.
\end{proof}

This recovers a property of the ordered group from a global functional
equation existence question. Theorem~\ref{hol:td:main:boundary} recovers the
same property from algebraic differential equations, with no dilation. The
workspace \eqref{hol:eq:gammainfty} shows why
the existence of nonpolynomial entire functions is by itself strictly weaker
than the existence of an order unit.

\subsection{Several meanings of transcendence, kept apart}

For the explicit $F$ the established conclusions are: transcendence over
$K_{\Gamma_\infty}(z)$ by Corollary~\ref{hol:cor:algebraic}; linear
differential independence of every finite family of derivatives; and linear
independence along every single nontorsion geometric dilation orbit. These are
substantial restrictions, and they do \emph{not} imply that no polynomial
relation $P(z,F,D_zF,\ldots,D_z^RF)=0$ exists. Nonlinear differential
equations lie outside these proofs, and a relation involving $F(q_1z)$ and
$F(q_2z)$ for multiplicatively independent parameters is not covered merely by
applying the single-$q$ theorem twice.

The nonlinear part adds one more conclusion by a separate proof: $F$ and
$D_zF$ are algebraically independent over $K_{\Gamma_\infty}(z)$
(Corollary~\ref{hol:nl:cor:infiniterank}), so every nonlinear relation with
$R\le1$ is excluded. The methods of that part exclude relations with $R\ge2$
only when their residue corner polynomial is nonzero, and give no statement
about several dilations. Because $\Gamma_\infty$ has no order unit, the
coefficient-field part removes both limits for this $F$:
Theorem~\ref{hol:cf:thm:Fomega} (source~12) makes all $(D_z^jF)(\lambda z)$,
for dilations $\lambda$ with pairwise quotients that are not roots of unity,
algebraically independent over $K_{\Gamma_\infty}(z)$. In particular $F$
satisfies no algebraic differential equation of any order, and no nonzero
relation, linear or not, links $F(q_1z)$ and $F(q_2z)$ for multiplicatively
independent $q_1,q_2$ (Corollary~\ref{hol:cf:cor:several}, a merge
observation).

\section{Coefficient fields, degenerate value groups, characteristic}
\label{hol:sec:fields}

The one-variable proofs work over any characteristic-zero coefficient field
with trivial valuation: the root-of-unity argument uses only finitely many
ordinary polynomial roots and needs no algebraic closedness, and the
integer-evaluation lemma needs only residue characteristic zero. All theorems
also hold over $\R((t^\Gamma))$, the generic-line proof of
Theorem~\ref{hol:main:multi} using uncountability of $\R$ in place of
uncountability of $\C$. The results therefore concern surreal as well as
surcomplex workspaces. More generally the core versions work over
$k((t^\Gamma))$ for any characteristic-zero field $k$, and, as
Section~\ref{hol:sub:divisibility} records, without divisibility of $\Gamma$.

No statement here is a theorem for arbitrary valued fields of positive residue
characteristic. In that setting Lemma~\ref{hol:lem:integer} fails and the
valuations of near-torsion geometric powers can behave differently, so any
comparison with the ultrametric literature must keep the coefficient field and
the residue characteristic visible. Example~\ref{hol:nl:ex:charp}, from
source~10, shows that characteristic zero is needed even for the single
linear equation $D_zf=0$: over $\mathbb F_p((t^\Q))$ a nonpolynomial strongly
entire series has zero derivative.

The nonlinear part (Convention~\ref{hol:nl:conv:field}) is proved by its
source over an arbitrary field of characteristic zero. That is a statement
about Sections~\ref{hol:nl:sec:setting}--\ref{hol:nl:sec:consequences} only;
the generality of the linear theorems is governed by the first paragraph of
this section and its own argument. The same holds for the coefficient-field
part (Convention~\ref{hol:cf:conv:field}), whose sources also work over an
arbitrary field of characteristic zero.

We assumed $\Gamma\ne0$ in order to speak of positive exterior scales. When
$\Gamma=0$ the field is $\C$ with the trivial valuation, strong evaluation at
$x=1$ already permits only finitely many nonzero coefficients because every
nonzero term has support $\{0\}$, and entire series are polynomials for that
immediate support reason.

\section{Quantitative certificates and symbolic representation}
\label{hol:sec:certificates}

\subsection{What one can extract from an annihilating equation}

The classification has a direct representational consequence. A system whose
nonpolynomial special functions are represented only by finite linear
differential equations cannot cover the nonpolynomial strong entire functions
of a Hahn workspace. Allowing linear dilation equations enlarges the class only
when a cofinal dilation scale is available. In the no-order-unit example, an
exact coefficient rule such as \eqref{hol:eq:explicitF} lies outside both pure
equation formats.

This does not rule out exact symbolic computation with such functions. It
specifies what a faithful representation must keep separate: the
coefficient-generation rule, the exponent group, the strong evaluation domain,
and the type of finite equation being asserted. A field-independent label such
as ``entire theta function'' discards relevant information; for $\Theta_p$ the
proof certificate is not merely $v(p)>0$, it includes that $v(p)$ is an order
unit of the selected group.

For a differential operator the reduction yields a concrete certificate.
Compute the finite shift polynomials $Q_s$ of \eqref{hol:eq:odecoeff} and
shift the least active index to zero. For each resulting $P_j$ record its
least coefficient valuation $\beta_j$ and the corresponding nonzero complex
residue polynomial $\overline{P_j}$. Record an integer $N$ beyond their
nonnegative integer roots and beyond the recurrence's starting index. Record
the threshold $B$ of Proposition~\ref{hol:prop:precursive}, or
\eqref{hol:eq:refinedthreshold} when $\Gamma$ is divisible. The escape
statements then certify that every formal solution strongly evaluable at some
$x\ne0$ with $v(x)\le-B$ has degree less than $N$.

For a valuation-zero dilation the same plan uses \eqref{hol:eq:qrec}: if
$\res(q)$ has infinite order use the minimum coefficient valuation, and if it
has order $h$ use \eqref{hol:eq:unitexpand}--\eqref{hol:eq:residuepoly} on
each residue class. For a nonzero-valuation dilation, finitely many affine
comparisons select the eventual winner of each $R_h(q^n)$; project their
constants modulo $\Delta_{\abs{v(q)}}$ to obtain the coarsened bound
\eqref{hol:eq:coarsebound0}, or \eqref{hol:eq:coarsebound} when $\Gamma$ is
divisible. When the parameter scale is noncofinal one can instead display
$\gamma\gg\abs{v(q)}$ and the explicit bad argument $t^{-(B+\gamma)}$ of
Proposition~\ref{hol:prop:noncofinal}.

These are mathematical certificate schemas relative to exact coefficient data
and order comparisons. They are \emph{not} a terminating algorithm for
arbitrary encoded surreals. Deciding whether an arbitrary complex residue is a
root of unity, comparing arbitrary surreal values, extracting arbitrary
supports, or deciding whether an abstract scale is cofinal can require
information not effectively available in a chosen representation, and none of
this has been made computable by the proofs.

\subsection{A valuation recurrence with zero coefficients allowed}

\begin{example}[Jump lengths, not just jumps]\label{hol:ex:sparse}
Suppose a sequence that is not eventually zero satisfies
\[
 t^Aa_n+(n-2)t^Ba_{n+1}+t^Ca_{n+3}=0\qquad(n\ge0),
\]
with $A,B,C\in\Gamma$. The middle coefficient vanishes at $n=2$, a finite
exceptional index; beyond it the three valuations are $A,B,C$. The escape
chain may move forward by one or by three, and no assumption that all $a_n$
are nonzero is needed. Corollary~\ref{hol:cor:boundedcost} excludes the region
$v(x)\le-\max\{A-B,\,A-C,\,0\}$ under
Convention~\ref{hol:conv:group} alone. \emph{If $\Gamma$ is divisible},
Corollary~\ref{hol:cor:periodic} gives the sharper excluded region
\[
 v(x)\le-\max\Bigl\{A-B,\ \frac{A-C}{3}\Bigr\}.
\]
The maximum must be divided by the actual jump length: replacing every jump by
one need not preserve the intended threshold, especially when the valuation
gaps have different signs. This is a worked instance of item (1) of
Section~\ref{hol:sub:divisibility}.
\end{example}

\subsection{Current Lean coverage and remaining dependency order}

The ledger \texttt{docs/FORMALIZATION.md} records the formalized escape
mechanism in \path{Surreal/HahnSeries/EscapeChain.lean}:
the nonincreasing-order obstruction, escape step and infinite chain,
evaluation exclusion, and the uniform bounded-cost corollary. It also covers
the univariate strong-domain definitions and the formal exponential
identity with its exact domain and attained exclusion boundary.
Part~(2) of Lemma~\ref{hol:lem:support} is mapped separately to
\path{Surreal/HahnSeries/NeumannWords.lean}.

\path{Surreal/HahnSeries/PartialThetaDomain.lean} covers both sufficient
criteria of Lemma~\ref{hol:lem:certificate}, the scale-subgroup facts,
the exact partial-theta domain and the theta identities. It also proves
\textup{(ii)}$\Rightarrow$\textup{(i)} of Theorem~\ref{hol:main:q}
and \textup{(i)}$\Rightarrow$\textup{(ii)} of
Corollary~\ref{hol:cor:detect}; the converse implications remain pending.

\path{Surreal/HahnSeries/UnitOrbitPolynomials.lean} covers integer
valuation stabilization, eventually affine geometric-orbit valuations,
the unit-orbit theorem and the bivariate index--dilation lemma.
\path{Surreal/Algebra/GenericConstantPoint.lean} proves
Lemma~\ref{hol:lem:generic}, including its least-exponent reduction, for
uncountable constants. \path{Surreal/HahnSeries/TorsionCovariance.lean}
covers inward stability, all of Proposition~\ref{hol:prop:torsion},
and the preservation half of Lemma~\ref{hol:lem:support}\,(3).
These results require no divisibility of the value group.

The remaining dependencies include the refined periodic threshold,
part~(1) and the reflection half of part~(3) of
Lemma~\ref{hol:lem:support}, finite-block grouping in
Lemma~\ref{hol:lem:certificate}, Proposition~\ref{hol:prop:closure},
and the several-variable strong-domain definition. The operator-to-recurrence
conversions and full differential and dilation classifications remain
pending. For the several-variable classification, constant-point avoidance
is proved; generic-line restriction, restriction of a rational
differential system, and the resulting D-finiteness argument are pending.

The nonlinear rigidity theorems in
Sections~\ref{hol:nl:sec:setting}--\ref{hol:nl:sec:consequences},
the coefficient-field theorem package in
Sections~\ref{hol:cf:sec:setting}--\ref{hol:cf:sec:families}, the
theta-descent part in
Sections~\ref{hol:td:sec:setting}--\ref{hol:td:sec:boundary}, the
order-three-threshold part in
Sections~\ref{hol:ot:sec:setting}--\ref{hol:ot:sec:theta}, and the
unit-dilation, Newton-rigidity, theta-hierarchy, single-scale,
polynomial-composition and recurrence parts in
Sections~\ref{hol:ud:sec:setting}--\ref{hol:rc:sec:surreal}, whose
suggested formalization orders are recorded in
Sections~\ref{hol:ud:sub:certificates}, \ref{hol:th:sub:third}
and~\ref{hol:pc:sub:boundaries} and
Questions~\ref{hol:nr:q:formal}, \ref{hol:fh:q:formal}
and~\ref{hol:ud:q:formal} (with its batch-33 status),
remain unformalized; the shared support and inward-stability prerequisites
above do not prove them. For the coefficient-field package source~12 suggests
the order:
valuation-rank inequality, proper-convex-subgroup obstruction, formal
coefficient recurrence, all-order differential corollary; the mixed
statements would also need a formal theorem implying the zero-set lemma
for nondegenerate exponential polynomials. For the order-three-threshold part
source~14 suggests two nearly independent layers: finite algebra
(dehomogenization, the logarithmic Euler identities, constants of rational
differentiation in characteristic zero and the two endpoint evaluations,
which give the finite endpoint lemma without surreal infrastructure), then
the Hahn layer (the all-scale criterion, transport through the divisible
hull, finite active-index sets and the estimate
\eqref{hol:ot:eq:gapbound}). The build and axiom audit cover
only imported declarations recorded in the ledger. The Python checks are
finite checks of formulas and examples; they do not establish any remaining
infinite statement.

%======================================================================
% Nonlinear part: source 10 (nonlinear rigidity). Labels use hol:nl:.
%======================================================================
\section{Nonlinear equations: finite initial polynomials}
\label{hol:nl:sec:setting}

This section and the next three contain the material of source~10. They
answer the first-order case of Question~\ref{hol:q:nonlinear}, which
sources~08 and~09 had left open, together with a nondegenerate higher-order
class and positive-weight Euler equations in several variables. The method is
not the escape chain of Section~\ref{hol:sec:escape}. The warning in
Remark~\ref{hol:rem:transcendence} stands: a nonlinear equation gives
coefficient convolutions, not a bounded-step linear recurrence. Instead, a
whole scaled series is reduced to a finite polynomial at one large scale, and
an exact leading-term computation gives the contradiction.

\subsection{Coefficient field, conventions and the merge's choices}
\label{hol:nl:sub:conventions}

\begin{convention}[Coefficient field of the nonlinear part]
\label{hol:nl:conv:field}
In Sections~\ref{hol:nl:sec:setting}--\ref{hol:nl:sec:consequences}, $k$ is an
arbitrary field of characteristic zero with the trivial valuation, $\Gamma$ is
as in Convention~\ref{hol:conv:group}, and
\[
 \KK=k((t^\Gamma)).
\]
The valuation $v$, leading coefficient $\lc$, residue $\res$, strong
summability, $\Dom(f)$ and entireness are defined exactly as in
Section~\ref{hol:sub:conventions} and Definition~\ref{hol:def:entire}, with
$k$ in place of $\C$, except that, following source~10, $\res(a)$ is defined
for every $a$ of nonnegative valuation as its coefficient at exponent zero, so
that $\res(a)=0$ when $v(a)>0$. Write $\EK$ for the strongly entire series in
$\KK[[z]]$. In several variables, strong entireness of
$f=\sum_\alpha a_\alpha x^\alpha$ means strong summability of the \emph{joint}
family $(a_\alpha b^\alpha)_{\alpha\in\N^m}$ at every $b\in\KK^m$. For $k=\C$
one has $\KK=\K$ and $\EK=\E$. The derivative $D_z$ kills $\KK$, and
$f'=D_zf$; no scalar surreal derivation, analytic continuation or
fine-topology limit is involved.
\end{convention}

\paragraph{Why the coefficient field differs.}
Source~10 proves every statement of the nonlinear part over an arbitrary
characteristic-zero field~$k$, and it uses neither algebraic nor real
closedness. Its hypothesis on the coefficient field is therefore weaker than
the hypothesis $k=\C$ under which
Theorems~\ref{hol:main:ode}--\ref{hol:main:multi} and the rest of
Sections~\ref{hol:sec:intro}--\ref{hol:sec:certificates} are stated. This
generality is recorded for the nonlinear part only. It is \emph{not}
transferred to the linear theorems: they remain stated over~$\C$, and their
possible extension to other coefficient fields is the separate remark of
Section~\ref{hol:sec:fields}, resting on its own argument. In the other
direction, the nonlinear part borrows from Section~\ref{hol:sec:support} only
facts about well-ordered subsets of~$\Gamma$ and finite coefficient sums
(Lemmas~\ref{hol:lem:support} and~\ref{hol:lem:certificate}), whose proofs do
not use any property of~$\C$; source~10 proves the facts it needs over~$k$
itself, except the positive-support lemma used in
Theorem~\ref{hol:nl:thm:nouniform}, which it imports as standard. The group $\Gamma$ is not assumed divisible, and no order unit, no
countable cofinal subset and no real-valued norm is used anywhere in the
nonlinear part.

\paragraph{Renamed symbols.}
Several symbols of source~10 would collide with symbols of
Sections~\ref{hol:sec:intro}--\ref{hol:sec:certificates}; they are renamed
here, in the new material only.
\begin{center}
\small
\begin{tabularx}{\textwidth}{@{}l l X@{}}
\toprule
Source~10 & Here & Reason\\
\midrule
$D=\max_\nu d_\nu$, $W$ & $d_P$, $w_P$ & $D_z$ is the derivative.\\
$H_P$ (residue corner polynomial) & $\chi_P$ & $H$ is an element of $\Gamma$
in Sections~\ref{hol:sec:support}--\ref{hol:sec:surreal}.\\
$B_P$ (degree cutoff, in $\Q$) & $\kappa_P$ & $B$, $B_L$, $B_{\mathrm{rec}}$
are elements of $\Gamma$.\\
$B$ (set of right sides) & $\mathcal R$ & The same.\\
$r$, $r_0$ (scale, ``radius'') & $\delta$, $\delta_0$ & $r$ indexes
derivative orders in \eqref{hol:eq:dlinear}; $t^{-\delta}$ is the notation
of Corollary~\ref{hol:cor:boundedcost}.\\
$h_r$, $p_r$ (Gauss value, initial polynomial) & $\gamma_\delta$,
$\phi_\delta$ & $h$ is a torsion order, $p_r$ an operator coefficient.\\
$M$ (an integer bound) & $n_\ast$ & $M(S)$ is a support monoid.\\
$E_q$ (Euler operator) & $\vartheta_q$ & $E_\eta$ is the formal exponential
of Proposition~\ref{hol:prop:expdomain}.\\
Theorems~A, B, C & Theorems~\ref{hol:nl:main:first}, \ref{hol:nl:thm:euler},
\ref{hol:nl:thm:corner} & Theorems~A--C of this report are different
theorems.\\
\bottomrule
\end{tabularx}
\end{center}
The shipped audit files \path{10-nonlinear-rigidity-PROOF_AUDIT.md} and
\path{10-nonlinear-rigidity-SOURCES_AND_SCOPE.md} keep the source's
notation and numbering.

\paragraph{Words with two meanings.}
\emph{Escape} keeps the meaning of Section~\ref{hol:sec:escape} (a chain
through a linear recurrence); source~10's lemma ``escape of active degree''
is a different statement and is called Lemma~\ref{hol:nl:lem:active} here.
An \emph{exclusion bound} is, as in Section~\ref{hol:sec:escape}, an element
$\delta_0\in\Gamma$ such that evaluation fails at the arguments $t^{-\delta}$
with $\delta\ge\delta_0$; the source's word ``radius'' is not used. A
\emph{corner} is always the finite set of terms of
Section~\ref{hol:nl:sub:corner}, never a full differential Newton polygon.
The \emph{residue polynomials} $\overline P$ and $R_r$ of
Section~\ref{hol:sec:orbits} are unrelated to the residue corner polynomial
$\chi_P$.

\paragraph{Results printed once.}
Source~10 reproves several facts that Sections~\ref{hol:sec:support}
and~\ref{hol:sec:surreal} already contain. They are printed once and cited:
the support calculus inside Lemma~\ref{hol:nl:lem:algebra}
(Lemma~\ref{hol:lem:support}\,(1),(3) and Proposition~\ref{hol:prop:closure});
its ``cofinal escape of coefficient minima'' lemma, which is
Lemma~\ref{hol:lem:certificate}\,(a); the positive-support geometric-series
lemma, which source~10 imports and Lemma~\ref{hol:lem:support}\,(2) proves;
the surreal embedding, which is Section~\ref{hol:sub:embedding} with the same
monomial convention $t^\gamma\mapsto\omega^{-\gamma}$; strong entireness of
the series \eqref{hol:eq:explicitF}, which is Theorem~\ref{hol:thm:explicit};
and the remark on $\Gamma=0$, which is in Section~\ref{hol:sec:fields}.
Everything else in source~10 is printed in
Sections~\ref{hol:nl:sec:setting}--\ref{hol:nl:sec:consequences}, with its
proofs.

\subsection{The coefficient-family algebra}\label{hol:nl:sub:algebra}

For a finite tuple $X=(X_1,\ldots,X_m)$ of indeterminates put
\[
 \Sfam_m=\Bigl\{F=\sum_\alpha b_\alpha X^\alpha\in\KK[[X]]:
 (b_\alpha)_\alpha\text{ is strongly summable}\Bigr\},
\]
and let $\Sfam_m^+$ be the subset whose coefficients all have nonnegative
valuation. The condition concerns the family of coefficients, not convergence
in the $X$-adic topology.

\begin{lemma}[Hahn-family algebra and its residue]\label{hol:nl:lem:algebra}
$\Sfam_m$ is a $\KK$-subalgebra of $\KK[[X]]$, stable under the formal partial
derivatives and under the Euler operators $\sum_\nu q_\nu X_\nu\partial_\nu$
with $q_\nu\in k$. Coefficientwise residue is a ring homomorphism
\[
 \rho:\Sfam_m^+\longrightarrow k[X],\qquad
 \rho\Bigl(\sum_\alpha b_\alpha X^\alpha\Bigr)=\sum_\alpha\res(b_\alpha)X^\alpha,
\]
whose values are polynomials, not arbitrary formal power series; $\rho$
commutes with these differential operators.
\end{lemma}

\begin{proof}
Closure under sums, scalar multiples and products is the support
calculation of Lemma~\ref{hol:lem:support}\,(1) and~(3), as in the proof of
Proposition~\ref{hol:prop:closure}. For strongly summable $(b_\alpha)$ and
$(c_\epsilon)$ a fixed exponent has finitely many representations as a sum of
two support exponents, each carried by finitely many indices, so the double
family $(b_\alpha c_\epsilon)$ is strongly summable, and regrouping by
$\alpha+\epsilon$ gives the Cauchy product: regrouping can remove supported
exponents by cancellation, but cannot introduce new ones or new contributors
at a fixed exponent. A scalar is a one-member Hahn family, so the same
argument gives scalar multiplication; a finite union handles addition.
Source~10 proves this over $k$ by
the same two-coordinate sequence argument; like the Hahn-field construction
itself, that argument uses ordinary set-theoretic choice. A formal derivative
or an Euler operator with coefficients in $k$ reindexes a subfamily and
multiplies its members by elements of the trivially valued field $k$; it
creates no new support exponent. For $F\in\Sfam_m^+$ the exponent zero lies
in only finitely many coefficient supports, so $\rho(F)$ is a polynomial. Each
$X$-coefficient of a product is a finite sum of products of coefficients, the
residue is additive and multiplicative on elements of nonnegative valuation,
and the local finiteness just proved justifies every interchange; so $\rho$ is
a ring homomorphism. Differentiation commutes with residue coefficientwise.
\end{proof}

For $0\ne F=\sum_\alpha b_\alpha X^\alpha\in\Sfam_m$ put
\[
 \gamma(F)=\min_{b_\alpha\ne0}v(b_\alpha),\qquad
 \inpoly(F)=\rho\bigl(t^{-\gamma(F)}F\bigr).
\]
The minimum exists because it is the least element of the nonempty
well-ordered union of supports. Then $t^{-\gamma(F)}F\in\Sfam_m^+$ and the
\emph{initial polynomial} $\inpoly(F)$ is nonzero.

\begin{corollary}[Gauss multiplicativity]\label{hol:nl:cor:gauss}
For nonzero $F,G\in\Sfam_m$,
\[
 \gamma(FG)=\gamma(F)+\gamma(G),\qquad
 \inpoly(FG)=\inpoly(F)\,\inpoly(G).
\]
\end{corollary}

\begin{proof}
Normalize both factors and apply Lemma~\ref{hol:nl:lem:algebra}. The product
of the two nonzero residues is nonzero because $k[X]$ is an integral domain,
so some coefficient of $t^{-\gamma(F)-\gamma(G)}FG$ has valuation exactly zero,
and none has negative valuation. This proves the asserted minimum. Taking
the residue of that normalization then gives the initial-polynomial identity.
\end{proof}

\subsection{Scaled initial polynomials}\label{hol:nl:sub:scaled}

Let $f=\sum_{n\ge0}a_nz^n\in\KK[[z]]$ and $\delta\in\Gamma$ with
$t^{-\delta}\in\Dom(f)$. Then
\[
 F_\delta(X)=f(t^{-\delta}X)=\sum_{n\ge0}a_nt^{-n\delta}X^n\in\Sfam_1.
\]
For $f\ne0$ put
\begin{equation}\label{hol:nl:eq:gaussdata}
 \gamma_\delta=\min_{a_n\ne0}\bigl(v(a_n)-n\delta\bigr),\qquad
 g_\delta=t^{-\gamma_\delta}F_\delta,\qquad
 \phi_\delta=\rho(g_\delta),\qquad
 N_\delta=\deg\phi_\delta.
\end{equation}
The \emph{active} indices are the $n$ attaining $\gamma_\delta$; there are
finitely many, $N_\delta$ is the largest, and $A_\delta=\lc(\phi_\delta)\in
k^\times$.
More explicitly,
\[
 \phi_\delta(X)=
 \sum_{v(a_n)-n\delta=\gamma_\delta}\lc(a_n)X^n.
\]
Here only nonzero $a_n$ are included. Each active index contributes at the
same Hahn exponent $\gamma_\delta$ in the scaled coefficient family, so
local finiteness makes this sum finite. Distinct indices remain distinct
powers of $X$; their nonzero leading coefficients cannot cancel as a
polynomial.

The element $\gamma_\delta$ is a \emph{coefficient Gauss value}. It need not
equal $v(f(t^{-\delta}))$: $\phi_\delta(1)$ can vanish although
$\phi_\delta\ne0$. The proofs below use identities in $\Sfam_1$ and in $k[X]$
only, never an identification of these two quantities, and never replace the
initial polynomial by its value at~$1$.
For example, with $\eta>0$ the finite family
$F(X)=1-X+t^\eta X^2$ has $\gamma(F)=0$ and $\inpoly(F)=1-X$,
but $F(1)=t^\eta$ has valuation $\eta>0$.

\begin{lemma}[Exact derivative scaling]\label{hol:nl:lem:derivatives}
For every $j\ge0$,
\begin{equation}\label{hol:nl:eq:scalederivative}
 (D_z^jf)(t^{-\delta}X)=t^{\gamma_\delta+j\delta}\,g_\delta^{(j)}(X),
 \qquad \rho\bigl(g_\delta^{(j)}\bigr)=\phi_\delta^{(j)}.
\end{equation}
If $N_\delta\ge j$, then $\phi_\delta^{(j)}$ has degree $N_\delta-j$ and
leading coefficient $A_\delta\fall{N_\delta}{j}$.
\end{lemma}

\begin{proof}
The scaling identity is a formal coefficient computation:
$D_X^jF_\delta(X)=t^{-j\delta}(D_z^jf)(t^{-\delta}X)$. The residue identity is
Lemma~\ref{hol:nl:lem:algebra}. The falling factorial is a nonzero element of
$k$ for $N_\delta\ge j$ because $k$ has characteristic zero.
\end{proof}

\begin{lemma}[Large active degrees]\label{hol:nl:lem:active}
Let $f$ be nonpolynomial, let $n_\ast\in\N$, and choose $m>n_\ast$ with
$a_m\ne0$. There is $\delta_1\in\Gamma_{>0}$ such that, whenever
$\delta\ge\delta_1$ and $t^{-\delta}\in\Dom(f)$, every active index exceeds
$n_\ast$. Moreover $\gamma_\delta\le v(a_m)-m\delta$.
\end{lemma}

\begin{proof}
Choose $\delta_1>0$ with
\begin{equation}\label{hol:nl:eq:lowindices}
 (m-n)\delta_1>v(a_m)-v(a_n)\qquad(0\le n\le n_\ast,\ a_n\ne0).
\end{equation}
There are finitely many inequalities, and they need no division in $\Gamma$:
take $\delta_1$ positive and larger than the absolute values of all right-hand
sides, since each multiplier $m-n$ is a positive integer. Such an element
exists: add any positive element of the nonzero ordered group to the maximum
of zero and those finitely many absolute values. This also covers an empty
list of low-index coefficients. The inequalities persist for
$\delta\ge\delta_1$ and say that each such low-index term has scaled valuation
$v(a_n)-n\delta$ strictly larger than that of the $m$th term. The last
assertion is the definition of $\gamma_\delta$.
\end{proof}

This is not a limit $N_\delta\to\infty$ along a countable cofinal sequence,
which need not exist; it concerns a final segment of $\Gamma_{>0}$ and only
those scales at which strong evaluation is available. It is also unrelated to
the escape chain of Lemma~\ref{hol:lem:escape}: no recurrence is used.

\subsection{The finite corner of a differential polynomial}
\label{hol:nl:sub:corner}

Let $P$ be a nonzero differential polynomial of order at most $s$, written in
combined-monomial form
\begin{equation}\label{hol:nl:eq:P}
 P(z,Y_0,\ldots,Y_s)=\sum_{\nu\in J}c_\nu z^{k_\nu}\prod_{j=0}^sY_j^{m_{\nu j}},
\end{equation}
with $J$ finite, every $c_\nu\in\KK^\times$, and identical monomials already
combined. Write $P[f]=P(z,f,D_zf,\ldots,D_z^sf)$. Define the integers
\begin{equation}\label{hol:nl:eq:dw}
 d_\nu=\sum_{j=0}^sm_{\nu j},\qquad
 w_\nu=k_\nu-\sum_{j=0}^sj\,m_{\nu j},\qquad
 d_P=\max_\nu d_\nu,\qquad
 w_P=\max_{d_\nu=d_P}w_\nu,
\end{equation}
and
\[
 \mathcal C_P=\{\nu:d_\nu=d_P,\ w_\nu=w_P\},\qquad
 \beta_P=\min_{\nu\in\mathcal C_P}v(c_\nu),\qquad
 \bar c_\nu=\res\bigl(t^{-\beta_P}c_\nu\bigr)\in k\quad(\nu\in\mathcal C_P).
\]
Thus $(d_P,w_P)$ is chosen lexicographically: first the largest total
differential degree, then the largest shift $w$. This is a finite corner
selection, not the construction of a differential Newton polygon. The
\emph{residue corner polynomial} and the \emph{full corner polynomial} are
\begin{equation}\label{hol:nl:eq:chiI}
 \chi_P(T)=\sum_{\nu\in\mathcal C_P}\bar c_\nu
 \prod_{j=0}^s\fall Tj^{m_{\nu j}}\in k[T],\qquad
 I_P(T)=\sum_{\nu\in\mathcal C_P}c_\nu
 \prod_{j=0}^s\fall Tj^{m_{\nu j}}\in\KK[T].
\end{equation}
In general $\chi_P$ can vanish although some $\bar c_\nu\ne0$, because
different products of falling factorials can be linearly dependent. If
$\chi_P\ne0$, then $I_P\ne0$, and
\begin{equation}\label{hol:nl:eq:rootcontainment}
 I_P(n)=0\ \Longrightarrow\ \chi_P(n)=0\qquad(n\in\N),
\end{equation}
because $t^{-\beta_P}I_P$ has coefficients of nonnegative valuation and residue
commutes with substitution of an ordinary integer.

\begin{lemma}[The exact corner identity]\label{hol:nl:lem:corner}
Let $\phi\in k[X]$ have degree $N\ge s$ and leading coefficient $A\ne0$, and
put
\[
 \Omega_\phi(X)=\sum_{\nu\in\mathcal C_P}\bar c_\nu X^{k_\nu}
 \prod_{j=0}^s\bigl(\phi^{(j)}(X)\bigr)^{m_{\nu j}}.
\]
Then $\deg\Omega_\phi\le d_PN+w_P$ and
\begin{equation}\label{hol:nl:eq:corneridentity}
 [X^{d_PN+w_P}]\,\Omega_\phi=A^{d_P}\chi_P(N).
\end{equation}
In particular $\chi_P(N)\ne0$ implies $\Omega_\phi\ne0$.
\end{lemma}

\begin{proof}
Each corner summand has degree at most
$k_\nu+\sum_jm_{\nu j}(N-j)=d_PN+w_P$. This common degree is a nonnegative
integer because $N\ge s\ge j$ and $k_\nu,m_{\nu j}\ge0$, even if $w_P<0$.
The coefficient of the summand in that degree is
$\bar c_\nu A^{d_P}\prod_j\fall Nj^{m_{\nu j}}$. Summing gives the formula.
\end{proof}

Only finitely many nonnegative integers are roots of a nonzero $\chi_P$, even
when $k$ is not algebraically closed. The active degree is an ordinary
integer, not an element of $\Gamma$, and no root extraction or division in
$\Gamma$ occurs in this step.

\section{First-order algebraic differential equations}\label{hol:nl:sec:first}

\subsection{A finite certificate excluding all large scales}

The following theorem carries the proof of Theorem~\ref{hol:nl:main:first}
and is also a practical audit object: its scale inequalities use only finitely
many coefficients of the solution. It applies to equations of any order.

\begin{theorem}[Nonlinear finite exclusion certificate]
\label{hol:nl:thm:certificate}
Let $P$ be as in \eqref{hol:nl:eq:P} with $d_P\ge1$ and $\chi_P\ne0$, and let
$f=\sum_na_nz^n$ be a nonpolynomial formal solution of $P[f]=0$. Choose an
integer $n_\ast\ge s$ at least as large as every nonnegative integer root of
$\chi_P$, and choose $m>n_\ast$ with $a_m\ne0$ and
\begin{equation}\label{hol:nl:eq:slope}
 (d_P-d_\nu)m+w_P-w_\nu>0\qquad(d_\nu<d_P).
\end{equation}
Every $\delta\in\Gamma_{>0}$ satisfying the finitely many strict inequalities
\begin{align}
 (m-n)\delta&>v(a_m)-v(a_n)
 &&(0\le n\le n_\ast,\ a_n\ne0),\label{hol:nl:eq:cert1}\\
 (w_P-w_\nu)\delta&>\beta_P-v(c_\nu)
 &&(d_\nu=d_P,\ w_\nu<w_P),\label{hol:nl:eq:cert2}\\
 \bigl((d_P-d_\nu)m+w_P-w_\nu\bigr)\delta
 &>\beta_P-v(c_\nu)+(d_P-d_\nu)v(a_m)
 &&(d_\nu<d_P)\label{hol:nl:eq:cert3}
\end{align}
has $t^{-\delta}\notin\Dom(f)$. Such $m$ and $\delta$ exist, and the
inequalities persist when $\delta$ increases. Consequently $\Dom(f)$ omits
$t^{-\delta}$ for every $\delta$ in a final segment of $\Gamma_{>0}$, and every
strongly entire solution of $P[f]=0$ is a polynomial.
\end{theorem}

\begin{proof}
An integer $n_\ast$ exists because $\chi_P\ne0$. For $d_\nu<d_P$ the
coefficient $d_P-d_\nu$ of $m$ in \eqref{hol:nl:eq:slope} is a positive
integer, so all sufficiently large $m$ satisfy it, and a nonpolynomial $f$ has
nonzero coefficients of arbitrarily large index. Every multiplier on the left
of \eqref{hol:nl:eq:cert1}--\eqref{hol:nl:eq:cert3} is then a positive
integer, so any $\delta$ that is positive and larger than the absolute values
of all right-hand sides satisfies them, as does every larger $\delta$.

Suppose $t^{-\delta}\in\Dom(f)$ for such a $\delta$, and form
$\gamma=\gamma_\delta$, $g=g_\delta$, $\phi=\phi_\delta$, $N=N_\delta$ as in
\eqref{hol:nl:eq:gaussdata}. Each inequality in \eqref{hol:nl:eq:cert1}
puts that low-index term strictly above the $m$th term in scaled valuation.
Thus no index at most $n_\ast$ is active, by the comparison in the proof of
Lemma~\ref{hol:nl:lem:active}, and
\begin{equation}\label{hol:nl:eq:hbound}
 N>n_\ast,\qquad \gamma\le v(a_m)-m\delta.
\end{equation}
All derivatives of $g$ lie in $\Sfam_1^+$. Substituting $z=t^{-\delta}X$, a
ring homomorphism $\KK[[z]]\to\KK[[X]]$, Lemma~\ref{hol:nl:lem:derivatives}
turns the $\nu$th summand of $P[f]$ into
\[
 c_\nu\,t^{d_\nu\gamma-w_\nu\delta}\,X^{k_\nu}\prod_j\bigl(g^{(j)}\bigr)^{m_{\nu j}}.
\]
Multiply the whole identity by $t^{-\lambda}$ with
$\lambda=\beta_P+d_P\gamma-w_P\delta$. The scalar in front of the $\nu$th
product then has valuation
\begin{equation}\label{hol:nl:eq:scalarvalue}
 v(c_\nu)-\beta_P+(d_\nu-d_P)\gamma+(w_P-w_\nu)\delta.
\end{equation}
For $\nu\in\mathcal C_P$ this is $v(c_\nu)-\beta_P\ge0$. For $d_\nu=d_P$ and
$w_\nu<w_P$ it is positive by \eqref{hol:nl:eq:cert2}. For $d_\nu<d_P$,
multiplying the bound on $\gamma$ in \eqref{hol:nl:eq:hbound} by the negative
integer $d_\nu-d_P$ reverses it and gives the lower bound
\[
 v(c_\nu)-\beta_P-(d_P-d_\nu)v(a_m)
 +\bigl((d_P-d_\nu)m+w_P-w_\nu\bigr)\delta>0
\]
by \eqref{hol:nl:eq:cert3}. Every polynomial product of the $g^{(j)}$ has
coefficients of nonnegative valuation. Multiplication by a positive-valuation
scalar makes its residue zero, including when the product itself is zero.
Thus the normalized identity lies in $\Sfam_1^+$, and only corner summands
can survive the residue map. Applying $\rho$ to the
formal identity gives
\begin{equation}\label{hol:nl:eq:residualequation}
 0=\sum_{\nu\in\mathcal C_P}\bar c_\nu X^{k_\nu}
 \prod_j\bigl(\phi^{(j)}\bigr)^{m_{\nu j}}=\Omega_\phi(X).
\end{equation}
But $N>n_\ast\ge s$ and $\chi_P(N)\ne0$, so by
Lemma~\ref{hol:nl:lem:corner} the coefficient of $X^{d_PN+w_P}$ in
$\Omega_\phi$ is $\lc(\phi)^{d_P}\chi_P(N)\ne0$, a contradiction. The last
assertion follows because a strongly entire series is evaluable at every
monomial argument.
\end{proof}

\paragraph{An explicit, division-free exclusion bound.}
Let $\mathcal R$ be the finite set of right-hand sides of
\eqref{hol:nl:eq:cert1}--\eqref{hol:nl:eq:cert3}; an empty list contributes
nothing. For any $\epsilon\in\Gamma_{>0}$,
\begin{equation}\label{hol:nl:eq:deltazero}
 \delta_0=\epsilon+\max\bigl(\{0\}\cup\{\abs b:b\in\mathcal R\}\bigr)
\end{equation}
satisfies all of them, since their multipliers are positive integers. This is
an explicit certificate; it is not claimed to be optimal. The proof uses the
full supports to form the initial polynomial but only finitely many leading
valuations to choose $\delta_0$. Leading valuations alone do \emph{not} in
general decide strong evaluation at a boundary argument, and the certificate
has no converse.

For $k=\C$, Lemma~\ref{hol:lem:inward} (inward stability) turns the monomial
conclusion into the exterior form used in Theorem~\ref{hol:main:ode}: a
nonpolynomial formal solution is then not strongly evaluable at any
$x\ne0$ with $v(x)\le-\delta_0$. This combination of the two statements is
made by the merge; source~10 states only the monomial form.

\begin{corollary}[Cofinal scales suffice]\label{hol:nl:cor:cofinal}
If $\chi_P\ne0$, a formal solution of $P[f]=0$ strongly evaluable at
$t^{-\delta}$ for a cofinal set of $\delta\in\Gamma_{>0}$ is a polynomial.
\end{corollary}

\begin{proof}
If $d_P=0$, then $P$ is a nonzero polynomial in $z$ alone and has no formal
solution. Otherwise a cofinal subset meets the final segment excluded by
Theorem~\ref{hol:nl:thm:certificate} for a nonpolynomial solution.
\end{proof}

\paragraph{Dependency audit.}
The only infinite-support step is Lemma~\ref{hol:nl:lem:algebra}, with the
existence and finiteness of a leading coefficient layer; the rest is finite
polynomial algebra. The proof makes none of the following replacements: a
higher-rank group by a real-valued norm; strong summation by convergence of
finite subsums; a nonsummable family by a cancelled expression; the residue
polynomial by its value at~$1$; a nonlinear coefficient recurrence by a
linear one. The three critical checkpoints are the minimum in
\eqref{hol:nl:eq:gaussdata}, the order reversal when \eqref{hol:nl:eq:hbound}
is multiplied by $d_\nu-d_P<0$, and the nonvanishing of $\chi_P(N)$.

\subsection{Why the corner never degenerates in order one; proof of
\texorpdfstring{Theorem~\ref{hol:nl:main:first}}{Theorem D}}
\label{hol:nl:sub:firstorder}

Write a first-order equation in combined-monomial form as
\begin{equation}\label{hol:nl:eq:firstP}
 P(z,Y,Y')=\sum_{k,i,j}c_{kij}z^kY^i(Y')^j.
\end{equation}
For a term, $d=i+j$ and $w=k-j$. At the corner, fixing $j$ forces
$i=d_P-j$ and $k=w_P+j$, so there is at most one corner term for each $j$, and
\eqref{hol:nl:eq:chiI} becomes
\begin{equation}\label{hol:nl:eq:firstchi}
 \chi_P(T)=\sum_{\substack{i+j=d_P\\ k-j=w_P}}\bar c_{kij}T^j,\qquad
 I_P(T)=\sum_{\substack{i+j=d_P\\ k-j=w_P}}c_{kij}T^j.
\end{equation}
At least one $\bar c_{kij}$ is nonzero by the choice of $\beta_P$, and its
power $T^j$ cannot cancel against another summand. Hence $\chi_P\ne0$.

\begin{proof}[Proof of Theorem~\ref{hol:nl:main:first}]
If $d_P=0$, then $P$ is a nonzero polynomial in $z$ alone and has no formal
solutions. Otherwise the preceding observation and
Theorem~\ref{hol:nl:thm:certificate} give polynomiality and the exclusion
statement, with $\delta_0$ as in \eqref{hol:nl:eq:deltazero}. A nonzero
algebraic relation between $f$ and $D_zf$ over $\KK(z)$ becomes a nonzero
polynomial of the form \eqref{hol:nl:eq:firstP} after denominators are
cleared. Multiplication by a common nonzero denominator preserves both the
nonzero relation and its vanishing in $\KK((z))$, even if that denominator
vanishes at $z=0$. Thus the polynomiality conclusion contradicts
nonpolynomiality of $f$, proving the independence statement.
\end{proof}

The uniqueness of a corner term for fixed $j$ is a structural fact of order
one. It is exactly what fails when several derivative orders can contribute
the same falling-factorial polynomial (Section~\ref{hol:nl:sub:vanishing}).

\subsection{Finite candidate degrees and the entire-solution locus}
\label{hol:nl:sub:degrees}

Once polynomiality is known, ordinary polynomial degree gives a sharper
finite candidate set that uses the \emph{full} coefficients. Assume $d_P\ge1$
and define the nonnegative rational number
\begin{equation}\label{hol:nl:eq:kappa}
 \kappa_P=\max\Bigl(\{0\}\cup\Bigl\{\frac{w_\nu-w_P}{d_P-d_\nu}:d_\nu<d_P\Bigr\}\Bigr)\in\Q.
\end{equation}
This divides ordinary integers in $\Q$, not elements of $\Gamma$.

\begin{theorem}[Finite candidate degrees]\label{hol:nl:thm:degrees}
For the first-order polynomial \eqref{hol:nl:eq:firstP} with $d_P\ge1$, every
nonconstant strongly entire solution has degree in the finite set
\begin{equation}\label{hol:nl:eq:candidates}
 \mathcal D_P=\{n\in\Z_{>0}:n\le\kappa_P\}\ \cup\ \{n\in\Z_{>0}:I_P(n)=0\}.
\end{equation}
The constant solutions are exactly the $c\in\KK$ with $P(z,c,0)=0$ as a
polynomial in $z$. The set \eqref{hol:nl:eq:candidates} is necessary; it does
not promise that each candidate occurs.
\end{theorem}

\begin{proof}
Polynomiality is Theorem~\ref{hol:nl:main:first}. Let $f=az^n+\cdots$ with
$n\ge1$ and $a\ne0$. A term of \eqref{hol:nl:eq:firstP} has degree
$k+in+j(n-1)=dn+w$ and leading coefficient $c_{kij}a^dn^j$. If $n>\kappa_P$,
every term with $d<d_P$ has degree below $d_Pn+w_P$, and so do the terms with
$d=d_P$, $w<w_P$. So the coefficient of $z^{d_Pn+w_P}$ in $P[f]$ is
$a^{d_P}I_P(n)$, which must vanish. $I_P$ is nonzero by
\eqref{hol:nl:eq:firstchi}, so its set of ordinary integer roots is finite.
The assertion on constants is immediate by substitution.
\end{proof}

\begin{remark}[What ``finite'' does and does not compute]
The set of ordinary integer roots of $I_P$ is well defined and finite for
arbitrary Hahn coefficients, but its existence is not a uniform algorithm for
arbitrary coefficient names: extracting a leading exponent, deciding
coefficient equality, or testing an arbitrary complex coefficient for exact
integrality can need information that no numerical approximation supplies.
For rational coefficients, or other exact representations with the needed
operations, \eqref{hol:nl:eq:candidates} is an actual finite algebraic
search. The shipped program uses only exact rational data.
\end{remark}

Replacing $I_P(n)=0$ by $\chi_P(n)=0$ gives a coarser necessary set, by
\eqref{hol:nl:eq:rootcontainment}. Keeping $I_P$ can be genuinely sharper.

\begin{example}[A resonance visible only in the first layer]
\label{hol:nl:ex:resonance}
Let $c=t^\eta$ with $\eta>0$, $N\in\Z_{>0}$, and $P=zY'-(N+c)Y$. Then
$\chi_P(T)=T-N$ but $I_P(T)=T-N-c$ has no ordinary nonnegative integer root.
The only formal solution in $\KK[[z]]$ is zero, since coefficient comparison
gives $(n-N-c)a_n=0$ for every $n$. A residue resonance alone is not a
polynomial solution.
\end{example}

\begin{corollary}[An algebraic entire-solution locus]\label{hol:nl:cor:locus}
For a first-order equation $P$ with $d_P\ge1$ put
$n_P=\max(\{0\}\cup\mathcal D_P)$. There is an explicitly defined affine
algebraic set $\mathcal V_P\subseteq\mathbb A_{\KK}^{n_P+1}$ whose
$\KK$-points are in bijection with all strongly entire formal solutions of
$P[f]=0$. The same equations parametrize all strongly entire solutions over
every Hahn-field extension $\KK\hookrightarrow\mathbb L=k'((t^{\Gamma'}))$
induced by an embedding of coefficient fields $k\hookrightarrow k'$ and an
embedding of ordered groups $\Gamma\hookrightarrow\Gamma'$.
\end{corollary}

\begin{proof}
With indeterminates $u_0,\ldots,u_{n_P}$ put $F_u(z)=\sum_{n\le n_P}u_nz^n$ and
expand $P(z,F_u,F_u')=\sum_\ell Q_\ell(u)z^\ell$ with
$Q_\ell\in\KK[u_0,\ldots,u_{n_P}]$. Let $\mathcal V_P$ be the common zero
locus of these finitely many $Q_\ell$. Its points are exactly the polynomial
solutions of degree at most $n_P$, each of which is strongly entire, and
Theorem~\ref{hol:nl:thm:degrees} shows there are no others. Under the stated
extension the monomial degrees and $\kappa_P$ are unchanged, and for an
ordinary integer $n$ the element $I_P(n)$ vanishes in $\mathbb L$ if and only
if it vanishes in $\KK$. Theorem~\ref{hol:nl:main:first} applies over
$\mathbb L$, whose coefficient field again has characteristic zero and whose
group is again nonzero, so the same degree bound and equations remain
exhaustive there.
\end{proof}

This is a finite-dimensional algebraic parametrization, not a claim that the
number of solutions is finite, and not a moduli statement over coefficient
rings with nilpotents. The extension class is exactly the one stated.

\subsection{Worked first-order equations}\label{hol:nl:sub:examples}

All identities use the formal derivative $D_z$. For $k=\C$ they are statements
about surcomplex Hahn workspaces; for $k=\R$ the real versions follow with the
non-real constant roots omitted. Consider
\begin{equation}\label{hol:nl:eq:riccati}
 D_zf=f^2+p(z),\qquad p\in\KK[z].
\end{equation}
The corner is the term $-Y^2$, so $d_P=2$, $w_P=0$ and $I_P=-1$. If $p$ has
degree $e\ge0$, then $\kappa_P=\max(0,e/2)$, so a nonconstant entire solution
has degree at most $\lfloor e/2\rfloor$. More precisely, a solution of degree
$n\ge1$ requires $\deg p=2n$ and the leading coefficient of $p$ to be minus
the square of that of $f$, since the top term of $f^2$ cannot be cancelled by
$D_zf$. These conditions are necessary only; the remaining coefficient
equations must still hold.

\begin{proposition}[Two exact Riccati classifications]\label{hol:nl:prop:riccati}
Over $\KK=\C((t^\Gamma))$, the strongly entire solutions of $D_zf=f^2+1$ are
exactly $f=i$ and $f=-i$, and the only strongly entire solution of
$D_zf=f^2+1-z^2$ is $f=z$.
\end{proposition}

\begin{proof}
In the first equation the degree bound excludes nonconstant solutions, and a
constant must satisfy $c^2+1=0$; in the field $\KK$, which contains $i$, the
roots are $\pm i$. In the second the degree bound is one and constants are
impossible. For $f=az+b$ coefficient comparison gives $a^2=1$, $2ab=0$ and
$a=b^2+1$, so $b=0$ and $a=1$, and $f=z$ is verified directly.
\end{proof}

The first equation is also \texttt{diff:cor:riccati} of the collection's
differential-equations report~\cite{repodifferential}, with the same two
solutions, but there $\partial u=1+u^2$ is solved in $\No[i]$ for a derivation
on surreal \emph{scalars}. The two statements concern different derivations
and different objects, and neither implies the other.

\begin{example}[The degree cutoff is attained]\label{hol:nl:ex:cutoff}
For every integer $m\ge1$ the equation $D_zf=f^2+mz^{m-1}-z^{2m}$ has the
strongly entire solution $f=z^m$. Here $I_P=-1$ and $\kappa_P=m$ exactly, so
the low-degree part of \eqref{hol:nl:eq:candidates} cannot in general be
replaced by $n<\kappa_P$.
\end{example}

\begin{example}[Arbitrarily large resonant degrees]\label{hol:nl:ex:resonant}
For $N\ge1$ the equation
\begin{equation}\label{hol:nl:eq:resonant}
 zf\,D_zf-Nf^2=0
\end{equation}
has $d_P=2$, $w_P=0$ and $I_P(T)=T-N$. Since $\KK[[z]]$ is an integral
domain, a nonzero solution satisfies $zD_zf=Nf$, and coefficient comparison
gives exactly $f=cz^N$ with $c\in\KK$ (including $c=0$). The exceptional
integer root is an actual degree, and no degree bound can depend only on the
positions of the monomials of $P$: the coefficient $N$ matters.
\end{example}

\subsection{The exclusion bound cannot be uniform}\label{hol:nl:sub:nouniform}

Theorem~\ref{hol:main:ode} has an exclusion bound depending only on the
operator. For nonlinear equations this is false, already for one fixed
equation of first order.

\begin{theorem}[A fixed Riccati equation defeats every uniform bound]
\label{hol:nl:thm:nouniform}
For $\lambda\in\Gamma$ let $c=t^\lambda$ and
\[
 f_\lambda(z)=\frac{c}{1-cz}=\sum_{n\ge0}c^{n+1}z^n\in\KK[[z]].
\]
Then $D_zf_\lambda=f_\lambda^2$, $f_\lambda$ is nonpolynomial, and
\begin{equation}\label{hol:nl:eq:riccatidomain}
 \Dom(f_\lambda)=\{0\}\cup\{x\in\KK^\times:v(x)+\lambda>0\}.
\end{equation}
Consequently, for every $\delta\in\Gamma$ some nonpolynomial solution of the
fixed equation $D_zf=f^2$ is strongly evaluable at $t^{-\delta}$, and no
exclusion bound depending on the equation alone excludes all its
nonpolynomial solutions.
\end{theorem}

\begin{proof}
The derivative identity is a coefficient comparison. For $x\ne0$ put $u=cx$,
so the $n$th evaluation term is $cu^n$. If $v(u)>0$, write
$u=b\,t^{v(u)}(1+w)$ with $b\in k^\times$ and $w$ zero or of positive
valuation; then $\supp(u^n)\subseteq nv(u)+M(\supp w)$, the leading exponents
$nv(u)$ increase strictly, and Lemma~\ref{hol:lem:certificate}\,(b), resting
on the positive-support monoid of Lemma~\ref{hol:lem:support}\,(2), gives
strong summability; multiplication by $c$ preserves it. (Source~10 cites the
Hahn--Neumann positive-support lemma here instead of proving it; no
Archimedean convergence of $n\,v(u)$ is used.) If $v(u)=0$, every $cu^n$ has
the nonzero coefficient $\lc(u)^n$ at exponent $\lambda$, so local finiteness
fails. If $v(u)<0$, the leading exponents $\lambda+nv(u)$ strictly decrease.
At $x=0$ only one member is nonzero. This proves
\eqref{hol:nl:eq:riccatidomain}. Given $\delta$, choose
$\lambda=\delta+\epsilon$ with any $\epsilon\in\Gamma_{>0}$, available
because $\Gamma\ne0$; no order unit is required. Then
$v(t^{-\delta})+\lambda=\lambda-\delta=\epsilon>0$.
\end{proof}

There is no conflict with Theorem~\ref{hol:nl:thm:certificate}, whose
finite data, and hence whose bound, may depend on the solution; for
$f_\lambda$ the exact boundary moves with $\lambda$. At the boundary
$v(x)=-\lambda$ geometric-series algebra cannot legalize infinitely many
contributions at one exponent. Nor is there a conflict with
Theorem~\ref{hol:main:ode}: each $f_\lambda$ also satisfies the linear
equation $(1-t^\lambda z)D_zf-t^\lambda f=0$, whose operator depends on
$\lambda$, exactly as the operator $D_z-t^\eta$ of the formal exponential of
Proposition~\ref{hol:prop:expdomain} depends on~$\eta$. (This comparison is
the merge's; it is a direct computation.)

\section{Higher order and positive-weight Euler equations}
\label{hol:nl:sec:higher}

\subsection{The nondegenerate higher-order class}

\begin{theorem}[Higher-order corner criterion and its boundary]
\label{hol:nl:thm:corner}
Let $P$ be a nonzero differential polynomial \eqref{hol:nl:eq:P} of any finite
order with $d_P\ge1$ and $\chi_P\ne0$. Then every strongly entire formal
solution of $P[f]=0$ is a polynomial, and every nonpolynomial formal solution
fails strong evaluation at $t^{-\delta}$ for all $\delta$ in a final segment
of $\Gamma_{>0}$, with the finite certificate of
Theorem~\ref{hol:nl:thm:certificate}. The condition $\chi_P\ne0$ holds
automatically in order one, but not in order two: a fixed second-order
equation can have polynomial solutions of every degree
(Proposition~\ref{hol:nl:prop:degenerate}).
\end{theorem}

\begin{proof}
The first two assertions are Theorem~\ref{hol:nl:thm:certificate}; the third
is Section~\ref{hol:nl:sub:firstorder} and
Proposition~\ref{hol:nl:prop:degenerate}.
\end{proof}

\begin{proposition}[Higher-order degree candidates]
\label{hol:nl:prop:higherdegrees}
Suppose $d_P\ge1$ and $\chi_P\ne0$ for an equation of order at most $s$. Every
nonzero strongly entire solution has degree in the finite set
\[
 \{0,\ldots,s-1\}\ \cup\ \{n\in\N:n\ge s,\ n\le\kappa_P\}\ \cup\
 \{n\in\N:n\ge s,\ I_P(n)=0\},
\]
with $\kappa_P$ as in \eqref{hol:nl:eq:kappa}. When $s=0$ the first set is
empty. Constant solutions can always be checked by substitution.
\end{proposition}

\begin{proof}
Polynomiality is Theorem~\ref{hol:nl:thm:certificate}. For degree $n\ge s$
every derivative leading-term formula of Lemma~\ref{hol:nl:lem:corner} applies,
and if also $n>\kappa_P$ the coefficient of $z^{d_Pn+w_P}$ in $P[f]$ is
$\lc(f)^{d_P}I_P(n)$, which must vanish; $I_P\ne0$ because $\chi_P\ne0$.
Degrees below $s$, where derivative leading terms can vanish, are listed
separately.
\end{proof}

\begin{proposition}[The first two Painlev\'e polynomial equations]
\label{hol:nl:prop:painleve}
In characteristic zero, the equation $D_z^2f=6f^2+z$ has no strongly entire
formal solution, and the equation $D_z^2f=2f^3+zf+\alpha$ with $\alpha\in\KK$
has a strongly entire solution exactly when $\alpha=0$, the only one then
being $f=0$.
\end{proposition}

\begin{proof}
The first equation has $d_P=2$, $w_P=0$, $\chi_P=-6$; the second has $d_P=3$,
$w_P=0$, $\chi_P=-2$. Theorem~\ref{hol:nl:thm:certificate} makes every entire
solution a polynomial. In the first equation a polynomial of degree $n\ge1$
gives $f^2$ of degree $2n$, larger than $\deg D_z^2f$ (when nonzero) and than
$\deg z=1$, so no cancellation is possible, and a constant cannot cancel $z$.
In the second, for $n\ge1$ one has $3n>n+1$ and $3n>\deg D_z^2f$, so $2f^3$
cannot cancel; a constant $c$ must satisfy $2c^3+zc+\alpha=0$ identically,
which forces $c=0$ and $\alpha=0$, and substitution proves sufficiency.
\end{proof}

The names only identify the displayed equations; no classification of their
classical meromorphic solutions is imported. The examples illustrate a class
beyond first order: a unique term of greatest total differential degree has a
nonzero corner whenever its falling-factorial product is nonzero. In
particular an equation $D_z^sf=Q(z,f)$ with $Q$ of degree at least two in $f$
has a nonzero corner and hence only polynomial strongly entire solutions; the
coefficient of the top power of $f$ may be a nonconstant polynomial in $z$,
whose highest $z$-term then selects the corner.

\subsection{A vanishing corner}\label{hol:nl:sub:vanishing}

Consider
\begin{equation}\label{hol:nl:eq:degenerate}
 z^2f\,D_z^2f+zf\,D_zf-z^2(D_zf)^2=0.
\end{equation}
All three terms have $d=2$ and $w=0$, and the residue corner polynomial is
\begin{equation}\label{hol:nl:eq:degeneratechi}
 T(T-1)+T-T^2=0,
\end{equation}
an identity in $k[T]$, not an exceptional integer resonance.

\begin{proposition}[Complete formal classification of the example]
\label{hol:nl:prop:degenerate}
The formal solutions of \eqref{hol:nl:eq:degenerate} in $\KK[[z]]$ are exactly
$0$ and $cz^n$ with $c\in\KK^\times$ and $n\in\N$. All are strongly entire,
and their degrees are unbounded for this single equation.
\end{proposition}

\begin{proof}
With $\vartheta=zD_z$, equation \eqref{hol:nl:eq:degenerate} reads
$f\,\vartheta^2f-(\vartheta f)^2=0$. For $f\ne0$ work in the Laurent series
field $\KK((z))$ and divide by $f^2$: $\vartheta(\vartheta f/f)=0$. The
constants of $\vartheta$ in $\KK((z))$ are exactly $\KK$, since
$\vartheta\sum b_nz^n=\sum nb_nz^n$ and characteristic zero kills no nonzero
integer. Hence $\vartheta f/f=c_0\in\KK$; if the first nonzero term of $f$
has exponent $n$, comparing lowest terms gives $c_0=n$. Then $\vartheta f=nf$
kills every coefficient except the $n$th. Substitution verifies the monomials
and zero.
\end{proof}

The example has two messages. A universal higher-order degree bound cannot
be obtained by declaring $\chi_P$ nonzero: it may vanish identically, and a
fixed equation can then have unbounded degrees. And the example is \emph{not}
a nonpolynomial entire solution, so it does not refute polynomial rigidity in
higher order; it locates a failure of this proof and of a stronger uniform
conclusion.

The corner is also not an invariant of the solution set. Multiplying an
equation by another differential polynomial, or passing to a differential
consequence, can change the corner data. For a concrete example put
\[
 Q=z^2Y_0Y_2+zY_0Y_1-z^2Y_1^2,\qquad
 P=Y_1,\qquad \widetilde P=Y_1(1+zQ).
\]
For every formal $f$, the factor $1+zQ[f]$ has constant coefficient one and
is a unit of $\KK[[z]]$. Hence $P[f]=0$ and $\widetilde P[f]=0$ have exactly
the same solutions, the constants. But $\chi_P(T)=T$, whereas the degree-three
corner of $\widetilde P$ gives
$\chi_{\widetilde P}(T)=T\bigl(T(T-1)+T-T^2\bigr)=0$.
The criterion $\chi_P\ne0$ is a
sufficient condition on the displayed combined polynomial: when $\chi_P=0$
another equation satisfied by the same function may have a nonzero corner, or
a logarithmic-derivative reduction as in
Proposition~\ref{hol:nl:prop:degenerate} may be available. A zero corner
cannot be repaired by retaining only monomial degrees, since the residue
coefficients and their exact cancellation are essential. The methods of this
section do not decide whether a higher-order equation with vanishing residue
corner has nonpolynomial strongly entire solutions. When $\Gamma$ has no
order unit, Theorem~\ref{hol:cf:main:entire} (sources~11 and~12) decides it
negatively in every order; Example~\ref{hol:cf:ex:corner} computes the
nonzero linearization multiplier of \eqref{hol:nl:eq:degenerate} along its
solutions. With an order unit the answer is different: the cleared
third-order equation of source~13 has residue corner $\chi_P=0$ and the
nonpolynomial strongly entire solution $\mathcal T_p$
(Theorem~\ref{hol:td:main:boundary}); its full corner polynomial $I_P$ is a
nonzero constant of positive valuation, so the hypothesis $\chi_P\ne0$ of
Theorems~\ref{hol:nl:thm:certificate} and~\ref{hol:nl:thm:corner} cannot be
weakened to $I_P\ne0$ (Proposition~\ref{hol:td:prop:corner} and
Remark~\ref{hol:td:rem:corner}). In order two the vanishing-corner case is
decided negatively for every $\Gamma$ by source~14
(Theorem~\ref{hol:ot:main:second}\,(a)): its endpoint lemma uses the whole
initial polynomial at infinitely many scales, and for
\eqref{hol:nl:eq:degenerate} it permits no nonmonomial initial polynomial at
all (Example~\ref{hol:ot:ex:degenerate}).

\subsection{Positive-weight Euler equations in several variables}
\label{hol:nl:sub:euler}

Fix positive integers $q=(q_1,\ldots,q_m)$, write
$\wt(\alpha)=q\cdot\alpha$ for $\alpha\in\N^m$, and let
\[
 \vartheta_q=\sum_{\nu=1}^mq_\nu x_\nu\frac{\partial}{\partial x_\nu},
 \qquad \vartheta_qx^\alpha=\wt(\alpha)x^\alpha.
\]
For $N\ge0$, a multi-index of weight at most $N$ satisfies
$0\le\alpha_\nu\le\lfloor N/q_\nu\rfloor$ in each coordinate. Thus at most
$\prod_\nu(1+\lfloor N/q_\nu\rfloor)$ multi-indices have that bound. This
finiteness, with the eigenvalue identity, is the essential feature of
positive weights.
Write a nonzero equation in combined-monomial form
\begin{equation}\label{hol:nl:eq:eulerP}
 P(x,Y,Z)=\sum_{\nu\in J}c_\nu x^{\alpha_\nu}Y^{i_\nu}Z^{j_\nu},
\end{equation}
to be read as $P(x,f,\vartheta_qf)=0$, and set
$d_\nu=i_\nu+j_\nu$, $w_\nu=\wt(\alpha_\nu)$, with $d_P$, $w_P$,
$\mathcal C_P$, $\beta_P$ and $\bar c_\nu$ as in
Section~\ref{hol:nl:sub:corner}. The corner polynomials now keep the spatial
variables:
\begin{equation}\label{hol:nl:eq:eulerchi}
 \chi(T,X)=\sum_{\nu\in\mathcal C_P}\bar c_\nu X^{\alpha_\nu}T^{j_\nu},\qquad
 I(T,X)=\sum_{\nu\in\mathcal C_P}c_\nu X^{\alpha_\nu}T^{j_\nu}.
\end{equation}
Both are nonzero: for fixed $(\alpha_\nu,j_\nu)$ the condition $d_\nu=d_P$
forces $i_\nu=d_P-j_\nu$, so no two combined corner terms give the same
monomial in $(X,T)$. Only finitely many integers $n$ make $\chi(n,X)$ the zero
polynomial, since they must all be roots of one nonzero $X$-coefficient of
$\chi$, and likewise for $I$.

\begin{theorem}[Positive-weight Euler rigidity]\label{hol:nl:thm:euler}
Let $f\in\KK[[x_1,\ldots,x_m]]$ satisfy $P(x,f,\vartheta_qf)=0$ for a nonzero
$P\in\KK[x_1,\ldots,x_m,Y,Z]$. If $f$ is nonpolynomial, it fails joint strong
evaluation at $(t^{-q_1\delta},\ldots,t^{-q_m\delta})$ for every sufficiently
large $\delta\in\Gamma_{>0}$. Consequently, if $f$ is strongly entire, or
merely jointly strongly evaluable at these points for a cofinal set of
$\delta\in\Gamma_{>0}$, then $f$ is a polynomial. Every fixed equation has a
finite set of possible positive weighted degrees
(Corollary~\ref{hol:nl:cor:eulerdegrees}).
\end{theorem}

\begin{proof}
If $d_P=0$, then $P$ is a nonzero polynomial in $x$ alone and there is no
solution; assume $d_P\ge1$. Choose $n_\ast\ge0$ with $\chi(n,X)\ne0$ for every integer $n>n_\ast$. Since
$f=\sum_\alpha a_\alpha x^\alpha$ is nonpolynomial and bounded-weight sets are
finite, choose $a_\mu\ne0$ of weight $n_\mu=\wt(\mu)>n_\ast$ so large that
$(d_P-d_\nu)n_\mu+w_P-w_\nu>0$ for $d_\nu<d_P$. Take $\delta>0$ satisfying
the finitely many inequalities
\begin{align*}
 (n_\mu-\wt(\alpha))\delta&>v(a_\mu)-v(a_\alpha)
 &&(\wt(\alpha)\le n_\ast,\ a_\alpha\ne0),\\
 (w_P-w_\nu)\delta&>\beta_P-v(c_\nu)
 &&(d_\nu=d_P,\ w_\nu<w_P),\\
 \bigl((d_P-d_\nu)n_\mu+w_P-w_\nu\bigr)\delta
 &>\beta_P-v(c_\nu)+(d_P-d_\nu)v(a_\mu)
 &&(d_\nu<d_P);
\end{align*}
as before they can be met in every nonzero ordered group and persist on a
final segment. Suppose the joint family is strongly summable at such a scale,
and put
\[
 F(X)=f(t^{-q_1\delta}X_1,\ldots,t^{-q_m\delta}X_m)\in\Sfam_m,\quad
 \gamma=\gamma(F),\quad g=t^{-\gamma}F,\quad \phi=\rho(g).
\]
Then $\gamma\le v(a_\mu)-n_\mu\delta$, and the first inequalities show that
$\phi$ has no monomial of weight at most $n_\ast$; its greatest weight $N$
exceeds $n_\ast$. Let $\phi_N\ne0$ be its weight-$N$ homogeneous part. The
Euler operator commutes with the scaling, so
$(\vartheta_qf)(t^{-q_1\delta}X_1,\ldots,t^{-q_m\delta}X_m)=t^\gamma\vartheta_qg$.
Normalizing the substituted equation by $t^{-\beta_P-d_P\gamma+w_P\delta}$,
the scalar computation \eqref{hol:nl:eq:scalarvalue} shows that only corner
terms survive the residue, whence
\begin{equation}\label{hol:nl:eq:eulerresidue}
 0=\sum_{\nu\in\mathcal C_P}\bar c_\nu X^{\alpha_\nu}\phi^{i_\nu}
 (\vartheta_q\phi)^{j_\nu}.
\end{equation}
The largest possible weight on the right is $d_PN+w_P$. Since the weight-$N$
part of $\vartheta_q\phi$ is $N\phi_N$, the homogeneous part of that weight is
\begin{equation}\label{hol:nl:eq:euleridentity}
 \phi_N(X)^{d_P}\,\chi(N,X),
\end{equation}
a product of two nonzero polynomials, hence nonzero. This contradicts
\eqref{hol:nl:eq:eulerresidue}.
\end{proof}

The proof uses \emph{joint} strong evaluation, not a scalar expression
obtained by first setting $x_1=\cdots=x_m$ and cancelling: a factor
$x_1-x_2$ can make a diagonal scalar expression zero while the ungrouped
monomial family is not summable. Keeping $X_1,\ldots,X_m$ in
\eqref{hol:nl:eq:euleridentity} avoids that error.

\begin{corollary}[Weighted-degree candidates]\label{hol:nl:cor:eulerdegrees}
With $\kappa_P$ defined by \eqref{hol:nl:eq:kappa} from the Euler data
$d_\nu,w_\nu$, every nonconstant strongly entire solution of
$P(x,f,\vartheta_qf)=0$ has greatest weight in the finite set
\[
 \{n\in\Z_{>0}:n\le\kappa_P\}\ \cup\ \{n\in\Z_{>0}:I(n,X)\equiv0\}.
\]
For a fixed equation all strongly entire solutions are therefore described by
finitely many polynomial equations in the coefficients of a bounded-weight
polynomial ansatz.
\end{corollary}

\begin{proof}
Polynomiality is Theorem~\ref{hol:nl:thm:euler}. For a polynomial of greatest
weight $N>\kappa_P$, the terms of lower differential degree have smaller total
weight, and the top-weight part of the equation is $f_N^{d_P}I(N,X)$, which
must vanish; so $I(N,X)=0$. The finite-root argument after
\eqref{hol:nl:eq:eulerchi} applies to $I$, and a bounded-weight ansatz has
finitely many coefficients.
\end{proof}

\begin{example}[An exactly described solution space]\label{hol:nl:ex:eulerspace}
For $N\ge0$ the solutions of $\vartheta_qf=Nf$ in $\KK[[x]]$ are exactly the
weighted homogeneous polynomials of weight $N$, since coefficient comparison
gives $(\wt(\alpha)-N)a_\alpha=0$. For $q=(1,2)$ this space has dimension
$\lfloor N/2\rfloor+1$ over $\KK$, with basis $x_1^{N-2j}x_2^j$,
$0\le j\le\lfloor N/2\rfloor$. The nonlinear equation
$f(\vartheta_qf-Nf)=0$ has the same solutions, because $\KK[[x]]$ is an
integral domain.
\end{example}

\paragraph{Why positive weights are needed.}
Let $\KK=\C((t^\Q))$ and $G(u)=\sum_{n\ge0}t^{n^2}u^n$, which is strongly
entire (Corollary~\ref{hol:nl:cor:theta}). The function $f(x,y)=G(y)$ is
nonpolynomial and strongly entire but satisfies $x\partial_xf=0$, so a zero
weight invalidates the theorem. For the mixed weights $(1,-1)$, the
nonpolynomial strongly entire function $f(x,y)=G(xy)$ satisfies
$(x\partial_x-y\partial_y)f=0$; at a fixed point its nonzero joint evaluation
terms are exactly $t^{n^2}(xy)^n$, which are strongly summable. Positivity is
therefore not a cosmetic substitute for nonzero weights, and a general
first-order partial differential equation need not have the rigidity of a
positive Euler operator.

\paragraph{Relation to Theorem~\ref{hol:main:multi}.}
Theorem~\ref{hol:main:multi} assumes that \emph{all} mixed partial
derivatives span a finite-dimensional space, and Section~\ref{hol:sec:multi}
notes that a single differential equation in several variables cannot force
polynomiality in general, since it may ignore variables; the zero-weight
example above is an instance. Theorem~\ref{hol:nl:thm:euler} instead assumes
one algebraic relation between $f$ and $\vartheta_qf$ with all weights
positive, with no finite-dimensional-span hypothesis. These are different
input criteria and different proof routes. Within the strongly entire class,
each criterion forces a polynomial, which then satisfies both criteria:
its mixed derivatives span a finite-dimensional space, and the relation
$Y-f(x)=0$ is an algebraic Euler relation.

\section{Nonlinear consequences, hypotheses and predecessors}
\label{hol:nl:sec:consequences}

\subsection{Two-jet independence of explicit series}\label{hol:nl:sub:jets}

With the surreal embedding of Section~\ref{hol:sub:embedding}, for instance
$\C((\omega^{-\Q}))$ is a set-sized surcomplex workspace to which
Theorem~\ref{hol:nl:main:first} applies, and its coefficients and admissible
arguments are actual elements of $\No[i]$. The independence conclusions,
however, concern \emph{formal functions}: each value $f(x)$ lies in $\KK$ and
is of course not transcendental over $\KK$.

\begin{corollary}[Two independent jets of the quadratic-exponent series]
\label{hol:nl:cor:theta}
Over $\KK=\C((t^\Q))$ let $G(z)=\sum_{n\ge0}t^{n^2}z^n$. Then $G$ is
nonpolynomial and strongly entire, and $G$ and $D_zG$ are algebraically
independent over $\KK(z)$.
\end{corollary}

\begin{proof}
For $x\ne0$ with $v(x)=\sigma\in\Q$ the $n$th evaluated member has valuation
$n^2+n\sigma$, which eventually exceeds every rational bound, so
Lemma~\ref{hol:lem:certificate}\,(a) gives strong summability even when $x$
has an arbitrary infinite Hahn tail; at $x=0$ the family is finite. All
coefficients are nonzero. Theorem~\ref{hol:nl:main:first} gives the
independence.
\end{proof}

The series is not new: it is the rank-one example of the
rank-one-Berkovich report~\cite{reporankone}, Example
\texttt{found:ex:internal} of the foundations report~\cite{repofoundations},
and a rescaled partial theta series: $G(z)=\Theta_{t^2}(tz)$ in the notation
of \eqref{hol:eq:theta}. The contribution of source~10 is the nonlinear two-jet
consequence only.

\begin{corollary}[An example without an order unit]
\label{hol:nl:cor:infiniterank}
In $K_{\Gamma_\infty}$, with $\Gamma_\infty$ as in \eqref{hol:eq:gammainfty},
the strongly entire series $F=\sum_{n\ge0}t^{e_n}z^n=\sum_{n\ge0}\omega^{-\omega^n}z^n$
of \eqref{hol:eq:explicitF} satisfies: $F$ and $D_zF$ are algebraically
independent over $K_{\Gamma_\infty}(z)$.
\end{corollary}

\begin{proof}
Strong entireness is Theorem~\ref{hol:thm:explicit}; source~10 proves it
again by the cofinal-valuation route recorded in the second paragraph of that
proof. Apply Theorem~\ref{hol:nl:main:first}.
\end{proof}

$F$ and its linear rigidity are in Theorem~\ref{hol:thm:explicit}; the new
statement is the two-jet strengthening. It answers the first-order case of
the last sentence of Question~\ref{hol:q:nonlinear} as posed by sources~08
and~09 (recorded in its status paragraph) with no coefficient
restriction, and says nothing about relations of order two or more;
Theorem~\ref{hol:cf:thm:Fomega} (source~12) answers that sentence in every
order, again with no coefficient restriction.

\paragraph{Extension of scalars versus preservation of entireness.}
The degree classification is stable under the Hahn-field extensions of
Corollary~\ref{hol:nl:cor:locus}, but entireness of a nonpolynomial series
need not be. Enlarge $\Q$ to $\Q+\Q\omega$ with its surreal order and view $G$
over the larger field. At $z=t^{-\omega}$ the member exponents are
$n^2-n\omega$, with successive differences $(2n+1)-\omega<0$, so the support
union contains an infinite descending sequence and the evaluation is not
strong. The same computation appears in the rank-one-Berkovich
report~\cite{reporankone}, and for $F$ in
Section~\ref{hol:sub:notallsurcomplex}. No theorem here claims that $G$ or $F$ is
entire on the proper class $\No[i]$; hypotheses and derivative are relative to
a specified set-sized workspace.

\subsection{Which hypotheses are used}\label{hol:nl:sub:hypotheses}

\begin{example}[Positive characteristic destroys the conclusion]
\label{hol:nl:ex:charp}
Let $p$ be prime and $\KK=\mathbb F_p((t^\Q))$. The nonpolynomial series
$U(z)=\sum_{n\ge0}t^{n^2}z^{pn}$ is strongly entire: at a point of valuation
$\sigma$ its member valuations are $n^2+pn\sigma$, and
Lemma~\ref{hol:lem:certificate}\,(a) applies, its proof being independent of
the coefficient field. Yet $D_zU=0$.
\end{example}

So the conclusion fails in positive characteristic even for one linear
equation. The characteristic-zero requirements of the proof are substantive:
ordinary integer eigenvalues stay distinct, nonzero falling factorials stay
nonzero, and a nonzero polynomial cannot vanish at all ordinary integers.

The hypothesis $\Gamma\ne0$ is needed for the scale argument; the degenerate
case $\Gamma=0$ is treated in Section~\ref{hol:sec:fields}. Under ordinary
complex analytic summation, by contrast, $\exp z$ is a nonpolynomial entire
function with $f'=f$; these are different notions of infinite sum, and the
classical exponential does not contradict the Hahn theorem. The coefficient
field may be $\R$, $\C$ or any field of characteristic zero; neither algebraic
nor real closedness is used, no completeness of a real-valued norm is used,
and $\Gamma$ needs neither an order unit nor a countable cofinal subset.
Strong summability supplies exactly the finite initial layer the argument
needs.

\subsection{An exact finite workflow}\label{hol:nl:sub:workflow}

For first-order input $P$ the reduction is as follows. Combine like monomials
and remove zero coefficients. If no $Y$ or $Y'$ remains, decide the polynomial
identity directly. Otherwise compute $d_P$, $w_P$, $I_P$ and $\kappa_P$ from
the finite term list, find the positive integer roots of $I_P$ in the chosen
exact representation of the coefficients, take the largest candidate degree,
substitute a generic polynomial of that degree, and solve the resulting
finite algebraic system. This procedure is complete for \emph{entire}
solutions whenever its finite coefficient operations and algebraic solving are
available. It is not a general algorithm for arbitrary surreal names, it does
not enumerate the local formal solutions, and it does not decide arbitrary
strong-support membership. A negative answer after an exhaustive exact solve
is a genuine nonexistence certificate for entire solutions; failure of a
numerical solver is not.

For a nonpolynomial candidate formal solution a second workflow builds the
certificate \eqref{hol:nl:eq:cert1}--\eqref{hol:nl:eq:cert3} from one
sufficiently high nonzero coefficient and finitely many low-index ones; a
claimed certificate must supply their exact leading valuations. The shipped
program checks finite algebraic identities and sample certificates; it does
not decide infinitary summability.

\subsection{Predecessors, repository comparison and priority}
\label{hol:nl:sub:predecessors}

\paragraph{Inherited ideas.}
Hahn fields, the positive-support summation lemma, Conway normal forms, Gauss
initial forms, Euler eigenvalues and indicial-polynomial reasoning are
established ideas and are not relabelled as new. The organizing method is
related in spirit to maximum-term and central-index arguments in the
Wiman--Valiron tradition, surveyed by Hayman~\cite{hayman}; here the finite
residue polynomial gives an exact identity rather than an analytic
approximation, and no real-valued growth estimate is imported. Source~10
consulted Hayman's survey through its publisher metadata and abstract only,
and uses no theorem from its full text.

Newton-polygon methods for nonlinear equations have a substantial history.
Cano and Fortuny Ayuso construct initial and indicial polynomials for
nonlinear $q$-difference equations and obtain support and growth results for
generalized-series solutions~\cite{cano}. Their operator and setting are not
those of a formal $z$-series with arbitrary-rank Hahn \emph{coefficients}
evaluated strongly at every workspace scale. The conceptual precedent is
important, and no invention of differential Newton polygons is claimed.

Non-Archimedean value distribution already restricts nonlinear functional
and differential equations. Hu and Luan work over a complete algebraically
closed field of characteristic zero with a nontrivial non-Archimedean absolute value and prove
Clunie-, Malmquist- and Mokhon'ko-type results~\cite{huluan}; source~10
consulted the parsed text of that paper for its field hypotheses and
Theorems~1.1, 1.3 and~1.5, reports that a remote screenshot service failed so
that no visual inspection of the PDF is claimed, and reproduces no figure or
table from it. Hu and Yang's monograph~\cite{huyang} is a broader
predecessor; only its publisher metadata and description were checked, so no
claim is made that all rank-one special cases were previously unknown.
Giansiracusa and Mereta develop differential enhancements of non-Archimedean
norms and a tropicalization of differential equations~\cite{giansiracusa};
the finite corner here records only the data this obstruction needs, and no
equivalence to their framework and no new tropical fundamental theorem is
asserted. Mantova and Matusinski's survey~\cite{mantovamatusinski} was
consulted for normal forms.

\paragraph{Repository comparison at source~10's pin.}
Source~10 compared its results with the collection at
\texttt{048b72c}; the comparison, with the corrections needed now, is as
follows.
\begin{center}
\small
\begin{tabularx}{\textwidth}{@{}>{\raggedright\arraybackslash}p{0.24\textwidth}
>{\raggedright\arraybackslash}X>{\raggedright\arraybackslash}X@{}}
\toprule
Material & Existing result or scope & Difference in the nonlinear part\\
\midrule
This report, Theorem~\ref{hol:main:ode} & Strongly entire D-finite series are
polynomials; the nonlinear case was the open Question~15.1, now
Question~\ref{hol:q:nonlinear}. & All first-order algebraic equations, a
higher-order corner criterion, and solution-dependent finite certificates.\\
This report, Theorem~\ref{hol:main:multi} & Finite-dimensional span of all
mixed derivatives forces polynomiality. & One algebraic relation in $f$ and
$\vartheta_qf$ suffices for positive weights; no finite-dimensional span.\\
Tail-spans report~\cite{repotail} & Particular Galois-supported constructions
have all-order differential independence, for scalar or fixed-domain analytic
derivations over specified smaller base fields. & A two-jet theorem for every
nonpolynomial strongly entire function over its own full Hahn field.\\
Documentation catalogue \texttt{docs/README.md} & Distinguishes workspaces, strong
summability, coefficient rings and scalar derivations. & Those distinctions
are kept; no global proper-class analytic claim is added.\\
\bottomrule
\end{tabularx}
\end{center}
The tail-span examples are neither weakened nor rediscovered: that report
proves stronger all-order independence for particular constructions under
different coefficient and function-class hypotheses, while the present
universal statement assumes strong entireness and concludes only two-jet
independence. (The coefficient-field part, from sources~11 and~12, now
concludes all-order independence for groups without an order unit; its
comparison with the tail-span report is in
Section~\ref{hol:cf:sub:continuum}.) Source~10 cited the quadratic-exponent series as occurring in
this report and in the catalogue; within the collection it occurs as the
rank-one-Berkovich example~\cite{reporankone}, as
\texttt{found:ex:internal}~\cite{repofoundations}, and as a case of
\eqref{hol:eq:theta}. This is the targeted comparison recorded at the pin,
not an absence claim about every report in the current collection. The
Riccati statement \texttt{diff:cor:riccati} of the
differential-equations report concerns a scalar derivation (compare
Proposition~\ref{hol:nl:prop:riccati}).

Source~10's repository audit covered the root README, the catalogue, contents
listings, this report's definitions, theorems and nonlinear-question passages,
and the README files of the tail-span, analytic-geometry and product-birthday
reports (the last two for context only). It was a targeted comparison, not a
line-by-line search of every report, Lean file, archive, branch or the Git
history, and no absence claim rests on a truncated preview.

\paragraph{Priority and proof status.}
Source~10's literature search, made on 22 September 2026, did not locate the
exact all-rank strong-family first-order theorem, the finite exclusion
certificate, or the positive-weight Euler formulation. Some searches were
poorly targeted, and no exhaustive MathSciNet, zbMATH, full-book or
citation-network review was made, so priority is provisional. The safest
description is a proposed new theorem package and a proved first-order answer
to a repository question, which is not a named historical conjecture; it is
not a certified first discovery of nonlinear non-Archimedean rigidity. The
written proofs are the evidence. The finite checks of
Section~\ref{hol:sec:verification} establish no arbitrary-support theorem;
the main nonlinear rigidity theorems remain unmapped in the current
implementation ledger. No independent human referee report accompanies the
delivered source package.

\section{The thirteen finite check suites}\label{hol:sec:verification}

Each source manuscript shipped an independent exact-arithmetic suite. The
suites of sources~08 and~09 were run on copies and passed, and neither
supersedes the other: they cover overlapping but distinct finite facts of the
linear theory, and they are kept side by side in \texttt{code/} for that
reason. The suite of source~10 checks finite facts of the nonlinear part
only; it was also run on a copy for this merge and reproduced its recorded
counts exactly. The suites of sources~11 and~12 check finite facts of the
coefficient-field part; both were run on copies for this merge and reproduced
their records exactly, apart from the recorded interpreter and package
versions. The suite of source~13 checks finite facts of the theta-descent
part; it was run on a copy for this merge and reproduced its record apart from
the interpreter version and the elapsed time. The suite of source~14 checks
finite facts of the order-three-threshold part; it was run on a copy for this
merge and reproduced its record apart from the interpreter version and the
elapsed time. The suites of sources~18 and~19 check finite facts of the
unit-dilation and Newton-rigidity parts; both were run on copies for this
merge, source~18's reproduced its record exactly and source~19's apart from
the recorded interpreter version. The suites of sources~16 and~17 check
finite facts of the theta-hierarchy and single-scale parts; both were run on
copies for this merge, source~16's reproduced its record apart from the
elapsed time and source~17's apart from the recorded interpreter version. The
suite of source~15 checks finite facts of the polynomial-composition part; it
was run on a copy for this merge and reproduced its record apart from the
recorded interpreter version. The suite of source~20 checks finite facts of
the recurrence part; it was run on a copy for this merge and reproduced its
record exactly, up to line endings. The
suites of sources~08, 09, 10, 12, 16, 18 and~20 use only the Python standard
library, integers, and \texttt{fractions.Fraction}; the suites of sources~11,
13, 14, 15, 17 and~19 use exact symbolic and rational arithmetic in SymPy, the
only third-party package any suite needs. None uses floating-point
mathematics or network access; those of sources~13, 14 and~16 record their
elapsed time.

\paragraph{\texttt{code/08-finite-recurrences-verify.py} --- 3{,}705 checks,
0 failures.}
It emphasises the operator-to-recurrence conversions in all three formats and
the higher-rank ordered-group arithmetic.
\begin{center}
\begin{tabularx}{\textwidth}{@{}X r@{}}
\toprule
Finite calculation & Checks\\
\midrule
Differential operator to coefficient recurrence & 648\\
Dilation operator to coefficient recurrence & 336\\
Mixed operator to coefficient recurrence & 276\\
Unit residues, parity, and finite exceptional cancellation & 400\\
Partial theta identities and finite quadratic estimates & 552\\
Higher-rank lexicographic valuation examples & 1{,}292\\
Sparse Airy recurrence and finite escape edges & 201\\
\midrule
Total & 3{,}705\\
\bottomrule
\end{tabularx}
\end{center}
The operator tests compare the grouped recurrence formulas against direct
finite polynomial differentiation, dilation and multiplication; the theta
tests compare Laurent monomial exponents exactly, including the $q^{-1}$-to-$q$
conversion of \eqref{hol:eq:inverseq}; the unit tests retain the exceptional
index $n=7$ rather than silently checking only generic values; and the
ordered-group checks use lexicographic integer tuples, not a real-valued
surrogate norm. This script takes a \emph{required} \verb|--output| path and
exits nonzero without one, which is deliberate: it prevents a re-run from
overwriting the delivered record.

\paragraph{\texttt{code/09-entire-hahn-holonomic-verification.py} --- 2{,}043
checks, 0 failures.}
It emphasises the boundary behaviour of the exclusion bounds and the exact
theta domain.
\begin{center}
\begin{tabularx}{\textwidth}{@{}X r@{}}
\toprule
Finite calculation & Checks\\
\midrule
Theta homogeneous equation, positive and negative valuation & 92\\
Theta first-order coefficient recurrence & 45\\
Dilation and differential forward recurrence conversions & 72\\
Unit orbit: parity, exceptional cancellation, residue polynomial & 504\\
Rank-two noncofinal affine minimum & 200\\
Theta excluded upper-scale argument & 200\\
Theta included coarse boundary & 170\\
Theta cofinal-scale growth & 120\\
Finite-coordinate entire witness growth & 480\\
Sparse recurrence exact identity & 50\\
Escape boundary: constant leading exponent, strict exterior descent & 50\\
Formal exponential recurrence & 60\\
\midrule
Total & 2{,}043\\
\bottomrule
\end{tabularx}
\end{center}
Two of these rows exist only in this suite and check the two boundary
phenomena that the report insists are different: ``theta included coarse
boundary'' checks that a zero coarse valuation coexists with positive growth
in the original lower-scale coordinate, which is
Section~\ref{hol:sub:coarseboundary}, and ``escape boundary'' checks the
constant-leading-exponent case that
Lemma~\ref{hol:lem:nonincreasing} rules out, which is what makes the excluded
region in Corollary~\ref{hol:cor:periodic} closed. Its \verb|--output|
argument is optional. Seeds are recorded in the shipped records.

\paragraph{\texttt{code/10-nonlinear-rigidity-verify.py} --- 2{,}113 checks,
0 failures.}
It checks the finite algebra of
Sections~\ref{hol:nl:sec:setting}--\ref{hol:nl:sec:consequences}.
\begin{center}
\begin{tabularx}{\textwidth}{@{}X r@{}}
\toprule
Finite calculation & Checks\\
\midrule
Finite Hahn initial layers and exclusion inequalities & 450\\
Automatic first-order nondegeneracy of the corner & 240\\
Univariate falling-factorial corner identities, orders $0$ to $3$ & 480\\
Lexicographic higher-rank ordered-group exclusion bounds & 500\\
Weighted Euler homogeneous-part identities & 360\\
Worked examples: Riccati, degree cutoff, resonance, Painlev\'e, vanishing
corner, Euler solution space & 83\\
\midrule
Total & 2{,}113\\
\bottomrule
\end{tabularx}
\end{center}
The group names in the record \path{data/10-nonlinear-rigidity-verification.json}
are \path{finite_hahn_certificate}, \path{first_order_nondegeneracy},
\path{higher_order_corner}, \path{ordered_group_radius},
\path{weighted_euler} and \path{worked_examples}, in the order of the
table; ``radius'' there is source~10's word for an exclusion bound. The
default seed is $20260922$. The \verb|--output| argument is required, and the
program refuses to overwrite an existing file. The vanishing-corner example
is tested for the monomial degrees $0,\ldots,20$ only. The counts are program
checks, not independently proved theorems.

\paragraph{\texttt{code/11-coarsening-differential-rigidity-verify.py} ---
8{,}673 cases in 8 groups, 0 failures.}
It checks finite facts of source~11's material and needs SymPy
(\path{data/11-coarsening-differential-rigidity-requirements.txt} pins
\texttt{sympy>=1.12,<2}; the record was made with Python~3.13.5 and
SymPy~1.14.0, the merge's rerun with Python~3.14.4 and SymPy~1.14.0).
\begin{center}
\begin{tabularx}{\textwidth}{@{}X r@{}}
\toprule
Group in \path{data/11-coarsening-differential-rigidity-verification.json}
& Cases\\
\midrule
\path{singular_tail_coefficient_isolation}: Taylor identity \eqref{hol:cf:eq:taylor}, orders $1$--$3$ & 12\\
\path{vanishing_corner_indicial_polynomial}: \eqref{hol:cf:eq:eulermultiplier} for $0\le m\le9$ & 10\\
\path{painleve_local_formal_solutions}: 24 coefficient equations, three initial conditions & 72\\
\path{common_factor_reduction_models}: numerator and denominator reduction & 3\\
\path{binary_prefix_coding}: code ranges and all pairs of length-seven branches & 8{,}136\\
\path{finite_reverse_lexicographic_tests}: growth and loss after an exterior scale & 434\\
\path{generator_value_cancellation}: Remark~\ref{hol:cf:rem:cancellation} & 1\\
\path{boundary_equations}: rational Riccati solution, characteristic~$p$ & 5\\
\midrule
Total & 8{,}673\\
\bottomrule
\end{tabularx}
\end{center}
Most of the count is prefix-pair comparisons; it is not a count of
independent theorems. The \verb|--output| argument is optional and defaults
to \path{verification.json} in the current directory, so a bare run can
overwrite a file of that name; use an explicit new path.

\paragraph{\texttt{code/12-coefficient-field-rigidity-verify.py} --- 9{,}308
checks, 0 failures.}
It checks finite facts of source~12's material with the standard library
only.
\begin{center}
\begin{tabularx}{\textwidth}{@{}X r@{}}
\toprule
Category in \path{data/12-coefficient-field-rigidity-verification.json} &
Checks\\
\midrule
coefficient recovery; formal solution coefficients & 387; 637\\
linearized recurrence identity; mixed derivative multiplier scaling & 389; 59\\
derivative-before-dilation convention \eqref{hol:cf:eq:jet} & 405\\
degenerate Euler multiplier \eqref{hol:cf:eq:eulermultiplier} & 1{,}500\\
exponential-polynomial annihilator & 257\\
distinct monomial exponent vectors & 256\\
exterior descending leading values; successive Archimedean levels & 462; 1{,}400\\
partial-theta exponent identity; partial-theta finite escape bounds & 200; 2{,}706\\
positive-characteristic derivative obstruction & 500\\
torsion even-function identity & 150\\
\midrule
Total & 9{,}308\\
\bottomrule
\end{tabularx}
\end{center}
The record also lists thirteen formal recurrence cases with their offsets
and sampled zero multipliers; the sampled zeros are not an effective
Skolem--Mahler--Lech bound. The \verb|--output| argument is required and an
existing file is refused unless \verb|--force| is given.

\paragraph{\texttt{code/13-theta-descent-verify.py} --- 1{,}622 checks, 0
failures.}
It checks finite facts of the theta-descent part with exact rational
arithmetic in SymPy (\path{data/13-theta-descent-requirements.txt} pins
\texttt{sympy>=1.12,<2}; the record was made with Python~3.13.5 and
SymPy~1.14.0 at precision $p^{128}$, the merge's rerun with Python~3.14.4 and
SymPy~1.14.0).
\begin{center}
\begin{tabularx}{\textwidth}{@{}X r@{}}
\toprule
Group in \path{data/13-theta-descent-verification.json} & Checks\\
\midrule
Chebyshev closed coefficients \eqref{hol:td:eq:chebcoef}; Chebyshev
reciprocal substitution \eqref{hol:td:eq:chebsub} & 30; 30\\
Taylor coefficient extraction: the coefficients of $C_n$ entering
\eqref{hol:td:eq:a0}--\eqref{hol:td:eq:am} & 495\\
Theta shift exponents \eqref{hol:td:eq:shift} & 81\\
Jacobi product $p$-coefficients \eqref{hol:td:eq:product} & 128\\
Third-order ODE $p$-coefficients \eqref{hol:td:eq:cleared} & 128\\
Exact jet identities: \eqref{hol:td:eq:jetpoly} and the system
\eqref{hol:td:eq:system2} & 6\\
Root-factor cancellations; root and rounding leading terms
\eqref{hol:td:eq:valder}--\eqref{hol:td:eq:valfloor} & 16; 48\\
Finite local-linearization leading terms \eqref{hol:td:eq:linearization} &
560\\
Lexicographic exclusion checks (Example~\ref{hol:td:ex:lex}) & 100\\
\midrule
Total & 1{,}622\\
\bottomrule
\end{tabularx}
\end{center}
The record's group names and its field \texttt{q\_precision\_exclusive} use
source~13's letter $q$ for the nome called $p$ here. The product and the
cleared equation have zero residual in every $p$-coefficient below $p^{128}$,
each checked as an exact polynomial in the variable; the Lambert constants use
the nome $p^2$. The \verb|--output| argument is required and an existing file
is overwritten, so pass a new path; \verb|--precision| defaults to $64$, and
the record uses $128$. For this merge the identities used by
Lemmas~\ref{hol:td:lem:substitution} and~\ref{hol:td:lem:psi3} and
Proposition~\ref{hol:td:prop:corner}, the expansions of
Section~\ref{hol:td:sub:ode}, the rational forms of
Theorems~\ref{hol:td:thm:ode} and~\ref{hol:td:thm:field}, and the cleared
equation and product modulo $p^{50}$ were also checked with an independent
SymPy script, which is not shipped; it also confirms that with the nome $p$ in
place of $p^2$ the residual of the cleared equation is nonzero.

\Needspace{24\baselineskip}
\paragraph{\texttt{code/14-order-three-threshold-verify.py} --- 600 checks, 0
failures.}
It checks finite facts of the order-three-threshold part with exact rational
and polynomial arithmetic in SymPy
(\path{data/14-order-three-threshold-requirements.txt} pins
\texttt{sympy>=1.13,<2}; the record was made with Python~3.13.5 and
SymPy~1.14.0, the merge's rerun with Python~3.14.4 and SymPy~1.14.0).
\begin{center}
\begin{tabularx}{\textwidth}{@{}X r@{}}
\toprule
Group in \path{data/14-order-three-threshold-verification.json} & Checks\\
\midrule
\path{symbolic_jets}: the coefficient $(z^2-4)^3Y_0^4$ of $Y_3^2$; after
\eqref{hol:ot:eq:eulersub}, $z$-degree six, the form \eqref{hol:ot:eq:thetatop}
and its residue \eqref{hol:ot:eq:thetainitial}; the binomial gap identity
$(\mathsf g^2-1)\mathsf v^2$; a second-order endpoint example & 6\\
\path{endpoint_cancellation}: removal of $V^{i_0}$, $0\le i_0\le4$, and the
resulting endpoint polynomial, as in Lemma~\ref{hol:ot:lem:endpoints} & 10\\
\path{finite_initial_classification}: all $242$ nonzero polynomials of degree
at most four with coefficients in $\{-1,0,1\}$ against
Theorem~\ref{hol:ot:thm:classification} & 242\\
\path{chebyshev}: identities \eqref{hol:td:eq:chebsub}
and~\eqref{hol:td:eq:chebcoef} & 72\\
\path{newton_corners}: at $\delta_n=(2n+1)\mu$, $0\le n\le80$, the active
indices $\{n,n+1\}$ and the value $-n(n+1)\mu$ for $\mathcal T_p$ & 162\\
\path{positive_characteristic}: derivative zero of $\sum z^{\mathfrak pn}$ modulo
$\mathfrak p\in\{2,3,5,7,11\}$ & 5\\
\path{theta_product_truncation}; \path{theta_ode_truncation}: the product
\eqref{hol:td:eq:product} and the cleared equation \eqref{hol:td:eq:cleared}
through $p^{47}$ & 48; 48\\
\path{theta_coefficient_values}: \eqref{hol:td:eq:coeffval} & 7\\
\midrule
Total & 600\\
\bottomrule
\end{tabularx}
\end{center}
The record's group names use source~14's word ``corner'' for a break and its
letter $q$ for the nome called $p$ here; its field
\texttt{cutoff\_exclusive} is $48$, so every $p$-coefficient below $p^{48}$ of
the product and of the cleared equation is compared as an exact polynomial in
$z$. The initial-form enumeration admits $26$ of the $242$ polynomials and is
not the proof of the classification. The script takes no arguments: it
writes \path{data/verification.json} in the parent of its own directory. In
this layout that is a new file beside the delivered record
\path{data/14-order-three-threshold-verification.json}, which it does not
touch, but a second run overwrites the first; run it on a copy.

\Needspace{24\baselineskip}
\paragraph{\texttt{code/18-unit-dilation-verify\_finite.py} --- 16{,}579
checks, 0 failures.}
It checks finite facts of the unit-dilation part with the standard library
only (exact integers and \texttt{fractions.Fraction}; fixed seed $20260923$;
Python~3.10 or later, as the source's README states).
\begin{center}
\begin{tabularx}{\textwidth}{@{}X r@{}}
\toprule
Category in \path{data/18-unit-dilation-finite_verification.json} & Checks\\
\midrule
\path{binomial_tail}: the tails of Example~\ref{hol:ud:ex:cancellation}, $n\le40$,
$d\le13$ & 975\\
\path{parity_valuation}: Example~\ref{hol:ex:unitperiod}, $1\le n\le100$ & 200\\
\path{theta_identity}; \path{differentiated_theta}: the coefficients of
\eqref{hol:eq:thetafirst} and \eqref{hol:ud:eq:thetamixed} for five rational
parameters & 205; 200\\
\path{two_prescribed_dilations}: the coefficients of
\eqref{hol:ud:eq:twoprescribed} for rational multipliers & 600\\
\path{operator_output_formula}; \path{reindexed_shift_formula}:
\eqref{hol:ud:eq:termcoefficient} and \eqref{hol:ud:eq:cj}--\eqref{hol:ud:eq:Rjl}
against direct application, 160 random operators & 3{,}187; 3{,}187\\
\path{finite_independence_matrix}: nonzero confluent Vandermonde determinants
& 12\\
\path{frobenius_boundary}: $\binom{\mathfrak p^a}i\equiv0\pmod{\mathfrak p}$ for
$0<i<\mathfrak p^a$, $\mathfrak p\le7$, $a\le3$ & 12\\
\path{escape_inequality}: $B+j\delta<0$ for $\delta=-B-g$, small integers &
1{,}440\\
\path{unbounded_polynomial_degree}: the coefficients of $(\sigma_2-2^a)z^a$,
$a,n\le80$ & 6{,}561\\
\midrule
Total & 16{,}579\\
\bottomrule
\end{tabularx}
\end{center}
Rational parameters check finite coefficient identities, not entireness. Two
categories check only elementary integer arithmetic: the $6{,}561$ checks of
\path{unbounded_polynomial_degree} compare an expression that is zero by
construction, and the $1{,}440$ checks of \path{escape_inequality} verify
$B+j\delta<0$ for small integers; together they are $8{,}001$ of the
$16{,}579$, which is why the count is no measure of proof coverage, as the
source says. The \verb|--output| argument is optional: without it the program
only prints its record; with it, it writes the given path, creating
directories and overwriting an existing file, so pass a new path. The record
names no interpreter version.

\Needspace{24\baselineskip}
\paragraph{\texttt{code/19-newton-rigidity-verify.py} --- 259 checks, 0
failures.}
It checks finite facts of the Newton-rigidity part with exact polynomial
arithmetic in SymPy (\path{data/19-newton-rigidity-requirements.txt} pins
\texttt{sympy==1.14.0}; the record was made with Python~3.13.5 and
SymPy~1.14.0, the merge's rerun with Python~3.14.4 and SymPy~1.14.0).
\begin{center}
\begin{tabularx}{\textwidth}{@{}X r@{}}
\toprule
Groups in \path{data/19-newton-rigidity-verification.json} & Checks\\
\midrule
\path{conic_homogeneous_rank}, \path{conic_kernel_dimension}: the conic kernel
of Remark~\ref{hol:nr:rem:conic}, degrees $1$--$9$ & 18\\
\path{twisted_cubic_gradient_injectivity}: Lemma~\ref{hol:nr:lem:secant} as a
rank computation, degrees $1$--$3$ & 3\\
\path{first_polar_double_factor}, \path{first_polar_remaining_degree}:
\eqref{hol:nr:eq:polarfactor} on a basis of curve-vanishing forms & 26\\
\path{secant_jacobian}, \path{cubic_polar_example}:
\eqref{hol:nr:eq:jacobian} and Example~\ref{hol:nr:ex:cubic} & 2\\
\path{conic_lowest_cross_term}: the coefficient $a_ma_n(n-m)^2$ of
Remark~\ref{hol:nr:rem:conic} & 60\\
\path{quartic_resultant}, \path{quartic_singular_on_twisted_cubic}:
\eqref{hol:nr:eq:resultant} and $\operatorname{Pol}_{\mathsf H_3}=0$ & 5\\
\path{quartic_binomials}, \path{fourth_order_binomials}: doubling binomials
solve $\mathsf H_3$ and $\mathsf Q_4$, $n\le12$ & 72\\
\path{fourth_order_pair_kernel}, \path{hidden_residual_z5}:
\eqref{hol:nr:eq:pairkernel} and \eqref{hol:nr:eq:residual} & 2\\
\path{quadratic_newton_transition}, \path{doubling_newton_transition}: finite
windows of the breaks of $G_1$ and $F_{\mathrm{sp}}$ & 51\\
\path{first_polar_lowest_coefficient}: \eqref{hol:nr:eq:bottom} for four forms
& 20\\
\midrule
Total & 259\\
\bottomrule
\end{tabularx}
\end{center}
The script takes no arguments and always writes \path{data/verification.json}
in the parent of its own directory, recording the interpreter and SymPy
versions. In this layout that is the same file that source~14's script
writes, a new file beside both delivered records, overwritten by each run;
run it on a copy. The source's \path{code/19-newton-rigidity-build.py} was
written for the delivery layout: it typesets an \texttt{article.tex} in its own
directory \texttt{code/}, where there is none, and its \verb|--verify| option
calls a \texttt{code/verify.py} relative to that directory, which is not
present. The shipped \path{data/19-newton-rigidity-build_audit.json} describes
the 24-page source manuscript, whose source and PDF, hashed there, are not
shipped.

\Needspace{24\baselineskip}
\paragraph{\texttt{code/16-theta-hierarchy-verify.py} --- exact comparisons,
status \texttt{PASS}.}
It checks finite facts of the theta-hierarchy part with the standard library
only, in $\Z[z][p]/(p^{128})$ with exact rationals; its record lists the
comparisons by kind and gives no single total.
\begin{center}
\begin{tabularx}{\textwidth}{@{}X r@{}}
\toprule
Field in \path{data/16-theta-hierarchy-verification.json} & Count\\
\midrule
Chebyshev recurrence against the closed formula \eqref{hol:td:eq:chebcoef} &
30\\
leading Taylor terms of the truncated descent (Proposition~\ref{hol:td:prop:coeff})
& 12\\
Jacobi product \eqref{hol:td:eq:product} against the descent, $p$-coefficients
below $p^{128}$ & 128\\
cleared order-three identity (source~16's Tate form), $p$-coefficients below
$p^{128}$ & 128\\
Gauss profile of Theorem~\ref{hol:th:thm:profile} and its bounds, rational
cases & 852\\
dominance in Example~\ref{hol:th:ex:two}, exact $\Q(\sqrt2)$ comparisons & 65\\
unique squarefree wedge assignment, $N=1,\ldots,8$ & 8\\
two-function constant-derivative model of the grid: candidates, nonzero jet
determinants & 81; 50\\
\bottomrule
\end{tabularx}
\end{center}
The last row is an illustrative finite specialization, not the compression
theorem, and a vanishing residual modulo $p^{128}$ checks normalization and
finite algebra, not an infinite identity. The script's \verb|--output| defaults
to \path{verification.json} in the current directory and overwrites it;
\verb|--precision| defaults to $128$, the value of the record; the record
names the elapsed time ($0.496$\,s) but no interpreter version. The shipped
\path{code/16-theta-hierarchy-Makefile} names the delivery paths
\texttt{article.tex} and \texttt{verify.py}: its default target would run
\texttt{latexmk} on whatever \texttt{article.tex} is in the current directory,
and its \texttt{verify} target calls a \texttt{verify.py} that is not present
under that name. \path{data/16-theta-hierarchy-build_audit.json} describes the
25-page source manuscript, which is not shipped.

\Needspace{24\baselineskip}
\paragraph{\texttt{code/17-single-scale-verify.py} --- 8{,}608 checks in 21
groups, 0 failures.}
It checks finite facts of the single-scale part with exact integer and sparse
polynomial arithmetic and SymPy (\path{data/17-single-scale-requirements.txt}
pins \texttt{sympy==1.14.0}; the record was made with Python~3.13.5 and
SymPy~1.14.0, the merge's rerun with Python~3.14.4 and SymPy~1.14.0).
\begin{center}
\begin{tabularx}{\textwidth}{@{}X r@{}}
\toprule
Groups in \path{data/17-single-scale-verification_results.json} & Checks\\
\midrule
\path{height_separation}: Lemma~\ref{hol:fh:lem:sep} in finite windows &
7{,}198\\
\path{coding_scale_cosets}; \path{commuting_jets}; \path{Euler_conversion} &
432; 256; 144\\
\path{intrinsic_monomials}; \path{intrinsic_triangular}:
\eqref{hol:fh:eq:bm}, Lemma~\ref{hol:fh:lem:bmtriangular} & 153; 9\\
\path{finite_valuation_domain}; \path{exterior_domain_obstruction}:
Theorem~\ref{hol:fh:thm:domain} & 124; 16\\
\path{polarization}; \path{polarization_orbits};
\path{polarization_cancellation}: Lemma~\ref{hol:fh:lem:polar},
Example~\ref{hol:fh:ex:polar} & 53; 53; 1\\
\path{binary_prefixes}; \path{private_prefixes} & 46; 32\\
\path{floor_cutoffs}; \path{floor_examples}: Theorem~\ref{hol:fh:thm:floor} &
45; 5\\
\path{omnific_code_support}; \path{jet_block_rank}; \path{finite_patterns} &
16; 12; 6\\
\path{coefficient_witnesses}; \path{necessary_threshold};
\path{polynomial_height_obstruction} & 5; 1; 1\\
\midrule
Total & 8{,}608\\
\bottomrule
\end{tabularx}
\end{center}
Most of the count is one lemma checked in many windows; it is not a count of
theorems. The record also lists three sparse coefficient witnesses (for
instance the coefficient $1$ of $p^{403200}z^{17}$ in
$\mathcal G\vartheta^2\mathcal G-(\vartheta\mathcal G)^2$) and the floor
examples. The \verb|--output| argument is optional and defaults to
\path{verification_results.json} \emph{beside the script}, here in
\texttt{code/}, overwriting any file of that name; pass a new path. The
source's \path{code/17-single-scale-build.py} looks for an \texttt{article.tex}
in its own directory \texttt{code/}, where there is none, and stops with an
error before writing anything; in the delivery layout it wrote into
\texttt{.build/} and replaced \texttt{article.pdf} and
\texttt{build\_audit.json} in place. \path{data/17-single-scale-build_audit.json}
and \path{data/17-single-scale-visual_audit.json} describe the 24-page source
manuscript, which is not shipped.

\Needspace{24\baselineskip}
\paragraph{\texttt{code/15-polynomial-composition-verify.py} --- 2{,}369 checks
in 12 groups, 0 failures.}
It checks finite facts of the polynomial-composition part with SymPy and exact
rational Hahn polynomials over value groups $\Z^2$ (lexicographic), with the
fixed seed $20260923$ (\path{data/15-polynomial-composition-requirements.txt}
pins \texttt{sympy==1.14.0}; the record was made with Python~3.13.5 and
SymPy~1.14.0, the merge's rerun with Python~3.14.4 and SymPy~1.14.0).
\begin{center}
\begin{tabularx}{\textwidth}{@{}X r@{}}
\toprule
Groups in \path{data/15-polynomial-composition-verification.json} & Checks\\
\midrule
\path{universal_encoding}: Theorem~\ref{hol:pc:thm:universal} & 648\\
\path{indicial_leading_terms}; \path{exact_resonance}:
Lemma~\ref{hol:pc:lem:polyleading}, Theorem~\ref{hol:pc:thm:bound} & 445; 208\\
\path{integer_convexity}: Lemma~\ref{hol:pc:lem:convex} & 360\\
\path{euler_projectors}: \eqref{hol:pc:eq:projector} & 228\\
\path{smith_certificates}: Theorem~\ref{hol:pc:thm:linarith} & 150\\
\path{rank_two_operator_profiles}; \path{rank_two_pullback_profiles};
\path{rank_two_gauss_products}: Propositions~\ref{hol:pc:prop:operatorexact},
\ref{hol:pc:prop:pullback}, Corollary~\ref{hol:nl:cor:gauss} & 90; 70; 35\\
\path{boundary_coefficient_checks}: Section~\ref{hol:pc:sub:boundaries} & 75\\
\path{sharp_degree_family}; \path{worked_equations}:
Section~\ref{hol:pc:sub:examples} & 48; 12\\
\midrule
Total & 2{,}369\\
\bottomrule
\end{tabularx}
\end{center}
The record also lists the degree bounds $1$, $3$, $1$ and $2$ of the four
worked equations of Section~\ref{hol:pc:sub:examples} with a bound. The
\verb|--output| argument is optional and defaults to
\path{verification.json} in the current directory, overwriting any file of
that name; pass a new path. The source's
\path{code/15-polynomial-composition-Makefile} was written for its delivery
layout: its default target runs \texttt{pdflatex} three times on an
\texttt{article.tex} in the current directory, rebuilding \texttt{article.pdf}
in place, its \texttt{verify} target calls a \texttt{verify.py} there, and its
\texttt{clean} target deletes auxiliary files with \texttt{rm -f}; do not run it
here. \path{data/15-polynomial-composition-BUILD_REPORT.json} describes the
24-page source manuscript, which is not shipped.

\Needspace{24\baselineskip}
\paragraph{\texttt{code/20-recurrences-verify.py} --- 4{,}605 checks in 12
families, 0 failures.}
It checks finite facts of the recurrence part with the standard library only
(exact integers and \texttt{fractions.Fraction}, finite Laurent polynomials
in $t$, lexicographic pairs for exponents $a\omega+b$; fixed seed $20260923$;
Python~3.10 or later, as the source states).
\begin{center}
\begin{tabularx}{\textwidth}{@{}X r@{}}
\toprule
Family in \path{data/20-recurrences-verification.txt} & Checks\\
\midrule
unit coefficient formula: \eqref{hol:rc:eq:coefficient} against direct
expansion, three cases, $n\le18$, exponents $-3$ to $24$ & 1{,}596\\
residual parity valuation: Example~\ref{hol:ex:unitperiod}, $1\le n\le80$ &
80\\
principal-unit cancellation: the case $d=2$ of
Example~\ref{hol:ud:ex:cancellation}, $n\le44$ & 45\\
unbounded coordinate transients: $\binom nd=0$ exactly for $n<d$
(Example~\ref{hol:rc:ex:transient}), $d\le18$, $n\le44$ & 810\\
higher-rank affine comparison: $\min\bigl((1,0),(0,n)\bigr)=(0,n)$,
$n\le200$ & 201\\
omnific coefficient coding: Theorem~\ref{hol:rc:thm:coding} for a rational
sequence and for the squares, finite windows & 281\\
Catalan algebraic identity; Catalan power coefficients; declining omnific
support: Example~\ref{hol:rc:ex:catalan} & 40; 88; 88\\
partial-theta operator: the coefficients of \eqref{hol:eq:thetahom} as
Laurent polynomials in the parameter, $n\le100$ & 101\\
mixed-operator coefficient indexing: \eqref{hol:ud:eq:termcoefficient} against
direct application, 35 random operators with multipliers $1,2,3$ & 875\\
positive-characteristic boundary: Example~\ref{hol:rc:ex:charp} for
$\mathfrak p\le7$, $n\le100$ & 400\\
\midrule
Total & 4{,}605\\
\bottomrule
\end{tabularx}
\end{center}
Integer multipliers check coefficient indexing, not entireness, and a finite
truncation of $\lambda_-$ or of $\mathfrak w_{\mathbf b}$ checks only the
coefficients that the omitted tail cannot reach. Four families check facts
that hold by construction: the $810$ binomial-vanishing checks, the $201$ and
$88$ lexicographic comparisons, and the $90$ positivity checks inside the
coding family; together they are $1{,}189$ of the $4{,}605$, so the count is no
measure of proof coverage, as the source says. The script takes no arguments
and prints its record to standard output only; the delivered record is that
output (the merge's rerun, under Python~3.14.4, matched it exactly after
normalizing the Windows line endings of the redirected output). The source's
\path{code/20-recurrences-build.sh} was written for its delivery layout: it
changes to its own directory, runs \texttt{pdflatex} three times on an
\texttt{article.tex} there and then \texttt{python3 verify.py >
verification.txt}, overwriting that record by redirection. Here its directory
is \texttt{code/}, which has neither file: run in this layout it stops at the
first \texttt{pdflatex} call and leaves \texttt{texput.log} in \texttt{code/};
do not run it here. Source~20 shipped no build record, and its 22-page
manuscript is not shipped.

\paragraph{What the checks do not do.}
No suite proves any infinite theorem. No finite prefix can prove strong
summability of an infinite family, cofinality of $(e_n)$, noncofinality of
$(n\epsilon)$, the existence of an infinite escape path, the eventual orbit
statements, the generic-line theorem, or the absence of every possible
annihilating equation. No infinite Hahn sum is evaluated, or numerically
approximated, by any of the programs. In particular a finite truncation of the
explicit series \eqref{hol:eq:explicitF} is a polynomial and is D-finite, so no
finite computation can witness its nonholonomicity; for the same reason no
finite computation can witness the two-jet independence of
Corollaries~\ref{hol:nl:cor:theta} and~\ref{hol:nl:cor:infiniterank}.
Source~10's suite in particular does not test arbitrary infinite supports,
cofinality, all Hahn elements, the nonexistence of all annihilating
equations, or historical novelty. The suites of sources~11 and~12 do not
prove the coefficient-field theorem in general, rationality descent for
infinite series, the Skolem--Mahler--Lech theorem, infinite support well
ordering, cofinality, or any algebraic independence statement, and no finite
computation can witness the independence of all jets of~$F$. Source~13's
suite does not verify strong summability of infinite supports, minimal
differential order or algebraic independence, the no-order-unit obstruction,
completeness of the infinite zero set, or priority, as its record states; a
finite truncation of $\mathcal T_p$ is a polynomial. Source~14's suite does
not establish the quantifiers over arbitrary Hahn supports, ordered groups,
scales or formal solutions, and its enumeration of initial polynomials is not
the classification proof. Source~18's suite does not prove well ordering of
arbitrary supports, strong summability at all field elements, the
Skolem--Mahler--Lech theorem, eventual bounds for exceptional zeros, the
nonexistence of an annihilator, or novelty; source~19's suite certifies no
infinite-support argument, no cofinality statement, nothing over arbitrary
fields, and no novelty, and a finite truncation of $F_{\mathrm{sp}}$ or of
$G_1$ is a polynomial. Source~16's suite proves no summability of infinite
supports, no algebraic independence of the full functions, not the
compression theorem, and no minimal order or cardinality. Source~17's suite
tests finite windows of the band lemma and finite examples; it quantifies over
no supports, equations or continuum families. Source~15's suite does not
establish summability of arbitrary infinite supports, the weighted theorem, the
universal encoding for all polynomials, or novelty; its Hahn profiles are
computed for polynomials over $\Z^2$ only. Source~20's suite does not prove
the Hahn support lemma, the Skolem--Mahler--Lech theorem, the infinite-support
minimum argument, Noetherian compression or the proper-class localization; its
coding checks are finite windows. What the suites
can detect is sign,
indexing, conversion and finite valuation-comparison mistakes. Everything else
rests on the written arguments.

The mathematical dependency chain is short.
\begin{longtable}{@{}>{\raggedright\arraybackslash}p{0.33\textwidth}
                    >{\raggedright\arraybackslash}p{0.63\textwidth}@{}}
\toprule
Conclusion & Load-bearing ingredients\\
\midrule
\endhead
Differential rigidity & Integer evaluation lemma; finite shift conversion;
escape chain and exclusion criterion.\\
Unit-dilation rigidity & Residue polynomial; finite binomial leading terms;
escape chain and exclusion criterion.\\
Noncofinal nonunit rigidity & Eventual order of finitely many affine values;
noncofinality witness; escape chain and exclusion criterion.\\
Coarsened exclusion & The same affine values, projected to
$\Gamma/\Delta_{\abs{v(q)}}$; strictness of the projected descent.\\
Positive dilation existence & Support monoid certificate and the eventual
strict increase of $h_n$; or, in the order-unit case, quadratic leading
valuations and the cofinal-valuation criterion.\\
Several-variable rigidity & Countable polynomial avoidance; finite-block
regrouping; rational differential system on a generic line.\\
First-order nonlinear rigidity & Hahn-family algebra and its residue
homomorphism; large active degrees; exact corner identity; uniqueness of a
corner term for each power of $Y'$.\\
Higher-order corner criterion & The same, with $\chi_P\ne0$ assumed.\\
Positive-weight Euler rigidity & The same algebra in $m$ variables; finite
bounded-weight sets; the eigenvalue identity of $\vartheta_q$.\\
All-order rigidity without an order unit & Minimal relation and its
linearization; one finitely generated coefficient field;
valuation-rank inequality; principal convex bound; exterior obstruction, or
coarsened reduction.\\
Mixed-jet independence & The same, with an exponential-polynomial multiplier
and the imported Skolem--Mahler--Lech theorem.\\
Rationality in $\FE$ & Finite coefficient field; common coarsened reduction of
numerator and denominator; rationality descent.\\
Continuum family & Joint-independence criterion; independent scale
monomials; binary-prefix codes; a separate cardinality bound.\\
Exact domain of $\mathcal T_p$ & Chebyshev coefficients; quadratic
positive-support certificate; descending leading exponents.\\
Order-three equation for $\mathcal T_p$ & Classical triple product; elliptic
pole cancellation; formal identity over $\Q$; ring maps to $\KK$.\\
No equation of order below three & Laurent ring and substitution
homomorphism; theta scaling law; fixed-field, additive-difference and
$u$-adic obstructions; $\psi_3\ne0$.\\
Degree two in the third derivative & Three-jet independence; odd pole order
of a cubic in the transcendental $\mathcal A$.\\
No linear equation for $\mathcal T_p$ & Transcendence of $\psi_1$ over the
elliptic field; highest power in normalized derivatives.\\
Omnific consequences & Linear-factor product; normal forms; local product
linearization.\\
Second-order rigidity (Theorem~\ref{hol:ot:main:second}) & All-scale
criterion; cofinal divisible extension; local finiteness and late breaks;
exterior homogeneous reduction by jet degree, then $z$-degree; nonzero
dehomogenization; removal of $V^{i_0}$; nonconstant logarithmic Euler
derivative in characteristic zero; the finite endpoint lemma.\\
Minimum order three & Second-order rigidity; the entire witness
$\mathcal T_p$ and its equation of order three.\\
Third-order gap certificate & Homogeneous logarithmic substitution; lowest
transverse degree; expansions at $0$ and $\infty$.\\
Theta initial forms and eventual convexity & Residue form of the exterior
equation; rational Riccati reduction; constants of rational differentiation;
exclusion of skipped and triply tied indices; local finiteness.\\
Periodic-affine valuations (Theorems~\ref{hol:ud:thm:unitvalues},
\ref{hol:ud:thm:affine}) & Common support envelope; one period from the
residues; the imported Skolem--Mahler--Lech theorem through
Lemma~\ref{hol:cf:lem:exppoly}; the least coefficient sequence that is not an
identity; eventual order of finitely many affine functions.\\
Relative-scale rigidity and dichotomy (Theorem~\ref{hol:ud:main:mixed}) &
Exact shift polynomials; periodic-affine valuations; noncofinality witness;
escape chain and exclusion criterion; the differentiated theta identity and
Theorem~\ref{hol:thm:thetadomain}.\\
Third order through cubic jet degree (Theorem~\ref{hol:nr:main:third}) &
Late breaks and the exterior form (as for Theorem~\ref{hol:ot:main:second});
bottom and top first-polar coefficients; the double factor; the eventual side
of integral roots through a $\Q$-linear functional; the secant Jacobian.\\
Exact classifications of $\mathsf H_3$ and $\mathsf Q_4$ & Resultant identity
and a common-root parameter; constants of $\vartheta$; the pair kernel
\eqref{hol:nr:eq:pairkernel}.\\
Independent theta scales (Theorem~\ref{hol:th:main:hierarchy}\,(a)) &
Coefficient values $\mu n^2$; exact Gauss profile; Gauss multiplicativity;
eventually affine Gauss values of polynomials; a unique dominant summand.\\
Unbounded order (Theorem~\ref{hol:th:main:hierarchy}\,(b)) & K\"ahler
criterion; invertible coefficient-to-jet change; unique squarefree wedge term;
finite-grid nonvanishing; separant argument for finite composita; the order-three
equation of $\mathcal T_p$.\\
The field $\mathcal D_{\mathrm{DA}}$ & Hamel basis (choice); countable well-ordered
real supports; finite composita.\\
Valuation route to order three & Laurent ring, shift $\mathsf S$, Gauss value of
radius one, eventually constant values of polynomials at integers,
quadratic growth of iterates.\\
Rapid-height freedom (Theorem~\ref{hol:fh:main:single}) & Height separation of
maximal-degree multisets; multicolour polarization; hereditary Zariski
density of $(n,\mathsf b_n)$; full-row-rank pattern change; finite coefficient
minors or constant extension.\\
Exact domain in $\No[i]$, codes, floors & Leading values $\mathsf b_n+n\gamma$;
positive support monoid; disjoint real cosets of powers of $M_\Omega$;
discreteness of $\Oz$.\\
Polynomial-composition rigidity (Theorem~\ref{hol:pc:main:composition}) &
Exact pullback profile; exact operator profile at large active degrees;
Gauss multiplicativity; integer convexity; superlinear escape.\\
Degree bounds and coefficient loci & Exact leading term through $I_L$; finite
integer resonance set; comparison of polynomial degrees.\\
Omnific modules, universality, undecidability & Smith normal form and the
constant-term retraction; Euler projectors; the MRDP theorem and Sun's theorem
(imported).\\
Observable algebra (Theorem~\ref{hol:rc:thm:algebra}) & Common support
envelope; finite binomial expansion with $\binom Ui\in k[U]$ in characteristic
zero.\\
Noetherian compression (Theorem~\ref{hol:rc:thm:compression}) & Observable
algebra; Hilbert's basis theorem; one torsion period for all monomials in the
residues; the imported Skolem--Mahler--Lech theorem through
Lemma~\ref{hol:cf:lem:exppoly}.\\
Surreal recurrence profiles, matrix powers & Algebraic closedness of $\No[i]$;
localization to a set-sized workspace; Jordan form of a finite companion
matrix; periodic-affine valuations; exterior powers.\\
Omnific coding (Theorem~\ref{hol:rc:thm:coding}) & The normal-form support
definition and monomial multiplication only; no Skolem--Mahler--Lech.\\
\bottomrule
\end{longtable}

All supports are sets. All operators are finite sums. All integer indices are
ordinary nonnegative integers. The coefficient valuation is trivial in
characteristic zero; the coefficient field is $\C$ in
Sections~\ref{hol:sec:intro}--\ref{hol:sec:certificates} and any
characteristic-zero field in the nonlinear, coefficient-field, theta-descent,
order-three-threshold, unit-dilation, Newton-rigidity,
polynomial-composition and recurrence parts. The
coefficient-field, unit-dilation and recurrence parts import one external
recurrence theorem without proof, the Skolem--Mahler--Lech theorem, used only
for relations involving at least two distinct dilations or multipliers, or
exponential polynomials with at least two bases. The
polynomial-composition part imports the MRDP theorem and Sun's theorem, used
only for Corollaries~\ref{hol:pc:cor:undecidable} and~\ref{hol:pc:cor:ten}.
The equations use the
external formal variable derivative. None of the proofs asserts convergence in the fine topology of
$\No[i]$, substitutes a real-valued norm for an arbitrary-rank valuation, or
transports a fixed-workspace evaluation to every surcomplex argument.

\section{Originality assessment and questions left open}
\label{hol:sec:status}

\subsection{Established background and proposed contribution}

The D-finite/P-recursive correspondence~\cite{stanley}, ordinary properties of
Hahn fields, the Hahn--Neumann support lemmas and Higman's finite-word
theorem~\cite{higman}, finite polynomial root bounds, the Conway normal-form
identification~\cite{bg,kks}, and partial theta series with their elementary
functional identity~\cite{andrewswarnaar} are established material. The
companion report's cofinality and scalar-extension work, including the group
$\Gamma_\infty$ and the distinction between order units and countable
cofinality, is prior project material~\cite{repoentire}. The rank-one
quadratic-exponent entire example is already in the
collection~\cite{reporankone}. We claim to originate none of these, and we do
not claim to have disproved any published theorem about a different analytic
function class.

The proposed research contribution is their combination into the sharp
arbitrary-rank dilation classification, Theorem~\ref{hol:main:q}: in
particular the necessity argument for all noncofinal scales, the unit case
with root-of-unity residue, the uniform exterior obstruction applying to every
solution of a fixed equation, the quantitative exclusion bounds with their two
different boundary behaviours, and the exact theta evaluation domain for an
arbitrary Hahn parameter. The mixed noncofinal extension
(Theorem~\ref{hol:thm:mixed}), generic-line rigidity
(Theorem~\ref{hol:main:multi}), and the explicit series excluded from every
nontorsion pure dilation equation (Theorem~\ref{hol:thm:explicit}) are further
consequences proved here.

The nonlinear part (source~10) inherits Gauss initial forms, Euler
eigenvalues and indicial-polynomial reasoning, and is related to
central-index, Newton-polygon, non-Archimedean value-distribution and tropical
differential methods (Section~\ref{hol:nl:sub:predecessors}). Its proposed
contribution is the precise arbitrary-rank strong-family formulation:
Theorem~\ref{hol:nl:main:first} with its finite exclusion certificate
(Theorem~\ref{hol:nl:thm:certificate}), the candidate degrees and the
entire-solution locus, the higher-order corner criterion, the positive-weight
Euler theorem, and the non-uniformity of the nonlinear exclusion bound. Its
priority is provisional in the sense recorded there.

The coefficient-field part (sources~11 and~12) inherits classical coefficient
recurrences, the finite-generation principle, the valuation-rank inequality
and the Skolem--Mahler--Lech theorem
(Section~\ref{hol:cf:sub:predecessors}). Its proposed contributions are the
all-order no-order-unit rigidity with its joint-dilation form
(Theorem~\ref{hol:cf:main:entire}), the coefficient-value-rank criterion and
the order-unit dichotomy (Theorems~\ref{hol:cf:main:rank}
and~\ref{hol:cf:main:sharp}), the meromorphic coarsening descent
(Theorem~\ref{hol:cf:thm:descent}) and the full-constant-field continuum
family (Theorem~\ref{hol:cf:thm:continuum}); its priority is equally
provisional.

The theta-descent part (source~13) inherits the Chebyshev identity, Jacobi's
triple product and the theta--elliptic differential identity, all classical,
the bilateral theta domain of the Hahn--Tate report, and the no-order-unit
direction of Theorem~\ref{hol:cf:main:entire}
(Section~\ref{hol:td:sub:predecessors}). Its proposed contributions are the
support-controlled descent to an entire series with its exact domain at every
scale, the exact boundary of Theorem~\ref{hol:td:main:boundary} with the
minimal-order certificate, the differential field and non-D-finiteness, and
the surreal and omnific consequences; its priority is equally provisional.

The order-three-threshold part (source~14) inherits the theta material of
source~13, the no-order-unit direction of Theorem~\ref{hol:cf:main:entire},
the Gauss initial polynomials of source~10 and classical logarithmic
derivatives (Section~\ref{hol:ot:sub:predecessors}). Its proposed
contributions are second-order rigidity for every value group
(Theorem~\ref{hol:ot:main:second}\,(a)) with its finite endpoint lemma, the
exact least order three, the third-order endpoint--gap obstruction, the
classification of theta initial polynomials and eventual Newton convexity;
its priority is equally provisional.

The unit-dilation part (source~18) inherits the support calculus, the escape
mechanism, the partial theta series and the imported Skolem--Mahler--Lech
theorem (Section~\ref{hol:ud:sub:predecessors}). Its proposed contributions
are the periodic-affine valuation theorem for Hahn exponential polynomials,
the bivariate unit-orbit estimate and the answer to
Question~\ref{hol:q:mixedunit}, and the relative-scale dichotomy of
Theorem~\ref{hol:ud:main:mixed}; its priority is equally provisional.

The Newton-rigidity part (source~19) inherits the all-scale criterion, the
corner and break techniques, and second-order rigidity, which it re-derives
independently (Section~\ref{hol:nr:sub:predecessors}). Its proposed
contributions are third-order rigidity through cubic jet degree
(Theorem~\ref{hol:nr:main:third}) with the first-polar, secant and
eventual-side lemmas, the constant-extension and algebraic two-state
consequences, and the exact classifications of two boundary equations; its
priority is equally provisional.

The theta-hierarchy part (source~16) inherits source~13's seed, equation and
minimal order three and the classical theta product and Tate identity
(Section~\ref{hol:th:sub:predecessors}). Its proposed contributions are
independence across scales, bounded polynomial jet compression, unbounded
minimal order, the transcendence degrees of the differentially algebraic
entire field for $\Gamma=\R$, integer sampling and a valuation route to order
three; its priority is equally provisional.

The single-scale part (source~17) inherits Gr\"onwall's single-series gap
criterion, the Berarducci--Mantova derivation, the omnific integers and the
bounded-support descent (Section~\ref{hol:fh:sub:predecessors}). Its proposed
contributions are the exact differential relation ideals of rapid-height
families, their differential transcendence over the full Hahn field at one
order-unit scale, and the surreal, domain and omnific realizations; its
priority is equally provisional.

The recurrence part (source~20) inherits the support lemma, the
Skolem--Mahler--Lech theorem, Hilbert's basis theorem, the escape mechanism and
the partial theta series, and re-derives source~18's valuation theorem
(Section~\ref{hol:rc:sub:predecessors}). Its proposed contributions are the
finitely generated algebra of coefficient observables, Noetherian compression
of coefficient conditions with one universal residual period, the recurrence
profiles over $\No[i]$, $\Oz$ and $\Oz[i]$, the reading of a known omnific
stream as a universal coefficient observable
(Theorem~\ref{hol:rc:main:observables}), and determinantal profiles of matrix
powers; its priority is equally provisional.

\subsection{A result-by-result attribution ledger}

\begin{longtable}{@{}>{\raggedright\arraybackslash}p{0.40\linewidth}
                    >{\raggedright\arraybackslash}p{0.55\linewidth}@{}}
\toprule
Result or ingredient & Status in this report\\
\midrule
\endhead
Hahn arithmetic, strong support sums, surreal normal forms & Standard
background; references and support proofs supplied.\\[4pt]
Higman's finite-word theorem & External standard ingredient; alternatively
import the classical Hahn--Neumann support lemma directly.\\[4pt]
D-finite / polynomial recurrence conversion & Classical; credited to Stanley
and derived explicitly for Hahn coefficients.\\[4pt]
Partial theta series and functional identity & Classical; not a novelty
claim.\\[4pt]
Escape chain and exclusion criterion, with the local-finiteness boundary &
Elementary support-theoretic lemma and quantitative formulation developed
here.\\[4pt]
Eventual affine and periodic orbit valuations & Proved here for arbitrary Hahn
coefficients; the finite residue-polynomial treatment covers near-torsion
units.\\[4pt]
Exact fixed-$q$ order-unit classification & Principal proposed contribution:
arbitrary rank, arbitrary Hahn parameter, both signs of $v(q)$.\\[4pt]
Exact theta evaluation domain & Support-controlled sharpness theorem; related
in spirit to the companion report's scalar-extension results, proved
independently here.\\[4pt]
Uniform exterior obstruction and common degree bound & Proved here; stronger
than eventual termination of each solution separately.\\[4pt]
Universal single-dilation nonholonomicity on no-order-unit workspaces &
Consequence of the classification, with an explicit entire witness.\\[4pt]
Mixed differential--dilation rigidity at noncofinal nonzero scales & Proved
here; the unit-valuation mixed case is answered without an order unit by
source~12 (Theorem~\ref{hol:cf:main:entire}), and for every $\Gamma$, for
linear equations, by source~18
(Corollary~\ref{hol:ud:cor:unitrigidity}).\\[4pt]
Several-variable rigidity and multivariate algebraic rigidity & Proved here;
the generic line is an existence device.\\[4pt]
Gauss initial forms, Euler eigenvalues, indicial polynomials, central-index
and Newton-polygon ideas & Established (source~10); credited to Hayman, Cano
and Fortuny Ayuso, Hu--Luan, Hu--Yang and Giansiracusa--Mereta as
predecessors in different settings.\\[4pt]
First-order nonlinear rigidity (Theorem~\ref{hol:nl:main:first}), finite
exclusion certificate, candidate degrees and entire-solution locus &
Proposed contribution of source~10, over any characteristic-zero coefficient
field; priority provisional.\\[4pt]
Higher-order corner criterion and positive-weight Euler rigidity & Proposed
by source~10; sufficient conditions only; the vanishing-corner case is
settled without an order unit by sources~11 and~12, and with one source~13's
equation of order three has vanishing residue corner and a nonpolynomial
entire solution.\\[4pt]
Solution-dependence of the nonlinear exclusion bound & Proved by source~10
with the explicit Riccati family of
Theorem~\ref{hol:nl:thm:nouniform}.\\[4pt]
Two-jet independence of $\sum t^{n^2}z^n$ and of
\eqref{hol:eq:explicitF} & New consequences (source~10); both series are
prior collection material.\\[4pt]
Classical coefficient recurrences, finite generation, valuation-rank
inequality, Skolem--Mahler--Lech & Established (sources~11 and~12); credited
to Hurwitz, Laohakosol--Kongsakorn--Ubolsri, Bell and Lech; the last imported
without proof.\\[4pt]
All-order rigidity without an order unit, with dilations
(Theorem~\ref{hol:cf:main:entire}) & Proposed by sources~11 and~12 (part~(b)
and the dilations by source~12, part~(c) by source~11); priority
provisional.\\[4pt]
Coefficient-value-rank criterion, order-unit dichotomy, rank-one family &
Proposed by source~12; sufficient conditions and one equivalence, as
stated.\\[4pt]
Coarsening descent, rationality in $\FE$, systems, Painlev\'e~I in $\FE$,
joint-independence criterion, continuum family & Proposed by source~11; the
polynomial half of the descent is the entire-function report's obstruction,
and neither the cardinal nor the almost-disjoint device is new.\\[4pt]
All-order independence of \eqref{hol:eq:explicitF} &
New consequence (sources~11 and~12); the series is prior material.\\[4pt]
Several multiplicatively independent dilations
(Corollary~\ref{hol:cf:cor:several}) & The merge's observation from
source~12's theorem; not claimed by source~12.\\[4pt]
Mixed jets in the fraction field
(Corollary~\ref{hol:cf:cor:meromorphicmixed}) &
Consequence added here from source~12's coefficient theorem and source~11's
descent; neither source states it.\\[4pt]
Chebyshev identity, Jacobi triple product, theta--elliptic differential
identity & Classical (source~13); DLMF, Sebbar and Ohyama credited; no
novelty claim for third-order theta relations.\\[4pt]
Bilateral theta domain & The Hahn--Tate report's theorem, specialized;
reproved by source~13, not claimed anew.\\[4pt]
Exact domain of $\mathcal T_p$, product and zeros, local valuations,
omnific rounding & Proposed by source~13, derived from the classical
product.\\[4pt]
Exact boundary (Theorem~\ref{hol:td:main:boundary}), minimal order three,
differential field, non-D-finiteness & Proposed by source~13; the
no-order-unit direction is Theorem~\ref{hol:cf:main:entire}\,(a); priority
provisional.\\[4pt]
Substitution homomorphism, global proof of $\psi_3\ne0$, corner data of the
cleared equation, clause~(iv) of Theorem~\ref{hol:cf:main:sharp},
Corollary~\ref{hol:td:cor:sharp} & Added by the merge from source~13's
results; the first two write out steps source~13 used.\\[4pt]
Second-order rigidity for every $\Gamma$ and exact least order three
(Theorem~\ref{hol:ot:main:second}), endpoint lemma, endpoint--gap
obstruction, theta initial forms, eventual Newton convexity, two-state systems
and polynomial-denominator quotients & Proposed by source~14; the witness of
order three is source~13's and the no-order-unit direction is
Theorem~\ref{hol:cf:main:entire}\,(a); priority provisional.\\[4pt]
Cofinal scalar extension, all-scale criterion, surreal full-class boundary &
Lemma~\ref{hol:ot:lem:extension} and
Proposition~\ref{hol:ot:prop:fullclass} are source~14's; the criterion is
Proposition~\ref{hol:cf:prop:criterion}, printed once.\\[4pt]
Exterior form as the residue corner in Euler jets; three jets of $G$ and
$\Theta_p$ & Added by the merge (Remark~\ref{hol:ot:rem:corner},
Corollary~\ref{hol:ot:cor:examples}).\\[4pt]
Periodic-affine valuations of Hahn exponential polynomials, the bivariate
unit-orbit estimate, mixed rigidity for noncofinal relative scales, the
relative-scale dichotomy (Theorem~\ref{hol:ud:main:mixed}) & Proposed by
source~18; the Skolem--Mahler--Lech theorem is imported; pure differential,
cyclic pure dilation and nonzero noncofinal mixed rigidity are this report's
earlier results, now special cases; priority provisional.\\[4pt]
Finite-word lemma with proof & The standard minimal-bad-sequence argument,
printed from source~18 (Lemma~\ref{hol:ud:lem:words}); not new.\\[4pt]
Second-order rigidity, re-derived & Source~19's Theorem~A is
Theorem~\ref{hol:ot:main:second}\,(a) of source~14, found independently the
same day; printed once, no priority between them asserted.\\[4pt]
Third order through cubic jet degree (Theorem~\ref{hol:nr:main:third}),
first-polar and factor criteria, secant and eventual-side lemmas,
constant-extension, omnific and algebraic two-state consequences, exact
classifications of $\mathsf H_3$ and $\mathsf Q_4$, the sparse series &
Proposed by source~19; priority provisional.\\[4pt]
First polar inside the endpoint--gap obstruction, a cubic outside
Corollary~\ref{hol:ot:cor:gapcriterion}, jet degree of the theta equation,
comparison of the two-state results & Added by the merge
(Remark~\ref{hol:nr:rem:gaps}, Corollary~\ref{hol:nr:cor:theta},
Section~\ref{hol:nr:sub:constants}).\\[4pt]
Named published open conjecture & None claimed solved; the nonlinear question,
answered in order one, without an order unit in all orders, with an order
unit by an example of order three, and in order two for every $\Gamma$, is a
question of this report, not a historical conjecture; so is the mixed
unit-valuation question answered by source~18.\\[4pt]
Seed, equation and minimal order three of $\mathcal T_p$ reused by source~16;
theta product, Tate identity & Source~13's and classical; credited, not
claimed.\\[4pt]
Independence across scales, exact Gauss profile, noncancellation certificate,
jet compression, unbounded order, $\trdeg$ and $\dtrdeg$ of
$\mathcal D_{\mathrm{DA}}$, integer sampling, valuation route to order three &
Proposed by source~16, for $\Gamma=\R$; priority provisional.\\[4pt]
Single-series gap differential transcendence; normal forms,
Berarducci--Mantova derivation, $\Oz$, bounded-support descent & Classical or
prior (Gr\"onwall via Ostrowski; \cite{bm}; \cite{repobst}); not claimed.\\[4pt]
Exact relation ideals of rapid-height families, pure $D_z$ freedom over the
full Hahn field, Hermite data, $\partial_{\mathrm{BM}}$-jets over the
bounded-support field, exact domain in $\No[i]$, omnific codes and floors &
Proposed by source~17; priority provisional.\\[4pt]
Identity of the Tate and Weierstrass seed equations, the known part of the
order spectrum, three fields of entire series, three routes to order three,
placement of single-scale freedom & Added by the merge
(Section~\ref{hol:th:sub:conventions}, Remarks~\ref{hol:th:rem:ordertwo},
\ref{hol:th:rem:fields}, \ref{hol:th:rem:thirdroute},
\ref{hol:fh:rem:placement}).\\[4pt]
Complex entire Mahler lemma, Mahler's method, normal forms and the
constant-term retraction, Smith normal forms, the MRDP theorem and Sun's
theorem & Classical or prior (\cite{bcr,nishioka,lmfactor,matiyasevich,sun};
\cite{repoodg}); credited, not claimed.\\[4pt]
Strict composition-weight rigidity for every $\Gamma$
(Theorem~\ref{hol:pc:main:composition}), exact-resonance degree bounds,
extension-stable coefficient loci, the matrix version, linear omnific
solution modules, the universal Euler-projector encoding and its
undecidability consequences & Proposed by source~15; priority
provisional.\\[4pt]
Removal of the order-unit dichotomy by a dominant argument of degree at least
two, identification of source~15's indicial polynomials, added Mahler
references & Added by the merge (Remark~\ref{hol:pc:rem:dichotomy},
Sections~\ref{hol:pc:sub:conventions} and~\ref{hol:pc:sub:predecessors}).\\[4pt]
Support envelope, unit-root and periodic-affine valuations, mixed unit
rigidity, re-derived; exact finite-group criterion & Source~20 found these
independently of source~18; printed once as source~18's
(Section~\ref{hol:rc:sub:conventions}); its finite-group criterion and its
convex-scale theorem are special cases of Theorem~\ref{hol:ud:main:mixed},
printed as Corollaries~\ref{hol:rc:cor:finitegroup}
and~\ref{hol:rc:cor:convex}.\\[4pt]
The omnific stream $\sum_m\mathbf b_m\omega^{\omega-m}$ & The omnific-notations
report's (\texttt{onot:eq:guardbasic}); credited by the merge, not claimed.\\[4pt]
One finitely generated observable algebra, Noetherian compression with one
residual period, recurrence profiles over $\No[i]$, $\Oz$ and $\Oz[i]$, the
decreasing omnific recurrence, the coding reading and the failure of a
coefficientwise Skolem--Mahler--Lech theorem, determinantal profiles &
Proposed by source~20 (Theorem~\ref{hol:rc:main:observables}); priority
provisional.\\[4pt]
Two periods, finite generation in the convex-hull criterion, the stream credit,
the exact positive-characteristic formula beside the unit-dilation part &
Added by the merge (Remarks~\ref{hol:rc:rem:periods},
\ref{hol:rc:rem:finitegeneration}, \ref{hol:rc:rem:stream},
Example~\ref{hol:rc:ex:charp}).\\[4pt]
Priority, peer review, formal verification & Priority and independent
refereeing are not certified. Lean covers support and escape lemmas,
exponential and partial-theta domains, polynomial-orbit valuations,
constant-point avoidance, inward stability and torsion covariance.
The full classifications and the theorem packages of the nonlinear,
coefficient-field, theta-descent, order-three-threshold, unit-dilation,
Newton-rigidity, theta-hierarchy, single-scale, polynomial-composition and
recurrence parts remain pending.\\
\bottomrule
\end{longtable}

The targeted literature checks of sources~08 and~09 located the standard
D-finite framework,
classical complex $q$-difference growth theory, $q$-holonomic degree growth in
a rational-function setting, and ultrametric $q$-difference theory, but did
not establish whether an identical arbitrary-rank strong-Hahn formulation
appears elsewhere. Several broad keyword searches returned irrelevant material
and were not treated as evidence of absence; a repository content search for
\texttt{holonomic} returning no matches is likewise not proof that no sentence
anywhere in the repository can be reformulated into one of these results.
Accordingly this is a proposed original result with proofs, not a certified
claim of bibliographic priority. Independent mathematical review is still
needed; machine-checked scope is the partial coverage recorded in
Section~\ref{hol:sec:certificates} and the formalization ledger. The
corresponding records for the nonlinear, coefficient-field,
theta-descent, order-three-threshold, unit-dilation, Newton-rigidity,
theta-hierarchy, single-scale, polynomial-composition and recurrence parts,
with their own search limits, are Sections~\ref{hol:nl:sub:predecessors},
\ref{hol:cf:sub:predecessors}, \ref{hol:td:sub:predecessors},
\ref{hol:ot:sub:predecessors}, \ref{hol:ud:sub:predecessors},
\ref{hol:nr:sub:predecessors}, \ref{hol:th:sub:predecessors},
\ref{hol:fh:sub:predecessors}, \ref{hol:pc:sub:predecessors}
and~\ref{hol:rc:sub:predecessors}.

\subsection{What the nonlinear part does not claim}
\label{hol:nl:sub:nonclaims}

Every limitation stated by source~10, in its manuscript or its two audit
files, is kept; they are collected here. They describe source~10's results;
where the coefficient-field or theta-descent part goes further, a
parenthesis says so. The limitations of sources~11 and~12 are listed
separately as (C1)--(C26) in Section~\ref{hol:cf:sub:nonclaims}, those of
source~13 as (T1)--(T19) in Section~\ref{hol:td:sub:nonclaims}, those of
source~14 as (O1)--(O18) in Section~\ref{hol:ot:sub:nonclaims}, those of
source~18 as (U1)--(U15) in Section~\ref{hol:ud:sub:nonclaims}, and those of
source~19 as (R1)--(R16) in Section~\ref{hol:nr:sub:nonclaims}, those of
source~16 as (Th1)--(Th16) in Section~\ref{hol:th:sub:nonclaims}, those of
source~17 as (F1)--(F19) in Section~\ref{hol:fh:sub:nonclaims}, and those of
source~15 as (P1)--(P18) in Section~\ref{hol:pc:sub:nonclaims}.
\begin{enumerate}[label=\textup{(N\arabic*)}]
\item Only the first-order case of Question~\ref{hol:q:nonlinear} is
answered, not the unrestricted higher-order question. (Sources~11 and~12
answer the case without an order unit in every order; source~13 shows that
with an order unit a nonpolynomial strongly entire series can satisfy an
equation of order three; source~14 answers order two for every $\Gamma$, and
source~19 order three with total jet degree at most three.)
\item Theorem~\ref{hol:nl:main:first} does not say that a nonpolynomial
entire function is differentially transcendental in all orders; independence
of $f$ and $D_zf$ does not exclude a relation involving $D_z^2f$.
(Theorem~\ref{hol:cf:main:entire}\,(a) says it when $\Gamma$ has no order
unit, and Theorem~\ref{hol:td:main:boundary} shows it fails with one;
Theorem~\ref{hol:ot:main:second} adds $D_z^2f$ to the independent jets for
every $\Gamma$.)
\item The criterion $\chi_P\ne0$ is sufficient, not a complete invariant of a
solution set, and the methods of source~10 do not decide equations of order
at least two with vanishing residue corner. Proposition~\ref{hol:nl:prop:degenerate} is not
a nonpolynomial entire solution and does not refute higher-order rigidity.
(Source~13's order-three equation, with $\chi_P=0$, does have one when
$\Gamma$ has an order unit; Proposition~\ref{hol:td:prop:corner}. Source~14
decides order two negatively for every $\Gamma$,
Theorem~\ref{hol:ot:main:second}, and source~19 order three in jet degree at
most three, Theorem~\ref{hol:nr:main:third}.)
\item The corner is a finite selection, not a full differential Newton
polygon; no invention of differential Newton polygons, no equivalence to
tropical differential algebra, and no new tropical fundamental theorem is
claimed.
\item The exclusion bound \eqref{hol:nl:eq:deltazero} is not claimed optimal;
the certificate has no converse; leading valuations alone do not decide
evaluation at a boundary argument.
\item Candidate degree sets are necessary, not realized; their finiteness is
not a uniform algorithm for arbitrary coefficient names, and the workflow of
Section~\ref{hol:nl:sub:workflow} is not an algorithm for arbitrary surreal
names, does not enumerate local formal solutions, does not decide
strong-support membership, and accepts no numerical solver failure as a
certificate.
\item Corollary~\ref{hol:nl:cor:locus} does not say that the number of
solutions is finite, is not a moduli statement over rings with nilpotents,
and covers only the extension class it states.
\item No theorem says that $G$, $F$ or any nonpolynomial series is entire on
the proper class $\No[i]$; entireness and the derivative are relative to one
set-sized workspace.
\item The independence statements concern formal functions; no statement
concerns algebraic independence of a value $f(x)$ from its own ambient field.
\item $D_z$ kills all Hahn coefficients and is not the Berarducci--Mantova
derivation; no scalar-derivation result is claimed.
\item The series $\sum t^{n^2}z^n$ and \eqref{hol:eq:explicitF}, and the
linear rigidity of the latter, are prior material; only their nonlinear
two-jet consequences are new. The tail-span report's all-order independence
results are neither weakened nor rediscovered.
\item Hahn fields, the positive-support summation lemma, Conway normal forms,
Gauss initial forms, Euler eigenvalues and indicial-polynomial reasoning are
established ideas, not relabelled as new; the positive-support lemma is
imported by source~10, and proved in this report by
Lemma~\ref{hol:lem:support}\,(2).
\item The precise all-rank statements are proposed contributions; historical
priority is not certified. The literature search of 22 September 2026 was
targeted and not exhaustive; Hayman's survey and the Hu--Yang monograph were
consulted through metadata only, and no claim is made that rank-one special
cases were previously unknown.
\item The repository comparison at \texttt{048b72c} was targeted: no
line-by-line audit of every report, Lean file, archive, branch or the Git
history, and no absence claim rests on a truncated preview.
\item The answered question is a question of this report, not a named
historical conjecture.
\item Source~10 supplied no Lean formalization or independent referee report
for the nonlinear part. The 2{,}113 finite checks are program checks, not theorems; they
do not test infinite supports, cofinality, all Hahn elements, nonexistence of
annihilating equations or novelty, and no infinite surreal sum is numerically
approximated. A successful PDF build is not a proof check.
\item The sequence arguments behind the support calculus use ordinary
set-theoretic choice.
\end{enumerate}

\subsection{Further mathematical questions}

\begin{question}[Differentially algebraic entire series of minimal order
three]\label{hol:q:nonlinear}
Suppose $\Gamma$ has an order unit. By Theorem~\ref{hol:ot:main:second} no
nonpolynomial $f\in\E$ satisfies a nonzero algebraic differential equation
over $\K(z)$ of order at most two, and by Theorem~\ref{hol:td:main:boundary}
order three occurs. Which nonpolynomial strongly entire series have minimal
differential order three? Are they all accounted for by support-controlled
algebraic descents of elliptic or theta constructions, or are there
fundamentally different families? More broadly, to what extent do finite
algebraic descents of theta or abelian functions, combined with elementary
operations preserving strong entireness, account for the nonpolynomial
strongly entire series that are differentially algebraic over $\K(z)$? The
same questions are open over $\KK=k((t^\Gamma))$ for every field $k$ of
characteristic zero.
\end{question}

\paragraph{Status.}
As posed by sources~08 and~09, the question asked about nonlinear algebraic
differential equations of every order, for every $\Gamma$, and ended by
asking under which hypotheses on $\Gamma$ linear rigidity can be strengthened
to differential algebraic rigidity in all orders, and whether the explicit
series \eqref{hol:eq:explicitF} avoids all equations of order at least two
under a natural coefficient restriction. Source~10 answers its first-order case
negatively, for every nonzero $\Gamma$ and every characteristic-zero
coefficient field (Theorem~\ref{hol:nl:main:first}): in order one no
hypothesis on $\Gamma$ beyond $\Gamma\ne0$ is needed, and
\eqref{hol:eq:explicitF} avoids every first-order equation with no
coefficient restriction (Corollary~\ref{hol:nl:cor:infiniterank}). It also
settles the higher-order equations with $\chi_P\ne0$
(Theorem~\ref{hol:nl:thm:corner}) and the relations between $f$ and a
positive-weight Euler derivative (Theorem~\ref{hol:nl:thm:euler}). After
source~10, order at least two with a vanishing residue corner stayed open.
Proposition~\ref{hol:nl:prop:degenerate} shows that such equations exist and
can have polynomial solutions of unbounded degree, so the study cannot rely on
a uniform degree bound for all polynomial solutions of an arbitrary
higher-order equation; it is not a nonpolynomial entire solution and gives no
negative answer. Source~10 suggests as a next step classifying the polynomial
solutions of a vanishing-corner residue equation beyond their leading monomial
and determining how those constraints propagate between scales.

Sources~11 and~12 then answered the question, negatively, in every order for
every $\Gamma$ without an order unit, whatever the residue corner and for
every characteristic-zero coefficient field
(Theorem~\ref{hol:cf:main:entire}\,(a)); source~11 extends the answer to
quotients of entire series (Theorem~\ref{hol:cf:main:entire}\,(c)). Their
method needs no degree bound: along the solutions of
\eqref{hol:nl:eq:degenerate} its linearization multiplier is $(T-m)^2\ne0$
(Example~\ref{hol:cf:ex:corner}). The last sentence of the question as posed
by sources~08 and~09 is answered positively for \eqref{hol:eq:explicitF} with no coefficient restriction
(Theorem~\ref{hol:cf:thm:Fomega}, source~12; source~11 for
$f_*=F-\omega^{-1}$). Its sentence on hypotheses is answered in one direction:
the absence of an order unit suffices. The question was therefore re-scoped to
groups with an order unit, in the wording recorded below, which merges source~11's
Question~9.1 and source~12's Question~10.1; both sources leave it open and
give no counterexample. There a finitely generated coefficient field can have
values whose convex hull is all of $\Gamma$, and the argument of
Section~\ref{hol:cf:sec:coefficients} stops. The partial theta series
$\Theta_p$ and $G$ are not counterexamples: $\Theta_p$ satisfies a dilation
equation, not a known algebraic differential one, and only the first-order
consequence is known for $G$ (Corollary~\ref{hol:nl:cor:theta}).
Theorem~\ref{hol:cf:main:sharp} shows that the order unit is exactly the
boundary of the \emph{mixed} statement; whether it is also the boundary of the
purely differential one was, at that stage, the question.

As re-scoped after sources~11 and~12, the question read: \emph{Suppose $\Gamma$
has an order unit. Can a nonpolynomial $f\in\E$ satisfy a nonzero algebraic
differential equation $P(z,f,D_zf,\ldots,D_z^sf)=0$ over $\K(z)$ of order
$s\ge2$? By Theorems~\ref{hol:nl:main:first} and~\ref{hol:nl:thm:corner}, such
an $f$ satisfies no nonzero equation of order at most one and no equation
with $\chi_P\ne0$; equivalently, can a nonpolynomial strongly entire function
satisfy an equation whose residue corner polynomial vanishes identically,
without satisfying any other equation to which
Theorem~\ref{hol:nl:thm:certificate} applies? More strongly, is every element
of $\operatorname{Frac}\E$ that is differentially algebraic over $\K(z)$
rational? Is the absence of an order unit necessary for differential algebraic
rigidity in all orders, or can linear rigidity be strengthened to it for every
nonzero $\Gamma$?} Source~13 (the theta-descent part) answers every sentence
of it, for every characteristic-zero coefficient field, including $k=\C$.
With an order unit $\mu$ and $p=t^\mu$, the series $\mathcal T_p$ of
\eqref{hol:td:eq:series} is nonpolynomial and strongly entire and satisfies an
explicit equation of order three and none of lower order
(Theorems~\ref{hol:td:main:boundary}, \ref{hol:td:thm:ode}
and~\ref{hol:td:thm:field}), so the answer to the first sentence is yes. Its
equation has $\chi_P=0$, and $\mathcal T_p$ satisfies no equation with
$\chi_P\ne0$ (Proposition~\ref{hol:td:prop:corner},
Remark~\ref{hol:td:rem:corner}), so the answer to the second is also yes.
$\mathcal T_p$ is a differentially algebraic element of $\operatorname{Frac}\E$
that is not rational, so the rationality question has a negative answer. The
absence of an order unit is necessary: linear rigidity
(Theorem~\ref{hol:main:ode}) cannot be strengthened to differential algebraic
rigidity for any $\Gamma$ with an order unit (Corollary~\ref{hol:td:cor:sharp}).
Source~13 leaves open whether order exactly two occurs and asks for a
classification (its Questions~11.1 and~11.2); the question was re-scoped to
these, in the wording recorded in the next paragraph. Neither $\Theta_p$ nor
$G$ is decided: whether either is differentially algebraic stays open.

As re-scoped after source~13, the question read: \emph{Suppose $\Gamma$ has an
order unit. Can a nonpolynomial $f\in\E$ satisfy a nonzero algebraic
differential equation $P(z,f,D_zf,\ldots,D_z^sf)=0$ over $\K(z)$ of order
$s=2$, that is, are $f$, $D_zf$ and $D_z^2f$ algebraically dependent over
$\K(z)$ for some nonpolynomial $f\in\E$? By Theorem~\ref{hol:nl:main:first}
no equation of order at most one is possible, and by
Theorem~\ref{hol:td:main:boundary} order three occurs, so the least order over
all examples is two or three. More broadly, to what extent do finite algebraic
descents of theta or abelian functions, combined with elementary operations
preserving strong entireness, account for the nonpolynomial strongly entire
series that are differentially algebraic over $\K(z)$? The same questions are
open over $\KK=k((t^\Gamma))$ for every field $k$ of characteristic zero.}
Source~14 (the order-three-threshold part) answers the first sentence
negatively, for every $\Gamma$, with or without an order unit, and for every
characteristic-zero coefficient field, including $k=\C$: a strongly entire
solution of a nonzero equation of order at most two is a polynomial
(Theorem~\ref{hol:ot:main:second}\,(a)), even when its residue corner
vanishes. So the least order is exactly three
(Theorem~\ref{hol:ot:main:second}\,(b)). Source~14 was written against the
repository at \texttt{b895e86}, before source~13 was merged, and answers
source~13's question \texttt{q:order2}, which is the same question
(Section~\ref{hol:ot:sub:conventions}). The classification sentence stays
open; source~14 poses it again as its Research question~1, which is merged
here, and the question is re-scoped to it, as printed above. For $\Theta_p$
and $G$ the answer gives three independent jets
(Corollary~\ref{hol:ot:cor:examples}); whether either satisfies an equation of
order at least three stays open. The eight further questions of source~14 are
Questions~\ref{hol:ot:q:lift}--\ref{hol:ot:q:omnific} below.

\emph{Status (batch~31).} Source~19 re-derives the order-two answer
independently (Section~\ref{hol:nr:sub:conventions}) and narrows the
classification from below: a series of minimal order three satisfies no
equation of order three and total jet degree at most three
(Theorem~\ref{hol:nr:main:third}), and the exterior form of every equation of
order three that it satisfies has vanishing first polar
(Corollary~\ref{hol:nr:cor:theta}); for $\mathcal T_p$ the least such jet
degree is four, five or six. Source~19's own Question~12.1, the order-unit
case of this question in all orders, is answered by source~13, and its
abstract's ``unresolved'' is stale against the current text. The
classification sentence stays open, and the question keeps its wording.

\emph{Status (batch~31, sources~16 and~17).} Source~16 shows, for $\Gamma=\R$,
that minimal orders are unbounded: polynomially weighted sums of theta descents
at reciprocal-independent scales have orders in $[N,3N]$ for every $N$, and the
differentially algebraic strongly entire series generate a field of
transcendence degree $2^{\aleph_0}$ and differential transcendence degree zero
(Theorem~\ref{hol:th:main:hierarchy}). These are finite algebraic combinations
of theta descents under operations preserving strong entireness, as the last
sentence of the question contemplates; so the minimal orders form an unbounded
set containing three (Remark~\ref{hol:th:rem:ordertwo}), and which orders above
three occur is Question~\ref{hol:th:q:spectrum}. Source~17 shows that the same
order-unit fields carry explicit nonpolynomial strongly entire series with
coefficients in $\Q(t^\mu)$ satisfying no algebraic differential equation
(Theorem~\ref{hol:fh:main:single}); its Question~12.1, this question in the
form posed before source~13, is answered by Theorem~\ref{hol:td:main:boundary}.
Neither source decides $\Theta_p$ or $G$ (Question~\ref{hol:fh:q:heights}). The
question keeps its wording.

The escape proof does not answer this, in any order. A nonlinear differential
relation produces convolution sums in the coefficients rather than a
bounded-step linear recurrence, and treating it as though it had the form
\eqref{hol:eq:generalrecurrence} would be an invalid extension. The
first-order answer uses the different route of
Section~\ref{hol:nl:sec:setting}, the no-order-unit answer the third route
of Section~\ref{hol:cf:sec:coefficients}, and the order-two answer a fourth,
initial polynomials at infinitely many scales
(Section~\ref{hol:ot:sub:endpoint}).

\begin{question}[Beyond a positive Euler operator]\label{hol:nl:q:vectorfields}
Which polynomial vector fields $V$ have the property that every strongly
entire $f$ satisfying a nonzero equation $P(x,f,Vf)=0$ is a polynomial?
\end{question}

Positive integer Euler weights give one family, and zero and mixed weights
give explicit counterexamples (Section~\ref{hol:nl:sub:euler}). A useful
extension would need both a finite-dimensional filtration and a sufficiently
nondegenerate action on its highest layers; diagonalization alone is not
enough if infinitely many monomials share a bounded eigenvalue. (Source~10.)

\begin{question}[Effective representations and minimal certificates]
\label{hol:nl:q:effective}
For which exact Hahn coefficient representations can the candidate degree set
and an optimal excluded valuation cut be computed uniformly?
\end{question}

The bound \eqref{hol:nl:eq:deltazero} is always a valid choice from exact
finite valuation data but can be far from minimal, and
Theorem~\ref{hol:nl:thm:nouniform} shows that an optimal cut must in general
depend on solution data. An algorithmic refinement must state how it receives
those data and how it decides coefficient cancellation. (Source~10; the
linear analogue is Question~\ref{hol:q:optimal}.)

\begin{question}[Mixed operators at unit valuation, with an order unit]
\label{hol:q:mixedunit}
Suppose $\Gamma$ has an order unit. For a nontorsion $q$ with $v(q)=0$, must
every strongly entire solution of a nonzero mixed equation
\eqref{hol:eq:mixed} be a polynomial?
\end{question}

As posed by sources~08 and~09 the question concerned every $\Gamma$. The
unresolved point in the escape method is uniform control of $v(P(n,q^n))$ for
a general bivariate $P$ in the unit case. The pure $q$-difference proof
controls $P(q^n)$ and the differential proof controls $P(n)$; neither
statement by itself handles their combination. This is a question suggested
by the argument, not asserted to be a previously catalogued open problem.

\paragraph{Status.}
When $\Gamma$ has no order unit, source~12 answers it positively and more
strongly: by Theorem~\ref{hol:cf:main:entire}\,(b) with
$\Lambda=\{q^i:i\in\Z\}$, a nonpolynomial $f\in\E$ satisfies no nonzero
algebraic relation, linear or not, among the $(D_z^jf)(q^iz)$, and a nonzero
equation \eqref{hol:eq:mixed} is such a relation
(Section~\ref{hol:cf:sub:questions}). The method bounds the coefficient field
instead of the valuations $v(P(n,q^n))$. The question is re-scoped to groups
with an order unit, where it stays open: the witness of
Theorem~\ref{hol:cf:main:sharp} uses a parameter of nonzero valuation, and
source~13's series $\mathcal T_p$, which gives nonlinear mixed relations for
every nontorsion $q$ (Corollary~\ref{hol:td:cor:sharp}), is not D-finite, and
source~13 says nothing about linear mixed equations.

\emph{Status (batch~31).} Source~18 answers the question positively, for every
$\Gamma$ with or without an order unit and over every characteristic-zero
coefficient field: for nontorsion $q$ with $v(q)=0$ every strongly entire
solution of a nonzero equation \eqref{hol:eq:mixed}, also with polynomial or
rational forcing, is a polynomial, with a uniform degree bound and exterior
obstruction for each equation (Corollary~\ref{hol:ud:cor:unitrigidity},
Theorem~\ref{hol:ud:main:mixed}). The obstruction named above, uniform control
of $v(P(n,q^n))$, is supplied by Corollary~\ref{hol:ud:cor:bivariate}: the
valuation is eventually periodic, with period dividing the order of
$\res(q)$ when that is a root of unity. The question is answered in the form
posed; the nonlinear analogue is settled negatively by
Corollary~\ref{hol:td:cor:sharp}, as recorded before.

\emph{Status (batch~33).} Source~20 re-derives the answer independently, for
multipliers in a torsion-free group, through the same valuation theorem for
Hahn exponential polynomials (Section~\ref{hol:rc:sub:mixed}); it was written
against the text before source~18 was merged, when the question was still
open there.

\begin{question}[Optimal exterior regions and exact solution domains]
\label{hol:q:optimal}
For a fixed differential or rigid dilation operator, can one compute the
largest common strong evaluation domain of its formal solution space from
finite valued operator data, rather than only a sufficient exclusion bound?
Can the coarse exclusion bound of Theorem~\ref{hol:thm:coarse} be supplemented
by support hypotheses that give the exact domain?
\end{question}

The thresholds supplied here are uniform and explicit once the valuation
certificates are available, but need not be optimal, and distinct solutions may
have cancellations or terminating tails that change their individual domains.
Proposition~\ref{hol:prop:expdomain} shows that the uncoarsened bound can be
exactly attained, and Theorem~\ref{hol:thm:thetadomain} exhibits one
sufficient pattern for an exact answer: a common positive support monoid
together with a controlled sequence of leading exponents. A full answer should
retain arbitrary Hahn tails rather than use only a numerical radius.

\begin{question}[Several independent dilations]\label{hol:q:severaldilations}
When $\Gamma$ has an order unit, what is the existence criterion for a
nonpolynomial strongly entire solution constrained by equations using several
multiplicatively independent dilations? Is an appropriate cofinality
condition on their joint valuation data sufficient, or do compatibility
relations impose additional obstructions? In the no-order-unit workspace
\eqref{hol:eq:gammainfty}, can one classify the finite linear relations among
$F(q_1z),\ldots,F(q_sz)$ when some quotient $q_i/q_j$ is a root of unity?
\end{question}

\paragraph{Status.}
As posed by sources~08 and~09 the first two sentences concerned every
$\Gamma$, and the last asked about parameters that are not powers of a common
element. In a group without an order unit the merge observes
(Corollary~\ref{hol:cf:cor:several}, from source~12's
Theorem~\ref{hol:cf:main:entire}\,(b); source~12 does not claim it) that no
nonpolynomial strongly entire series satisfies any nonzero equation, linear or
not, in finitely many derivatives at dilations from the group generated by
multiplicatively independent parameters, and that $F(q_1z),\ldots,F(q_sz)$
satisfy only the trivial linear relation whenever no quotient $q_i/q_j$ is a
root of unity, whether or not the $q_i$ are powers of a common element. The
question is re-scoped to what this leaves: groups with an order unit, and
parameters with a root-of-unity quotient.

\emph{Status (batch~31).} For linear equations source~18 answers the
existence half of the first sentence, for every $\Gamma$, in the following
form. For a finite set of multipliers with no root-of-unity quotient, some
nonzero mixed linear equation supported in the set has a nonpolynomial
strongly entire solution if and only if $\abs{v(\lambda_i/\lambda_j)}$ is an
order unit for some pair; two multipliers and first derivatives then suffice,
and otherwise every such equation, with polynomial forcing, has only
polynomial entire solutions (Theorems~\ref{hol:ud:thm:rigidity}
and~\ref{hol:ud:thm:dichotomy}). So the appropriate cofinality condition is
the relative one, and a common absolute scale is irrelevant. This is a
partial answer (Remark~\ref{hol:ud:rem:quantifiers}): it concerns the
existence of \emph{some} equation, not the solvability of a prescribed
operator (Question~\ref{hol:ud:q:prescribed}) or of simultaneous equations
(Question~\ref{hol:ud:q:simultaneous}), and it keeps the exclusion of
root-of-unity quotients (Question~\ref{hol:ud:q:torsionblocks}). The question
stays open in those respects, and its last sentence, on root-of-unity
quotients in $\Gamma_\infty$, is untouched.

\emph{Status (batch~31, source~15).} If an equation also contains a nonzero
operator applied to $f(P(z))$ with $\deg P\ge2$ strictly above the degrees of
all other arguments, every strongly entire solution is a polynomial, for
every $\Gamma$ and whatever the dilations
(Corollary~\ref{hol:pc:cor:linear}, Remark~\ref{hol:pc:rem:dichotomy}). The
question concerns arguments of degree one and is otherwise unaffected.

The single-dilation theorem identifies the exact obstruction when only one
scale is repeatedly added. With several dilation directions, distinct
monomials can have equal valuation slopes and their leading coefficients can
produce nontrivial recurrence phenomena of their own. A meaningful formulation
must also quotient out identities caused by repeated parameters and any
genuine algebraic relations among them. No multidilation classification is
implied by Theorem~\ref{hol:main:q}.

The next three questions come from sources~11 and~12.

\begin{question}[An intrinsic descent criterion; source~11]
\label{hol:cf:q:descent}
For a fixed Hahn field $\KK$ and a subfield $L\subseteq\KK$, can the equality
$\FE\cap L((z))=L(z)$ be characterized solely in terms of the valued
extension $\KK/L$?
\end{question}

Containment of $v(L^\times)$ in a proper convex subgroup is sufficient
(Theorem~\ref{hol:cf:thm:descent}) and applies to many infinitely generated
fields; it is not claimed necessary. A criterion using rational rank alone
will not capture all cases: what matters to the reduction is whether the
coefficient values reach every Archimedean scale of $\Gamma$.

\begin{question}[A finite-rank replacement; source~12]
\label{hol:cf:q:finiterank}
Which additional properties of a finitely generated coefficient field, of its
valuation, or of the recurrence \eqref{hol:cf:eq:recurrence} force a linear
upper bound on the relevant nonzero coefficient values, and so exclude strong
entireness even when those values are cofinal?
\end{question}

Bounding the multiplier $\Xi(n)$ is not enough: a nonlinear recurrence can
involve cancellations among many products of earlier coefficients, and an
extension must control them or replace them by another invariant. An answer
would bear on Question~\ref{hol:q:nonlinear}. The series $\mathcal T_p$ of
Theorem~\ref{hol:td:main:boundary} is differentially algebraic and strongly
entire, with a finitely generated coefficient field and coefficient values
$m^2\mu$ that are cofinal and not linearly bounded; with an order unit, any
such property must therefore fail for it (Section~\ref{hol:td:sub:consistency}).

\begin{question}[Torsion blocks; source~12]\label{hol:cf:q:torsion}
When some dilations have root-of-unity quotients, can the coefficient-field
conclusion of Theorem~\ref{hol:cf:thm:coefficients}\,(b) be recovered
separately on the finitely many residue classes of Taylor indices that
survive the resonances?
\end{question}

The even-series example of Section~\ref{hol:cf:sub:boundaries} shows that no
unconditional version for all coefficients exists; a blockwise theorem would
need explicit conditions on which residue classes are constrained, not only
on the nonzero part of the multiplier.

The next question comes from source~13 (its Question~11.3).

\begin{question}[Effective strong domains for differential solutions]
\label{hol:td:q:domains}
For an explicitly represented formal solution of an algebraic differential
equation, can its maximal strong domain be recovered from a finite
differential system and coefficient data?
\end{question}

For $\mathcal T_p$ the answer is the coarsened valuation ring
$\mathcal O_{\Delta_\mu}$ (Theorem~\ref{hol:td:thm:domain}), whose boundary
depends on the convex subgroup generated by one scale. General solutions can
have much more complicated supports. An algorithm needs a representation of
those supports and a decision procedure for the relevant coefficient
cancellations; inspecting the leading differential monomials cannot suffice
in the vanishing-corner case. Compare Questions~\ref{hol:nl:q:effective}
and~\ref{hol:q:optimal}.

The next eight questions come from source~14 (its Research questions~2--9;
its Research question~1 is merged into Question~\ref{hol:q:nonlinear}).

\begin{question}[Lifting admissible initial polynomials; source~14]
\label{hol:ot:q:lift}
Given an equation of order three and an unbounded sequence of initial
polynomials satisfying its exterior form, what compatibility conditions
guarantee that they are the initial polynomials of one strongly entire
solution?
\end{question}

A criterion would have to control both coefficient lifting and global
support: local formal solutions at separate scales need not assemble into one
series. $\mathcal T_p$ is a fully worked positive case, with the initial
polynomials $X^n(1+X)$ at its breaks.

\emph{Status (batch~31).} Source~19 gives a concrete negative instance
(Example~\ref{hol:nr:ex:sparse}): the initial polynomials of the strongly
entire series $F_{\mathrm{sp}}$ satisfy the exterior forms of $\mathsf H_3$ and
of $\mathsf Q_4$ at every scale, yet $F_{\mathrm{sp}}$ solves neither
equation, the first failure being the coefficient $-126\,t^{17}$ of $z^5$ in
$\mathsf Q_4(F_{\mathrm{sp}})$. So agreement of initial polynomials at every
break is not a sufficient compatibility condition, even when those
polynomials come from one entire series. Question~\ref{hol:nr:q:layers} asks
for the higher-layer conditions; the question stays open.

\begin{question}[The entire solutions of the theta equation; source~14]
\label{hol:ot:q:thetasolutions}
For the constants \eqref{hol:td:eq:constants}, determine all strongly entire
solutions of \eqref{hol:td:eq:cleared}. Which are scalar multiples,
coordinate transforms or other members of a theta family, and which come from singular
branches of the cleared equation?
\end{question}

Theorem~\ref{hol:ot:thm:convexity} forces eventually full coefficient support
and strictly increasing consecutive coefficient differences. The next step
would use the lower exterior terms, discarded by the exterior form, to
constrain those differences and the leading coefficients; the cleared form may
have exceptional solutions not represented in logarithmic coordinates, which
need separate treatment. Since \eqref{hol:td:eq:jetpoly} is homogeneous in
the jets, every constant multiple of $\mathcal T_p$ is a solution.

\begin{question}[Effective third-order exclusion; source~14]
\label{hol:ot:q:effective}
For which coefficient fields and equations can boundedness of the
nonnegative-integral endpoint projection of $\Upsilon_{\mathfrak h_P}(T,G)=0$
be decided? Can a finite library of classifications of residue polynomials
extend the certificate of Corollary~\ref{hol:ot:cor:gapcriterion}?
\end{question}

At order two the construction is finite algebra once the leading coefficient
values are accessible: dehomogenize, substitute, remove a power of $V$ and
inspect the remainder on $V=0$. Order three introduces a genuine two-integer
constraint. Any algorithmic claim must include computability of the Hahn
coefficients of the input and of their exact leading terms; compare
Question~\ref{hol:nl:q:effective}.

\emph{Status (batch~31).} Source~19 extends the certificate in one direction:
when the exterior form has a nonzero first polar, that is $i_0=1$, the two
endpoint conditions cannot both hold far out, whatever the upper-endpoint set
of Corollary~\ref{hol:ot:cor:gapcriterion} is
(Proposition~\ref{hol:nr:prop:firstpolar}, Remark~\ref{hol:nr:rem:gaps}); the
test applies factor by factor (Corollary~\ref{hol:nr:cor:factor}), and it
settles every equation of order three and total jet degree at most three
(Theorem~\ref{hol:nr:main:third}). The quartic $\mathsf H_3$, with $i_0=2$ and
$\Upsilon=(G+T)(2G-T)$, shows that the test fails in degree four; its formal
solutions are classified exactly (Theorem~\ref{hol:nr:thm:H3}). The question
stays open for $i_0\ge2$.

\begin{question}[Higher-order break invariants; source~14]
\label{hol:ot:q:higher}
Can the expansion transverse to constant logarithmic derivatives be organized
into a hierarchy of gap and multiplicity spectra for orders four and above,
with rigidity theorems at each level?
\end{question}

At higher orders several logarithmic derivatives vanish to the same first
local order, with leading ratios determined by powers of an endpoint gap;
secondary terms could record further support spacings. When the resulting
constraints are complete is open.

\emph{Status (batch~31).} Source~19's quadratic $\mathsf Q_4$ of order four is
a first worked case: its cross coefficient is the pair kernel
$(m-n)^2(m-2n)(n-2m)$, and its formal solutions are exactly the monomials and
the doubling binomials (Theorem~\ref{hol:nr:thm:Q4}); Question~\ref{hol:nr:q:jetcurve}
asks for the forms with such sparse solutions in general.

\begin{question}[Meromorphic second-order solutions; source~14]
\label{hol:ot:q:meromorphic}
Beyond polynomial denominators, which algebraic differential equations of
order two have nonrational quotients of strongly entire series as solutions?
Can the break method bound the density or the valuation distribution of
their poles?
\end{question}

Elliptic meromorphic functions suggest that polynomiality cannot simply be
extended to every quotient of entire series when an order unit exists; the
comparison involves separate Newton data for numerator and denominator,
including cancellation and common zeros. Without an order unit,
Theorem~\ref{hol:cf:main:entire}\,(c) already makes every differentially
algebraic element of $\FE$ rational.

Source~19 asks the same question (its Question~12.4), for equations of order
at most two and elements of $\FE$ in the order-unit case, and notes that the
initial object of a quotient is rational rather than polynomial, so the
first- and last-support arguments need replacement.

\begin{question}[Partial differential systems; source~14]
\label{hol:ot:q:pde}
Is there a useful several-variable analogue of three-jet independence for
series strongly summable at every point of $\KK^d$? Which finite-dimensional
rational differential systems force polynomiality of all strongly entire
observables?
\end{question}

Restriction to lines may settle special cases, as for
Theorem~\ref{hol:main:multi}, but algebraic relations among several partial
derivatives do not in general reduce to a relation among the first two
derivatives on every line; a several-variable analysis of initial forms is
needed.

Source~19 asks the same question (its Question~12.6) for a positive-weight
Euler operator: a leading weighted component is then a homogeneous polynomial
rather than a multiple of one monomial, the conic factorization suggests a
route, and diagonal specialization can hide cancellation, so a multivariable
break statement should come first.

\begin{question}[Positive residue characteristic; source~14]
\label{hol:ot:q:charp}
Can a replacement for second-order rigidity be formulated with Hasse or
divided derivatives, after excluding series that are constant for the
Frobenius and specifying the behaviour of integer valuations?
\end{question}

The example of Section~\ref{hol:ot:sub:surreal} and
Example~\ref{hol:nl:ex:charp} show that ordinary derivatives cannot suffice.
A replacement must distinguish the characteristic of the field, that of the
residue field and the notion of summation.

Source~19 asks the same question (its Question~12.8), adding that neither the
conic cross coefficient nor the double factor of the first polar transfers
verbatim, and that equal characteristic zero, mixed characteristic and
positive characteristic must be distinguished; source~18's version for the
unit theorem is Question~\ref{hol:ud:q:charp}.

\begin{question}[Omnific approximation of differential zero sets; source~14]
\label{hol:ot:q:omnific}
For explicitly represented entire solutions of minimal order three, which
zero sets admit uniform approximation by omnific integers, and how do endpoint
gaps constrain the valuations of the rounding errors?
\end{question}

The zeros of $\mathcal T_{\omega^{-1}}$ have exact rounding errors
$\omega^{-(2m-1)}$ (Corollary~\ref{hol:td:cor:omnific}). For other entire
solutions the coefficient valuations are not enough: residue directions and
higher terms can move individual zeros.

The next eleven questions come from source~18 (its Research
questions~11.2--11.12; its Research question~11.1 is answered in the current
text, Remark~\ref{hol:ud:rem:nonlinearq}).

\begin{question}[Prescribed operators at a cofinal quotient; source~18]
\label{hol:ud:q:prescribed}
For a specified mixed operator $L$ whose multipliers cross the top valuation
scale, which conditions characterize the existence of a nonpolynomial
strongly entire solution of $Lf=0$?
\end{question}

Theorem~\ref{hol:ud:thm:dichotomy} lets one choose $L$ and does not answer
this. A first target is a two-multiplier operator of differential order one:
its recurrence has explicit valuation slopes, and one must decide whether
compatible solutions have coefficient valuations escaping fast enough, not
merely whether a formal solution exists.

\begin{question}[Simultaneous equations; source~18]\label{hol:ud:q:simultaneous}
Can a nonpolynomial strongly entire series satisfy two independently
prescribed mixed equations whose cofinal dilation parameters are
multiplicatively independent? Which compatibility conditions on their
recurrences are decisive?
\end{question}

The witness for one pair of multipliers need not satisfy an equation for
another pair. This sharpens Question~\ref{hol:q:severaldilations}; a starting
point is to compare the two induced recurrences before seeking a general
difference--differential Galois theory.

\begin{question}[A torsion-block criterion for linear mixed rigidity;
source~18]\label{hol:ud:q:torsionblocks}
When some multiplier quotients are roots of unity, can
Theorem~\ref{hol:ud:thm:rigidity} be replaced by a finite block criterion
deciding which residue classes of Taylor indices are forced to terminate?
\end{question}

Theorem~\ref{hol:ud:thm:affine} already allows identically zero progressions;
the difficulty is that recurrence jumps can move between classes on which
the leftmost coefficient vanishes. A finite directed graph of residue classes
with nonvanishing shift data on its edges is a plausible tool, but every
surviving infinite path needs an escape inequality; checking classes in
isolation is not enough. Question~\ref{hol:cf:q:torsion} asks the analogous
question for the coefficient-field conclusion.

\emph{Status (batch~33).} Source~20 poses the same question (its
Question~12.2) for a finitely generated dilation group with torsion: decompose
formal series by degree modulo the torsion order, account for the sector
shifts caused by derivatives and polynomial coefficients, and distinguish a
genuinely constrained sector from the free even series of
Section~\ref{hol:rc:sub:boundaries}. It stays open.

\begin{question}[Several variables without an ellipticity shortcut;
source~18]\label{hol:ud:q:several}
For jointly strongly entire series in several variables, which finite systems
of partial differential--dilation equations admit an analogous exact
relative-scale obstruction?
\end{question}

A single partial derivative may annihilate an arbitrary function of the
other variables, so the one-variable statement does not extend verbatim. The
hypothesis should constrain all escaping multi-indices, perhaps through
finitely many positive-weight directions or a rank condition; restriction to
a generic line, as for Theorem~\ref{hol:main:multi}, would need the operator
system and joint strong summability to be preserved. Compare
Question~\ref{hol:ot:q:pde}.

\emph{Status (batch~33).} Source~20's Question~12.12, on finite linear
substitutions in several variables and matrix dilations, is the neighbouring
Question~\ref{hol:rc:q:substitutions}.

\begin{question}[Optimal exterior cuts for mixed operators; source~18]
\label{hol:ud:q:optimal}
Given the exact periodic-affine valuation data of a mixed operator, can one
compute the largest common strong evaluation domain of its formal solutions?
\end{question}

The jump-sensitive inequalities of Theorem~\ref{hol:thm:escapeexclusion} are
sufficient; optimality would require knowing which recurrence edges actual
solutions realize and how their nonleading supports behave, and valuation data
alone may not decide it. A first problem is an operator with monomial
coefficients and a fixed finite jump graph. This is the mixed form of
Question~\ref{hol:q:optimal}.

\begin{question}[Effective coefficient representations; source~18]
\label{hol:ud:q:effective}
For which finitely specified Hahn subfields can the support envelope, the
first coefficient sequence that is not an identity, and all exceptional
indices in Theorem~\ref{hol:ud:thm:affine} be computed?
\end{question}

A first target is Laurent or Puiseux polynomial data with finite-order residue
quotients, where the zero tests reduce to polynomial calculations once a
support bound is proved; arbitrary residue bases need a separate exact
zero-set analysis. No universal algorithm follows from the qualitative
theorem.

\emph{Status (batch~33).} Source~20 poses the same question (its
Question~12.1) for Hahn inputs with algebraic coefficients and finitely
generated, explicitly represented supports, and asks to separate the search
for the first coefficient sequence that is not an identity from the location
of the zeros of that exponential polynomial, which the existence proof does not
eliminate. It adds a target of the same kind:
Theorem~\ref{hol:rc:thm:compression} gives a finite equivalent subfamily of a
set-sized family of coefficient conditions only existentially
(Section~\ref{hol:rc:sub:boundaries}). Both stay open.

\begin{question}[Valuation profiles beyond Hahn fields; source~18]
\label{hol:ud:q:beyond}
Which equicharacteristic-zero valued fields have the periodic-affine property
of Theorem~\ref{hol:ud:thm:affine} for exponential polynomials, without a
specified Hahn expansion?
\end{question}

A support-preserving embedding into a Hahn field would transfer the result
when available, but finding sufficient hypotheses for it is part of the
problem; an intrinsic proof needs substitutes for the well-ordered set of
candidate coefficients and for the least coefficient sequence that is not an
identity. No embedding theorem beyond the explicit surreal workspaces is
asserted.

\emph{Status (batch~33).} Source~20's Question~12.5 is the same question: which
embedding or completion hypotheses preserve the conclusion without unjustified
summability assumptions, so that the genuinely Hahn-theoretic part of the
theorem is isolated. It stays open.

\begin{question}[A positive-characteristic replacement for the unit theorem;
source~18]\label{hol:ud:q:charp}
Can Frobenius-sensitive index sets and a modified notion of differential
operator give a useful replacement for Theorem~\ref{hol:ud:thm:unitvalues} and
Corollary~\ref{hol:ud:cor:unitrigidity} in characteristic $\mathfrak p$?
\end{question}

Ordinary derivatives cannot suffice (Section~\ref{hol:ud:sub:examples}).
Hasse derivatives and automatic zero-set patterns~\cite{derksen} are natural
ingredients, but unbounded valuations along $\mathfrak p$-power indices must be
treated explicitly. Compare Question~\ref{hol:ot:q:charp}.

\begin{question}[Scalar extension of the obstruction; source~18]
\label{hol:ud:q:extension}
For an ordered-group extension $\Gamma\subseteq\Gamma'$, which entire solution
spaces survive the passage from $k((t^\Gamma))$ to $k((t^{\Gamma'}))$?
\end{question}

Rigidity survives whenever the relative scales stay noncofinal; a cofinal
quotient in $\Gamma$ may stop being cofinal in $\Gamma'$, and its theta witness
then loses entireness. An exact domain theorem should track full support
certificates, not only the extension of the coefficient field; the
enlargement of $G_\mu$ in Section~\ref{hol:ud:sub:surreal} is a first example.

\begin{question}[Arithmetic restrictions on omnific operators; source~18]
\label{hol:ud:q:omnific}
Do operators with coefficients in an omnific subring admit stronger bounds on
polynomial solutions, or more effective valuation certificates, than general
Hahn-coefficient operators?
\end{question}

The theorems allow such coefficients but use neither their one-sided supports
nor their integer constant terms. A first task is to study finite-support
omnific coefficients and determine which cancellation profiles in
Theorem~\ref{hol:ud:thm:affine} occur; any improvement must account for
$\sigma_2-2^a$ (Example~\ref{hol:ud:ex:degree}), which already has ordinary
integer coefficients.

\emph{Status (batch~33).} Source~20 continues the question on the side of
sequences rather than operators. The leading exponent of a recurrence over
$\Oz$ or $\Oz[i]$ is eventually affine on finitely many progressions, two
sufficing when the characteristic roots are real
(Theorem~\ref{hol:rc:thm:surreal}); for a real-surreal unit-root recurrence,
support conditions are decided by finitely many coefficients along the orbit,
with hitting sets periodic modulo $\mathsf M_{\mathrm{tor}}$ up to a finite set
(Corollary~\ref{hol:rc:cor:tails}); but one fixed Conway coefficient of the
purely infinite orbit $y_{n+1}=\omega y_n$ can be any real sequence
(Theorem~\ref{hol:rc:thm:coding}). Full omnific membership is
Question~\ref{hol:rc:q:membership}. The question about operators stays open.

\begin{question}[A proof-assistant implementation with an explicit trust
boundary; source~18]\label{hol:ud:q:formal}
Can the finite-order-residue case and the two-multiplier witness be
formalized first, and the general theorem once the Skolem--Mahler--Lech input
is available in verified form or as an explicit hypothesis?
\end{question}

This separates a support and order component, a finite algebraic component
and a number-theoretic zero-set component, and keeps finite testing, imported
theorems and kernel-checked proof apart.

\emph{Status (batch~33).} Source~20's Question~12.11 proposes the same
separation for its chain: formalize Theorem~\ref{hol:rc:thm:algebra} with
Lemma~\ref{hol:ud:lem:envelope}, the argument by the least coefficient sequence
that is not an identity, and the escape lemma, keeping the
Skolem--Mahler--Lech theorem as a declared imported theorem until it is
available; the repository's support and recurrence-escape infrastructure would
need a theorem-label audit before being reused as checked coverage. No Lean
work is added.

The next six questions come from source~19 (its Questions~12.2, 12.3, 12.5,
12.7, 12.9 and~12.10; its Question~12.1 is answered by
Theorem~\ref{hol:td:main:boundary}, and its Questions~12.4, 12.6 and~12.8 are
Questions~\ref{hol:ot:q:meromorphic}, \ref{hol:ot:q:pde}
and~\ref{hol:ot:q:charp}).

\begin{question}[Higher residue layers and differential prolongation;
source~19]\label{hol:nr:q:layers}
Is there a finite collection of higher-layer compatibility conditions that
excludes every nonpolynomial entire lift of an exterior form whose initial
polynomial solutions are not eventually monomials?
\end{question}

Example~\ref{hol:nr:ex:sparse} is a concrete test: all initial polynomials of
$F_{\mathrm{sp}}$ satisfy both boundary equations, while the coefficient
\eqref{hol:nr:eq:residual} detects a later obstruction. In the quartic case,
differentiating a common-root parameter produced the decisive identity
\eqref{hol:nr:eq:parameter}; a general version might combine breaks with
differential elimination and separants. This refines
Question~\ref{hol:ot:q:lift}.

\begin{question}[Jet-curve ideals and sparse polynomial solutions; source~19]
\label{hol:nr:q:jetcurve}
For the rational normal curve $\mathbf m_r(T)=(1,T,\ldots,T^r)$, classify the
homogeneous forms $\mathfrak h$ for which
$\mathfrak h(\phi,\vartheta_X\phi,\ldots,\vartheta_X^r\phi)=0$ has nonmonomial
polynomial solutions of arbitrarily large lowest degree.
\end{question}

Relevant invariants include the multiplicity of vanishing along the curve,
first and higher polars and secant varieties. The conic's principal ideal is
unusually rigid, and $\mathsf H_3$ shows that singularity along the twisted
cubic allows unbounded doubling supports. An effective criterion on
$\mathfrak h$ is the useful outcome, not merely a classification of its zero
set. Compare Questions~\ref{hol:ot:q:effective} and~\ref{hol:ot:q:higher}.

\begin{question}[Systems in three or more variables; source~19]
\label{hol:nr:q:systems}
Which strongly entire solutions of finite-dimensional rational or algebraic
first-order systems must be polynomial? In particular, with an order unit, do
three states suffice to produce a nonpolynomial strongly entire observable?
\end{question}

The planar result follows from three-jet independence. Without an order unit
every number of states is covered (Corollary~\ref{hol:cf:cor:systems}); with
one, four states suffice to produce $\mathcal T_p$
(Section~\ref{hol:nr:sub:constants}), and three are open. Dimension three
would need higher-order independence or a bound on the degree of an
eliminated relation, since Theorem~\ref{hol:nr:main:third} applies only to
relations of total degree at most three and such a bound cannot be assumed
from the dimension. (The second sentence is the merge's re-scoping.)

\begin{question}[Effective certificates and bounds; source~19]
\label{hol:nr:q:certificates}
Given a finitely described differential equation and coefficient profile, can
one compute an exterior non-entireness certificate for every nonpolynomial
formal solution?
\end{question}

The corner inequalities are finite and explicit but can depend on a high
nonzero coefficient of the particular solution, and uniform degree bounds are
impossible in general by \eqref{hol:nl:eq:degenerate}. A realistic target is a
solution-dependent certificate with a finite exceptional-degree analysis
where the monomial test does not vanish identically. Compare
Questions~\ref{hol:nl:q:effective} and~\ref{hol:ot:q:effective}.

\begin{question}[Omnific normalization with presentations; source~19]
\label{hol:nr:q:omnific}
For finitely presented surcomplex constants, when can the multiplier
$\omega^\rho$ of Lemma~\ref{hol:nr:lem:omnificclear} be computed with a useful
bound on the complexity of $\rho$?
\end{question}

Existence uses only a set-sized bound on supports; effective normalization
needs representations of those supports and an algorithm bounding them,
without confusing coefficient clearance with effective global surreal
computation.

\begin{question}[Formal verification and independent review; source~19]
\label{hol:nr:q:formal}
Can the conic and first-polar arguments be checked in Lean using only the
repository's stated Hahn interfaces, with every constant extension and
universe change explicit?
\end{question}

The finite algebra (conic kernel, bottom and top first-polar coefficients,
double factor, eventual side, secant Jacobian) is small enough to isolate
first; the break construction, the sign-sensitive inequality of the exterior
reduction and the assembly would follow, and finite descent should precede
the transfer to $\No[i]$. Independent review of the break construction and the
eventual-side lemma should precede broader claims.

The next nine questions come from source~16 (its Questions~11.1 and
11.3--11.10; its Question~11.2 is answered by
Theorem~\ref{hol:ot:main:second}, Remark~\ref{hol:th:rem:ordertwo}).

\begin{question}[The exact order spectrum; source~16]\label{hol:th:q:spectrum}
Which integers occur as $\operatorname{dord}(f)$ for nonpolynomial $f\in\EK$
differentially algebraic over $\KK(z)$? Does every integer $r\ge3$ occur?
\end{question}

Theorem~\ref{hol:th:main:hierarchy} gives unboundedness and realizes the
dimensions $d_N$; neither identifies the spectrum, which would need control of
derivative interactions beyond leading Gauss growth. By
Theorem~\ref{hol:ot:main:second} the spectrum lies in $\{3,4,\ldots\}$, and it
contains $3$ (Remark~\ref{hol:th:rem:ordertwo}, merge). Compare
Question~\ref{hol:q:nonlinear}.

\begin{question}[Joint jet independence; source~16]\label{hol:th:q:joint}
For scales with $\Q$-independent reciprocals, are the $3N$ series
$\mathcal T_{t^{\mu_i}},D_z\mathcal T_{t^{\mu_i}},D_z^2\mathcal T_{t^{\mu_i}}$,
$1\le i\le N$, algebraically independent over $\KK(z)$, that is, is $d_N=3N$?
\end{question}

A positive answer would give exact orders $3N$ through
Theorem~\ref{hol:th:thm:exactdimension}. The one-seed difference argument does
not apply scale by scale: dilating the shared variable changes all seeds at
once.

\begin{question}[Commensurable scales; source~16]\label{hol:th:q:commensurable}
Classify the algebraic relations among $\mathcal T_{t^{\mu_1}},\ldots,
\mathcal T_{t^{\mu_N}}$ when the reciprocals are $\Q$-dependent. Is there a
criterion by theta transformations or isogenies describing the exceptional
configurations exactly?
\end{question}

The present condition is sufficient, not necessary: its failure lets two
monomials of a relation share their quadratic growth, and produces no
relation. A classification must use the bounded and linear terms or further
elliptic structure.

\begin{question}[Efficient compression; source~16]\label{hol:th:q:compression}
Can the bound $(N+1)^{N^2}$ be reduced substantially, can simpler weights be
prescribed for the theta family, and can a good candidate be certified from a
finite computable presentation of the relevant K\"ahler differentials?
\end{question}

The universal argument ignores the structure of theta functions and is
probably wasteful for them; an effective version needs a certified
representation of jet relations, not a numerical search.

\begin{question}[Smaller exponent groups; source~16]\label{hol:th:q:smaller}
Over $\C((t^\Q))$, $\C((t^\Z))$ or another rank-one group, are the minimal
orders of differentially algebraic strongly entire series still unbounded?
Can large independent families be produced without $\Q$-independent real
growth constants?
\end{question}

The mechanism needs rationally independent reciprocal scales; embedding a
smaller field in $\KK$ does not answer the question, since the base field and
the coefficient field of the constructed series change.

\begin{question}[Beyond quadratic profiles; source~16]\label{hol:th:q:profiles}
Which differentially algebraic strongly entire series have profiles
$\gamma_\delta(f)=-c\delta^\varpi+o(\delta^\varpi)$ with $\varpi\ne2$, and can a
Legendre-transform analysis give further independent families and sharper
order bounds?
\end{question}

Theorem~\ref{hol:th:thm:growth} is ready for any $\varpi>1$; what is missing
are differentially algebraic seeds. Fast coefficient decay alone gives entire
series, not differential algebraicity (compare
Theorem~\ref{hol:fh:thm:entire}).

\begin{question}[Arbitrary value groups; source~16]\label{hol:th:q:groups}
Which version of the hierarchy survives in higher-rank Hahn fields with an
order unit, and what replaces the real-valued growth comparison? Can a
finite dominance certificate be stated in the ordered group without embedding
a subgroup in $\R$?
\end{question}

A generalization should isolate the dominating Archimedean quotient and check
that coefficient values in lower convex layers do not spoil the comparison.

\begin{question}[Intrinsic derivations and global domains; source~16]
\label{hol:th:q:intrinsic}
Is there an order hierarchy for an intrinsic surreal derivation that moves the
coefficients? Separately, which identities survive beyond the real-exponent
workspace without inadmissible strong sums?
\end{question}

These are distinct problems, one changing the derivation and the other the
domain; neither is solved by renaming $t$.

\begin{question}[Omnific sampling beyond the ordinary integers; source~16]
\label{hol:th:q:omnific}
For which explicit infinite subsets of $\Oz$ does a nonzero entire series of a
fixed Hahn workspace have only finitely many zeros? Can the residue argument of
Theorem~\ref{hol:th:thm:sampling} be replaced by finitely many residue tests
at varying leading scales?
\end{question}

Theorem~\ref{hol:th:thm:sampling} uses distinct ordinary residues; general
omnific inputs can share a residue or leave the workspace, and any extension
must fix both scales and domain.

The next ten questions come from source~17 (its Questions~12.2--12.11; its
Question~12.1 is answered in this report, Remark~\ref{hol:fh:rem:universal}).

\begin{question}[Pure differential freedom below rapid ratios; source~17]
\label{hol:fh:q:heights}
Which height sequences make every infinite subseries
$\sum_{n\in A}p^{\mathsf b_n}z^n$ differentially transcendental over $\KK(z)$
for $D_z$?
\end{question}

Quadratic heights are the instructive case: the series are entire, the mixed
theorem is obstructed by Proposition~\ref{hol:fh:prop:heat}, and for
$A=\N_{\ge2}$ the series is $G_\mu$ without its first two terms, a rescaled
partial theta series, whose differential
algebraicity is the open point recorded after
Corollary~\ref{hol:ot:cor:examples} and in Question~\ref{hol:q:nonlinear}.
Fixed-ratio exponential heights also fall outside the band argument.

\begin{question}[An exact mixed-height criterion; source~17]
\label{hol:fh:q:mixed}
Is mixed differential freedom characterized by a combination of
bounded-degree additive separation of the $\mathsf b_n$ and Zariski density of
the points $(n,\mathsf b_n)$?
\end{question}

The hypotheses are sufficient, not necessary; a useful theorem should
distinguish failure of the proof from an actual equation.

\begin{question}[Infinite coefficient alphabets; source~17]
\label{hol:fh:q:alphabets}
For pattern vectors $\mathbf v_n$ taking infinitely many values, what replaces
the finite pattern matrix?
\end{question}

$\mathbf v_n=(1,n)$ shows that the invariant must see differential shifts; one
candidate is the eventual linear geometry of the vectors
$n^a\mathsf b_n^b\mathbf v_n$, uniformly in the derivative orders.

\begin{question}[Effective relation reconstruction; source~17]
\label{hol:fh:q:effective}
Given a computable finite-alphabet stream and an oracle for its eventual
patterns, can the generators of Theorem~\ref{hol:fh:thm:relations} and small
coefficient witnesses be produced with practical complexity bounds?
\end{question}

The factorial threshold is explicit but far from economical; finite-index
certificates may be much smaller.

\begin{question}[Larger differential base fields; source~17]
\label{hol:fh:q:bases}
Which logarithmic, exponential or differentially algebraic extensions of the
bounded-support base of Corollary~\ref{hol:fh:cor:boundedjets} preserve the
independence of all intrinsic jets?
\end{question}

The endpoint is a field containing the witnesses; a sharp intermediate
boundary would connect support geometry with differential closure.

\begin{question}[Intrinsic differentiation of the codes; source~17]
\label{hol:fh:q:codes}
For a given coding scale $M_\Omega$, is there a natural differential base
containing its intrinsic derivatives over which the $\mathfrak z_A$ remain
differentially independent?
\end{question}

Multiplication by a differential unit preserves independence over a suitable
base, but enlarging the base may destroy the original independence.

\begin{question}[Infinite interpolation; source~17]\label{hol:fh:q:interpolation}
Which infinite discrete collections of prescribed jets admit a continuum of
independent strongly entire solutions in one workspace?
\end{question}

The supports forced by interpolation, not only the linear equations, must be
controlled.

\begin{question}[Specialization and chain rules; source~17]
\label{hol:fh:q:specialization}
For which substitutions $z=z(\omega)$ can mixed formal independence be
transferred to intrinsic numerical independence over a specified smaller
field?
\end{question}

Both derivations must be tracked, and specializations that collapse the
external variable, such as ordinary constants, excluded.

\begin{question}[Formalization of the finite core; source~17]
\label{hol:fh:q:formal}
Can the band lemma, multicolour polarization, the finite-pattern kernel
theorem and the constant-extension minor lemma be formalized in Lean
independently of surreal numbers, with the Hahn and Conway interpretation as a
separate layer?
\end{question}

\begin{question}[Arithmetic of the codes; source~17]
\label{hol:fh:q:arithmetic}
What can be said about divisibility, irreducibility or ideal membership of the
$\mathfrak z_A$ in $\Oz$?
\end{question}

Algebraic independence over $\mathcal F^{\mathrm{bd}}(M_\Omega)$ is not a
factorization theorem; positive support and zero constant term interact with
omnific divisibility in ways the field-theoretic result does not determine.

The next eight questions come from source~15 (its Questions~12.1--12.8), in
the notation of Section~\ref{hol:pc:sec:setting}; they are research questions
formulated there, not claims of literature-wide open status.

\begin{question}[Critical composition weight; source~15]
\label{hol:pc:q:critical}
Classify the equations \eqref{hol:pc:eq:main} in which some competing
monomials have weight exactly $\bar sd$. Can the leading coefficient, the
indicial factors and the next face of the coefficient Newton polygon replace
the strict inequality in significant subclasses?
\end{question}

The proof fails at one identified place: the coefficient of $\gamma_\sigma(f)$
in \eqref{hol:pc:eq:dominance} becomes zero, and the remaining expression must
be controlled rather than the strict argument reused. The elementary equation
$f(z^d)=f(z)^d$ is classified in Example~\ref{hol:pc:ex:critical}; its
polynomial perturbations and multi-term critical identities are the natural
next tests.

\begin{question}[Several top products; source~15]
\label{hol:pc:q:topproducts}
If the largest-degree pullback occurs in a differential polynomial rather than
in a single product of nonzero linear differential transforms, which explicit
nonvanishing conditions on its indicial or initial polynomial still force an
exact top Gauss profile?
\end{question}

The corner-cancellation issues for ordinary nonlinear equations are those of
Sections~\ref{hol:nl:sec:first}--\ref{hol:nr:sec:boundary}; the new ingredient is
their interaction with distinct polynomial arguments. The question is not the
unrestricted higher-order problem.

\begin{question}[Matrix differential pullbacks; source~15]
\label{hol:pc:q:matrix}
Does Theorem~\ref{hol:pc:thm:matrix} extend to matrix differential operators
whose maximal-shift indicial matrix has nonzero determinant? What is the
criterion when different rows have different shifts?
\end{question}

A promising route introduces row and column shifts, normalizes a matrix
initial form and tracks the integer roots of its determinant; lower-shift
couplings must be accounted for, and entrywise use of the scalar result does
not suffice.

\begin{question}[Several variables; source~15]
\label{hol:pc:q:variables}
For a polynomial map $P:\KK^r\to\KK^r$, can a positive weight vector and a
strictly expanding Newton-polytope face replace the scalar degree $d$? Which
finite-map or noncancellation assumptions make the pullback of the active
initial polynomial injective?
\end{question}

The scalar proof uses the injective substitution $X\mapsto\bar cX^d$;
multivariable leading forms can satisfy algebraic relations, a genuine
obstruction.

\begin{question}[Meromorphic solutions; source~15]
\label{hol:pc:q:meromorphic}
For a quotient of strongly entire series satisfying a differential--Mahler
equation, when does the conclusion improve from a finite coefficient reduction
to rationality? Can pole divisors and pullback degrees replace the single
Gauss profile?
\end{question}

Classical Mahler theory gives a rationality precedent~\cite{bcr}; arbitrary
rank and strong summation need a separate divisor and support argument.

\begin{question}[Minimal universal encoders; source~15]
\label{hol:pc:q:encoders}
What are the least differential order, polynomial degree bound and number of
pullback arguments of an undecidable omnific-coefficient subclass? Can the
Euler projectors be replaced by a smaller sparse operator without losing an
explicit bound?
\end{question}

The bound ten of Corollary~\ref{hol:pc:cor:ten} is inherited from a known
integer encoding and not claimed optimal; a lower bound would have to separate
analytic rigidity from arithmetic universality.

\begin{question}[Arithmetic subclasses beyond linear equations; source~15]
\label{hol:pc:q:arithmetic}
Which nonlinear subclasses have decidable omnific solvability despite
Corollary~\ref{hol:pc:cor:undecidable}? Natural targets are quadratic
coefficient loci of fixed dimension, equations with a prescribed rational
parametrization, and cases in which positivity forces ordinary coefficients.
\end{question}

A general algorithm must avoid the universal encoder explicitly; a finite
degree bound alone is not such a condition.

\begin{question}[Formal verification and certified solvers; source~15]
\label{hol:pc:q:formal}
Can the scalar weighted theorem and the universal encoder be formalized with
the repository's Hahn summability and polynomial infrastructure? Can a solver
emit a compact certificate of operator shifts, indicial polynomials, a degree
bound and a verified finite coefficient system?
\end{question}

The dependency order proposed by source~15 is in
Section~\ref{hol:pc:sub:boundaries}. Its suggested research order: first the
matrix-indicial criterion, then critical-weight equations with one lower
argument; the several-variable and meromorphic problems need further
geometric infrastructure. The universal encoder supplies adversarial tests
for symbolic solvers: a program claiming to decide all nonlinear omnific
equations of this class must be incorrect, however well it computes degree
bounds on examples.

The next eight questions come from source~20 (its Questions~12.3, 12.4,
12.6--12.10 and~12.12, the fourth re-scoped by the merge), in the notation of
Section~\ref{hol:rc:sec:setting}; its Questions~12.1, 12.2, 12.5 and~12.11 are
recorded in the statuses of Questions~\ref{hol:ud:q:effective},
\ref{hol:ud:q:torsionblocks}, \ref{hol:ud:q:beyond} and~\ref{hol:ud:q:formal}.

\begin{question}[Nonunit coefficient tests without arbitrary coding;
source~20]\label{hol:rc:q:nonunit}
Which restrictions on the initial support or on the observed coefficient cut
restore eventual periodicity of coefficient tests for recurrences whose roots
are not units?
\end{question}

Finite support, supports in a fixed Archimedean subgroup and one-sided
valuation hypotheses are natural starting points.
Theorem~\ref{hol:rc:thm:coding} shows that arbitrary reverse well-ordered
omnific support cannot be allowed without further restrictions.

\begin{question}[Full omnific-membership sets; source~20]
\label{hol:rc:q:membership}
For real-surreal $y_n$ whose characteristic roots are units,
Section~\ref{hol:rc:sub:compression} controls the absence of negative Conway
exponents. Which further hypotheses make the integer constant-term condition,
and so the set $\{n:y_n\in\Oz\}$, effectively classifiable?
\end{question}

This is the intersection of a support-theoretic problem with an arithmetic
integrality problem for an exponential polynomial, not another application of
coefficientwise algebraicity.

\begin{question}[Multivariate time; source~20]\label{hol:rc:q:multivariate}
For $\mathbf n\in\N^{\bar r}$, what are the coefficients and valuations of
$\sum_jR_j(\mathbf n)\prod_i\lambda_{ij}^{n_i}$?
\end{question}

The coefficient-algebra reduction still suggests finitely generated
observable algebras, but one-variable arithmetic progressions must be replaced
by a finer description of lattice subsets; polynomial conditions in several
integer variables warn against a naive semilinearity assertion.

\begin{question}[Nonlinear mixed equations at unit dilation; source~20,
re-scoped]\label{hol:rc:q:nonlinear}
Suppose $\Gamma$ has an order unit and $q\in\KK^\times$ is a nontorsion unit.
What are the least differential order and the least degree of a nonzero
algebraic relation over $\KK(z)$ among finitely many $(D_z^rf)(q^\ell z)$,
involving at least two powers of $q$, that has a nonpolynomial strongly entire
solution?
\end{question}

As posed by source~20 the question asked for the smallest nonlinear
differential--dilation classes admitting such solutions with an order unit,
measured by differential order and nonlinear degree, noting that its linear
theorem does not extrapolate to nonlinear equations. In the current text
unrestricted nonlinear rigidity is already refuted
(Remark~\ref{hol:ud:rem:nonlinearq}): $\mathcal T_p$ satisfies a relation of
order three, of degree six, with no dilation, and for relations without
dilations the least order is three (Theorem~\ref{hol:ot:main:second}) and the
least degree at order three is at least four (Theorem~\ref{hol:nr:main:third});
which series attain order three is Question~\ref{hol:q:nonlinear}. The merge
re-scopes source~20's question to what these do not decide: whether a unit
dilation lowers the order or the degree needed.

\begin{question}[Derivations of the coefficient field; source~20]
\label{hol:rc:q:derivations}
Replace $D_z$, which fixes $\KK$, by an operator that also uses a specified
series derivation of $\KK$, such as a compatible restriction of the
Berarducci--Mantova derivation. Which valuation-cost estimates survive?
\end{question}

Derivatives of coefficients introduce new support dynamics that the present
normal-form computation does not cover.

\begin{question}[Definability of the coding observation; source~20]
\label{hol:rc:q:definability}
In which expansions of $(\No,\Oz)$ can the coefficient map $[\omega^\omega]$
and the orbit $n\mapsto\omega^n\mathfrak w$ be defined uniformly?
\end{question}

The coding theorem becomes a model-theoretic interpretation only after these
definability questions are answered; the normal-form identity alone
interprets no second-order arithmetic in the pure ring.

\begin{question}[Uniform complexity bounds for finite observations;
source~20]\label{hol:rc:q:complexity}
For a prescribed finite set of coefficient coordinates and bounded polynomial
degrees, can one bound the dimension, the recurrence order, or the number of
possible exceptional zeros of the resulting observation sequences?
\end{question}

Such bounds can be useful even when the last exceptional index is not
effectively known; by Example~\ref{hol:rc:ex:transient} any bound must depend
on the complexity of the observation.

\begin{question}[Beyond scalar commuting dilations; source~20]
\label{hol:rc:q:substitutions}
What replaces the valuation theory for finite linear substitutions in several
variables and for finite-dimensional matrix dilations?
\end{question}

The one-variable proof relies on scalar exponential-polynomial coefficients;
noncommuting substitutions may need an orbit theorem for matrix products
rather than a direct reuse of the Skolem--Mahler--Lech theorem, and
simultaneously triangularizable cases are a tractable first step. Compare
Questions~\ref{hol:ud:q:several} and~\ref{hol:pc:q:variables}.

These questions are directions suggested by the proved statements, not claims
that they were already posed in the published literature, and their
formulations are deliberately narrower than a general theory of nonlinear
surcomplex differential equations.

%======================================================================
% Coefficient-field part: sources 11 and 12. Labels use hol:cf:.
%======================================================================
\section{All-order rigidity without an order unit: setting and valuations}
\label{hol:cf:sec:setting}

This section and the next two contain the material of sources~11 and~12,
two manuscripts delivered together after the rest of this report was
written. They answer Question~\ref{hol:q:nonlinear} in every differential
order when $\Gamma$ has no order unit, including the equations with vanishing
residue corner that Sections~\ref{hol:nl:sec:setting}--\ref{hol:nl:sec:consequences}
leave open, and source~12 answers the same case of
Question~\ref{hol:q:mixedunit}. The method is neither the escape chain of
Section~\ref{hol:sec:escape} nor the finite initial polynomial of
Section~\ref{hol:nl:sec:setting}. A finite relation forces all Taylor
coefficients of a solution into one finitely generated field
(Theorem~\ref{hol:cf:thm:coefficients}); such a field has a value group of
finite rational rank (Lemma~\ref{hol:cf:lem:valuerank}), which is trapped in a
proper convex subgroup when $\Gamma$ has no order unit
(Lemma~\ref{hol:cf:lem:principal}); and strong entireness needs coefficient
values at cofinally many scales (Proposition~\ref{hol:cf:prop:criterion}).
The warning of Remark~\ref{hol:rem:transcendence} stands: nothing here turns
a nonlinear equation into a bounded-step linear recurrence with controlled
valuations. The argument uses only that one scalar multiplier is eventually
nonzero, and no valuation of it.

\subsection{Sources, conventions and the merge's choices}
\label{hol:cf:sub:conventions}

\begin{convention}[Coefficient field of the coefficient-field part]
\label{hol:cf:conv:field}
In Sections~\ref{hol:cf:sec:setting}--\ref{hol:cf:sec:families}, $k$ is an
arbitrary field of characteristic zero with the trivial valuation, $\Gamma$ is
as in Convention~\ref{hol:conv:group}, $\KK=k((t^\Gamma))$, and $v$, strong
summability, $\Dom(f)$ and $\EK$ are as in
Convention~\ref{hol:nl:conv:field}. Write $\FE$ for the fraction field of
$\EK$ inside the formal Laurent field $\KK((z))$; the word
\emph{meromorphic} in this part means membership of $\FE$, and nothing else.
The derivative $D_z$ kills $\KK$ and extends to $\KK((z))$ by the quotient
rule. A series $u$ is \emph{differentially algebraic} over a differential
subfield $B\subseteq\KK((z))$ if $P(u,D_zu,\ldots,D_z^ru)=0$ for some finite
$r$ and some nonzero $P\in B[Y_0,\ldots,Y_r]$, and \emph{differentially
transcendental} otherwise. For $\lambda\in\KK^\times$ and $j\ge0$ put
\begin{equation}\label{hol:cf:eq:jet}
 \jet\lambda jf=(D_z^jf)(\lambda z)
 =\sum_{n\ge j}a_n\fall nj\lambda^{n-j}z^{n-j}
 \qquad\Bigl(f=\sum_na_nz^n\Bigr),
\end{equation}
differentiation \emph{before} dilation. The same notation is used for
$u\in\KK((z))$, using formal Laurent substitution. Since
$D_z^j\sigma_\lambda f=\lambda^j\sigma_\lambda D_z^jf$, the operator
$D_z^r\sigma_q^\ell$ of \eqref{hol:eq:mixed}, which dilates first, differs from
$\jet{q^\ell}r$ by the nonzero scalar $q^{\ell r}$. This changes no algebraic
independence statement, but recurrence coefficients must use one convention
throughout. An infinite family is algebraically independent when every
finite subfamily is. For a subgroup $\Delta'\subseteq\Gamma$,
$\rrk\Delta'=\dim_\Q(\Delta'\otimes_\Z\Q)$ is its rational rank; ordered
groups are torsion free, so they embed in their divisible hulls, and no
divisibility of $\Gamma$ is assumed.
\end{convention}

For $k=\C$ one has $\KK=\K$ and $\EK=\E$. As for the nonlinear part, the
generality in $k$ is recorded for this part only and is \emph{not}
transferred to Theorems~\ref{hol:main:ode}--\ref{hol:main:multi}.

\paragraph{Two sources, one spine.}
Sources~11 and~12 were written independently and prove the same five steps.
\begin{enumerate}[label=\textup{(\arabic*)}]
\item Along an actual solution, a minimal relation is linearized. The
evaluated partial derivatives have an offset $\theta$ and a nonzero
\emph{linearization multiplier} $\Xi$, the quadratic remainder has
$z$-order at least $2(n-r)>n+\theta$, and so all Taylor coefficients lie in
one field generated by finitely many of them and finitely many scalars
(source~11's Theorem~4.1 and Corollary~4.3; source~12's Theorems~4.1
and~5.2).
\item The valuation-rank inequality (source~11's Lemma~5.1, source~12's
Lemma~3.1).
\item A subgroup of finite rational rank lies in a principal convex subgroup
$\Delta_\beta$, which is proper when there is no order unit (source~11's
Lemma~5.2, source~12's Lemma~3.2).
\item An exterior obstruction at $t^{-\eta}$ with $\eta>\Delta_\beta$
(source~11's Proposition~5.4, source~12's Proposition~3.4).
\item The same worked example: along the solutions $az^m$ of the
vanishing-corner equation \eqref{hol:nl:eq:degenerate} the linearization
multiplier is a nonzero multiple of $(T-m)^2$.
\end{enumerate}
Both assume only characteristic zero, $\Gamma\ne0$ set-sized, no divisibility,
and no order unit where that is used; neither has weaker hypotheses than the
other. The spine is printed once, with both sources credited at each
statement. \textbf{Source~12 is the base} of this part: its all-order theorem
with dilations (Theorem~\ref{hol:cf:main:entire}\,(b) here) contains, for the
single dilation $1$, source~11's statement about entire functions; it answers
the no-order-unit case of Question~\ref{hol:q:mixedunit}; and its order-unit
dichotomy (Theorem~\ref{hol:cf:main:sharp}) gives the exact boundary of the
mixed statement. Source~11 contributes what source~12 lacks: coarsening
descent for arbitrary, not necessarily finitely generated, coefficient fields
(Theorem~\ref{hol:cf:thm:descent}); rationality of every differentially
algebraic element of $\FE$ (Theorem~\ref{hol:cf:main:entire}\,(c)); the
rational-system and Painlev\'e corollaries; the joint-independence criterion
(Theorem~\ref{hol:cf:thm:indcriterion}); and a continuum-sized independent
family (Theorem~\ref{hol:cf:thm:continuum}). Two of source~11's proofs differ
genuinely from source~12's and are kept as marked second routes:
its choice of a minimal relation (Remark~\ref{hol:cf:rem:separant}) and its
proof of polynomiality by coarsened reduction (proof of
Theorem~\ref{hol:cf:main:entire}). Every proof taken from source~11 was
re-read against the statement it is attached to: source~11 has no stronger
standing hypothesis than source~12; its cardinal equality uses $k=\C$, while
its lower bound holds for every $k$; and its coefficient principle is
relative to a Laurent field $k_0((z))$, a form kept because the
joint-independence criterion needs it.

\paragraph{Renamed symbols.}
The following symbols of the two sources would collide with symbols of
Sections~\ref{hol:sec:intro}--\ref{hol:sec:status} or of
\path{docs/NOTATION.md}. They are renamed here, in the new material only.
\begin{center}
\small
\begin{tabularx}{\textwidth}{@{}>{\raggedright\arraybackslash}p{0.2\textwidth}
>{\raggedright\arraybackslash}p{0.17\textwidth}
>{\raggedright\arraybackslash}p{0.17\textwidth}X@{}}
\toprule
Source~11 & Source~12 & Here & Reason\\
\midrule
$D$ & $D$ & $D_z$ & $D$ is a scalar derivation in the notation guide;
Definition~\ref{hol:def:operators}.\\
$K$, $\mathcal E_\Gamma$ & $K$, $\mathcal E$ & $\KK$, $\EK$ &
Convention~\ref{hol:nl:conv:field}; $\E$ is over $\C$.\\
$\mathcal M_\Gamma=\operatorname{Frac}\mathcal E_\Gamma$ & --- & $\FE$ &
$M(S)$ is a support monoid.\\
$H$, $H_\delta$ & $H$, $H_h$ & $\Delta$, $\Delta_\beta$ &
\eqref{hol:eq:delta}; $H\in\Gamma$ in Section~\ref{hol:sub:tables}.\\
$w$, $\mathcal O_w$, $\mathfrak m_w$, $L$, $\res$ & $w_H$ & $v_\Delta$,
$\mathcal O_\Delta$, $\mathfrak m_\Delta$, $k_\Delta$, $\red_\Delta$ &
$w_P$ is a corner shift; $\res$ is the residue at $v$.\\
$s$, $A_j$, $C_j$, $Q(T)$ & $\ell$, $S_{ij}$, $c_{ij}$, $B(T)$ & $\theta$,
$\Pi_{ij}$, $\xi_{ij}$, $\Xi(T)$ & $s$ is an order in
Question~\ref{hol:q:nonlinear}, $\ell$ a dilation power in
\eqref{hol:eq:mixed}, $B$ a threshold, $Q_s$ a shift polynomial.\\
$F$ (a subfield) & $L$, $F$ & $L$, $K_0$, $k_0$ & $F$ is the series
\eqref{hol:eq:explicitF}.\\
$F_\alpha$, $\sigma$, $c(\sigma)$ & --- & $\mathcal F_\alpha$, $b$,
$\code(b)$ & $F$ as above; $\sigma_q$ is a dilation.\\
$B_i$ & --- & $k_{J,i}$ & $B$ is a threshold.\\
$\Gamma_\infty=\bigoplus_{m\ge1}\Q e_m$ & $\bigoplus_{n\ge0}\Q e_n$ &
\eqref{hol:eq:gammainfty} & One group, indexed from $0$.\\
$f_*$ & $F_\omega$ & $F$ & See below.\\
$\Theta$ over $\Q$ & $\Theta_u$, $u=t^h$ & $G$, $\Theta_p$ &
Corollary~\ref{hol:nl:cor:theta}; \eqref{hol:eq:theta}.\\
--- & $\rho(f)$ & $\cvr(f)$ & $\rho$ is the residue of
Lemma~\ref{hol:nl:lem:algebra}.\\
--- & $E_h$ & $E_\eta$ & Proposition~\ref{hol:prop:expdomain}.\\
--- & $F_\tau$ & $R_\tau$ & $F$ as above.\\
--- & $s$, $\lambda,\mu$ & $m$, $\lambda,\lambda'$ & $s$ is an order; $\mu$
a dilation valuation.\\
Theorems~1.1--1.3 & Theorems~A, B, C & Theorems~\ref{hol:cf:thm:descent},
\ref{hol:cf:main:entire}, \ref{hol:cf:thm:continuum};
\ref{hol:cf:main:entire}\,(b), \ref{hol:cf:main:rank},
\ref{hol:cf:main:sharp} & Theorems~A--D of this report are different
theorems.\\
\bottomrule
\end{tabularx}
\end{center}
The shipped audit files \path{11-coarsening-differential-rigidity-*} and
\path{12-coefficient-field-rigidity-SOURCE_AUDIT.md} keep the sources'
notation and numbering.

\paragraph{Words with two meanings.}
\emph{No order unit} is not a synonym of \emph{non-Archimedean},
\emph{higher rank} or \emph{infinite rational rank}
(Section~\ref{hol:cf:sub:sharp}). \emph{Rank one} means Archimedean; the
\emph{rational rank} $\rrk$ is a $\Q$-dimension. A \emph{dilation} is always
the argument dilation $f(z)\mapsto f(\lambda z)$ with $\lambda\in\KK^\times$,
never an automorphism of the scalar field. \emph{Residue} keeps its meaning
at $v$ ($\res$, $\rho$); the residue at a coarsened valuation is written
$\red_\Delta$. The \emph{linearization multiplier} $\Xi$ depends on a
solution; it is not the residue corner polynomial $\chi_P$ or the full corner
polynomial $I_P$ of \eqref{hol:nl:eq:chiI}, which depend on the equation
alone, and not a shift polynomial $Q_s$ or $P_s$ of
Sections~\ref{hol:sec:ode}--\ref{hol:sec:mixed}. Source~12's lemma ``escape of
leading values is sufficient'' is Lemma~\ref{hol:lem:certificate}\,(a), not the
escape chain.

\paragraph{Results printed once.}
Sources~11 and~12 reprove several facts that the report already contains.
They are printed once and cited: the definition of strong summability
(Definition~\ref{hol:def:entire}); the sufficiency half of the all-scale
criterion (Lemma~\ref{hol:lem:certificate}\,(a)); closure of the entire class
under derivatives and dilations (Proposition~\ref{hol:prop:closure}, whose
proof uses no property of $\C$); the surreal normal-form embedding
(Section~\ref{hol:sub:embedding}, with the same convention
$t^\gamma\mapsto\omega^{-\gamma}$); the group $\Gamma_\infty$ and strong
entireness of $F$ (Theorem~\ref{hol:thm:explicit}); the failure of that
entireness after a dominating scale is adjoined
(Section~\ref{hol:sub:notallsurcomplex}); the domain of the formal exponential
(Proposition~\ref{hol:prop:expdomain}); and the failure in positive
characteristic (Example~\ref{hol:nl:ex:charp}). Everything else in the two
sources is printed in
Sections~\ref{hol:cf:sec:setting}--\ref{hol:cf:sec:families}, with proofs.

\paragraph{Where each numbered result went.}
Source~11: Theorems~1.1--1.3 are Theorem~\ref{hol:cf:thm:descent},
Theorem~\ref{hol:cf:main:entire}\,(a),(c) and
Theorem~\ref{hol:cf:thm:continuum}; Lemma~2.2 is
Proposition~\ref{hol:cf:prop:criterion}, Remark~2.3 the paragraph after it,
Proposition~2.4 and Lemma~2.5 are Proposition~\ref{hol:cf:prop:ring} and
Lemma~\ref{hol:cf:lem:min}; Section~3 is Section~\ref{hol:cf:sub:descent};
Theorem~4.1 is Theorem~\ref{hol:cf:thm:coefficients}\,(a) with
Remark~\ref{hol:cf:rem:separant}, Remark~4.2 is
Remark~\ref{hol:cf:rem:singular}, Corollary~4.3 is
Corollary~\ref{hol:cf:cor:finitegen}, Example~4.4 is
Example~\ref{hol:cf:ex:corner}, and Example~4.5 follows it;
Lemmas~5.1--5.2, Remark~5.3 and Proposition~5.4 are
Lemmas~\ref{hol:cf:lem:valuerank}--\ref{hol:cf:lem:principal},
Remark~\ref{hol:cf:rem:cancellation} and
Proposition~\ref{hol:cf:prop:exterior}; Corollaries~6.1--6.3 and Example~6.4
are Corollaries~\ref{hol:cf:cor:relative}--\ref{hol:cf:cor:painleve} and the
paragraph after the last; Section~7 is Section~\ref{hol:cf:sub:continuum};
Section~8 is Section~\ref{hol:cf:sub:surreal}; Questions~9.1 and~9.3 are
Questions~\ref{hol:q:nonlinear} (re-scoped) and~\ref{hol:cf:q:descent}, and
Example~9.2 is discussed in Section~\ref{hol:cf:sub:sharp}. Source~12:
Theorems~A--C are Theorems~\ref{hol:cf:main:entire}\,(a),(b),
\ref{hol:cf:main:rank} and~\ref{hol:cf:main:sharp}; Proposition~2.3,
Corollary~2.4 and Proposition~2.5 are
Proposition~\ref{hol:cf:prop:criterion}, Corollary~\ref{hol:cf:cor:cofinality}
and Proposition~\ref{hol:cf:prop:ring}; Lemmas~3.1--3.2, Remark~3.3,
Proposition~3.4 and Corollaries~3.5--3.6 are
Lemmas~\ref{hol:cf:lem:valuerank}--\ref{hol:cf:lem:principal},
Remark~\ref{hol:cf:rem:cancellation}, Proposition~\ref{hol:cf:prop:exterior}
and Corollaries~\ref{hol:cf:cor:finitefield}--\ref{hol:cf:cor:cvr};
Theorems~4.1 and~5.2 are Theorem~\ref{hol:cf:thm:coefficients}\,(a) and~(b),
Remarks~4.2 and~5.3 are Remarks~\ref{hol:cf:rem:singular}
and~\ref{hol:cf:rem:existential}, Corollary~4.3 is
Theorem~\ref{hol:cf:main:entire}\,(a), Corollary~4.4 is
Corollary~\ref{hol:cf:cor:DAexterior}, Lemma~5.1 is
Lemma~\ref{hol:cf:lem:exppoly} and Corollary~5.4 is
Corollary~\ref{hol:cf:cor:free}; Proposition~6.1 and Theorem~6.2 are
Proposition~\ref{hol:cf:prop:theta} and Theorem~\ref{hol:cf:thm:rankone};
Theorem~7.1 and Example~7.2 are Theorem~\ref{hol:cf:thm:Fomega} and the
paragraph after it; Propositions~8.1 and~8.2 are
Proposition~\ref{hol:cf:prop:scalar} and the whole-class failure of
Section~\ref{hol:cf:sub:surreal}; Section~9 is
Section~\ref{hol:cf:sub:boundaries}; Questions~10.1--10.3 are
Questions~\ref{hol:q:nonlinear} (re-scoped), \ref{hol:cf:q:finiterank}
and~\ref{hol:cf:q:torsion}.

\paragraph{The explicit series, and an indexing discrepancy.}
Source~12's $F_\omega$ is exactly $F$ of \eqref{hol:eq:explicitF}, over the
same group indexed from~$0$. Source~11 indexes the group from~$1$, maps $e_m$
to $\omega^m$, and writes $f_*=\sum_{n\ge1}\omega^{-\omega^n}z^n$, calling it
``this exact entire example'' of the report. That is off by a constant: the
term $n=0$ of $F$ is $t^{e_0}=\omega^{-1}$, so $F=f_*+\omega^{-1}$. The
discrepancy is harmless, because $F$ and $f_*$ differ by an element of the
field, their derivatives agree, and a polynomial $P(Y_0,Y_1,\ldots)$
vanishes on the jets of $F$ exactly when $P(Y_0+\omega^{-1},Y_1,\ldots)$
vanishes on those of $f_*$. It is disclosed here and not corrected in the
shipped audit files. Source~11's branch family uses only codes at least~$2$,
so it is printed over the group \eqref{hol:eq:gammainfty} unchanged.

\subsection{Strong entireness at every scale}\label{hol:cf:sub:criterion}

\begin{proposition}[All-scale coefficient criterion; sources~11 and~12]
\label{hol:cf:prop:criterion}
For $f=\sum_na_nz^n\in\KK[[z]]$ the following are equivalent:
\begin{enumerate}[label=\textup{(\roman*)}]
\item $f\in\EK$;
\item the family $(a_nt^{n\gamma})_{n\ge0}$ is strongly summable for every
$\gamma\in\Gamma$;
\item for all $\gamma,\beta\in\Gamma$, only finitely many $n$ with $a_n\ne0$
satisfy $v(a_n)+n\gamma\le\beta$.
\end{enumerate}
\end{proposition}

\begin{proof}
(i)$\Rightarrow$(ii) is immediate. For (ii)$\Rightarrow$(iii) fix $\gamma$
and suppose that infinitely many nonzero coefficients have
$A_n:=v(a_n)+n\gamma\le\beta$. Strong summability at $t^\gamma$ gives a common
lower bound $\alpha$ for all $A_n$, the least element of the union of the
supports. Choose any $\eta>0$ in $\Gamma$ and set
$\epsilon=\max(\beta-\alpha,0)+\eta$. Then $\epsilon>0$ and
$\epsilon>\beta-\alpha$, without division or an order-unit assumption.
For two selected indices $n'<n$,
\[
 (A_n-n\epsilon)-(A_{n'}-n'\epsilon)\le\beta-\alpha-(n-n')\epsilon<0,
\]
so at the more exterior argument $t^{\gamma-\epsilon}$ the leading exponents
of the selected summands form a strictly decreasing sequence, contradicting
well-ordered support. For (iii)$\Rightarrow$(i), a nonzero $x$ with
$v(x)=\gamma$ gives summands $a_nx^n$ of valuation $v(a_n)+n\gamma$, and
Lemma~\ref{hol:lem:certificate}\,(a), whose proof uses no property of $\C$,
applies; at $x=0$ only one summand is nonzero.
\end{proof}

The sufficiency half is Lemma~\ref{hol:lem:certificate}\,(a). Source~11
describes the criterion as part of the collection's support framework and
source~12 as part of its entire-function programme~\cite{repoentire}; both
prove it, and it is printed here once, with the necessity half, which
Sections~\ref{hol:sec:intro}--\ref{hol:sec:certificates} did not need. The
necessity proof uses a second, more exterior argument: strong summability at
one argument does not by itself give cofinal growth. In $k((t^\Q))$ the family
$(t^{n/(n+1)})_{n\ge1}$ is strongly summable although its valuations are
bounded (source~12), and Remark~\ref{hol:rem:positive} gives a higher-rank
version. Replacing the all-scale hypothesis by evaluation on one ordinary
complex disk would therefore be an error (source~11).

\begin{corollary}[When the statements are nonvacuous; source~12]
\label{hol:cf:cor:cofinality}
A nonpolynomial $f\in\EK$ exists if and only if $\Gamma$ has countable
cofinality.
\end{corollary}

\begin{proof}
If $f$ is nonpolynomial, (iii) with $\gamma=0$ makes the values of its
infinitely many nonzero coefficients a countable cofinal subset; a nonzero
ordered abelian group has no largest element, so its cofinality is not finite.
Conversely choose a nondecreasing cofinal sequence of positive $\delta_n$,
$n\ge1$, and put $a_0=1$, $a_n=t^{n\delta_n}$. Given $\gamma,\beta$, eventually
$\delta_n+\gamma>\max(\beta,0)$, so $n(\delta_n+\gamma)>\beta$, and
Proposition~\ref{hol:cf:prop:criterion} applies.
\end{proof}

The equivalence is background from the entire-function
report~\cite{repoentire} (Section~\ref{hol:sub:position}); source~12's proof is
printed because the next sections use it. Theorems~\ref{hol:cf:main:entire}
and~\ref{hol:cf:main:sharp} therefore have nonpolynomial instances exactly for
groups of countable cofinality; in uncountable cofinality they remain true but
are vacuous.

\begin{proposition}[Differential algebra structure; source~11]
\label{hol:cf:prop:ring}
$\EK$ is a differential $\KK$-subalgebra of $\KK[[z]]$, stable under every
dilation $\sigma_\lambda$, $\lambda\in\KK^\times$. Its fraction field $\FE$
is a differential subfield of $\KK((z))$.
\end{proposition}

\begin{proof}
Closure under sums, products, $D_z$ and $\sigma_\lambda$ is
Proposition~\ref{hol:prop:closure}, whose proof uses no property of $\C$.
Through Proposition~\ref{hol:cf:prop:criterion} it can also be read off the
weighted values: for a product, if $A_i=v(a_i)+i\gamma$ and
$B_j=v(b_j)+j\gamma$ have lower bounds $m_A,m_B$, then $A_i+B_j\le\beta$
forces $A_i\le\beta-m_B$ and $B_j\le\beta-m_A$, which only finitely many pairs
satisfy. If a nonzero product coefficient at degree $n$ has weighted value
at most $\beta$, one of its finitely many nonzero terms $a_i b_j$, $i+j=n$,
must have weighted value at most $\beta$. Thus only finitely many degrees
can occur, even when some terms cancel. A zero factor is immediate.
For $\jet\lambda jf$ the coefficient of $z^n$ is
$a_{n+j}\fall{n+j}j\lambda^n$, whose weighted value at $\gamma$ is that of
$a_{n+j}$ at $v(\lambda)+\gamma$ minus the constant $j(v(\lambda)+\gamma)$.
Here the falling factorial is a nonzero element of the trivially valued
prime field, so it contributes no valuation (source~12). The quotient rule
extends $D_z$ to $\FE$, and every nonzero
power series is invertible in $\KK((z))$, so $\FE$ embeds there.
\end{proof}

\begin{lemma}[A minimum coefficient value; source~11]\label{hol:cf:lem:min}
If $0\ne h=\sum_nh_nz^n\in\EK$, some $h_j$ satisfies $v(h_j)\le v(h_n)$ for
every nonzero $h_n$.
\end{lemma}

\begin{proof}
Take any nonzero coefficient value $\beta$. By
Proposition~\ref{hol:cf:prop:criterion} with $\gamma=0$ only finitely many
coefficients have value at most $\beta$; the least of them is the minimum.
\end{proof}

\subsection{Value groups of finitely generated fields}
\label{hol:cf:sub:valuation}

\begin{lemma}[Valuation-rank inequality; sources~11 and~12]
\label{hol:cf:lem:valuerank}
Let $k_0\subseteq L$ be fields and $v$ a valuation on $L$, with values in an
ordered abelian group, trivial on $k_0^\times$. If $x_1,\ldots,x_d\in L^\times$
have values linearly independent over $\Q$, they are algebraically independent
over $k_0$. Consequently
\begin{equation}\label{hol:cf:eq:valuerank}
 \rrk v(L^\times)\le\trdeg_{k_0}L.
\end{equation}
\end{lemma}

\begin{proof}
For distinct $\alpha,\alpha'\in\N^d$ the monomials $x^\alpha$ and
$x^{\alpha'}$ have different values, since equality would be an integer
relation among the $v(x_i)$. A nonzero coefficient from $k_0$ has value zero.
So the finitely many nonzero terms of a nonzero polynomial evaluated at the
$x_i$ have pairwise distinct values, and the unique term of least value
cannot be cancelled. For the rank inequality, choose a rational basis of
$v(L^\times)\otimes_\Z\Q$ from the images of elements of $v(L^\times)$,
and lift those values to elements of $L^\times$. Every polynomial uses only
finitely many of the lifts, so the preceding argument proves that the whole
lifted family is algebraically independent. Its cardinality is the rational
rank, giving the inequality also when either rank is infinite.
\end{proof}

Source~12 states the lemma for $k\subseteq L\subseteq\KK$, source~11 for
$k_0=\Q$ and any valuation trivial on $\Q$; both cases are used below.

\begin{lemma}[A principal convex bound; sources~11 and~12]
\label{hol:cf:lem:principal}
Let $\Delta'\subseteq\Gamma$ be a subgroup of finite rational rank. Either
$\Delta'=0$, or $\Delta'\subseteq\Delta_\beta$ for some $\beta>0$, where
$\Delta_\beta$ is the convex subgroup \eqref{hol:eq:delta}; in the second case
$\beta$ can be chosen in $\Delta'$, and $\Delta_\beta$ is then the smallest
convex subgroup containing $\Delta'$. The subgroup $\Delta_\beta$ is proper
exactly when $\beta$ is not an order unit, so it is proper for every $\beta$
when $\Gamma$ has no order unit.
\end{lemma}

\begin{proof}
Choose $\delta_1,\ldots,\delta_d\in\Delta'$ whose images form a basis of
$\Delta'\otimes\Q$ and put $\beta=\sum_i\abs{\delta_i}>0$, an element of
$\Delta'$. For $\delta\in\Delta'$, clearing denominators gives
$m\delta=\sum_in_i\delta_i$ with a positive integer $m$ and integers $n_i$,
so $\abs\delta\le m\abs\delta\le(\sum_i\abs{n_i})\beta$. Hence
$\Delta'\subseteq\Delta_\beta$, and $\Delta_\beta$, the smallest convex
subgroup containing $\beta\in\Delta'$, lies in every convex subgroup
containing $\Delta'$. The last assertion is the definition of an order unit.
\end{proof}

No division in $\Gamma$ occurs. Source~11 takes
$\delta=\eta+\sum_i\abs{\delta_i}$ with an arbitrary $\eta>0$, which covers
$\Delta'=0$ by the same formula.

\begin{remark}[Field values, not generator values; sources~11 and~12]
\label{hol:cf:rem:cancellation}
The bound \eqref{hol:cf:eq:valuerank} concerns the value group of the whole
field. It would be false to use the values of a chosen list of generators:
for $0<\alpha<\beta$ with $\alpha,\beta$ rationally independent, both
generators of $\Q(t^\alpha,t^\alpha+t^\beta)$ have value $\alpha$, but their
difference has value $\beta$ (source~11). Conversely, one generator $b\in\KK$
may have infinitely many rationally independent exponents in its support
without contradicting $\rrk v(k(b)^\times)\le1$: the field $k(b)$ contains
rational functions of $b$, not the truncations of its Hahn expansion
(source~12). This is why the coefficients of the equations below may be
arbitrary Hahn elements.
\end{remark}

For a convex subgroup $\Delta\subseteq\Gamma$ the quotient $\Gamma/\Delta$ is
an ordered abelian group, and
\[
 v_\Delta(a)=v(a)+\Delta\in\Gamma/\Delta\qquad(a\in\KK^\times),
 \qquad v_\Delta(0)=+\infty,
\]
is a valuation on the same field $\KK$, a coarsening of $v$, as in
Theorem~\ref{hol:thm:coarse}; it is not a coefficientwise map to a smaller
Hahn field. Write $\mathcal O_\Delta=\{a:v_\Delta(a)\ge0\}$,
$\mathfrak m_\Delta=\{a:v_\Delta(a)>0\}$, $k_\Delta=\mathcal O_\Delta/\mathfrak
m_\Delta$ and $\red_\Delta:\mathcal O_\Delta\to k_\Delta$. No splitting of
$\Gamma\to\Gamma/\Delta$ and no identification of $k_\Delta$ with a Hahn field
is used. Write $\eta>\Delta$ when $\eta>\delta$ for every $\delta\in\Delta$;
by convexity every positive element outside $\Delta$ satisfies this, and
$v_\Delta(x)<0$ means exactly $v(x)=-\eta$ with $\eta>\Delta$.

\begin{proposition}[Coarsened exterior obstruction; sources~11 and~12]
\label{hol:cf:prop:exterior}
Suppose all nonzero coefficients of a nonpolynomial $f=\sum_na_nz^n\in\KK[[z]]$
have values in a proper convex subgroup $\Delta$. Then
\begin{equation}\label{hol:cf:eq:exterior}
 \Dom(f)\subseteq\{0\}\cup\{x\in\KK^\times:v_\Delta(x)\ge0\};
\end{equation}
in particular $t^{-\eta}\notin\Dom(f)$ for every $\eta>\Delta$.
\end{proposition}

\begin{proof}
Let $v(x)=-\eta$ with $\eta>\Delta$. For nonzero coefficients at indices
$m>n$,
$v(a_mx^m)-v(a_nx^n)=\bigl(v(a_m)-v(a_n)\bigr)-(m-n)\eta<0$, because the first
difference lies in $\Delta$. The leading exponents of the infinitely many
nonzero summands therefore strictly decrease, and
Lemma~\ref{hol:lem:nonincreasing} applies.
\end{proof}

Source~11 states this for $x=t^{-\eta}$ and coefficients in a field whose
values lie in $\Delta$; source~12 states \eqref{hol:cf:eq:exterior}. The
inclusion is necessary, not sufficient: points with $v_\Delta(x)\ge0$ need not
lie in $\Dom(f)$. For $k=\C$ it is consistent with inward stability
(Lemma~\ref{hol:lem:inward}).

\begin{corollary}[Finite transcendence degree; source~12]
\label{hol:cf:cor:finitefield}
Assume $\Gamma$ has no order unit. If all coefficients of $f\in\EK$ lie in a
subfield $L\supseteq k$ of $\KK$ with $\trdeg_kL<\infty$, then $f\in\KK[z]$.
\end{corollary}

\begin{proof}
By Lemmas~\ref{hol:cf:lem:valuerank} and~\ref{hol:cf:lem:principal},
$v(L^\times)$ lies in $\{0\}$ or in some $\Delta_\beta$, and both are proper.
Proposition~\ref{hol:cf:prop:exterior} excludes a nonpolynomial $f$.
\end{proof}

For $f=\sum_na_nz^n$ define the \emph{coefficient-value rank}
\begin{equation}\label{hol:cf:eq:cvr}
 \cvr(f)=\dim_\Q\left(\operatorname{span}_\Q
 \{v(a_n)\otimes1:a_n\ne0\}\right).
\end{equation}
The span is taken in $\Gamma\otimes_\Z\Q$, and $\cvr(f)$ is its dimension.
It is the rational rank of the values of the Taylor coefficients, not the
number of support exponents of any single coefficient.

\begin{corollary}[Necessary infinite coefficient-value rank; source~12]
\label{hol:cf:cor:cvr}
If $\Gamma$ has no order unit and $f\in\EK$ is not a polynomial, then
$\cvr(f)=\aleph_0$ and $\trdeg_kk(a_0,a_1,\ldots)=\aleph_0$.
\end{corollary}

\begin{proof}
If $\cvr(f)$ were finite, the group generated by the coefficient values would
lie in a proper $\Delta_\beta$ by Lemma~\ref{hol:cf:lem:principal},
contradicting Proposition~\ref{hol:cf:prop:exterior}. There are countably
many coefficients, which bounds the rank; Lemma~\ref{hol:cf:lem:valuerank}
gives the lower bound for the transcendence degree and countable generation
the upper bound.
\end{proof}

\section{The coefficient-field theorem and its consequences}
\label{hol:cf:sec:coefficients}

\subsection{One finitely generated field for every finite relation}
\label{hol:cf:sub:coefficients}

Coefficient recurrences for algebraic differential equations go back to
Hurwitz~\cite{hurwitz}, and Lemma~3 of Laohakosol, Kongsakorn and
Ubolsri~\cite{lku} isolates a high coefficient through a polynomial in its
index; Krattenthaler and Rivoal treat the arithmetic of such
coefficients~\cite{kr}, and Ait El Manssour, Sattelberger and Teguia Tabuguia
the general framework of D-algebraic functions~\cite{ast}. Only the
field-theoretic consequence is needed here. Both sources supply a
self-contained formal linearization proof over arbitrary characteristic-zero
fields that allows singular solutions, and neither advertises the
finite-generation principle or the valuation-rank inequality as a discovery.

For several dilations one classical theorem is imported and not proved: the
\emph{Skolem--Mahler--Lech theorem}, that the zero set of a linear recurrence
sequence over a field of characteristic zero is a finite set together with
finitely many infinite arithmetic progressions. Source~12 uses it in the
general-field form of Bell's Theorem~1.1~\cite{bell}, credited there to
Lech~\cite{lech}; no generalized dynamical form of the theorem is used.

\begin{lemma}[Nondegenerate exponential polynomials; source~12]
\label{hol:cf:lem:exppoly}
Let $K_0$ be a field of characteristic zero and let
$\lambda_1,\ldots,\lambda_m\in K_0^\times$ with $\lambda_i/\lambda_{i'}$ not a
root of unity for $i\ne i'$. If $p_1,\ldots,p_m\in K_0[T]$ are not all zero,
then
\begin{equation}\label{hol:cf:eq:exppoly}
 \Xi(n)=\sum_{i=1}^m\lambda_i^n\,p_i(n)
\end{equation}
is nonzero for all sufficiently large $n\in\N$.
\end{lemma}

\begin{proof}
Discard the zero $p_i$, and let $E$ be the shift $(Eu)(n)=u(n+1)$. The
sequence $\lambda_i^np_i(n)$ is annihilated by $(E-\lambda_i)^{\deg p_i+1}$, so
the product of these annihilators kills $\Xi$: they commute because they
are polynomials in the same shift $E$. Thus $\Xi$ satisfies a linear
recurrence with constant coefficients in $K_0$, and
by the imported theorem its zero set is a finite set together with finitely
many infinite arithmetic progressions. No infinite progression can occur.
If $\Xi(a+bu)=0$ for all $u\ge0$, with $b\ge1$, then
$\sum_i(\lambda_i^b)^u\lambda_i^ap_i(a+bu)=0$ for all $u$, and the bases
$\alpha_i=\lambda_i^b$ are pairwise distinct by the root-of-unity hypothesis.
Polynomial-exponential sequences with pairwise distinct nonzero bases are
linearly independent: to isolate the base $\alpha_i$, apply
$\prod_{i'\ne i}(E-\alpha_{i'})^{d_{i'}+1}$, where $d_{i'}$ is the degree of the
polynomial at $\alpha_{i'}$; this kills the other terms, and on
$\alpha_i^up(u)$ each factor keeps the degree and multiplies the leading
coefficient by $\alpha_i-\alpha_{i'}\ne0$, leaving a nonzero power times a
nonzero polynomial, which cannot vanish at every natural number in
characteristic zero. Hence every $p_i(a+bT)$, and so every $p_i$, is zero, a
contradiction.
\end{proof}

For $m=1$ the lemma is the finiteness of the roots of a nonzero polynomial,
since $\lambda_1^n\ne0$. When only the $j=0$ jets occur and
$\lambda_i=q^{m_i}$ with distinct integers $m_i$ and a nontorsion $q$, the
multiplier is a nonzero Laurent polynomial evaluated at the distinct elements
$q^n$, again without the imported theorem. The general mixed-dilation
argument uses the imported theorem.

\begin{theorem}[Finite coefficient field; sources~11 and~12]
\label{hol:cf:thm:coefficients}
Let $k_0\subseteq K_0$ be fields of characteristic zero and
$u=\sum_{n\ge0}a_nz^n\in K_0[[z]]$, with $D_z$ killing $K_0$.
\begin{enumerate}[label=\textup{(\alph*)}]
\item \textup{(Sources~11 and~12.)} If $u$ is differentially algebraic over
$k_0((z))$, there is $N\ge0$ with $a_n\in k_0(a_0,\ldots,a_{N-1})$ for every
$n$.
\item \textup{(Source~12.)} Let $\lambda_1,\ldots,\lambda_m\in K_0^\times$
have pairwise quotients that are not roots of unity. If finitely many of the
series $\jet{\lambda_i}ju$ are algebraically dependent over $K_0(z)$, then all
$a_n$ lie in one subfield $L\subseteq K_0$ finitely generated over $k_0$,
generated by the $\lambda_i$, the finitely many coefficients of a suitable
polynomial relation, and finitely many $a_n$.
\end{enumerate}
\end{theorem}

\begin{proof}
Treat both parts at once, first choosing a relation involving derivative
orders at most $r$. In case~(a) put $m=1$, $\lambda_1=1$ and
$R=k_0[[z]]$, and multiply the relation by a sufficiently large power of $z$
to remove the finitely many Laurent poles of its coefficients. In case~(b)
put $R=K_0[z]$ and clear a common polynomial denominator of its finitely many
rational-function coefficients. In either case there is a nonzero
polynomial over $R$ in variables $Y_{ij}$, $1\le i\le m$, $0\le j\le r$,
vanishing at $Y_{ij}=\jet{\lambda_i}ju$. Among these relations choose $P$ of
least total degree in the $Y_{ij}$. Its degree is positive, since a nonzero
element of $R$ does not vanish in $K_0[[z]]$. Put
\[
 \Pi_{ij}=\frac{\partial P}{\partial Y_{ij}}
 \bigl((\jet{\lambda_i}ju)_{i,j}\bigr)\in K_0[[z]].
\]
A nonzero polynomial partial has nonzero evaluation, by minimality of the
total degree, and characteristic zero gives at least one. Put
\begin{equation}\label{hol:cf:eq:theta}
 \theta=\min_{\Pi_{ij}\ne0}\bigl(\ordz\Pi_{ij}-j\bigr),\qquad
 \xi_{ij}=\coeffz{\theta+j}{\Pi_{ij}},
\end{equation}
with $\xi_{ij}=0$ when $\Pi_{ij}=0$ or $\theta+j<0$, so $\theta\ge-r$ and some
$\xi_{ij}\ne0$, and put
\begin{equation}\label{hol:cf:eq:multiplier}
 \Xi(n)=\sum_{i,j}\xi_{ij}\lambda_i^{n-j}\fall nj
 =\sum_i\lambda_i^np_i(n),\qquad
 p_i(T)=\sum_j\xi_{ij}\lambda_i^{-j}\fall Tj.
\end{equation}
Falling factorials of distinct degrees are linearly independent, so
$p_i\ne0$ whenever some $\xi_{ij}\ne0$, and
Lemma~\ref{hol:cf:lem:exppoly} makes $\Xi(n)\ne0$ for all large $n$. In
case~(a), $\Xi(T)=\sum_j\xi_{1j}\fall Tj$ is a nonzero polynomial, and the
lemma is not needed.

Let $T_n=\sum_{m'<n}a_{m'}z^{m'}$ and $g_n=u-T_n$. By \eqref{hol:cf:eq:jet},
$\ordz\jet{\lambda_i}jg_n\ge n-j\ge n-r$ for $n\ge r$. Taylor expansion of $P$
at the jets of $u$ gives
\begin{equation}\label{hol:cf:eq:taylor}
 P\bigl((\jet{\lambda_i}jT_n)_{i,j}\bigr)
 =-\sum_{i,j}\Pi_{ij}\,\jet{\lambda_i}jg_n+R_n,\qquad
 \ordz R_n\ge2(n-r),
\end{equation}
because every term of $R_n$ has at least two factors $\jet{\lambda_i}jg_n$
and its other factors have nonnegative order; this holds even when the
evaluated partials vanish to high order. For $n>\theta+2r$ one has
$2(n-r)>n+\theta$, so $R_n$ contributes nothing in degree $n+\theta$.
For each nonzero $\Pi_{ij}$ the coefficient contributed by $a_{m'}$ to
degree $n+\theta$ is
\[
 a_{m'}\fall{m'}j\lambda_i^{m'-j}
 [z^{n+\theta-m'+j}]\Pi_{ij},\qquad m'\ge n.
\]
If $m'>n$, the displayed exponent is below $\theta+j$ and hence below
$\ordz\Pi_{ij}$, or is negative, so this coefficient is zero. At $m'=n$
it is $a_n\fall nj\lambda_i^{n-j}\xi_{ij}$. Terms with $\Pi_{ij}=0$
contribute nothing. Thus only the stated multiplier of $a_n$ remains:
\begin{equation}\label{hol:cf:eq:recurrence}
 a_n\,\Xi(n)=-\coeffz{n+\theta}{P\bigl((\jet{\lambda_i}jT_n)_{i,j}\bigr)}.
\end{equation}
Put $N_0=\theta+2r\ge r$. Let $L_0=k_0$ in case~(a); in case~(b), let
$L_0$ be generated over $k_0$ by the $\lambda_i$ and the finitely many
scalar coefficients of $P$. Replacing $u$ by $T_{N_0+1}$ changes every jet
of order at most $r$ only in degrees at least
$N_0+1-r=\theta+r+1$. Polynomial substitution with coefficients in $R$
preserves this congruence. Since $\xi_{ij}$ is a coefficient in degree
$\theta+j\le\theta+r$, it is determined by that truncation; thus
$\xi_{ij}\in L_0(a_0,\ldots,a_{N_0})$.

Choose $N>N_0$ so that $\Xi(n)\ne0$ for every $n\ge N$, and put
$L=L_0(a_0,\ldots,a_{N-1})$. Its generators contain the required prefix,
so $\Xi(n)\in L^\times$ for $n\ge N$. If all $a_{m'}$ with $m'<n$
already lie in $L$, then the right side of
\eqref{hol:cf:eq:recurrence} lies in $L$: in case~(a) every coefficient
of the finitely many power series occurring as coefficients of $P$ lies
in $k_0$, and in case~(b) all scalar coefficients of $P$ lie in $L_0$.
Division by $\Xi(n)$ gives $a_n\in L$, completing the induction.
In case~(a), $L=k_0(a_0,\ldots,a_{N-1})$; in case~(b), $L$ is finitely
generated over $k_0$.
\end{proof}

\begin{remark}[Second route: source~11's minimal relation]
\label{hol:cf:rem:separant}
For part~(a), source~11 chooses a relation of least order $r$ and then least
degree in $Y_r$. Its separant $\partial P/\partial Y_r$ is a nonzero
polynomial in characteristic zero, and its evaluation is nonzero, since
otherwise it would be a relation of smaller order, or of the same order and
smaller degree in $Y_r$. It then defines the offset and multiplier from the
nonzero evaluated partials exactly as in \eqref{hol:cf:eq:theta}. Source~12's
least total degree makes \emph{every} nonzero partial evaluate to a nonzero
series. The proof needs only one nonzero evaluated partial, so either choice
works; they are the two sources' routes to the same recurrence.
\end{remark}

\begin{remark}[No nonsingular initial condition; sources~11 and~12]
\label{hol:cf:rem:singular}
The theorem assumes no evaluated partial nonzero at $z=0$, no equation solved
for the highest derivative, no irreducibility and no nonzero residue corner. A nonzero
evaluated partial is used as a formal series, not through its constant term;
the offset $\theta$ absorbs its vanishing order, and the finitely many natural
zeros of $\Xi$ are placed in the initial segment. No implicit-function
theorem at an ordinary point is invoked.
\end{remark}

\begin{remark}[Existential, not an algorithm; source~12]
\label{hol:cf:rem:existential}
The invoked form of Skolem--Mahler--Lech supplies no effective bound on the
last exceptional index, and no
effective equality, leading-exponent access or computable representation of
the coefficients of $P$ is assumed. Part~(b) is an existence statement. In
part~(a) the multiplier is explicitly a polynomial, but computational claims
would still have to say how its scalar coefficients are represented.
\end{remark}

\begin{example}[Linearization along the vanishing-corner example; sources~11
and~12]\label{hol:cf:ex:corner}
For the equation \eqref{hol:nl:eq:degenerate},
$P=z^2Y_0Y_2+zY_0Y_1-z^2Y_1^2$, whose residue corner polynomial vanishes
identically, take the solution $u=az^m$ with $a\ne0$ and $m\ge0$
(Proposition~\ref{hol:nl:prop:degenerate}). Writing $\Pi_j$ for $\Pi_{1j}$, the
evaluated partials are
\[
 \Pi_0=am^2z^m,\qquad\Pi_1=a(1-2m)z^{m+1},\qquad\Pi_2=az^{m+2},
\]
so $\theta=m$ (for $m=0$ after discarding $\Pi_0=0$) and
\begin{equation}\label{hol:cf:eq:eulermultiplier}
 \Xi(T)=a\bigl(m^2+(1-2m)T+T(T-1)\bigr)=a(T-m)^2\ne0.
\end{equation}
Source~12 takes $a=1$; source~11 carries~$a$. The multiplier is nonzero although
$\chi_P=0$: it is computed along a solution, $\chi_P$ from the equation
alone, and its only exceptional index $n=m$ moves with the solution. This is
why the argument needs no degree bound for polynomial solutions, which this
equation does not have. The displayed $P$ is not of least degree for $az^m$,
which also satisfies $zD_zu-mu=0$; the computation shows that the recurrence
\eqref{hol:cf:eq:recurrence} can be run for this $P$ directly, since one of its
evaluated partials is nonzero. Neither source infers the general theorem from
the example.
\end{example}

\begin{example}[The initial segment is not bounded by the order; source~11]
For an integer $N\ge0$, the equation $zD_zu-Nu=0$ has the solution $az^N$ with
arbitrary $a$, and $\Xi(T)=T-N$. Over a field not containing $a$, the initial
segment must reach degree $N$, so the differential order alone does not bound
its length.
\end{example}

\begin{corollary}[Finitely many scalar parameters; source~11]
\label{hol:cf:cor:finitegen}
Let $K_0$ be a field of characteristic zero.
If $u\in K_0((z))$ is differentially algebraic over $K_0(z)$, all its Laurent
coefficients lie in a subfield of $K_0$ finitely generated over $\Q$.
\end{corollary}

\begin{proof}
Clear denominators in a relation; its finitely many scalar coefficients
$c_1,\ldots,c_q$ generate $k_0=\Q(c_1,\ldots,c_q)$. Write $u=z^{-m'}y$ with a
power series $y$, taking $m'\ge0$. Leibniz's formula expresses the jets of
$u$ through order $r$ as a triangular $k_0(z)$-linear combination of the
jets of $y$ through that order, with nonzero diagonal $z^{-m'}$.
This invertible substitution preserves nonzero polynomial relations, so
$y$ is differentially algebraic over $k_0(z)\subseteq k_0((z))$. Theorem~\ref{hol:cf:thm:coefficients}\,(a)
puts the coefficients of $y$, hence of $u$, in $k_0$ with finitely many
initial coefficients adjoined.
\end{proof}

\subsection{Coarsening descent}\label{hol:cf:sub:descent}

This subsection is source~11's. Its polynomial half is a coefficient-field
form of the entire-function report's obstruction for noncofinal extensions
of the workspace~\cite{repoentire}; the fraction-field half is the part that
needs a new argument.

\begin{lemma}[Polynomial reduction of an entire series; source~11]
\label{hol:cf:lem:reduction}
Let $\Delta\subsetneq\Gamma$ be convex. If $g=\sum_ng_nz^n\in\EK$ has all
coefficients in $\mathcal O_\Delta$, its coefficientwise reduction
$\overline g=\sum_n\red_\Delta(g_n)z^n$ is a polynomial in $k_\Delta[z]$. Every
nonzero $h\in\EK$ has a scalar multiple $h/c$ with coefficients in
$\mathcal O_\Delta$ and nonzero reduction.
\end{lemma}

\begin{proof}
Because $\Delta$ is proper, choose a positive element $\eta$ outside it;
convexity gives $\eta>\Delta$. By Proposition~\ref{hol:cf:prop:criterion} with
$\gamma=0$, eventually $v(g_n)>\eta$, so $v_\Delta(g_n)>0$ and
$\red_\Delta(g_n)=0$. For the second assertion take $c=h_j$ as in
Lemma~\ref{hol:cf:lem:min}: all coefficients of $h/c$ have nonnegative $v$,
and the $j$th is~$1$.
\end{proof}

\begin{lemma}[Rationality needs no new constants; source~11]
\label{hol:cf:lem:rationaldescent}
For any extension of fields $L\subseteq L'$,
$L'(z)\cap L((z))=L(z)$ inside $L'((z))$.
\end{lemma}

\begin{proof}
After multiplying by a power of $z$, let $u=\sum_{n\ge0}u_nz^n\in L[[z]]$ lie
in $L'(z)$. Being regular at zero, $u=p/q$ with $p,q\in L'[z]$ and $q(0)=1$;
put $d=\deg q$ and take $N\ge\max(d,\deg p)$. Then
$u_n+q_1u_{n-1}+\cdots+q_du_{n-d}=0$ for $n>N$. This is a linear system with
coefficients in $L$ in the finitely many unknowns $q_1,\ldots,q_d$.
The augmented rows lie in $L^{d+1}$, so finitely many of them span all the
rows. Solving that finite subsystem is therefore equivalent to solving the
whole system. Gaussian elimination uses coefficients in $L$, and cannot
produce an inconsistent row because a solution exists in $L'$. Setting its
free variables to zero gives a solution $b_1,\ldots,b_d\in L$.
Then $\widetilde q=1+\sum_{j=1}^d b_jz^j\in L[z]$ satisfies
$[z^n](\widetilde q u)=0$ for all $n>N$, so
$\widetilde q u\in L[z]$ and $u\in L(z)$. This includes $d=0$, when there
are no unknowns and $u$ is already a polynomial.
\end{proof}

\begin{theorem}[Coarsening descent; source~11]\label{hol:cf:thm:descent}
Let $L\subseteq\KK$ be a subfield and $\Delta\subsetneq\Gamma$ a convex
subgroup with $v(L^\times)\subseteq\Delta$. Then, inside $\KK((z))$,
\begin{equation}\label{hol:cf:eq:descent}
 \EK\cap L[[z]]=L[z],\qquad \FE\cap L((z))=L(z).
\end{equation}
\end{theorem}

\begin{proof}
Let $u=\sum_nu_nz^n\in\EK\cap L[[z]]$. Every nonzero $u_n$ has value in
$\Delta$, below any fixed $\eta>\Delta$, so by
Proposition~\ref{hol:cf:prop:criterion} only finitely many $u_n$ are nonzero.

Next let $u\in\FE\cap L[[z]]$, say $u=g/h$ with $g,h\in\EK$ and $h\ne0$.
Normalize $h$ by $c$ as in Lemma~\ref{hol:cf:lem:reduction}. Since $v_\Delta$
vanishes on $L^\times$, $\red_\Delta$ embeds $L$ as a subfield $\overline
L\subseteq k_\Delta$: every nonzero element of $L$ is a unit of
$\mathcal O_\Delta$, and so has nonzero residue. In the formal identity
\begin{equation}\label{hol:cf:eq:reductionidentity}
 g/c=(h/c)\,u
\end{equation}
every coefficient on the right is a \emph{finite} sum of products of elements
of $\mathcal O_\Delta$, so $g/c$ has coefficients in $\mathcal O_\Delta$; both
$g/c$ and $h/c$ lie in $\EK$. By Lemma~\ref{hol:cf:lem:reduction} their
reductions are polynomials, the second nonzero, and reducing
\eqref{hol:cf:eq:reductionidentity} coefficientwise gives
$\overline{h/c}\;\overline u=\overline{g/c}$ in $k_\Delta[[z]]$. So
$\overline u\in k_\Delta(z)\cap\overline L[[z]]$, and
Lemma~\ref{hol:cf:lem:rationaldescent} gives $\overline u\in\overline L(z)$.
Write $\overline u=\overline p/\overline q$ with
$\overline p,\overline q\in\overline L[z]$ and $\overline q\ne0$.
Lift their finitely many coefficients through the isomorphism
$L\to\overline L$ to $p,q\in L[z]$, with $q\ne0$. Every coefficient of
$qu-p$ belongs to $L$ and has zero residue; injectivity on $L$ therefore
gives $qu=p$, hence $u\in L(z)$. This lifts a rational identity over the
embedded field, not arbitrary residue coefficients from $k_\Delta$.
Finally, for $u\in\FE\cap L((z))$ choose $m'\ge0$ so that
$z^{m'}u\in L[[z]]$, apply the power-series case, and divide by $z^{m'}$.
The inclusions $\supseteq$ hold because all polynomials over $L$ are strongly
entire and their quotients belong to $\FE$.
\end{proof}

\begin{remark}[Numerator and denominator together; source~11]
The coefficients of $g$ and $h$ need not lie in $L$, and neither need be a
polynomial. The proof works because their \emph{common} coarsened reduction
gives a polynomial equation for the quotient; asserting that $g,h\in L[[z]]$
would be unjustified.
\end{remark}

\begin{example}[The formal exponential; source~11]
With $L=k$ and $\Delta=\{0\}$, Theorem~\ref{hol:cf:thm:descent} gives
$\FE\cap k((z))=k(z)$ for every nonzero $\Gamma$. The formal series
$\sum_nz^n/n!$ is not rational: if $r=p/q\ne0$ with $p,q\in k[z]$
satisfied $D_zr=r$, then $p'q-pq'=pq$. The left side has degree at most
$\deg p+\deg q-1$ (or is zero), while the right side has degree
$\deg p+\deg q$. This is impossible, so the formal exponential is not a
quotient of strongly entire series. This says nothing about the complex
exponential, which belongs to a different analytic category.
\end{example}

\subsection{Proofs of Theorems~\ref{hol:cf:main:entire}--\ref{hol:cf:main:sharp}
and further consequences}\label{hol:cf:sub:proofs}

\begin{proof}[Proof of Theorem~\ref{hol:cf:main:rank} \textup{(source~12)}]
A dependence uses finitely many $\lambda\in\Lambda$ and derivative orders.
Theorem~\ref{hol:cf:thm:coefficients}\,(b) with $k_0=k$ and $K_0=\KK$ puts
every $a_n$ in a subfield $L\subseteq\KK$ finitely generated over $k$. By
Lemma~\ref{hol:cf:lem:valuerank}, $v(L^\times)$ has finite rational rank, and
it contains every coefficient value, contradicting $\cvr(f)=\infty$.
\end{proof}

\begin{proof}[Proof of Theorem~\ref{hol:cf:main:entire}]
(b) (Source~12.) By Corollary~\ref{hol:cf:cor:cvr}, $\cvr(f)=\aleph_0$, and
Theorem~\ref{hol:cf:main:rank} applies. For $q$ of infinite order and
$\Lambda=\{q^i:i\in\Z\}$, a quotient $q^{i-i'}$ is a root of unity only if
$i=i'$.

(a) (Sources~11 and~12.) Take $\Lambda=\{1\}$ in (b). This case needs only
Theorem~\ref{hol:cf:thm:coefficients}\,(a), with $k_0$ generated over $k$ by
the finitely many scalar coefficients of a relation, and
Corollary~\ref{hol:cf:cor:finitefield}; no Skolem--Mahler--Lech. An
algebraic dependence among $f,D_zf,\ldots,D_z^rf$ becomes, after clearing
denominators, a nonzero equation $P(z,f,\ldots,D_z^rf)=0$ with $P\in\KK[z,Y]$,
and conversely. \emph{Second route (source~11):} a differentially algebraic
$f\in\EK$ has, by Corollary~\ref{hol:cf:cor:finitegen}, its coefficients in a
field $L$ finitely generated over $\Q$; by
Lemmas~\ref{hol:cf:lem:valuerank} (with $k_0=\Q$)
and~\ref{hol:cf:lem:principal}, $v(L^\times)$ lies in a proper convex
subgroup; and the polynomial half of Theorem~\ref{hol:cf:thm:descent} gives
$f\in L[z]$.

(c) (Source~11.) Let $u\in\FE$ be differentially algebraic over $\KK(z)$. As
in the second route, its Laurent coefficients lie in a field $L$ with
$v(L^\times)$ in a proper convex subgroup $\Delta$, and the fraction-field
half of Theorem~\ref{hol:cf:thm:descent} gives $u\in L(z)\subseteq\KK(z)$.
Conversely $\KK(z)\subseteq\FE$, since polynomials are entire, and rational
functions are algebraic, hence differentially algebraic, over $\KK(z)$.
\end{proof}

\begin{corollary}[A free differential algebra of dilates; source~12]
\label{hol:cf:cor:free}
Under the hypotheses of Theorem~\ref{hol:cf:main:entire}\,(b), the
homomorphism from the differential polynomial algebra
$\KK(z)\{X_\lambda:\lambda\in\Lambda\}$, with $D_z(z)=1$ and $D_z(\KK)=0$, to
$\KK((z))$ sending $X_\lambda$ to $f(\lambda z)$ is injective. For
$\Lambda=\{q^i:i\in\Z\}$ this gives countably many differentially
algebraically independent functions.
\end{corollary}

\begin{proof}
$D_z^jX_\lambda$ goes to $\lambda^j\jet\lambda jf$, and multiplying members of
an algebraically independent family by nonzero elements of the base field
keeps it independent.
\end{proof}

This is stronger than differential transcendence of each dilate separately:
for $q=2$, no nonzero polynomial over $\KK(z)$ relates $f(z)$, $D_zf(z)$,
$f(2z)$, $(D_z^2f)(4z)$ and any further finite selection from the family.

\begin{corollary}[Relative differential-algebraic closedness; source~11]
\label{hol:cf:cor:relative}
If $\Gamma$ has no order unit, no element of $\FE\setminus\KK(z)$ is
differentially algebraic over $\KK(z)$.
\end{corollary}

This is Theorem~\ref{hol:cf:main:entire}\,(c) restated. It does not say that
$\KK(z)$ is differentially closed; it describes only the extension
$\FE/\KK(z)$.

\begin{corollary}[Finite rational differential systems; source~11]
\label{hol:cf:cor:systems}
Assume $\Gamma$ has no order unit, and let $u_1,\ldots,u_{m'}\in\FE$ satisfy
$D_zu_i=R_i(z,u_1,\ldots,u_{m'})$ with $R_i\in\KK(z)(Y_1,\ldots,Y_{m'})$, each
denominator nonzero after substitution. Then every $u_i$ lies in $\KK(z)$, and
every strongly entire $u_i$ is a polynomial.
\end{corollary}

\begin{proof}
The field $\KK(z)(u_1,\ldots,u_{m'})$ is stable under $D_z$: the equations
put each $D_zu_i$ in the field, and the quotient rule handles every rational
expression in the $u_i$. Its transcendence degree over $\KK(z)$ is at most
$m'$, so the $m'+1$ elements $u_i,D_zu_i,\ldots,D_z^{m'}u_i$ are
algebraically dependent. Apply Theorem~\ref{hol:cf:main:entire}\,(c), and (a)
for the entire ones.
\end{proof}

Corollary~\ref{hol:cor:systems} is the linear case for every $\Gamma\ne0$;
this corollary allows nonlinear rational systems and quotients of entire
series, but needs the absence of an order unit.

\begin{corollary}[Painlev\'e~I in the fraction field; source~11]
\label{hol:cf:cor:painleve}
If $\Gamma$ has no order unit, the equation $D_z^2u=6u^2+z$ has no solution in
$\FE$.
\end{corollary}

\begin{proof}
A solution would be rational by Theorem~\ref{hol:cf:main:entire}\,(c), and
$u=0$ fails. A nonzero rational $u$ has a leading term $cz^d$ at infinity. If
$d\ge1$, then $6u^2$ has degree $2d$, above the degree $1$ of $z$ and the
degree at most $d-2$ of $D_z^2u$. If $d\le0$, then $z$ has degree $1$ while
$u^2$ and $D_z^2u$ have degree at most $0$. No cancellation is possible.
\end{proof}

Proposition~\ref{hol:nl:prop:painleve} already excludes strongly entire
solutions for every $\Gamma\ne0$; the corollary extends this to quotients of
entire series when there is no order unit. It is not a statement about local
formal solutions, which exist for all $a_0,a_1\in\KK$ through
$(n+1)(n+2)a_{n+2}=6\sum_{j=0}^na_ja_{n-j}+\mathbf 1_{n=1}$, nor about
classical complex solutions. The conclusion cannot be strengthened from
rational to constant: for every $c\in\KK$, $u=1/(z-c)\in\FE$ satisfies
$D_zu=-u^2$, and at $c=0$ it has a pole at the origin, which is why the Laurent
form of Theorem~\ref{hol:cf:thm:descent} is needed (source~11).

\begin{corollary}[A solution-dependent exterior certificate; source~12]
\label{hol:cf:cor:DAexterior}
Let $f\in\KK[[z]]$ be a nonpolynomial formal series, differentially algebraic
over $\KK(z)$, and let $L\supseteq k$ be a field finitely generated over $k$
that contains its coefficients, as supplied by
Theorem~\ref{hol:cf:thm:coefficients}\,(a) with $k_0$ generated over $k$ by
the scalar coefficients of a relation. If the smallest convex subgroup
$\Delta$ containing $v(L^\times)$ is proper, every $x$ with $v_\Delta(x)<0$ lies
outside $\Dom(f)$. If $\Gamma$ has no order unit, $\Delta$ is always proper.
\end{corollary}

\begin{proof}
Lemmas~\ref{hol:cf:lem:valuerank} and~\ref{hol:cf:lem:principal} make
$\Delta=\{0\}$ or $\Delta=\Delta_\beta$ with $\beta\in v(L^\times)$; apply
Proposition~\ref{hol:cf:prop:exterior}.
\end{proof}

The certificate is structural, not a uniform algorithm for arbitrary Hahn
names: it can need the values of rational functions of the finitely many
generators of $L$, which their raw supports do not supply. It differs from
Theorem~\ref{hol:nl:thm:certificate}, which needs $\chi_P\ne0$ and finitely
many leading valuations of the solution but no hypothesis on convex
subgroups.

\subsection{The order-unit boundary}\label{hol:cf:sub:sharp}

\begin{proposition}[A one-generator coefficient field; source~12]
\label{hol:cf:prop:theta}
Let $\mu>0$ be an order unit of $\Gamma$ and $p=t^\mu$. The partial theta
series $\Theta_p$ of \eqref{hol:eq:theta} is strongly entire and
nonpolynomial, satisfies $\Theta_p(z)-z\Theta_p(pz)-1=0$, has all its
coefficients in $k(p)$, and $p$ has infinite multiplicative order.
\end{proposition}

\begin{proof}
All coefficients are nonzero and $v(p^{m'})=m'\mu\ne0$ for $m'\ne0$. Given
$\gamma,\beta$, choose positive integers $A,B$ with $\gamma\ge-A\mu$ and
$\beta\le B\mu$; then
$v(p^{n(n-1)/2})+n\gamma\ge\bigl(n(n-1)/2-An\bigr)\mu>\beta$ for large $n$,
and Proposition~\ref{hol:cf:prop:criterion} applies. The functional identity
is \eqref{hol:eq:thetafirst}.
\end{proof}

For $k=\C$ this is also a case of Theorem~\ref{hol:thm:thetadomain}; the
short proof above works over any $k$.

\begin{proof}[Proof of Theorem~\ref{hol:cf:main:sharp} \textup{(source~12)}]
Proposition~\ref{hol:cf:prop:theta} gives (i)$\Rightarrow$(ii) and
(i)$\Rightarrow$(iii): the identity $\Theta_p(z)-z\Theta_p(pz)-1=0$ is the
nonzero polynomial $Y_0-zY_1-1$ vanishing at $(\jet10\Theta_p,\jet p0\Theta_p)$.
If (ii) holds without an order unit, Corollary~\ref{hol:cf:cor:finitefield}
contradicts nonpolynomiality, so (ii)$\Rightarrow$(i). A dependence as in
(iii) involves finitely many powers of $q$ and derivative orders, and
Theorem~\ref{hol:cf:thm:coefficients}\,(b) puts every coefficient in a field
finitely generated over $k$, so (iii)$\Rightarrow$(ii).
\end{proof}

Theorem~\ref{hol:main:q} and Corollary~\ref{hol:cor:detect} detect an order
unit by linear dilation equations with a cofinal parameter;
Theorem~\ref{hol:cf:main:sharp} detects it by coefficient fields and by mixed
jets, with the same partial theta witness. As proved by source~12, neither
said that an order unit produces a nonpolynomial strongly entire
\emph{differentially} algebraic series: the witness's equation contains a
dilation. That was Question~\ref{hol:q:nonlinear} before source~13;
clause~(iv), proved in Section~\ref{hol:td:sub:proofs} from
Theorem~\ref{hol:td:main:boundary}, now says it, with the witness
$\mathcal T_p$.

\paragraph{What ``no order unit'' does and does not mean (sources~11 and~12).}
It is not a synonym of non-Archimedean or higher rank: $\Q\oplus\Q$ ordered
lexicographically, with the first coordinate dominant, has two Archimedean
classes, and $(1,0)$ is an order unit; so has every finite lexicographic
product. An Archimedean group, $\R$ for instance, can have infinite rational
rank and an order unit, and an infinite-rank group has an order unit when it
has a largest Archimedean class. The increasing direct sum
\eqref{hol:eq:gammainfty} has none. The proofs use exactly that every
$\Delta_\beta$ is proper. With an order unit a finitely generated coefficient
field can already have values whose convex hull is all of $\Gamma$; then
Lemma~\ref{hol:cf:lem:principal} forces no proper coarsening, and the argument
stops for that definite reason. The series $G$ of
Corollary~\ref{hol:nl:cor:theta} over $\C((t^\Q))$ is the example of
source~11: strongly entire, with coefficients in $\Q(t)$ whose value group
$\Z$ has convex hull $\Q$. The collection already contains $G$ and its
first-order consequence, and neither source makes a new all-order claim for
it.

\begin{theorem}[An explicit rank-one family; source~12]
\label{hol:cf:thm:rankone}
Let $k=\C$, $\Gamma=\R$, and let $\tau>1$ be a transcendental real number. The
series $R_\tau(z)=\sum_{n\ge0}t^{\tau^n}z^n$ is strongly entire over
$\C((t^\R))$, and for every $\Lambda\subseteq\C((t^\R))^\times$ with pairwise
quotients that are not roots of unity, all $(D_z^jR_\tau)(\lambda z)$ are
algebraically independent over $\C((t^\R))(z)$.
\end{theorem}

\begin{proof}
For real $\gamma$, $\tau^n+n\gamma\to+\infty$, so
Proposition~\ref{hol:cf:prop:criterion} gives strong entireness. A nontrivial
rational relation among distinct $\tau^n$ would be a nonzero polynomial
equation for $\tau$ over $\Q$, so $\cvr(R_\tau)=\infty$, and
Theorem~\ref{hol:cf:main:rank} applies.
\end{proof}

The group $\R$ has an order unit, so Theorem~\ref{hol:cf:main:entire} does not
apply; the theorem uses the infinite rational rank of $\R$ through
Theorem~\ref{hol:cf:main:rank}. It says nothing about a series whose
coefficient values span only $\Q\mu$, such as $\Theta_p$ or $G$, and it does
not settle the purely differential problem over $\C((t^\Q))$. (Source~17's
lacunary series $\mathcal G_A$ have coefficient values spanning $\Q\mu$ and are
differentially transcendental over every order-unit field, including
$\C((t^\Q))$, Theorem~\ref{hol:fh:main:single}; $\Theta_p$ and $G$ remain
undecided.) Through
\eqref{hol:eq:embedding} its surreal realization is
$\sum_{n\ge0}\omega^{-\tau^n}z^n$ on the image of $\C((t^\R))$ in $\No[i]$,
with independence over the full scalar field of that workspace.

\section{Explicit families, surreal realizations and boundaries}
\label{hol:cf:sec:families}

\subsection{The explicit series in all orders}\label{hol:cf:sub:explicit}

\begin{theorem}[Full mixed-jet independence of the explicit series;
source~12]\label{hol:cf:thm:Fomega}
Let $F$ be the strongly entire series \eqref{hol:eq:explicitF} over
$K_{\Gamma_\infty}$. For every $\Lambda\subseteq K_{\Gamma_\infty}^\times$ with
pairwise quotients that are not roots of unity, the family
$\{(D_z^jF)(\lambda z):\lambda\in\Lambda,\ j\ge0\}$ is algebraically independent
over $K_{\Gamma_\infty}(z)$. In particular $F$ satisfies no nonzero algebraic
differential equation of any order over $K_{\Gamma_\infty}(z)$.
\end{theorem}

\begin{proof}
Strong entireness is Theorem~\ref{hol:thm:explicit}; source~12 reproves it by
the highest-coordinate argument recorded there. The values $e_n$ are
rationally independent, so $\cvr(F)=\infty$ and
Theorem~\ref{hol:cf:main:rank} applies; Theorem~\ref{hol:cf:main:entire}
gives the same, since $\Gamma_\infty$ has no order unit.
\end{proof}

Source~11 derives the purely differential part for $f_*=F-\omega^{-1}$ from
Theorem~\ref{hol:cf:main:entire}\,(a); by the remark on indexing in
Section~\ref{hol:cf:sub:conventions} this is the same statement. The theorem
upgrades the linear rigidity of Theorem~\ref{hol:thm:explicit} and the
two-jet independence of Corollary~\ref{hol:nl:cor:infiniterank}; the series
itself is not new. For example, all formal series $(D_z^jF)(q^iz)$, $i\in\Z$, $j\ge0$, for any one of
$q=2$, $q=1+t^{e_0}$ or $q=t^{e_0}$ are jointly algebraically
independent over $K_{\Gamma_\infty}(z)$: the first two parameters have
valuation zero, and the third has a positive valuation that is not an order
unit (source~12).

\subsection{Joint independence and a continuum-sized family}
\label{hol:cf:sub:continuum}

This subsection is source~11's.

\begin{theorem}[Infinite independent coefficient supply; source~11]
\label{hol:cf:thm:indcriterion}
Let $k_0\subseteq K_0$ be fields of characteristic zero and
$f_i=\sum_{n\ge0}a_{i,n}z^n\in K_0[[z]]$ for $i\in I$. For every nonempty
finite $J\subseteq I$ and $i\in J$ put
$k_{J,i}=k_0(a_{i',n}:i'\in J\setminus\{i\},\ n\ge0)$, and suppose
\begin{equation}\label{hol:cf:eq:infcoeff}
 \trdeg_{k_{J,i}}k_{J,i}(a_{i,n}:n\ge0)=\infty
\end{equation}
for all such $J$ and $i$. Then the jets $D_z^jf_i$, $i\in I$, $j\ge0$, are
algebraically independent over $K_0(z)$.
\end{theorem}

\begin{proof}
Otherwise choose an inclusion-minimal nonempty finite $J$ whose members have
algebraically dependent jets, and fix $i\in J$. By minimality the jets of the
other members are algebraically independent over $K_0(z)$. Take a nonzero
relation, clear its denominators in $z$, and let $c_1,\ldots,c_q\in K_0$ be its
scalar coefficients. As a polynomial in the jets of $f_i$ its coefficients
are polynomials in $z$ and in the jets of the other members, and they do not
all vanish after substitution, because that substitution is injective by
the independence just noted. So $f_i$ satisfies a nonzero differential
equation with coefficients in $k_{J,i}(c_1,\ldots,c_q)[[z]]$, and
Theorem~\ref{hol:cf:thm:coefficients}\,(a) puts all $a_{i,n}$ in
$k_{J,i}(c_1,\ldots,c_q,a_{i,0},\ldots,a_{i,N-1})$, of transcendence degree at
most $q+N$ over $k_{J,i}$, contradicting \eqref{hol:cf:eq:infcoeff}.
\end{proof}

The criterion does not require disjoint coefficient fields; finite overlaps
are harmless. It assumes nothing about the Hahn supports of constants in a
relation, which may be arbitrary elements of $K_0$.

\begin{lemma}[Independent scale monomials; source~11]
\label{hol:cf:lem:monomials}
The monomials $t^{e_0},t^{e_1},\ldots$ in $k((t^{\Gamma_\infty}))$ are
algebraically independent over $k$.
\end{lemma}

\begin{proof}
A polynomial in finitely many of them is a finite sum of monomials
$t^{\sum_jm_je_j}$; distinct exponent vectors give distinct elements of
$\Gamma_\infty$, so no cancellation occurs.
\end{proof}

For a binary word $b=(b_1,\ldots,b_n)$ with $n\ge1$ put
\begin{equation}\label{hol:cf:eq:code}
 \code(b)=2^n+\sum_{j=1}^nb_j2^{n-j}\in[2^n,2^{n+1}-1].
\end{equation}
Distinct words have distinct codes, codes strictly increase along an infinite
branch, and two distinct branches share only finitely many prefixes, so their
code sets meet in a finite set. The almost-disjoint prefix device is standard
and also appears, with a different coefficient construction, in the
tail-span report~\cite{repotail}. For $\alpha\in\{0,1\}^{\N_{>0}}$ put
\begin{equation}\label{hol:cf:eq:branches}
 \mathcal F_\alpha(z)=\sum_{n\ge1}t^{e_{\code(\alpha|n)}}z^n,\qquad
 \alpha|n=(\alpha_1,\ldots,\alpha_n).
\end{equation}
Each nonzero coefficient is a Hahn monomial with coefficient $1$ in the
prime field.

\begin{proposition}[Source~11]\label{hol:cf:prop:branchentire}
Every $\mathcal F_\alpha$ is strongly entire over $k((t^{\Gamma_\infty}))$.
\end{proposition}

\begin{proof}
Fix $\gamma,\beta\in\Gamma_\infty$, both supported on
$\{e_0,\ldots,e_{n_\ast}\}$. Once $\code(\alpha|n)>n_\ast$, the largest nonzero coordinate of
$e_{\code(\alpha|n)}+n\gamma-\beta$ is at index $\code(\alpha|n)$ and equals
$1$, so it is positive. Proposition~\ref{hol:cf:prop:criterion} applies.
\end{proof}

\begin{proposition}[Source~11]\label{hol:cf:prop:branchind}
Over $K=k((t^{\Gamma_\infty}))$, all $D_z^j\mathcal F_\alpha$, with
$\alpha\in\{0,1\}^{\N_{>0}}$ and $j\ge0$, are algebraically independent over
$K(z)$.
\end{proposition}

\begin{proof}
Given finitely many distinct branches and one of them, take the largest
common-prefix length with the other branches, or $0$ if there are none.
Beyond that length every code of the chosen branch occurs on none of the
others. So the chosen series has infinitely many coefficient monomials
$t^{e_c}$ indexed by codes absent from all other coefficient fields, and by
Lemma~\ref{hol:cf:lem:monomials} they are algebraically independent over the
field generated by the other branches' monomials. Indeed, a putative
polynomial relation uses only finitely many private monomials and rational
functions of finitely many other monomials. Clearing denominators would
contradict the algebraic independence of that finite combined family.
Condition \eqref{hol:cf:eq:infcoeff} holds with $k_0=k$, and
Theorem~\ref{hol:cf:thm:indcriterion} applies.
\end{proof}

\begin{theorem}[Maximal differential transcendence in one workspace;
source~11]\label{hol:cf:thm:continuum}
Let $k=\C$ and $K_{\Gamma_\infty}=\C((t^{\Gamma_\infty}))$. The family
$(\mathcal F_\alpha)_{\alpha\in\{0,1\}^{\N_{>0}}}$ of strongly entire series has
all its jets jointly algebraically independent over $K_{\Gamma_\infty}(z)$,
and
\begin{equation}\label{hol:cf:eq:dtrdeg}
 \dtrdeg_{K_{\Gamma_\infty}(z)}\operatorname{Frac}\mathcal E_{\Gamma_\infty}
 =2^{\aleph_0}.
\end{equation}
\end{theorem}

\begin{proof}
Propositions~\ref{hol:cf:prop:branchentire} and~\ref{hol:cf:prop:branchind}
give a differentially independent family indexed by the $2^{\aleph_0}$
branches, so the differential transcendence degree is at least $2^{\aleph_0}$.
Conversely $\Gamma_\infty$ is countable and $|\C|=2^{\aleph_0}$, so
$|K_{\Gamma_\infty}|\le|\C^{\Gamma_\infty}|=2^{\aleph_0}$ and
$|K_{\Gamma_\infty}[[z]]|=2^{\aleph_0}$; no differentially independent subset
of the fraction field can be larger.
\end{proof}

The lower bound holds for every $k$; the equality is stated, as by source~11,
for $k=\C$. Equivalently, the free differential polynomial algebra
$K_{\Gamma_\infty}(z)\{X_\alpha\}$ embeds in the fraction field by
$X_\alpha\mapsto\mathcal F_\alpha$. This is cardinal maximality only: the family
is not claimed maximal under inclusion, nor to be a differential
transcendence basis.

The tail-span report already constructs continuum-sized differentially
independent families~\cite{repotail}, using algebraic coefficient extensions
and Galois sign changes over smaller fields, and its analytic construction
lives on a fixed ordinary disk rather than being strongly entire at every Hahn
scale. Here the coefficients are exponent monomials, the functions are
strongly entire on one full Hahn field, and the independence is over that
whole field as field of constants. Neither the cardinal $2^{\aleph_0}$ nor the
almost-disjoint combinatorics is claimed new.

\subsection{Surreal realization, scalar extension, whole-class failure}
\label{hol:cf:sub:surreal}

Through \eqref{hol:eq:embedding}, with $e_j\mapsto\omega^j$, the branch family
becomes
\[
 \mathcal F_\alpha(z)=\sum_{n\ge1}\omega^{-\omega^{\code(\alpha|n)}}z^n,
\]
and at every argument of the embedded workspace its evaluation is an actual
surcomplex number. Everything is set-sized: the value group is countable,
the field has the cardinality of the continuum, and no class-sized sum or
unrestricted class quantification is used (source~11).

\begin{proposition}[Independence survives extension of constants;
source~12]\label{hol:cf:prop:scalar}
Let $E/\KK$ be a field extension and $(g_\alpha)$ a family in $\KK[[z]]$
algebraically independent over $\KK(z)$. Under the coefficientwise inclusion
$\KK[[z]]\subseteq E[[z]]$ the family is algebraically independent over $E(z)$.
\end{proposition}

\begin{proof}
A relation over $E(z)$ among finitely many $g_\alpha$ gives, after clearing
denominators, a nonzero $Q\in E[z,X_1,\ldots,X_{m'}]$ vanishing at them. Choose
a $\KK$-basis $b_1,\ldots,b_d$ of the $\KK$-span of its coefficients and write
$Q=\sum_\nu b_\nu Q_\nu$ with $Q_\nu\in\KK[z,X]$, some $Q_\nu\ne0$. Each
$Q_\nu(z,g)$ has coefficients in $\KK$, so linear independence of the
$b_\nu$ over $\KK$ forces every $Q_\nu(z,g)=0$, and a nonzero $Q_\nu$ is a
relation over $\KK(z)$.
\end{proof}

The proposition concerns fixed formal functions and extensions of the
constant field. It adjoins no new functional variable, says nothing about
evaluation domains, and does not cover new dilation parameters taken from the
extension. Entireness behaves quite differently. Adjoin a dominating
level, $\Gamma'=\Gamma_\infty\oplus\Q e_*$ with $e_*>\Gamma_\infty$, realized by
$e_*=\omega^\omega$; the new group has the order unit $e_*$. At
$x=t^{-e_*}=\omega^{\omega^\omega}$ the evaluation family of $F$ is not
strongly summable (Section~\ref{hol:sub:notallsurcomplex}; source~12 repeats
this computation), and neither is that of any $\mathcal F_\alpha$, whose
leading exponents $e_{\code(\alpha|n)}-ne_*$ strictly decrease since
$e_{\code(\alpha|n+1)}-e_{\code(\alpha|n)}-e_*<0$ (source~11). So the formal
independences survive the larger field by
Proposition~\ref{hol:cf:prop:scalar}, while entireness does not: two different
kinds of stability. Nothing in this part is entire on the proper class
$\No[i]$.

The derivative $D_z$ kills every scalar, including every $\omega^{-\gamma}$,
and $D_z(z)=1$. The Berarducci--Mantova derivation acts on surreal scalars and
is a different operation~\cite{bm}. The tail-span report's all-order
independence results for particular constructions, other base fields and
other derivations~\cite{repotail} are not restated here, and neither kind of
statement is transferred to the other's derivation. The independence
statements concern formal functions: at a point $x\in\KK$ every value $f(x)$
lies in $\KK$ and cannot be transcendental over it, and at $z=0$ the $j$th jet
has the value $j!\,a_j\in\KK$ (sources~11 and~12).

\subsection{Boundary examples}\label{hol:cf:sub:boundaries}

\paragraph{Positive characteristic (sources~11 and~12).}
If $k$ has characteristic $p>0$, the series $\sum_{n\ge0}t^{e_n}z^{pn}$ over
$k((t^{\Gamma_\infty}))$ is strongly entire, by the highest-coordinate
argument and Lemma~\ref{hol:lem:certificate}\,(a), whose proof does not use the
characteristic, yet has $D_z$-derivative zero. This is the phenomenon of
Example~\ref{hol:nl:ex:charp} in the no-order-unit workspace; characteristic
zero is essential even for order-one linear equations, and the finite
coefficient-field principle fails there. No claim about alternative
derivations in positive characteristic is made.

\paragraph{Torsion-related dilations (source~12).}
Let $\zeta\in k$ be a root of unity of order $m'\ge2$ and, over
$k((t^{\Gamma_\infty}))$, let $g(z)=F(z^{m'})=\sum_nt^{e_n}z^{m'n}$. The same highest-coordinate argument
makes $g$ strongly entire, and $\cvr(g)=\infty$, yet $g(\zeta z)-g(z)=0$; for
$\zeta=-1$, an even series relates its dilates by $1$ and $-1$. The
non-root-of-unity condition therefore cannot be dropped from
Theorems~\ref{hol:cf:main:entire}\,(b) and~\ref{hol:cf:main:rank}; compare
Section~\ref{hol:sec:torsion}.

\paragraph{Local evaluation is not entireness (source~12).}
The formal exponential $E_\eta$ of Proposition~\ref{hol:prop:expdomain}, whose
proof works over any $k$, satisfies $D_zE_\eta=t^\eta E_\eta$, has its
coefficients in $k(t^\eta)$, and is strongly evaluable exactly on
$\{0\}\cup\{v(x)>-\eta\}$. In higher rank the leading exponents $nv(y)$ of its
evaluation need not escape cofinally when $v(y)>0$, which is why the
positive-support lemma, not the cofinal criterion, decides that domain
(Remark~\ref{hol:rem:positive}).

\paragraph{Finite coefficient-value rank is not a converse (source~12).}
Theorem~\ref{hol:cf:main:rank} is a sufficient condition. A series with
$\cvr(f)<\infty$ may still be differentially transcendental, and finite
$\cvr(f)$ does not by itself put the coefficients in a finitely generated
field, since residues and tails can carry further algebraic complexity. The
implications proved are
\[
 \text{a mixed relation}\ \Longrightarrow\
 \text{a finitely generated coefficient field}\ \Longrightarrow\
 \cvr(f)<\infty,
\]
not their converses.

An explicit example, added here, is
\[
 f(z)=F(z)+\frac1{1-z}
 =\sum_{n\ge0}(1+t^{e_n})z^n
 \quad\text{over }K_{\Gamma_\infty}.
\]
Every coefficient has value zero, so $\cvr(f)=0$, while subtracting $1$
recovers the algebraically independent monomials of
Lemma~\ref{hol:cf:lem:monomials}; its coefficient field has infinite
transcendence degree. This $f$ is not strongly entire, by
Proposition~\ref{hol:cf:prop:criterion} with $\gamma=\beta=0$.
Nevertheless
\[
 (D_z^jf)(\lambda z)=(D_z^jF)(\lambda z)
       +\frac{j!}{(1-\lambda z)^{j+1}}.
\]
For every finite selection of jets the added terms belong to
$K_{\Gamma_\infty}(z)$, so translation by them is a polynomial-algebra
automorphism. The full mixed independence of
Theorem~\ref{hol:cf:thm:Fomega} therefore holds for $f$ with the same
restriction on the dilation quotients, despite its zero coefficient-value
rank.

\subsection{Consequences for the questions of Section~\ref{hol:sec:status}}
\label{hol:cf:sub:questions}

\paragraph{Question~\ref{hol:q:nonlinear}.}
Theorem~\ref{hol:cf:main:entire}\,(a) (sources~11 and~12) answers it for every
$\Gamma$ without an order unit, in every order: no nonpolynomial strongly
entire series satisfies a nonzero algebraic differential equation over
$\KK(z)$, whatever its residue corner, for every characteristic-zero $k$,
without divisibility, with arbitrary Hahn coefficients in the equation.
Theorem~\ref{hol:cf:main:entire}\,(c) (source~11) extends this to quotients of
entire series. The last sentence of the question as posed by sources~08
and~09 is answered for $F$ with no
coefficient restriction (Theorem~\ref{hol:cf:thm:Fomega}, source~12; source~11
for $f_*$). The question's sentence about hypotheses on $\Gamma$ is answered
in one direction: the absence of an order unit suffices. With an order unit
it stayed open after sources~11 and~12 and was re-scoped to that case;
source~13 has since answered it there (Theorem~\ref{hol:td:main:boundary}),
source~14 has excluded order two for every $\Gamma$
(Theorem~\ref{hol:ot:main:second}), and it is now re-scoped to the
classification of series of minimal order three in
Section~\ref{hol:sec:status}.

\paragraph{Question~\ref{hol:q:mixedunit}.}
Theorem~\ref{hol:cf:main:entire}\,(b) with $\Lambda=\{q^i:i\in\Z\}$
(source~12) excludes, for every nontorsion $q$ including those with $v(q)=0$,
every nonzero algebraic relation among the $\jet{q^i}jf$ of a nonpolynomial
$f\in\EK$ when $\Gamma$ has no order unit. Since
$D_z^r\sigma_q^\ell f=q^{\ell r}\jet{q^\ell}rf$, a nonzero equation
\eqref{hol:eq:mixed} is such a relation, and so are nonlinear mixed
relations. With an order unit the unit-valuation case stayed open after
sources~11 and~12; the witness of Theorem~\ref{hol:cf:main:sharp} uses the
parameter $t^\mu$ of nonzero valuation and says nothing there. Source~18 has
since answered it for linear equations and every $\Gamma$
(Corollary~\ref{hol:ud:cor:unitrigidity}). Theorem~\ref{hol:thm:mixed} and
Theorem~\ref{hol:cf:main:entire}\,(b) do not contain each other: the first
covers every $\Gamma\ne0$ but only nonzero noncofinal $v(q)$, the second every
nontorsion $q$ but only groups without an order unit. Theorem~\ref{hol:ud:thm:rigidity}
contains the first and, for linear relations, the second's cyclic case, and is
contained in neither (Remark~\ref{hol:ud:rem:contains}).

\paragraph{Question~\ref{hol:q:severaldilations}: a merge observation.}
Source~12 does not claim this question. The following consequence of its
theorem was checked for this merge.

\begin{corollary}[Several multiplicatively independent dilations; merge
observation]\label{hol:cf:cor:several}
Suppose $\Gamma$ has no order unit, and let $q_1,\ldots,q_s\in\KK^\times$ be
multiplicatively independent, that is, $q_1^{m_1}\cdots q_s^{m_s}=1$ with
$m_i\in\Z$ only for $m_1=\cdots=m_s=0$. Let $\Lambda$ be the subgroup of
$\KK^\times$ they generate. Then for every nonpolynomial $f\in\EK$ the family
$\{\jet\lambda jf:\lambda\in\Lambda,\ j\ge0\}$ is algebraically independent
over $\KK(z)$. In particular no nonpolynomial strongly entire series satisfies
a nonzero equation, linear or not, in finitely many
$(D_z^jf)(q_1^{m_1}\cdots q_s^{m_s}z)$. For the series $F$ of
\eqref{hol:eq:explicitF}, the only linear relation over
$K_{\Gamma_\infty}(z)$ among $F(\lambda_1z),\ldots,F(\lambda_{s'}z)$ is the
trivial one whenever $\lambda_1,\ldots,\lambda_{s'}\in K_{\Gamma_\infty}^\times$
have pairwise quotients that are not roots of unity, for instance when they
are distinct elements of a group generated by multiplicatively independent
parameters.
\end{corollary}

\begin{proof}
$\Lambda$ is free abelian on the $q_i$, hence torsion free. For distinct
$\lambda,\lambda'\in\Lambda$ the quotient $\lambda/\lambda'$ is a nonidentity
element of $\Lambda$, so no positive power of it is $1$: it is not a root of
unity. Apply Theorem~\ref{hol:cf:main:entire}\,(b), and
Theorem~\ref{hol:cf:thm:Fomega} for the last clause.
\end{proof}

So the first sentence of Question~\ref{hol:q:severaldilations} has, in a
group without an order unit, the answer that no nonpolynomial strongly
entire solution exists, and its last sentence has the answer that only the
trivial relation occurs whenever no quotient of parameters is a root of
unity. The existence criterion in groups with an order unit, and relations
among parameters with a root-of-unity quotient, remained open, and the
question was re-scoped accordingly in Section~\ref{hol:sec:status}; source~18
has since given the existence criterion for some linear equation supported in
a prescribed finite set of multipliers (Theorem~\ref{hol:ud:thm:dichotomy}),
recorded in the status of the question.

\paragraph{Mixed jets in the fraction field.}
The coefficient theorem of source~12 and the fraction-field descent of
source~11 also give the following consequence. Neither source states it;
it is added here by combining their results.

\begin{corollary}[Mixed rigidity in the fraction field]
\label{hol:cf:cor:meromorphicmixed}
Suppose $\Gamma$ has no order unit, and let $u\in\FE\setminus\KK(z)$.
If $\Lambda\subseteq\KK^\times$ has no root-of-unity quotient between
distinct elements, the Laurent series
$\{(D_z^ju)(\lambda z):\lambda\in\Lambda,\ j\ge0\}$ are algebraically
independent over $\KK(z)$.
\end{corollary}

\begin{proof}
Choose $m'\ge0$ so that $y=z^{m'}u\in\KK[[z]]$; also $y\in\FE$.
If the stated family were dependent, a nonzero relation would involve
finitely many parameters and derivative orders. Adjoin the lower orders
through the largest order $r$ used, and for each parameter use Leibniz's
formula, with differentiation before dilation:
\[
 (D_z^ju)(\lambda z)=
 \sum_{\ell=0}^j\binom j\ell\fall{-m'}{j-\ell}
 (\lambda z)^{-m'-j+\ell}(D_z^\ell y)(\lambda z).
\]
For each $\lambda$ this is a triangular change of jet variables over
$\KK(z)$ with nonzero diagonal $(\lambda z)^{-m'}$. The combined change
is therefore an automorphism of the finite polynomial algebra of jet
variables, so it carries the relation to a nonzero relation among the jets
of $y$. Clear its rational denominators and apply
Theorem~\ref{hol:cf:thm:coefficients}\,(b), with $k_0=k$ and $K_0=\KK$.
All coefficients of $y$ lie in a field $L\subseteq\KK$ finitely generated
over $k$. Lemmas~\ref{hol:cf:lem:valuerank} and~\ref{hol:cf:lem:principal}
put $v(L^\times)$ in a proper convex subgroup. Since
$y\in\FE\cap L[[z]]$, the fraction-field part of
Theorem~\ref{hol:cf:thm:descent} gives $y\in L(z)$. Thus
$u=z^{-m'}y\in\KK(z)$, a contradiction.
\end{proof}

The excluded class is the rational functions: a nonpolynomial rational
function still satisfies a polynomial relation over $\KK(z)$. Taking
$\Lambda=\{1\}$ recovers the purely differential statement in
Theorem~\ref{hol:cf:main:entire}\,(c). This consequence uses the same
characteristic-zero and quotient hypotheses as the mixed coefficient
theorem. The no-order-unit hypothesis supplies the proper convex subgroup
needed for descent.

\subsection{Predecessors, repository comparison and priority}
\label{hol:cf:sub:predecessors}

\paragraph{Literature.}
Source~11 consulted the NUMDAM bibliographic page of Hurwitz's
paper~\cite{hurwitz} and its attribution in~\cite{lku}, without a complete
fresh reading of Hurwitz's proof; it read the parsed text of
Laohakosol--Kongsakorn--Ubolsri and viewed its pages 404, 408 and 409, where
Lemma~3 isolates a high coefficient through a polynomial in its index; it
consulted the abstract of Krattenthaler--Rivoal~\cite{kr} as arithmetic
context, not as an imported theorem, and the framework
of~\cite{ast} for context, using none of its software. Source~12 read the
general characteristic-zero form of the Skolem--Mahler--Lech theorem in Bell's
Theorem~1.1~\cite{bell}, imports only that classical theorem and not Bell's
affine-variety generalization, and checked Lech's paper~\cite{lech} through
publisher metadata only, without reconstructing its proof. Both consulted
Hu--Luan~\cite{huluan}, whose setting is an algebraically closed field
of characteristic zero, complete for a nontrivial real-valued non-Archimedean
absolute value; its hypotheses are not transferred to arbitrary-rank Hahn fields, its results are
neither inferred nor replaced, and no rank-one question is claimed resolved.
For normal forms source~11 used Mantova--Matusinski, Section~2.2~\cite{mantovamatusinski},
and source~12 Bournez--Guilmant, Corollary~2.7~\cite{bgexp}, with the same
convention $t^a\mapsto\omega^{-a}$ as \eqref{hol:eq:embedding}.

\paragraph{Repository comparison at the pin.}
Both sources compared their results with the collection at
\texttt{465a54b}, at which this report consisted of
Sections~\ref{hol:sec:intro}--\ref{hol:sec:status} and the conclusion and was
unchanged up to the merge. Their statements about it were checked and are
accurate as of that pin: Question~\ref{hol:q:nonlinear} left order at least two
with vanishing residue corner open, and the README named
$z^2ff''+zff'-z^2(D_zf)^2=0$; Question~\ref{hol:q:mixedunit} is the
unit-valuation mixed case (source~12's phrase ``including unit-valued
parameters'' refers to exactly that); the source~10 audit file still says
``Question~15.1'', as source~11 notes; the report already contained the
sufficiency half of the all-scale criterion, the series $F$ and its domain
warning, first-order independence and the polynomial Painlev\'e exclusions.
The entire-function report~\cite{repoentire} contains the polynomial
intersection obstruction $\mathcal E_\Delta\cap K_\Gamma[[Z]]=K_\Gamma[Z]$ for
noncofinal extensions, coarsened support methods and the distinction between
countable cofinality and order units; the tail-span report~\cite{repotail}
contains continuum-sized families built with binary-prefix almost-disjoint
sets. Since the pin, the statements of this report that called the
vanishing-corner case open have been re-scoped as recorded in
Section~\ref{hol:sec:status}.

Source~11 inspected the root README, the catalogue, directory listings, this
report's README, the source~10 audit, and the READMEs of the entire-function
and tail-span reports; the \path{docs/new} listing then contained only a
README. Source~12 inspected the root README, the catalogue, this report's
README (rereading its lines 131--182) and the tail-span README, and
identified the foundations and entire-function reports through the
catalogue. Neither was a line-by-line audit of the manuscripts, Lean
declarations, archived deliveries, branches or history.

\paragraph{Priority.}
Both literature searches were made on 22 September 2026, were targeted, and
included no exhaustive MathSciNet, zbMATH, Zentralblatt, full-book or
citation-network review; some searches returned irrelevant or inaccessible
material, and no failed search is used as evidence of absence. The inspected
sources did not show the meromorphic descent theorem with its no-order-unit
consequence, the full-constant-field entire branch family, the all-order
no-order-unit and joint-dilation conclusions, the coefficient-value-rank
criterion or the order-unit dichotomy. These are proposed contributions with
proofs, not certified historical firsts, and the answered questions are
questions of this report, not named published conjectures.

\subsection{What the coefficient-field part does not claim}
\label{hol:cf:sub:nonclaims}

Every limitation stated by source~11 or source~12, in its manuscript, README
or audit files, is kept. The tag says which source states it.
\begin{enumerate}[label=\textup{(C\arabic*)}]
\item \emph{[11, 12]} The purely differential case with an order unit is not
settled, and no counterexample is given. Partial theta, $\Theta_p$ or $G$,
is not a counterexample to purely differential rigidity: its equation
contains a dilation, and no new all-order claim is made for it. (Source~13
settles the case with the series $\mathcal T_p$ of
Theorem~\ref{hol:td:main:boundary}; $\Theta_p$ and $G$ remain undecided.)
\item \emph{[11, 12]} The delivered sources supply no Lean formalization of
this theorem package or independent referee report. Source~11 performed no
Lean build of the repository. The finite checks provide no computer-algebra
certification of an infinite statement.
\item \emph{[11, 12]} Historical priority is not certified; the literature
searches were targeted, with the access limits recorded in
Section~\ref{hol:cf:sub:predecessors}, and a failed search is not evidence.
\item \emph{[11, 12]} The repository comparison was targeted, not line by
line; there is no blanket claim that no equivalent result occurs elsewhere
in the collection.
\item \emph{[11, 12]} Not claimed new: classical coefficient recurrences
(Hurwitz, Laohakosol--Kongsakorn--Ubolsri; Krattenthaler--Rivoal as context),
the finite-generation principle, the valuation-rank inequality, Hahn algebra
and Hahn--Neumann support, Conway normal forms, the Skolem--Mahler--Lech
theorem, partial theta series, the all-scale criterion, the polynomial half of
the descent theorem, continuum cardinality and almost-disjoint combinatorics,
and the series $F$ (or $f_*$) with its two-jet consequence.
\item \emph{[11, 12]} \emph{Entire} means strong evaluation on one set-sized
field, not on the class $\No[i]$ and not ordinary convergence of constant
coefficients; entireness fails after an exterior scale is adjoined, with the
explicit witness $x=\omega^{\omega^\omega}$.
\item \emph{[11, 12]} $D_z$ is not the Berarducci--Mantova derivation; nothing
is said about scalar surreal differential equations, and nothing is
transferred to the tail-span report's derivations.
\item \emph{[11, 12]} Independence concerns formal functions, not point
values.
\item \emph{[11]} Containment of $v(L^\times)$ in a proper convex subgroup is
sufficient for descent, not claimed necessary
(Question~\ref{hol:cf:q:descent}); a criterion using rational rank alone will
not suffice.
\item \emph{[11]} Corollary~\ref{hol:cf:cor:relative} does not say that
$\KK(z)$ is differentially closed.
\item \emph{[11]} The continuum family is cardinally maximal only; it is not
maximal under inclusion and not a differential transcendence basis.
\item \emph{[11]} Corollary~\ref{hol:cf:cor:painleve} concerns $\FE$; local
formal solutions exist, and nothing is said about classical complex
solutions.
\item \emph{[11]} ``Meromorphic'' means membership of $\FE$, not a sheaf and
not a fine-topology notion.
\item \emph{[11, 12]} Characteristic zero is essential, and no claim is made
about alternative derivations in positive characteristic.
\item \emph{[11, 12]} There is no uniform algorithm for extracting a convex
subgroup from Hahn names; Corollary~\ref{hol:cf:cor:DAexterior} is structural
and Theorem~\ref{hol:cf:thm:coefficients}\,(b) existential.
\item \emph{[11, 12]} Hu--Luan's rank-one hypotheses are not transferred;
Hu--Luan is context only, and no rank-one question is claimed resolved.
\item \emph{[11, 12]} The finite checks prove no arbitrary-support,
continuum, recurrence zero-set, cofinality or transcendence statement.
\item \emph{[11]} No font files or third-party texts are distributed.
\item \emph{[12]} The answered questions are questions of this research
programme, not named historical or published conjectures.
\item \emph{[12]} The Skolem--Mahler--Lech theorem is imported through Bell's
Theorem~1.1; Bell's affine generalization is not used, Lech's proof was not
reconstructed, and the zero multipliers sampled by the checks are not an
effective bound.
\item \emph{[12]} Theorem~\ref{hol:cf:main:rank} is sufficient only; finite
$\cvr(f)$ does not give a finitely generated coefficient field.
\item \emph{[12]} Proposition~\ref{hol:cf:prop:scalar} covers extension of
constants only, not new dilation parameters from the extension.
\item \emph{[12]} Theorem~\ref{hol:cf:thm:rankone} does not settle the purely
differential problem over $\C((t^\Q))$.
\item \emph{[12]} Dilations with a root-of-unity quotient are excluded
(Question~\ref{hol:cf:q:torsion}).
\item \emph{[12]} In groups of uncountable cofinality there is no
nonpolynomial strongly entire series, so statements requiring one are
vacuous (Corollary~\ref{hol:cf:cor:cofinality}). The coefficient-value-rank
criterion for arbitrary formal series has no such restriction.
\item \emph{[merge]} Corollary~\ref{hol:cf:cor:several} is a consequence
assembled from source~12. Corollary~\ref{hol:cf:cor:meromorphicmixed} is an
additional consequence of sources~11 and~12 together. Neither is attributed
to the delivered sources as a stated result, and no historical priority is
claimed for either deduction.
\end{enumerate}

%======================================================================
% Theta-descent part: source 13. Labels use hol:td:.
%======================================================================
\section{Theta descent: an entire series with an equation of order three}
\label{hol:td:sec:setting}

This section and the next two contain the material of source~13, delivered
after the coefficient-field part had been merged. They answer the order-unit
case of Question~\ref{hol:q:nonlinear} as it stood after that merge: when
$\Gamma$ has an order unit, a nonpolynomial strongly entire series satisfies a
nonzero algebraic differential equation of order three and none of lower order
(Theorem~\ref{hol:td:main:boundary}). With
Theorem~\ref{hol:cf:main:entire}\,(a) this makes the order unit the exact
boundary of differential algebraic rigidity. The ingredients are classical: a
bilateral theta series, Jacobi's triple product and the theta--elliptic
differential identity. What source~13 adds is their descent to an ordinary
power series by the substitution $z=u+u^{-1}$, the support control that makes
the descended series strongly entire exactly when its scale is an order unit,
and a difference-field proof that no equation of order below three exists.
The method is neither the escape chain, nor the finite initial polynomial, nor
the finite coefficient field; the last is used only for the converse
direction, which is Theorem~\ref{hol:cf:main:entire}\,(a).

\subsection{Source, conventions and the merge's choices}
\label{hol:td:sub:conventions}

\begin{convention}[Setting of the theta-descent part]\label{hol:td:conv:field}
In Sections~\ref{hol:td:sec:setting}--\ref{hol:td:sec:boundary}, $k$ is an
arbitrary field of characteristic zero with the trivial valuation, $\Gamma$ is
as in Convention~\ref{hol:conv:group}, $\KK=k((t^\Gamma))$, and $v$, $\lc$,
$\res$, strong summability, $\Dom(f)$, $\EK$, $\FE$ and $D_z$ are as in
Conventions~\ref{hol:nl:conv:field} and~\ref{hol:cf:conv:field}. Throughout,
$\mu\in\Gamma$ is positive, $p=t^\mu$, and $\Delta_\mu$ is the principal convex
subgroup \eqref{hol:eq:delta}; $\mu$ is an order unit only where this is
stated, and then $\Delta_\mu=\Gamma$. The power-series variable is $z$, while
$x$ and $y$ denote elements of $\KK$, and $u$ is the variable of the Laurent
series of Section~\ref{hol:td:sub:laurent}. Write $\vartheta=zD_z$, as in the
proof of Proposition~\ref{hol:nl:prop:degenerate}, and $\vartheta_u=uD_u$. The
elliptic constants below use the nome $p^2$, not $p$.
\end{convention}

For $k=\C$ one has $\KK=\K$ and $\EK=\E$. As for the two previous parts, the
generality in $k$ is recorded for this part only.

\paragraph{The source.}
Source~13 is \emph{Theta Descent and the Exact Boundary of Differential
Rigidity} (21 pages, 23 September 2026). It consulted the repository at
\begin{center}\small\texttt{4cdeaec43ff1692ed2ad9f5bdf761a41872975a5},\end{center}
at which this report was exactly the five-source text into which it is merged
here; its references to Question~19.1 (\texttt{hol:q:nonlinear}), to its
statement and status passage, to Theorem~E and to the coefficient-field part
were accurate there. It assumes characteristic zero, $\Gamma\ne0$ set-sized,
and, where stated, that the scale $\mu$ is an order unit. It uses no
divisibility, no rank hypothesis, no real or algebraic closedness and no
countable cofinal subset. Its standing hypotheses are therefore those of the
coefficient-field part, and every proof taken from it was re-read against the
statement it is attached to.

\paragraph{Renamed symbols.}
The following symbols of source~13 would collide with symbols of the earlier
parts. They are renamed here, in the new material only.
{\small
\begin{longtable}{@{}>{\raggedright\arraybackslash}p{0.27\textwidth}
>{\raggedright\arraybackslash}p{0.2\textwidth}
>{\raggedright\arraybackslash}p{0.45\textwidth}@{}}
\toprule
Source~13 & Here & Reason\\
\midrule
\endhead
variable $x$ & $z$ & Definition~\ref{hol:def:operators}; $x$ denotes
points. (Source~13's $\KK=k((t^\Gamma))$ and $\EK$ are this report's.)\\
$h$, $q=t^h$ & $\mu$, $p=t^\mu$ & $h$ is a torsion order and $q$ a dilation
parameter; $\mu$ and $p=t^\mu$ are the notation of
Theorem~\ref{hol:cf:main:sharp}.\\
$H_h$, $\mathcal D_h$ & $\Delta_\mu$, $\mathcal O_{\Delta_\mu}$ &
\eqref{hol:eq:delta}; $\mathcal D_P$ is a candidate set
\eqref{hol:nl:eq:candidates}; $\mathcal O_\Delta$ as in
Section~\ref{hol:cf:sub:valuation}.\\
$F_q$ & $\mathcal T_p$ & $F$ is the series \eqref{hol:eq:explicitF}.\\
$\Theta_q$ (bilateral) & $\Psi_p$ & $\Theta_p$ is the partial theta series
\eqref{hol:eq:theta}.\\
$r_m$, $P_q$, $j_m$, $e$, $\eta$ & $x_m$, $\mathcal P_p$, $2m-1$, $y$,
$\varpi$ & $r$ indexes derivative orders, $P$ is a differential polynomial,
$e_n$ and $\eta$ are elements of $\Gamma$.\\
$w_\gamma(F_q)$ & $\gamma_\delta$ & The coefficient Gauss value
\eqref{hol:nl:eq:gaussdata}, with $\delta=-\gamma$.\\
$Q=q^2$, $s_r(Q)$, $c$ & $p^2$, $\mathcal Q_r(p^2)$, $\mathfrak c$ & $Q_s$ is
a shift polynomial, $s$ an order, $c_\nu$ a corner coefficient.\\
$A_f$, $H_f$, $J_f$, $B_f$ & $\mathcal A[f]$, $\mathcal H[f]$,
$\mathcal K[f]$, $\mathcal B[f]$ & $A_\delta$ is a leading coefficient,
$H\in\Gamma$, $J$ indexes monomials in \eqref{hol:nl:eq:P}, $B$ is a
threshold, $\jet\lambda j$ a jet.\\
$S_Q(w)$, $W(z)$, $a$, $b$ & $\mathcal Z(w)$, $\mathcal W(y)$,
$\mathfrak a_2$, $\mathfrak a_4$ & $z$ is the power-series variable; $a_n$
are Taylor coefficients.\\
$\delta=u\,\mathrm d/\mathrm d u$, $s=u-u^{-1}$ & $\vartheta_u$, $\check z$ & $\delta$ is
an element of $\Gamma$; $s$ is an order.\\
$\sigma$, $\tau$ & $\mathsf S$, $\mathsf S_a$ & $\sigma_q$ is an argument
dilation; $\tau$ is the parameter of $R_\tau$.\\
$\mathcal A$, $\mathcal M=\operatorname{Frac}\mathcal A$ & $\mathcal U_\mu$,
$\operatorname{Frac}\mathcal U_\mu$ & $\mathcal A[f]$ as above; $M(S)$ is a
support monoid.\\
$L$, $A$, $V$ (in $u$); $Z$ & $\psi_1$, $\psi_2$, $\psi_3$; $V$ & $L$ is a
coefficient field; $z$ as above.\\
$\ell$, $B=A'$, $R(T)$, $\mu=q^{-1}u^{-1}$ & $\mathcal L$, $\mathcal A_1$,
$\Phi(T)$, written out & $\ell$ is a dilation power; $R_\tau$ is a series;
$\mu$ is the scale.\\
$\vartheta=x\,\mathrm d/\mathrm d x$; $\No(i)$ & $\vartheta=zD_z$; $\No[i]$ & The same
operator; not $\vartheta_q$. The collection writes $\No[i]$.\\
Theorem~1.2 & Theorem~\ref{hol:td:main:boundary} & Theorems~A--G of this
report are different theorems.\\
\bottomrule
\end{longtable}}
The integer subscripts of the classical invariants $g_2,g_3$ distinguish them
from the scaled series $g_\delta$ of \eqref{hol:nl:eq:gaussdata}. The shipped
audit files \path{13-theta-descent-PROOF_AUDIT.md} and
\path{13-theta-descent-SOURCES_AND_SCOPE.md} keep the source's notation,
including $F_q$, $\Theta$, $q$, $h$, $L$, $A$ and $\sigma$.

\paragraph{Words with two meanings.}
The \emph{order} of an algebraic differential equation is the highest
derivative occurring in it; an \emph{order unit} is as in
Definition~\ref{hol:def:orderunit}; the \emph{multiplicative order} of $q$ is
as in Theorem~\ref{hol:main:q}; and $\operatorname{ord}_u$ is the $u$-adic
order of a rational function. The \emph{minimal order} of a series is the
least order of a nonzero algebraic differential equation over $\KK(z)$ that it
satisfies. \emph{Descent} in this part means the Chebyshev substitution
$z=u+u^{-1}$, not the coarsening descent of
Theorem~\ref{hol:cf:thm:descent}. Three theta series occur and are not
interchangeable: the partial series $\Theta_p$ of \eqref{hol:eq:theta}, which
satisfies a dilation equation; the bilateral series $\Psi_p$ of
\eqref{hol:td:eq:psi}; and the Hahn--Tate report's series
$\sum_{n\in\Z}(-1)^nq^{n(n-1)/2}u^n$~\cite{repotate}, of which $\Psi_p(u)$ is
the value at the nome $p^2$ and the argument $-pu$. A \emph{corner} keeps the
meaning of Section~\ref{hol:nl:sub:corner}: source~13's ``vanishing corner''
is the vanishing of the residue corner polynomial $\chi_P$, and the full
corner polynomial $I_P$ need not vanish with it
(Proposition~\ref{hol:td:prop:corner}). A \emph{zero} of $\mathcal T_p$ is a
point $x\in\Dom(\mathcal T_p)$ with $\mathcal T_p(x)=0$.

\paragraph{Results printed once.}
Source~13 reproves several facts that the report already contains. They are
printed once and cited: the definitions of an order unit and of strong
summability (Definitions~\ref{hol:def:orderunit} and~\ref{hol:def:entire});
its cofinal-valuation lemma, which is Lemma~\ref{hol:lem:certificate}\,(a);
the positive-support lemma, which is Lemma~\ref{hol:lem:support}\,(2); the
surreal normal-form embedding \eqref{hol:eq:embedding}; the strong domain of
the bilateral theta series, which is the Hahn--Tate report's theta-domain
theorem~\cite[\texttt{tate:thm:theta}]{repotate} after the substitution just
described; and its whole Section~7, the no-order-unit direction, which is
Theorem~\ref{hol:cf:main:entire}\,(a) with
Theorem~\ref{hol:cf:thm:coefficients}\,(a),
Lemmas~\ref{hol:cf:lem:valuerank} and~\ref{hol:cf:lem:principal} and
Proposition~\ref{hol:cf:prop:exterior}; its Euler-derivative presentation of
that argument is recorded in Remark~\ref{hol:td:rem:euler}. For $k=\C$ the
non-D-finiteness of $\mathcal T_p$ is also a case of
Theorem~\ref{hol:main:ode}; source~13's separate proof, valid for every $k$,
is kept (Proposition~\ref{hol:td:prop:nondfinite}). Everything else in
source~13 is printed in
Sections~\ref{hol:td:sec:setting}--\ref{hol:td:sec:boundary}, with its
proofs.

\paragraph{Where each numbered result went.}
Theorem~1.2 is Theorem~\ref{hol:td:main:boundary};
Lemma~2.3 is Lemma~\ref{hol:td:lem:quadratic}; Lemma~3.1,
Proposition~3.2 and Theorem~3.3 are Lemma~\ref{hol:td:lem:cheb},
Proposition~\ref{hol:td:prop:coeff} and Theorem~\ref{hol:td:thm:domain};
Theorem~4.1 and Proposition~4.2 are Theorem~\ref{hol:td:thm:zeros} and
Proposition~\ref{hol:td:prop:local}; Theorem~5.1 is
Theorem~\ref{hol:td:thm:ode}, and Remark~5.2 is
Remark~\ref{hol:td:rem:corner}; Lemma~6.1, Theorem~6.2, Corollary~6.3 and
Theorem~6.4 are Lemma~\ref{hol:td:lem:rational},
Theorem~\ref{hol:td:thm:independent}, Corollary~\ref{hol:td:cor:order} and
Theorem~\ref{hol:td:thm:field}, and Section~6.5 is
Proposition~\ref{hol:td:prop:nondfinite}; Lemmas~7.1--7.2 and Theorem~7.3 are
printed once as above; Corollary~8.1 is Corollary~\ref{hol:td:cor:omnific};
Examples~9.1--9.3 and Corollary~9.4 are
Examples~\ref{hol:td:ex:discrete}--\ref{hol:td:ex:cofinality} and
Corollary~\ref{hol:td:cor:affine}; Questions~11.1 and~11.2 are merged into the
re-scoped Question~\ref{hol:q:nonlinear}, and Question~11.3 is
Question~\ref{hol:td:q:domains}; Appendix~A is
\eqref{hol:td:eq:jetpoly} and the expansions after
\eqref{hol:td:eq:constants}; Appendix~B is merged into the dependency table of
Section~\ref{hol:sec:verification}.

\paragraph{Two steps written out, and other choices.}
The review of source~13 for this merge found its main line correct and two
steps used but not written. Both are supplied here, marked as the merge's.
First, the minimal-order proof substitutes $z=u+u^{-1}$ into strongly entire
series and uses that this is a ring homomorphism into a Laurent ring,
compatible with the derivations; Lemma~\ref{hol:td:lem:substitution} states
and proves this. Second, source~13 proved $\vartheta_u\psi_2\ne0$ from
``formal meromorphic germs'' at the simple zero $u_0=-p$ of $\Psi_p$, a local
expansion inside $\operatorname{Frac}\mathcal U_\mu$ that it did not define;
Lemma~\ref{hol:td:lem:psi3} replaces this by the global identity
$\Psi_p^3\,\vartheta_u\psi_2=p(u-u^{-1})+O(p^2)$. The passage from the complex
elliptic identity to $\KK$ is written as three ring maps
(Section~\ref{hol:td:sub:formal}), which is the content of source~13's
``support-controlled specialization''. The merge also computes the corner data
of the cleared equation (Proposition~\ref{hol:td:prop:corner}), adds
clause~(iv) to Theorem~\ref{hol:cf:main:sharp}, and records the consequences
for Theorem~\ref{hol:cf:main:entire} (Corollary~\ref{hol:td:cor:sharp}).
Source~13's theorems, proofs and non-claims are otherwise printed as
delivered, in the renamed notation.

\subsection{Chebyshev descent and the exact Taylor coefficients}
\label{hol:td:sub:chebyshev}

The monic Chebyshev polynomials $C_n\in\Z[z]$ are defined by
\begin{equation}\label{hol:td:eq:cheb}
 C_0=2,\qquad C_1=z,\qquad C_{n+1}=zC_n-C_{n-1}\quad(n\ge1),
\end{equation}
and $\mathcal T_p$ is the series \eqref{hol:td:eq:series}. In ordinary
notation $C_n(z)=2T_n(z/2)$; the normalization $C_0=2$ is why the constant
term of $\mathcal T_p$ is written separately.

\begin{lemma}[Chebyshev identities; source~13]\label{hol:td:lem:cheb}
For $n\ge1$, $C_n$ is monic of degree $n$ and
\begin{align}
 C_n(u+u^{-1})&=u^n+u^{-n},\label{hol:td:eq:chebsub}\\
 C_n(z)&=\sum_{j=0}^{\lfloor n/2\rfloor}
 (-1)^j\frac{n}{n-j}\binom{n-j}{j}z^{n-2j}.
 \label{hol:td:eq:chebcoef}
\end{align}
All the coefficients in the second formula are integers.
\end{lemma}

\begin{proof}
The first identity follows from the initial values and the same three-term
recurrence on both sides. The recurrence also proves monicity and
integrality. Substituting the displayed coefficients into the recurrence and
using Pascal's identity proves \eqref{hol:td:eq:chebcoef}; the initial cases
are $C_1=z$ and $C_2=z^2-2$.
\end{proof}

Thus
\[
 \mathcal T_p(z)=1+pz+p^4(z^2-2)+p^9(z^3-3z)+p^{16}(z^4-4z^2+2)
 +p^{25}(z^5-5z^3+5z)+\cdots.
\]

\begin{proposition}[Exact Taylor coefficients; source~13]
\label{hol:td:prop:coeff}
For every $\mu>0$ write $\mathcal T_p=\sum_{m\ge0}a_mz^m$. Then
\begin{align}
 a_0&=1+2\sum_{j\ge1}(-1)^jp^{4j^2},\label{hol:td:eq:a0}\\
 a_m&=\sum_{j\ge0}(-1)^j\frac{m+2j}{m+j}\binom{m+j}{j}p^{(m+2j)^2}
 \qquad(m\ge1),\label{hol:td:eq:am}
\end{align}
and
\begin{equation}\label{hol:td:eq:coeffval}
 v(a_0)=0,\qquad v(a_m)=m^2\mu,\quad\lc(a_m)=1\qquad(m\ge1).
\end{equation}
No coefficient vanishes, so $\mathcal T_p$ is not a polynomial.
\end{proposition}

\begin{proof}
For each $m$ collect the terms $n=m+2j$ of \eqref{hol:td:eq:chebcoef}. Their
exponents $(m+2j)^2\mu$ increase strictly, so the family is strongly summable
for every $\mu>0$ (Lemma~\ref{hol:lem:certificate}\,(b) with $S=\{\mu\}$) and
each exponent carries one term. For $m\ge1$ the first coefficient is one and
every later term has strictly larger valuation. The constant term follows from
$C_{2j}(0)=2(-1)^j$ and $C_{2j+1}(0)=0$.
\end{proof}

\subsection{The exact strong domain}\label{hol:td:sub:domain}

\begin{lemma}[A quadratic support certificate; source~13]
\label{hol:td:lem:quadratic}
Let $\mu>0$ and $x\in\KK$ with $x=0$ or $v(x)\ge-N\mu$ for some $N\in\N$. The
family
\[
 \bigl(p^{n^2}x^{n-2j}\bigr)_{n\ge1,\ 0\le j\le\lfloor n/2\rfloor}
\]
is strongly summable, and so is every family obtained from it by multiplying
members by ordinary integers or omitting members.
\end{lemma}

\begin{proof}
For $x=0$ only the members with $n-2j=0$ are nonzero, and their exponents
$n^2\mu$ are distinct. Otherwise write
\[
 x=a\,t^{-N\mu+\beta}(1+\varepsilon_x),\qquad a\in k^\times,\quad\beta\ge0,
 \quad\varepsilon_x=0\text{ or }v(\varepsilon_x)>0.
\]
The support of a member lies in
$[n^2-N(n-2j)]\mu+(n-2j)\beta+M(\supp\varepsilon_x)$. For $n>N$ the integer
$c_{n,j}=n^2-N(n-2j)\ge n(n-N)$ is positive and tends to infinity with $n$,
uniformly in $j$. Apply Lemma~\ref{hol:lem:support}\,(2) to the well-ordered
positive alphabet
$S=\{\mu\}\cup(\{\beta\}\setminus\{0\})\cup\supp\varepsilon_x$. After the
binomial expansion of $(1+\varepsilon_x)^{n-2j}$, every support element of a
member with $n>N$ is represented by a word with at least $c_{n,j}$ letters
$\mu$. A fixed exponent has only finitely many words, so the number of their
letters $\mu$ is bounded, which bounds $n$ and hence $j$; each of the finitely
many remaining members is one Hahn element. The members with $n\le N$ are
finitely many. The same argument shows that the union of the supports is well
ordered. Multiplication by ordinary integers and omission enlarge neither
supports nor multiplicities. No cofinality of $(n\mu)$ is assumed.
\end{proof}

\begin{theorem}[Exact strong domain; source~13]\label{hol:td:thm:domain}
For every $\mu>0$, not necessarily an order unit,
\begin{equation}\label{hol:td:eq:domain}
 \Dom(\mathcal T_p)=\mathcal O_{\Delta_\mu}
 =\{0\}\cup\{x\in\KK^\times:v(x)\ge-N\mu\text{ for some }N\in\N\},
\end{equation}
the valuation ring of the coarsened valuation $v_{\Delta_\mu}$. Consequently
$\mathcal T_p\in\EK$ if and only if $\mu$ is an order unit of $\Gamma$.
\end{theorem}

\begin{proof}
For $x\in\mathcal O_{\Delta_\mu}$, Lemma~\ref{hol:td:lem:quadratic} applies to
the full expansion $(p^{n^2}c_{n,j}x^{n-2j})_{n,j}$, with $c_{n,j}$ the
integers of \eqref{hol:td:eq:chebcoef}. Regrouping it by powers of $x$ gives
exactly $(a_mx^m)_{m\ge0}$. Here and below we use the standard associativity
of strong sums: if a strongly summable family is partitioned into arbitrary
subfamilies, each subfamily is strongly summable, and the family of their sums
is strongly summable with the same sum, because the supports of the partial
sums lie in the union of the original supports and only finitely many
subfamilies contain a member contributing at a fixed exponent.

If $x\notin\mathcal O_{\Delta_\mu}$, put $\gamma=v(x)$, so $\gamma<-N\mu$ for
every $N$. By \eqref{hol:td:eq:coeffval} the members with $m\ge1$ have leading
exponents $m^2\mu+m\gamma$, whose consecutive differences $(2m+1)\mu+\gamma$
are negative. They strictly decrease, and
Lemma~\ref{hol:lem:nonincreasing} excludes strong summability.

A coset in $\Gamma/\Delta_\mu$ is nonnegative exactly when its representative
lies in $\Delta_\mu$ or is positive and outside it, that is, when it is
bounded below by some $-N\mu$; this identifies the valuation ring. It equals
$\KK$ exactly when every negative value is so bounded, that is, when
$\Delta_\mu=\Gamma$.
\end{proof}

When $\mu$ is an order unit there is the shorter route: for fixed $\gamma$ the
values $m^2\mu+m\gamma$ escape cofinally, and
Lemma~\ref{hol:lem:certificate}\,(a) applies. The support proof is needed for
the noncofinal boundary, where strong summability and convergence in the full
valuation topology differ (Remark~\ref{hol:rem:positive}). When $\mu$ is not an
order unit, the coefficient values of $\mathcal T_p$ lie in the proper convex
subgroup $\Delta_\mu$, and Theorem~\ref{hol:td:thm:domain} shows that the
necessary inclusion of Proposition~\ref{hol:cf:prop:exterior} is then an
equality for this series.

\subsection{The bilateral series, the product and the zeros}
\label{hol:td:sub:bilateral}

Put
\begin{equation}\label{hol:td:eq:psi}
 \Psi_p(u)=\sum_{n\in\Z}p^{n^2}u^n .
\end{equation}
At $u\in\KK^\times$ this family is strongly summable exactly when
$v(u)\in\Delta_\mu$: both tails are covered by the argument of
Lemma~\ref{hol:td:lem:quadratic} when $\abs{v(u)}\le N\mu$, and otherwise one
of the sequences $n^2\mu\pm nv(u)$ strictly decreases. Since
$\Psi_p(u)=\sum_n(-1)^n(p^2)^{n(n-1)/2}(-pu)^n$, this is the Hahn--Tate theta
domain~\cite[\texttt{tate:thm:theta}]{repotate} at the nome $p^2$ and the
argument $-pu$, printed once; source~13 describes it as the symmetric
specialization of that domain. On this domain, and as identities of formal
series in Section~\ref{hol:td:sub:formal}, reindexing gives
\begin{align}
 \Psi_p(u^{-1})&=\Psi_p(u),\label{hol:td:eq:inverse}\\
 \Psi_p(p^2u)&=p^{-1}u^{-1}\Psi_p(u),\label{hol:td:eq:shift}\\
 \mathcal T_p(u+u^{-1})&=\Psi_p(u).\label{hol:td:eq:descent}
\end{align}
For \eqref{hol:td:eq:shift} the exponent identity is $n^2+2n=(n+1)^2-1$. For
\eqref{hol:td:eq:descent}, $x=u+u^{-1}$ has $v(x)\ge-\abs{v(u)}$, so
Lemma~\ref{hol:td:lem:quadratic} applies at $x$; grouping the full expansion
by $n$ gives $(p^{n^2}C_n(x))_n=(p^{n^2}(u^n+u^{-n}))_n$ by
\eqref{hol:td:eq:chebsub}, and grouping it by powers of $x$ gives the Taylor
family. No square root or inverse cosine is used to define $\mathcal T_p$, and
no branch is chosen.

Jacobi's triple product, in the normalization of DLMF~20.5.9~\cite{dlmftheta},
is
\[
 \sum_{n\in\Z}q^{n^2}u^n
 =\prod_{m\ge1}(1-q^{2m})(1+q^{2m-1}u)(1+q^{2m-1}u^{-1}).
\]
It is classical, not a new theorem about theta functions. Pairing the
reciprocal factors gives the following form.

\begin{theorem}[Product and complete zero set; source~13]
\label{hol:td:thm:zeros}
Let $\mu>0$ and $\mathcal P_p=\prod_{m\ge1}(1-p^{2m})$. For every
$x\in\mathcal O_{\Delta_\mu}$,
\begin{equation}\label{hol:td:eq:product}
 \mathcal T_p(x)=\mathcal P_p\prod_{m\ge1}\bigl(1+p^{2m-1}x+p^{4m-2}\bigr),
\end{equation}
where the products are strong products, that is, strong sums of their
finite-subset expansions. The scalar $\mathcal P_p$ is a unit of valuation
zero and leading coefficient one. The zeros of $\mathcal T_p$ in its domain are
exactly
\begin{equation}\label{hol:td:eq:roots}
 x_m=-\bigl(p^{2m-1}+p^{-(2m-1)}\bigr)\qquad(m=1,2,\ldots),
\end{equation}
and every zero is simple. The conclusion is already split over
$k((t^\Gamma))$; no algebraic closedness is used.
\end{theorem}

\begin{proof}
First let $p$ be a formal variable. At each $p$-degree both sides of the triple
product have Laurent-polynomial coefficients in $u$, and the identity holds in
$\Q[u,u^{-1}][[p]]$. Since
$(1+p^{2m-1}u)(1+p^{2m-1}u^{-1})=1+p^{2m-1}(u+u^{-1})+p^{4m-2}$ and the
substitution $z\mapsto u+u^{-1}$ is injective on $\Q[z]$ (the leading power
$u^d$ of a polynomial of degree $d$ cannot cancel), the paired identity
$1+\sum_{n\ge1}p^{n^2}C_n(z)=\mathcal P_p\prod_m(1+p^{2m-1}z+p^{4m-2})$ holds
in $\Q[z][[p]]$.

Now let $p=t^\mu$ and $v(x)\ge-N\mu$. For $2m-1>N$ both $p^{2m-1}x$ and
$p^{4m-2}$ have positive valuation, and their supports, with the factor labels
retained, lie in a positive support monoid generated by $\mu$ and the support
of the normalized $x$ as in Lemma~\ref{hol:td:lem:quadratic}; every nonempty
choice from a late factor uses at least one letter $\mu$, so a fixed exponent
bounds the number and the indices of the late factors involved.
Lemma~\ref{hol:lem:support}\,(2) therefore makes the finite-subset expansion of
the tail product strongly summable, and the finitely many early factors
multiply it by a finite family. The same holds for $\mathcal P_p$. Expand both
sides of the paired identity at $x$ into families of terms $e\,p^kx^j$ with
$e\in\Z$: the Taylor side is strongly summable by
Lemma~\ref{hol:td:lem:quadratic}, the product side by the argument just given.
Grouping either family by the formal monomial $p^kz^j$ gives the same family,
by the identity in $\Q[z][[p]]$, so the two sums agree. The tail product lies
in $1+\mathfrak m_v$ and is nonzero.

Each factor of \eqref{hol:td:eq:product} equals $p^{2m-1}(x-x_m)$, and the
roots have the distinct valuations $-(2m-1)\mu$. At a fixed point all
sufficiently late factors form a nonzero unit product, so a zero comes from
one of finitely many early factors; there are no other zeros. For simplicity,
split off the $m$th factor $\Pi_m(z)=1+p^{2m-1}z+p^{4m-2}$: the formal identity
gives $\mathcal T_p=\Pi_m\cdot G_m$ in $\KK[[z]]$, with $G_m$ the image of the
complementary product, which is strongly evaluable on $\mathcal
O_{\Delta_\mu}$ by the same argument. Differentiating and evaluating at the
nonzero point $x_m$, where $\Pi_m$ vanishes, gives
$(D_z\mathcal T_p)(x_m)=p^{2m-1}G_m(x_m)$, a nonzero strong product.
\end{proof}

For $k=\R$ every zero is real and negative; for $k=\C$ there are no further
nonreal zeros. These conclusions come from the displayed factorization, not
from a claim that non-Archimedean entire functions split in general.

\paragraph{Valuations at the zeros.}
Write $\mathcal T_p'=D_z\mathcal T_p$. At $x=x_m$ the factors with $k<m$ have
valuation $-2(m-k)\mu$ and leading coefficient $-1$, and those with $k>m$ have
valuation zero and leading coefficient one. Hence
\begin{equation}\label{hol:td:eq:rootderivative}
 v\bigl(\mathcal T_p'(x_m)\bigr)=(-m^2+3m-1)\mu,\qquad
 \lc\bigl(\mathcal T_p'(x_m)\bigr)=(-1)^{m-1},
\end{equation}
which is also a direct certificate that the derivative does not vanish.

\begin{proposition}[Exact linearization on the nearest-root ball; source~13]
\label{hol:td:prop:local}
If $y\ne0$ and $v(y)>-(2m-1)\mu$, then
\begin{equation}\label{hol:td:eq:linearization}
 \mathcal T_p(x_m+y)=\mathcal T_p'(x_m)\,y\,(1+\varpi),\qquad
 \varpi=0\text{ or }v(\varpi)>0.
\end{equation}
Therefore $v(\mathcal T_p(x_m+y))=(-m^2+3m-1)\mu+v(y)$.
\end{proposition}

\begin{proof}
For $k\ne m$ the valuation of $x_m-x_k$ is the smaller of the two root
valuations and is at most $-(2m-1)\mu$, so $v(y/(x_m-x_k))>0$. Dividing
\eqref{hol:td:eq:product} at $x_m+y$ by $\mathcal T_p'(x_m)\,y$ gives
\[
 \frac{\mathcal T_p(x_m+y)}{\mathcal T_p'(x_m)\,y}
 =\prod_{k\ne m}\Bigl(1+\frac y{x_m-x_k}\Bigr).
\]
The finitely many early factors lie in $1+\mathfrak m_v$. In the tail,
$1/(x_m-x_k)$ has valuation $(2k-1)\mu$; expanding after normalizing its
leading monomial gives the same positive support certificate as for the
original product. The strong product therefore lies in $1+\mathfrak m_v$.
\end{proof}

\paragraph{Newton profile.}
In the notation \eqref{hol:nl:eq:gaussdata}, at a scale $\delta$ with
$t^{-\delta}\in\Dom(\mathcal T_p)$, that is $\delta\le N\mu$ for some $N$, the
coefficient Gauss value is
$\gamma_\delta=\min_{n\ge0}(n^2\mu-n\delta)$. The minimum exists because the
members with $n>N$ exceed the member $n=0$, and consecutive differences are
$(2n+1)\mu-\delta$. There is one active index $n$ when
$(2n-1)\mu<\delta<(2n+1)\mu$ (with the evident change for $n=0$), and two, $n$
and $n+1$, when $\delta=(2n+1)\mu$; the root $x_{n+1}$ has exactly the valuation
$-(2n+1)\mu$. The quadratic Newton profile and the linear-factor product give
the same complete shell classification, without a general Newton-polygon
theorem.

\section{An equation of order three, and none of lower order}
\label{hol:td:sec:order}

\subsection{The cleared third-order equation}\label{hol:td:sub:ode}

For $r\in\{1,3,5\}$ put
\begin{equation}\label{hol:td:eq:lambert}
 \mathcal Q_r(p^2)=\sum_{n\ge1}\frac{n^rp^{2n}}{1-p^{2n}}
 =\sum_{N\ge1}\Bigl(\sum_{d\mid N}d^r\Bigr)p^{2N},
\end{equation}
and
\begin{equation}\label{hol:td:eq:constants}
 \mathfrak c=\frac1{12}-2\mathcal Q_1(p^2),\qquad
 g_2=\frac1{12}+20\mathcal Q_3(p^2),\qquad
 g_3=-\frac1{216}+\frac73\mathcal Q_5(p^2).
\end{equation}
The powers are supported on positive multiples of $2\mu$, and each exponent
has finitely many divisors, so the series exist for every $\mu>0$. They begin
\begin{align*}
 \mathfrak c&=\tfrac1{12}-2p^2-6p^4-8p^6-14p^8+O(p^{10}),\\
 g_2&=\tfrac1{12}+20p^2+180p^4+560p^6+1460p^8+O(p^{10}),\\
 g_3&=-\tfrac1{216}+\tfrac73p^2+77p^4+\tfrac{1708}3p^6+\tfrac{7399}3p^8
 +O(p^{10}),
\end{align*}
where $O(p^{10})$ denotes a formal series divisible by $p^{10}$, not a
numerical error bound. Put
\begin{equation}\label{hol:td:eq:Phi}
 \Phi(T)=4(\mathfrak c-T)^3-g_2(\mathfrak c-T)-g_3 .
\end{equation}

For a nonzero element $f$ of a differential field extending $\KK(z)$ with
derivation $D_z$, put
\begin{align}
 \mathcal A[f]&=(z^2-4)\Bigl(\frac{D_z^2f}f-\frac{(D_zf)^2}{f^2}\Bigr)
 +z\frac{D_zf}f,\label{hol:td:eq:A}\\
 \mathcal H[f]&=(z^2-4)\bigl(fD_z^2f-(D_zf)^2\bigr)+zfD_zf,
 \label{hol:td:eq:H}\\
 \mathcal K[f]&=fD_z\mathcal H[f]-2D_zf\,\mathcal H[f],\qquad
 \mathcal B[f]=\mathfrak cf^2-\mathcal H[f],\label{hol:td:eq:KB}
\end{align}
so that $\mathcal A[f]=\mathcal H[f]/f^2$ and
$D_z\mathcal A[f]=\mathcal K[f]/f^3$. These are rational differential
expressions; no logarithm of a field element is used.

\begin{theorem}[Cleared third-order equation; source~13]\label{hol:td:thm:ode}
For every $\mu>0$ the series $f=\mathcal T_p$ satisfies, in $\KK[[z]]$,
\begin{equation}\label{hol:td:eq:cleared}
 (z^2-4)\mathcal K[f]^2-4\mathcal B[f]^3+g_2\mathcal B[f]f^4+g_3f^6=0 .
\end{equation}
Its left side is $P(z,f,D_zf,D_z^2f,D_z^3f)$ for the nonzero polynomial
\begin{equation}\label{hol:td:eq:jetpoly}
 P=(z^2-4)\mathrm K^2-4\mathrm B^3+g_2\mathrm BY_0^4+g_3Y_0^6\in\KK[z,Y_0,Y_1,Y_2,Y_3],
\end{equation}
where
$\mathrm H=(z^2-4)(Y_0Y_2-Y_1^2)+zY_0Y_1$,
$\mathrm H_1=(z^2-4)(Y_0Y_3-Y_1Y_2)+2z(Y_0Y_2-Y_1^2)+Y_0Y_1+z(Y_1^2+Y_0Y_2)$
(the total derivative of $\mathrm H$),
$\mathrm K=Y_0\mathrm H_1-2Y_1\mathrm H$ and $\mathrm B=\mathfrak cY_0^2-\mathrm
H$. The polynomial $P$ is homogeneous of degree six in $Y_0,\ldots,Y_3$, and
its coefficient of $Y_3^2$ is $(z^2-4)^3Y_0^4$, so no values of $\mathfrak c$,
$g_2$, $g_3$ make it zero. In $\KK((z))$ the equivalent rational form is
\begin{equation}\label{hol:td:eq:rationalODE}
 (z^2-4)\bigl(D_z\mathcal A[f]\bigr)^2=\Phi\bigl(\mathcal A[f]\bigr).
\end{equation}
Both are formal identities. The cleared one can be evaluated at every point of
$\Dom(\mathcal T_p)$, including the zeros of $f$ and $z=\pm2$, since
evaluation is a ring homomorphism on series evaluable at a common point and no
division occurs.
\end{theorem}

The proof occupies the next two subsections. The coefficient of $D_z^3f$ in
$\mathcal K[f]$ is $(z^2-4)f^2$, which gives the leading coefficient in $Y_3$.

\subsection{The classical elliptic identity}\label{hol:td:sub:elliptic}

Source~13 supplies the normalization calculation because a sign or nome
convention changes all three constants. Theta--elliptic differential
identities are classical~\cite{dlmfelliptic,sebbar}, and Tate coordinates give
another route to the same normalization~\cite{woodbury,repotate}.

In this subsection only, $p$ is an ordinary real number with $0<p<1$, and
$\Psi(u)=\sum_{n\in\Z}p^{n^2}u^n$ for $u\in\C^\times$. Let
$\psi_1=\vartheta_u\Psi/\Psi$ and $\psi_2=\vartheta_u\psi_1$. Differentiating
the product away from its zeros gives
\begin{equation}\label{hol:td:eq:S}
 \psi_2=-\mathcal Z(-pu),\qquad
 \mathcal Z(w)=\sum_{n\in\Z}\frac{p^{2n}w}{(1-p^{2n}w)^2}.
\end{equation}
For negative $n$ the summand can be rewritten through
$w/(1-w)^2=w^{-1}/(1-w^{-1})^2$, so this is a normally convergent meromorphic
series on $\C^\times$, invariant under $w\mapsto p^2w$ and $w\mapsto w^{-1}$.
Put
\[
 \mathcal W(y)=\mathcal Z(e^y)-2\mathcal Q_1(p^2)+\tfrac1{12}.
\]
It is meromorphic with periods $2\pi i$ and $\log p^2$, and its poles are
exactly the points of the period lattice, all double. Near $y=0$ the summand
$n=0$ is
\[
 \frac{e^y}{(1-e^y)^2}=y^{-2}-\frac1{12}+\frac{y^2}{240}-\frac{y^4}{6048}+O(y^6),
\]
and pairing the summands $n$ and $-n$ for $n\ge1$ gives
$2\sum_{d\ge1}dp^{2nd}\cosh(dy)$. Consequently
\begin{equation}\label{hol:td:eq:Wlaurent}
 \mathcal W(y)=y^{-2}+\mathfrak a_2y^2+\mathfrak a_4y^4+O(y^6),\qquad
 \mathfrak a_2=\tfrac1{240}+\mathcal Q_3(p^2),\quad
 \mathfrak a_4=-\tfrac1{6048}+\tfrac1{12}\mathcal Q_5(p^2).
\end{equation}
The elliptic function $(\mathcal W')^2-4\mathcal W^3+20\mathfrak a_2\mathcal
W+28\mathfrak a_4$ has no pole: its $y^{-6}$, $y^{-2}$ and constant terms
cancel by \eqref{hol:td:eq:Wlaurent}, and periodicity treats every lattice
point. It is bounded on a compact period parallelogram, hence constant, and
its limit at $0$ is zero. As $20\mathfrak a_2=g_2$ and $28\mathfrak a_4=g_3$,
\[
 (\mathcal W')^2=4\mathcal W^3-g_2\mathcal W-g_3 .
\]
With $e^y=-pu$ one has $\vartheta_u=\mathrm d/\mathrm d y$ and $\mathcal W=\mathfrak
c-\psi_2$, so, writing $\psi_3=\vartheta_u\psi_2$,
\begin{equation}\label{hol:td:eq:ellipticdelta}
 \psi_3^2=\Phi(\psi_2)=4(\mathfrak c-\psi_2)^3-g_2(\mathfrak c-\psi_2)-g_3 .
\end{equation}
The proof uses only the classical product, a Laurent expansion and the
constancy of an elliptic function without poles.

\subsection{From the complex identity to the Hahn field}
\label{hol:td:sub:formal}

Now let $p$ be an indeterminate again, and work in the ring $\Q[u,u^{-1}][[p]]$
of power series in $p$ whose coefficients are Laurent polynomials in $u$. The
series $\Psi=\sum_np^{n^2}u^n=1+p(u+u^{-1})+\cdots$ is a unit there, so
$\psi_1,\psi_2,\psi_3$ belong to it, and $\mathfrak c,g_2,g_3\in\Q[[p]]$. For
$u$ in the annulus $\frac12\le\abs u\le2$ and $0<p<\epsilon_0$, with
$\epsilon_0$ small, all these formal expansions converge absolutely to the
complex functions of Section~\ref{hol:td:sub:elliptic}: $\abs{\Psi-1}<\frac12$,
so $1/\Psi$ is a convergent geometric series. Hence, for each such $u$, the
power series in $p$ obtained from $\psi_3^2-\Phi(\psi_2)$ vanishes on
$(0,\epsilon_0)$, so all its coefficients vanish there; each coefficient is a
Laurent polynomial vanishing on an annulus, hence zero. Thus
\eqref{hol:td:eq:ellipticdelta} is an identity in $\Q[u,u^{-1}][[p]]$, over
$\Q$, and not merely for particular complex nomes.

\emph{From $u$ to $z$.} The map $\iota:\Q[z][[p]]\to\Q[u,u^{-1}][[p]]$,
$\sum_kP_k(z)p^k\mapsto\sum_kP_k(u+u^{-1})p^k$, is an injective ring
homomorphism. With $\check z=u-u^{-1}$,
\begin{equation}\label{hol:td:eq:chain}
 \vartheta_u(u+u^{-1})=\check z,\qquad \vartheta_u\check z=u+u^{-1},\qquad
 \check z^2=\iota(z^2-4),\qquad \vartheta_u(\iota g)=\check z\,\iota(D_zg),
\end{equation}
the last by the chain rule, coefficientwise in $p$. The formal series
$\mathcal T=1+\sum_{n\ge1}p^{n^2}C_n(z)\in\Q[z][[p]]$ has $\iota(\mathcal
T)=\Psi$ by \eqref{hol:td:eq:chebsub}. From \eqref{hol:td:eq:chain} one
computes $\iota(\mathcal H[\mathcal T])=\Psi\vartheta_u^2\Psi-(\vartheta_u\Psi)^2
=\Psi^2\psi_2$, then $\check z\,\iota(\mathcal K[\mathcal T])=\Psi^3\psi_3$
and $\iota(\mathcal B[\mathcal T])=\Psi^2(\mathfrak c-\psi_2)$. Hence
$\iota$ maps the left side of \eqref{hol:td:eq:cleared} for $\mathcal T$ to
$\Psi^6\bigl(\psi_3^2-\Phi(\psi_2)\bigr)=0$, and injectivity gives
\eqref{hol:td:eq:cleared} in $\Q[z][[p]]$.

\emph{From $\Q[[p]]$ to $\KK$.} For every $\mu>0$ the map
$s_\mu:\Q[[p]]\to\KK$, $\sum_kc_kp^k\mapsto\sum_kc_kt^{k\mu}$, is an injective
ring homomorphism: the support lies in the well-ordered set $\{k\mu\}$, each
exponent comes from one $k$, and the Cauchy product is preserved. Applied
coefficientwise it is a ring homomorphism $\Q[[p]][[z]]\to\KK[[z]]$ commuting
with $D_z$, and $\Q[z][[p]]\subseteq\Q[[p]][[z]]$ as subrings of $\Q[[p,z]]$.
It maps the formal $\mathcal T$ to $\mathcal T_p$, whose coefficients are
\eqref{hol:td:eq:a0}--\eqref{hol:td:eq:am}, and the formal constants to
\eqref{hol:td:eq:constants}. This proves \eqref{hol:td:eq:cleared} in
$\KK[[z]]$ for every $\mu>0$. Dividing by $f^6$ in $\KK((z))$ gives
\eqref{hol:td:eq:rationalODE}.

Inverses are taken only in $\Q[u,u^{-1}][[p]]$, where $\Psi$ is a unit, and
only the cleared identity is specialized. This avoids the incorrect claim that
the $p$-adic geometric expansion of $1/\Psi_p(u)$ converges at every Hahn
argument, a point source~13 makes explicitly; in a Hahn function field the
inverse is taken as a fraction.

\subsection{The corner of the cleared equation}\label{hol:td:sub:corner}

\begin{proposition}[Corner data of the cleared equation; merge]
\label{hol:td:prop:corner}
Let $P$ be the polynomial \eqref{hol:td:eq:jetpoly} with the constants
\eqref{hol:td:eq:constants}, written in the combined-monomial form
\eqref{hol:nl:eq:P}. In the notation of Section~\ref{hol:nl:sub:corner},
$d_P=6$, $w_P=0$, $\beta_P=0$, and
\[
 I_P(T)=-\Phi(0)=-\bigl(4\mathfrak c^3-g_2\mathfrak c-g_3\bigr),\qquad
 \chi_P(T)=0 .
\]
The constant $4\mathfrak c^3-g_2\mathfrak c-g_3$ equals
$-4p^2-48p^4-144p^6-448p^8-600p^{10}+O(p^{12})$; so $I_P$ is a nonzero
constant of valuation $2\mu$ and leading coefficient $4$, while the residue
corner polynomial vanishes identically.
\end{proposition}

\begin{proof}
Every monomial of $P$ has total degree six in the $Y_j$, so $d_P=6$. For the
weight $w_\nu=k_\nu-\sum_jjm_{\nu j}$ of \eqref{hol:nl:eq:dw}, every term of
$\mathrm H$ has weight at most $0$, with weight-zero part
$\mathrm H_0=z^2(Y_0Y_2-Y_1^2)+zY_0Y_1$ (the terms
$-4(Y_0Y_2-Y_1^2)$ have weight $-2$). The total derivative lowers the weight of
each term by exactly one, so the terms of $\mathrm H_1$ have weight at most
$-1$, with weight $-1$ part the total derivative $\mathrm H_0'$ of $\mathrm
H_0$. Hence every monomial of $P$ has weight at most $0$, and its weight-zero
part is
\[
 P_0=z^2\bigl(Y_0\mathrm H_0'-2Y_1\mathrm H_0\bigr)^2
 -4\bigl(\mathfrak cY_0^2-\mathrm H_0\bigr)^3
 +g_2\bigl(\mathfrak cY_0^2-\mathrm H_0\bigr)Y_0^4+g_3Y_0^6 .
\]
The substitution $z\mapsto1$, $Y_j\mapsto\fall Tj$ is a ring homomorphism
$\KK[z,Y]\to\KK[T]$ sending a monomial of weight $0$ to the product of
falling factorials in \eqref{hol:nl:eq:chiI}. It sends $\mathrm H_0$ to
$T(T-1)-T^2+T=0$ and $\mathrm H_0'=z^2(Y_0Y_3-Y_1Y_2)+2z(Y_0Y_2-Y_1^2)+Y_0Y_1
+z(Y_1^2+Y_0Y_2)$ to
\[
 T(T-1)(T-2)-T^2(T-1)+2\bigl(T(T-1)-T^2\bigr)+T+T^2+T(T-1)=0,
\]
so $P_0$ goes to $-4\mathfrak c^3+g_2\mathfrak c+g_3=-\Phi(0)$. From
$\mathfrak c=\frac1{12}-2p^2+O(p^4)$, $g_2=\frac1{12}+20p^2+O(p^4)$ and
$g_3=-\frac1{216}+\frac73p^2+O(p^4)$ one finds
$4\mathfrak c^3=\frac1{432}-\frac16p^2+O(p^4)$ and
$g_2\mathfrak c=\frac1{144}+\frac32p^2+O(p^4)$, so
$4\mathfrak c^3-g_2\mathfrak c-g_3=-4p^2+O(p^4)$; the further terms were
computed with SymPy. In particular the coefficient of the weight-zero monomial
$Y_0^6$ in $P$ is $-\Phi(0)\ne0$, so $P_0\ne0$, $w_P=0$, and $I_P=-\Phi(0)$.
The monomial $z^6Y_0^4Y_3^2$ of $P_0$ has coefficient $1$, and every
coefficient of $P$ lies in $\Z[\mathfrak c,g_2,g_3]\subseteq s_\mu(\Q[[p]])$,
which has nonnegative valuation; so $\beta_P=0$. Since $\res$ is a ring
homomorphism on elements of nonnegative valuation and the falling-factorial
products have integer coefficients, $\chi_P$ is obtained by applying $\res$
to the coefficients of $I_P$, and $\res(-\Phi(0))=0$ because
$v(\Phi(0))=2\mu>0$.
\end{proof}

The corner consists of $19$ monomials (SymPy). The computation has a
transparent form: the weight-zero part $\mathrm H_0$ of $\mathcal H$ is
exactly the left side of the vanishing-corner equation
\eqref{hol:nl:eq:degenerate}, whose residue corner polynomial is the identity
\eqref{hol:nl:eq:degeneratechi}.

\begin{remark}[The vanishing corner; source~13 and merge]
\label{hol:td:rem:corner}
The corner criterion of Theorem~\ref{hol:nl:thm:corner} is a sufficient
rigidity criterion, not a theorem for every higher-order equation, and the
cleared equation necessarily lies outside it. Source~13 records the
cancellation directly: on a trial monomial $f=z^N$ the top-degree part of
$\mathcal A[f]$ vanishes, because $z^2(-N/z^2)+z(N/z)=0$, so retaining only
separate derivative degrees and ignoring their cancellation would miss the
example. Proposition~\ref{hol:td:prop:corner} makes this exact and adds three
consequences. (1) The equation is consistent with
Theorems~\ref{hol:nl:thm:certificate} and~\ref{hol:nl:thm:corner}: they require
$\chi_P\ne0$, and here $\chi_P=0$. (2) When $\mu$ is an order unit,
$\mathcal T_p$ satisfies no nonzero equation with $d_P\ge1$ and
$\chi_P\ne0$ at all (and a nonzero $P$ with $d_P=0$ is a nonzero element of
$\KK[z]$, which no series satisfies), since by Theorem~\ref{hol:nl:thm:corner}
such an equation has only polynomial strongly entire solutions; so a nonpolynomial strongly entire series can satisfy an
equation whose residue corner vanishes identically without satisfying any
equation to which Theorem~\ref{hol:nl:thm:certificate} applies. (3) The
hypothesis $\chi_P\ne0$ of those theorems cannot be weakened to $I_P\ne0$:
here the full corner polynomial is a nonzero constant of positive valuation,
while the residue corner vanishes. The containment
\eqref{hol:nl:eq:rootcontainment} and
Proposition~\ref{hol:nl:prop:higherdegrees}, which assume $\chi_P\ne0$, are
unaffected. None of this contradicts the first-order or D-finite rigidity
theorems.
\end{remark}

\subsection{The Laurent ring and the substitution
\texorpdfstring{$z=u+u^{-1}$}{z = u + 1/u}}\label{hol:td:sub:laurent}

In this subsection and the next two, $\mu$ is an order unit. This supplies a
common Laurent function field for all nonzero arguments; the lower bound
concerns formal functions, not the algebraic independence of values in their
own ambient field.

Let $\mathcal U_\mu$ be the set of formal Laurent series
$b=\sum_{m\in\Z}b_mu^m$ with $b_m\in\KK$ such that, for all
$\gamma,\beta\in\Gamma$, only finitely many $m$ with $b_m\ne0$ satisfy
$v(b_m)+m\gamma\le\beta$. This is source~13's condition
``$v(b_m)+m\gamma\to+\infty$ as $\abs m\to\infty$'', written in the form of
Proposition~\ref{hol:cf:prop:criterion}\,(iii). For each $\gamma$ the weighted
values $v(b_m)+m\gamma$ of the nonzero coefficients then have a least element,
attained by finitely many $m$.

\begin{lemma}[The Laurent ring; source~13]\label{hol:td:lem:laurent}
\begin{enumerate}[label=\textup{(\alph*)}]
\item $\mathcal U_\mu$ is an integral domain containing $\KK[u,u^{-1}]$, with
the product $\bigl(\sum b_mu^m\bigr)\bigl(\sum c_mu^m\bigr)=\sum_m
\bigl(\sum_{m_1+m_2=m}b_{m_1}c_{m_2}\bigr)u^m$, the inner sums being strong
sums.
\item $\vartheta_u\bigl(\sum b_mu^m\bigr)=\sum mb_mu^m$ is a derivation of
$\mathcal U_\mu$ vanishing on $\KK$, and
$\mathsf S\bigl(\sum b_mu^m\bigr)=\sum b_mp^{2m}u^m$ is a ring automorphism
fixing $\KK$ and commuting with $\vartheta_u$. Both extend to
$\operatorname{Frac}\mathcal U_\mu$.
\item $\Psi_p$ and all $\vartheta_u^j\Psi_p$ belong to $\mathcal U_\mu$, and
$\Psi_p\ne0$.
\end{enumerate}
\end{lemma}

\begin{proof}
(a) With $\gamma=0$ the coefficient values of $b$ and $c$ are bounded below
and escape cofinally, so for fixed $m$ the family $(b_{m_1}c_{m-m_1})_{m_1}$ is
strongly summable by Lemma~\ref{hol:lem:certificate}\,(a). For fixed $\gamma$,
if the product coefficient $d_m\ne0$ has $v(d_m)+m\gamma\le\beta$, some term
has $(v(b_{m_1})+m_1\gamma)+(v(c_{m_2})+m_2\gamma)\le\beta$, and the lower
bounds of the two weighted families leave finitely many pairs, as in the proof
of Proposition~\ref{hol:cf:prop:ring}; so the product lies in $\mathcal
U_\mu$. Commutativity, associativity and distributivity follow from the
associativity of strong sums applied to the families $(b_{m_1}c_{m_2}e_{m_3})$.
For integrality, let $w_b$ and $w_c$ be the least coefficient values of
nonzero $b$ and $c$ and $\operatorname{in}(b)=\sum_{v(b_m)=w_b}\lc(b_m)u^m\in
k[u,u^{-1}]$, a nonzero Laurent polynomial. The coefficient of
$t^{w_b+w_c}u^m$ in $bc$ is the coefficient of $u^m$ in
$\operatorname{in}(b)\operatorname{in}(c)$, because every term has value at
least $w_b+w_c$. Since $k[u,u^{-1}]$ is a domain, $bc\ne0$.

(b) Multiplying coefficients by ordinary integers preserves the defining
condition, and Leibniz's rule holds coefficientwise since $m=m_1+m_2$. For
$\mathsf S$, $v(b_mp^{2m})+m\gamma=v(b_m)+m(\gamma+2\mu)$, the inverse uses
$p^{-2m}$, and $p^{2m}=p^{2m_1}p^{2m_2}$ with
Lemma~\ref{hol:lem:support}\,(3) gives multiplicativity. Both operators act
diagonally on $u^m$, so they commute. An automorphism and a derivation of a
domain extend to its fraction field.

(c) If $\abs\gamma\le N\mu$ and $\beta\le B\mu$, then
$v(p^{n^2})+n\gamma\ge(n^2-N\abs n)\mu>\beta$ for all but finitely many $n$;
multiplying coefficients by $n^j$ changes no valuation.
\end{proof}

\begin{lemma}[Substitution of $u+u^{-1}$; merge]\label{hol:td:lem:substitution}
For $f=\sum_na_nz^n\in\EK$ and $m\in\Z$ the family
$\bigl(a_n\binom n{(n-m)/2}\bigr)_{n\ge\abs m,\ n\equiv m\,(2)}$ is strongly
summable. Write $b_m$ for its sum and $f^\natural=\sum_mb_mu^m$. Then:
\begin{enumerate}[label=\textup{(\alph*)}]
\item $f^\natural\in\mathcal U_\mu$;
\item $f\mapsto f^\natural$ is a $\KK$-algebra homomorphism
$\EK\to\mathcal U_\mu$, equal to the substitution $z\mapsto u+u^{-1}$ on
$\KK[z]$;
\item $\vartheta_u(f^\natural)=\check z\,(D_zf)^\natural$, where $\check
z=u-u^{-1}$;
\item $\mathcal T_p^\natural=\Psi_p$.
\end{enumerate}
\end{lemma}

\begin{proof}
Write $\beta_{n,m}=\binom n{(n-m)/2}$ when $n\ge\abs m$ and $n\equiv m\pmod
2$, and $\beta_{n,m}=0$ otherwise; this is the coefficient of $u^m$ in
$(u+u^{-1})^n$. By Proposition~\ref{hol:cf:prop:criterion}\,(iii) with
$\gamma=0$, the coefficient values of $f$ escape cofinally, and the nonzero
integers $\beta_{n,m}$ have value zero, so Lemma~\ref{hol:lem:certificate}\,(a)
gives strong summability.

(a) Fix $\gamma,\beta$ and put $\gamma'=-\abs\gamma$. If $b_m\ne0$ and
$v(b_m)+m\gamma\le\beta$, some $n\ge\abs m$ with $a_n\ne0$ has
$v(a_n)+m\gamma\le\beta$, and then
$v(a_n)+n\gamma'\le v(a_n)+m\gamma\le\beta$ because $n\gamma'\le\abs
m\gamma'\le m\gamma$. By Proposition~\ref{hol:cf:prop:criterion}\,(iii) at
$\gamma'$ only finitely many $n$ qualify, and each serves only the $m$ with
$\abs m\le n$.

(b) Linearity over $\KK$ holds because strong sums are linear. For
multiplicativity fix $m$ and consider the family
$\bigl(a_ic_j\beta_{i,m_1}\beta_{j,m_2}\bigr)$ indexed by $(i,j,m_1,m_2)$
with $m_1+m_2=m$, where $g=\sum c_jz^j\in\EK$. The double family $(a_ic_j)$ is
strongly summable, as in the proof of Lemma~\ref{hol:nl:lem:algebra}; each
pair $(i,j)$ carries finitely many $(m_1,m_2)$, and nonzero integer factors
change neither supports nor multiplicities. Grouping by $(m_1,m_2)$ and then by
$m_1$ gives the coefficient of $u^m$ in $f^\natural g^\natural$. Grouping by
$(i,j)$ gives $a_ic_j\beta_{i+j,m}$, because
$(u+u^{-1})^i(u+u^{-1})^j=(u+u^{-1})^{i+j}$, and grouping those by $n=i+j$
gives $\bigl(\sum_{i+j=n}a_ic_j\bigr)\beta_{n,m}$, whose sum is the coefficient
of $u^m$ in $(fg)^\natural$. By associativity of strong sums the two totals
agree. On polynomials all families are finite.

(c) Since $\vartheta_u(u+u^{-1})^n=n(u+u^{-1})^{n-1}(u-u^{-1})$, one has
$m\beta_{n,m}=n(\beta_{n-1,m-1}-\beta_{n-1,m+1})$. The coefficient of $u^m$
in $\vartheta_u(f^\natural)$ is $mb_m=\sum_na_nm\beta_{n,m}$, and that of
$\check z\,(D_zf)^\natural$ is
$\sum_nna_n(\beta_{n-1,m-1}-\beta_{n-1,m+1})$; they agree termwise, and
$D_zf\in\EK$ by Proposition~\ref{hol:cf:prop:ring}.

(d) With $c_{N,j}$ the integers of \eqref{hol:td:eq:chebcoef}, the coefficient
of $u^m$ in $\mathcal T_p^\natural$ is the sum of the family
$\bigl(p^{N^2}c_{N,j}\beta_{N-2j,m}\bigr)_{N,j}$, grouped by $n=N-2j$ (the
constant term included with $N=0$). Its exponents $N^2\mu$ each come from
finitely many $j$, so it is strongly summable. Grouping instead by $N$ gives
$p^{N^2}$ times the coefficient of $u^m$ in $C_N(u+u^{-1})=u^N+u^{-N}$,
which is nonzero only for $N=\abs m$. The sum is $p^{m^2}$.
\end{proof}

Injectivity of $f\mapsto f^\natural$ is not needed. By (c) and
Proposition~\ref{hol:cf:prop:ring}, $(D_z^rf)^\natural=(\check
z^{-1}\vartheta_u)^rf^\natural$ in $\operatorname{Frac}\mathcal U_\mu$, which
equals $\check z^{-r}\vartheta_u^rf^\natural$ plus a combination of lower
$\vartheta_u^jf^\natural$ with coefficients in $\Q(u)$. For $\mathcal T_p$,
\begin{equation}\label{hol:td:eq:jets}
 \Psi_p=\mathcal T_p^\natural,\qquad
 \vartheta_u\Psi_p=\check z\,(D_z\mathcal T_p)^\natural,\qquad
 \vartheta_u^2\Psi_p=\check z^2(D_z^2\mathcal T_p)^\natural
 +(u+u^{-1})(D_z\mathcal T_p)^\natural .
\end{equation}

\subsection{Three independent theta quantities}\label{hol:td:sub:independent}

\begin{lemma}[Three rational-function obstructions; source~13]
\label{hol:td:lem:rational}
Let $E$ be a field of characteristic zero and $a\in E^\times$ not a root of
unity, and let $\mathsf S_a$ be the automorphism of $E(u)$ with
$\mathsf S_au=au$ fixing $E$.
\begin{enumerate}[label=\textup{(\alph*)}]
\item If $\mathsf S_ar=r$, then $r\in E$.
\item For $d\in E^\times$ there is no $r\in E(u)$ with $\mathsf S_ar-r=d$.
\item On $E(u,V)$, with $V$ transcendental, extend $\mathsf S_a$ by
$\mathsf S_aV=V-1$. Then $\mathsf S_a$ preserves the $u$-adic order
$\operatorname{ord}_u$ of every nonzero element.
\end{enumerate}
\end{lemma}

\begin{proof}
Expand a rational function as a formal Laurent series at $u=0$. For (a), each
coefficient of degree $n\ne0$ is multiplied by $a^n\ne1$, so invariance leaves
only degree zero. For (b), the coefficient of $u^0$ in $r(au)-r(u)$ is zero,
also when $r$ has a pole at $0$. For (c), a lowest nonzero coefficient in
$E(V)((u))$ is sent to the nonzero coefficient obtained by $V\mapsto V-1$,
multiplied by a nonzero power of $a$; its degree is unchanged.
\end{proof}

In $\operatorname{Frac}\mathcal U_\mu$ put $\Psi=\Psi_p$,
$\psi_1=\vartheta_u\Psi/\Psi$, $\psi_2=\vartheta_u\psi_1$ and
$\psi_3=\vartheta_u\psi_2$; the field $\KK(u)$ lies in it and is preserved by
$\mathsf S$. By \eqref{hol:td:eq:shift},
\begin{equation}\label{hol:td:eq:sigmajets}
 \mathsf S\Psi=p^{-1}u^{-1}\Psi,\qquad
 \mathsf S\psi_1=\psi_1-1,\qquad
 \mathsf S\psi_2=\psi_2 .
\end{equation}

\begin{lemma}[Transfer, and a global certificate for $\psi_3\ne0$; merge]
\label{hol:td:lem:psi3}
\begin{enumerate}[label=\textup{(\alph*)}]
\item For every polynomial $\Pi$ in four variables with coefficients in
$\Q[[p]]$ and every $m\in\Z$, the coefficient of $u^m$ in
\[
 \Pi(\Psi,\vartheta_u\Psi,\vartheta_u^2\Psi,\vartheta_u^3\Psi),
\]
computed in $\mathcal U_\mu$ with the coefficients of $\Pi$ mapped by
$s_\mu$, is the image under $s_\mu$ of the same coefficient computed in
$\Q[u,u^{-1}][[p]]$. In particular
$\psi_3^2=\Phi(\psi_2)$ in $\operatorname{Frac}\mathcal U_\mu$.
\item In $\mathcal U_\mu$,
\[
 \Psi^3\psi_3=\mathcal N:=\Psi^2\vartheta_u^3\Psi
 -3\Psi\,\vartheta_u\Psi\,\vartheta_u^2\Psi+2(\vartheta_u\Psi)^3,
\]
and the coefficient of $u$ in $\mathcal N$ has valuation $\mu$ and leading
coefficient $1$. In particular $\psi_3\ne0$.
\end{enumerate}
\end{lemma}

\begin{proof}
(a) The four series $\vartheta_u^j\Psi=\sum_nn^jp^{n^2}u^n$ lie in
$\mathcal U_\mu$ (Lemma~\ref{hol:td:lem:laurent}). By associativity of strong
sums, the coefficient of $u^m$ on the $\mathcal U_\mu$ side is the strong sum
of the family of terms $e\,t^{(k+n_1^2+\cdots+n_r^2)\mu}$, $e\in\Z$, indexed by
a monomial of $\Pi$ of degree $r$, a term $ep^k$ of its coefficient, and a
tuple $(n_1,\ldots,n_r)$ with $n_1+\cdots+n_r=m$. The exponent $K\mu$ is
carried by the finitely many members with $k+\sum n_i^2=K$, whose integer
coefficients add up to the coefficient of $u^mp^K$ on the formal side. This is
the assertion. Applied to the cleared form of \eqref{hol:td:eq:ellipticdelta},
$\Psi^6(\psi_3^2-\Phi(\psi_2))$, which is a polynomial in the four series
with coefficients in $\Z[\mathfrak c,g_2,g_3]$ and vanishes in
$\Q[u,u^{-1}][[p]]$ (Section~\ref{hol:td:sub:formal}), it gives the identity
in $\operatorname{Frac}\mathcal U_\mu$, because $\Psi\ne0$ there.

(b) The identity is the quotient rule applied twice. By (a), the coefficient
of $u$ in $\mathcal N$ is the image under the injective ring homomorphism
$s_\mu$ of the coefficient of $u$ in the same expression computed in
$\Q[u,u^{-1}][[p]]$. There $\Psi=1+p(u+u^{-1})+O(p^4)$, $\vartheta_u\Psi=\vartheta_u^3\Psi
=p(u-u^{-1})+O(p^4)$ and $\vartheta_u^2\Psi=p(u+u^{-1})+O(p^4)$, so
\[
 \mathcal N=p(u-u^{-1})+O(p^2).
\]
The coefficient of $u$ is $p+O(p^2)$; SymPy gives
$p-8p^3+20p^5-70p^9+64p^{11}+O(p^{12})$. Its image $t^\mu+\cdots$ is nonzero,
so $\mathcal N\ne0$, and since $\mathcal U_\mu$ is a domain and $\Psi\ne0$,
$\psi_3=\mathcal N/\Psi^3\ne0$.
\end{proof}

Source~13's proof of this step expanded $\psi_1$ at the simple zero
$u_0=-p$ of $\Psi$ as $u_0/(u-u_0)+O(1)$ and read off
$\psi_3=2u_0^3/(u-u_0)^3+O((u-u_0)^{-2})$; the lemma replaces that local
argument, which used expansions in $\operatorname{Frac}\mathcal U_\mu$ that
were not defined. Neither route uses a transcendence theorem for complex theta
values.

\begin{theorem}[Three independent theta quantities; source~13]
\label{hol:td:thm:independent}
Let $\mu$ be an order unit. The elements $\Psi,\psi_1,\psi_2$ of
$\operatorname{Frac}\mathcal U_\mu$ are algebraically independent over
$\KK(u)$. Equivalently, $\Psi,\vartheta_u\Psi,\vartheta_u^2\Psi$ are
algebraically independent over $\KK(u)$.
\end{theorem}

\begin{proof}
\emph{$\psi_2$ is transcendental over $\KK(u)$.} Otherwise let $\pi(T)$ be its
monic minimal polynomial. Since $\mathsf S\psi_2=\psi_2$ and $\mathsf S$
preserves $\KK(u)$, applying $\mathsf S$ to the coefficients gives the same
monic minimal polynomial, so by Lemma~\ref{hol:td:lem:rational}\,(a), with
$E=\KK$ and $a=p^2$, all its coefficients lie in $\KK$. Applying $\vartheta_u$
to $\pi(\psi_2)=0$ gives $\pi'(\psi_2)\psi_3=0$. By separability in
characteristic zero $\pi'(\psi_2)\ne0$, contradicting
Lemma~\ref{hol:td:lem:psi3}.

\emph{$\psi_1$ is transcendental over $\KK(u,\psi_2)$.} This field is $E(u)$
with $E=\KK(\psi_2)$ fixed by $\mathsf S$. Suppose the monic minimal polynomial
of $\psi_1$ over it has degree $d\ge1$ and begins $T^d+bT^{d-1}+\cdots$.
Applying $\mathsf S$ and then replacing $T$ by $T-1$ gives another monic
polynomial vanishing at $\psi_1$, because $\mathsf S\psi_1=\psi_1-1$; its next
coefficient is $\mathsf Sb-d$. Uniqueness gives $\mathsf Sb-b=d$, contradicting
Lemma~\ref{hol:td:lem:rational}\,(b).

\emph{$\Psi$ is transcendental over $\KK(u,\psi_2,\psi_1)$.} Let its
hypothetical monic minimal polynomial have degree $d$ and constant coefficient
$b_0$, nonzero because $\Psi\ne0$. Applying $\mathsf S$, substituting
$p^{-1}u^{-1}T$ for $T$ and dividing by $(p^{-1}u^{-1})^d$ gives another monic
polynomial vanishing at $\Psi$, so $\mathsf Sb_0=(p^{-1}u^{-1})^db_0$. By the
first two steps the field is the rational function field
$\KK(\psi_2)(u,\psi_1)$, on which $\mathsf S$ acts as in
Lemma~\ref{hol:td:lem:rational}\,(c) with $V=\psi_1$. The left side has the
$u$-adic order of $b_0$, the right side that order minus $d$, which is
impossible.

Finally $\vartheta_u\Psi=\psi_1\Psi$ and
$\vartheta_u^2\Psi=(\psi_2+\psi_1^2)\Psi$, and these formulas and their
rational inverses identify the fields generated by the two triples.
\end{proof}

\begin{corollary}[Exact minimal order; source~13]\label{hol:td:cor:order}
Let $\mu$ be an order unit. Then $\mathcal T_p$, $D_z\mathcal T_p$ and
$D_z^2\mathcal T_p$ are algebraically independent over $\KK(z)$, and the least
order of a nonzero algebraic differential equation for $\mathcal T_p$ over
$\KK(z)$ is three.
\end{corollary}

\begin{proof}
Suppose that a nonzero polynomial $P$ in $Y_0,Y_1,Y_2$ with coefficients in
$\KK[z]$, obtained after clearing denominators, satisfies
$P(z,\mathcal T_p,D_z\mathcal T_p,D_z^2\mathcal T_p)=0$. Apply
Lemma~\ref{hol:td:lem:substitution}\,(b): the image of this identity in
$\mathcal U_\mu$ is $P^\natural(\Psi,(D_z\mathcal T_p)^\natural,(D_z^2\mathcal
T_p)^\natural)=0$, where $P^\natural$ has the coefficients of $P$ with
$u+u^{-1}$ substituted for $z$. By \eqref{hol:td:eq:jets} the three $u$-jets
are obtained from $\Psi,\vartheta_u\Psi,\vartheta_u^2\Psi$ by a triangular
linear change with diagonal $1,\check z^{-1},\check z^{-2}$, invertible over
$\KK(u)$ because $\check z\ne0$. The polynomial $P^\natural$ is nonzero,
because $z\mapsto u+u^{-1}$ is injective on $\KK[z]$, so this is a nonzero
relation among $\Psi,\vartheta_u\Psi,\vartheta_u^2\Psi$ over $\KK(u)$,
contradicting Theorem~\ref{hol:td:thm:independent}. This proves the lower
bound; Theorem~\ref{hol:td:thm:ode} gives the upper bound.
\end{proof}

\subsection{The differential field and non-D-finiteness}\label{hol:td:sub:field}

Let $\mu$ be an order unit and write, in $\KK((z))$,
\[
 f=\mathcal T_p,\qquad \mathcal L=D_zf/f,\qquad \mathcal A=\mathcal A[f],\qquad
 \mathcal A_1=D_z\mathcal A .
\]
The first three are algebraically independent over $\KK(z)$, being an
invertible rational change of $f,D_zf,D_z^2f$.

\begin{theorem}[Differential-field presentation; source~13]
\label{hol:td:thm:field}
Let $\mu$ be an order unit. The differential field generated by $\mathcal T_p$
over $\KK(z)$ is
\begin{equation}\label{hol:td:eq:field}
 \KK(z)\langle\mathcal T_p\rangle=\KK(z)(f,\mathcal L,\mathcal A,\mathcal A_1),
 \qquad (z^2-4)\mathcal A_1^2=\Phi(\mathcal A).
\end{equation}
It has transcendence degree three over $\KK(z)$ and degree two over the
rational function field $\KK(z)(f,\mathcal L,\mathcal A)$. Its derivation is
\begin{align}
 D_zf&=\mathcal Lf,& D_z\mathcal L&=\frac{\mathcal A-z\mathcal L}{z^2-4},
 \label{hol:td:eq:system1}\\
 D_z\mathcal A&=\mathcal A_1,&
 D_z\mathcal A_1&=\frac{-12(\mathfrak c-\mathcal A)^2+g_2-2z\mathcal A_1}{2(z^2-4)}.
 \label{hol:td:eq:system2}
\end{align}
Consequently the minimal polynomial of $D_z^3\mathcal T_p$ over
$\KK(z)(\mathcal T_p,D_z\mathcal T_p,D_z^2\mathcal T_p)$ has degree two, and
\eqref{hol:td:eq:cleared} is that minimal relation up to a nonzero factor.
\end{theorem}

\begin{proof}
Regard $\Phi(\mathcal A)/(z^2-4)$ as a rational function of the transcendental
$\mathcal A$ over $\KK(z)(f,\mathcal L)$. Its order at $\mathcal A=\infty$ is
$-3$, an odd integer, so it is not a square, and adjoining $\mathcal A_1$ is a
quadratic extension. In particular $\mathcal A_1\ne0$, as also follows from
$\Phi\ne0$ and the transcendence of $\mathcal A$.

The formulas for $D_zf$ and $D_z\mathcal L$ follow from the definitions, since
$\mathcal A=(z^2-4)D_z\mathcal L+z\mathcal L$. Differentiate
$(z^2-4)\mathcal A_1^2=\Phi(\mathcal A)$ and use $D_z\mathcal A=\mathcal A_1$
and $\Phi'(T)=-12(\mathfrak c-T)^2+g_2$; cancelling $\mathcal A_1\ne0$ gives
\eqref{hol:td:eq:system2}. So all derivatives of the four generators remain in
the field, and conversely the generators are rational in
$f,D_zf,D_z^2f,D_z^3f$; the field is the differential field generated by $f$.
Finally $\mathcal A_1=\mathcal K[f]/f^3$ is affine linear in $D_z^3f$ with
nonzero leading coefficient $(z^2-4)/f$, so the degree-two assertion for
$\mathcal A_1$ is the same assertion for $D_z^3f$, and clearing powers of $f$
introduces only a nonzero factor.
\end{proof}

The poles at $z=\pm2$ of this rational system describe a coordinate chart,
not poles of the entire $f$; the globally meaningful equation is the cleared
identity \eqref{hol:td:eq:cleared}. The cubic in the elliptic coordinate
$\mathfrak c-\mathcal A$ accounts for the algebraic part of the differential
field, while $\mathcal L$ and $f$ carry the logarithmic and multiplicative
theta information.

\begin{proposition}[Non-D-finiteness; source~13]\label{hol:td:prop:nondfinite}
Let $\mu$ be an order unit. Then $\mathcal T_p$ satisfies no nonzero linear
differential equation over $\KK(z)$.
\end{proposition}

\begin{proof}
Minimal algebraic order three does not by itself exclude a linear equation of
higher order, so a separate argument is needed. In the $u$-coordinate,
\eqref{hol:td:eq:ellipticdelta} holds in $\operatorname{Frac}\mathcal U_\mu$,
$\psi_3^2=\Phi(\psi_2)$, and,
after differentiation and cancellation of $\psi_3\ne0$,
$\vartheta_u\psi_3=\frac12\Phi'(\psi_2)$
(Lemma~\ref{hol:td:lem:psi3}). Hence $\KK(\psi_2,\psi_3)$ is closed
under $\vartheta_u$ and algebraic over $\KK(\psi_2)$, and by
Theorem~\ref{hol:td:thm:independent} $\psi_1$ is transcendental over
$\KK(u,\psi_2,\psi_3)$. Inductively,
\[
 \frac{\vartheta_u^n\Psi}{\Psi}=\psi_1^n+\text{a polynomial of degree less
 than }n\text{ in }\psi_1\text{ over }\KK(\psi_2,\psi_3):
\]
apply $\vartheta_u$, add multiplication by $\psi_1$, and use
$\vartheta_u\psi_1=\psi_2$. A nontrivial $\KK(u)$-linear combination of the
$\vartheta_u^n\Psi$, divided by $\Psi$, has a nonzero highest power of the
transcendental $\psi_1$, so it does not vanish. A nonzero linear operator
$\sum_{r\le R}\rho_r(z)D_z^r$ with $\rho_R\ne0$ becomes, through
Lemma~\ref{hol:td:lem:substitution} and the triangular formula after it, a
$\KK(u)$-linear combination of $\vartheta_u^r\Psi$, $r\le R$, whose coefficient
of $\vartheta_u^R\Psi$ is $\rho_R(u+u^{-1})\check z^{-R}\ne0$.
\end{proof}

For $k=\C$ the proposition is also a case of Theorem~\ref{hol:main:ode}, since
$\mathcal T_p$ is then a nonpolynomial element of $\E$; the argument above
works for every $k$ and uses no escape chain.

\section{The exact boundary, surreal consequences and scope}
\label{hol:td:sec:boundary}

\subsection{Proof of Theorem~\ref{hol:td:main:boundary}, and Theorems~E
and~G}\label{hol:td:sub:proofs}

\begin{proof}[Proof of Theorem~\ref{hol:td:main:boundary}]
(ii)$\Rightarrow$(i) is the contrapositive of
Theorem~\ref{hol:cf:main:entire}\,(a) (sources~11 and~12; source~13 proves it
again, Remark~\ref{hol:td:rem:euler}). For (i)$\Rightarrow$(iii) let $\mu$ be
an order unit and $p=t^\mu$. Theorem~\ref{hol:td:thm:domain} gives
$\mathcal T_p\in\EK$ and Proposition~\ref{hol:td:prop:coeff} nonpolynomiality;
Theorem~\ref{hol:td:thm:ode}, Corollary~\ref{hol:td:cor:order} and
Theorem~\ref{hol:td:thm:field} give an equation of order three, algebraic
independence of $\mathcal T_p,D_z\mathcal T_p,D_z^2\mathcal T_p$ and degree two
of $D_z^3\mathcal T_p$ over the field they generate.
(iii)$\Rightarrow$(ii) is immediate. Non-D-finiteness is
Proposition~\ref{hol:td:prop:nondfinite}. Finally $\mathcal T_p$ is not in
$\KK(z)$, since it is transcendental over $\KK(z)$.
\end{proof}

\begin{remark}[Source~13's proof of the no-order-unit direction]
\label{hol:td:rem:euler}
Section~7 of source~13 proves (ii)$\Rightarrow$(i) self-contained, so that its
classification does not rest on a citation. It is the argument of
Theorem~\ref{hol:cf:main:entire}\,(a) and is printed once there. Source~13
writes a relation with the Euler operator $\vartheta=zD_z$, which is a
triangular rational change of the ordinary derivatives, chooses it of least
total degree in the jets, and linearizes along the solution: if
$\pi_j=\partial P/\partial Y_j$ evaluated at $(f,\vartheta f,\ldots)$, $\theta'$
is the least $z$-order of a nonzero $\pi_j$ and
$I(T)=\sum_j[z^{\theta'}]\pi_j\,T^j$, then, because $\vartheta$ does not lower
$z$-order, $I(n)a_n=-[z^{n+\theta'}]P(z,f_{<n},\vartheta f_{<n},\ldots)$ for
$n>\theta'$. With the Euler operator the linearization multiplier is a
polynomial in the index. Its Lemmas~7.1 and~7.2 are
Theorem~\ref{hol:cf:thm:coefficients}\,(a) and
Lemmas~\ref{hol:cf:lem:valuerank}--\ref{hol:cf:lem:principal}, and its
Theorem~7.3, via the exterior obstruction of
Proposition~\ref{hol:cf:prop:exterior}, is
Theorem~\ref{hol:cf:main:entire}\,(a). Source~13 credits this direction to the
coefficient-field part and claims no new finite-generation principle.
\end{remark}

\begin{proof}[Proof of Theorem~\ref{hol:cf:main:sharp}\,\textup{(iv)}]
(i)$\Leftrightarrow$(iv) is (i)$\Leftrightarrow$(ii) of
Theorem~\ref{hol:td:main:boundary}; with the proof of
Theorem~\ref{hol:cf:main:sharp} given in Section~\ref{hol:cf:sub:sharp}, all
four conditions are equivalent.
\end{proof}

\begin{corollary}[Theorem~E is sharp; merge]\label{hol:td:cor:sharp}
For $\Gamma$ as in Convention~\ref{hol:conv:group}:
\begin{enumerate}[label=\textup{(\alph*)}]
\item every nonpolynomial $f\in\EK$ is differentially transcendental over
$\KK(z)$ if and only if $\Gamma$ has no order unit;
\item the elements of $\FE$ that are differentially algebraic over $\KK(z)$
are exactly those of $\KK(z)$ if and only if $\Gamma$ has no order unit.
\end{enumerate}
If $\Gamma$ has an order unit $\mu$, then for $f=\mathcal T_p$ and every
$q\in\KK^\times$ of infinite multiplicative order, including $v(q)=0$, the
family of Theorem~\ref{hol:cf:main:sharp}\,\textup{(iii)} is algebraically
dependent over $\KK(z)$.
\end{corollary}

\begin{proof}
(a) is Theorem~\ref{hol:td:main:boundary}, (i)$\Leftrightarrow$(ii). For (b),
Theorem~\ref{hol:cf:main:entire}\,(c) gives the equality without an order unit;
with one, $\mathcal T_p\in\EK\subseteq\FE$ is differentially algebraic and not
in $\KK(z)$. The last assertion holds because the subfamily with $i=0$ is
already dependent (Theorem~\ref{hol:td:thm:ode}).
\end{proof}

\subsection{Consistency with the rest of the report}
\label{hol:td:sub:consistency}

The new part contradicts no earlier statement; the points of contact are these.
\begin{itemize}
\item \emph{Theorem~\ref{hol:cf:main:entire}.} Without an order unit
$\mathcal T_p$ is not entire: $\Dom(\mathcal T_p)=\mathcal O_{\Delta_\mu}\ne\KK$
(Theorem~\ref{hol:td:thm:domain}). Theorem~E(a) is thus not contradicted, and
by Corollary~\ref{hol:td:cor:sharp} its hypothesis is necessary.
\item \emph{Theorem~\ref{hol:nl:main:first} and the corner criterion.} The
minimal order is three, not one; and $\chi_P=0$ for the cleared equation, as
Theorem~\ref{hol:nl:thm:corner} requires of an equation with a nonpolynomial
entire solution, while $I_P\ne0$ (Proposition~\ref{hol:td:prop:corner},
Remark~\ref{hol:td:rem:corner}).
\item \emph{Theorem~\ref{hol:cf:main:rank}.} The coefficient values $m^2\mu$
span $\Q\mu$, so $\cvr(\mathcal T_p)=1$ and Theorem~F does not apply; it is a
sufficient condition only (Section~\ref{hol:cf:sub:boundaries}).
Theorem~\ref{hol:cf:thm:rankone} says nothing about series whose coefficient
values span $\Q\mu$, and $\mathcal T_p$ is such a series with an algebraic
differential relation.
\item \emph{Theorem~\ref{hol:cf:main:sharp}.} $\mathcal T_p$ is a second
witness for (ii) and (iii), and a stronger one for (iii): its jets are already
dependent at $i=0$, so (iii) holds for every $q$ of infinite order, including
$v(q)=0$. Its coefficients lie in a field finitely generated over $k$, by
Theorem~\ref{hol:cf:thm:coefficients}\,(a) with $k_0=k(\mathfrak c,g_2,g_3)$,
as (ii) requires.
\item \emph{Theorem~\ref{hol:cf:thm:descent}.} The coefficient values lie in
$\Z\mu$, whose smallest containing convex subgroup $\Delta_\mu$ is all of
$\Gamma$ exactly when $\mu$ is an order unit; the descent theorem needs a
proper convex subgroup and does not apply.
\item \emph{The partial theta series.} $\Theta_p$ and $G$
(Proposition~\ref{hol:cf:prop:theta}, Corollary~\ref{hol:nl:cor:theta}) are
untouched: source~13 explains why a bilateral descent is used instead, since
the dilation equation of $\Theta_p$ is not an algebraic differential equation.
Whether $\Theta_p$ or $G$ is differentially algebraic stays open.
\item \emph{Question~\ref{hol:q:mixedunit}.} It concerns \emph{linear} mixed
equations at unit valuation. $\mathcal T_p$ is not D-finite, and source~13
says nothing about linear mixed equations, so the question stayed open after
source~13; source~18 has since answered it positively
(Corollary~\ref{hol:ud:cor:unitrigidity}), and the nonlinear analogue is
settled by Corollary~\ref{hol:td:cor:sharp}.
\item \emph{Question~\ref{hol:cf:q:finiterank}.} $\mathcal T_p$ is
differentially algebraic and strongly entire, with a finitely generated
coefficient field whose values $m^2\mu$ are cofinal and grow faster than any
linear bound. With an order unit, therefore, no property shared by the
coefficient fields of all differentially algebraic series can force a linear
bound. The question asks for additional properties and stays open.
\item \emph{The Hahn--Tate report.} $\Psi_p$ is that report's theta series at
the nome $p^2$ and the argument $-pu$; its domain and product
(\texttt{tate:thm:theta}, \texttt{tate:prop:theta-product}, the latter for
$0\le v(u)<v(q)$) are specialized here, not claimed anew~\cite{repotate}.
\end{itemize}

\subsection{Surreal, surcomplex and omnific consequences}
\label{hol:td:sub:surreal}

Through \eqref{hol:eq:embedding} take $\Gamma=\R$ and $\mu=1$. Then
$\mathcal T_t$ is an actual nonpolynomial strongly entire function on the
set-sized surcomplex field $\C((\omega^{-\R}))\subseteq\No[i]$, with
coefficients in the smaller field $\Q((\omega^{-1}))$; the same holds on
$\C((\omega^{-\Z}))$ and $\C((\omega^{-\Q}))$, and algebraic closedness of the
workspace is not needed. The derivative $D_z$ still kills the coefficients,
including $p=\omega^{-1}$, so this is a statement about functions whose values
and coefficients are surreal or surcomplex, not about the scalar derivation of
$\No$.

\paragraph{Not entire on the class.}
With $p=\omega^{-1}$ the same coefficient formula can be evaluated at every
surreal or surcomplex point whose leading exponent is bounded above by an
ordinary integer; the support proof is pointwise and uses only set-sized
supports. It fails at $x=\omega^\omega$: the leading exponent of $a_mx^m$ in
Conway's convention is $m\omega-m^2$, and these increase strictly, so they are
not the decreasing well-ordered support of a strong surreal sum. The example
is therefore not strongly entire on $\No$ or $\No[i]$, as for $F$ in
Section~\ref{hol:sub:notallsurcomplex}. More generally, enlarging $\Gamma$ to
$\Gamma'$ while keeping $p=t^\mu$ changes the domain exactly to
$\{x:v(x)\ge-N\mu\text{ for some }N\}$ in the larger field: an extension can
turn an order unit into a non-order-unit and destroy entireness without
changing a coefficient or the equation.

\paragraph{Omnific integers.}
The omnific integers are $\Oz=\Z\oplus\Pi$, where $\Pi$ consists of the
surreals with purely positive Conway exponents~\cite[\texttt{odg:def:rings}]{repoodg}:
a nonzero term with negative exponent excludes membership. $\Oz$ is
discretely ordered with $1$ the least positive element, and a surreal that is
not an omnific integer lies between two consecutive ones, its omnific floor
and ceiling; for the numbers below this is read off their normal forms.

\begin{corollary}[Zeros next to omnific integers; source~13]
\label{hol:td:cor:omnific}
For $p=\omega^{-1}$ the zeros of $\mathcal T_p$ in its strong domain are
$x_m=-\omega^{2m-1}-\omega^{-(2m-1)}$, $m\ge1$. None is an omnific integer, and
\begin{equation}\label{hol:td:eq:rounding}
 \lfloor x_m\rfloor_{\Oz}=-\omega^{2m-1}-1,\qquad
 \lceil x_m\rceil_{\Oz}=-\omega^{2m-1}.
\end{equation}
So $\mathcal T_{\omega^{-1}}(x)=0$ has no omnific-integer solution in its strong
domain, although it has infinitely many simple real surreal solutions there.
\end{corollary}

\begin{proof}
The zeros are those of Theorem~\ref{hol:td:thm:zeros}. The last summand of
$x_m$ has negative exponent and nonzero coefficient, so $x_m\notin\Oz$. Since
$0<\omega^{-(2m-1)}<1$, $x_m$ lies strictly between the two consecutive
omnific integers of \eqref{hol:td:eq:rounding}, and there is no other zero.
\end{proof}

This concerns a strongly evaluated infinite-series equation, not a polynomial
Diophantine statement: the differential equation of $\mathcal T_p$ is
polynomial in its jets, but $\mathcal T_p(x)=0$ is not a polynomial equation in
$x$.

\paragraph{Rounding can magnify an infinitesimal error.}
Let $x_m^+=-\omega^{2m-1}$ and $x_m^-=x_m^+-1$, the ceiling and floor of
$x_m$. Then $x_m^+-x_m=\omega^{-(2m-1)}$ is infinitesimal and
$x_m^--x_m=-1+\omega^{-(2m-1)}$ is finite with leading coefficient $-1$, and
both satisfy the hypothesis of Proposition~\ref{hol:td:prop:local}. With the
valuation normalized by $v(\omega^{-1})=1$,
\begin{align}
 v\bigl(\mathcal T_p'(x_m)\bigr)&=-m^2+3m-1,&
 \lc\bigl(\mathcal T_p'(x_m)\bigr)&=(-1)^{m-1},\label{hol:td:eq:valder}\\
 v\bigl(\mathcal T_p(x_m^+)\bigr)&=-m^2+5m-2,&
 \lc\bigl(\mathcal T_p(x_m^+)\bigr)&=(-1)^{m-1},\label{hol:td:eq:valceil}\\
 v\bigl(\mathcal T_p(x_m^-)\bigr)&=-m^2+3m-1,&
 \lc\bigl(\mathcal T_p(x_m^-)\bigr)&=(-1)^m.\label{hol:td:eq:valfloor}
\end{align}
\begin{center}
\begin{tabular}{rrrr}
\toprule
$m$ & $v(\mathcal T_p'(x_m))$ & $v(\mathcal T_p(x_m^+))$ & $v(x_m^+-x_m)$\\
\midrule
1 & 1 & 2 & 1\\
2 & 1 & 4 & 3\\
3 & $-1$ & 4 & 5\\
4 & $-5$ & 2 & 7\\
5 & $-11$ & $-2$ & 9\\
6 & $-19$ & $-8$ & 11\\
\bottomrule
\end{tabular}
\end{center}
For every $m\ge5$, $\mathcal T_p(x_m^+)$ is infinite although $x_m^+$ differs
from a genuine zero by an infinitesimal. This violates no local continuity:
the derivative at these increasingly large roots is itself large. The exact
local formula records both scales; an argument based only on the small input
difference would lose the dominant effect.

\subsection{Examples}\label{hol:td:sub:examples}

\begin{example}[The simplest discrete field; source~13]\label{hol:td:ex:discrete}
For $\KK=\Q((t))$ and $p=t$, the series
$1+tz+t^4(z^2-2)+t^9(z^3-3z)+\cdots$ is strongly entire, nonpolynomial,
differentially algebraic of minimal order three and not D-finite. It needs no
uncountable constant field, no real closure and no divisible exponents.
\end{example}

\begin{example}[One group, two scales; source~13]\label{hol:td:ex:lex}
Let $\Gamma=\Z\oplus_{\mathrm{lex}}\Z$ with the first coordinate dominant.
The scale $\mu=(1,0)$ gives a strongly entire witness; $\mu=(0,1)$ does not:
its strong domain is the valuation ring for the projection to the first
coordinate, and $t^{(-1,0)}$ is excluded. The existence classification
quantifies over the scale; it does not say that every positive $\mu$ works.
\end{example}

\begin{example}[Countable cofinality without an order unit; source~13]
\label{hol:td:ex:cofinality}
Order $\Gamma=\bigoplus_{j\ge0}\Z e_j$ by the largest nonzero coordinate. The
sequence $(e_j)$ is cofinal, but every fixed $\mu>0$, and every ordinary
multiple of it, is dominated by some $e_N$. By
Theorem~\ref{hol:td:main:boundary}, or Theorem~\ref{hol:cf:main:entire}\,(a),
no nonpolynomial strongly entire series is differentially algebraic. This
does not exclude nonpolynomial strongly entire series; it says that any such
series has no algebraic differential relation over the full coefficient field.
The group is the integer analogue of \eqref{hol:eq:gammainfty}, over which the
series $F$ of \eqref{hol:eq:explicitF} is strongly entire.
\end{example}

\begin{corollary}[Stability under affine changes; source~13]
\label{hol:td:cor:affine}
Let $\mu$ be an order unit, $a\in\KK^\times$ and $b\in\KK$. Then
$g(z)=\mathcal T_p(az+b)$ is strongly entire and has minimal order three, and
its zeros are exactly the $(x_m-b)/a$, all simple.
\end{corollary}

\begin{proof}
For a fixed input, $v(ax+b)$ is bounded below by a multiple of $-\mu$, and the
quadratic support bound, expanded binomially, justifies the Taylor
coefficients of $g$ and their evaluation. By the regrouping argument of
Lemma~\ref{hol:td:lem:substitution}\,(b), $f\mapsto f(az+b)$ is a
$\KK$-algebra endomorphism of $\EK$ with inverse $f\mapsto f((z-b)/a)$, and
$D_z^j(f(az+b))=a^j(D_z^jf)(az+b)$. A nonzero affine substitution is an
automorphism of $\KK(z)$ and scales the successive derivatives by nonzero
powers of $a$, so it preserves the algebraic independence of the first three
jets and the order of an equation. The zeros follow by substitution.
\end{proof}

There are therefore many explicit entire nonlinear solutions, not one isolated
series. The argument classifies neither all differentially algebraic entire
series nor the possible minimal orders. (For $\Gamma=\R$ the minimal orders
are unbounded, by source~16's Theorem~\ref{hol:th:main:hierarchy}\,(b).) Order two, which source~13 left open
for other series, is excluded for every series by
Theorem~\ref{hol:ot:main:second} (source~14); the classification is
Question~\ref{hol:q:nonlinear}.

\subsection{Predecessors, repository comparison and priority}
\label{hol:td:sub:predecessors}

\paragraph{Literature.}
Source~13 consulted the NIST DLMF product and elliptic
sections~\cite{dlmftheta,dlmfelliptic}; the Tate-curve lecture notes hosted by
M.~Woodbury~\cite{woodbury}, whose \S16 (printed page~32, the Tate equation
and invariant differential) it inspected as a PDF image and whose
normalization agrees with the constants \eqref{hol:td:eq:constants}, the
hosting attribution not being a claim that Woodbury wrote every part; the
publisher abstract and bibliography of Sebbar's paper~\cite{sebbar}, whose
abstract explicitly discusses classical third-order theta differential
relations, without its full text; the publication record of Ohyama's
paper~\cite{ohyama}, not used in place of a proof; and publication records of
Mantova--Matusinski~\cite{mantovamatusinski} for normal forms. Several broad
searches returned irrelevant results and are not evidence of absence; no
exhaustive MathSciNet or zbMATH review, survey of book treatments or
citation-network review was made. In particular no claim is made that Sebbar
or Ohyama lacks a related minimal-order theorem.

\paragraph{Repository comparison at the pin.}
Source~13 read the repository at \texttt{4cdeaec} through a connector: the
root and documentation inventories, this report's README, its nonlinear corner
criterion and limitations, the vanishing-corner example, the first-order
limitation, its attribution and dependency sections, and the question and
status passage of Question~\ref{hol:q:nonlinear} (source lines 4135--4202 of
the article; its audit file gives the read range as about 4125--4290); and the
README of the Hahn--Tate report, not its full article. Its statements were
checked for this merge and are accurate at the pin: Question~19.1 asked for
existence with an order unit, for rationality of differentially algebraic
elements of $\operatorname{Frac}\E$, and whether the absence of an order unit
is necessary; its status settled the case without an order unit and gave no
counterexample with one; and the Hahn--Tate README contains the bilateral
theta domain on $U_q$. A repository content search that returned no matches
was followed by direct file reads and not used as proof of absence. This was a
targeted comparison, not an audit of every file, archived manuscript, branch or
commit. Since the pin the statements of this report that called the
order-unit case open have been updated, and Question~\ref{hol:q:nonlinear} has
been re-scoped, as recorded in Section~\ref{hol:sec:status}.

\paragraph{Priority.}
The triple product, the Chebyshev identity and the theta--elliptic identity
are classical, and no novelty is claimed for a third-order theta relation; the
counterexample is a classical function, not a new theta identity. The
no-order-unit direction is Theorem~\ref{hol:cf:main:entire}\,(a), and the
bilateral theta domain and uniformization belong to the Hahn--Tate report. The
proposed contribution of source~13 is the support-controlled descent to an
entire series, its exact domain at every scale, the existence classification
of Theorem~\ref{hol:td:main:boundary}, the minimal-order certificate with the
differential field and non-D-finiteness, and the stated surreal and omnific
consequences. Whether the arbitrary-rank formulation, the descent as an
entire-function counterexample or the three-stage minimal-order proof already
occurs in the literature needs a more exhaustive bibliographic review; the
answered question is a question of this report, not a named historical
conjecture.

\subsection{What the theta-descent part does not claim}
\label{hol:td:sub:nonclaims}

Every limitation stated by source~13, in its manuscript, README, audit files
or verification record, is kept.
\begin{enumerate}[label=\textup{(T\arabic*)}]
\item No Lean formalization, no independent referee report and no certified
bibliographic priority accompany source~13; its compiled and visually
inspected PDF is no evidence of mathematical correctness.
\item The classical theta, triple-product, Chebyshev and elliptic identities
are not new, and no novelty is claimed for third-order theta differential
relations; the counterexample is a classical function.
\item The no-order-unit direction is prior repository material
(Theorem~\ref{hol:cf:main:entire}\,(a)), re-proved by source~13 and not
claimed.
\item The bilateral theta domain and the Hahn--Tate uniformization are not
claimed anew.
\item $\mathcal T_p$ is not strongly entire on $\No$ or $\No[i]$; it fails at
$x=\omega^\omega$. Entireness is relative to one specified Hahn field, and an
extension of the value group can destroy it.
\item $D_z$ kills every Hahn coefficient, including $p=\omega^{-1}$, and is
not the Berarducci--Mantova derivation.
\item The finite checks do not verify strong summability of infinite or
arbitrary supports, transcendence, algebraic independence or minimal order,
the no-order-unit obstruction or the classification, completeness of the
infinite zero set, novelty, or correctness in general; they are not a formal
verification.
\item No classification of strongly entire differentially algebraic series is
given, and whether minimal order two occurs for some other series is open; a
lower-order equation for $\mathcal T_p$ itself is excluded, which is a
different assertion. (Source~14 has since excluded order two for every series,
Theorem~\ref{hol:ot:main:second}; the non-claim is kept as source~13's
statement, and the classification stays open.) The geometric route suggested by
Theorem~\ref{hol:td:thm:field} is a research direction, not a theorem.
\item ``No omnific zero'' concerns the strong domain only; it is not a
polynomial Diophantine statement, and no values are assigned outside the
domain.
\item The local valuation formulas are not uniform across the roots; a small
input error does not imply a small output error.
\item No transcendence theorem for numerical theta values is used or claimed;
the independence statements, and all fields of meromorphic functions in the
proofs, concern formal functions, not values at a point.
\item The minimal-order, differential-field and non-D-finiteness arguments
assume that $\mu$ is an order unit. The parameter is the monomial $p=t^\mu$;
no result is generalized to a parameter with an arbitrary higher-order tail.
\item Sebbar's paper and Ohyama's paper were not read in full, so no absence
claim is made about them; broad searches are not evidence of absence, and no
exhaustive MathSciNet or zbMATH review was made.
\item The repository comparison at \texttt{4cdeaec} was targeted, not a full
audit of every file, archived manuscript, branch or commit; the Hahn--Tate
report was compared through its README.
\item The answered question is a question of this report; no named historical
conjecture is claimed solved.
\item The existence classification quantifies over the scale: not every
positive $\mu$ gives an entire series (Example~\ref{hol:td:ex:lex}).
\item Without an order unit the conclusion excludes differentially algebraic
nonpolynomial entire series only, not nonpolynomial entire series
(Example~\ref{hol:td:ex:cofinality}).
\item The poles of the rational system at $z=\pm2$ are those of a coordinate
chart; the globally meaningful equation is the cleared one.
\item \emph{[merge]} Lemmas~\ref{hol:td:lem:substitution}
and~\ref{hol:td:lem:psi3}, Proposition~\ref{hol:td:prop:corner}, clause~(iv) of
Theorem~\ref{hol:cf:main:sharp}, Corollary~\ref{hol:td:cor:sharp} and the
consistency statements of Section~\ref{hol:td:sub:consistency} are the
merge's. The first two write out steps that source~13 used; no historical
priority is claimed for any of them.
\end{enumerate}

%======================================================================
% Order-three-threshold part: source 14. Labels use hol:ot:.
%======================================================================
\section{Second-order rigidity: finite endpoint spectra at exterior breaks}
\label{hol:ot:sec:setting}

This section and the next contain the material of source~14, delivered after
the theta-descent part had been merged. They answer the question left open by
source~13 and re-scoped in Section~\ref{hol:sec:status} after it: no
nonpolynomial strongly entire series satisfies a nonzero algebraic
differential equation of order at most two, for every value group and every
coefficient field of characteristic zero
(Theorem~\ref{hol:ot:main:second}\,(a)). With
Theorem~\ref{hol:td:main:boundary} this makes the least order exactly three
when $\Gamma$ has an order unit (Theorem~\ref{hol:ot:main:second}\,(b)). The
method is the finite initial polynomial of
Sections~\ref{hol:nl:sec:setting}--\ref{hol:nl:sec:consequences}, applied not
at one scale but at infinitely many scales where the initial polynomial is not
a monomial, and read through logarithmic Euler derivatives; it succeeds
exactly where the corner criterion of Theorem~\ref{hol:nl:thm:corner} is
silent, since an equation of order two may have $\chi_P=0$
(Proposition~\ref{hol:nl:prop:degenerate}). It is neither the escape chain,
nor the finite coefficient field, nor a difference-field argument.

\subsection{Source, conventions and the merge's choices}
\label{hol:ot:sub:conventions}

\begin{convention}[Setting of the order-three-threshold part]
\label{hol:ot:conv:field}
In Sections~\ref{hol:ot:sec:setting}--\ref{hol:ot:sec:theta}, $k$ is an
arbitrary field of characteristic zero with the trivial valuation, $\Gamma$ is
as in Convention~\ref{hol:conv:group}, $\KK=k((t^\Gamma))$, and $v$, $\lc$,
$\res$, strong summability, $\Dom(f)$, $\EK$, $\FE$, $D_z$, the coefficient
family algebra $\Sfam_1$ with its residue homomorphism $\rho$, and the Gauss
data $\gamma_\delta$, $\phi_\delta$, $N_\delta$ of \eqref{hol:nl:eq:gaussdata}
are as in Conventions~\ref{hol:nl:conv:field} and~\ref{hol:cf:conv:field} and
Section~\ref{hol:nl:sub:scaled}. Write $\vartheta=zD_z$, as in
Convention~\ref{hol:td:conv:field}, and $\vartheta_X=X\,\mathrm d/\mathrm dX$
on $k(X)$. Where a statement needs a divisible group it says so; it is then
applied to the divisible hull $\Gamma_\Q=\Gamma\otimes_\Z\Q$ and
$\KK_\Q=k((t^{\Gamma_\Q}))$ (Lemma~\ref{hol:ot:lem:extension}). $U,V,W,T$
and $G$ are indeterminates of polynomial rings over $k$ in this part, except
that in Corollary~\ref{hol:ot:cor:examples} $G$ is the series of
Corollary~\ref{hol:nl:cor:theta}; $E_0,\ldots,E_s$ are Euler-jet variables.
\end{convention}

For $k=\C$ one has $\KK=\K$ and $\EK=\E$. As for the three previous parts,
the generality in $k$ is recorded for this part only.

\paragraph{The source.}
Source~14 is \emph{The Order-Three Threshold for Entire Hahn--Surcomplex
Functions: Second-order rigidity, Newton-corner spectra, and sharp theta
examples} (26 pages, 23 September 2026). It consulted the repository at
\begin{center}\small\texttt{b895e8672990e8b5a56f97dd7246f9dd5f86c771}.\end{center}
At that pin the article and README of this report were exactly the five-source
text (sources~08--12); source~13's audit files, programs and records had been
placed in this directory but its manuscript had not yet been merged. Source~14
read this report's README there, and took source~13 from the author's own
file library as an unpublished manuscript, with its baseline \texttt{4cdeaec}
and its question \texttt{q:order2}, ``whether a different example could have
order two''. Its statements about this report (first-order nonlinear rigidity
and the all-order obstruction without an order unit) were accurate at the pin.
Since then source~13 has been merged as Sections~\ref{hol:td:sec:setting}--%
\ref{hol:td:sec:boundary}, and \texttt{q:order2} is its Question~11.1, merged
into Question~\ref{hol:q:nonlinear}; the question source~14 answers is
therefore the one recorded there, and source~14's description of source~13 as
a private project source that is not a published paper describes the state at
its pin. Source~14 assumes characteristic zero and an arbitrary ordered
abelian group; it passes to the divisible hull by a cofinal extension and
uses no divisibility, rank, algebraic or real closedness or order-unit
hypothesis in its main theorem. Its standing hypotheses are therefore those of
the coefficient-field and theta-descent parts, and every proof taken from it
was re-read against the statement it is attached to.

\paragraph{Renamed symbols.}
The following symbols of source~14 would collide with symbols of the earlier
parts. They are renamed here, in the new material only.
{\small
\begin{longtable}{@{}>{\raggedright\arraybackslash}p{0.27\textwidth}
>{\raggedright\arraybackslash}p{0.2\textwidth}
>{\raggedright\arraybackslash}p{0.45\textwidth}@{}}
\toprule
Source~14 & Here & Reason\\
\midrule
\endhead
$D=\mathrm d/\mathrm dz$, $f'$, $f''$ & $D_z$, $D_zf$, $D_z^2f$ & The
notation of Convention~\ref{hol:nl:conv:field}. (Source~14's $\KK=k((t^\Gamma))$,
$\EK$ and $\vartheta=zD_z$ are this report's.)\\
radius $r$; $\alpha_f(r)$; $\phi_r$; $I_f(r)$; $m_f(r)$, $N_f(r)$ & scale
$\delta$; $\gamma_\delta$; $\phi_\delta$; $\operatorname{Act}_\delta$;
$\underline N_\delta$, $N_\delta$ & The Gauss data \eqref{hol:nl:eq:gaussdata};
$r$ indexes derivative orders, and $I_P$ is the full corner polynomial.\\
Newton corner & break (Definition before
Lemma~\ref{hol:ot:lem:localfinite}) & A \emph{corner} is the finite term
selection of Section~\ref{hol:nl:sub:corner}.\\
restricted ring $\mathcal O\langle X\rangle_v$, $\red$ & $\Sfam_1^+$, $\rho$ &
Lemma~\ref{hol:nl:lem:algebra}; $\red_\Delta$ is a coarsened residue.\\
order bound $\ell$; blocks $P_{d,j}$; $\mu$, $M$, $C$ & $s$; $P_{d,w}$;
$\beta$, $n_1$, $C_1$ & $\ell$ is a dilation power; $\mu$ is a scale; $M(S)$
is a support monoid.\\
$H$ (exterior homogeneous polynomial) & $\mathfrak h_P$ & $H\in\Gamma$, and
$\mathcal H[f]$ is a theta expression.\\
$Q(U,V)$, $Q(U,V,W)$; $s$, $R$; $Q_s$ & $\mathfrak q$; $i_0$, $\mathfrak r$;
$\mathfrak q_{i_0}$ & $Q_s$ is a shift polynomial, $s$ an order, $R$ an order
bound.\\
$S_H$ & $\operatorname{Sp}(\mathfrak h)$ & $S$ is a support alphabet.\\
$u=\vartheta_X\phi/\phi$, $V_\phi$, $W_\phi$; $h=W/V$, $r$ & $\mathsf u_\phi$,
$\mathsf v_\phi$, $\mathsf w_\phi$; $\mathsf k$, $\mathsf r$ & $u$ is the
Laurent variable of Section~\ref{hol:td:sub:laurent}.\\
$\chi_H(T,G)$ & $\Upsilon_{\mathfrak h}(T,G)$ & $\chi_P$ is the residue corner
polynomial.\\
upper gap $g$, lower gap $e$ & $\mathsf g^+$, $\mathsf g^-$ & $g_2,g_3$ are
theta invariants, $e_n\in\Gamma$.\\
$h$, $q=t^h$, $Q=q^2$, $s_r(Q)$, $c$ & $\mu$, $p=t^\mu$, $p^2$,
$\mathcal Q_r(p^2)$, $\mathfrak c$ & As for source~13
(Section~\ref{hol:td:sub:conventions}).\\
$F_q$; $H_f$, $J_f$, $B_f$ & $\mathcal T_p$; $\mathcal H[f]$, $\mathcal K[f]$,
$\mathcal B[f]$ & \eqref{hol:td:eq:series}, \eqref{hol:td:eq:H},
\eqref{hol:td:eq:KB}.\\
$r_n=v(a_{n+1})-v(a_n)$ & $\delta_n$ & It is a break scale.\\
$F=g/p$; characteristic $p$ & $F_0=g/\mathsf p$; $\mathfrak p$ & $F$ is the
series \eqref{hol:eq:explicitF}; $p=t^\mu$.\\
differential field $L$; state variables $u,v$ & $\mathfrak L\subseteq\mathfrak
F$; $u_1,u_2$ & $L$ is a coefficient subfield, $M(S)$ a support monoid.\\
Stirling numbers (implicit) & $\varsigma_{j,i}$ & $s$ is an order.\\
$\Gamma_\Q$; $\No(i)$ & $\Gamma_\Q$; $\No[i]$ & The collection writes
$\No[i]$.\\
Theorems~1.1 and~1.3 & Theorem~\ref{hol:ot:main:second}\,(a) and~(b) &
Theorems~A--H of this report are different theorems.\\
\bottomrule
\end{longtable}}
The shipped audit file \path{14-order-three-threshold-SOURCE_AUDIT.md} keeps
the source's notation.

\paragraph{Words with two meanings.}
A \emph{break} is a scale $\delta$ at which the initial polynomial
$\phi_\delta$ is not a monomial; source~14 calls it a Newton corner. A
\emph{corner} keeps the meaning of Section~\ref{hol:nl:sub:corner}, and the
residue corner polynomial $\chi_P$ of a differential polynomial is not the
endpoint--gap polynomial $\Upsilon_{\mathfrak h}$ of a homogeneous one;
Remark~\ref{hol:ot:rem:corner} gives their exact relation. \emph{Exterior}
means, as elsewhere in this report, large $\delta$, that is arguments
$t^{-\delta}$ of large absolute valuation. The \emph{endpoint spectrum}
$\operatorname{Sp}(\mathfrak h)$ is a finite set of nonnegative integers, not
the spectrum of an operator. \emph{Order} and \emph{minimal order} are as in
Section~\ref{hol:td:sub:conventions}: the order of an equation is the highest
derivative in it. A \emph{gap} is a difference of exponents of $\phi$ at one
end. $\mathsf u_\phi,\mathsf v_\phi,\mathsf w_\phi$ are rational functions of
$X$ attached to a polynomial $\phi$, while $U,V,W$ are indeterminates.

\paragraph{Results printed once.}
Source~14 reproves several facts that the report already contains. They are
printed once and cited: the definitions of an order unit and of strong
entireness (Definitions~\ref{hol:def:orderunit} and~\ref{hol:def:entire});
its all-scale coefficient criterion (Lemma~2.2), which is
Proposition~\ref{hol:cf:prop:criterion}, proved there by the same two-scale
argument; the restricted-series ring and its reduction homomorphism, which are
$\Sfam_1^+$ and $\rho$ of Lemma~\ref{hol:nl:lem:algebra}, with the
commutation of Euler derivatives and reduction in
Lemma~\ref{hol:nl:lem:derivatives}; the first assertion of its Lemma~3.2,
which is Lemma~\ref{hol:nl:lem:active}; its Proposition~5.2, whose equation
$z^2(ff''-(f')^2)+zff'=0$ is \eqref{hol:nl:eq:degenerate} and whose
classification is Proposition~\ref{hol:nl:prop:degenerate}; its
Corollary~6.4 on Painlev\'e~I and~II, which is
Proposition~\ref{hol:nl:prop:painleve} (source~14's second route is recorded in
Remark~\ref{hol:ot:rem:painleve}); its Section~8, the theta example, whose
entireness, coefficients, cleared equation and derivation are
Theorem~\ref{hol:td:thm:domain}, Proposition~\ref{hol:td:prop:coeff},
Theorem~\ref{hol:td:thm:ode} and
Sections~\ref{hol:td:sub:elliptic}--\ref{hol:td:sub:formal}, except that its
minimality argument is a genuinely different route and is kept
(Remark~\ref{hol:ot:rem:secondroute}); its Section~10, the no-order-unit
direction, which is Theorem~\ref{hol:cf:main:entire}\,(a) in the Euler
presentation already recorded in Remark~\ref{hol:td:rem:euler} (its
Lemmas~10.1--10.2 and Theorem~10.3 are
Theorem~\ref{hol:cf:thm:coefficients}\,(a),
Lemmas~\ref{hol:cf:lem:valuerank}--\ref{hol:cf:lem:principal} and
Proposition~\ref{hol:cf:prop:exterior}); its Examples~10.4 and~10.5, which
are Example~\ref{hol:td:ex:cofinality} and, over $\R\oplus_{\mathrm{lex}}\R$
instead of $\Z\oplus_{\mathrm{lex}}\Z$, Example~\ref{hol:td:ex:lex}; the
surreal embedding \eqref{hol:eq:embedding}; the surreal witness, product,
zeros and omnific rounding of its Section~11.3, which are
Section~\ref{hol:td:sub:surreal}, Theorem~\ref{hol:td:thm:zeros} and
Corollary~\ref{hol:td:cor:omnific}; its Example~12.1 in positive
characteristic, which is Example~\ref{hol:nl:ex:charp} with the scale $1$ of
$\Q$ replaced by an order unit of $\Gamma$; and the $\Gamma=0$ remark, which
is in Section~\ref{hol:sec:fields}. Source~14 credits the theta material and
the no-order-unit direction as prior and claims neither.

\paragraph{Where each numbered result went.}
Theorem~1.1 is Theorem~\ref{hol:ot:main:second}\,(a), and Theorem~1.3 is
Theorem~\ref{hol:ot:main:second}\,(b); Lemma~2.4 is
Lemma~\ref{hol:ot:lem:extension}; Lemma~3.1 is
Lemma~\ref{hol:ot:lem:localfinite}; the second assertion of Lemma~3.2 is
Lemma~\ref{hol:ot:lem:breaks}; Lemma~4.1 and Remark~4.2 are
Lemma~\ref{hol:ot:lem:exterior} and Example~\ref{hol:ot:ex:degenerate};
Lemma~5.1 is Lemma~\ref{hol:ot:lem:endpoints};
Corollaries~6.1--6.3 are
Corollaries~\ref{hol:ot:cor:threejets}--\ref{hol:ot:cor:meromorphic};
Definition~7.1, Theorem~7.2 and Corollary~7.3 are
Definition~\ref{hol:ot:def:gaps}, Theorem~\ref{hol:ot:thm:gaps} and
Corollary~\ref{hol:ot:cor:gapcriterion}; Theorems~9.1 and~9.2 are
Theorems~\ref{hol:ot:thm:classification} and~\ref{hol:ot:thm:convexity};
Proposition~11.1 is Proposition~\ref{hol:ot:prop:fullclass}; Research
question~1 is merged into the re-scoped Question~\ref{hol:q:nonlinear}, and
Research questions~2--9 are
Questions~\ref{hol:ot:q:lift}--\ref{hol:ot:q:omnific}; Appendix~A is
\eqref{hol:td:eq:jetpoly} with \eqref{hol:ot:eq:eulersub}; Appendix~B and the
proof checkpoints of Section~14.1 are merged into the dependency table of
Section~\ref{hol:sec:verification}.

\paragraph{The merge's choices.}
The review of source~14 for this merge found its proofs correct. Three points
are written out here and marked as the merge's. First, the existence of a
late break (Lemma~\ref{hol:ot:lem:breaks}) is proved by naming the break
scale, the least of finitely many quotients, rather than by the phrase ``a
change of the minimum occurs at an equality''. Second, the exterior form
$\mathfrak h_P$ of Lemma~\ref{hol:ot:lem:exterior} is identified with the
residue corner of Section~\ref{hol:nl:sub:corner} written in Euler jets, and
its values on monomials are $\chi_P$
(Remark~\ref{hol:ot:rem:corner}); this is what connects source~14's endpoint
lemma to the corner criterion of source~10. Third, Corollary~\ref{hol:ot:cor:threejets}
and Corollary~\ref{hol:ot:cor:systems} are stated inside an explicit
differential field containing $\KK((z))$, which is how source~14's phrase
``every $f\in L\cap\mathcal E_\Gamma(k)$'' must be read, and
Corollary~\ref{hol:ot:cor:examples} applies Theorem~\ref{hol:ot:main:second}
to the series $G$ and $\Theta_p$ of earlier parts. Source~14's theorems,
proofs and non-claims are otherwise printed as delivered, in the renamed
notation.

\subsection{Breaks and the exterior form}\label{hol:ot:sub:breaks}

\begin{lemma}[Cofinal scalar extension; source~14]\label{hol:ot:lem:extension}
Let $\Gamma\subseteq\Lambda$ be an extension of ordered abelian groups with
$\Gamma$ cofinal in $\Lambda$. Every $f\in\EK$ is strongly entire over
$k((t^\Lambda))$. In particular this holds for $\Lambda=\Gamma_\Q$, in which
$\Gamma$ is cofinal.
\end{lemma}

\begin{proof}
Write $f=\sum_na_nz^n$ and use Proposition~\ref{hol:cf:prop:criterion}\,(iii)
over $\Lambda$. Given $\gamma,\beta\in\Lambda$, choose $\gamma'\in\Gamma$ with
$\gamma'\le\gamma$ and $\beta'\in\Gamma$ with $\beta'\ge\beta$; since $\Gamma$
is cofinal in the group $\Lambda$, it is also coinitial. If $n\ge0$ and
$v(a_n)+n\gamma\le\beta$, then $v(a_n)+n\gamma'\le\beta'$, which holds for only
finitely many $n$ with $a_n\ne0$. For $\Lambda=\Gamma_\Q$, the torsion-free
group $\Gamma$ embeds in $\Gamma\otimes_\Z\Q$, and an element $g/m$ with
$g\in\Gamma$ and $m\ge1$ satisfies $\abs{g/m}\le\abs g$.
\end{proof}

A series is a polynomial over $\KK_\Q$ exactly when it is one over $\KK$, and
a polynomial relation over $\KK(z)$ remains a nonzero relation over
$\KK_\Q(z)$. So a polynomiality conclusion proved over $\KK_\Q$ descends to
$\KK$. This is the only use of divisibility in the part.

Let $f=\sum_na_nz^n\in\EK$ be nonzero. For every $\delta\in\Gamma$ the Gauss
value $\gamma_\delta$ and the initial polynomial $\phi_\delta$ of
\eqref{hol:nl:eq:gaussdata} exist. Put
\begin{equation}\label{hol:ot:eq:active}
 \operatorname{Act}_\delta=\{n:a_n\ne0,\ v(a_n)-n\delta=\gamma_\delta\},\qquad
 \underline N_\delta=\min\operatorname{Act}_\delta,\qquad
 N_\delta=\max\operatorname{Act}_\delta ,
\end{equation}
so that $\underline N_\delta=\operatorname{ord}_X\phi_\delta$, the $X$-adic order, and
$N_\delta=\deg\phi_\delta$. A scale $\delta$ is a \emph{break} of $f$ if
$\underline N_\delta<N_\delta$, that is, if $\phi_\delta$ is not a monomial.
Lemma~\ref{hol:nl:lem:algebra} gives
$\rho\bigl(t^{-\gamma_\delta}(\vartheta^jf)(t^{-\delta}X)\bigr)
=\vartheta_X^j\phi_\delta$ for every $j$, because $\vartheta^jf$ scales like
$f$ with each coefficient multiplied by the integer $n^j$, and Euler
operators commute with $\rho$. A derivative on the right may vanish; no
equality of individual valuations is asserted.

\begin{lemma}[Local finiteness of active indices; source~14]
\label{hol:ot:lem:localfinite}
Let $f\in\EK$ be nonzero and $\delta_1\le\delta_2$ in $\Gamma$. Only finitely
many indices lie in $\operatorname{Act}_\delta$ for some
$\delta\in[\delta_1,\delta_2]$. If $\delta<\delta'$, $n\in\operatorname{Act}_\delta$
and $n'\in\operatorname{Act}_{\delta'}$, then $n\le n'$.
\end{lemma}

\begin{proof}
Fix $a_j\ne0$. On the interval its weighted value $v(a_j)-j\delta$ is at most
$C=v(a_j)-j\delta_1$. By Proposition~\ref{hol:cf:prop:criterion}\,(iii) at
$\gamma=-\delta_2$ all but finitely many nonzero $a_n$ have
$v(a_n)-n\delta_2>C$, and then $v(a_n)-n\delta\ge v(a_n)-n\delta_2>C$ for
$\delta\le\delta_2$, so such $n$ is never active. For the monotonicity, the
two minimality inequalities $v(a_n)-n\delta\le v(a_{n'})-n'\delta$ and
$v(a_{n'})-n'\delta'\le v(a_n)-n\delta'$ add to
$(n'-n)(\delta'-\delta)\ge0$.
\end{proof}

\begin{lemma}[Late breaks; source~14]\label{hol:ot:lem:breaks}
Let $\Gamma$ be divisible and $f\in\EK$ nonpolynomial. For every $\delta_*\in\Gamma$
and every $B\in\N$ there is a break $\delta\ge\delta_*$ of $f$ with
$\underline N_\delta>B$.
\end{lemma}

\begin{proof}
By Lemma~\ref{hol:nl:lem:active} there is $\delta_B$ such that every active
index at every $\delta\ge\delta_B$ exceeds $B$. Let
$\delta_1=\max(\delta_*,\delta_B)$. If $\delta_1$ is a break we are done;
otherwise $\operatorname{Act}_{\delta_1}=\{N\}$. Lemma~\ref{hol:nl:lem:active}
again gives $\delta_2>\delta_1$ at which every active index exceeds $N$. By
Lemma~\ref{hol:ot:lem:localfinite} the set $F$ of indices active somewhere on
$[\delta_1,\delta_2]$ is finite, and on that interval $\gamma_\delta$ is the
minimum of the finitely many affine functions $v(a_n)-n\delta$, $n\in F$. For
$n\in F$ with $n>N$ put
\[
 \delta^{(n)}=\frac{v(a_n)-v(a_N)}{n-N}\in\Gamma,
\]
using divisibility; then $v(a_N)-N\delta\le v(a_n)-n\delta$ exactly when
$\delta\le\delta^{(n)}$. Since $N$ is active at $\delta_1$, every
$\delta^{(n)}\ge\delta_1$, and since some $n>N$ in $F$ is active at $\delta_2$
while $N$ is not, $\delta^{(n)}<\delta_2$ for that $n$. Let $\delta^\star$ be
the least $\delta^{(n)}$ (the merge names this scale). For $n\in F$ with
$n<N$, activity of $N$ at $\delta_1$ and $\delta^\star\ge\delta_1$ give
$v(a_N)-N\delta^\star\le v(a_n)-n\delta^\star$; for $n>N$ the same holds
because $\delta^\star\le\delta^{(n)}$, with equality for an index attaining
$\delta^\star$. Since the least weighted value on the interval is attained in
$F$, $N$ and that index are both active at $\delta^\star$, which is a break,
and $\delta^\star\ge\delta_1\ge\delta_B$ makes every active index there exceed
$B$.
\end{proof}

The proof uses only finite lower envelopes on bounded intervals: no
completeness of $\Gamma$, no real-valued norm, no countable cofinal sequence
and no order unit.

Let $P\in\KK[z,Y_0,\ldots,Y_s]$ be a nonzero differential polynomial of order
at most $s$. Since $z^jD_z^j=\vartheta(\vartheta-1)\cdots(\vartheta-j+1)$, the
substitution
\begin{equation}\label{hol:ot:eq:stirling}
 Y_j\longmapsto z^{-j}\sum_{i=0}^j\varsigma_{j,i}E_i\qquad(0\le j\le s),
\end{equation}
with $\varsigma_{j,i}$ the signed Stirling numbers of the first kind, so that
$\sum_i\varsigma_{j,i}T^i=\fall Tj$, rewrites $P(z,f,D_zf,\ldots)$ as a Laurent
polynomial in $z$ in the Euler jets $E_i=\vartheta^if$. After multiplication by
a power of $z$ it becomes $\widetilde P(z,f,\vartheta f,\ldots,\vartheta^sf)$
with $0\ne\widetilde P\in\KK[z,E_0,\ldots,E_s]$; the change is invertible and
preserves total jet degree. Decompose
\[
 \widetilde P=\sum_{d,w}z^wP_{d,w}(E_0,\ldots,E_s),
\]
with $P_{d,w}$ homogeneous of total degree $d$ in the jets, choose the largest
$d$ occurring, then the largest $w$ with $P_{d,w}\ne0$, let $\beta$ be the
least valuation of a coefficient of $P_{d,w}$, and put
\begin{equation}\label{hol:ot:eq:exterior}
 \mathfrak h_P=\res\bigl(t^{-\beta}P_{d,w}\bigr)\in k[Y_0,\ldots,Y_s],
\end{equation}
applying $\res$ to coefficients and renaming $E_i$ as $Y_i$. It is a nonzero
homogeneous polynomial, the \emph{exterior form} of $P$.

\begin{lemma}[Exterior homogeneous reduction; source~14]
\label{hol:ot:lem:exterior}
Let $f\in\EK$ be nonpolynomial with $P(z,f,D_zf,\ldots,D_z^sf)=0$ for a nonzero
$P$ of order at most $s$. Then $\mathfrak h_P$ has positive total degree, and
there is $\delta_0\in\Gamma_{>0}$ such that for every $\delta\ge\delta_0$
\[
 \mathfrak h_P\bigl(\phi_\delta,\vartheta_X\phi_\delta,\ldots,
 \vartheta_X^s\phi_\delta\bigr)=0\quad\text{in }k[X].
\]
No divisibility of $\Gamma$ is used.
\end{lemma}

\begin{proof}
If $d=0$, then $\widetilde P$ is a nonzero element of $\KK[z]$, which cannot
vanish as a series; so $d\ge1$. Because $f$ is nonpolynomial there is $n_1$ with
$a_{n_1}\ne0$ and $(d-d')n_1+w-w'>0$ for every block $P_{d',w'}\ne0$ with
$d'<d$; put $C_1=v(a_{n_1})$, so that $\gamma_\delta\le C_1-n_1\delta$ for
every $\delta$.

Evaluate at $z=t^{-\delta}X$ and normalize by $t^{\beta+d\gamma_\delta-w\delta}$.
Each factor $(\vartheta^if)(t^{-\delta}X)$ equals $t^{\gamma_\delta}$ times an
element of $\Sfam_1^+$, so a monomial of a block $z^{w'}P_{d',w'}$ with
coefficient $b\ne0$ becomes $t^{v(b)-w'\delta+d'\gamma_\delta}$ times an element
of $\Sfam_1^+$. Relative to the normalization its exponent is at least
\begin{equation}\label{hol:ot:eq:gapbound}
 v(b)-\beta-(d-d')\gamma_\delta+(w-w')\delta
 \ \ge\ v(b)-\beta-(d-d')C_1+\bigl((d-d')n_1+w-w'\bigr)\delta .
\end{equation}
For $d'<d$ the coefficient of $\delta$ is a positive integer by the choice of
$n_1$; for $d'=d$ and $(d',w')\ne(d,w)$ it is $w-w'>0$. There are finitely many
monomials, so the right side is positive for all $\delta\ge\delta_0$ once
$\delta_0>0$ exceeds the absolute values of the finitely many constants; no
division occurs. The monomials of $P_{d,w}$ itself with coefficients of value
above $\beta$ also lose their residue. Applying the ring homomorphism $\rho$ of
Lemma~\ref{hol:nl:lem:algebra} to the normalized identity $0=\widetilde P(\ldots)$
leaves $X^w\mathfrak h_P(\phi_\delta,\vartheta_X\phi_\delta,\ldots)=0$, and
$X^w$ cancels in $k[X]$.
\end{proof}

\begin{remark}[The exterior form is the residue corner in Euler jets; merge]
\label{hol:ot:rem:corner}
Write $P$ in the combined-monomial form \eqref{hol:nl:eq:P}. By
\eqref{hol:ot:eq:stirling} the monomial
$c_\nu z^{k_\nu}\prod_jY_j^{m_{\nu j}}$ becomes
$c_\nu z^{w_\nu}\prod_j\bigl(\sum_i\varsigma_{j,i}E_i\bigr)^{m_{\nu j}}$, with the
weight $w_\nu$ of \eqref{hol:nl:eq:dw}. The substitution
$Y_j\mapsto\sum_i\varsigma_{j,i}Y_i$ is a unitriangular, hence invertible, linear
change of variables over $\Z$ and preserves total degree, so the selection of
Lemma~\ref{hol:ot:lem:exterior} picks $d=d_P$ and the $z$-power of weight $w_P$:
its block is the image of $\sum_{\nu\in\mathcal C_P}c_\nu\prod_jY_j^{m_{\nu j}}$.
Its least coefficient value is $\beta_P$, because the image of the nonzero
residue vector $(\bar c_\nu)$ under the invertible integral change is nonzero.
Hence
\[
 \mathfrak h_P(Y_0,\ldots,Y_s)=\sum_{\nu\in\mathcal C_P}\bar c_\nu
 \prod_{j=0}^s\Bigl(\sum_{i=0}^j\varsigma_{j,i}Y_i\Bigr)^{m_{\nu j}},\qquad
 \mathfrak h_P(1,T,T^2,\ldots,T^s)=\chi_P(T),
\]
the second because $\sum_i\varsigma_{j,i}T^i=\fall Tj$. On a monomial $\phi=X^N$ one
has $\vartheta_X^iX^N=N^iX^N$, so
$\mathfrak h_P(X^N,\vartheta_XX^N,\ldots)=\chi_P(N)X^{d_PN}$: the monomial test
of the exterior form is the residue corner polynomial, and the corner
criterion of Theorem~\ref{hol:nl:thm:corner} uses exactly this test. When
$\chi_P=0$ the exterior form vanishes on every monomial but is still a
nonzero polynomial, and source~14 keeps the whole initial polynomial at
breaks. For the equation \eqref{hol:nl:eq:degenerate}, $\mathfrak
h_P=Y_0Y_2-Y_1^2$; for the cleared theta equation \eqref{hol:td:eq:cleared},
$\beta_P=0$ and the dehomogenized exterior form is
\eqref{hol:ot:eq:thetainitial} below, consistent with
Proposition~\ref{hol:td:prop:corner}.
\end{remark}

\begin{example}[The cancellation that must not be ignored; source~14]
\label{hol:ot:ex:degenerate}
For $\mathfrak h=Y_0Y_2-Y_1^2$ one has $\mathfrak h(X^N,\vartheta_XX^N,
\vartheta_X^2X^N)=0$ for every $N$, so substituting a trial monomial and
assuming that the resulting polynomial in $N$ is nonzero proves nothing. This
$\mathfrak h$ is the exterior form of \eqref{hol:nl:eq:degenerate}, whose
formal solutions are exactly $0$ and the monomials $cz^n$
(Proposition~\ref{hol:nl:prop:degenerate}, which is source~14's
Proposition~5.2). Lemma~\ref{hol:ot:lem:endpoints} below gives
$\operatorname{Sp}(\mathfrak h)=\varnothing$ for it: no nonmonomial initial
polynomial is permitted at any exterior break, which is stronger than the
identically vanishing monomial test and consistent with unbounded monomial
solutions. Polynomiality is therefore not a uniform degree bound: a fixed
equation can have polynomial solutions of every degree.
\end{example}

\subsection{The endpoint lemma and the proof of Theorem~\ref{hol:ot:main:second}\,(a)}
\label{hol:ot:sub:endpoint}

For a nonzero $\phi\in k[X]$ put $\mathsf u_\phi=\vartheta_X\phi/\phi$,
$\mathsf v_\phi=\vartheta_X\mathsf u_\phi$ and
$\mathsf w_\phi=\vartheta_X\mathsf v_\phi$ in $k(X)$. Then
\[
 \frac{\vartheta_X^2\phi}\phi=\mathsf u_\phi^2+\mathsf v_\phi,\qquad
 \frac{\vartheta_X^3\phi}\phi=\mathsf u_\phi^3+3\mathsf u_\phi\mathsf v_\phi
 +\mathsf w_\phi .
\]

\begin{lemma}[Finite endpoint spectrum; source~14]\label{hol:ot:lem:endpoints}
Let $0\ne\mathfrak h\in k[Y_0,Y_1,Y_2]$ be homogeneous. There is a finite set
$\operatorname{Sp}(\mathfrak h)\subseteq\N$ such that every nonmonomial
$\phi\in k[X]$ with $\mathfrak h(\phi,\vartheta_X\phi,\vartheta_X^2\phi)=0$ has
both its $X$-adic order $\operatorname{ord}_X\phi$ and $\deg\phi$ in
$\operatorname{Sp}(\mathfrak h)$.
\end{lemma}

\begin{proof}
Put
\begin{equation}\label{hol:ot:eq:dehom}
 \mathfrak q(U,V)=\mathfrak h(1,U,U^2+V).
\end{equation}
It is nonzero: setting $Y_0=1$ is injective on homogeneous polynomials of a
fixed degree, since the missing power of $Y_0$ is determined by the other
exponents, and $Y_2\mapsto U^2+V$ is a triangular automorphism. Let $V^{i_0}$
be the largest power of $V$ dividing $\mathfrak q$ and write
\begin{equation}\label{hol:ot:eq:factor}
 \mathfrak q=V^{i_0}\mathfrak r,\qquad \mathfrak r(U,0)\ne0 .
\end{equation}
By homogeneity $\mathfrak q(\mathsf u_\phi,\mathsf v_\phi)=0$. The function
$\mathsf u_\phi$ is not constant: $\vartheta_X\phi=b\phi$ with $b\in k$ forces
$n=b$ for every exponent $n$ of $\phi$, and in characteristic zero there is at
most one such exponent. The constants of $\mathrm d/\mathrm dX$ in $k(X)$ are
$k$, so $\mathsf v_\phi\ne0$, and cancelling $\mathsf v_\phi^{i_0}$ in the field
$k(X)$ gives $\mathfrak r(\mathsf u_\phi,\mathsf v_\phi)=0$. With
$m=\operatorname{ord}_X\phi$ and $N=\deg\phi$, the expansions at $0$ and at $\infty$ are
\[
 \mathsf u_\phi=m+O(X),\quad\mathsf v_\phi=O(X);\qquad
 \mathsf u_\phi=N+O(X^{-1}),\quad\mathsf v_\phi=O(X^{-1}),
\]
and their constant terms give $\mathfrak r(m,0)=\mathfrak r(N,0)=0$. The
nonzero polynomial $\mathfrak r(T,0)$ has finitely many roots in $k$, and
distinct nonnegative integers are distinct in $k$; let
$\operatorname{Sp}(\mathfrak h)$ be its set of roots in $\N$.
\end{proof}

\begin{proof}[Proof of Theorem~\ref{hol:ot:main:second}\,\textup{(a)}]
Let $f\in\EK$ satisfy a nonzero relation of order at most two over $\KK(z)$;
clearing denominators gives a nonzero $P\in\KK[z,Y_0,Y_1,Y_2]$. Suppose $f$ is
nonpolynomial. By Lemma~\ref{hol:ot:lem:extension}, $f$ is strongly entire
over $\KK_\Q$ and satisfies the same relation there, so we may assume $\Gamma$
divisible. Lemma~\ref{hol:ot:lem:exterior} gives $\mathfrak h_P\ne0$ and
$\delta_0$. Choose $B\in\N$ larger than every element of
$\operatorname{Sp}(\mathfrak h_P)$ ($B=0$ if it is empty).
Lemma~\ref{hol:ot:lem:breaks} gives a break $\delta\ge\delta_0$ with
$\underline N_\delta>B$; its initial polynomial is nonmonomial, satisfies
$\mathfrak h_P(\phi_\delta,\vartheta_X\phi_\delta,\vartheta_X^2\phi_\delta)=0$,
and has $\operatorname{ord}_X\phi_\delta=\underline N_\delta\notin
\operatorname{Sp}(\mathfrak h_P)$, contradicting
Lemma~\ref{hol:ot:lem:endpoints}. So $f$ is a polynomial over $\KK_\Q$, hence
over $\KK$. The equivalent formulation is the same statement for the
algebraic dependence of $f,D_zf,D_z^2f$ over $\KK(z)$.
\end{proof}

The proof uses characteristic zero twice: nonzero integer multipliers keep
valuation zero under differentiation, and infinitely many degree indices stay
distinct in the residue field. The finite endpoint set is what the corner
criterion lacks when $\chi_P=0$: at order two the degenerate part of an
equation leaves only the finitely many roots of $\mathfrak r(T,0)$ for both
endpoints of every nonmonomial initial polynomial, while a nonpolynomial
entire series has breaks with arbitrarily large endpoints.

\subsection{Consequences and the proof of Theorem~\ref{hol:ot:main:second}\,(b)}
\label{hol:ot:sub:consequences}

\begin{corollary}[Three-jet transcendence; source~14]
\label{hol:ot:cor:threejets}
For nonpolynomial $f\in\EK$,
\[
 \trdeg_{\KK(z)}\KK(z)(f,D_zf,D_z^2f)=3 .
\]
More generally, let $\mathfrak F$ be a
differential field extension of $\KK((z))$ whose derivation extends $D_z$, and
$\mathfrak L\subseteq\mathfrak F$ a differential subfield containing $\KK(z)$ with
$\trdeg_{\KK(z)}\mathfrak L\le2$. Every $f\in\mathfrak L\cap\EK$ is a polynomial.
\end{corollary}

\begin{proof}
The first assertion is Theorem~\ref{hol:ot:main:second}\,(a). For the second,
$f,D_zf,D_z^2f\in\mathfrak L$ are algebraically dependent over $\KK(z)$ in $\mathfrak F$, and a
polynomial relation among elements of $\KK((z))$ holds in $\mathfrak F$ exactly when it
holds in $\KK((z))$.
\end{proof}

\begin{corollary}[Rational two-state systems; source~14]
\label{hol:ot:cor:systems}
Let $\mathfrak F$ be as in Corollary~\ref{hol:ot:cor:threejets} and let $u_1,u_2\in\mathfrak F$
satisfy $D_zu_1=R_1(z,u_1,u_2)$ and $D_zu_2=R_2(z,u_1,u_2)$ with
$R_1,R_2\in\KK(z,Y_1,Y_2)$, all denominators nonzero after substitution.
Every element of $\KK(z,u_1,u_2)$ that lies in $\EK$ is a polynomial.
\end{corollary}

\begin{proof}
The system makes $\KK(z,u_1,u_2)$ a differential field, by the quotient rule,
of transcendence degree at most two over $\KK(z)$; apply
Corollary~\ref{hol:ot:cor:threejets}.
\end{proof}

The state variables themselves need not be entire; the corollary restricts
the strongly entire observables of the field. A nonpolynomial strongly entire
function therefore cannot be generated by a rational nonautonomous system with
two state variables over the full Hahn coefficient field, for every $\Gamma$.
Corollary~\ref{hol:cf:cor:systems} (source~11) allows any number of states and
quotients of entire series but needs the absence of an order unit, and
Corollary~\ref{hol:cor:systems} is the linear case. With an order unit,
$\mathcal T_p$ is an observable of the rational system with the four states
$f,\mathcal L,\mathcal A,\mathcal A_1$ of Theorem~\ref{hol:td:thm:field}, so
there the number two cannot be replaced by four; three states are not decided
here (an observation of the merge). Arbitrary
transcendental right-hand sides remove the finite differential-field
hypothesis and are outside the corollary.

\begin{corollary}[Polynomial-denominator meromorphic rigidity; source~14]
\label{hol:ot:cor:meromorphic}
Let $g\in\EK$, $0\ne\mathsf p\in\KK[z]$ and $F_0=g/\mathsf p\in\KK((z))$. If $F_0$
satisfies a nonzero algebraic differential equation of order at most two over
$\KK(z)$, then $F_0\in\KK(z)$.
\end{corollary}

\begin{proof}
The relations $g=\mathsf pF_0$, $D_zg=D_z\mathsf p\,F_0+\mathsf pD_zF_0$ and
$D_z^2g=D_z^2\mathsf p\,F_0+2D_z\mathsf p\,D_zF_0+\mathsf pD_z^2F_0$ form an invertible
triangular linear change of jets over $\KK(z)$, so $g$ satisfies a nonzero
relation of order at most two, and Theorem~\ref{hol:ot:main:second}\,(a)
makes it a polynomial.
\end{proof}

This concerns quotients with a polynomial denominator only; it does not
assert that every object called meromorphic with finitely many poles in a
Hahn analytic category has such a representation. For the fraction field $\FE$
without an order unit, Theorem~\ref{hol:cf:main:entire}\,(c) gives rationality
in every order.

\begin{remark}[Painlev\'e~I and~II: a second route; source~14]
\label{hol:ot:rem:painleve}
Source~14's Corollary~6.4 is Proposition~\ref{hol:nl:prop:painleve}, with the
same degree comparison for polynomials. Source~10 reaches polynomiality
through the nonzero corners $\chi_P=-6$ and $\chi_P=-2$; source~14 reaches it
through Theorem~\ref{hol:ot:main:second}\,(a), since both equations have order
two. The equations of DLMF~\S32.2~\cite{dlmfpainleve} are read as formal
equations over $\KK$; nothing is said about local formal solutions, which
exist, or about classical complex Painlev\'e transcendents.
\end{remark}

\begin{corollary}[The explicit series of earlier parts; merge]
\label{hol:ot:cor:examples}
The series $G=\sum_{n\ge0}t^{n^2}z^n$ over $\C((t^\Q))$ of
Corollary~\ref{hol:nl:cor:theta}, and, for an order unit $\mu$ and $p=t^\mu$,
the partial theta series $\Theta_p\in\EK$ of
Proposition~\ref{hol:cf:prop:theta}, have three algebraically independent jets
$f,D_zf,D_z^2f$ over the respective $\KK(z)$.
\end{corollary}

\begin{proof}
Both are nonpolynomial and strongly entire (Corollary~\ref{hol:nl:cor:theta},
Proposition~\ref{hol:cf:prop:theta}); apply
Theorem~\ref{hol:ot:main:second}\,(a).
\end{proof}

Whether $G$ or $\Theta_p$ satisfies an algebraic differential equation of
order at least three is not decided; the corollary only raises the known
lower bound from two jets (Corollary~\ref{hol:nl:cor:theta}) to three.

\begin{proof}[Proof of Theorem~\ref{hol:ot:main:second}\,\textup{(b)}]
Let $\mu$ be an order unit and $p=t^\mu$. By
Theorem~\ref{hol:td:thm:domain} and Proposition~\ref{hol:td:prop:coeff},
$\mathcal T_p$ is a nonpolynomial element of $\EK$, and by
Theorem~\ref{hol:td:thm:ode} it satisfies a nonzero algebraic differential
equation of order three. By part~(a) no nonpolynomial element of $\EK$
satisfies one of order at most two; an equation of order zero is included,
being of order at most two. So the least order is three. Without an order
unit, Theorem~\ref{hol:cf:main:entire}\,(a) excludes every finite order.
\end{proof}

When $\Gamma=0$ every strongly entire series is a polynomial, by strong
summability at $x=1$ (Section~\ref{hol:sec:fields}); source~14 includes this
case in its statement, and Convention~\ref{hol:conv:group} excludes it here.

\begin{remark}[A second route to minimal order three for $\mathcal T_p$;
source~14]\label{hol:ot:rem:secondroute}
The lower bound of Corollary~\ref{hol:td:cor:order}, that
$\mathcal T_p,D_z\mathcal T_p,D_z^2\mathcal T_p$ are algebraically independent
over $\KK(z)$ when $\mu$ is an order unit, is also a case of
Theorem~\ref{hol:ot:main:second}\,(a), since $\mathcal T_p$ is nonpolynomial
and strongly entire. Source~14 proves minimality this way, ``without the
separate difference-field minimality argument''. The two routes are genuinely
different: source~13's argument uses the Laurent ring, the substitution
homomorphism and three difference-field obstructions for $u\mapsto p^2u$, and
also yields the degree-two statement of Theorem~\ref{hol:td:thm:field} and
non-D-finiteness (Proposition~\ref{hol:td:prop:nondfinite}), which
Theorem~\ref{hol:ot:main:second} does not give; source~14's argument uses no
property of $\mathcal T_p$ beyond nonpolynomial strong entireness. Both are
kept. The strict inequality of orders also shows again that
$\mathcal T_p$ satisfies no equation to which
Theorem~\ref{hol:nl:thm:certificate} applies (Remark~\ref{hol:td:rem:corner}),
but not conversely.
\end{remark}

\section{Order three: endpoint gaps, theta breaks and scope}
\label{hol:ot:sec:theta}

\subsection{The endpoint--gap polynomial}\label{hol:ot:sub:gaps}

At order three the endpoint lemma has a precise analogue with one more
integer parameter; this is what an equation of order three can use, and the
theta equation uses it.

\begin{definition}[Endpoint gaps; source~14]\label{hol:ot:def:gaps}
For a nonmonomial $\phi\in k[X]$ with highest exponent $N$ and second-highest
exponent $N-\mathsf g^+$, the positive integer $\mathsf g^+$ is its
\emph{upper gap}; with lowest exponent $m$ and second-lowest exponent
$m+\mathsf g^-$, the positive integer $\mathsf g^-$ is its \emph{lower gap}.
The pairs $(N,\mathsf g^+)$ and $(m,\mathsf g^-)$ are its endpoint--gap
pairs.
\end{definition}

\begin{theorem}[Third-order endpoint--gap obstruction; source~14]
\label{hol:ot:thm:gaps}
Let $0\ne\mathfrak h\in k[Y_0,Y_1,Y_2,Y_3]$ be homogeneous and
\begin{equation}\label{hol:ot:eq:dehom3}
 \mathfrak q(U,V,W)=\mathfrak h(1,U,U^2+V,U^3+3UV+W).
\end{equation}
Let $i_0$ be the least total degree in $(V,W)$ of a nonzero term of
$\mathfrak q$ and $\mathfrak q_{i_0}$ the sum of the terms of that degree. Then
\begin{equation}\label{hol:ot:eq:upsilon}
 \Upsilon_{\mathfrak h}(T,G)=\mathfrak q_{i_0}(T,1,-G)\in k[T,G]
\end{equation}
is nonzero, and every nonmonomial $\phi\in k[X]$ with
$\mathfrak h(\phi,\vartheta_X\phi,\vartheta_X^2\phi,\vartheta_X^3\phi)=0$
satisfies $\Upsilon_{\mathfrak h}(N,\mathsf g^+)=0$ and
$\Upsilon_{\mathfrak h}(m,-\mathsf g^-)=0$ for its endpoint--gap pairs.
\end{theorem}

\begin{proof}
As in Lemma~\ref{hol:ot:lem:endpoints}, $\mathfrak q\ne0$. The polynomial
$\mathfrak q_{i_0}$ is homogeneous of degree $i_0$ in $(V,W)$, and $V=1$,
$W=-G$ is injective on such polynomials, so $\Upsilon_{\mathfrak h}\ne0$. By
homogeneity and the identities before Lemma~\ref{hol:ot:lem:endpoints},
$\mathfrak q(\mathsf u_\phi,\mathsf v_\phi,\mathsf w_\phi)=0$. At infinity
write $\phi=aX^N(1+bX^{-\mathsf g^+}+\cdots)$ with $ab\ne0$; then, in
$k((X^{-1}))$,
\[
 \mathsf u_\phi=N-\mathsf g^+bX^{-\mathsf g^+}+\cdots,\qquad
 \mathsf v_\phi=(\mathsf g^+)^2bX^{-\mathsf g^+}+\cdots,\qquad
 \mathsf w_\phi/\mathsf v_\phi=-\mathsf g^++O(X^{-1}).
\]
Writing each homogeneous part as
$\mathfrak q_i(\mathsf u_\phi,\mathsf v_\phi,\mathsf w_\phi)=\mathsf v_\phi^i
\mathfrak q_i(\mathsf u_\phi,1,\mathsf w_\phi/\mathsf v_\phi)$, the parts with
$i>i_0$ have $X^{-1}$-order above $i_0\mathsf g^+$, and the coefficient of
$X^{-i_0\mathsf g^+}$ is $\bigl((\mathsf g^+)^2b\bigr)^{i_0}
\mathfrak q_{i_0}(N,1,-\mathsf g^+)$; replacing $\mathsf u_\phi$ by $N$ changes
only higher orders. So $\Upsilon_{\mathfrak h}(N,\mathsf g^+)=0$. At $0$
write $\phi=aX^m(1+bX^{\mathsf g^-}+\cdots)$: then $\mathsf u_\phi\to m$,
$\mathsf v_\phi$ has order $\mathsf g^-$ and
$\mathsf w_\phi/\mathsf v_\phi\to\mathsf g^-$, and the same argument gives
$\mathfrak q_{i_0}(m,1,\mathsf g^-)=\Upsilon_{\mathfrak h}(m,-\mathsf g^-)=0$.
\end{proof}

At order two the lowest transverse part is $A(U)V^{i_0}$, whose
specialization is $A(T)$: no gap variable remains, which is
Lemma~\ref{hol:ot:lem:endpoints}. At order three a polynomial such as $G^2-1$
permits arbitrarily large endpoints while fixing the gap to one; this is what
happens for the theta equation (Section~\ref{hol:ot:sub:thetaform}).

\begin{corollary}[A third-order rigidity certificate; source~14]
\label{hol:ot:cor:gapcriterion}
Let $P$ have order at most three, with exterior form $\mathfrak h_P$ as in
Lemma~\ref{hol:ot:lem:exterior}. If the set
\[
 \{N\in\N:\ \Upsilon_{\mathfrak h_P}(N,\mathsf g)=0\text{ for some integer }
 1\le\mathsf g\le N\}
\]
is bounded, every strongly entire solution of $P[f]=0$ is a polynomial.
\end{corollary}

\begin{proof}
Pass to $\Gamma_\Q$ as in the proof of Theorem~\ref{hol:ot:main:second}\,(a).
A nonpolynomial entire solution has breaks $\delta\ge\delta_0$ with
$\underline N_\delta$, hence $N_\delta$, arbitrarily large
(Lemma~\ref{hol:ot:lem:breaks}), and the upper endpoint--gap pair
$(N_\delta,\mathsf g^+)$ of $\phi_\delta$ has $1\le\mathsf g^+\le N_\delta$;
Theorem~\ref{hol:ot:thm:gaps} and Lemma~\ref{hol:ot:lem:exterior} contradict the
boundedness.
\end{proof}

This is a sufficient certificate, not a decision algorithm for third-order
equations: deciding the integral zero set of a bivariate polynomial is a
separate arithmetic problem, and an unbounded zero set of
$\Upsilon_{\mathfrak h_P}$ is necessary for escaping the obstruction, not
sufficient for the existence of an entire solution.

\subsection{The exterior form of the theta equation}
\label{hol:ot:sub:thetaform}

Let $P$ be the polynomial \eqref{hol:td:eq:jetpoly} with the constants
\eqref{hol:td:eq:constants}; it is homogeneous of degree six in the jets. In
Euler jets $E_j$ the substitution \eqref{hol:ot:eq:stirling} reads
\begin{equation}\label{hol:ot:eq:eulersub}
 Y_0=E_0,\quad Y_1=E_1/z,\quad Y_2=(E_2-E_1)/z^2,\quad
 Y_3=(E_3-3E_2+2E_1)/z^3 ,
\end{equation}
and after multiplication by $z^6$ the highest $z$-power is six. Its block,
dehomogenized by $E_0=1$, $E_1=U$, $E_2=U^2+V$, $E_3=U^3+3UV+W$, is
\begin{equation}\label{hol:ot:eq:thetatop}
 W^2-4(\mathfrak c-V)^3+g_2(\mathfrak c-V)+g_3
 =W^2+4V^3-12\mathfrak cV^2+(12\mathfrak c^2-g_2)V-4\mathfrak c^3
 +\mathfrak cg_2+g_3 .
\end{equation}
Directly: in Euler jets the top part of $\mathcal H[f]$ is
$f\vartheta^2f-(\vartheta f)^2$, that is $E_0E_2-E_1^2$, which dehomogenizes to
$V$; the top part of $z\mathcal K[f]$ is
$f\vartheta\mathcal H_0-2\vartheta f\,\mathcal H_0$ with $\mathcal H_0$ that
part, which dehomogenizes to $W$; so $\mathcal B[f]$ contributes
$\mathfrak c-V$ and $(z^2-4)\mathcal K[f]^2$ contributes $W^2$. The
coefficient of
$W^2$ is $1$ and all coefficients lie in $\Z[\mathfrak c,g_2,g_3]$, of
nonnegative valuation, so $\beta=0$ in \eqref{hol:ot:eq:exterior}. With the
residues $\res\mathfrak c=\frac1{12}$, $\res g_2=\frac1{12}$ and
$\res g_3=-\frac1{216}$ the constant and linear terms cancel and the
dehomogenized exterior form is
\begin{equation}\label{hol:ot:eq:thetainitial}
 W^2-V^2+4V^3 .
\end{equation}
Its value at $V=W=0$ is $0$, which is $\chi_P=0$ in
Proposition~\ref{hol:td:prop:corner}, by Remark~\ref{hol:ot:rem:corner}. Its
transverse part of least degree is $W^2-V^2$, so
Theorem~\ref{hol:ot:thm:gaps} gives $\Upsilon(T,G)=G^2-1$: every exterior upper
and lower gap of a solution is one. The Euler-jet conversion and
\eqref{hol:ot:eq:thetainitial} are checked symbolically by source~14's
program.

\subsection{Theta breaks and eventual Newton convexity}
\label{hol:ot:sub:convexity}

\begin{theorem}[Polynomial solutions of the theta exterior equation;
source~14]\label{hol:ot:thm:classification}
Let $0\ne\phi\in k[X]$. Then
$\mathsf w_\phi^2-\mathsf v_\phi^2+4\mathsf v_\phi^3=0$ if and only if
$\phi=bX^m$ or $\phi=bX^m(1+aX)$ with $m\in\N$ and $a,b\in k^\times$.
\end{theorem}

\begin{proof}
If $\mathsf v_\phi=0$, then $\mathsf u_\phi$ is constant and $\phi$ is a
monomial, as in the proof of Lemma~\ref{hol:ot:lem:endpoints}. Suppose
$\mathsf v_\phi\ne0$ and put $\mathsf k=\mathsf w_\phi/\mathsf v_\phi$. The
equation becomes $\mathsf k^2=1-4\mathsf v_\phi$. If $\mathsf k=0$ then
$\mathsf v_\phi=\frac14$, impossible because the constant term of
$\vartheta_X\mathsf u_\phi$ in its Laurent expansion at $0$ is zero. Applying
$\vartheta_X$ and using $\mathsf w_\phi=\mathsf k\mathsf v_\phi$ gives
$2\mathsf k\vartheta_X\mathsf k=-4\mathsf k\mathsf v_\phi$, so
$\vartheta_X\mathsf k=(\mathsf k^2-1)/2$; and $\mathsf k\ne\pm1$, since either
would give $\mathsf v_\phi=0$. For $\mathsf r=(\mathsf k-1)/(\mathsf k+1)$ a
direct computation gives $\vartheta_X\mathsf r=\mathsf r$, so
$\mathrm d(\mathsf r/X)/\mathrm dX=0$ and $\mathsf r=b_1X$ with $b_1\in
k^\times$. With $a=-b_1$,
\[
 \mathsf k=\frac{1-aX}{1+aX},\qquad \mathsf v_\phi=\frac{aX}{(1+aX)^2},\qquad
 \vartheta_X\Bigl(\mathsf u_\phi-\frac{aX}{1+aX}\Bigr)=0 .
\]
The constant is $m=\operatorname{ord}_X\phi$, by evaluation at $0$, so
$\phi/(X^m(1+aX))$ has zero logarithmic derivative and is a nonzero constant.
Conversely a monomial has $\mathsf v_\phi=\mathsf w_\phi=0$, and for
$\phi=bX^m(1+aX)$
\[
 \mathsf u_\phi=m+\frac{aX}{1+aX},\quad
 \mathsf v_\phi=\frac{aX}{(1+aX)^2},\quad
 \mathsf w_\phi=\frac{aX(1-aX)}{(1+aX)^3},
\]
which satisfy the equation.
\end{proof}

For a binomial with gap $\mathsf g$, $\phi=bX^m(1+aX^{\mathsf g})$, the same
computation with $aX^{\mathsf g}$ in place of $aX$ gives
$\mathsf w_\phi^2-\mathsf v_\phi^2+4\mathsf v_\phi^3
=(\mathsf g^2-1)\mathsf v_\phi^2$: gap one is accepted and every larger
binomial gap is rejected (source~14, Appendix~A).

\begin{theorem}[Eventual strict Newton convexity; source~14]
\label{hol:ot:thm:convexity}
Let $\mathfrak c,g_2,g_3\in\KK$ have nonnegative valuation and the residues
$\frac1{12},\frac1{12},-\frac1{216}$, and let a nonpolynomial
$f=\sum_na_nz^n\in\EK$ satisfy the cleared equation \eqref{hol:td:eq:cleared}
with these constants. Then:
\begin{enumerate}[label=\textup{(\roman*)}]
\item after passing to $\Gamma_\Q$, every break $\delta\ge\delta_0$ of $f$ has
initial polynomial $bX^m(1+aX)$ with $a,b\in k^\times$;
\item $a_n\ne0$ for every sufficiently large $n$;
\item the differences $\delta_n=v(a_{n+1})-v(a_n)\in\Gamma$ are strictly
increasing for all sufficiently large $n$ and tend cofinally to $+\infty$ in
$\Gamma$.
\end{enumerate}
\end{theorem}

\begin{proof}
Pass to $\Gamma_\Q$ (Lemma~\ref{hol:ot:lem:extension}). The exterior form of
the equation, dehomogenized, is \eqref{hol:ot:eq:thetainitial}, by the
computation of Section~\ref{hol:ot:sub:thetaform}, which uses only the
residues and the nonnegative valuations of the constants; so, beyond the
$\delta_0$ of Lemma~\ref{hol:ot:lem:exterior}, every initial polynomial
satisfies $\mathsf w^2-\mathsf v^2+4\mathsf v^3=0$, and
Theorem~\ref{hol:ot:thm:classification} excludes the monomial case at a break,
proving~(i). Beyond $\delta_0$, therefore, a break has exactly two active
indices $n$ and $n+1$, with no skipped index and no three-way tie, and between
breaks the active index is single. By Lemma~\ref{hol:ot:lem:localfinite}
only finitely many indices are active on a bounded interval, and a change of
the active index happens only at a break (as in the proof of
Lemma~\ref{hol:ot:lem:breaks}); so the active index increases by exactly one
at each break, and by Lemma~\ref{hol:nl:lem:active} it exceeds every bound.
Every sufficiently large integer is therefore active at some scale, and its
coefficient is nonzero, proving~(ii). The break at which $n$ is replaced by
$n+1$ is the solution of $v(a_n)-n\delta=v(a_{n+1})-(n+1)\delta$, namely
$\delta=\delta_n$, which lies in $\Gamma$ without any division. These scales
increase by the monotonicity in Lemma~\ref{hol:ot:lem:localfinite}, strictly
because a repeated break would carry three active indices, and they are
cofinal because an upper bound would put infinitely many active indices in one
bounded interval. Nonvanishing of coefficients, inequalities between their
values and cofinality in $\Gamma_\Q$ descend to $\Gamma$, since $\Gamma$ is
cofinal in $\Gamma_\Q$.
\end{proof}

For $f=\mathcal T_p$ with $\mu$ an order unit, \eqref{hol:td:eq:coeffval}
gives the realization $\delta_n=(2n+1)\mu$,
$\gamma_{\delta_n}=-n(n+1)\mu$ and $\phi_{\delta_n}=X^n(1+X)$ for $n\ge0$: at
$\delta_n$ the weighted value of degree $j$ is $(j^2-(2n+1)j)\mu$, least exactly
at $j=n,n+1$. This is the Newton profile of Section~\ref{hol:td:sub:bilateral},
and it makes the order-three escape visible without analytic asymptotics:
unbounded degrees coexist with the fixed gap one.
Theorem~\ref{hol:ot:thm:convexity} is a necessary condition on solutions, not
a classification of the entire solutions of the theta equation: it does not
force $v(a_n)=n^2\mu$, does not determine coefficients from their values, and
does not assert that every eventually strictly convex value sequence occurs.

\subsection{Surreal consequences and boundaries}\label{hol:ot:sub:surreal}

Through \eqref{hol:eq:embedding}, Theorem~\ref{hol:ot:main:second} applies to
every set-sized surreal or surcomplex field $k((\omega^{-\Gamma}))$ with
$\Gamma\subseteq\No$ and $k=\R$ or $\C$: there no nonpolynomial strongly
entire series satisfies an algebraic differential equation of order at most
two, whatever the rank of $\Gamma$, and algebraic closedness of the
coefficient field is not needed. For $\Gamma=\Q$ or $\R$ and $p=\omega^{-1}$
the series $\mathcal T_p$ of Section~\ref{hol:td:sub:surreal} has three
algebraically independent jets over the field of rational functions and an
explicit equation of order three. The following proposition, from source~14,
states once for every series what Section~\ref{hol:sub:notallsurcomplex},
Section~\ref{hol:td:sub:surreal} and Corollary~\ref{hol:nl:cor:theta} show for
particular ones.

\begin{proposition}[No nonpolynomial series is entire on the class;
source~14]\label{hol:ot:prop:fullclass}
A power series $\sum_na_nz^n$ with coefficients in $\No[i]$ that is strongly
summable at every point of $\No[i]$ has only finitely many nonzero
coefficients.
\end{proposition}

\begin{proof}
Suppose infinitely many $a_n$ are nonzero, and let $v(a)$ be the negative of
the largest Conway exponent of $a\ne0$. The values $v(a_n)$ of the nonzero
coefficients and their pairwise differences form a set of surreals, so there
is a positive surreal $\gamma$ larger than the absolute value of every such
difference. At $x=\omega^\gamma$ the leading values of the terms are
$v(a_n)-n\gamma$, and along increasing indices with $a_n\ne0$ they decrease
strictly, since $v(a_{n'})-v(a_n)<\gamma\le(n'-n)\gamma$. They cannot lie in a
well-ordered union of Hahn supports, contradicting strong summability at $x$.
\end{proof}

This is a size-and-domain obstruction, not an instance of second-order
rigidity: entireness is always relative to one set-sized field, $\mathcal T_p$
is entire on its field because $\mu$ is an order unit there, and enlarging to
every surreal scale destroys that property. The zeros of $\mathcal T_p$ and
their omnific floors and ceilings are those of Theorem~\ref{hol:td:thm:zeros}
and Corollary~\ref{hol:td:cor:omnific}, inherited; source~14's new contribution
there is the universal lower bound for the function that produces them and the
necessary Newton geometry of Theorem~\ref{hol:ot:thm:convexity} for every other
entire solution of its equation.

\paragraph{Boundaries recorded by source~14.}
In characteristic $\mathfrak p>0$ and with an order unit $\mu$, the series
$\sum_{n\ge0}t^{n^2\mu}z^{\mathfrak pn}$ is nonpolynomial and strongly entire with
derivative zero, so Theorem~\ref{hol:ot:main:second} fails already in order
one; this is Example~\ref{hol:nl:ex:charp} with $\Q$ replaced by any group with
an order unit. Mixed-characteristic $\mathfrak p$-adic fields are not covered merely
because they have characteristic zero: their residue characteristic and
integer valuations differ from the hypotheses used here. The ordinary
exponential $\sum z^n/n!$ solves $D_zf=f$ formally but is not strongly
summable at $z=1$, since infinitely many terms have exponent zero; allowing
ordinary convergence inside a fixed coefficient layer would change the
function class and need a different theorem. The proof classifies neither
functions patched from local formulas nor analytic identities in the fine
topology of $\No[i]$. It passes only to a cofinal divisible extension and never
reduces the group to rank one; the no-order-unit argument uses finite rational
rank of a coefficient field, not of the ambient group, and a group of high rank
can have an order unit while one of countable cofinality can lack one
(Examples~\ref{hol:td:ex:lex} and~\ref{hol:td:ex:cofinality}).

\subsection{Predecessors, repository comparison and priority}
\label{hol:ot:sub:predecessors}

\paragraph{Literature.}
Source~14 consulted NIST DLMF \S20.5 for the triple product and logarithmic
derivatives~\cite{dlmftheta}, \S23.3 for the Weierstrass differential
equation~\cite{dlmfelliptic}, and \S32.2 for the Painlev\'e
equations~\cite{dlmfpainleve}; the Mantova--Matusinski
survey~\cite{mantovamatusinski} for normal forms, not as a source for the
second-order theorem; and the publisher abstract of Sebbar's
paper~\cite{sebbar}, which explicitly treats classical third-order theta
relations, without its full text. It also consulted the Wikipedia article on
surreal numbers for orientation only. General searches for second-order Hahn
and non-Archimedean rigidity returned unrelated results, including Hahn
difference operators; these are not evidence of absence. No exhaustive
MathSciNet or zbMATH review and no full subscription-literature review was
made. For this merge the DLMF sections and the Mantova--Matusinski survey were
checked, and the bibliographic data of Sebbar's paper were checked against
its Crossref record; its content was not checked by this merge, and nothing
in this part rests on it.

\paragraph{Repository comparison at the pin.}
Source~14 read the repository at \texttt{b895e86} through a connector: the
root README and tree, \texttt{docs/README.md}, this report's README and
neighbouring report inventories. The formalization ledger did not return text,
so no assertion about it is made, and the new proof is not represented as
covered by the Lean build. It took source~13 from the author's file library
(\texttt{article(20260923-193210).tex}, baseline \texttt{4cdeaec}). Its
statements were checked for this merge and are accurate at the pin, as
recorded in Section~\ref{hol:ot:sub:conventions}. This was a targeted
comparison, not an audit of every file, branch or archived draft.

\paragraph{Priority.}
The triple product, the theta--elliptic identity and third-order theta
relations are classical; the Chebyshev descent of the theta series, its
equation, zeros and omnific rounding are source~13's; the no-order-unit
obstruction and the finite-generation principle are the coefficient-field
part's. The proposed contribution of source~14 is universal second-order
rigidity for every ordered value group and every characteristic-zero
coefficient field (Theorem~\ref{hol:ot:main:second}\,(a)), with the endpoint
lemma at degenerate breaks, the resulting exact minimum order
(Theorem~\ref{hol:ot:main:second}\,(b)), the third-order endpoint--gap
obstruction, the classification of theta initial polynomials and eventual
Newton convexity, and the stated consequences for two-state systems and
polynomial-denominator quotients. Whether an equivalent result exists in the
literature needs a more exhaustive review; the answered question is a
question of this report, not a named classical conjecture.

\subsection{What the order-three-threshold part does not claim}
\label{hol:ot:sub:nonclaims}

Every limitation stated by source~14, in its manuscript, README, audit file or
verification record, is kept.
\begin{enumerate}[label=\textup{(O\arabic*)}]
\item No Lean formalization, no independent referee report and no certified
bibliographic priority accompany source~14; no new Lean file is presented as
verified, and the repository's build and axiom audit do not cover this part.
Independent review of the exterior reduction and the endpoint lemma is the
most useful next check.
\item Only the question of this report is resolved; no named classical
conjecture is claimed solved.
\item The repository comparison at \texttt{b895e86} was targeted; the
formalization ledger was not read, and no audit of every branch or archived
draft was made.
\item Entireness is relative to one specified Hahn field; $\mathcal T_p$ is
not strongly entire on $\No[i]$, and no nonpolynomial series is
(Proposition~\ref{hol:ot:prop:fullclass}).
\item $D_z$ acts on the external variable and kills every Hahn coefficient; it
is not the Berarducci--Mantova derivation.
\item The theta product and the elliptic identity are classical, and the theta
descent with its equation, zeros and omnific rounding and the no-order-unit
obstruction are prior; source~14 reproves them for completeness and claims
none. ``No omnific zero'' concerns the strong domain only.
\item Theorem~\ref{hol:ot:thm:convexity} is necessary only: it gives no
sufficiency and no classification of entire solutions, does not force
$v(a_n)=n^2\mu$, does not determine coefficients from their values, and does
not assert that every eventually strictly convex sequence occurs.
\item Corollary~\ref{hol:ot:cor:gapcriterion} is a sufficient certificate, not
a decision algorithm for third-order equations; an unbounded zero set of
$\Upsilon_{\mathfrak h_P}$ is necessary, not sufficient, for an entire
solution.
\item Corollary~\ref{hol:ot:cor:meromorphic} concerns quotients $g/\mathsf p$ with a
polynomial denominator only; no global division theorem in a Hahn analytic
category is asserted.
\item Corollary~\ref{hol:ot:cor:systems} concerns rational systems; arbitrary
transcendental right-hand sides are outside it.
\item Nothing is claimed about local formal solutions or about classical
complex Painlev\'e functions; the all-scale requirement, not local existence
theory, does the work.
\item Positive characteristic is excluded, with a counterexample, and
mixed-characteristic $\mathfrak p$-adic fields are not covered.
\item Theorem~\ref{hol:ot:main:second} gives no uniform degree bound for
polynomial solutions of a fixed equation
(Example~\ref{hol:ot:ex:degenerate}).
\item The proof does not reduce every group to rank one; it classifies neither
functions patched from local formulas nor identities in the fine topology,
and a different summation convention is a different function class.
\item The 600 finite checks, with the theta product and cleared equation
compared modulo $p^{48}$ (source~14's $q^{48}$), do not certify the
quantifiers over arbitrary supports, value groups, scales or formal solutions;
the classification of Theorem~\ref{hol:ot:thm:classification} is proved for
rational functions, not extrapolated from the finite enumeration, and no
floating-point agreement is used as evidence.
\item Sebbar's paper was seen through its publisher abstract only, and broad
searches are not evidence of absence; no exhaustive MathSciNet or zbMATH
review was made.
\item The research questions of source~14
(Questions~\ref{hol:q:nonlinear} and~\ref{hol:ot:q:lift}--\ref{hol:ot:q:omnific})
are not asserted solved.
\item \emph{[merge]} Remark~\ref{hol:ot:rem:corner},
Corollary~\ref{hol:ot:cor:examples}, the named break scale in the proof of
Lemma~\ref{hol:ot:lem:breaks}, the reading of
Corollaries~\ref{hol:ot:cor:threejets} and~\ref{hol:ot:cor:systems} inside a
differential field containing $\KK((z))$, and the four-state observation after
Corollary~\ref{hol:ot:cor:systems} are the merge's; no historical priority is
claimed for them.
\end{enumerate}

%======================================================================
% Unit-dilation part: source 18. Labels use hol:ud:.
%======================================================================
\section{Unit-dilation rigidity: valuations of Hahn exponential polynomials}
\label{hol:ud:sec:setting}

This section and the next contain the material of source~18, delivered after
the order-three-threshold part had been merged. They answer
Question~\ref{hol:q:mixedunit} positively for every value group: for a
nontorsion $q$ with $v(q)=0$, every strongly entire solution of a nonzero
mixed equation \eqref{hol:eq:mixed} is a polynomial
(Corollary~\ref{hol:ud:cor:unitrigidity}). More generally, for a finite set of
multipliers with no root-of-unity quotient, a nonpolynomial strongly entire
solution of \emph{some} nonzero mixed linear equation exists exactly when the
absolute valuation of one quotient is an order unit
(Theorem~\ref{hol:ud:main:mixed}). The missing ingredient named in the status
of Question~\ref{hol:q:mixedunit}, uniform control of $v(P(n,q^n))$ for a
bivariate $P$ at unit valuation, is Corollary~\ref{hol:ud:cor:bivariate}; it
comes from a valuation form of the Skolem--Mahler--Lech theorem for
exponential polynomials with arbitrary Hahn coefficients
(Theorems~\ref{hol:ud:thm:unitvalues} and~\ref{hol:ud:thm:affine}). The
recurrence is then closed by the escape chain of
Section~\ref{hol:sec:escape}; neither the finite coefficient field of the
coefficient-field part nor any nonlinear method is used.

\subsection{Source, conventions and the merge's choices}
\label{hol:ud:sub:conventions}

\begin{convention}[Setting of the unit-dilation part]
\label{hol:ud:conv:field}
In Sections~\ref{hol:ud:sec:setting}--\ref{hol:ud:sec:mixed}, $k$ is an
arbitrary field of characteristic zero with the trivial valuation, $\Gamma$ is
as in Convention~\ref{hol:conv:group}, $\KK=k((t^\Gamma))$, and $v$, $\lc$,
$\res$, strong summability, $\Dom(f)$, $\EK$, $D_z$ and $\sigma_\lambda f(z)=
f(\lambda z)$ are as in Conventions~\ref{hol:nl:conv:field}
and~\ref{hol:cf:conv:field}. A \emph{unit} is an element of valuation zero,
not merely a nonzero element. Every $\lambda\in\KK$ with $v(\lambda)=0$ has a
unique expression
\begin{equation}\label{hol:ud:eq:unitdecomp}
 \lambda=\zeta(1+\tau),\qquad \zeta=\res(\lambda)\in k^\times,\qquad
 \tau=0\ \text{or}\ v(\tau)>0 .
\end{equation}
The multipliers of a mixed operator are $\lambda_1,\ldots,\lambda_m\in\KK^\times$,
as the dilation parameters of Theorem~\ref{hol:cf:main:entire}\,(b); $q$ is a
single dilation parameter, as in Definition~\ref{hol:def:operators}.
\end{convention}

For $k=\C$ one has $\KK=\K$ and $\EK=\E$. As for the four previous parts,
the generality in $k$ is recorded for this part; unlike them, this part also
contains Theorem~\ref{hol:main:ode}, the negative direction of
Theorem~\ref{hol:main:q} and Theorem~\ref{hol:thm:mixed} as special cases over
every such $k$ (Remark~\ref{hol:ud:rem:contains}).

\paragraph{The source.}
Source~18 is \emph{Unit-Dilation Rigidity at Surreal Scales: A valuation form
of Skolem--Mahler--Lech and an exact relative-scale criterion for mixed
differential--dilation equations} (26 pages, A4, 23 September 2026), delivered
as the archive \texttt{Surreal\_Unit\_Dilation\_Rigidity} (batch~31,
manuscript~05). It consulted the repository at
\begin{center}\small\texttt{bcac55ae6fd3f2b354e568ba1d2e94d496a2cc59},\end{center}
where this report's \texttt{article.tex} is the blob
\texttt{86fc4925b4c3997b984894e6b4756e327f27a4c6}; its shipped file
\path{18-unit-dilation-SOURCES_AND_SCOPE.md} lists the line ranges it read. At
that pin this report was the five-source text (sources~08--12): neither the
theta-descent nor the order-three-threshold part had been merged. Source~18's
description of this report was accurate there: Question~\ref{hol:q:mixedunit}
was open with an order unit, the mixed theorem at unit valuation was not
asserted, and ``higher-order nonlinear equations with an order unit'' were
open. The last point is stale now. Theorem~\ref{hol:td:main:boundary} and
Theorem~\ref{hol:ot:main:second} have since settled that question up to the
classification of Question~\ref{hol:q:nonlinear}, and source~18's first
research question is answered negatively in the current text
(Remark~\ref{hol:ud:rem:nonlinearq}). Source~18 claims neither the pure
differential theorem, the cyclic pure-dilation criterion, the mixed theorem at
nonzero noncofinal valuation, the partial theta series, nor the escape
mechanism; it credits them to this report.

\paragraph{Renamed symbols.}
The following symbols of source~18 would collide with symbols of the earlier
parts. They are renamed here, in the new material only.
{\small
\begin{longtable}{@{}>{\raggedright\arraybackslash}p{0.27\textwidth}
>{\raggedright\arraybackslash}p{0.2\textwidth}
>{\raggedright\arraybackslash}p{0.45\textwidth}@{}}
\toprule
Source~18 & Here & Reason\\
\midrule
\endhead
$K=k((t^\Gamma))$, $\mathcal E_\Gamma(k)$, $D$ & $\KK$, $\EK$, $D_z$ & $\K$
is $\C((t^\Gamma))$ in this report.\\
multipliers $u_1,\ldots,u_m$; $u=c(1+\epsilon)$ & $\lambda_1,\ldots,\lambda_m$;
$\lambda=\zeta(1+\tau)$ & $u$ is the Laurent variable of
Section~\ref{hol:td:sub:laurent}; $\epsilon\in\Gamma$ below; $c$ is a
coefficient.\\
$A(n)=\sum_jP_j(n)u_j^n$; $P_j$ & $\mathcal Y(n)=\sum_jR_j(n)\lambda_j^n$;
$R_j$ & $P$ is a differential polynomial; $R_h$ plays the same role in
Lemma~\ref{hol:lem:qrec}.\\
$A_j(n)$ (recurrence coefficients), $Q_{j,\ell}$ & $c_j(n)$, $R_{j,\ell}$ & As
in Lemma~\ref{hol:lem:escape}.\\
$Q_{j,\gamma}$ (coefficient polynomials); envelope $T$ & $r_{j,\gamma}$;
$\Sigma$ & $T$ is an indeterminate.\\
forward shift $S$ & $\mathrm{Sh}$ & $S$ is a support alphabet, as here.\\
period $M$; residue class $r$ & $\mathsf M$; $b$ & $M(S)$ is the positive
monoid; $r$ is a derivative order.\\
slopes $\lambda_r=v(u_j)$ & $\varpi_b$ & $\lambda$ now names a multiplier.\\
relative scale $\mu$ & $\operatorname{rs}(\lambda_1,\ldots,\lambda_m)$ & $\mu$
is an order unit or a positive scale, with $p=t^\mu$.\\
noncofinality witness $\eta$ & $\gamma$ of \eqref{hol:eq:dominates} & $\eta$ is
$v(\tau)$ in Theorem~\ref{hol:thm:unitorbit}.\\
operator coefficients $c_{\ell,r,k}$, $z^k$; output index $m$ &
$\varkappa_{\ell,r,i}$, $z^i$; $\bar n$ & $k$ is the coefficient field; $m$
counts multipliers.\\
right side $h$, $\deg h\le e$; binomial index $h$ & $\mathsf h$,
$\deg\mathsf h\le\bar e$; $i$ & $h$ is the order of a root of unity, as in
Theorem~\ref{hol:thm:unitorbit}.\\
$\Theta_q$, $q$ with $v(q)$ an order unit & $\Theta_p$, $p$ &
\eqref{hol:eq:theta}; $q$ is kept for a dilation parameter.\\
order unit $H$; $H_*$; $F=\sum t^{n^2H}z^n$; $G=\sum t^{n^2H}z^{hn}$ & $\mu$;
$e_*$; $G_\mu$; $G_\mu(z^h)$ & $F$ is \eqref{hol:eq:explicitF}; $G$ is the
series of Corollary~\ref{hol:nl:cor:theta}, which is $G_1$ over $\C((t^\Q))$.\\
$\Delta$ (maximal proper convex subgroup) & $\Delta^{\mathrm{top}}$ &
$\Delta_\beta$ is the least convex subgroup containing $\beta$.\\
field $F$ of its Definition~3.1 & $K_0$ & $F$ is \eqref{hol:eq:explicitF};
$K_0$ as in Theorem~\ref{hol:cf:thm:coefficients}.\\
Theorems~4.2, 4.7, 6.4, 7.2 & Theorems~\ref{hol:ud:thm:unitvalues},
\ref{hol:ud:thm:affine}, \ref{hol:ud:thm:rigidity},
\ref{hol:ud:thm:dichotomy} & Theorems~A--N of this report are different
theorems; Theorem~\ref{hol:ud:main:mixed} collects the last two.\\
\bottomrule
\end{longtable}}
The shipped file \path{18-unit-dilation-SOURCES_AND_SCOPE.md} keeps the
source's notation.

\paragraph{Words with two meanings.}
A \emph{unit} is a valuation-zero element (Convention~\ref{hol:ud:conv:field}),
and a \emph{unit dilation} is $\sigma_q$ with $v(q)=0$. An \emph{order unit}
is Definition~\ref{hol:def:orderunit}; the two words are unrelated. The
\emph{relative scale} of a list of multipliers is
$\operatorname{rs}(\lambda_1,\ldots,\lambda_m)=\max_{i,j}\abs{v(\lambda_i)-
v(\lambda_j)}$ (Definition~\ref{hol:ud:def:rs}); it is \emph{noncofinal} if it
is $0$ or not an order unit. \emph{Nondegenerate} describes an exponential
polynomial none of whose base quotients is a root of unity; it has nothing
to do with the nondegenerate corners of Section~\ref{hol:nl:sec:higher}. A
\emph{root-of-unity quotient} concerns the actual multipliers $\lambda_i/\lambda_j$,
not their residues; the residues may be roots of unity
(Example~\ref{hol:ex:unitperiod}).

\paragraph{Results printed once.}
Source~18 reproves several facts the report already contains. They are
printed once and cited: its definition of strong summability and entireness
(Definition~\ref{hol:def:entire}); its Lemma~2.2 (sums of well-ordered sets)
and Lemma~2.4 (the positive monoid), which are
Lemma~\ref{hol:lem:support}\,(1) and~(2), except that its proof of the
finite-word lemma is printed below as Lemma~\ref{hol:ud:lem:words}, since this
report imports that lemma from Higman; its Lemmas~2.5 and~2.6, which are
Lemma~\ref{hol:lem:certificate}\,(a) and Lemma~\ref{hol:lem:nonincreasing};
the invariance of $\EK$ under $\sigma_\lambda$ and $D_z$, which is
Proposition~\ref{hol:prop:closure}; its Lemma~3.2 and Corollary~3.4
(polynomial--exponential independence, finitely many zeros of a nonzero
nondegenerate exponential polynomial, also on every arithmetic progression),
which are Lemma~\ref{hol:cf:lem:exppoly} and its proof; its Theorem~3.3, the
classical Skolem--Mahler--Lech theorem, imported there as here
(Section~\ref{hol:cf:sub:coefficients}); its Example~4.8, which is the third
example of Remark~\ref{hol:rem:shortcut}; its Section~5 (Theorem~5.1,
Corollary~5.2, Proposition~5.3), which is Lemma~\ref{hol:lem:escape},
Theorem~\ref{hol:thm:escapeexclusion} and Corollary~\ref{hol:cor:boundedcost},
the last with the closed region $v(x)\le-B^+$ where source~18 states the open
one $v(x)<-B$; the entireness half of its Lemma~7.1, which is
Theorem~\ref{hol:thm:thetadomain} by the route of
Remark~\ref{hol:rem:thetasecond}; the embedding \eqref{hol:eq:embedding}; its
parity example (Section~9.1), which is Example~\ref{hol:ex:unitperiod}; its
torsion example (Section~9.5), which is the periodic-support construction of
Section~\ref{hol:sec:torsion}; its positive-characteristic series, which is
Example~\ref{hol:nl:ex:charp}; and its remark on $\Gamma=0$, which is in
Section~\ref{hol:sec:fields}.

\paragraph{Where each numbered result went.}
Lemma~2.3 is Lemma~\ref{hol:ud:lem:words}; Lemma~3.5 is
Lemma~\ref{hol:ud:lem:period}; Lemma~4.1, Theorem~4.2, Remark~4.3,
Corollary~4.4 and Remark~4.5 are Lemma~\ref{hol:ud:lem:envelope},
Theorem~\ref{hol:ud:thm:unitvalues}, Remark~\ref{hol:ud:rem:uniform},
Corollary~\ref{hol:ud:cor:bivariate} and Remark~\ref{hol:ud:rem:elementary};
Lemma~4.6 and Theorem~4.7 are Lemma~\ref{hol:ud:lem:lines} and
Theorem~\ref{hol:ud:thm:affine}; Lemma~6.1, Corollary~6.2, Definition~6.3,
Theorem~6.4, Corollaries~6.5 and~6.6 are Lemma~\ref{hol:ud:lem:operator},
Corollary~\ref{hol:ud:cor:faithful}, Definition~\ref{hol:ud:def:rs},
Theorem~\ref{hol:ud:thm:rigidity} and
Corollaries~\ref{hol:ud:cor:unitrigidity} and~\ref{hol:ud:cor:linind};
Lemma~7.1, Theorem~7.2, Remark~7.3, Corollary~7.4 and Proposition~7.5 are
Lemma~\ref{hol:ud:lem:theta}, Theorem~\ref{hol:ud:thm:dichotomy},
Remark~\ref{hol:ud:rem:quantifiers}, Corollary~\ref{hol:ud:cor:cyclic} and
Proposition~\ref{hol:ud:prop:top}; Corollary~8.1 is
Corollary~\ref{hol:ud:cor:surcomplex}; Proposition~10.1 is
Proposition~\ref{hol:ud:prop:finitespace}; Research question~11.1 is recorded
in Remark~\ref{hol:ud:rem:nonlinearq}, and Research
questions~11.2--11.12 are Questions~\ref{hol:ud:q:prescribed}--%
\ref{hol:ud:q:formal}. Theorem~\ref{hol:ud:main:mixed} in the introduction
combines Theorems~6.4 and~7.2.

\paragraph{The merge's choices.}
The review of source~18 for this merge found its proofs correct. The merge
adds four points, each marked. First, Remark~\ref{hol:ud:rem:contains}
records that Theorem~\ref{hol:ud:thm:rigidity} contains
Theorem~\ref{hol:main:ode}, the negative direction of
Theorem~\ref{hol:main:q} and Theorem~\ref{hol:thm:mixed}, over every
characteristic-zero $k$, and compares it with
Theorem~\ref{hol:cf:main:entire}\,(b). Second, the exterior region of
Theorem~\ref{hol:ud:thm:rigidity} is closed, by
Corollary~\ref{hol:cor:boundedcost}. Third,
Remark~\ref{hol:ud:rem:nonlinearq} records the answer in the current text to
source~18's nonlinear question. Fourth, Corollary~\ref{hol:ud:cor:cyclic}
is read beside Theorem~\ref{hol:main:q}: adding derivatives creates no new
existence case. Source~18's theorems, proofs, examples and non-claims are
otherwise printed as delivered, in the renamed notation.

\subsection{Finite words and exponential polynomials}
\label{hol:ud:sub:sml}

Lemma~\ref{hol:lem:support}\,(2) imports the finite-word lemma from
Higman~\cite{higman}. Source~18 proves the case it needs, by the usual
minimal-bad-sequence argument; the proof is printed here because it makes the
positive-monoid lemma self-contained. It is standard, not new.

\begin{lemma}[Finite-word lemma for a well-ordered alphabet; source~18]
\label{hol:ud:lem:words}
Let $S$ be well ordered, and order finite words over $S$ by subsequence
embedding with coordinatewise increase of letters. Every infinite sequence of
words contains an earlier word that embeds in a later one.
\end{lemma}

\begin{proof}
Suppose a bad sequence exists, that is, one with no such pair. Choose a bad
sequence $w_0,w_1,\ldots$ so that, at each position, $w_n$ has the least
possible length among the choices that allow a bad continuation of the prefix
already chosen. No word is empty. Write $w_n=w_n'a_n$ with $a_n$ its last
letter, and choose $i_0<i_1<\cdots$ with $a_{i_0}\le a_{i_1}\le\cdots$, which
is possible in a well-ordered set. By the minimality at position $i_0$ the
sequence $w_0,\ldots,w_{i_0-1},w'_{i_0},w'_{i_1},\ldots$ is not bad. An
embedding between two words of the prefix contradicts badness of the original
sequence; an embedding of a prefix word into some $w'_{i_j}$ is an embedding
into $w_{i_j}$; and an embedding $w'_{i_j}\hookrightarrow w'_{i_l}$ with $j<l$
extends, by $a_{i_j}\le a_{i_l}$, to $w_{i_j}\hookrightarrow w_{i_l}$. Every
case is a contradiction.
\end{proof}

With it, the proof of Lemma~\ref{hol:lem:support}\,(2) needs no import: a
strictly decreasing sequence of word sums would give an embedding of an
earlier word into a later one, which cannot decrease the sum; and if
infinitely many distinct nondecreasingly sorted words had one sum, an
embedding between two of them would force equal lengths, since inserted
letters are positive, and then equal letters, since the sums agree.

An \emph{exponential polynomial} over a field $K_0$ of characteristic zero is a
sequence $\mathcal Y(n)=\sum_{j=1}^mR_j(n)\lambda_j^n$ with $R_j\in K_0[U]$ and
$\lambda_j\in K_0^\times$; it is \emph{nondegenerate} if the bases are distinct
and no quotient $\lambda_i/\lambda_j$, $i\ne j$, is a root of unity. A
constant-coefficient recurrence annihilates it (a product of powers of
$\mathrm{Sh}-\lambda_j$, with $\mathrm{Sh}$ the forward shift), so the
Skolem--Mahler--Lech theorem applies over $k$ and over $\KK$. Source~18 imports
it in the characteristic-zero form of Lech~\cite{lech} and Bell's
account~\cite{belldoc}; the report imports it through Bell's
Theorem~1.1~\cite{bell}. Its consequence used below is
Lemma~\ref{hol:cf:lem:exppoly}: a nonzero nondegenerate exponential polynomial
has only finitely many zeros, and so has its restriction to every arithmetic
progression, the restriction being again nonzero and nondegenerate; the proof
of that lemma also shows that polynomial--exponential sequences with distinct
bases are linearly independent on every tail.

\begin{lemma}[One period removes all residue torsion; source~18]
\label{hol:ud:lem:period}
For $\zeta_1,\ldots,\zeta_m\in k^\times$ let $\mathsf M$ be the least common
multiple of the orders of those quotients $\zeta_i/\zeta_j$ that are roots of
unity ($\mathsf M=1$ if there are none), or any multiple of it. After equal
bases are combined, every sequence
\[
 l\longmapsto\sum_jr_j(\mathsf Ml+b)\,\zeta_j^{\mathsf Ml+b},\qquad r_j\in k[U],
\]
is either identically zero or nondegenerate. The same $\mathsf M$ serves for
every choice of the polynomials $r_j$ and every $b$.
\end{lemma}

\begin{proof}
The bases in $l$ are $\zeta_j^{\mathsf M}$. If $(\zeta_i/\zeta_j)^{\mathsf M}$
is a root of unity, so is $\zeta_i/\zeta_j$, whose order then divides
$\mathsf M$; so $\zeta_i^{\mathsf M}$ and $\zeta_j^{\mathsf M}$ are equal or
have a nontorsion quotient. Combine equal bases.
\end{proof}

The uniformity in the last sentence is essential below. There may be
infinitely many Hahn coefficient sequences to consider, but they all use the
same finite list of residue bases: the proof never takes a maximum of
infinitely many unrelated exceptional zero sets or periods.

\subsection{The valuation theorem}\label{hol:ud:sub:valuation}

\begin{lemma}[A common support envelope for unit orbits; source~18]
\label{hol:ud:lem:envelope}
Let $R_j\in\KK[U]$ and let $\lambda_j=\zeta_j(1+\tau_j)$ be units,
$1\le j\le m$. There is one well-ordered set $\Sigma\subseteq\Gamma$
containing the support of
\[
 \mathcal Y(n)=\sum_jR_j(n)\lambda_j^n
\]
for every $n\ge0$, and for each $\gamma\in\Sigma$ there are polynomials
$r_{j,\gamma}\in k[U]$ with
\begin{equation}\label{hol:ud:eq:coefficientep}
 [t^\gamma]\,\mathcal Y(n)=\sum_j\zeta_j^n\,r_{j,\gamma}(n)\qquad(n\ge0).
\end{equation}
\end{lemma}

\begin{proof}
Let $C$ be the union of the supports of the coefficients of all $R_j$ and $S$
the union of the supports of the nonzero $\tau_j$: a finite union of
well-ordered sets, and a well-ordered subset of $\Gamma_{>0}$. If all $R_j$
vanish take $\Sigma=\varnothing$; otherwise $\Sigma=C+M(S)$, well ordered by
Lemma~\ref{hol:lem:support}. For $n\ge0$,
$(1+\tau_j)^n=\sum_{i=0}^n\binom ni\tau_j^i$, and every resulting support lies
in $\Sigma$. Fix $\gamma\in\Sigma$ and write $R_j(U)=\sum_a\rho_{j,a}U^a$. The
coefficient in question is
\begin{equation}\label{hol:ud:eq:coefficientexpansion}
 \sum_{j,a,i}\zeta_j^nn^a\binom ni\,[t^\gamma]\bigl(\rho_{j,a}\tau_j^i\bigr).
\end{equation}
Although $i$ is displayed as unbounded, only finitely many $i$ and finitely
many support decompositions contribute at the fixed $\gamma$, by
Lemma~\ref{hol:lem:support}\,(1) and~(2), and the finite list does not depend
on $n$. Each $U^a\binom Ui$ is a polynomial over $k$, since $i!$ is
invertible, and $\binom ni=0$ for $i>n$, so extending the sum in $i$ changes no
value at a nonnegative integer. Collecting the finitely many polynomials gives
\eqref{hol:ud:eq:coefficientep}.
\end{proof}

No bound on $\deg r_{j,\gamma}$ uniform in $\gamma$ is asserted, and none need
hold. The period of Lemma~\ref{hol:ud:lem:period}, by contrast, does not
depend on those degrees.

\begin{theorem}[Eventual periodicity at unit valuation; source~18]
\label{hol:ud:thm:unitvalues}
In Lemma~\ref{hol:ud:lem:envelope}, choose $\mathsf M$ for the residues
$\zeta_j$ by Lemma~\ref{hol:ud:lem:period}. For each
$b\in\{0,\ldots,\mathsf M-1\}$ exactly one of the following holds:
\begin{enumerate}[label=\textup{(\roman*)}]
\item $\mathcal Y(\mathsf Ml+b)=0$ for every $l\ge0$;
\item there is $\beta_b\in\Gamma$ with $v(\mathcal Y(\mathsf Ml+b))=\beta_b$
for all sufficiently large $l$.
\end{enumerate}
If the representation is nonzero and no quotient $\lambda_i/\lambda_j$ of the
actual multipliers is a root of unity, only (ii) occurs; then
$\mathcal Y(n)\ne0$ for all large $n$, and its eventual valuations lie in a
finite subset of $\Gamma$.
\end{theorem}

\begin{proof}
Fix $b$. For every $\gamma\in\Sigma$ the restriction of
\eqref{hol:ud:eq:coefficientep} to $n=\mathsf Ml+b$ is the zero sequence or a
nondegenerate exponential polynomial (Lemma~\ref{hol:ud:lem:period}), and in
the second case it has only finitely many zeros
(Lemma~\ref{hol:cf:lem:exppoly}). If every such coefficient sequence is
identically zero, (i) holds. Otherwise the set of $\gamma\in\Sigma$ for which
$l\mapsto[t^\gamma]\mathcal Y(\mathsf Ml+b)$ is not identically zero has a
least element $\beta_b$. Every coefficient below $\beta_b$ vanishes
\emph{identically}, at every index of the progression, not merely eventually;
at $\beta_b$ there are finitely many exceptional zeros, and beyond them
$\beta_b$ is exactly the least exponent of $\mathcal Y(\mathsf Ml+b)$.

For the last assertion, the restriction of $\mathcal Y$ to a progression has
the bases $\lambda_j^{\mathsf M}$ and the nonzero polynomial factors
$\lambda_j^bR_j(\mathsf MU+b)$. The hypothesis on the actual quotients makes
the bases distinct, and the independence statement in the proof of
Lemma~\ref{hol:cf:lem:exppoly}, applied over $\KK$, excludes an identically
zero restriction. The finitely many progressions have one common tail.
\end{proof}

\begin{remark}[Why a naive coefficientwise argument would be invalid;
source~18]\label{hol:ud:rem:uniform}
It would not be enough to say that every coefficient sequence is eventually
stable and then to take a maximum of infinitely many exceptional indices: such
a maximum need not exist. The proof selects the least coefficient sequence
that is not an identity; all earlier ones vanish at \emph{every} index of the
progression, and only one finite exceptional set is used on each progression.
\end{remark}

\begin{corollary}[The bivariate unit-orbit estimate; source~18]
\label{hol:ud:cor:bivariate}
Let $q\in\KK^\times$ have $v(q)=0$ and infinite multiplicative order. For every
nonzero $P(U,X)\in\KK[U,X]$, $P(n,q^n)$ is eventually nonzero and
$v(P(n,q^n))$ is eventually periodic. If $\res(q)$ is not a root of unity the
valuation is eventually constant; if $\res(q)$ has finite order $h$, a period
dividing $h$ suffices.
\end{corollary}

\begin{proof}
Write $P(U,X)=\sum_jP_j(U)X^j$. The multipliers $q^j$ with $P_j\ne0$ have no
root-of-unity quotient, since $q$ is nontorsion, and their residues are powers
of $\res(q)$. If $\res(q)$ is not a root of unity no nontrivial residue
quotient is one, and $\mathsf M=1$; if it has order $h$, $\mathsf M=h$ is
admissible in Lemma~\ref{hol:ud:lem:period}. Apply
Theorem~\ref{hol:ud:thm:unitvalues}.
\end{proof}

This is the estimate named in the status of Question~\ref{hol:q:mixedunit}.
It extends Theorem~\ref{hol:thm:unitorbit}, which treats $P(q^n)$ for a
univariate $P$, to the bivariate orbit $P(n,q^n)$; the univariate case of
part~(b) there has the elementary proof printed there, and the parity example
of Example~\ref{hol:ex:unitperiod} shows that periodicity, not constancy, is
the right conclusion.

\begin{remark}[An elementary subcase; source~18]
\label{hol:ud:rem:elementary}
When $\res(q)$ has finite order $h$, every coefficient sequence on a class
modulo $h$ in the preceding proof is an ordinary polynomial in the new index
times a nonzero constant, and the only zero-set input needed is the
finiteness of the roots of a nonzero polynomial. The full
Skolem--Mahler--Lech theorem is needed for general residue bases, not for this
subcase. The actual $q$ must still be nontorsion, to exclude a zero
progression.
\end{remark}

\begin{lemma}[Eventually ordered affine functions; source~18]
\label{hol:ud:lem:lines}
For finitely many functions $n\mapsto\beta_i+n\varpi_i$ into an ordered abelian
group, the pairwise comparisons stabilize as $n$ runs through $\N$. If the
slopes $\varpi_i$ are pairwise distinct, all comparisons are eventually strict
and one index is the unique eventual minimum. The same holds along any fixed
arithmetic progression.
\end{lemma}

\begin{proof}
A difference $\beta+n\varpi$ has constant sign if $\varpi=0$, and is strictly
monotone, crossing zero at most once, otherwise. Choose one tail on which all
finitely many comparisons have stabilized. On a progression the slope becomes
a positive integer multiple and only the intercept changes.
\end{proof}

This is the argument inside Lemma~\ref{hol:lem:affine}, stated for an
arbitrary finite list; like that lemma it uses no Archimedean comparison.

\begin{theorem}[Periodic-affine valuations of Hahn exponential polynomials;
source~18]\label{hol:ud:thm:affine}
Let $R_j\in\KK[U]$, $\lambda_j\in\KK^\times$ and
$\mathcal Y(n)=\sum_jR_j(n)\lambda_j^n$. There is $\mathsf M\ge1$ such that for
each class $b$ modulo $\mathsf M$ either
\begin{equation}\label{hol:ud:eq:zeroprogression}
 \mathcal Y(n)=0\quad\text{for every }n\ge0\text{ with }n\equiv b
 \pmod{\mathsf M},
\end{equation}
or there are $\beta_b\in\Gamma$ and $\varpi_b\in\{v(\lambda_1),\ldots,
v(\lambda_m)\}$ with
\begin{equation}\label{hol:ud:eq:affine}
 v(\mathcal Y(n))=\beta_b+n\varpi_b\quad\text{for all sufficiently large }
 n\equiv b\pmod{\mathsf M}.
\end{equation}
If the representation is nonzero and no quotient $\lambda_i/\lambda_j$ is a
root of unity, no zero progression occurs; in particular $\mathcal Y(n)$ is
eventually nonzero. $\mathsf M$ may be chosen from the residues of the
normalized multipliers $t^{-v(\lambda_j)}\lambda_j$ alone, independently of the
coefficients of the $R_j$.
\end{theorem}

\begin{proof}
Put $\hat\lambda_j=t^{-v(\lambda_j)}\lambda_j$, a unit, and group the indices
with equal valuation:
\[
 \mathcal Y(n)=\sum_\varpi t^{n\varpi}\mathcal Y_\varpi(n),\qquad
 \mathcal Y_\varpi(n)=\sum_{j:\,v(\lambda_j)=\varpi}R_j(n)\hat\lambda_j^n .
\]
Choose one period $\mathsf M$ for all residues $\res(\hat\lambda_j)$ by
Lemma~\ref{hol:ud:lem:period}; it is admissible for each group. On a class
$b$, Theorem~\ref{hol:ud:thm:unitvalues} makes every $\mathcal Y_\varpi$ either
identically zero there or of eventually constant valuation $\beta_{\varpi,b}$.
If all groups vanish identically on the class,
\eqref{hol:ud:eq:zeroprogression} holds. Otherwise discard the zero groups; the
rest have, on a common tail of the class, the valuations
$n\varpi+\beta_{\varpi,b}$ with pairwise distinct slopes, and
Lemma~\ref{hol:ud:lem:lines} gives a unique least one on a further tail. A
unique least-valuation summand cannot cancel, which is
\eqref{hol:ud:eq:affine}. The assertion without zero progressions follows
from the independence statement of Lemma~\ref{hol:cf:lem:exppoly} over $\KK$
on a progression, as in Theorem~\ref{hol:ud:thm:unitvalues}.
\end{proof}

The smallest slope need not win eventually: in the third example of
Remark~\ref{hol:rem:shortcut}, which is source~18's Example~4.8, the intercept
$e_1$ dominates every multiple of the slope $e_0$. The proof uses the eventual
winner of finitely many affine functions, not the rule ``take the smallest
slope''.

\paragraph{What is, and is not, effective (source~18).}
The theorem is an existence statement for a finite valuation description.
Given arbitrary Hahn coefficients it provides no algorithm for the first
coefficient sequence that is not an identity, for the eventual winning
indices, or for the last exceptional zero. In the finite-order residue case
the relevant zero sets are polynomial roots, but selecting the coefficients
still needs exact support information. No decision procedure for arbitrary
surreal inputs follows.

\section{Mixed rigidity and the relative-scale dichotomy}
\label{hol:ud:sec:mixed}

\subsection{Operator normal form and its recurrence}
\label{hol:ud:sub:operator}

The operators are $D_z$ and $\sigma_\lambda$, with
$D_z\sigma_\lambda=\lambda\,\sigma_\lambda D_z$. Write a finite operator as
\begin{equation}\label{hol:ud:eq:L}
 L=\sum_{\ell=1}^m\sum_{r=0}^{R}p_{\ell,r}(z)\,D_z^r\sigma_{\lambda_\ell},\qquad
 p_{\ell,r}(z)=\sum_{i=0}^{d_L}\varkappa_{\ell,r,i}z^i\in\KK[z],
\end{equation}
with $\lambda_\ell\ne0$ and some $\varkappa_{\ell,r,i}\ne0$; $D_z^r\sigma_\lambda$
is kept in that order, not silently replaced by $\sigma_\lambda D_z^r$. Since
$D_z^r\sigma_\lambda=\lambda^r\sigma_\lambda D_z^r$, the terms are the jets
$\lambda^r\jet\lambda rf$ of \eqref{hol:cf:eq:jet}, and \eqref{hol:eq:mixed} is
the case $\lambda_\ell=q^\ell$.

\begin{lemma}[Exact shift polynomials; source~18]\label{hol:ud:lem:operator}
Let $f=\sum_na_nz^n$. For an output degree $\bar n\ge d_L$, the term
$\varkappa_{\ell,r,i}z^iD_z^r\sigma_{\lambda_\ell}f$ contributes
\begin{equation}\label{hol:ud:eq:termcoefficient}
 \varkappa_{\ell,r,i}\,\fall{\bar n+r-i}{r}\,\lambda_\ell^{\bar n+r-i}\,
 a_{\bar n+r-i}
\end{equation}
to $[z^{\bar n}]Lf$. Let $s_0$ be the least active shift $r-i$ and put
$n=\bar n+s_0$. After grouping equal coefficient indices, the high
coefficients of $Lf=\mathsf h$ read
\begin{equation}\label{hol:ud:eq:recurrence}
 c_0(n)a_n+\sum_{j=1}^Jc_j(n)a_{n+j}=[z^{n-s_0}]\,\mathsf h,
\end{equation}
with
\begin{align}
 c_j(n)&=\sum_{\ell=1}^m\lambda_\ell^n\,R_{j,\ell}(n),\label{hol:ud:eq:cj}\\
 R_{j,\ell}(U)&=\lambda_\ell^j\sum_{r-i=s_0+j}\varkappa_{\ell,r,i}\,
 \fall{U+j}{r}.\label{hol:ud:eq:Rjl}
\end{align}
The active $j$ run from $0$ to a finite $J$; for each active $j$ some
$R_{j,\ell}$ is nonzero, in particular some $R_{0,\ell}$. The right side
vanishes once $\bar n>\deg\mathsf h$.
\end{lemma}

\begin{proof}
The monomial $a_Nz^N$ becomes $a_N\lambda_\ell^N\fall Nrz^{N-r}$ under
$D_z^r\sigma_{\lambda_\ell}$, and multiplication by $z^i$ gives degree
$N-r+i$; setting it equal to $\bar n$ gives \eqref{hol:ud:eq:termcoefficient},
and $\bar n\ge d_L$ ensures $N\ge r$. With $r-i=s_0+j$ one has $N=n+j$, which
gives \eqref{hol:ud:eq:cj}--\eqref{hol:ud:eq:Rjl}. For fixed $j$ and $\ell$
distinct contributing $r$ give the monic polynomials $\fall{U+j}r$ of distinct
degrees, so a nonempty combination with nonzero coefficients is nonzero, and
for fixed $r$ and $j$ there is only one $i$.
\end{proof}

For $\lambda_\ell=q^\ell$ this is the conversion \eqref{hol:eq:mixedshift} of
Theorem~\ref{hol:thm:mixed}, grouped by multiplier instead of by powers of $X$.

\begin{corollary}[No hidden zero operator; source~18]
\label{hol:ud:cor:faithful}
If the $\lambda_\ell$ are distinct, a normal form \eqref{hol:ud:eq:L} with some
nonzero coefficient is a nonzero operator on $\KK[[z]]$.
\end{corollary}

\begin{proof}
At the least active shift, \eqref{hol:ud:eq:cj} is a nonzero
polynomial--exponential representation with distinct bases, so $c_0(n)\ne0$
for arbitrarily large $n$ by the independence statement of
Lemma~\ref{hol:cf:lem:exppoly}; for such $n$ the coefficient of degree $n-s_0$
in $L(z^n)$ is $c_0(n)\ne0$.
\end{proof}

\subsection{Relative scales and rigidity}\label{hol:ud:sub:rigidity}

\begin{definition}[Relative scale; source~18]\label{hol:ud:def:rs}
For a nonempty finite list of multipliers put
\[
 \operatorname{rs}(\lambda_1,\ldots,\lambda_m)=\max_{i,j}\abs{v(\lambda_i)-
 v(\lambda_j)}\in\Gamma_{\ge0}.
\]
It is $0$ for one multiplier. It is \emph{noncofinal} if it is $0$ or is not
an order unit.
\end{definition}

If $\operatorname{rs}>0$ is not an order unit, \eqref{hol:eq:dominates}
supplies $\gamma>0$ with $n\operatorname{rs}<\gamma$ for every $n\ge0$. For a
finite list, noncofinality of every $\abs{v(\lambda_i/\lambda_j)}$ is
equivalent to noncofinality of their maximum.

\begin{theorem}[Rigidity for noncofinal relative scales; source~18]
\label{hol:ud:thm:rigidity}
Suppose no quotient $\lambda_i/\lambda_j$, $i\ne j$, is a root of unity and
$\operatorname{rs}(\lambda_1,\ldots,\lambda_m)$ is noncofinal. Let $L$ be a
nonzero normal form \eqref{hol:ud:eq:L}. For every $\bar e\in\N$ there are
$N\ge1$ and $B\in\Gamma_{\ge0}$ such that, for every $\mathsf h\in\KK[z]$ with
$\deg\mathsf h\le\bar e$ and every formal solution of $Lf=\mathsf h$:
\begin{enumerate}[label=\textup{(\roman*)}]
\item strong evaluation at one $x\ne0$ with $v(x)<-B$ forces $\deg f<N$;
\item every strongly entire solution is a polynomial of degree less than $N$;
\item every nonpolynomial formal solution fails strong evaluation throughout
that common exterior region.
\end{enumerate}
The polynomial conclusion also holds for rational operator coefficients and a
rational right side, read as an identity of formal Laurent series. By
Corollary~\ref{hol:cor:boundedcost} the region in \textup{(i)} and
\textup{(iii)} may be taken closed, $v(x)\le-B$ \emph{(merge)}.
\end{theorem}

\begin{proof}
Apply Lemma~\ref{hol:ud:lem:operator}. For each active $j$, $c_j(n)$ is a
nonzero exponential polynomial with bases among the $\lambda_\ell$, so
Theorem~\ref{hol:ud:thm:affine} gives one period $\mathsf M$ and one tail on
which, for $n\equiv b\pmod{\mathsf M}$,
\begin{equation}\label{hol:ud:eq:cjvaluations}
 v(c_j(n))=\beta_{j,b}+n\varpi_{j,b},\qquad
 \varpi_{j,b}\in\{v(\lambda_1),\ldots,v(\lambda_m)\},
\end{equation}
and all these $c_j(n)$, in particular $c_0(n)$, are nonzero there. Inactive
$j$ give identically zero coefficients and are discarded; if no $j\ge1$ is
active the recurrence terminates directly. Otherwise let $B_0\ge0$ bound the
finitely many differences $\beta_{0,b}-\beta_{j,b}$, $j\ge1$. Then
\[
 v(c_0(n))-v(c_j(n))=\beta_{0,b}-\beta_{j,b}+n(\varpi_{0,b}-\varpi_{j,b})
 \le B_0+n\operatorname{rs}(\lambda_1,\ldots,\lambda_m).
\]
If the relative scale is $0$ put $B=B_0$; otherwise put $B=B_0+\gamma$ with
$\gamma$ from \eqref{hol:eq:dominates}. This is a uniform cost bound. Enlarge
$N$ past the finite valuation exceptions and so that $n-s_0\ge d_L$ and
$n-s_0>\bar e$. Lemma~\ref{hol:lem:escape} with
Theorem~\ref{hol:thm:escapeexclusion} and Corollary~\ref{hol:cor:boundedcost}
proves (i)--(iii): a nonzero $a_n$ with $n\ge N$ starts an escape chain whose
evaluated leading exponents do not increase at such $x$. For rational data,
multiply the equation on the left by a common nonzero polynomial denominator:
the operator stays a nonzero normal form, the right side becomes a
polynomial, and no denominator is differentiated.
\end{proof}

\begin{remark}[What the theorem contains; merge]\label{hol:ud:rem:contains}
With $m=1$ and $\lambda_1=1$ the relative scale is $0$, and
Theorem~\ref{hol:ud:thm:rigidity} is Theorem~\ref{hol:main:ode}, with its
uniform exterior obstruction and degree bound, over every
characteristic-zero $k$. With $R=0$ and $\lambda_\ell=q^\ell$ it is the
negative direction of Theorem~\ref{hol:main:q}: the relative scale is a
multiple of $\abs{v(q)}$, and is noncofinal exactly when (ii) there fails. With
$\lambda_\ell=q^\ell$ and $v(q)\ne0$ noncofinal it is
Theorem~\ref{hol:thm:mixed}. The proofs of this report reach those three
theorems through Theorem~\ref{hol:thm:unitorbit},
Lemma~\ref{hol:lem:affine} and Lemma~\ref{hol:lem:bivariate}; source~18 reaches
all of them, and the unit case, through Theorem~\ref{hol:ud:thm:affine}. The
comparison with Theorem~\ref{hol:cf:main:entire}\,(b) is this: that theorem
excludes every algebraic relation, linear or not, among the mixed jets when
$\Gamma$ has no order unit, which is then the case of every list of
multipliers; Theorem~\ref{hol:ud:thm:rigidity} excludes only linear relations,
but for every $\Gamma$ with a noncofinal relative scale, with a uniform degree
bound and exterior region. Neither contains the other.
\end{remark}

\begin{corollary}[Mixed equations at unit valuation; source~18]
\label{hol:ud:cor:unitrigidity}
Let $q\in\KK^\times$ be nontorsion with $v(q)=0$. Every strongly entire
solution of a nonzero equation
$\sum_{r,\ell}p_{r,\ell}(z)D_z^r\sigma_q^\ell f(z)=\mathsf h(z)$ with
$p_{r,\ell},\mathsf h\in\KK[z]$ is a polynomial, for every $\Gamma$. Rational
coefficients are allowed as in Theorem~\ref{hol:ud:thm:rigidity}, and each
fixed equation has a uniform degree bound and a uniform exterior obstruction
for its nonpolynomial formal solutions.
\end{corollary}

\begin{proof}
The distinct multipliers $q^\ell$ that occur have nontorsion quotients and
valuation zero, so the relative scale is $0$; apply
Theorem~\ref{hol:ud:thm:rigidity}. Negative powers of $q$ change nothing.
\end{proof}

For $\mathsf h=0$ and $k=\C$ this is Question~\ref{hol:q:mixedunit}, answered
positively, with no order-unit hypothesis. The decisive estimate is
Corollary~\ref{hol:ud:cor:bivariate}; no coefficient-field theorem and no
nonlinear method is used.

\begin{corollary}[Mixed linear independence; source~18]
\label{hol:ud:cor:linind}
Let $f\in\EK$ be nonpolynomial. For every finite list of multipliers with no
root-of-unity quotient and a noncofinal relative scale, the family
$\{D_z^r\sigma_{\lambda_j}f:1\le j\le m,\ r\ge0\}$ is linearly independent over
$\KK(z)$; equivalently, so is $\{\jet{\lambda_j}rf\}_{j,r}$.
\end{corollary}

\begin{proof}
Clearing denominators in a nontrivial relation gives a nonzero normal form
annihilating $f$, and Theorem~\ref{hol:ud:thm:rigidity} makes $f$ a
polynomial. The nonzero scalars $\lambda_j^r$ do not affect independence.
\end{proof}

This is linear independence of formal functions over a field of rational
functions. It is not algebraic independence of values at a point, and it does
not exclude polynomial relations among these functions; for those see
Theorem~\ref{hol:cf:main:entire}\,(b) without an order unit and
Corollary~\ref{hol:td:cor:sharp} with one.

\subsection{The dichotomy for a finite set of multipliers}
\label{hol:ud:sub:dichotomy}

\begin{lemma}[A mixed equation for partial theta; source~18]
\label{hol:ud:lem:theta}
If $p\in\KK^\times$ and $v(p)>0$ is an order unit, then $\Theta_p$ of
\eqref{hol:eq:theta} is a nonpolynomial element of $\EK$ and satisfies
\begin{equation}\label{hol:ud:eq:thetamixed}
 D_z\Theta_p-\sigma_p\Theta_p-z\,D_z\sigma_p\Theta_p=0 .
\end{equation}
\end{lemma}

\begin{proof}
Entireness and nonpolynomiality are Theorem~\ref{hol:thm:thetadomain}, whose
order-unit half source~18 proves by the route of
Remark~\ref{hol:rem:thetasecond}. Differentiate \eqref{hol:eq:thetafirst},
$\Theta_p(z)=1+z\Theta_p(pz)$, with $p$ constant for $D_z$: the product rule
gives $D_z\Theta_p=\sigma_p\Theta_p+z\,D_z(\sigma_p\Theta_p)$.
\end{proof}

\begin{theorem}[Exact relative-scale dichotomy; source~18]
\label{hol:ud:thm:dichotomy}
Let $\lambda_1,\ldots,\lambda_m\in\KK^\times$ with no quotient
$\lambda_i/\lambda_j$, $i\ne j$, a root of unity. The following are
equivalent.
\begin{enumerate}[label=\textup{(\roman*)}]
\item There are a nonpolynomial $f\in\EK$ and a nonzero operator
\eqref{hol:ud:eq:L} with $Lf=0$.
\item $\abs{v(\lambda_i/\lambda_j)}$ is an order unit for some $i,j$.
\item There are such $f$ and $L$ using only two of the $\lambda_\ell$ and
derivatives of order at most one.
\end{enumerate}
If \textup{(ii)} fails, every nonpolynomial $f\in\EK$ has the independence
property of Corollary~\ref{hol:ud:cor:linind} for the whole list.
\end{theorem}

\begin{proof}
If (ii) fails the relative scale is noncofinal and
Theorem~\ref{hol:ud:thm:rigidity} excludes (i). Assume (ii) and order the pair
so that $p=\lambda_j/\lambda_i$ has positive order-unit valuation. Put
$f(z)=\Theta_p(z/\lambda_i)$, nonpolynomial and strongly entire since
dilation preserves $\EK$. Then $\sigma_{\lambda_i}f=\Theta_p$ and
$\sigma_{\lambda_j}f=\sigma_p\Theta_p$, and \eqref{hol:ud:eq:thetamixed} becomes
\begin{equation}\label{hol:ud:eq:twoprescribed}
 \bigl(D_z\sigma_{\lambda_i}-\sigma_{\lambda_j}-zD_z\sigma_{\lambda_j}\bigr)f=0,
\end{equation}
a nonzero normal form by Corollary~\ref{hol:ud:cor:faithful}. This is (iii),
and (iii)$\Rightarrow$(i) is immediate.
\end{proof}

\begin{proof}[Proof of Theorem~\ref{hol:ud:main:mixed}]
Part~(a) is Theorem~\ref{hol:ud:thm:rigidity}, part~(b) is
Theorem~\ref{hol:ud:thm:dichotomy}, and the last assertion is
Corollary~\ref{hol:ud:cor:unitrigidity}.
\end{proof}

\begin{remark}[The quantifiers matter; source~18]\label{hol:ud:rem:quantifiers}
Theorem~\ref{hol:ud:thm:dichotomy} classifies the existence of \emph{at least
one} nonzero mixed equation supported in a prescribed finite set of
multipliers. It does not say that every operator using a cofinal quotient has
a nonpolynomial entire solution, it does not decide the solvability of a
specified equation, and it does not classify simultaneous solutions of
several independently prescribed equations. The positive construction may
give zero coefficients to all but two of the listed multipliers. This is the
precise sense in which it bears on Question~\ref{hol:q:severaldilations}.
\end{remark}

\begin{corollary}[Cyclic mixed equations; source~18]\label{hol:ud:cor:cyclic}
Let $q\in\KK^\times$ be nontorsion. A nonpolynomial strongly entire solution
of some nonzero polynomial-coefficient equation in $D_z$ and nonnegative
powers of $\sigma_q$ exists if and only if $v(q)\ne0$ and $\abs{v(q)}$ is an
order unit. In the negative case every such equation, also with polynomial
forcing, has only polynomial strongly entire solutions. The same holds with
negative powers of $\sigma_q$ admitted.
\end{corollary}

\begin{proof}
The valuation differences of a finite list of powers of $q$ are integer
multiples of $v(q)$, and a nonzero integer multiple of a positive element is
an order unit exactly when the element is. If $v(q)=0$ or $\abs{v(q)}$ is
noncofinal, Theorem~\ref{hol:ud:thm:rigidity} applies to every finite list.
In the positive case apply Theorem~\ref{hol:ud:thm:dichotomy} to $\{1,q\}$;
its equation uses only $\sigma_1$ and $\sigma_q$, whatever the sign of $v(q)$.
\end{proof}

With derivatives removed, the existence half is Theorem~\ref{hol:main:q},
whose witness is a pure dilation equation. So adding finite derivative orders
creates \emph{no new existence case} (merge's reading): the criterion of
Theorem~\ref{hol:main:q} is also the criterion for mixed equations in $D_z$
and $\sigma_q$, and the unit case missing before this part is now included.

\paragraph{Why a common absolute scale does not matter (source~18).}
Replacing every $\lambda_j$ by $c\lambda_j$ with one $c\ne0$ gives
$\sum p_{j,r}D_z^r\sigma_{c\lambda_j}f=\sum p_{j,r}D_z^r\sigma_{\lambda_j}
(\sigma_cf)$, and $\sigma_c$ is a bijection of $\EK$ preserving polynomiality,
so the existence problem is unchanged. An absolute-scale criterion would
therefore be wrong: one multiplier of arbitrarily large valuation gives an
ordinary differential equation for the entire series $\sigma_\lambda f$ and
cannot by itself produce a nonpolynomial strongly entire solution.

\begin{proposition}[The top Archimedean scale; source~18]
\label{hol:ud:prop:top}
Suppose $\Gamma$ has an order unit and put
$\Delta^{\mathrm{top}}=\{\gamma\in\Gamma:\gamma=0$ or $\abs\gamma$ is not an
order unit$\}$. Then $\Delta^{\mathrm{top}}$ is the unique maximal proper
convex subgroup of $\Gamma$ and $\Gamma/\Delta^{\mathrm{top}}$ is Archimedean.
The negative case of Theorem~\ref{hol:ud:thm:dichotomy} holds exactly when all
$v(\lambda_j)$ have the same image in $\Gamma/\Delta^{\mathrm{top}}$.
\end{proposition}

\begin{proof}
For $\gamma_1,\gamma_2\in\Delta^{\mathrm{top}}$ not both zero,
$\abs{\gamma_1+\gamma_2}\le2\rho$ with $\rho=\max(\abs{\gamma_1},\abs{\gamma_2})$
not an order unit, and $2\rho$ is then not one either; negatives and convexity are immediate. An order unit lies
outside, so the subgroup is proper, and a convex subgroup containing an order
unit is $\Gamma$, since every element lies between two integer multiples of
it; so every proper convex subgroup lies in $\Delta^{\mathrm{top}}$. For
positive classes $\bar\xi_1,\bar\xi_2$ of the quotient a positive
representative $\xi_1$ is an order unit, and $n\xi_1>\xi_2$ for some $n$. Finally
$v(\lambda_i)-v(\lambda_j)\in\Delta^{\mathrm{top}}$ says exactly that its
absolute value is not an order unit.
\end{proof}

These ordered-group facts are elementary background, not a new
classification. For an order unit $\mu$, $\Delta^{\mathrm{top}}$ is the kernel
$\{\gamma:n\abs\gamma<\mu\text{ for all }n\ge1\}$ of the real-valued
coarsening cited in Remark~\ref{hol:rem:coarsening}. The interpretation:
under the no-root-of-unity restriction, a nonpolynomial strongly entire mixed
relation exists exactly when the multipliers are separated at the top
Archimedean scale.

\subsection{Surreal, surcomplex and omnific coefficients}
\label{hol:ud:sub:surreal}

Through \eqref{hol:eq:embedding} every statement above is a statement about
series with surreal or surcomplex coefficients in one fixed set-sized
workspace $K_\Gamma\subseteq\No[i]$, with $\Gamma\subseteq\No$.

\begin{corollary}[A surcomplex no-annihilator theorem; source~18]
\label{hol:ud:cor:surcomplex}
Let $\Gamma\subseteq\No$ be a nonzero set-sized subgroup, $f\in K_\Gamma[[z]]$
nonpolynomial and strongly entire on $K_\Gamma$, and $q\in K_\Gamma^\times$
nontorsion with $v(q)=0$. There is no nonzero finite relation
$\sum_{r,\ell}R_{r,\ell}(z)(D_z^rf)(q^\ell z)=0$ with
$R_{r,\ell}\in K_\Gamma(z)$. The same holds for every finite list of
multipliers with a noncofinal relative scale and no root-of-unity quotient.
\end{corollary}

\begin{proof}
Corollary~\ref{hol:ud:cor:linind}; the factors $q^{r\ell}$ relating
$D_z^r\sigma_q^\ell f$ to $(D_z^rf)(q^\ell z)$ are absorbed in the rational
functions.
\end{proof}

\paragraph{An example with an order unit (source~18).}
Let $\Gamma=\Q\mu\oplus\Q\epsilon$ with $\mu>n\epsilon>0$ for every $n\ge1$,
realized by $\mu=\omega$, $\epsilon=1$. Then $\mu$ is an order unit and
$\epsilon$ is not. Put
\[
 G_\mu(z)=\sum_{n\ge0}t^{n^2\mu}z^n=\sum_{n\ge0}\omega^{-n^2\omega}z^n,\qquad
 q=1+t^\epsilon=1+\omega^{-1}.
\]
The valuations $n^2\mu+nv(x)$ are cofinal for every $x\ne0$, so $G_\mu$ is
nonpolynomial and strongly entire by Lemma~\ref{hol:lem:certificate}\,(a). The
unit $q$ is nontorsion, since $(1+t^\epsilon)^n-1$ has leading term
$nt^\epsilon$. By Corollary~\ref{hol:ud:cor:linind}, all finite collections
of derivatives of $G_\mu(z),G_\mu(qz),G_\mu(q^2z),\ldots$ are linearly
independent over $K_\Gamma(z)$. This lies in the order-unit regime, where
Theorem~\ref{hol:cf:main:entire}\,(b) does not apply. At a different
dilation scale there is a relation: $G_\mu(z)=\Theta_{p}(t^\mu z)$ with
$p=t^{2\mu}$, so $G_\mu$ satisfies the dilation equation \eqref{hol:eq:thetahom}
rescaled; excluding every unit-dilation mixed equation does not exclude one at
a cofinal nonzero dilation valuation.

\paragraph{Omnific coefficients (source~18).}
Under \eqref{hol:eq:embedding}, omnific coefficients in the workspace have
Hahn support in $\Gamma_{\le0}$, with an integer coefficient at exponent $0$;
Gaussian omnific coefficients have a Gaussian-integer constant term. The
proofs place no boundedness or positivity condition on operator coefficients,
so coefficients in $\Oz\cap K_\Gamma$ or $\Oz[i]\cap K_\Gamma$ are allowed, in
the real or complex workspace; infinite omnific coefficients do not evade the
theorem. Every set of surreal coefficients lies in some set-sized workspace:
the union of their normal-form supports is a set, the group it generates is
set sized, and one nonzero exponent may be adjoined if needed. This chooses a
common workspace; it does \emph{not} give strong entireness there, which
remains a hypothesis and can be lost on enlarging the group. No result on
prime omnific integers, their integral closure, automorphisms or
Diophantine definability is claimed; the arithmetic connection is only that
full, possibly infinite normal forms are legitimate coefficients.

\paragraph{Enlargement (source~18).}
Enlarge the group of the example by $e_*$ with $e_*>n\mu$ for every $n$, for
instance $e_*=\omega^2$. At $x=t^{-e_*}$ the $n$th term of $G_\mu$ has
valuation $n^2\mu-ne_*$, with successive differences $(2n+1)\mu-e_*<0$, so the
evaluation family is not strongly summable
(Lemma~\ref{hol:lem:nonincreasing}): $G_\mu$ is entire on the original
workspace but not on this enlargement, let alone on $\No[i]$. This is the
computation of Section~\ref{hol:sub:notallsurcomplex} and after
Corollary~\ref{hol:nl:cor:infiniterank}, and no nonpolynomial series is
entire on $\No[i]$ (Proposition~\ref{hol:ot:prop:fullclass}).

\subsection{Examples and necessary boundaries}\label{hol:ud:sub:examples}

\paragraph{Periodic valuations from a torsion residue.}
Source~18's first example, $q=-(1+t^\epsilon)$ with $v(q^n-1)=0$ for odd $n$
and $=\epsilon$ for even $n\ge1$, is Example~\ref{hol:ex:unitperiod}: even
for one nontorsion parameter, Theorem~\ref{hol:ud:thm:unitvalues} needs
periodicity, not eventual constancy. At $n=0$ the value is zero, an isolated
exceptional index.

\begin{example}[Cancellation below a prescribed coefficient; source~18]
\label{hol:ud:ex:cancellation}
For $q=1+t^\epsilon$ and an integer $d\ge1$ put
$\mathcal Y_d(n)=q^n-\sum_{i=0}^{d-1}\binom nit^{i\epsilon}$, an exponential
polynomial with the bases $q$ and $1$. For $n\ge d$ its leading term is
$\binom ndt^{d\epsilon}$, so $v(\mathcal Y_d(n))=d\epsilon$; for $n<d$ it is
$0$. All coefficient sequences below $d\epsilon$ vanish identically. This is
the first-nonidentity-coefficient selection of
Theorem~\ref{hol:ud:thm:unitvalues} at work: inspecting the residue of
$\mathcal Y_d(n)$ would never find the valuation.
\end{example}

\begin{example}[A large common absolute scale is not enough; source~18]
\label{hol:ud:ex:commonscale}
In the group of the example of Section~\ref{hol:ud:sub:surreal} take
$\lambda_1=t^\mu$, $\lambda_2=t^\mu(1+t^\epsilon)$ and
$\lambda_3=2t^\mu(1+t^{2\epsilon})$. All have the cofinal valuation $\mu$ and
their relative valuations are zero. Their quotients are nontorsion: a
quotient with residue $2$ or $\frac12$ is not a root of unity in
characteristic zero, and $1+t^\epsilon$ is nontorsion by its leading term. So
no nonpolynomial strongly entire series satisfies a nonzero mixed equation
using this list; the common cofinal factor is irrelevant.
\end{example}

\begin{example}[Crossing the cofinality boundary; source~18]
\label{hol:ud:ex:crossing}
In the same group, for $\lambda_1=1$ and $\lambda_2=t^\epsilon$ the quotient
valuation is positive but noncofinal, and all mixed entire solutions are
polynomials. Replacing $\lambda_2$ by $t^\mu$ changes the answer:
$\Theta_{t^\mu}$ is nonpolynomial and entire and satisfies
\[
 D_z\Theta_{t^\mu}(z)-\Theta_{t^\mu}(t^\mu z)
 -zD_z\bigl(\Theta_{t^\mu}(t^\mu z)\bigr)=0 .
\]
So ``nonzero valuation'' is not the boundary, even for a first-order equation
with two multipliers; compare Example~\ref{hol:ex:ranktwo} for pure
dilations.
\end{example}

\paragraph{Actual root-of-unity quotients cannot be ignored (source~18).}
If $k$ contains a root of unity $\zeta\ne1$ of order $h$ and $\mu$ is an order
unit, then $G_\mu(z^h)=\sum_nt^{n^2\mu}z^{hn}$ is nonpolynomial and strongly
entire, and $(\sigma_\zeta-1)G_\mu(z^h)=0$, although $\sigma_\zeta-1$ is not
zero on $\KK[[z]]$ (it sends $z$ to $(\zeta-1)z$) and the relative scale of
$1,\zeta$ is zero. This is the periodic-support construction of
Proposition~\ref{hol:prop:torsion}\,(a) read as a mixed equation. The
hypothesis on root-of-unity quotients in
Theorems~\ref{hol:ud:thm:rigidity} and~\ref{hol:ud:thm:dichotomy} is therefore a
real restriction, not a device to avoid a syntactically zero operator; the
mechanism is an unconstrained residue class of Taylor indices.

\paragraph{Positive characteristic breaks both steps (source~18).}
Over $\mathbb F_{\mathfrak p}((t))$ the unit $q=1+t$ satisfies
$q^{\mathfrak p^a}-1=t^{\mathfrak p^a}$ for $a\ge0$, so $v(q^{\mathfrak p^a}-1)$ is
unbounded and the finite-range conclusion of
Theorem~\ref{hol:ud:thm:unitvalues} fails. The series of
Example~\ref{hol:nl:ex:charp} is nonpolynomial and strongly entire with
derivative zero, while $D_z$ is not zero ($D_zz=1$). So characteristic zero is
essential to the coefficient-polynomial argument and to differential
rigidity; the positive-characteristic zero-set theory~\cite{derksen} is a
different subject.

\begin{example}[No degree bound from the shape of the operator; source~18]
\label{hol:ud:ex:degree}
For $q=2\in k$ and every $a\in\N$, $L_a=\sigma_2-2^a$ has the strongly entire
solution $z^a$. Each $L_a$ has differential order zero, two dilation terms and
constant coefficients, yet the degrees of their polynomial solutions are
unbounded. So the degree bound of Theorem~\ref{hol:ud:thm:rigidity} depends on
the exact operator data, not only on the number of terms, derivative orders
and coefficient degrees; compare Proposition~\ref{hol:nl:prop:degenerate} for
a single nonlinear equation.
\end{example}

The trivial group $\Gamma=0$, where every strongly entire series is a
polynomial by evaluation at $z=1$, is excluded to keep the scale language
meaningful; it is elementary, not an uncovered counterexample
(Section~\ref{hol:sec:fields}).

\subsection{Certificates and computation}\label{hol:ud:sub:certificates}

For a fixed polynomial operator $L$ the proof of
Theorem~\ref{hol:ud:thm:rigidity} reduces to finite data \emph{once the
eventual valuation descriptions are established}: the coefficients
$\varkappa_{\ell,r,i}$, the least active shift $s_0$, the polynomials
$R_{j,\ell}$ of \eqref{hol:ud:eq:Rjl}, one period $\mathsf M$, one common tail
index $N$, the pairs $(\beta_{j,b},\varpi_{j,b})$ of
\eqref{hol:ud:eq:cjvaluations}, either $\operatorname{rs}=0$ or a witness
$\gamma$ for \eqref{hol:eq:dominates}, and a bound on the finitely many
intercept differences. These give $B$, hence the common bad argument
$t^{-B-\rho'}$ for any $\rho'>0$. The support envelope and the
first-nonidentity coefficient are part of the certificate's \emph{proof}, not
information recoverable from a long finite prefix: a late exceptional zero of
an exponential polynomial can invalidate an unjustified choice of $N$.

\begin{proposition}[The entire solution space is finite; source~18]
\label{hol:ud:prop:finitespace}
Under Theorem~\ref{hol:ud:thm:rigidity} the strongly entire solutions of $Lf=0$
form a finite-dimensional $\KK$-vector space inside $\KK[z]_{<N}$; those of
$Lf=\mathsf h$ form the empty set or an affine translate of it, inside the
same space. Once a valid $N$ is supplied, all of them are determined by a
finite linear system over $\KK$.
\end{proposition}

\begin{proof}
Every entire solution lies in $\KK[z]_{<N}$, and every polynomial is entire.
If the coefficients of $L$ have degree at most $d_L$, then $L$ maps
$\KK[z]_{<N}$ into polynomials of degree at most $N-1+d_L$, so $Lf=\mathsf h$ is
finitely many linear equations in the $N$ unknown coefficients. No algorithm
for arbitrary encoded Hahn elements is assumed.
\end{proof}

\paragraph{A formalization route proposed by source~18.}
The report's Lean ledger records modules for escape chains, positive-support
words, partial-theta domains and unit-orbit polynomials, which source~18
names as natural reuse points; no Lean build was run and no declaration is
added. The proposed order is: the operator-to-recurrence conversion and
polynomial--exponential independence; the common support envelope; the
first-nonidentity coefficient argument relative to a clearly named zero-set
theorem; the relative-scale cost bound and the two-dilation theta witness.
The Skolem--Mahler--Lech input should either be proved in the library or
appear openly as a hypothesis of a conditional theorem, never as a hidden
axiom of a result advertised as kernel checked. The finite-order-residue case
(Remark~\ref{hol:ud:rem:elementary}) needs only root finiteness and is a
smaller first target that still includes nontorsion units such as
$1+\omega^{-1}$ in an order-unit workspace.

\begin{remark}[Source~18's nonlinear question; merge]
\label{hol:ud:rem:nonlinearq}
Source~18's Research question~11.1 asks: in a workspace with an order unit
and for a nontorsion unit $q$, must every strongly entire solution of a
nonzero algebraic relation over $\KK(z)$ among finitely many
$(D_z^rf)(q^\ell z)$ be a polynomial, after genuine algebraic identities among
the displayed expressions are excluded? At its pin that was open; in the
current text the answer is no. With an order unit $\mu$ and $p=t^\mu$, the
nonpolynomial strongly entire series $\mathcal T_p$ satisfies a nonzero
algebraic relation among $f,D_zf,D_z^2f,D_z^3f$, the terms with $\ell=0$
(Theorem~\ref{hol:td:thm:ode}). It is no identity among the displayed
expressions: it is a nonzero polynomial in the jet variables, and it fails for
$f=z$, where its coefficient of $z^2$ is $64-192\mathfrak c$, of residue
$48$. So the family of Theorem~\ref{hol:cf:main:sharp}\,(iii) is dependent for
every nontorsion $q$, including $v(q)=0$ (Corollary~\ref{hol:td:cor:sharp}).
Source~18's suggested first test case, a second-order differential polynomial
without dilations whose leading residue expression vanishes identically, is
decided by Theorem~\ref{hol:ot:main:second}\,(a): every strongly entire
solution of an equation of order at most two is a polynomial. What remains
open with an order unit is the classification of
Question~\ref{hol:q:nonlinear}; the question is therefore not printed as an
open question.
\end{remark}

\subsection{Predecessors, repository comparison and priority}
\label{hol:ud:sub:predecessors}

\paragraph{Literature.}
Source~18 imports the Skolem--Mahler--Lech theorem in its characteristic-zero
form, citing Lech~\cite{lech} as the primary source (bibliographic metadata
verified) and Bell's account~\cite{belldoc}, whose accessible text it read,
including the first two pages; it credits Neumann~\cite{neumann} for the
support calculus, whose needed lemma it proves; Gonshor~\cite{gonshor} for
the normal-form embedding (publisher metadata checked, the monograph not
reread); Fuchs--Heintze~\cite{fuchsheintze} as related work on growth of
recurrences over one-variable function fields, whose nondegeneracy hypothesis
and valuation-bound problem differ and which is not an input; and
Derksen~\cite{derksen} as background for a positive-characteristic
direction. It consulted the Wikipedia article on surreal numbers for
orientation only. Its literature searches were targeted, not an exhaustive
priority audit. For this merge the bibliographic data of these references
were checked; their content was not rechecked, and no proof here rests on
Fuchs--Heintze or Derksen.

\paragraph{Repository comparison at the pin.}
Source~18 read this report's \texttt{article.tex} at \texttt{bcac55a} in the
line ranges listed in its scope file, including the opening statements, the
mixed noncofinal theorem and its limitation, the certificate and
formalization notes and the questions, comparing against
\texttt{hol:thm:mixed}, \texttt{hol:q:mixedunit},
\texttt{hol:q:severaldilations} and \texttt{hol:q:nonlinear}; it also read the
repository overview and selected openings of other reports. Some responses
were truncated, and no absence claim rests on their unreturned parts. There
was no line-by-line audit of the repository, the companion archives, the Lean
modules or the history. Its statements about this report were accurate at
the pin, as recorded in Section~\ref{hol:ud:sub:conventions}.

\paragraph{Priority.}
Pure differential rigidity, the cyclic pure-dilation criterion, mixed
rigidity at a nonzero noncofinal valuation, the partial theta series, the
escape mechanism and the stronger nonlinear statements without an order unit
are this report's earlier material and are not claimed by source~18. Its
proposed contributions are the periodic-affine valuation theorem for Hahn
exponential polynomials with arbitrary well-ordered supports
(Theorems~\ref{hol:ud:thm:unitvalues} and~\ref{hol:ud:thm:affine}), the
bivariate unit-orbit estimate and the answer to
Question~\ref{hol:q:mixedunit}, with uniform degree bounds and polynomial
forcing (Theorem~\ref{hol:ud:thm:rigidity},
Corollary~\ref{hol:ud:cor:unitrigidity}), and the relative-scale dichotomy
for a finite set of multipliers with its two-multiplier witness
(Theorem~\ref{hol:ud:thm:dichotomy}). Historical novelty of the valuation
theorem is not certified, and no rank-one special case is claimed to be new;
the answered question is a question of this report, not a named historical
conjecture.

\emph{Batch~33 (merge).} Source~20 found the support envelope, the unit and
periodic-affine valuation theorems, the operator reduction and the mixed
unit-valuation corollary independently, on the same day, against the text
before this part was merged. They are printed once, here, and its narrower
criteria are Corollaries~\ref{hol:rc:cor:convex}
and~\ref{hol:rc:cor:finitegroup} (Section~\ref{hol:rc:sub:mixed}); no priority
between the two sources is asserted.

\subsection{What the unit-dilation part does not claim}
\label{hol:ud:sub:nonclaims}

Every limitation stated by source~18, in its manuscript, README, scope file or
verification record, is kept.
\begin{enumerate}[label=\textup{(U\arabic*)}]
\item No independent referee report and no Lean formalization accompany
source~18; no Lean build was run, no declaration was added, and the
repository's axiom audit does not cover this part. A PDF build and a finite
test log are not refereeing.
\item The Skolem--Mahler--Lech theorem is imported, not proved anew; no
effective bound on exceptional zeros is given.
\item Historical novelty of the valuation theorem and of its full scope is not
certified; no rank-one special case is claimed to be historically new, and
independence from all earlier rank-one results is not claimed. The answered
question is a question of this report at the pin \texttt{bcac55a}, not a named
historical conjecture.
\item No nonlinear rigidity is claimed: source~18 solves neither the
higher-order nonlinear order-unit question nor its own nonlinear question
(Remark~\ref{hol:ud:rem:nonlinearq} records that the latter is answered
negatively in the current text by source~13's series).
\item No unconditional torsion-block theorem is claimed: the hypothesis that
no quotient of actual multipliers is a root of unity is essential
(Section~\ref{hol:ud:sub:examples}), and a finite block criterion is open
(Question~\ref{hol:ud:q:torsionblocks}).
\item Theorem~\ref{hol:ud:thm:dichotomy} concerns the existence of \emph{some}
equation supported in a prescribed set of multipliers; it does not decide
prescribed operators or simultaneous systems
(Remark~\ref{hol:ud:rem:quantifiers}).
\item The proofs are existence proofs: no algorithm finds the first
nonidentity coefficient, the eventual winners or the last exceptional zero,
no effective zero set is computed, and no decision procedure for arbitrary
surreal inputs follows.
\item The exterior regions and degree bounds are uniform for a fixed operator
but are not claimed to be exact domains, and they depend on the exact
operator data (Example~\ref{hol:ud:ex:degree}).
\item Entireness is relative to one set-sized Hahn field: nothing is entire on
the proper class $\No[i]$, and entireness can be lost on enlarging the group.
$D_z$ kills every coefficient and is not the Berarducci--Mantova derivation;
no statement concerns that derivation or power series entire on all of
$\No[i]$.
\item Omnific integers enter only as permissible coefficients: no
factorization, integral-closure, automorphism or definability result is
claimed.
\item Linear independence of formal functions is not algebraic independence
of values at a point, and does not exclude polynomial relations.
\item Positive characteristic is excluded, with counterexamples to both the
valuation theorem and differential rigidity.
\item The 16{,}579 finite assertions check finite identities and examples.
They do not prove well ordering of arbitrary supports, strong summability at
all field elements, the Skolem--Mahler--Lech theorem, finite exceptional sets,
the nonexistence of an annihilator, or novelty, and their number is not a
measure of proof coverage.
\item The repository comparison was partial and targeted: selected line
ranges of this report at the pin, the overview and selected openings of other
reports; no audit of all reports, archives, Lean modules or history, and no
absence claim from truncated responses. The literature search was targeted.
\item \emph{[merge]} Remarks~\ref{hol:ud:rem:contains}
and~\ref{hol:ud:rem:nonlinearq}, the closed exterior region in
Theorem~\ref{hol:ud:thm:rigidity}, and the reading of
Corollary~\ref{hol:ud:cor:cyclic} beside Theorem~\ref{hol:main:q} are the
merge's; no historical priority is claimed for them.
\end{enumerate}

%======================================================================
% Newton-rigidity part: source 19. Labels use hol:nr:.
%======================================================================
\section{Order three through cubic jet degree: first polars on the twisted
cubic}\label{hol:nr:sec:setting}

This section and the next contain the material of source~19, delivered after
the order-three-threshold part had been merged. Its first theorem, second-order
rigidity for every value group and every differential degree, is
Theorem~\ref{hol:ot:main:second}\,(a), which source~19 found independently of
source~14; it is printed once, in Section~\ref{hol:ot:sec:setting}. The new
material starts one order higher. At order three the endpoint--gap
obstruction of Theorem~\ref{hol:ot:thm:gaps} constrains an exterior initial
polynomial at both ends; source~19 shows that when the exterior form has a
nonzero \emph{first polar} along the twisted cubic of monomial Euler jets, the
two constraints push the next exponent to opposite sides of the endpoint and
cannot both hold far out, and that every nonzero cubic form vanishing on that
curve has a nonzero first polar. Hence no nonpolynomial strongly entire series
satisfies an equation of order three and total jet degree at most three
(Theorem~\ref{hol:nr:main:third}). Source~19 adds external constants and
Gaussian-omnific coefficients, algebraic two-state systems, exact
classifications of two boundary equations beyond the first-polar test, and a
sparse strongly entire series whose initial polynomials satisfy both boundary
equations at every scale although the series satisfies neither.

\subsection{Source, conventions and the merge's choices}
\label{hol:nr:sub:conventions}

\begin{convention}[Setting of the Newton-rigidity part]
\label{hol:nr:conv:field}
In Sections~\ref{hol:nr:sec:setting}--\ref{hol:nr:sec:boundary}, $k$, $\Gamma$,
$\KK$, $v$, $\res$, $\EK$, $D_z$, $\vartheta=zD_z$, $\vartheta_X$, the Gauss data
$\gamma_\delta$, $\phi_\delta$, $\operatorname{Act}_\delta$,
$\underline N_\delta$, $N_\delta$, breaks, the exterior form $\mathfrak h_P$ of
\eqref{hol:ot:eq:exterior}, its dehomogenization $\mathfrak q$ and the
integer $i_0$ are as in Convention~\ref{hol:ot:conv:field} and
Section~\ref{hol:ot:sub:breaks}. For a homogeneous
$\mathfrak h\in k[Y_0,\ldots,Y_r]$ write
$\mathbf m_r(T)=(1,T,\ldots,T^r)$ for the \emph{monomial jet curve} (a conic
for $r=2$, the twisted cubic for $r=3$) and
$\chi_{\mathfrak h}(T)=\mathfrak h(\mathbf m_r(T))$; by
Remark~\ref{hol:ot:rem:corner}, $\chi_{\mathfrak h_P}=\chi_P$. A constant
extension is a field extension $\mathbb L/\KK$, with $D_z$ extended to
$\mathbb L((z))$ by killing $\mathbb L$; entireness always refers to $\KK$.
\end{convention}

For $k=\C$ one has $\KK=\K$ and $\EK=\E$; the generality in $k$ is recorded for
this part only.

\paragraph{The source.}
Source~19 is \emph{Low-Order Nonlinear Rigidity at Surreal Scales: Newton
transitions, Euler-jet geometry, and omnific coefficient descent} (24 pages,
US letter, 23 September 2026), delivered as the archive
\texttt{surreal\_newton\_rigidity} (batch~31, manuscript~08). It consulted the
repository at
\begin{center}\small\texttt{3d40856953f4bb0f5d069e45a3b4bab6eda33040},\end{center}
through a connector: the root README (partially), \texttt{docs/README.md}, this
report's directory and README, the complete files
\path{10-nonlinear-rigidity-PROOF_AUDIT.md} and
\path{11-coarsening-differential-rigidity-SOURCES_AND_SCOPE.md}, and the first
hundred lines of \texttt{article.tex}; the formalization ledger returned no
content, and the repository was not built. At that pin this report was the
five-source text (sources~08--12): neither source~13 nor source~14 had been
merged. Source~19's description was accurate there: first-order rigidity for
every value group, all-order rigidity without an order unit, and the
nonlinear question, Question~\ref{hol:q:nonlinear} (formerly~15.1), open with
an order unit. It is stale now in three places. Its principal claim, the
unrestricted second-order case, is Theorem~\ref{hol:ot:main:second}\,(a),
merged from source~14 before source~19 arrived; the two were written on the
same day, each against a text of this report that did not contain the other,
and neither cites the other. Its
abstract's ``the unrestricted all-order, order-unit case is not resolved'' and
its Question~12.1 (``is every strongly entire differentially algebraic series a
polynomial?'' with an order unit) are answered negatively by
Theorem~\ref{hol:td:main:boundary}: $\mathcal T_p$ is a nonpolynomial strongly
entire series with an equation of order three. Its statement that its own
second-order theorem is the principal proposed contribution is therefore
recorded as an independent re-derivation.

\paragraph{Renamed symbols.}
The following symbols of source~19 would collide with symbols of the earlier
parts. They are renamed here, in the new material only.
{\small
\begin{longtable}{@{}>{\raggedright\arraybackslash}p{0.27\textwidth}
>{\raggedright\arraybackslash}p{0.2\textwidth}
>{\raggedright\arraybackslash}p{0.45\textwidth}@{}}
\toprule
Source~19 & Here & Reason\\
\midrule
\endhead
$\partial=d/dz$; $\mathsf E=z\,d/dz$, $\mathsf E_X$ & $D_z$; $\vartheta$,
$\vartheta_X$ & The notation of Conventions~\ref{hol:nl:conv:field}
and~\ref{hol:td:conv:field}; $E_i$ are Euler-jet variables. (Source~19's
$\K=k((t^\Gamma))$ and $\mathcal E_\Gamma(k)$ are $\KK$ and $\EK$.)\\
$\alpha_n$, $I_\delta$, $F_\delta$ & $v(a_n)$, $\operatorname{Act}_\delta$,
$t^{-\gamma_\delta}f(t^{-\delta}X)$ & $F$ is \eqref{hol:eq:explicitF}; $I_P$
is the full corner polynomial.\\
Newton transition & break & As in Section~\ref{hol:ot:sub:conventions}.\\
$D,W,\mu$; residue corner $H$ & $d$, $w$, $\beta$; $\mathfrak h_P$ &
Lemma~\ref{hol:ot:lem:exterior}; $\mu$ is an order unit, $H\in\Gamma$.\\
$c_2(T)$, $c_3(T)$, $c_r(T)$ & $\mathbf m_2$, $\mathbf m_3$, $\mathbf m_r$ &
$\mathfrak c$ is a theta constant, $c_\nu$ corner coefficients.\\
$S=Y_0Y_2-Y_1^2$; $H=S^eQ$ & $\mathfrak s$; $\mathfrak h=\mathfrak s^e
\mathfrak h_1$ & $S$ is a support alphabet; $Q$ a cofactor.\\
$\chi_H$; first polar $A_H$, $a$, $b$ & $\chi_{\mathfrak h}$;
$\operatorname{Pol}_{\mathfrak h}$, $a_{\mathfrak h}$, $b_{\mathfrak h}$ &
$\chi_P$ is the residue corner polynomial, $\mathcal A$ a theta expression,
$A(U)$ a transverse part.\\
$H_3$, $Q_4$, $B(m,n)$ & $\mathsf H_3$, $\mathsf Q_4$, $\mathsf B(m,n)$ & $H$,
$Q_s$, $B$ and $\mathcal B$ have other meanings.\\
$A(N),B(N)$, $N$; $U=\mathsf Ef/f$ & $\mathsf A_1(\mathsf n)$,
$\mathsf A_2(\mathsf n)$, $\mathsf n$; $\mathsf u_f$ & $N$ is an index;
$\mathsf u_\phi$ as in Section~\ref{hol:ot:sub:endpoint}.\\
$L/\K$ (constants); $F=L(z,f,g)$, $F_1$; $q(T)$ & $\mathbb L/\KK$;
$\mathfrak M$, $\mathfrak M_1$; $\mathsf m(T)$ & $L$ is a coefficient
subfield; $q$ a dilation.\\
$u$ (order unit); $c_n$; $\Theta(z)=1+t^uz\Theta(t^{2u}z)$ & $\mu$; $b_n$;
$G_\mu(z)=\sum t^{n^2\mu}z^n$ & $u$ is the Laurent variable; source~19's
$\Theta$ is $\Theta_{p^2}(pz)$ with $p=t^\mu$, not $\Theta_p$ of
\eqref{hol:eq:theta}.\\
$F=\sum t^{e_n}z^n$ & $F$ & The same series \eqref{hol:eq:explicitF}.\\
field $F$ of Section~10; $p=a_mz^m+a_{2m}z^{2m}$ & $K_0$; $\mathsf d_0$ & $F$
is \eqref{hol:eq:explicitF}, as in the row above; $K_0$ as in
Theorem~\ref{hol:cf:thm:coefficients}; $p=t^\mu$.\\
characteristic $p$ & $\mathfrak p$ & $p=t^\mu$.\\
Theorems~A, B, C & Theorem~\ref{hol:ot:main:second}\,(a),
Theorem~\ref{hol:nr:main:third}, Corollary~\ref{hol:nr:cor:external}
and Theorem~\ref{hol:nr:thm:systems} & Theorems~A--N of this report are different
theorems.\\
\bottomrule
\end{longtable}}
The shipped audit files \path{19-newton-rigidity-PROOF_AUDIT.md} and
\path{19-newton-rigidity-SOURCES_AND_SCOPE.md} keep the source's notation.

\paragraph{Words with two meanings.}
Source~19's \emph{residue corner} is the exterior form $\mathfrak h_P$, not
the residue corner polynomial $\chi_P$; its \emph{Newton transition} is a
break. The \emph{first polar} $\operatorname{Pol}_{\mathfrak h}$ is a
polynomial in two variables attached to a homogeneous form, not the full corner
polynomial $I_P$. \emph{Degree} of an equation means its total degree in the
jet variables $Y_0,\ldots,Y_r$ (source~19: differential degree); the degree
in $z$ is unrestricted. \emph{Order} is the highest derivative, as in
Section~\ref{hol:td:sub:conventions}.

\paragraph{Results printed once.}
Source~19's Theorem~A, its principal claim, is
Theorem~\ref{hol:ot:main:second}\,(a), with identical hypotheses (any $k$ of
characteristic zero, any $\Gamma$, any degree, any residue corner), and its
proof is in substance source~14's: Theorem~3.2 (arbitrarily exterior
nonmonomial initial polynomials) is Lemma~\ref{hol:ot:lem:breaks}, with the
break found by a crossing slope in a finite window below one prescribed scale
rather than on an interval; Proposition~4.1 (exterior corner reduction) is
Lemma~\ref{hol:ot:lem:exterior}, with the same sign-sensitive inequality
\eqref{hol:ot:eq:gapbound}; and its Lemma~5.1 and Proposition~5.3 are
Lemma~\ref{hol:ot:lem:endpoints} in homogeneous coordinates
(Remark~\ref{hol:nr:rem:conic}). Also printed once: its all-scale criterion
(Proposition~2.2), which is Proposition~\ref{hol:cf:prop:criterion}; the closure
of $\EK$ (Corollary~2.3), which is Proposition~\ref{hol:prop:closure} with
$\vartheta=zD_z$; the cofinal scalar extension (Lemma~2.4), which is
Lemma~\ref{hol:ot:lem:extension}; the escape of the active support
(Lemma~3.1), which is Lemma~\ref{hol:nl:lem:active}; its Remark~5.4, whose
equation $z^2ff''+zff'-z^2f'^2=0$ is \eqref{hol:nl:eq:degenerate}, with the
classification of Proposition~\ref{hol:nl:prop:degenerate}; the finite linear
descent (Lemma~7.1), which is Proposition~\ref{hol:cf:prop:scalar}, whose
basis proof also preserves the total degree; the full-class obstruction
(Proposition~7.6), which is Proposition~\ref{hol:ot:prop:fullclass}; the
embedding \eqref{hol:eq:embedding}; and its examples in characteristic
$\mathfrak p$ and with $\sum z^n/n!$ (Section~9.3), which are
Example~\ref{hol:nl:ex:charp}, the positive-characteristic paragraph of
Section~\ref{hol:ot:sub:surreal} and \eqref{hol:eq:monomialexp}. The quadratic
profile of its Example~3.4 (for $G_1$ at $\delta=2n+1$ the initial polynomial
is $X^n+X^{n+1}$) is the Newton profile recorded after
Theorem~\ref{hol:ot:thm:convexity} for $\mathcal T_p$.

\paragraph{Where each numbered result went.}
Theorem~A is Theorem~\ref{hol:ot:main:second}\,(a), and Theorem~B is
Theorem~\ref{hol:nr:main:third}; Definition~6.1 and Lemmas~6.2, 6.3 are
Definition~\ref{hol:nr:def:polar} and Lemmas~\ref{hol:nr:lem:polar},
\ref{hol:nr:lem:side}; Proposition~6.4, Lemma~6.5, Remark~6.6,
Corollary~6.7 and Example~6.8 are Proposition~\ref{hol:nr:prop:firstpolar},
Lemma~\ref{hol:nr:lem:secant}, Remark~\ref{hol:nr:rem:exactuse},
Corollary~\ref{hol:nr:cor:factor} and Example~\ref{hol:nr:ex:cubic};
Corollaries~7.2, 7.3, Lemma~7.4 and Corollary~7.5 are
Corollaries~\ref{hol:nr:cor:external}, \ref{hol:nr:cor:surcomplex},
Lemma~\ref{hol:nr:lem:omnificclear} and Corollary~\ref{hol:nr:cor:omnific};
Theorem~8.1 and Corollary~8.2 are Theorem~\ref{hol:nr:thm:systems} and
Corollary~\ref{hol:nr:cor:planar}; Proposition~9.1 is
Proposition~\ref{hol:nr:prop:family}; Theorems~10.1 and~10.2 are
Theorems~\ref{hol:nr:thm:H3} and~\ref{hol:nr:thm:Q4}, and Section~10.3 is
Example~\ref{hol:nr:ex:sparse}; Question~12.1 is answered
(Section~\ref{hol:nr:sub:conventions}); Questions~12.4, 12.6 and~12.8 are
Questions~\ref{hol:ot:q:meromorphic}, \ref{hol:ot:q:pde}
and~\ref{hol:ot:q:charp}, asked by source~14 in the same form, with status
notes recording source~19's versions; Questions~12.2, 12.3, 12.5, 12.7, 12.9
and~12.10 are Questions~\ref{hol:nr:q:layers}--\ref{hol:nr:q:formal}.

\paragraph{The merge's choices.}
The review of source~19 for this merge found its proofs correct; the
identities it states were rechecked with an independent SymPy script, which is
not shipped. The merge adds four points, each marked. First,
Remark~\ref{hol:nr:rem:gaps} places the first-polar criterion inside the
endpoint--gap obstruction of Theorem~\ref{hol:ot:thm:gaps}: for $\chi=0$ the
first polar is nonzero exactly when $i_0=1$, the two linear conditions of
source~19 are the two endpoint conditions of Theorem~\ref{hol:ot:thm:gaps},
and an explicit cubic shows that Corollary~\ref{hol:ot:cor:gapcriterion},
which uses the upper endpoint only, does not contain
Theorem~\ref{hol:nr:main:third}. Second, Corollary~\ref{hol:nr:cor:theta}
applies Theorem~\ref{hol:nr:main:third} to $\mathcal T_p$ and to the theta
exterior form. Third, Example~\ref{hol:nr:ex:sparse} is read as a negative
instance for Question~\ref{hol:ot:q:lift}. Fourth, the difference between
Theorem~\ref{hol:nr:thm:systems} and Corollary~\ref{hol:ot:cor:systems} is
stated. Source~19's theorems, proofs, examples and non-claims are otherwise
printed as delivered, in the renamed notation.

\begin{remark}[Source~19's route to Theorem~\ref{hol:ot:main:second}\,(a);
source~19]\label{hol:nr:rem:conic}
Source~19 proves the endpoint lemma in homogeneous coordinates. Modulo
$\mathfrak s=Y_0Y_2-Y_1^2$ the monomials $Y_0^{d-b}Y_1^b$ ($0\le b\le d$) and
$Y_1^{d-c}Y_2^c$ ($1\le c\le d$) span degree $d$, and their images
$1,T,\ldots,T^{2d}$ under $Y\mapsto\mathbf m_2(T)$ are independent; so a
homogeneous $\mathfrak h$ with $\chi_{\mathfrak h}=0$ is divisible by
$\mathfrak s$, and every nonzero homogeneous $\mathfrak h$ is uniquely
$\mathfrak s^e\mathfrak h_1$ with $\chi_{\mathfrak h_1}\ne0$. Since
$\mathfrak s(1,U,U^2+V)=V$, this $e$ is the $i_0$ of
\eqref{hol:ot:eq:factor} and $\mathfrak h_1(1,U,U^2+V)=\mathfrak r$, so the
endpoint spectrum $\operatorname{Sp}(\mathfrak h)$ is the set of nonnegative
integral roots of $\chi_{\mathfrak h_1}=\mathfrak r(T,0)$. Source~19 then
argues at the top degree only: if $\deg\phi$ exceeds these roots, the
coefficient of degree $(\deg\mathfrak h_1)\deg\phi$ in
$\mathfrak h_1(\phi,\vartheta_X\phi,\vartheta_X^2\phi)$ is
$\lc(\phi)^{\deg\mathfrak h_1}\chi_{\mathfrak h_1}(\deg\phi)\ne0$, so there
is no such $\phi$ if $e=0$, and
$\mathfrak s(\phi,\vartheta_X\phi,\vartheta_X^2\phi)=0$ if $e>0$; and a nonzero
polynomial with $\mathfrak s(\phi,\ldots)=0$ is a monomial, because for its
first two exponents $m<n$ the coefficient of $X^{m+n}$ in
$\phi\vartheta_X^2\phi-(\vartheta_X\phi)^2$ is $a_ma_n(n-m)^2\ne0$. Source~14
uses logarithmic Euler derivatives at both endpoints instead. The two routes
prove the same lemma; the variant is recorded, not printed as a separate
theorem.
\end{remark}

\subsection{The first polar and third-order rigidity}
\label{hol:nr:sub:polar}

\begin{definition}[First polar along the twisted cubic; source~19]
\label{hol:nr:def:polar}
For a homogeneous $\mathfrak h\in k[Y_0,Y_1,Y_2,Y_3]$ of degree $d\ge1$ put
\begin{equation}\label{hol:nr:eq:polar}
 \operatorname{Pol}_{\mathfrak h}(T,U)=\sum_{j=0}^3
 \frac{\partial\mathfrak h}{\partial Y_j}\bigl(\mathbf m_3(T)\bigr)\,U^j .
\end{equation}
\end{definition}

\begin{lemma}[The double factor; source~19]\label{hol:nr:lem:polar}
If $\chi_{\mathfrak h}=0$, there are $a_{\mathfrak h},b_{\mathfrak h}\in k[T]$
with
\begin{equation}\label{hol:nr:eq:polarfactor}
 \operatorname{Pol}_{\mathfrak h}(T,U)=(U-T)^2\bigl(a_{\mathfrak h}(T)U
 +b_{\mathfrak h}(T)\bigr),
\end{equation}
not both zero if $\operatorname{Pol}_{\mathfrak h}\ne0$.
\end{lemma}

\begin{proof}
Euler's identity gives $\operatorname{Pol}_{\mathfrak h}(T,T)=
d\,\chi_{\mathfrak h}(T)=0$, and differentiating $\chi_{\mathfrak h}(T)=0$ gives
$\partial_U\operatorname{Pol}_{\mathfrak h}(T,U)|_{U=T}=0$. Division in
$k[T][U]$ gives the factor $(U-T)^2$, and the $U$-degree is at most three.
\end{proof}

\begin{lemma}[An eventual side for integral roots; source~19]
\label{hol:nr:lem:side}
Let $a,b\in k[T]$ be not both zero. There is $B\in\N$ such that for integers
$m>B$ every integral solution $n\ne m$ of $a(m)n+b(m)=0$ lies on one fixed side
of $m$: always $n>m$, or always $n<m$, the side not depending on $m$. There
may be no solutions.
\end{lemma}

\begin{proof}
If $a=0$, then $b(m)\ne0$ for all but finitely many $m$, and eventually there
are no solutions. Let $a\ne0$ and $C(T)=Ta(T)+b(T)$. If $C=0$, then away from
the roots of $a$ the only root is $n=m$, which is excluded. Let $a,C\ne0$ and
let $V\subseteq k$ be the finite-dimensional $\Q$-span of their coefficients.
There is a $\Q$-linear $\ell:V\to\Q$ for which neither coefficientwise image
$a_*,C_*\in\Q[T]$ is zero: the functionals killing all coefficients of $a$,
or all of $C$, form two proper subspaces of $V^*$, and a vector space over the
infinite field $\Q$ is not the union of two proper subspaces. Applying $\ell$
to $a(m)(n-m)+C(m)=0$ gives $a_*(m)(n-m)+C_*(m)=0$. For large $m$ both
$a_*(m)$ and $C_*(m)$ have fixed nonzero signs, so $n-m=-C_*(m)/a_*(m)$ has a
fixed strict sign. This is a statement about integers, not an ordering of $k$.
\end{proof}

\begin{proposition}[The first-polar criterion; source~19]
\label{hol:nr:prop:firstpolar}
Let $0\ne\mathfrak h\in k[Y_0,Y_1,Y_2,Y_3]$ be homogeneous, and suppose
$\chi_{\mathfrak h}\ne0$, or $\chi_{\mathfrak h}=0$ and
$\operatorname{Pol}_{\mathfrak h}\ne0$. There is $B_{\mathfrak h}\in\N$ such
that every nonzero $\phi\in k[X]$ with
\[
 \mathfrak h(\phi,\vartheta_X\phi,\vartheta_X^2\phi,\vartheta_X^3\phi)=0,
 \qquad \operatorname{ord}_X\phi>B_{\mathfrak h},
\]
is a monomial.
\end{proposition}

\begin{proof}
If $\chi_{\mathfrak h}\ne0$, the top coefficient of the expression is
$\lc(\phi)^d\chi_{\mathfrak h}(\deg\phi)\ne0$ once $\deg\phi$ exceeds the
integral roots of $\chi_{\mathfrak h}$, and there is no such solution. Let
$\chi_{\mathfrak h}=0\ne\operatorname{Pol}_{\mathfrak h}$ and let $\phi$ be
nonmonomial, with first two exponents $m<n$, last two $\ell<M$, and
coefficients $a_j$. Start with all $d$ factors at degree $m$; one replacement by
$n$ gives exactly the first polar, and two replacements, or one by a larger
exponent, raise the degree further. So the coefficient of $X^{(d-1)m+n}$ is
\begin{equation}\label{hol:nr:eq:bottom}
 a_m^{d-1}a_n\operatorname{Pol}_{\mathfrak h}(m,n),
\end{equation}
and symmetrically the coefficient of $X^{(d-1)M+\ell}$ is
$a_M^{d-1}a_\ell\operatorname{Pol}_{\mathfrak h}(M,\ell)$. Both vanish; by
\eqref{hol:nr:eq:polarfactor} and characteristic zero,
$a_{\mathfrak h}(m)n+b_{\mathfrak h}(m)=0$ and
$a_{\mathfrak h}(M)\ell+b_{\mathfrak h}(M)=0$. Once $m$, and so $M\ge m$,
exceeds the bound of Lemma~\ref{hol:nr:lem:side}, these are incompatible:
$n>m$ but $\ell<M$, on opposite sides.
\end{proof}

\begin{lemma}[No cubic is singular along the twisted cubic; source~19]
\label{hol:nr:lem:secant}
If a homogeneous $\mathfrak h\in k[Y_0,Y_1,Y_2,Y_3]$ of positive degree at
most three has all first partial derivatives zero on $\mathbf m_3(T)$, then
$\mathfrak h=0$.
\end{lemma}

\begin{proof}
Extend $k$ to an algebraic closure. By Euler's identity also
$\chi_{\mathfrak h}=0$, since the degree is nonzero in $k$. Restrict
$\mathfrak h$ to secants: $G(A,B)=\mathfrak h(A\mathbf m_3(s)+B\mathbf m_3(t))$
is homogeneous of degree at most three in $(A,B)$, and its value and first
transverse derivative vanish at $B=0$, so $B^2\mid G$; likewise $A^2\mid G$;
hence $G=0$. The secant map
$(A,B,s,t)\mapsto(A+B,As+Bt,As^2+Bt^2,As^3+Bt^3)$ has Jacobian determinant
\begin{equation}\label{hol:nr:eq:jacobian}
 -AB(s-t)^4\ne0,
\end{equation}
so in characteristic zero its coordinates are algebraically independent (a
minimal-degree relation, differentiated, would have all partial derivatives
zero by the invertible Jacobian), and no nonzero polynomial vanishes on its
image. So $\mathfrak h=0$.
\end{proof}

\begin{remark}[The exact use of the lemma; source~19]
\label{hol:nr:rem:exactuse}
$\operatorname{Pol}_{\mathfrak h}=0$ says that every coefficient in $U$ of
\eqref{hol:nr:eq:polar} vanishes, that is, every first partial vanishes on the
curve. So a nonzero homogeneous $\mathfrak h$ of degree at most three with
$\chi_{\mathfrak h}=0$ has $\operatorname{Pol}_{\mathfrak h}\ne0$.
\end{remark}

\begin{proof}[Proof of Theorem~\ref{hol:nr:main:third}]
Suppose a nonpolynomial $f\in\EK$ satisfies a nonzero equation of order at
most three and total jet degree at most three. Pass to $\KK_\Q$
(Lemma~\ref{hol:ot:lem:extension}) and form the exterior form $\mathfrak h_P$,
nonzero homogeneous of degree between one and three, since the Euler
substitution \eqref{hol:ot:eq:stirling} preserves total jet degree
(Lemma~\ref{hol:ot:lem:exterior}). If $\chi_{\mathfrak h_P}=0$,
Lemma~\ref{hol:nr:lem:secant} gives $\operatorname{Pol}_{\mathfrak h_P}\ne0$.
Either way Proposition~\ref{hol:nr:prop:firstpolar} applies: beyond the
$\delta_0$ of Lemma~\ref{hol:ot:lem:exterior} and the scale of
Lemma~\ref{hol:nl:lem:active} beyond which every active index exceeds
$B_{\mathfrak h_P}$, every initial polynomial is a monomial, while
Lemma~\ref{hol:ot:lem:breaks} gives a break there. For the general assertion,
apply Proposition~\ref{hol:nr:prop:firstpolar} to $\mathfrak h_P$, or
Corollary~\ref{hol:nr:cor:factor} to its irreducible factors, in the same
argument.
\end{proof}

\begin{corollary}[A factor criterion beyond cubic degree; source~19]
\label{hol:nr:cor:factor}
For an equation of order at most three, factor $\mathfrak h_P$ into
homogeneous irreducible factors over $k$. If each factor $\mathfrak g$ has
$\chi_{\mathfrak g}\ne0$ or $\operatorname{Pol}_{\mathfrak g}\ne0$, every
strongly entire solution is a polynomial.
\end{corollary}

\begin{proof}
Take the largest of the bounds of Proposition~\ref{hol:nr:prop:firstpolar} for
the finitely many factors. Since $k[X]$ is a domain, an initial polynomial
solving the product equation solves one factor equation, and the argument
above applies. Factoring first removes repeated factors, whose unfactored first
polar can vanish identically.
\end{proof}

\begin{example}[Two ends forcing incompatible next exponents; source~19]
\label{hol:nr:ex:cubic}
For $\mathfrak h=Y_0^2Y_3-4Y_0Y_1Y_2+3Y_1^3$ one has $\chi_{\mathfrak h}=0$ and
$\operatorname{Pol}_{\mathfrak h}(T,U)=(U-T)^2(U-2T)$. A nonmonomial
polynomial solution with lowest exponent $m$ would need its next exponent to
be $2m$, and at its largest exponent $M$ the next smaller one would also have
to be $2M$, which is impossible. This case shows the two-end argument without
the rational functional of Lemma~\ref{hol:nr:lem:side}.
\end{example}

\begin{remark}[Relation to the endpoint--gap obstruction; merge]
\label{hol:nr:rem:gaps}
Let $\mathfrak h$ be homogeneous in $Y_0,\ldots,Y_3$ with $\chi_{\mathfrak h}=0$,
and write $\mathfrak h_j$ for $\partial\mathfrak h/\partial Y_j$ evaluated on
$\mathbf m_3(T)$. In \eqref{hol:ot:eq:dehom3} the linear part of $\mathfrak q$
in $(V,W)$ is $V(\mathfrak h_2+3T\mathfrak h_3)+W\mathfrak h_3$ at $U=T$, and
the relations $\mathfrak h_0+T\mathfrak h_1+T^2\mathfrak h_2+T^3\mathfrak h_3=0$
(Euler) and $\mathfrak h_1+2T\mathfrak h_2+3T^2\mathfrak h_3=0$ (derivative of
$\chi_{\mathfrak h}$) show that it vanishes exactly when all $\mathfrak h_j$
vanish. So, when $\chi_{\mathfrak h}=0$, $\operatorname{Pol}_{\mathfrak h}\ne0$
exactly when $i_0=1$, and then, comparing coefficients in
\eqref{hol:nr:eq:polarfactor}, $a_{\mathfrak h}=\mathfrak h_3$,
$b_{\mathfrak h}=\mathfrak h_2+2T\mathfrak h_3$ and
\[
 \Upsilon_{\mathfrak h}(T,G)=a_{\mathfrak h}(T)(T-G)+b_{\mathfrak h}(T).
\]
The conditions $\Upsilon_{\mathfrak h}(N,\mathsf g^+)=0$ and
$\Upsilon_{\mathfrak h}(m,-\mathsf g^-)=0$ of Theorem~\ref{hol:ot:thm:gaps} are
then exactly source~19's $a_{\mathfrak h}(N)\ell+b_{\mathfrak h}(N)=0$ with
$\ell=N-\mathsf g^+$, and $a_{\mathfrak h}(m)n+b_{\mathfrak h}(m)=0$ with
$n=m+\mathsf g^-$. The new ingredients of Theorem~\ref{hol:nr:main:third} are
therefore Lemma~\ref{hol:nr:lem:side}, which makes one of the two endpoint
conditions eventually unsatisfiable, and Lemma~\ref{hol:nr:lem:secant}, which
forces $i_0\le1$ in degree at most three. Corollary~\ref{hol:ot:cor:gapcriterion}
tests the upper endpoint only and does not contain the theorem. For the cubic
\[
 \mathfrak h=2Y_0^2Y_3-5Y_0Y_1Y_2+3Y_1^3,\qquad
 \chi_{\mathfrak h}=0,\qquad
 \operatorname{Pol}_{\mathfrak h}=(U-T)^2(2U-T),\qquad
 \Upsilon_{\mathfrak h}=T-2G,
\]
the upper condition allows the gap $\mathsf g^+=N/2$ for every even $N$, so the
set in Corollary~\ref{hol:ot:cor:gapcriterion} is unbounded, while the lower
condition asks for the next exponent $m/2<m$ and has no solution; the equation
$2f^2\vartheta^3f-5f\vartheta f\,\vartheta^2f+3(\vartheta f)^3=0$ is covered by
Theorem~\ref{hol:nr:main:third} and not by that corollary. For
Example~\ref{hol:nr:ex:cubic}, $\Upsilon_{\mathfrak h}=-(G+T)$ and the upper
condition alone suffices. These identities were checked with SymPy.
\end{remark}

\subsection{External constants and planar systems}
\label{hol:nr:sub:constants}

\begin{corollary}[External constants; source~19]\label{hol:nr:cor:external}
For nonpolynomial $f\in\EK$ and every field extension $\mathbb L/\KK$, the
series $f,D_zf,D_z^2f$ are algebraically independent over $\mathbb L(z)$, and
$f,D_zf,D_z^2f,D_z^3f$ satisfy no nonzero polynomial relation of total degree
at most three over $\mathbb L(z)$.
\end{corollary}

\begin{proof}
A relation over $\mathbb L(z)$ gives one over $\KK(z)$ of no larger total
degree: in the proof of Proposition~\ref{hol:cf:prop:scalar} each $Q_\nu$ has
total degree at most that of $Q$ (source~19 uses a $\KK$-linear functional on
the finite span of the coefficients instead of a basis). Apply
Theorems~\ref{hol:ot:main:second}\,(a) and~\ref{hol:nr:main:third}.
\end{proof}

This is stronger than reapplying the analytic theorems over a larger Hahn
field, which would need entireness there and can fail; the finite algebraic
descent needs no analytic hypothesis.

\begin{corollary}[External surcomplex coefficients; source~19]
\label{hol:nr:cor:surcomplex}
Let $\Gamma\subseteq\No$ be a set-sized subgroup and $f\in\C((t^\Gamma))[[z]]$
nonpolynomial and strongly entire on that workspace. No nonzero polynomial
relation with finitely many coefficients in $\No[i](z)$ holds among
$f,D_zf,D_z^2f$, and none of total degree at most three among
$f,\ldots,D_z^3f$.
\end{corollary}

\begin{proof}
The workspace and the finitely many coefficients of a relation lie in a
set-sized subfield of $\No[i]$; apply Corollary~\ref{hol:nr:cor:external}. No
proper-class polynomial ring or tensor product is used.
\end{proof}

\begin{lemma}[Finite monomial clearance; source~19]
\label{hol:nr:lem:omnificclear}
For finitely many $c_1,\ldots,c_s\in\No[i]$ there is a surreal $\rho>0$ with
$\omega^\rho c_j\in\Oz[i]$ for all $j$; the real and imaginary parts of these
products have constant term $0$.
\end{lemma}

\begin{proof}
The union of the normal-form supports of the real and imaginary parts is a
set; take $\rho>0$ larger than the negative of each of its exponents.
Translating by $\rho$ makes every exponent positive and preserves the order of
the support, which is the normal-form description of $\Oz$ with constant term
$0$, as in \texttt{odg:def:rings}~\cite{repoodg}. If the union is empty any
$\rho>0$ works.
\end{proof}

\begin{corollary}[Omnific differential-relation obstruction; source~19]
\label{hol:nr:cor:omnific}
For every nonpolynomial strongly entire $f$ in a surcomplex Hahn workspace,
$P(z,f,D_zf,D_z^2f)\ne0$ for every nonzero $P\in\Oz[i][z,Y_0,Y_1,Y_2]$, and the
same holds through $D_z^3f$ when the total jet degree is at most three.
Conversely every surcomplex polynomial differential relation with finitely
many coefficients becomes a Gaussian-omnific one after clearing
$z$-denominators and multiplying by one monomial $\omega^\rho$.
\end{corollary}

\begin{proof}
The obstruction is a case of Corollary~\ref{hol:nr:cor:surcomplex}; the
conversion is Lemma~\ref{hol:nr:lem:omnificclear}, and multiplication by a
nonzero scalar preserves order and degree.
\end{proof}

Monomial clearance is elementary normal-form algebra, not a new arithmetic
theorem. The consequence concerns finitely many coefficients of formal
relations: it says nothing about factorization of omnific integers and is not
algebraic independence of a value $f(x)$ from the field containing it.

\begin{theorem}[Algebraic two-state systems; source~19]
\label{hol:nr:thm:systems}
Let $\mathbb L/\KK$ be a field extension and $f,g\in\EK$. If $D_zf$ and $D_zg$
are algebraic over $\mathfrak M=\mathbb L(z,f,g)\subseteq\mathbb L((z))$, then
$f,g\in\KK[z]$.
\end{theorem}

\begin{proof}
$\mathfrak M_1=\mathfrak M(D_zf,D_zg)$ is finite algebraic over $\mathfrak M$,
so $\trdeg_{\mathbb L(z)}\mathfrak M_1\le2$, and derivatives of elements of
$\mathfrak M$ lie in $\mathfrak M_1$. Let $\mathsf m(T)\in\mathfrak M[T]$ be the
minimal polynomial of $D_zf$; it is separable in characteristic zero, so
$\mathsf m'(D_zf)\ne0$. Differentiating $\mathsf m(D_zf)=0$ gives
$D_z^2f=-\mathsf m^{D}(D_zf)/\mathsf m'(D_zf)\in\mathfrak M_1$, with
$\mathsf m^D$ obtained by differentiating the coefficients. Likewise
$D_z^2g\in\mathfrak M_1$. So $f,D_zf,D_z^2f$ lie in a field of transcendence
degree at most two over $\mathbb L(z)$ and are dependent, and
Corollary~\ref{hol:nr:cor:external} makes $f$ a polynomial; likewise $g$.
\end{proof}

The separant is nonzero because it comes from a minimal polynomial, not from
an arbitrary singular presentation of an implicit equation, and
$\mathfrak M_1$ is in fact closed under $D_z$.

\begin{corollary}[Rational planar systems; source~19]
\label{hol:nr:cor:planar}
Every strongly entire pair $f,g$ with $D_zf=R_1(z,f,g)$ and $D_zg=R_2(z,f,g)$,
$R_1,R_2\in\mathbb L(z,Y_1,Y_2)$, denominators nonzero after substitution, is
polynomial. The coefficients may be external surcomplex constants.
\end{corollary}

These statements assert neither the existence of polynomial solutions, nor
finiteness of their number, nor a uniform bound on their degrees; the
strongly entire solution locus contains no genuinely transcendental pair, and
\eqref{hol:nl:eq:degenerate} has polynomial solutions of unbounded degree.

\paragraph{Comparison with Corollary~\ref{hol:ot:cor:systems} (merge).}
The two statements overlap without containing each other.
Corollary~\ref{hol:ot:cor:systems} (source~14) needs a \emph{rational} system
over $\KK(z)$ but lets the states $u_1,u_2$ be arbitrary elements of a
differential field, not necessarily entire, and concludes for every entire
element of $\KK(z,u_1,u_2)$. Theorem~\ref{hol:nr:thm:systems} (source~19) needs
the states $f,g$ themselves to be entire but allows \emph{algebraic}
dependence of $D_zf,D_zg$ over $\mathbb L(z,f,g)$ and external constants
$\mathbb L\supseteq\KK$. Both rest on three-jet independence
(Corollary~\ref{hol:ot:cor:threejets}), the second through
Proposition~\ref{hol:cf:prop:scalar}. With an order unit, $\mathcal T_p$ is an
observable of a four-state rational system (Theorem~\ref{hol:td:thm:field}), so
neither statement extends to four states; three states are open
(Question~\ref{hol:nr:q:systems}).

\section{Exact classifications beyond the first-polar test, and scope}
\label{hol:nr:sec:boundary}

The first-polar criterion is sufficient, not necessary. The following two
equations show its boundary without producing a counterexample to rigidity.
In Sections~\ref{hol:nr:sub:H3} and~\ref{hol:nr:sub:Q4}, $K_0$ is any field of
characteristic zero, $f\in K_0[[z]]$, and $\vartheta=z\,d/dz$; no Hahn or
entire hypothesis is used.

\subsection{A quartic equation of order three}\label{hol:nr:sub:H3}

Put
\begin{equation}\label{hol:nr:eq:H3}
 \mathsf H_3=2Y_0^2Y_3^2-13Y_0Y_1Y_2Y_3+9Y_0Y_2^3+9Y_1^3Y_3-7Y_1^2Y_2^2 .
\end{equation}
All its first partial derivatives vanish on $\mathbf m_3(T)$, so
$\operatorname{Pol}_{\mathsf H_3}=0$ and Proposition~\ref{hol:nr:prop:firstpolar}
does not apply; and $\mathsf H_3$ annihilates the Euler jets of $az^n+bz^{2n}$
for all $a,b,n$, which have arbitrarily large lowest degree.

\begin{theorem}[Exact formal solutions of the quartic; source~19]
\label{hol:nr:thm:H3}
The formal power series with $\mathsf H_3(f,\vartheta f,\vartheta^2f,
\vartheta^3f)=0$ are exactly $0$, the monomials $az^m$ ($m\in\N$) and the
binomials $az^n+bz^{2n}$ with $n\ge1$ and $ab\ne0$.
\end{theorem}

\begin{proof}
Write $Y_j=\vartheta^jf$ and, in an auxiliary variable $\mathsf n$,
\[
 \mathsf A_1(\mathsf n)=2Y_0\mathsf n^2-3Y_1\mathsf n+Y_2,\qquad
 \mathsf A_2(\mathsf n)=2Y_1\mathsf n^2-3Y_2\mathsf n+Y_3 .
\]
Expanding the Sylvester determinant gives
\begin{equation}\label{hol:nr:eq:resultant}
 \operatorname{Res}_{\mathsf n}(\mathsf A_1,\mathsf A_2)=2\mathsf H_3 .
\end{equation}
Let $f\ne0$ and put $\mathfrak s_f=f\vartheta^2f-(\vartheta f)^2$. If
$\mathfrak s_f=0$, the first-two-exponents argument of
Remark~\ref{hol:nr:rem:conic}, valid for formal series, or
Proposition~\ref{hol:nl:prop:degenerate}, makes $f$ a monomial. Let
$\mathfrak s_f\ne0$. By \eqref{hol:nr:eq:resultant}, $\mathsf A_1$ and
$\mathsf A_2$ have a common root in an algebraic closure of $K_0((z))$
($\mathsf A_1$ is genuinely quadratic since $f\ne0$; a degree drop in
$\mathsf A_2$ does not affect this). The combination
$Y_0\mathsf A_2-Y_1\mathsf A_1=-3\mathfrak s_f\mathsf n+(Y_0Y_3-Y_1Y_2)$ shows
that the root is
$\mathsf n=(f\vartheta^3f-\vartheta f\,\vartheta^2f)/(3\mathfrak s_f)\in
K_0((z))$. Applying $\vartheta$ to $\mathsf A_1(\mathsf n)=0$, with
$\vartheta Y_j=Y_{j+1}$,
\begin{equation}\label{hol:nr:eq:parameter}
 \mathsf A_2(\mathsf n)+(4f\mathsf n-3\vartheta f)\,\vartheta\mathsf n=0 ,
\end{equation}
so $\vartheta\mathsf n=0$ or $4f\mathsf n=3\vartheta f$. In the first case
$\mathsf n\in K_0$, the constants of $\vartheta$ on $K_0((z))$, and the
coefficients of $\mathsf A_1(\mathsf n)=0$ read $(j-\mathsf n)(j-2\mathsf n)
a_j=0$: at most two coefficients are nonzero, and if two, their degrees are
$\mathsf n$ and $2\mathsf n$ with $\mathsf n$ a positive integer. In the second
case substitution gives $8f\vartheta^2f-9(\vartheta f)^2=0$, that is
$8\vartheta\mathsf u_f-\mathsf u_f^2=0$ for $\mathsf u_f=\vartheta f/f$. With
$m=\operatorname{ord}_zf$, $\mathsf u_f=m+O(z)$ and the constant term forces
$m=0$; if $\mathsf u_f\ne0$, its first positive degree $r$ gives a nonzero term
of degree $r$ in $8\vartheta\mathsf u_f$ while $\mathsf u_f^2$ starts at degree
$2r$. So $\mathsf u_f=0$, $f$ is constant and $\mathfrak s_f=0$, a
contradiction. Substitution verifies the listed solutions.
\end{proof}

\subsection{A quadratic equation of order four}\label{hol:nr:sub:Q4}

Put $\mathsf Q_4=-2Y_0Y_4+9Y_1Y_3-7Y_2^2$. For distinct monomial degrees $m,n$
its cross coefficient is governed by
\begin{equation}\label{hol:nr:eq:pairkernel}
 \mathsf B(m,n)=-2(m^4+n^4)+9(mn^3+nm^3)-14m^2n^2=(m-n)^2(m-2n)(n-2m).
\end{equation}

\begin{theorem}[Exact formal solutions of the order-four quadratic;
source~19]\label{hol:nr:thm:Q4}
The solutions in $K_0[[z]]$ of
$-2f\vartheta^4f+9\vartheta f\,\vartheta^3f-7(\vartheta^2f)^2=0$ are exactly
those of Theorem~\ref{hol:nr:thm:H3}: zero, monomials and doubling binomials.
\end{theorem}

\begin{proof}
Monomials are solutions since $-2+9-7=0$. If $f$ has two terms or more, let
$m<n$ be its first two exponents; the coefficient at $z^{m+n}$ is
$a_ma_n\mathsf B(m,n)$, so \eqref{hol:nr:eq:pairkernel} forces $n=2m$ and
$m>0$. For $\mathsf d_0=a_mz^m+a_{2m}z^{2m}$ the same factorization gives
$\mathsf Q_4(\mathsf d_0)=0$. If $f-\mathsf d_0\ne0$ has first exponent
$r>2m$, the coefficient at $z^{m+r}$ of $\mathsf Q_4(f)$ is $a_ma_r\mathsf
B(m,r)$: the other cross terms start later and the square of $f-\mathsf d_0$
starts at $2r>m+r$. No factor of $\mathsf B(m,r)$ vanishes for $r>2m$.
\end{proof}

Theorems~\ref{hol:nr:thm:H3} and~\ref{hol:nr:thm:Q4} are exact positive
results, not counterexamples to polynomial rigidity. They locate an obstruction
to a \emph{proof technique}: their exterior equations admit nonmonomial
polynomial solutions of arbitrarily large lowest degree. In the language of
Theorem~\ref{hol:ot:thm:gaps} (merge), $\mathsf H_3$ has $i_0=2$ and
$\Upsilon_{\mathsf H_3}(T,G)=(G+T)(2G-T)$, which accepts the upper gap $N/2$
and the lower gap $m$ of a doubling binomial at every scale; $\mathsf Q_4$ has
order four, outside that theorem.

\begin{example}[Valid initial equations at every scale, and no identity;
source~19]\label{hol:nr:ex:sparse}
Over a rank-one workspace with the order unit $1$, for instance
$\KK=k((t^\Q))$, put
\begin{equation}\label{hol:nr:eq:sparse}
 F_{\mathrm{sp}}(z)=\sum_{j\ge0}t^{4^j}z^{2^j}.
\end{equation}
It is strongly entire, since $4^j-2^j\delta\to+\infty$ for every $\delta$. At
the break $\delta=3\cdot2^j$ the active degrees are $2^j$ and $2^{j+1}$ and
$\phi_\delta=X^{2^j}+X^{2^{j+1}}$; between breaks the initial polynomial is a
monomial. So every initial polynomial satisfies both $\mathsf H_3=0$ and
$\mathsf Q_4=0$, yet $F_{\mathrm{sp}}$ satisfies neither: only the degrees $1$
and $4$ contribute to $z^5$ in $\mathsf Q_4(F_{\mathrm{sp}})$, and
\begin{equation}\label{hol:nr:eq:residual}
 [z^5]\,\mathsf Q_4(F_{\mathrm{sp}})=\mathsf B(1,4)\,t^{1+16}=-126\,t^{17}\ne0,
\end{equation}
while Theorem~\ref{hol:nr:thm:H3} excludes $\mathsf H_3(F_{\mathrm{sp}})=0$.
Vanishing leading terms at all scales therefore do not yield an actual
differential identity: higher support layers and the compatibility of
successive breaks carry information that no single initial equation sees. For
Question~\ref{hol:ot:q:lift} this is a concrete negative instance (merge's
reading): an unbounded sequence of initial polynomials, each satisfying the
exterior form of an equation, need not be the initial polynomials of any
solution of that equation, although here they are those of one strongly entire
series.
\end{example}

\begin{corollary}[The theta equation and its jet degree; merge]
\label{hol:nr:cor:theta}
Suppose $\Gamma$ has an order unit $\mu$ and $p=t^\mu$. Every nonzero algebraic
differential equation of order at most three over $\KK(z)$ satisfied by
$\mathcal T_p$ has order exactly three and total jet degree at least four; the
equation of Theorem~\ref{hol:td:thm:ode} has degree six, so the least such
degree is four, five or six. The exterior form \eqref{hol:ot:eq:thetainitial}
of that equation has $i_0=2$ and vanishing first polar, as
Proposition~\ref{hol:nr:prop:firstpolar} requires of every equation of order at
most three with a nonpolynomial strongly entire solution.
\end{corollary}

\begin{proof}
$\mathcal T_p$ is nonpolynomial and strongly entire
(Theorem~\ref{hol:td:main:boundary}); Theorem~\ref{hol:ot:main:second}\,(a)
excludes order at most two and Theorem~\ref{hol:nr:main:third} total degree at
most three. The equation \eqref{hol:td:eq:jetpoly} is homogeneous of degree six
in the jets. The least transverse degree of $W^2-V^2+4V^3$ is two, and with
$\chi=0$ the first polar vanishes exactly when $i_0\ge2$
(Remark~\ref{hol:nr:rem:gaps}); if an equation of order at most three had an
exterior form with $\chi\ne0$ or a nonzero first polar, all its strongly entire
solutions would be polynomials, by the proof of
Theorem~\ref{hol:nr:main:third}.
\end{proof}

Which of four, five and six is the least degree for $\mathcal T_p$ is not
decided here: it depends on whether \eqref{hol:td:eq:jetpoly} is irreducible
over $\KK(z)$, which was not checked.

\subsection{Examples and the role of the hypotheses}
\label{hol:nr:sub:examples}

\begin{proposition}[A large order-unit family; source~19]
\label{hol:nr:prop:family}
Let $\mu>0$ be an order unit of $\Gamma$ and let each $b_n\in\KK$ be zero or
of valuation zero. Then $f(z)=\sum_{n\ge0}b_nt^{n^2\mu}z^n$ is strongly entire.
If infinitely many $b_n$ are nonzero, its first three jets are algebraically
independent over every constant extension $\mathbb L(z)$, and its first four
jets satisfy no polynomial relation of total degree at most three over
$\mathbb L(z)$.
\end{proposition}

\begin{proof}
Given $\delta,\beta\in\Gamma$ choose positive integers $r,s$ with $\delta\le
r\mu$ and $\beta\le s\mu$; then $n^2\mu-n\delta\ge(n^2-rn)\mu>\beta$ for large
$n$. Apply Proposition~\ref{hol:cf:prop:criterion} and
Corollary~\ref{hol:nr:cor:external}.
\end{proof}

With $\Gamma=\Q\subseteq\No$ and $\mu=1$, the surcomplex reading is
$\sum_nb_n\omega^{-n^2}z^n$ on $\C((\omega^{-\Q}))$, and the coefficients
$b_n$, among zero and the units, are arbitrary: the conclusion is not about
one named special function. For $b_n=1$ the series is $G_\mu$, with
$G_\mu(z)=1+t^\mu zG_\mu(t^{2\mu}z)$, that is $G_\mu(z)=\Theta_{p^2}(pz)$ for
$p=t^\mu$; this dilation example is prior material
(Corollary~\ref{hol:nl:cor:theta}, Section~\ref{hol:ud:sub:surreal}), and
three independent jets for $G_1$ are Corollary~\ref{hol:ot:cor:examples}. It
illustrates why differential and dilation equations have different rigidity
thresholds.

\paragraph{A higher-rank comparison (source~19).}
Over $\Gamma_\infty$ of \eqref{hol:eq:gammainfty}, which has no order unit, the
series $F$ of \eqref{hol:eq:explicitF} is strongly entire, and
Theorems~\ref{hol:ot:main:second} and~\ref{hol:nr:main:third} apply to it; but
Theorem~\ref{hol:cf:thm:Fomega} is already stronger for it, in every order.
The comparison explains the division of labour and adds no second novelty
claim. If $\Gamma$ has no countable cofinal subset, every strongly entire
series is a polynomial, since by Proposition~\ref{hol:cf:prop:criterion} at
$\delta=0$ the countably many nonzero coefficient values would be cofinal;
the order-unit examples are not affected by this size obstruction.

\paragraph{The literature boundary (source~19).}
Non-Archimedean functional-equation theory~\cite{huluan} is context: its
absolute-value hypotheses and particular equations are not transferred, and
no theorem is imported from it. The general D-algebraic framework~\cite{ast}
is context; the low-order obstructions here use strong entireness in
addition.

\subsection{Predecessors, repository comparison and priority}
\label{hol:nr:sub:predecessors}

\paragraph{Literature.}
Source~19 consulted the Mantova--Matusinski survey~\cite{mantovamatusinski},
Section~2.2, for normal forms and the $\omega$-map and to distinguish scalar
surreal derivations from $D_z$; Hu--Luan~\cite{huluan} for comparison with
non-Archimedean functional equations; and Ait El Manssour, Sattelberger and
Teguia Tabuguia~\cite{ast} for general differential-algebraic context. No
theorem of these is used in its proofs. It opened the Wikipedia article on
surreal numbers for orientation only and located related coefficient-arithmetic
literature not used as a dependency. Several targeted queries returned
irrelevant results, which support no absence claim; no exhaustive MathSciNet,
zbMATH, book, thesis or citation-network audit was made.

\paragraph{Repository comparison at the pin.}
Source~19 read the files listed in Section~\ref{hol:nr:sub:conventions} at
\texttt{3d40856} through a connector; a code search for the question label
returned nothing, and the question was located through the scope documents; a
container clone failed for lack of network access. This was not a
line-by-line audit of the 91-page report as it then stood, of the 57 reports
then catalogued, of archived submissions, history or Lean declarations, and
directory inventories are not used to prove absence. Its statements were
accurate at the pin; the stale ones are recorded in
Section~\ref{hol:nr:sub:conventions}.

\paragraph{Priority.}
Second-order rigidity is Theorem~\ref{hol:ot:main:second}\,(a), merged from
source~14; source~19's proof is an independent re-derivation written the same
day against a text of this report that did not contain source~14, and no
priority between the two is asserted. The
all-scale criterion, the corner technique, the no-order-unit results and the
partial-theta witness are credited as prior. The proposed contributions of
source~19 that are new to this report are third-order rigidity through cubic
jet degree (Theorem~\ref{hol:nr:main:third}) with the first-polar and factor
criteria, the secant lemma and the eventual-side lemma; the external-constant,
Gaussian-omnific and algebraic two-state consequences; and the exact
classifications of $\mathsf H_3$ and $\mathsf Q_4$ with the sparse series.
Whether equivalent results exist in the literature needs a more exhaustive
review, and no named conjecture is claimed solved.

\subsection{What the Newton-rigidity part does not claim}
\label{hol:nr:sub:nonclaims}

Every limitation stated by source~19, in its manuscript, README, audit files or
verification record, is kept.
\begin{enumerate}[label=\textup{(R\arabic*)}]
\item No all-order rigidity with an order unit is claimed (and none holds:
Theorem~\ref{hol:td:main:boundary}); the unrestricted order-unit question was
not resolved by source~19.
\item No meromorphic order-unit analogue is proved: for quotients of entire
series with an order unit nothing beyond
Corollary~\ref{hol:ot:cor:meromorphic} is claimed.
\item All-order differential independence is not inferred from three-jet
independence.
\item No independence statement concerns point values: the relations are
formal, and $f(x)\in\KK$ for every $x$.
\item No omnific factorization or arithmetic structure theorem is claimed; the
omnific consequence concerns finitely many coefficients of formal relations.
\item No uniform degree bound for polynomial solutions of a fixed equation is
claimed; \eqref{hol:nl:eq:degenerate} has solutions of every degree, and the
systems results assert neither existence, nor finiteness, nor bounded
degrees of polynomial solutions.
\item Entireness is relative to one Hahn field; nothing is claimed about
entireness on the proper class $\No[i]$, where no nonpolynomial series is
entire (Proposition~\ref{hol:ot:prop:fullclass}).
\item $D_z$ acts on the external variable and kills every scalar; it is not the
Berarducci--Mantova derivation, and no theorem about that derivation is
claimed.
\item External constants change the coefficients allowed in a relation, not
the domain of strong summation.
\item The first-polar criterion and the factor criterion are sufficient, not
necessary (Theorems~\ref{hol:nr:thm:H3} and~\ref{hol:nr:thm:Q4}).
\item $\mathsf H_3$, $\mathsf Q_4$ and $F_{\mathrm{sp}}$ are not counterexamples to
polynomial rigidity; an obstruction to a proof technique is not a
counterexample.
\item The 259 finite assertions check polynomial identities and finite Newton
windows; they certify nothing about infinite supports, cofinality, arbitrary
fields or novelty, and a large count is no substitute for the written proofs.
\item No Lean proof or checked Lean theorem is supplied, the repository was not
built by source~19, its formalization claims are not certificates for these
results, and no independent referee report or historical priority is claimed.
\item The repository and literature comparisons were targeted, not exhaustive,
as recorded in Section~\ref{hol:nr:sub:predecessors}.
\item Positive characteristic and mixed characteristic are excluded: the
proofs use $v(n)=0$ for nonzero integers, and the theorems are not statements
about $\mathfrak p$-adic analytic fields.
\item \emph{[merge]} Remark~\ref{hol:nr:rem:gaps} with its cubic,
Corollary~\ref{hol:nr:cor:theta}, the reading of
Example~\ref{hol:nr:ex:sparse} for Question~\ref{hol:ot:q:lift}, and the
comparison of Theorem~\ref{hol:nr:thm:systems} with
Corollary~\ref{hol:ot:cor:systems} are the merge's; no historical priority is
claimed for them.
\end{enumerate}

%======================================================================
% Theta-hierarchy part: source 16. Labels use hol:th:.
%======================================================================
\section{Independent theta scales and unbounded differential order}
\label{hol:th:sec:setting}

This section and the next contain the material of source~16, delivered in
batch~31 with sources~15, 17, 18 and~19, and merged after the order-three-threshold part. They
take source~13's descended theta series as given and ask how many
differentially algebraic strongly entire series there are, and how large
their minimal orders can be. Over $\KK=k((t^\R))$ with $k=\R$ or $\C$ the
answer is: the descents $\mathcal T_{t^\mu}$ for scales $\mu$ with
$\Q$-independent reciprocals are algebraically independent, a family of
cardinality $2^{\aleph_0}$; polynomially weighted sums of them have
arbitrarily large minimal differential order; and the field generated by all
differentially algebraic strongly entire series has transcendence degree
$2^{\aleph_0}$ and differential transcendence degree zero
(Theorem~\ref{hol:th:main:hierarchy}). The mechanism is the exact quadratic
Gauss profile of $\mathcal T_{t^\mu}$ together with a finite-grid compression
lemma of differential algebra. Source~16 also gives a third route to the
minimal order three of $\mathcal T_p$ (Section~\ref{hol:th:sub:third}).

\subsection{Source, conventions and the merge's choices}
\label{hol:th:sub:conventions}

\begin{convention}[Setting of the theta-hierarchy part]
\label{hol:th:conv:field}
In Sections~\ref{hol:th:sec:setting}--\ref{hol:th:sec:scope}, unless a
statement says otherwise, $\Gamma=\R$, $k$ is $\R$ or $\C$, $\KK=k((t^\R))$,
and $\EK$, $\FE$, $v$, $\res$, $D_z$ and the Gauss value
$\gamma_\delta(f)=\min_n(v(a_n)-n\delta)$ of \eqref{hol:nl:eq:gaussdata}, here
for real $\delta$ and extended multiplicatively to $\FE$, are as in
Conventions~\ref{hol:nl:conv:field} and~\ref{hol:cf:conv:field}. For a real
$\mu>0$ the series $\mathcal T_{t^\mu}$ is \eqref{hol:td:eq:series} with $p=t^\mu$;
every positive real is an order unit of $\R$. The abstract results of
Section~\ref{hol:th:sub:compression} hold in any differential field of
characteristic zero; there $\partial$ is the derivation and $\mathrm d$ the
universal derivation of K\"ahler differentials, which are unrelated to the
partial derivatives $\partial_i$ of Section~\ref{hol:sec:multi}.
\end{convention}

\paragraph{The source.}
Source~16 is \emph{Independent Theta Families and Unbounded Differential Order
in Hahn--Surcomplex Analysis: Valuation growth, bounded polynomial jet
compression, and a continuum-sized differential-algebraic function field}
(25 pages, A4, 23 September 2026), delivered as the archive
\texttt{Surreal\_Theta\_Hierarchy} (batch~31, manuscript~02). It records the
reference snapshot
\begin{center}\small\texttt{0fffc26c670c15c7db451ed899f98d4384f9df1e},\end{center}
taken while the repository was changing. At that snapshot this report's
article and README were the five-source text (sources~08--12), and source~13's
two audit files had been placed in this directory without its manuscript being
merged; source~16 read those audits and took the seed, its equation and its
minimal order three from them, crediting all three to source~13. Its
description was accurate at the snapshot; its statement that ``the repository
explicitly leaves the order-two alternative unresolved'' and its Question~11.2
are stale now, answered by Theorem~\ref{hol:ot:main:second}. Source~16 works
over $\Gamma=\R$ only, with $k=\C$ and, by the same proofs, $k=\R$; its
asymptotics are real-valued, and no statement of this part is transferred to
other value groups.

\paragraph{Renamed symbols.}
The following symbols of source~16 would collide with symbols of the earlier
parts. They are renamed here, in the new material only.
{\small
\begin{longtable}{@{}>{\raggedright\arraybackslash}p{0.27\textwidth}
>{\raggedright\arraybackslash}p{0.2\textwidth}
>{\raggedright\arraybackslash}p{0.45\textwidth}@{}}
\toprule
Source~16 & Here & Reason\\
\midrule
\endhead
function variable $X$; $\partial=\mathrm d/\mathrm dX$; $\mathbb F=\mathbb K(X)$ & $z$;
$D_z$; $\KK(z)$ & The notation of Convention~\ref{hol:nl:conv:field};
$X$ is the variable of initial polynomials.\\
$\mathbb K=\C((t^\R))$, $\mathbb K_\R=\R((t^\R))$; $\mathcal E_{\mathbb K}$ & $\KK$ with $k=\C$, $\R$;
$\EK$ & The same objects; only the real-coefficient field is renamed.\\
scale $h$; $q=t^h$; $Q=q^2$; $s_k$ & $\mu$; $p=t^\mu$; $p^2$;
$\mathcal Q_k(p^2)$ & As in Section~\ref{hol:td:sub:conventions}; $q$ is a
dilation parameter.\\
$F_h$ & $\mathcal T_{t^\mu}$ & The series \eqref{hol:td:eq:series}; $F$ is
\eqref{hol:eq:explicitF}.\\
Laurent variable $z$; $\Theta=\sum q^{n^2}z^n$; $D=z\,\mathrm d/\mathrm dz$ &
$u$; $\Psi_p$; $\vartheta_u$ & Section~\ref{hol:td:sub:laurent}; $\Theta_p$ is
the partial theta series.\\
$G=D\Theta/\Theta$; $H=-DG-2s_1$; $Y=z-z^{-1}$ & $\psi_1$;
$\mathsf x=-\psi_2-2\mathcal Q_1(p^2)$; $\check z$ & $G$, $H$ have other
meanings; $\mathsf x$ is a Tate $x$-coordinate.\\
$\Delta=X^2-4$; $A$, $B$ of its (3.3) & $z^2-4$; $\mathcal B[f]-\frac1{12}f^2$,
$\mathcal H[f]$ & $\Delta_\mu$ is a convex subgroup; \eqref{hol:td:eq:H},
\eqref{hol:td:eq:KB}.\\
$a_4$, $a_6$ (Tate) & kept & Related to $g_2,g_3$ in
Section~\ref{hol:th:sub:conventions}.\\
radius parameter $r$; $w_r$ & $\delta$; $\gamma_\delta$ & The Gauss data of
Section~\ref{hol:nl:sub:scaled}; $r$ is a derivative order.\\
$\operatorname{ord}_\partial(f)$ & $\operatorname{dord}(f)$ & $\operatorname{ord}_z$ is the
$z$-adic order.\\
$\mathfrak c=2^{\aleph_0}$ & $2^{\aleph_0}$ & $\mathfrak c$ is a theta constant.\\
$F\subseteq L$, $s$, $a_i$, $p_i(s)$, $c_{ij}$ (compression) &
$\mathfrak K_0\subseteq\mathfrak K_1$, $\mathsf s$, $y_i$, $\mathsf P_i(\mathsf
s)$, $\mathsf c_{ij}$ & $L$ is a coefficient subfield, $a_n$ Taylor
coefficients, $p=t^\mu$.\\
$W$, $U_{ir}$, $\lambda_i$ & $\mathsf W$, $\mathsf U_{ir}$, $\varepsilon_i$ &
$W,U$ are indeterminates, $\lambda$ a multiplier.\\
primes $p_i$ & $\mathrm p_i$ & $p=t^\mu$.\\
$\boldsymbol m$, $p_{\boldsymbol m}$, $A_{\boldsymbol m}$, $B_{\boldsymbol m}$,
$\delta_{\boldsymbol m}$, $b_{\boldsymbol m}$, $c_{\boldsymbol m}$, $R_0$, $R$ &
$\mathbf e$, $P_{\mathbf e}$, $\kappa_{\mathbf e}$, $\kappa'_{\mathbf e}$,
$\kappa_*-\kappa_{\mathbf e}$, $\bar d_{\mathbf e}$, $\bar\alpha_{\mathbf e}$,
$\delta^{(0)}$, $\delta^{\mathrm{cert}}$ & Radii are scales $\delta$ here.\\
Hamel set $B$ & $\mathfrak V$ & $B$ is a threshold.\\
$\mathcal A$, $\mathcal D$ & $\mathcal A_{\mathrm{DA}}$, $\mathcal D_{\mathrm{DA}}$
& $\mathcal A[f]$ is a theta expression.\\
$G$, $\mu=w_0(G)$, $\overline G$ (sampling) & $g$, $\gamma_0(g)$, $\phi_0$ &
$\phi_0$ is the initial polynomial at the scale $0$.\\
Laurent-entire ring $\mathcal L$; $w$; $\sigma$ & $\mathcal U_\mu$;
$\gamma^{\mathcal U}$; $\mathsf S$ & Section~\ref{hol:td:sub:laurent}.\\
Theorems~1.1--1.3 & Theorem~\ref{hol:th:main:hierarchy} & Theorems~A--N of
this report are different theorems.\\
\bottomrule
\end{longtable}}
The shipped files \path{16-theta-hierarchy-PROOF_AUDIT.md} and
\path{16-theta-hierarchy-SOURCES_AND_SCOPE.md} keep the source's notation.

\paragraph{Words with two meanings.}
The \emph{minimal differential order} of $f$, written $\operatorname{dord}(f)$,
is the least $r$ with $f,D_zf,\ldots,D_z^rf$ algebraically dependent over
$\KK(z)$; it is the
least order of Section~\ref{hol:td:sub:conventions}, not the degree of an
equation and not the growth order of a complex entire function. A
\emph{radius} $t^{-\delta}$ is a Hahn monomial, not a real radius.
\emph{Compression} is the finite-grid lemma of
Section~\ref{hol:th:sub:compression}, not a primitive-element theorem.
\emph{Independent} without qualifier means algebraically independent over
$\KK(z)$ \emph{without} derivatives; differential independence is stated as
such.

\paragraph{Results printed once.}
Source~16 reproves several facts already in the report. They are printed
once and cited: its escape lemma (Lemma~2.1), which is
Lemma~\ref{hol:lem:certificate}\,(a); its definition of entireness
(Definition~2.2), which is Definition~\ref{hol:def:entire} in the form of
Proposition~\ref{hol:cf:prop:criterion}; closure and evaluation
(Proposition~2.3), which is Proposition~\ref{hol:prop:closure}; Gauss
multiplicativity (Lemma~2.5), which is Corollary~\ref{hol:nl:cor:gauss}; the
coefficients of the seed (Proposition~3.1), which are
Proposition~\ref{hol:td:prop:coeff}; its product (3.4), which is
\eqref{hol:td:eq:product}; its explicit equation (Proposition~3.2), which is
\eqref{hol:td:eq:cleared} (see below); its Proposition~9.1, the failure at
$\omega^\omega$, which is in Section~\ref{hol:td:sub:surreal}; its Laurent
ring and shift (Appendix~A.1), which are Lemma~\ref{hol:td:lem:laurent} and
\eqref{hol:td:eq:shift}; its Lemma~A.1, whose proof is that of
Lemma~\ref{hol:lem:integer} over any valued field of residue characteristic
zero; and the embedding \eqref{hol:eq:embedding}.

Source~16 states the seed equation in the Tate normalization, with the
constants $a_4=-5\mathcal Q_3(p^2)$, $a_6=-(5\mathcal Q_3(p^2)+7\mathcal
Q_5(p^2))/12$ and $\mathsf A=\mathcal B[f]-f^2/12$:
\[
 (z^2-4)\bigl(f\,D_z\mathsf A-2D_zf\,\mathsf A\bigr)^2=4\mathsf A^3+\mathsf
 A^2f^2+4a_4\mathsf Af^4+4a_6f^6 .
\]
Since $f\,D_z\mathsf A-2D_zf\,\mathsf A=-\mathcal K[f]$,
$g_2=1/12-4a_4$ and $g_3=-1/216+a_4/3-4a_6$, this is
\eqref{hol:td:eq:cleared} as a polynomial identity in $z$ and the jets (the
merge's check, with SymPy), and its leading coefficient $(z^2-4)^3f^4$ in
$(D_z^3f)^2$ agrees with Theorem~\ref{hol:td:thm:ode}. The Tate cubic
$y^2+xy=x^3+a_4x+a_6$ is the Weierstrass form after the shift $x\mapsto
x-1/12$; source~16 imports it from the course notes~\cite{woodbury}, as
source~13 did.

\paragraph{Where each numbered result went.}
Theorems~1.1--1.3 are Theorem~\ref{hol:th:main:hierarchy}\,(a)--(c);
Lemma~2.6 is Lemma~\ref{hol:th:lem:poly}; Theorems~4.1, 4.2,
Proposition~4.4 and Example~4.5 are Theorems~\ref{hol:th:thm:profile},
\ref{hol:th:thm:growth}, Proposition~\ref{hol:th:prop:threshold} and
Example~\ref{hol:th:ex:two}; Lemmas~5.1, 5.2, Theorem~5.3, Remark~5.4,
Example~5.5 and Corollary~5.6 are Lemmas~\ref{hol:th:lem:kahler},
\ref{hol:th:lem:grid}, Theorem~\ref{hol:th:thm:compression},
Remark~\ref{hol:th:rem:specialization}, Example~\ref{hol:th:ex:small} and
Corollary~\ref{hol:th:cor:entirecompression}; Definition~6.1, Lemma~6.2,
Corollaries~6.3, 6.4, Remark~6.5 and Theorem~6.6 are
Definition~\ref{hol:th:def:dord}, Lemma~\ref{hol:th:lem:ordertrdeg},
Corollaries~\ref{hol:th:cor:composita}, \ref{hol:th:cor:nobound},
Remark~\ref{hol:th:rem:limitation} and Theorem~\ref{hol:th:thm:exactdimension};
Lemmas~7.1, 7.3 and Corollary~7.2 are Lemmas~\ref{hol:th:lem:hamel},
\ref{hol:th:lem:card} and Corollary~\ref{hol:th:cor:continuum};
Theorem~8.1 and Corollaries~8.2, 8.3 are Theorem~\ref{hol:th:thm:sampling}
and Corollaries~\ref{hol:th:cor:samplingrelations},
\ref{hol:th:cor:jetsampling}; Theorem~A.2 and Corollary~A.3 are
Theorem~\ref{hol:th:thm:seedthree} and Corollary~\ref{hol:th:cor:seedminimal};
Questions~11.1 and 11.3--11.10 are
Questions~\ref{hol:th:q:spectrum}--\ref{hol:th:q:omnific}, and Question~11.2 is
answered (Remark~\ref{hol:th:rem:ordertwo}); Appendix~B is merged into the
non-claims of Section~\ref{hol:th:sub:nonclaims} and the dependency table of
Section~\ref{hol:sec:verification}.

\paragraph{The merge's choices.}
The review of source~16 for this merge found its proofs correct. The merge
adds four points, each marked. First, the identity of the two seed equations
above. Second, Remark~\ref{hol:th:rem:ordertwo}: with
Theorems~\ref{hol:ot:main:second} and~\ref{hol:nr:main:third}, the minimal
orders of nonpolynomial differentially algebraic strongly entire series over
$\KK$ form a subset of $\{3,4,\ldots\}$ that contains $3$ and is unbounded, and
the compressed series of Theorem~\ref{hol:th:main:hierarchy}\,(b) have order at
least $\max(N,3)$. Third, Remark~\ref{hol:th:rem:fields} compares
Theorem~\ref{hol:th:main:hierarchy}\,(c) with Theorem~\ref{hol:cf:thm:continuum}
and with source~17's Corollary~\ref{hol:fh:cor:entirecard}. Fourth,
Remark~\ref{hol:th:rem:thirdroute} compares the three routes to minimal order
three. Source~16's theorems, proofs, examples and non-claims are otherwise
printed as delivered, in the renamed notation.

\subsection{Gauss profiles and independence across scales}
\label{hol:th:sub:profiles}

\begin{lemma}[Polynomial growth; source~16]\label{hol:th:lem:poly}
If $P=\sum_{i=0}^dc_iz^i\in\KK[z]$ with $c_d\ne0$, then
$\gamma_\delta(P)=v(c_d)-d\delta$ for all sufficiently large real $\delta$;
for a nonzero rational function $\gamma_\delta$ is eventually an affine
function of $\delta$.
\end{lemma}

\begin{proof}
For each $i<d$ with $c_i\ne0$ the top term has strictly smaller weighted value
once $\delta>(v(c_d)-v(c_i))/(d-i)$; there are finitely many such
inequalities. Subtract for numerator and denominator.
\end{proof}

\begin{theorem}[Exact quadratic profile; source~16]\label{hol:th:thm:profile}
For real $\mu>0$ and $\delta\ge0$,
\[
 \gamma_\delta(\mathcal T_{t^\mu})=\min_{n\ge0}(\mu n^2-\delta n)
 =-\frac{\delta^2}{4\mu}+\mu\,\operatorname{dist}\Bigl(\frac\delta{2\mu},
 \Z_{\ge0}\Bigr)^2,
\]
so $-\delta^2/(4\mu)\le\gamma_\delta(\mathcal T_{t^\mu})\le-\delta^2/(4\mu)+\mu/4$;
for $\delta<0$, $\gamma_\delta(\mathcal T_{t^\mu})=0$.
\end{theorem}

\begin{proof}
By \eqref{hol:td:eq:coeffval}, $v(a_n)=\mu n^2$. Complete the square,
$\mu n^2-\delta n=\mu(n-\delta/(2\mu))^2-\delta^2/(4\mu)$, and use that a
nonnegative real lies within $\frac12$ of a nonnegative integer. For
$\delta<0$ the expression is positive for $n>0$ and zero at $n=0$.
\end{proof}

The corners are at $\delta=(2n+1)\mu$, where $n$ and $n+1$ are active: these
are the breaks $\delta_n$ of Theorem~\ref{hol:ot:thm:convexity}, with
$\gamma_{\delta_n}=-n(n+1)\mu$. The bounded oscillation lets distinct growth
constants be compared uniformly, without special sequences of scales.

\begin{theorem}[Superlinear growth separation; source~16]
\label{hol:th:thm:growth}
Let $(f_i)_{i\in I}$ be nonzero elements of $\EK$ and $\varpi>1$ real, and
suppose $\gamma_\delta(f_i)=-c_i\delta^\varpi+o(\delta^\varpi)$ as
$\delta\to+\infty$, with positive reals $c_i$ linearly independent over $\Q$.
Then $(f_i)_{i\in I}$ is algebraically independent over $\KK(z)$.
\end{theorem}

\begin{proof}
A relation involves finitely many $f_i$; clear denominators to get
$\sum_{\mathbf e\in E}P_{\mathbf e}(z)\prod_{i=1}^Nf_i^{e_i}=0$ with $E\subset
\Z_{\ge0}^N$ finite and every $P_{\mathbf e}\ne0$. By
Corollary~\ref{hol:nl:cor:gauss} and Lemma~\ref{hol:th:lem:poly} the summand
indexed by $\mathbf e$ has $\gamma_\delta=-(\sum_ie_ic_i)\delta^\varpi+
o(\delta^\varpi)+O(\delta)$. The numbers $\sum_ie_ic_i$ are pairwise distinct,
by the $\Q$-independence, so one is largest, and since $\varpi>1$ its summand
has strictly the least Gauss value for all large $\delta$; a finite sum with a
unique least term is nonzero.
\end{proof}

\begin{proof}[Proof of Theorem~\ref{hol:th:main:hierarchy}\,\textup{(a)}]
Apply Theorem~\ref{hol:th:thm:growth} with $\varpi=2$ and $c_i=1/(4\mu_i)$,
using the bounds of Theorem~\ref{hol:th:thm:profile}. Each
$\mathcal T_{t^{\mu_i}}$ is a nonpolynomial element of $\EK$ of minimal order
three (Theorem~\ref{hol:td:main:boundary}); the cardinality assertion is
Corollary~\ref{hol:th:cor:continuum}.
\end{proof}

The coefficients of a relation may be arbitrary elements of $\KK$: a fixed
polynomial in $z$ has finitely many coefficients, each with one valuation, so
Lemma~\ref{hol:th:lem:poly} still gives only linear growth. That is why the
quadratic separation survives the whole Hahn coefficient field.

\begin{proposition}[An explicit noncancellation threshold; source~16]
\label{hol:th:prop:threshold}
In a relation as in the proof of Theorem~\ref{hol:th:thm:growth} with
$f_i=\mathcal T_{t^{\mu_i}}$, put $\kappa_{\mathbf e}=\frac14\sum_ie_i/\mu_i$,
$\kappa'_{\mathbf e}=\frac14\sum_ie_i\mu_i$, $d_{\mathbf e}=\deg P_{\mathbf e}$
and $\alpha_{\mathbf e}=v(\lc P_{\mathbf e})$, let $\mathbf e_*$ be the unique
maximizer of $\kappa_{\mathbf e}$, write $\kappa_*,d_*,\alpha_*,\kappa'_*$ for
its values, and let $\delta^{(0)}$ be the largest of $0$ and the numbers
$(\alpha_{\mathbf e}-v(c_{\mathbf e,i}))/(d_{\mathbf e}-i)$ over all
coefficients $c_{\mathbf e,i}\ne0$ of $P_{\mathbf e}=\sum_ic_{\mathbf e,i}z^i$
with $i<d_{\mathbf e}$. For $\mathbf e\ne\mathbf e_*$ put
$\bar d_{\mathbf e}=\abs{d_*-d_{\mathbf e}}$ and
$\bar\alpha_{\mathbf e}=\abs{\alpha_{\mathbf e}-\alpha_*}+\kappa'_*$. Then
every real $\delta\ge\delta^{\mathrm{cert}}$, where
\[
 \delta^{\mathrm{cert}}=1+\max\Bigl(\{\delta^{(0)}\}\cup\Bigl\{
 \frac{2\bar d_{\mathbf e}}{\kappa_*-\kappa_{\mathbf e}},\
 \sqrt{\frac{2\bar\alpha_{\mathbf e}}{\kappa_*-\kappa_{\mathbf e}}}:
 \mathbf e\ne\mathbf e_*\Bigr\}\Bigr),
\]
makes the $\mathbf e_*$ summand the unique one of least Gauss value; in
particular the expression is nonzero.
\end{proposition}

\begin{proof}
For $\delta>\delta^{(0)}$ the summand $\mathbf e$ has Gauss value
$-\kappa_{\mathbf e}\delta^2-d_{\mathbf e}\delta+\alpha_{\mathbf e}+
\epsilon_{\mathbf e}(\delta)$ with $0\le\epsilon_{\mathbf e}(\delta)\le
\kappa'_{\mathbf e}$, by Theorem~\ref{hol:th:thm:profile}. Its difference from
the $\mathbf e_*$ value is at least $(\kappa_*-\kappa_{\mathbf e})\delta^2-
\bar d_{\mathbf e}\delta-\bar\alpha_{\mathbf e}$, and the two conditions in
$\delta^{\mathrm{cert}}$ make each of the last two terms smaller than half of
the first. A single summand needs only $\delta^{(0)}$.
\end{proof}

This certifies a specified polynomial expression; it is not a decision
procedure for differential relations.

\begin{example}[Two explicit scales; source~16]\label{hol:th:ex:two}
Take $\mu_1=1$, $\mu_2=1/\sqrt2$ and $P=Y_1^2+(z+t^{-7})Y_2+1$. The growth
constants of the three summands are $\frac12$, $\frac{\sqrt2}4$ and $0$, so
the $Y_1^2$ summand dominates; $\delta^{(0)}=7$, and $\delta\ge16$ is a valid
certificate (the exact value is about $14.7$). So
$\mathcal T_t^2+(z+t^{-7})\mathcal T_{t^{1/\sqrt2}}+1\ne0$, with the infinite
coefficient $t^{-7}$ allowed.
\end{example}

\subsection{Bounded polynomial jet compression}
\label{hol:th:sub:compression}

The valuation argument yields independent \emph{functions}, not one function
of high order; a plain sum can lose information. The following is pure
differential algebra, related in purpose to differential primitive-element
theorems~\cite{pogudin2014,pogudin2018} but different in conclusion: it keeps
a prescribed number of independent jets, with an explicit grid of integer
weights, and does not claim a generator of the whole extension.

\begin{lemma}[K\"ahler criterion; source~16]\label{hol:th:lem:kahler}
Let $\mathfrak K_0\subseteq\mathfrak K_1$ be fields of characteristic zero and
$y_1,\ldots,y_N\in\mathfrak K_1$. They are algebraically independent over
$\mathfrak K_0$ if and only if $\mathrm dy_1\wedge\cdots\wedge\mathrm dy_N\ne0$
in $\bigwedge^N_{\mathfrak K_1}\Omega_{\mathfrak K_1/\mathfrak K_0}$.
\end{lemma}

\begin{proof}
If they are independent, extend them to a transcendence basis; the coordinate
derivations of the rational function field extend to $\mathfrak K_1$ by
separability and give a dual system to the $\mathrm dy_i$. Conversely, if they
are dependent, every $y$ outside a maximal independent subfamily is algebraic
over it, and differentiating its minimal polynomial (nonzero derivative by
separability) makes $\mathrm dy$ a combination of the others, so the wedge is
zero.
\end{proof}

\begin{lemma}[Finite-grid nonvanishing; source~16]\label{hol:th:lem:grid}
A nonzero polynomial in $M$ variables over a field of characteristic zero, of
degree at most $N$ in each variable, is nonzero at some point of
$\{0,1,\ldots,N\}^M$.
\end{lemma}

\begin{proof}
For $M=1$ this is the root bound. Otherwise write the polynomial in the last
variable; vanishing on the grid makes every specialization of the first $M-1$
variables have $N+1$ roots, so every coefficient vanishes on the smaller
grid, and induction applies.
\end{proof}

\begin{theorem}[Bounded polynomial jet compression; source~16]
\label{hol:th:thm:compression}
Let $(\mathfrak K_0,\partial)\subseteq(\mathfrak K_1,\partial)$ be differential
fields of characteristic zero, $\mathsf s\in\mathfrak K_0$ with
$\partial\mathsf s=1$, and $y_1,\ldots,y_N\in\mathfrak K_1$ algebraically
independent over $\mathfrak K_0$. There are
$\mathsf P_i(\mathsf s)=\sum_{j=0}^{N-1}\mathsf c_{ij}\mathsf s^j$ with
$\mathsf c_{ij}\in\{0,1,\ldots,N\}$ such that for
$g=\sum_i\mathsf P_i(\mathsf s)y_i$ the elements
$g,\partial g,\ldots,\partial^{N-1}g$ are algebraically independent over
$\mathfrak K_0$. At most $(N+1)^{N^2}$ coefficient arrays need be considered.
\end{theorem}

\begin{proof}
Treat the $\mathsf c_{ij}$ first as indeterminates, constant for $\partial$
and killed by $\mathrm d$, which is relative to $\mathfrak K_0$. The wedge
\[
 \mathsf W(\mathbf{\mathsf c})=\mathrm dg\wedge\mathrm d(\partial g)\wedge\cdots
 \wedge\mathrm d(\partial^{N-1}g)\in\Bigl(\bigwedge\nolimits^N_{\mathfrak K_1}
 \Omega_{\mathfrak K_1/\mathfrak K_0}\Bigr)[\mathbf{\mathsf c}]
\]
has total degree $N$, each factor being linear in the parameters. Replace them
by the jet parameters $\mathsf U_{ir}=\partial^r\mathsf P_i(\mathsf s)$,
$0\le r<N$; this change is invertible over $\mathfrak K_0$, triangular with the
factorials $0!,\ldots,(N-1)!$ on the diagonal, so it suffices to show
$\mathsf W\ne0$ in the new variables. By Leibniz,
\[
 \mathrm d(\partial^jg)=\sum_{i=1}^N\sum_{r=0}^j\binom jr\mathsf U_{ir}\,
 \mathrm d(\partial^{j-r}y_i),
\]
with no differentials of the $\mathsf U_{ir}$, whose specialized values lie in
$\mathfrak K_0$. Specialize $\mathsf U_{ir}=\varepsilon_i$ if $r=i-1$ and $0$
otherwise, with new indeterminates $\varepsilon_i$. The $j$th factor then
involves only $\varepsilon_1,\ldots,\varepsilon_{j+1}$, so the squarefree
monomial $\varepsilon_1\cdots\varepsilon_N$ arises only when the $j$th factor
supplies $\varepsilon_{j+1}\mathrm dy_{j+1}$, with binomial coefficient one; its
coefficient is $\mathrm dy_1\wedge\cdots\wedge\mathrm dy_N\ne0$ by
Lemma~\ref{hol:th:lem:kahler}. So $\mathsf W$ is a nonzero polynomial. A
$\mathfrak K_1$-linear functional on the finite span of its coefficients not
killing a nonzero one gives a nonzero scalar polynomial of degree at most $N$
in each of the $N^2$ variables; Lemma~\ref{hol:th:lem:grid} gives a point of
$\{0,\ldots,N\}^{N^2}$ where it, hence $\mathsf W$, is nonzero, and
Lemma~\ref{hol:th:lem:kahler} concludes.
\end{proof}

\begin{remark}[Why the formal jet specialization is legitimate; source~16]
\label{hol:th:rem:specialization}
The specialization $\mathsf U_{ir}=\varepsilon_i\mathbf 1_{r=i-1}$ only proves
that a polynomial is nonzero; it is not a claim that one polynomial in
$\mathsf s$ has independently prescribed derivatives. Invertibility of the
coefficient-to-jet change over $\mathfrak K_0$ is what makes the test valid.
The final specialization is in the original coefficients, where the grid gives
actual integer polynomials.
\end{remark}

\begin{example}[The first two cases; source~16]\label{hol:th:ex:small}
For $N=1$ the grid contains $g=y_1$. For $N=2$, among the $81$ candidates
$(\mathsf c_{10}+\mathsf c_{11}\mathsf s)y_1+(\mathsf c_{20}+\mathsf c_{21}
\mathsf s)y_2$ with $\mathsf c_{ij}\in\{0,1,2\}$ at least one has $g,\partial
g$ independent; an unweighted sum $y_1+y_2$ is not asserted to work. The grid
is uniform: it works for every independent pair in the stated setting.
\end{example}

\begin{corollary}[An entire-preserving version; source~16]
\label{hol:th:cor:entirecompression}
If $y_1,\ldots,y_N\in\EK$ are algebraically independent over $\KK(z)$, one of
the candidates of Theorem~\ref{hol:th:thm:compression} with $\mathsf s=z$ is
strongly entire and has its first $N$ jets algebraically independent over
$\KK(z)$; with real Hahn coefficients, so does the candidate.
\end{corollary}

\begin{proof}
Apply the theorem in $\FE$; polynomial weights and finite sums preserve
$\EK$ (Proposition~\ref{hol:prop:closure}), and integer weights preserve the
real subfield.
\end{proof}

\subsection{Minimal differential order and the unbounded hierarchy}
\label{hol:th:sub:order}

\begin{definition}[Minimal differential order; source~16]\label{hol:th:def:dord}
For a differential field extension $\mathfrak K_0\subseteq\mathfrak K_1$ and
$f\in\mathfrak K_1$ differentially algebraic over $\mathfrak K_0$, put
$\operatorname{dord}(f/\mathfrak K_0)=\min\{r\ge0:f,\partial f,\ldots,
\partial^rf$ are algebraically dependent over $\mathfrak K_0\}$; write
$\operatorname{dord}(f)$ when $\mathfrak K_0=\KK(z)$, and
$\mathfrak K_0\langle f\rangle$ for the differential field generated by $f$.
\end{definition}

A rational function has order zero. For a nonpolynomial $f\in\EK$ the number
$\operatorname{dord}(f)$ is the least order of Theorem~\ref{hol:ot:main:second}
and Section~\ref{hol:td:sub:conventions}.

\begin{lemma}[Order equals transcendence degree; source~16]
\label{hol:th:lem:ordertrdeg}
In characteristic zero, if $f$ is differentially algebraic over
$\mathfrak K_0$ and $r=\operatorname{dord}(f/\mathfrak K_0)$, then
$\mathfrak K_0\langle f\rangle=\mathfrak K_0(f,\partial f,\ldots,\partial^rf)$
and its transcendence degree over $\mathfrak K_0$ is $r$.
\end{lemma}

\begin{proof}
$f,\ldots,\partial^{r-1}f$ are independent and $\partial^rf$ is algebraic over
the field they generate; differentiate its minimal polynomial, whose
derivative at $\partial^rf$ is nonzero by separability, and solve for
$\partial^{r+1}f$ inside $\mathfrak K_0(f,\ldots,\partial^rf)$, which is
therefore closed under $\partial$. This is the argument of
Theorem~\ref{hol:nr:thm:systems} and of the finite differential fields of
Theorem~\ref{hol:td:thm:field}; it is standard.
\end{proof}

\begin{corollary}[Finite composita; source~16]\label{hol:th:cor:composita}
If $f_1,\ldots,f_N$ are differentially algebraic over $\mathfrak K_0$ of orders
$r_1,\ldots,r_N$, then $\mathfrak K_0\langle f_1,\ldots,f_N\rangle$ is a finitely
generated field extension of transcendence degree at most $r_1+\cdots+r_N$,
and each of its elements is differentially algebraic of order at most that
transcendence degree.
\end{corollary}

\begin{proof}
The compositum of the fields of Lemma~\ref{hol:th:lem:ordertrdeg} is closed
under $\partial$; any $d+1$ jets of an element of a field of transcendence
degree $d$ are dependent.
\end{proof}

\begin{proof}[Proof of Theorem~\ref{hol:th:main:hierarchy}\,\textup{(b)}]
Take $\mu_i=1/\sqrt{\mathrm p_i}$ for distinct primes $\mathrm p_i$. The
$\mathcal T_{t^{\mu_i}}$ are independent by part~(a), and
Corollary~\ref{hol:th:cor:entirecompression} gives an entire candidate $g$
with $N$ independent jets, so $\operatorname{dord}(g)\ge N$. Each seed has an
equation of order three (Theorem~\ref{hol:td:thm:ode}), so by
Corollary~\ref{hol:th:cor:composita} $g$ is differentially algebraic with
$\operatorname{dord}(g)\le3N$; being of order at least $N\ge1$ it is not
rational, hence not a polynomial.
\end{proof}

The square roots of distinct primes are linearly independent over $\Q$: in
the multiquadratic field $\Q(\sqrt{\mathrm p_1},\ldots,\sqrt{\mathrm p_N})$ the
sign of any one square root can be changed while the others are fixed, and
applying such a change to a linear relation and subtracting isolates a
coefficient. No nonconstructive choice is needed for finite $N$.

\begin{corollary}[No finite upper bound; source~16]\label{hol:th:cor:nobound}
For every $M\ge0$ some strongly entire differentially algebraic series over
$\KK$, for $k=\R$ and for $k=\C$, has minimal differential order greater than
$M$.
\end{corollary}

\begin{remark}[An important limitation; source~16]\label{hol:th:rem:limitation}
The interval $[N,3N]$ does not show that every integer in it occurs; $N=2$
does not produce an entire series of order two, and the argument does not
prove that the first three jets of all the seeds are jointly independent.
\end{remark}

\begin{remark}[Order two, and the known part of the spectrum; merge]
\label{hol:th:rem:ordertwo}
Source~16's Question~11.2, whether a nonpolynomial strongly entire series of
minimal order two exists, is answered negatively by
Theorem~\ref{hol:ot:main:second}\,(a), for every $\Gamma$. Hence over $\KK$
the minimal orders of nonpolynomial differentially algebraic strongly entire
series form a subset of $\{3,4,5,\ldots\}$ that contains $3$
(Theorem~\ref{hol:td:main:boundary}) and is unbounded
(Corollary~\ref{hol:th:cor:nobound}); the compressed series of
Theorem~\ref{hol:th:main:hierarchy}\,(b) have order at least $\max(N,3)$, and
by Theorem~\ref{hol:nr:main:third} every equation of order three that such a
series satisfies has total jet degree at least four. The dimension $d_N$ of
Theorem~\ref{hol:th:thm:exactdimension} is at least $\max(N,3)$. Which orders
above three occur stays open (Question~\ref{hol:th:q:spectrum}).
\end{remark}

\begin{theorem}[Exact order at the compositum dimension; source~16]
\label{hol:th:thm:exactdimension}
For scales with $\Q$-independent reciprocals let
$\mathfrak L_N=\KK(z)\langle\mathcal T_{t^{\mu_1}},\ldots,
\mathcal T_{t^{\mu_N}}\rangle$ and $d_N=\trdeg_{\KK(z)}\mathfrak L_N$. Then
$N\le d_N\le3N$, and there is $u_N\in\mathfrak L_N\cap\EK$ with
$\operatorname{dord}(u_N)=d_N$: a bounded polynomial combination, with weights
of degree below $d_N$ and coefficients in $\{0,\ldots,d_N\}$, of a
transcendence basis chosen from $\{\mathcal T_{t^{\mu_i}},D_z\mathcal
T_{t^{\mu_i}},D_z^2\mathcal T_{t^{\mu_i}}\}_{i\le N}$. Moreover $\mathfrak L_N$ is
a finite algebraic extension of $\KK(z)\langle u_N\rangle$.
\end{theorem}

\begin{proof}
Each third derivative is algebraic over the field of the displayed finite list
(Theorem~\ref{hol:td:thm:ode}), so by Lemma~\ref{hol:th:lem:ordertrdeg}
$\mathfrak L_N$ is finite over that field and a maximal independent subfamily
of the list is a transcendence basis, of size $d_N$. Compress it
(Theorem~\ref{hol:th:thm:compression}) inside $\mathfrak L_N$: the result has
$d_N$ independent jets and, by Corollary~\ref{hol:th:cor:composita}, order at
most $d_N$. The two fields have the same finite transcendence degree, and
$\mathfrak L_N$ is finitely generated over the smaller one.
\end{proof}

This realizes an intrinsically defined dimension as an exact order without
computing $d_N$, and does not identify $u_N$ as a differential primitive
element.

\section{Continuum fields, integer sampling, a third route to order three,
and scope}\label{hol:th:sec:scope}

\subsection{Continuum many independent differentially algebraic series}
\label{hol:th:sub:continuum}

\begin{lemma}[A positive Hamel set; source~16]\label{hol:th:lem:hamel}
There is $\mathfrak V\subset\R_{>0}$ of cardinality $2^{\aleph_0}$ linearly
independent over $\Q$.
\end{lemma}

\begin{proof}
Take a Hamel basis of $\R$ over $\Q$ and replace negative elements by their
negatives; an infinite-dimensional space over a countable field has
cardinality equal to its dimension, so the basis has size $2^{\aleph_0}$.
\end{proof}

\begin{corollary}[The independent family; source~16]
\label{hol:th:cor:continuum}
$\{\mathcal T_{t^{1/b}}:b\in\mathfrak V\}$ is algebraically independent over
$\KK(z)$, consists of nonpolynomial strongly entire series of minimal order
three, and has cardinality $2^{\aleph_0}$.
\end{corollary}

\begin{proof}
Every finite subfamily satisfies Theorem~\ref{hol:th:main:hierarchy}\,(a).
\end{proof}

The finite examples use explicit prime scales; the uncountable family uses
the axiom of choice through a Hamel basis, and no continuum hypothesis or
large-cardinal axiom.

Let $\mathcal A_{\mathrm{DA}}\subseteq\EK$ be the set of strongly entire
series that are differentially algebraic over $\KK(z)$. By
Proposition~\ref{hol:prop:closure} and Corollary~\ref{hol:th:cor:composita}
it is a differential subring containing $\KK[z]$, and its fraction field is
the differential field $\mathcal D_{\mathrm{DA}}\subseteq\FE$ of
Theorem~\ref{hol:th:main:hierarchy}\,(c).

\begin{lemma}[Cardinality; source~16]\label{hol:th:lem:card}
$\abs{\KK}=\abs{\EK}=2^{\aleph_0}$.
\end{lemma}

\begin{proof}
A well-ordered subset of $\R$ is countable (choose a rational between each
element and its successor), so there are at most $(2^{\aleph_0})^{\aleph_0}$
supports and coefficient assignments; a power series is a countable sequence
of elements of $\KK$. Constants give the lower bounds.
\end{proof}

\begin{proof}[Proof of Theorem~\ref{hol:th:main:hierarchy}\,\textup{(c)}]
Corollary~\ref{hol:th:cor:continuum} gives the lower bound $2^{\aleph_0}$ and
Lemma~\ref{hol:th:lem:card} the upper bound for the transcendence degree.
Each element of $\mathcal D_{\mathrm{DA}}$ lies in a differential field
generated by finitely many elements of $\mathcal A_{\mathrm{DA}}$, so is
differentially algebraic (Corollary~\ref{hol:th:cor:composita}); the
differential transcendence degree is zero. Finitely many differential
generators would give finite transcendence degree.
\end{proof}

Algebraic independence concerns relations without derivatives; differential
algebraicity allows relations involving derivatives of the same element.
The family is as large as possible in the first sense and consists of
constrained elements in the second; this difference is what compression uses.

\begin{remark}[Three fields of entire series; merge]\label{hol:th:rem:fields}
Theorem~\ref{hol:cf:thm:continuum} gives
$\dtrdeg\FE=2^{\aleph_0}$ over $\C((t^{\Gamma_\infty}))$, a group without an
order unit, where every nonpolynomial entire series is differentially
transcendental. Over $\KK=k((t^\R))$, which has an order unit,
Theorem~\ref{hol:th:main:hierarchy}\,(c) shows that the differentially
algebraic part $\mathcal D_{\mathrm{DA}}$ of $\FE$ is already of transcendence
degree $2^{\aleph_0}$, and source~17's family (Theorem~\ref{hol:fh:main:single}
with $p=t$) shows that $\FE$ itself has $\dtrdeg_{\KK(z)}\FE=2^{\aleph_0}$ there
too: its lower bound applies because $1$ is an order unit of $\R$, and
Lemma~\ref{hol:th:lem:card} gives the upper bound. Source~17 states this for
$k((t^\Q))$ (Corollary~\ref{hol:fh:cor:entirecard}); the case of $k((t^\R))$ is
the merge's.
\end{remark}

\subsection{Sampling on ordinary integers}\label{hol:th:sub:sampling}

\begin{theorem}[Residue-polynomial sampling; source~16]
\label{hol:th:thm:sampling}
Let $0\ne g=\sum_na_nz^n\in\EK$ and let $\phi_0\in k[z]$ be its initial
polynomial at the scale $0$, $\phi_0=\sum_n\res(t^{-\gamma_0(g)}a_n)z^n$
\eqref{hol:nl:eq:gaussdata}. Then $\phi_0$ is a nonzero polynomial, the zeros
of $g$ in $k\subseteq\KK$ are zeros of $\phi_0$, and in particular
$\#\{m\in\Z:g(m)=0\}\le\deg\phi_0$.
\end{theorem}

\begin{proof}
By entireness at $\delta=0$ only finitely many coefficients have the least
value $\gamma_0(g)$, and at least one does. For $c\in k$, $v(c)\ge0$, and the
residue of the normalized value $t^{-\gamma_0(g)}g(c)$ is $\phi_0(c)$: terms of
positive valuation contribute nothing, and finitely many of valuation zero
remain. So $g(c)=0$ forces $\phi_0(c)=0$.
\end{proof}

This is a quantitative form of Proposition~\ref{hol:prop:faithful}, which
says that a strongly entire series vanishing at infinitely many ordinary
complex arguments is zero; here the number of such zeros is bounded by the
degree of one initial polynomial.

\begin{corollary}[No fixed relation on infinitely many integers; source~16]
\label{hol:th:cor:samplingrelations}
For scales with $\Q$-independent reciprocals and a nonzero
$P\in\KK[z,Y_1,\ldots,Y_N]$, $P(m,\mathcal T_{t^{\mu_1}}(m),\ldots,
\mathcal T_{t^{\mu_N}}(m))\ne0$ for all but finitely many integers $m$; for
finitely many polynomials the exceptional set is common.
\end{corollary}

\begin{proof}
$P(z,\mathcal T_{t^{\mu_1}},\ldots)$ is a nonzero element of $\EK$
(Theorem~\ref{hol:th:main:hierarchy}\,(a), Proposition~\ref{hol:prop:closure});
apply Theorem~\ref{hol:th:thm:sampling}.
\end{proof}

\begin{corollary}[Jet-sampling rigidity; source~16]
\label{hol:th:cor:jetsampling}
For a series $g$ of Theorem~\ref{hol:th:main:hierarchy}\,(b) and every nonzero
$P\in\KK[z,Y_0,\ldots,Y_{N-1}]$,
$P(m,g(m),D_zg(m),\ldots,D_z^{N-1}g(m))\ne0$ for all but finitely many integers
$m$.
\end{corollary}

These statements concern a \emph{fixed} polynomial evaluated at varying
integers. They do not say that the values at one point are algebraically
independent over $\KK$: every value lies in $\KK$, and $Y-\mathcal T_{t^\mu}(m)$
vanishes at the point $m$. Ordinary integers are omnific integers, so the
corollaries give a concrete omnific test set against functional identities;
they are not a classification of Diophantine equations over $\Oz$ and use no
accumulation of integers in a surreal topology.

\subsection{A third route to minimal order three}\label{hol:th:sub:third}

Source~16 reproves, for $\Gamma=\R$ and $k=\C$, the lower bound of
Corollary~\ref{hol:td:cor:order}, by a valuation argument. Let $p=t^\mu$ with
real $\mu>0$, let $\mathcal U_\mu$ and $\mathsf S$ be as in
Lemma~\ref{hol:td:lem:laurent} (source~16's Laurent-entire ring is exactly
$\mathcal U_\mu$ for $\Gamma=\R$), and let $\gamma^{\mathcal U}$ be the Gauss
valuation of radius one, $\gamma^{\mathcal U}(\sum b_mu^m)=\min_mv(b_m)$; it is
multiplicative on $\mathcal U_\mu$, since the normalized reduction of a product
is the product of two nonzero Laurent polynomials, and extends to
$\operatorname{Frac}\mathcal U_\mu$, with residue characteristic zero and
$\gamma^{\mathcal U}(u)=0$. With $\Psi=\Psi_p$, $\psi_1=\vartheta_u\Psi/\Psi$ and
$\mathsf x=-\vartheta_u\psi_1-2\mathcal Q_1(p^2)$, \eqref{hol:td:eq:shift} gives
\[
 \mathsf S^m\Psi=p^{-m^2}u^{-m}\Psi,\qquad \mathsf S\psi_1=\psi_1-1,\qquad
 \mathsf S\mathsf x=\mathsf x ,
\]
and $\gamma^{\mathcal U}(\Psi)=0$, $\mathsf x=-p(u+u^{-1})+{}$terms of Gauss
value at least $2\mu$, so $\vartheta_u\mathsf x=-p(u-u^{-1})+\cdots\ne0$ and
$\mathsf x$ is transcendental over $\KK$ (an element algebraic over the
constants $\KK$ of $\vartheta_u$ has zero derivative). For a nonzero
$R\in\operatorname{Frac}\mathcal U_\mu[T]$, $\gamma^{\mathcal U}(R(m))$ is
constant for all but finitely many $m\in\N$, by the proof of
Lemma~\ref{hol:lem:integer} (source~16's Lemma~A.1).

\begin{theorem}[Three theta coordinates; source~16]
\label{hol:th:thm:seedthree}
$\Psi$, $\psi_1$ and $\mathsf x$ are algebraically independent over $\KK(u)$.
\end{theorem}

\begin{proof}
A relation can be written $\sum_{i,j}u^j\Psi^iP_{ij}(\psi_1,\mathsf x)=0$ with
finitely many distinct pairs $(i,j)$, $i\ge0$, and nonzero
$P_{ij}\in\KK[U,V]$. For each, $P_{ij}(\psi_1-T,\mathsf x)$ is a nonzero
polynomial in $T$: its leading coefficient is, up to sign, the leading
coefficient in $U$ of $P_{ij}(U,\mathsf x)$, a nonzero polynomial in the
transcendental $\mathsf x$. Apply $\mathsf S^m$: the $(i,j)$ summand becomes
$p^{-im^2+2jm}u^{j-im}\Psi^iP_{ij}(\psi_1-m,\mathsf x)$, of Gauss value
$-i\mu m^2+2j\mu m+\beta_{ij}$ for all large $m$, with $\beta_{ij}$ fixed. The
pair with the largest $i$ and, among those, the smallest $j$ has the uniquely
least value for all large $m$, so the transformed sum is nonzero, a
contradiction.
\end{proof}

\begin{corollary}[Minimal order three; source~16, credited conclusion]
\label{hol:th:cor:seedminimal}
$\mathcal T_{t^\mu}$, $D_z\mathcal T_{t^\mu}$ and $D_z^2\mathcal T_{t^\mu}$ are
algebraically independent over $\KK(z)$, and
$\operatorname{dord}(\mathcal T_{t^\mu})=3$.
\end{corollary}

\begin{proof}
Substitute $z=u+u^{-1}$, $\check z=u-u^{-1}$
(Lemma~\ref{hol:td:lem:substitution}). With $f=\mathcal T_{t^\mu}$ one has
$\Psi=f^\natural$, $\psi_1=\check z(D_zf)^\natural/f^\natural$ and
$\mathsf x=-\check z^2\bigl(D_z^2f/f-(D_zf/f)^2\bigr)^\natural-(u+u^{-1})
(D_zf/f)^\natural-2\mathcal Q_1(p^2)$, by \eqref{hol:td:eq:jets},, an invertible rational change of the three coordinates
over $\KK(u)$; Theorem~\ref{hol:th:thm:seedthree} excludes every relation, and
Theorem~\ref{hol:td:thm:ode} supplies one at the next jet.
\end{proof}

The argument separates finitely many terms by the growth of iterates of a
difference automorphism; it concerns functions, not values of theta
functions.

\begin{remark}[Three routes to minimal order three; merge]
\label{hol:th:rem:thirdroute}
The lower bound of Corollary~\ref{hol:td:cor:order} now has three proofs.
Source~13's (Theorem~\ref{hol:td:thm:independent}) uses the same shift
$\mathsf S$ and the same three coordinates, and three field-theoretic
obstructions on rational functions: fixed elements are constants, an additive
difference equation has no rational solution, and $\mathsf S$ preserves
$u$-adic order (Lemma~\ref{hol:td:lem:rational}), with the global certificate
$\psi_3\ne0$ of Lemma~\ref{hol:td:lem:psi3}. Source~16's replaces those three
steps by one valuation argument: iterating $\mathsf S$ makes the Gauss values
of the summands quadratic in the number of iterations, and a unique least one
survives; the needed nonvanishing is $\vartheta_u\mathsf x\ne0$, read off from
the leading term. Source~14's (Remark~\ref{hol:ot:rem:secondroute}) uses no
property of $\mathcal T_p$ beyond nonpolynomial strong entireness. Source~16's
route is stated for $\Gamma=\R$ and $k=\C$ only, as delivered; the other two
hold for every $\Gamma$ with an order unit and every $k$ of characteristic
zero. All three are kept.
\end{remark}

\paragraph{The surreal boundary and the derivation (source~16).}
Through $t^a\mapsto\omega^{-a}$ the field $\R((t^\R))$ is the surreal subfield
with real normal-form exponents, and $p=t^\mu$ becomes $\omega^{-\mu}$; every
real $\mu>0$ is cofinal in $\R$, which is why quadratic coefficient values
make these series entire on that field. At $\omega^\omega$ the evaluation
family of $\mathcal T_{t^\mu}$ is not strongly summable
(Section~\ref{hol:td:sub:surreal}), so nothing here is entire on $\No$ or
$\No[i]$. The derivative $D_z$ fixes $\KK$; an intrinsic surreal derivation
such as that of Berarducci and Mantova~\cite{bm} need not fix $p$, and
replacing $D_z$ by it changes the equations and invalidates the order
calculations. Values at omnific inputs need not be omnific; the only omnific
consequence claimed is the sampling statement above.

\paragraph{What the bounded grid makes effective (source~16).}
For given $N$ the scales $1/\sqrt{\mathrm p_i}$, the coefficient range and
the degree bound are explicit, so all $(N+1)^{N^2}$ candidates can be
listed. The proof guarantees that one has $N$ independent jets but supplies
no way to recognize it from finite truncations: $\mathsf W$ lives in a module of
K\"ahler differentials of a function field without an effective presentation
here, and Proposition~\ref{hol:th:prop:threshold} handles polynomial
expressions in the seeds, not mixed jet expressions.

\paragraph{A formalization plan proposed by source~16.}
The most reusable lemma is Theorem~\ref{hol:th:thm:compression}, with no
surreal input: the K\"ahler criterion, the invertible coefficient-to-jet
matrix, the unique squarefree coefficient and grid nonvanishing. The analytic
part needs the entire-series ring with its Gauss valuation, the coefficient
values of the seed, a finite dominant-term lemma and the separant argument;
the theta equation can first be a named imported theorem. No Lean code or
certificate accompanies source~16.

\subsection{Predecessors, repository comparison and priority}
\label{hol:th:sub:predecessors}

\paragraph{Literature.}
Source~16 consulted DLMF~\S20.5 and \S23.6~\cite{dlmftheta,dlmfelliptic}; the
Tate equation, coefficients and invariant differential on printed page~32 of
the course notes hosted by M.~Woodbury~\cite{woodbury}, inspected as an image
(it attributes the notes to a course of J.~Ellenberg; the hosting attribution
does not identify the author of every part); Pogudin's primitive-element
papers~\cite{pogudin2014,pogudin2018} for context; and the Mantova--Matusinski
survey~\cite{mantovamatusinski} for normal forms and intrinsic derivations. It
consulted the Wikipedia article on surreal numbers for orientation only. No
exhaustive MathSciNet or zbMATH search, citation-network review or referee
report was made, and irrelevant keyword searches were not taken as evidence
of novelty. A close variant of the compression argument, or a classical
antecedent of the hierarchy in differential algebra or non-Archimedean
function theory, may exist. For this merge the Pogudin preprints'
bibliographic data were checked; their content was not.

\paragraph{Repository comparison at the snapshot.}
Source~16 read, through a connector, the root and \texttt{docs/} inventories,
this report's directory and README and source~13's two audit files, at a
recorded snapshot \texttt{0fffc26} taken while the repository was changing;
this is not an assertion that every read was pinned to that commit, nor an
audit of all reports, archives, branches or commits. A search returning no
match was not treated as evidence of absence. Its statements were accurate
at the snapshot, as recorded in Section~\ref{hol:th:sub:conventions}; its
first research direction, theta descent as a counterexample to order-unit
rigidity, was dropped when source~13's audit showed it was already done.

\paragraph{Priority.}
The seed, its coefficient values, its equation and its minimal order three are
source~13's; the theta product and the Tate identity are classical and
imported. The proposed contributions of source~16 are independence across
reciprocal-independent scales with the exact profile and the finite-radius
certificate, bounded polynomial jet compression, unbounded minimal order with
the bounds $[N,3N]$ and the exact-dimension refinement, the transcendence
degrees of $\mathcal D_{\mathrm{DA}}$, integer-sampling rigidity, and the
valuation route to minimal order three; priority is provisional.

\subsection{What the theta-hierarchy part does not claim}
\label{hol:th:sub:nonclaims}

Every limitation stated by source~16, in its manuscript, README, audit files or
verification record, is kept.
\begin{enumerate}[label=\textup{(Th\arabic*)}]
\item Not every candidate of Theorem~\ref{hol:th:main:hierarchy}\,(b) is
claimed to work, and no algorithm recognizes a working candidate.
\item The bound $3N$ is not shown sharp, $d_N$ is not computed, and the full
order spectrum is not determined; not every integer in $[N,3N]$ is claimed to
occur.
\item Order two is not decided by the case $N=2$ (it is decided, negatively,
by Theorem~\ref{hol:ot:main:second}, not by source~16).
\item Joint independence of all $3N$ seed jets is not proved.
\item The finite checks do not decide differential dependence, prove
summability of infinite supports, algebraic independence of the full
functions, the compression theorem, minimal orders or cardinalities.
\item Nothing is entire on the full class $\No$ or $\No[i]$.
\item No point values are claimed algebraically independent over their own
coefficient field.
\item The seed, its equation and its minimal order three are credited as prior
(source~13); the theta product and the Tate identity are imported classical
inputs.
\item Priority is provisional: a primitive-element-type antecedent of the
compression lemma, or of the hierarchy, may exist.
\item No independent referee report and no Lean formalization.
\item The repository comparison was targeted, at a snapshot recorded while the
repository was changing.
\item The continuum family uses choice (a Hamel basis), not the continuum
hypothesis or a large-cardinal axiom.
\item The asymptotics are real-valued: the part is proved for $\Gamma=\R$ only,
and the third route to order three for $\Gamma=\R$, $k=\C$.
\item No named historical conjecture is claimed solved, and no complete
classification.
\item The derivative is the external $D_z$; nothing is claimed for an intrinsic
surreal derivation.
\item \emph{[merge]} The identification of the two seed equations,
Remarks~\ref{hol:th:rem:ordertwo}, \ref{hol:th:rem:fields}
and~\ref{hol:th:rem:thirdroute} are the merge's; no historical priority is
claimed for them.
\end{enumerate}

%======================================================================
% Single-scale freedom part: source 17. Labels use hol:fh:.
%======================================================================
\section{Differential freedom at a single scale: exact relation ideals}
\label{hol:fh:sec:setting}

This section and the next contain the material of source~17, delivered in
batch~31 with sources~15, 16, 18 and~19. The coefficient-field part explains
differential transcendence without an order unit by coefficients whose values
reach every Archimedean level; Theorem~\ref{hol:cf:main:rank} does it for a
single series by infinite coefficient-value rank. Source~17 supplies the
opposite mechanism, \emph{support}: the lacunary series
$\mathcal G_A(p,z)=\sum_{n\in A}p^{\mathsf b_n}z^n$, with heights
$\mathsf b_n$ whose successive ratios tend to infinity, have all Taylor
coefficients in $\Q(p)$ and coefficient values spanning $\Q\mu$ when
$p=t^\mu$, yet every infinite subseries is strongly entire, nonpolynomial and
differentially transcendental over the full field $\KK(z)$ whenever $\mu$ is an
order unit (Theorem~\ref{hol:fh:main:single}). More strongly, the whole
differential relation ideal of every finite family is determined: after
finite polynomial corrections it is generated by the linear relations among
the membership patterns that occur infinitely often. The same finite core
gives numerical surreals with independent Berarducci--Mantova jets, an exact
strong-evaluation domain in $\No[i]$, omnific codes and floor profiles.

\subsection{Source, conventions and the merge's choices}
\label{hol:fh:sub:conventions}

\begin{convention}[Setting of the single-scale part]
\label{hol:fh:conv:field}
In Sections~\ref{hol:fh:sec:setting}--\ref{hol:fh:sec:realizations}, $k$ is a
field of characteristic zero and $n\ge2$ throughout. In the formal sections
$p$ and $z$ are algebraically independent indeterminates, with the commuting
derivations $\vartheta_p=p\,\mathrm d/\mathrm dp$ and $\vartheta=z\,D_z$ on
$k(p,z)$, $\vartheta_pz=0$ and $\vartheta p=0$; over a larger coefficient field
$E(z)$ the derivation $D_z$ kills all of $E$. In the Hahn sections $\Gamma$ has
an order unit $\mu$, $\KK=k((t^\Gamma))$, $p=t^\mu$, and $v$, $\EK$, $\FE$ and
strong summability are as in Conventions~\ref{hol:nl:conv:field}
and~\ref{hol:cf:conv:field}. A family of derivatives is \emph{differentially
algebraically independent} if every finite list of distinct iterated
derivatives is algebraically independent over the stated base.
\end{convention}

\paragraph{The source.}
Source~17 is \emph{Differential Freedom at a Single Surreal Scale:
Factorial-height entire functions, exact relation ideals, and omnific codes}
(24 pages, A4, 23 September 2026), delivered as the archive
\texttt{single\_scale\_surreal\_research} (batch~31, manuscript~03). It
consulted the repository at
\begin{center}\small\texttt{bcac55ae6fd3f2b354e568ba1d2e94d496a2cc59},\end{center}
the pin of source~18, at which this report was the five-source text
(sources~08--12). It read this report's directory, catalogue entry and the
source~12 audit, and the READMEs of the tail-span~\cite{repotail} and
bounded-support~\cite{repobst} reports. Its description of this report was
accurate there. Its statement that ``the general question whether every
differentially algebraic strongly entire function in an order-unit Hahn field
is a polynomial remains open'' and its Question~12.1 are stale now: they are
answered negatively by Theorem~\ref{hol:td:main:boundary}, with least order
three by Theorem~\ref{hol:ot:main:second} (Remark~\ref{hol:fh:rem:universal}).
Source~17 does not claim that question; its results concern explicit families.

\paragraph{Renamed symbols.}
The following symbols of source~17 would collide with symbols of the earlier
parts. They are renamed here, in the new material only.
{\small
\begin{longtable}{@{}>{\raggedright\arraybackslash}p{0.27\textwidth}
>{\raggedright\arraybackslash}p{0.2\textwidth}
>{\raggedright\arraybackslash}p{0.45\textwidth}@{}}
\toprule
Source~17 & Here & Reason\\
\midrule
\endhead
formal variable $q$; $q=t^\eta$; order unit $\eta$ & $p$; $p=t^\mu$; $\mu$ &
$q$ is a dilation parameter; $p=t^\mu$ as in
Section~\ref{hol:td:sub:conventions}; $\eta=v(\tau)$ in
Theorem~\ref{hol:thm:unitorbit}.\\
$\delta=q\,\mathrm d/\mathrm dq$; $\Theta=z\,\mathrm d/\mathrm dz$;
$D=\mathrm d/\mathrm dz$ & $\vartheta_p$; $\vartheta$; $D_z$ & $\delta$ is a
scale; $\Theta_p$ is the partial theta series.\\
$S_A(q)$; $F_A(q,z)$; heights $b_n$, $\boldsymbol b$ & $\mathsf N_A(p)$;
$\mathcal G_A(p,z)$; $\mathsf b_n$ & $\Sfam$ is the coefficient-family
algebra, $F$ is \eqref{hol:eq:explicitF}, $\mathcal F_\alpha$ the branch
family; $b$ is a derivative order.\\
$K=k((t^\Gamma))$, $\mathcal E(K)$ (cofinally entire) & $\KK$, $\EK$ & The same
class (Proposition~\ref{hol:cf:prop:criterion}).\\
patterns $v_n$, $T_v$, $\mathcal P_\infty$, $E_0$; $p(q,z)$, $s(q)$; $V$, $W$,
$d$; $\lambda_j$; $L$ & $\mathbf v_n$, $A_{\mathbf v}$, $\mathcal P_\infty$,
$A_{\mathrm{fin}}$; $\mathsf c(p,z)$, $\mathsf c_0(p)$; $\mathsf V$, its column
span, $d_{\mathrm{pat}}$; $\boldsymbol\ell_j$; $\mathsf L$ & $p=t^\mu$;
$V$, $W$ are indeterminates; $\lambda$ is a multiplier; $L$ an operator.\\
$D$, $H$, $T$, $U$, $k$, $B_T$ (Lemma~3.1) & $\bar d$, $\bar h$, $\mathbf T$,
$\mathbf U$, $n_*$, $\mathsf b_{\mathbf T}$ & $D_z$ is the derivative, $H\in\Gamma$,
$T$ an indeterminate, $k$ the coefficient field.\\
$L$, $Q$, $J$, $U_{i,h}$, $X_{i,h}$, $W_Q$ (polarization) & $\mathfrak K$,
$\mathfrak P$, $\mathcal J$, $\varepsilon_{i,h}$, $\mathbf x_{i,h}$,
$\mathsf W_{\mathfrak P}$ & $Q_s$, $\mathcal Q_r$ and $W$ have other
meanings.\\
$\mathcal D$ (domain in $\No[i]$) & $\mathcal O_{\langle1\rangle}$ & The valuation
ring of the coarsening by the convex hull of $\Z$, as
$\mathcal O_{\Delta_\mu}$ in Theorem~\ref{hol:td:thm:domain}.\\
$\mathcal B_\Gamma(k)$, $\mathcal F_\Gamma(k)$; $H=\R((t^\R))$ & $\mathcal
B^{\mathrm{bd}}_\Gamma(k)$, $\mathcal F^{\mathrm{bd}}_\Gamma(k)$;
$\R((t^\R))$ & $\mathcal B[f]$ and $\mathcal F_\alpha$ have other
meanings.\\
$\xi_A=\sum\omega^{-b_n}$; $\partial_{\mathrm{BM}}$ & $\mathfrak x_A$;
$\partial_{\mathrm{BM}}$ & $\xi_{ij}$ are leading coefficients; the tail-span
report's $\xi_A$ are different numbers.\\
$\Lambda$, $M=\omega^\Lambda$; $Z_A$; $I_A(a)$ & $\Omega$,
$M_\Omega=\omega^\Omega$; $\mathfrak z_A$; $\mathsf I_A(a)$ & $\Lambda$ is a set of
dilations, $I_P$ the full corner polynomial.\\
nodes $a_i$, multiplicities $m_i$, values $c_{i,j}$, $M$, $H_\alpha$, $P$,
$G_\alpha$ (Hermite) & $x_i$, $e_i$, $\mathsf y_{ij}$, $\bar e$,
$\mathsf R_\alpha$, $\mathsf R_0$, $\mathcal G^*_\alpha$ & $a_n$ are Taylor
coefficients.\\
polynomial height $h(n)$ & $\mathsf h(n)$ & $h$ is a root-of-unity order.\\
$\mathfrak c=2^{\aleph_0}$ & $2^{\aleph_0}$ & $\mathfrak c$ is a theta constant.\\
Theorems~4.2, 4.3, 5.1, 7.2 & Theorems~\ref{hol:fh:thm:numeric},
\ref{hol:fh:thm:mixed}, \ref{hol:fh:thm:relations}, \ref{hol:fh:thm:entire} &
Theorem~\ref{hol:fh:main:single} collects them; Theorems~A--N of this report
are different theorems.\\
\bottomrule
\end{longtable}}
The shipped files \path{17-single-scale-PROOF_AUDIT.md} and
\path{17-single-scale-SOURCES_AND_SCOPE.md} keep the source's notation.

\paragraph{Words with two meanings.}
A \emph{rapid height sequence} is Definition~\ref{hol:fh:def:rapid}; it has
nothing to do with the scale $\mu$. A \emph{pattern} is the membership vector
$\mathbf v_n$ of an index. \emph{Cofinally entire} (source~17) is exactly
$\EK$. \emph{Differential independence} of the numerical series $\mathsf N_A$
refers to $\vartheta_p$, or to $\partial_{\mathrm{BM}}$ after the substitution
$p=\omega^{-1}$, never to $D_z$; that of $\mathcal G_A$ refers to $D_z$ (pure)
or to $\vartheta,\vartheta_p$ (mixed), over the base stated each time. The
\emph{support} mechanism of this part is not the \emph{coefficient} mechanism
of Sections~\ref{hol:cf:sec:setting}--\ref{hol:cf:sec:families} or of the
tail-span report.

\paragraph{Results printed once.}
Source~17's definition of the entire class (Definition~2.2) is $\EK$ in the
form of Proposition~\ref{hol:cf:prop:criterion}; its Lemma~7.1 (the entire
ring) is Proposition~\ref{hol:prop:closure}; its Euler--Stirling conversion
(2.1) is \eqref{hol:ot:eq:stirling}; its Hahn--Neumann support facts are
Lemma~\ref{hol:lem:support}; its constant-extension step (Lemma~6.1 and
Proposition~6.2) is the argument of Proposition~\ref{hol:cf:prop:scalar},
over $\Q(p)$ instead of $\KK$, applied to the family of all jets, source~17's proof via an invertible finite coefficient
minor being a second route to the linear case; its normal-form embedding is
\eqref{hol:eq:embedding}; its bounded-support descent (Lemma~8.3) is the
integer-lattice case of \texttt{bst:thm:descent}~\cite{repobst}, reproved
below for completeness; its binary-prefix family is the almost-disjoint device
already used for $\mathcal F_\alpha$ (Section~\ref{hol:cf:sub:continuum}); and
its floor profile (Theorem~10.4) is a case of \texttt{odg:thm:floor}~\cite{repoodg},
printed with its short proof.

\paragraph{Where each numbered result went.}
Definition~2.1 is Definition~\ref{hol:fh:def:rapid}; Lemma~3.1,
Corollary~3.2, Lemma~3.3, Example~3.4 and Lemma~3.5 are
Lemma~\ref{hol:fh:lem:sep}, Corollary~\ref{hol:fh:cor:bands},
Lemma~\ref{hol:fh:lem:polar}, Example~\ref{hol:fh:ex:polar} and
Lemma~\ref{hol:fh:lem:density}; Definition~4.1, Theorems~4.2, 4.3,
Corollary~4.4 and Remark~4.5 are Definition~\ref{hol:fh:def:private},
Theorems~\ref{hol:fh:thm:numeric}, \ref{hol:fh:thm:mixed},
Corollary~\ref{hol:fh:cor:every} and Remark~\ref{hol:fh:rem:nonlinear};
Theorem~5.1, Corollary~5.2, Examples~5.3, 5.4 and Corollary~5.5 are
Theorem~\ref{hol:fh:thm:relations}, Corollary~\ref{hol:fh:cor:jetdegree},
Examples~\ref{hol:fh:ex:union}, \ref{hol:fh:ex:finite} and
Corollary~\ref{hol:fh:cor:alphabet}; Theorem~6.3 and Warning~6.4 are
Theorem~\ref{hol:fh:thm:allscalars} and Remark~\ref{hol:fh:rem:constext};
Theorem~7.2, Corollary~7.3 and Theorem~7.4 are
Theorem~\ref{hol:fh:thm:entire}, Corollary~\ref{hol:fh:cor:entirecard} and
Theorem~\ref{hol:fh:thm:hermite}; Lemma~8.1, Theorem~8.2, Lemma~8.3,
Corollary~8.4 and Warning~8.5 are Lemma~\ref{hol:fh:lem:bmtriangular},
Theorem~\ref{hol:fh:thm:surreal}, Lemma~\ref{hol:fh:lem:bounded},
Corollary~\ref{hol:fh:cor:boundedjets} and Remark~\ref{hol:fh:rem:absolute};
Theorem~9.1, Propositions~9.2, 9.3 and Warning~9.4 are
Theorem~\ref{hol:fh:thm:domain}, Propositions~\ref{hol:fh:prop:domainring},
\ref{hol:fh:prop:faithful} and Remark~\ref{hol:fh:rem:three}; Lemma~10.1,
Theorem~10.2, Warning~10.3 and Theorem~10.4 are Lemma~\ref{hol:fh:lem:scale},
Theorem~\ref{hol:fh:thm:codes}, Remark~\ref{hol:fh:rem:codes} and
Theorem~\ref{hol:fh:thm:floor}; Proposition~11.1 is
Proposition~\ref{hol:fh:prop:heat}; Warning~1.1 is merged into
Section~\ref{hol:fh:sub:predecessors}; Question~12.1 is answered
(Remark~\ref{hol:fh:rem:universal}) and Questions~12.2--12.11 are
Questions~\ref{hol:fh:q:heights}--\ref{hol:fh:q:arithmetic}; Appendices~A
and~B are merged into the dependency table and suite description of
Section~\ref{hol:sec:verification}, and Appendix~C into
Section~\ref{hol:fh:sub:predecessors}.

\paragraph{The merge's choices.}
The review of source~17 for this merge found its proofs correct. The merge
adds three points, each marked. First, Remark~\ref{hol:fh:rem:universal}
records the answer to source~17's universal question. Second,
Remark~\ref{hol:fh:rem:placement} places Theorem~\ref{hol:fh:thm:entire}
beside Theorems~\ref{hol:cf:main:rank}, \ref{hol:cf:main:sharp},
\ref{hol:td:main:boundary} and~\ref{hol:th:main:hierarchy}, and records that it
fills the gap noted after Theorem~\ref{hol:cf:thm:rankone}. Third,
Remark~\ref{hol:fh:rem:domains} compares the common domain of
Theorem~\ref{hol:fh:thm:domain} with that of $\mathcal T_p$. Source~17's
theorems, proofs, examples and non-claims are otherwise printed as delivered,
in the renamed notation.

\subsection{Rapid heights and the finite separation mechanism}
\label{hol:fh:sub:finite}

\begin{definition}[Rapid heights; source~17]\label{hol:fh:def:rapid}
A \emph{rapid height sequence} is a strictly increasing sequence of positive
integers $(\mathsf b_n)_{n\ge2}$ with $\mathsf b_n/\mathsf b_{n-1}\to+\infty$,
for instance $\mathsf b_n=n!$. For $A\subseteq\N_{\ge2}$ put
\[
 \mathsf N_A(p)=\sum_{n\in A}p^{\mathsf b_n}\in k[[p]],\qquad
 \mathcal G_A(p,z)=\sum_{n\in A}p^{\mathsf b_n}z^n\in k[p][[z]] .
\]
\end{definition}

Then $\vartheta_p^b\mathsf N_A=\sum_{n\in A}\mathsf b_n^bp^{\mathsf b_n}$ and
$\vartheta^a\vartheta_p^b\mathcal G_A=\sum_{n\in A}n^a\mathsf b_n^bp^{\mathsf b_n}
z^n$. By \eqref{hol:ot:eq:stirling} and its analogue in $p$, finite-jet
independence and its transcendence degree are the same for Euler and ordinary
derivatives, the changes being triangular and invertible over $k(p,z)$.

\begin{lemma}[Height separation; source~17]\label{hol:fh:lem:sep}
Fix rapid heights and integers $\bar d\ge1$, $\bar h\ge0$. For all large $N$:
if $\mathbf T$ is a multiset of exactly $\bar d$ indices, all at least $N$, and
$\mathbf U\ne\mathbf T$ a multiset of at most $\bar d$ indices, then
$\abs{\sum_{n\in\mathbf T}\mathsf b_n-\sum_{n\in\mathbf U}\mathsf b_n}>\bar h$.
For $\mathsf b_n=n!$ it suffices that $N\ge\max(3,4\bar d)$ and $N!/2>\bar h$.
\end{lemma}

\begin{proof}
If the multiplicities agree at every index at least $N$, then $\mathbf U$
already contains all $\bar d$ entries of $\mathbf T$ and has no room for more,
so $\mathbf U=\mathbf T$. Hence there is a largest differing index $n_*\ge N$.
There the multiplicity difference is at least one in absolute value, the
lower indices contribute at most $2\bar d\,\mathsf b_{n_*-1}$, and higher
indices cancel, so the difference is at least $\mathsf b_{n_*}-2\bar d\,
\mathsf b_{n_*-1}$. Choose $N$ with $\mathsf b_n>4\bar d\,\mathsf b_{n-1}$ for
$n\ge N$ and $\mathsf b_N/2>\bar h$. For $n!$ and $n\ge4\bar d$ the bound is at
least $n!/2$.
\end{proof}

That $\mathbf T$ has \emph{exactly} the maximal degree is essential: it stops a
competitor from matching every high index and adding low ones.

\begin{corollary}[Isolated bands; source~17]\label{hol:fh:cor:bands}
With $\mathsf b_{\mathbf T}=\sum_{n\in\mathbf T}\mathsf b_n$, the integer
interval $[\mathsf b_{\mathbf T},\mathsf b_{\mathbf T}+\bar h]$ is disjoint from
$[\mathsf b_{\mathbf U},\mathsf b_{\mathbf U}+\bar h]$ for every different
$\mathbf U$ of size at most $\bar d$.
\end{corollary}

\begin{lemma}[Multicolour polarization; source~17]\label{hol:fh:lem:polar}
Let $\mathfrak K$ have characteristic zero, $\mathcal J\subseteq\N^s$ be finite,
and $\mathfrak P\ne0$ be a polynomial in variables $Y_{i,\alpha}$
($1\le i\le m$, $\alpha\in\mathcal J$), multihomogeneous of degree $d_i$ in the
variables of colour $i$. For each colour take $d_i$ slots with scalars
$\varepsilon_{i,h}$ and $s$-tuples $\mathbf x_{i,h}$, and substitute
$Y_{i,\alpha}=\sum_h\varepsilon_{i,h}\mathbf x_{i,h}^\alpha$. The coefficient
$\mathsf W_{\mathfrak P}$ of $\prod_{i,h}\varepsilon_{i,h}$ is a nonzero
polynomial in the slot tuples.
\end{lemma}

\begin{proof}
A monomial of $\mathfrak P$ contributes the sum over all assignments of its
exponent vectors, colour by colour, to the labelled slots, with positive
integer multiplicities. Different monomials differ in the multiset of exponent
vectors of some colour, so their orbits of slot monomials are disjoint, and
characteristic zero keeps the multiplicities nonzero.
\end{proof}

When the indices put in the slots are distinct, $\mathsf W_{\mathfrak P}$ is
exactly the coefficient obtained by taking each once in the series product;
distinctness is enforced below, and no repeated-index factorial is divided
out.

\begin{example}[A cancellation-sensitive case; source~17]\label{hol:fh:ex:polar}
For $\mathfrak P=Y_0Y_2-Y_1^2$, $\mathsf W_{\mathfrak P}(x_1,x_2)=(x_1-x_2)^2$:
a single repeated index gives zero, two distinct indices detect it. This is
the conic $\mathfrak s$ of Remark~\ref{hol:nr:rem:conic}, whose cross
coefficient $(n-m)^2$ drives the same phenomenon there.
\end{example}

\begin{lemma}[Hereditary density; source~17]\label{hol:fh:lem:density}
For rapid heights, on every infinite set of indices a nonzero polynomial over
a field of characteristic zero vanishes at only finitely many of the points
$\mathsf b_n\in\mathbb A^1$, $n\in\mathbb A^1$, and $(n,\mathsf b_n)\in\mathbb
A^2$.
\end{lemma}

\begin{proof}
The one-dimensional cases follow from distinctness. Eventually $\mathsf
b_n\ge2\mathsf b_{n-1}$, so $\mathsf b_n$ dominates every power of $n$; for
$0\ne P=\sum_{j\le e}P_j(X)Y^j\in\Q[X,Y]$ with $P_e\ne0$ the term
$P_e(n)\mathsf b_n^e$ dominates the others in absolute value for large $n$.
For general coefficients write $P=\sum_je_jP_j$ over a $\Q$-basis $(e_j)$ of
the span of its coefficients; some $P_j\ne0$, and $P(n,\mathsf b_n)=0$ forces
$P_j(n,\mathsf b_n)=0$.
\end{proof}

Products of such sets are Zariski dense by induction on the number of tuples,
and an infinite index set splits into any finite number of disjoint infinite
sets, so distinct indices can be chosen for all slots, arbitrarily far out,
with $\mathsf W_{\mathfrak P}\ne0$; this holds over $\mathfrak K=k(p)$.

\subsection{Joint independence}\label{hol:fh:sub:independence}

\begin{definition}[Finite private tails; source~17]\label{hol:fh:def:private}
A family $\mathcal A$ of infinite subsets of $\N_{\ge2}$ has the \emph{finite
private-tail property} if for distinct $A_1,\ldots,A_m\in\mathcal A$ every
$A_i\setminus\bigcup_{j\ne i}A_j$ is infinite; pairwise almost-disjoint
families have it.
\end{definition}

\begin{theorem}[Numerical all-order independence; source~17]
\label{hol:fh:thm:numeric}
For rapid heights and a family $\mathcal A$ with finite private tails, the
elements $\vartheta_p^b\mathsf N_A$, $A\in\mathcal A$, $b\ge0$, are
algebraically independent over $k(p)$.
\end{theorem}

\begin{proof}
A relation gives, after clearing denominators, a nonzero $P\in k[p][Y_{i,j}]$
of top total degree $\bar d$ and $p$-degree $\bar h$; $\bar d=0$ is impossible
since $p$ is transcendental over $k$. Take a nonzero multihomogeneous part
$\mathfrak P$ of degree $\bar d$ with colour degrees $d_i$, take $N$ from
Lemma~\ref{hol:fh:lem:sep}, split each colour's private indices beyond $N$
into $d_i$ infinite sets, and by Lemmas~\ref{hol:fh:lem:polar}
and~\ref{hol:fh:lem:density} over $k(p)$, with slot weights $\mathsf b_n$,
choose $\bar d$ distinct private indices beyond $N$ with polarized
coefficient $C(p)\ne0$ in $k[p]$, of degree at most $\bar h$. Let $\mathbf T$
be this multiset. Every contribution to $P$ evaluated comes from an index
multiset of size at most $\bar d$ followed by a shift in $[0,\bar h]$; by
Corollary~\ref{hol:fh:cor:bands} only $\mathbf T$ reaches the band
$[\mathsf b_{\mathbf T},\mathsf b_{\mathbf T}+\bar h]$, the private indices force
the colour degrees $d_i$, and the contribution there is
$p^{\mathsf b_{\mathbf T}}C(p)\ne0$.
\end{proof}

\begin{theorem}[Mixed differential independence; source~17]
\label{hol:fh:thm:mixed}
For the same heights and family, all mixed jets
$\vartheta^a\vartheta_p^b\mathcal G_A$, $A\in\mathcal A$, $a,b\ge0$, are
algebraically independent over $k(p,z)$, and so are all ordinary mixed
partial derivatives.
\end{theorem}

\begin{proof}
From a nonzero $P\in k[p,z][Y_{i,a,b}]$ of top degree $\bar d$ and $p$-degree
$\bar h$, choose $z^{j_0}$ and a colour multidegree of total degree $\bar d$
with nonzero part $\mathfrak P\in k[p][Y_{i,a,b}]$, polarize with the
two-coordinate weights $(n,\mathsf b_n)$, and choose $\bar d$ distinct private
target indices beyond the separation threshold with polarized weight
$C(p)\ne0$. With $\mathsf b_{\mathbf T}=\sum_{\mathbf T}\mathsf b_n$ and
$M=j_0+\sum_{\mathbf T}n$, the band $[\mathsf b_{\mathbf T},\mathsf b_{\mathbf T}+
\bar h]$ of the coefficient of $z^M$ receives only the target multiset, which
forces the colours and the power $z^{j_0}$, and equals
$p^{\mathsf b_{\mathbf T}}C(p)$. Each $z$-coefficient is a polynomial in $p$, so
the argument takes place in an ordinary polynomial ring; the triangular
changes give the ordinary derivatives.
\end{proof}

\begin{corollary}[Every infinite subseries; source~17]\label{hol:fh:cor:every}
For infinite $A$, $\mathsf N_A$ is differentially transcendental over $k(p)$,
and $\mathcal G_A$ satisfies no nonzero algebraic partial differential equation
in $\vartheta,\vartheta_p$ over $k(p,z)$, in particular no algebraic ordinary
differential equation in $D_z$ over $k(p,z)$.
\end{corollary}

\begin{remark}[Why nonlinear equations are handled; source~17]
\label{hol:fh:rem:nonlinear}
The selected contribution uses the top \emph{total} degree of the whole
relation, not its highest derivative; polarization handles cancellation among
nonlinear monomials, and the band lemma excludes every lower-degree competitor
at once. No first-order, separant, linearization or residue hypothesis is
imposed. For a single ordinary series the conclusion is classical
(Section~\ref{hol:fh:sub:predecessors}).
\end{remark}

\subsection{The exact relation ideal}\label{hol:fh:sub:relations}

Fix arbitrary $A_1,\ldots,A_m\subseteq\N_{\ge2}$. For a nonzero pattern
$\mathbf v\in\{0,1\}^m$ let $A_{\mathbf v}=\{n:\mathbf v_n=\mathbf v\}$, where
$\mathbf v_n=(\mathbf 1_{A_1}(n),\ldots,\mathbf 1_{A_m}(n))$, and
$\mathcal P_\infty=\{\mathbf v\ne0:A_{\mathbf v}\text{ infinite}\}$. The union
$A_{\mathrm{fin}}$ of the finite nonzero pattern classes is finite; put
$\mathsf c(p,z)=\sum_{n\in A_{\mathrm{fin}}}\mathbf v_np^{\mathsf b_n}z^n$ and
$\mathsf c_0(p)=\sum_{n\in A_{\mathrm{fin}}}\mathbf v_np^{\mathsf b_n}$, let
$\mathsf V$ be the matrix with the columns $\mathbf v\in\mathcal P_\infty$, and
$d_{\mathrm{pat}}=\operatorname{rank}\mathsf V$, the same over every field of
characteristic zero. With $\boldsymbol{\mathcal G}=(\mathcal G_{A_i})_i$ and
$\boldsymbol{\mathsf N}=(\mathsf N_{A_i})_i$,
\begin{equation}\label{hol:fh:eq:patterns}
 \boldsymbol{\mathcal G}=\mathsf c+\mathsf V(\mathcal G_{A_{\mathbf v}})_{\mathbf
 v\in\mathcal P_\infty},\qquad
 \boldsymbol{\mathsf N}=\mathsf c_0+\mathsf V(\mathsf N_{A_{\mathbf v}})_{\mathbf
 v\in\mathcal P_\infty},
\end{equation}
where the $A_{\mathbf v}$ are pairwise disjoint and infinite, so
Theorems~\ref{hol:fh:thm:numeric} and~\ref{hol:fh:thm:mixed} apply to them.

\begin{theorem}[Complete relation classification; source~17]
\label{hol:fh:thm:relations}
Let $\boldsymbol\ell_1,\ldots,\boldsymbol\ell_{m-d_{\mathrm{pat}}}\in\Q^m$ be a
basis of the left kernel of $\mathsf V$. The kernel of evaluation of the
differential polynomial ring over $k(p)$ at $\boldsymbol{\mathsf N}$ is the
differential ideal generated by $\boldsymbol\ell_j\cdot(Y-\mathsf c_0(p))$; the
kernel over $k(p,z)$ at $\boldsymbol{\mathcal G}$, with the derivations
$\vartheta,\vartheta_p$, is the differential ideal generated by
$\boldsymbol\ell_j\cdot(Y-\mathsf c(p,z))$. There are no further nonlinear
relations; for $d_{\mathrm{pat}}=0$ all functions are the corrections, and for
$d_{\mathrm{pat}}=m$ the family is differentially algebraically independent.
\end{theorem}

\begin{proof}
Choose $\mathsf L\in\mathrm{GL}_m(\Q)$ whose last $m-d_{\mathrm{pat}}$ rows span
the left kernel. After the affine change $X=\mathsf L(Y-\mathsf c)$ the last
coordinates vanish and the first $d_{\mathrm{pat}}$ are linear forms of full
row rank in the pattern functions, acting separately on each block of their
independent jets; a full-row-rank substitution is injective on polynomial
rings (a right inverse of the matrix gives a left inverse). So the kernel is
generated by the jets of the last coordinates. The numerical case is the same
with $\mathsf c_0$ and $\vartheta_p$.
\end{proof}

\begin{corollary}[Exact finite-jet degrees; source~17]
\label{hol:fh:cor:jetdegree}
For every finite family,
\begin{align*}
 \trdeg_{k(p)}k(p)(\vartheta_p^b\mathsf N_{A_i}:i\le m,\,b\le s)
 &=(s+1)d_{\mathrm{pat}},\\
 \trdeg_{k(p,z)}k(p,z)(\vartheta^a\vartheta_p^b\mathcal G_{A_i}:i\le m,\,a\le r,\,
 b\le s)&=(r+1)(s+1)d_{\mathrm{pat}},
\end{align*}
and the same with ordinary derivatives in the corresponding rectangles.
\end{corollary}

Equivalently, $d_{\mathrm{pat}}$ is the rank of the indicator classes of the
$A_i$ in $\Q^{\N_{\ge2}}/\Q^{(\N_{\ge2})}$: a rational combination of
indicators is eventually zero exactly when it annihilates every recurring
pattern. Up to finite corrections, the differential dependence geometry of
$A\mapsto\mathcal G_A$ is ordinary linear dependence modulo finite sets.

\begin{example}[Union and intersection; source~17]\label{hol:fh:ex:union}
$\mathcal G_{A\cup B}+\mathcal G_{A\cap B}=\mathcal G_A+\mathcal G_B$, and
likewise for $\mathsf N$. If $A\setminus B$, $B\setminus A$ and $A\cap B$ are
infinite, the four series have pattern rank three, and this identity with its
derivatives generates their whole relation ideal; the family of all infinite
subsets is not independent.
\end{example}

\begin{example}[Finite changes; source~17]\label{hol:fh:ex:finite}
If $B=A\cup\{n_0\}$ with $n_0\notin A$, then $\mathcal G_B-\mathcal G_A=
p^{\mathsf b_{n_0}}z^{n_0}$: the ideal keeps this affine correction, and after it
there is exactly one differential generator when $A$ is infinite.
\end{example}

\begin{corollary}[Finite alphabets; source~17]\label{hol:fh:cor:alphabet}
The classification holds when the patterns $\mathbf v_n$ range over a fixed
finite subset of $k^m$, with the recurring vectors, the corrections and
$k$-linear rank.
\end{corollary}

Finiteness of the alphabet matters: for $\mathbf v_n=(1,n)$ the second series
is $\vartheta$ of the first, and no finite-pattern description applies.

\subsection{The full scalar field}\label{hol:fh:sub:scalars}

\begin{theorem}[All scalars may enter the equation; source~17]
\label{hol:fh:thm:allscalars}
Let $E$ have characteristic zero and contain an element $p$ transcendental
over its prime field, and regard $\mathcal G_A$ in $E[[z]]$. For rapid heights,
the pure $D_z$ relation ideal over $E(z)$ of every finite family is exactly the
ideal generated by the $\boldsymbol\ell_j\cdot(Y-\mathsf c(p,z))$ and their
$D_z$-derivatives. In particular every infinite subseries is differentially
transcendental over $E(z)$, and a finite-private-tail family is jointly so.
\end{theorem}

\begin{proof}
The pattern functions have coefficients in $\Q(p)\subseteq E$ and all their
$D_z$-jets are algebraically independent over $\Q(p)(z)$
(Theorem~\ref{hol:fh:thm:mixed} with $k=\Q$ and the triangular change); the
proof of Proposition~\ref{hol:cf:prop:scalar}, with $\Q(p)$ in place of $\KK$,
carries this to $E(z)$. Source~17 proves
this step by writing a relation as an $E$-linear relation among $\Q(p)$-linearly
independent series $z^j\prod(D_z^af_i)^{e_{i,a}}$ and choosing an invertible
finite minor of their coefficient matrix over $\Q(p)$, whose determinant stays
nonzero in $E$. The full-row-rank argument of
Theorem~\ref{hol:fh:thm:relations} then gives the ideal.
\end{proof}

\begin{remark}[What constant extension does not say; source~17]
\label{hol:fh:rem:constext}
The theorem freezes all of $E$ when differentiating in the external variable
$z$. It does not make the numbers $\mathsf N_A(p)$ independent over $E=k((p))$,
which contains them, and it equips $E$ with no derivation extending
$\vartheta_p$: the mixed theorem has base $k(p,z)$ only. A finite relation
with coefficients in a proper class, such as $\No[i]$, involves finitely many
coefficients, which with $p$ generate a set-sized field to which the theorem
applies.
\end{remark}

\section{Entire, surreal and omnific realizations of single-scale freedom}
\label{hol:fh:sec:realizations}

\subsection{Strongly entire series at one cofinal scale}
\label{hol:fh:sub:entire}

\begin{theorem}[Entire differential freedom; source~17]
\label{hol:fh:thm:entire}
Let $\Gamma$ have an order unit $\mu$, $k$ have characteristic zero,
$\KK=k((t^\Gamma))$ and $p=t^\mu$. Every $\mathcal G_A$ with rapid heights lies
in $\EK$; for each finite family its pure differential relations over $\KK(z)$
are those of Theorem~\ref{hol:fh:thm:allscalars}; in particular every infinite
subseries is a nonpolynomial strongly entire series satisfying no algebraic
ordinary differential equation over $\KK(z)$ of any order.
\end{theorem}

\begin{proof}
For $n\in A$ the coefficient value is $\mathsf b_n\mu$. For $\gamma\in\Gamma$
choose $m$ with $\gamma\ge-m\mu$; then $\mathsf b_n\mu+n\gamma\ge(\mathsf
b_n-mn)\mu$, and $\mathsf b_n/n\to\infty$ with $\mu$ an order unit makes this
cofinal, so $\mathcal G_A\in\EK$ by Proposition~\ref{hol:cf:prop:criterion}.
The monomial $p$ is transcendental over the prime field; apply
Theorem~\ref{hol:fh:thm:allscalars} with $E=\KK$.
\end{proof}

\begin{remark}[Where the theorem sits; merge]\label{hol:fh:rem:placement}
The coefficient values $\mathsf b_n\mu$ of $\mathcal G_A$ span $\Q\mu$, so its
coefficient-value rank is one and Theorem~\ref{hol:cf:main:rank} does not
apply; the remark after Theorem~\ref{hol:cf:thm:rankone} noted that that
theorem ``says nothing about a series whose coefficient values span only
$\Q\mu$'' and ``does not settle the purely differential problem over
$\C((t^\Q))$''. Theorem~\ref{hol:fh:thm:entire} supplies, over $\C((t^\Q))$
and every order-unit field, explicit nonpolynomial strongly entire series of
coefficient-value rank one that are differentially transcendental over
$\KK(z)$. Their coefficients lie in $\Q(p)$, a field of transcendence degree one
over $\Q$, so they are further witnesses for condition~(ii) of
Theorem~\ref{hol:cf:main:sharp}: a finitely generated coefficient field is
necessary for differential algebraicity
(Theorem~\ref{hol:cf:thm:coefficients}\,(a)) but, with an order unit, not
sufficient. Beside Theorem~\ref{hol:td:main:boundary} and
Theorem~\ref{hol:th:main:hierarchy}, the same fields therefore carry
nonpolynomial strongly entire series of every minimal order in a set that
contains $3$ and is unbounded (for $\Gamma=\R$), and series of no finite order.
They do not decide $\Theta_p$ or $G$, whose heights $n^2$ are not rapid
(Proposition~\ref{hol:fh:prop:heat}, Question~\ref{hol:fh:q:heights}).
\end{remark}

For a binary sequence $\alpha=(\alpha_1,\alpha_2,\ldots)$ put
$n_\ell(\alpha)=2^\ell+\sum_{j\le\ell}\alpha_j2^{\ell-j}$ and $A_\alpha=\{n_\ell(\alpha):
\ell\ge1\}$. Each $n_\ell(\alpha)$ is $1$ followed by the first $\ell$ digits of
$\alpha$ in binary, different lengths lie in disjoint dyadic intervals, and
$A_\alpha\cap A_\beta$ is finite for $\alpha\ne\beta$; so the family has finite
private tails. Each $A_\alpha$ is computable relative to $\alpha$; uniform
computability of all members is not claimed.

\begin{corollary}[Maximal differential transcendence degree; source~17]
\label{hol:fh:cor:entirecard}
For $\KK=\R((t^\Q))$ or $\C((t^\Q))$, with $p=t=\omega^{-1}$ in the surreal
realization, $\dtrdeg_{\KK(z)}\FE=2^{\aleph_0}$, with the lower bound supplied
by $(\mathcal G_{A_\alpha})_{\alpha\in2^\N}$, all of whose coefficients lie in
$\Q(t)$.
\end{corollary}

\begin{proof}
Joint differential independence gives the lower bound;
$\abs{\KK[[z]]}=(2^{\aleph_0})^{\aleph_0}$ gives the upper bound.
\end{proof}

This is the order-unit counterpart of Theorem~\ref{hol:cf:thm:continuum}; for
$\KK=k((t^\R))$ see Remark~\ref{hol:th:rem:fields}. A family of maximal
cardinality is not claimed to be a differential transcendence basis.

\begin{theorem}[Freedom inside every finite-jet fibre; source~17]
\label{hol:fh:thm:hermite}
In either field of Corollary~\ref{hol:fh:cor:entirecard}, fix distinct nodes
$x_1,\ldots,x_s\in\KK$, multiplicities $e_i\ge1$ and values
$\mathsf y_{ij}\in\KK$, $0\le j<e_i$. There are $2^{\aleph_0}$ jointly
differentially algebraically independent $\mathcal G^*_\alpha\in\EK$ with
$D_z^j\mathcal G^*_\alpha(x_i)=\mathsf y_{ij}$ for all $i,j,\alpha$.
\end{theorem}

\begin{proof}
With $\bar e=\sum e_i$, the map $\KK[z]_{<\bar e}\to\prod_i\KK^{e_i}$,
$\mathsf R\mapsto(D_z^j\mathsf R(x_i))$, is injective, a polynomial in its kernel
being divisible by $\prod(z-x_i)^{e_i}$, hence bijective. Let $\mathsf
R_\alpha$ interpolate the jets of $\mathcal G_{A_\alpha}$ and $\mathsf R_0$ the
values $\mathsf y_{ij}$, and put $\mathcal G^*_\alpha=\mathcal G_{A_\alpha}-
\mathsf R_\alpha+\mathsf R_0$. For finitely many $\alpha$ the translation is an
invertible change of variables over $\KK(z)$.
\end{proof}

\subsection{Numerical surreals and the intrinsic derivation}
\label{hol:fh:sub:actual}

The normalized Berarducci--Mantova derivation $\partial_{\mathrm{BM}}$ has real
constants, is strongly additive and satisfies $\partial_{\mathrm{BM}}\omega=1$
\cite{bm}. On the integer-exponent Laurent field with $p=\omega^{-1}$ it
therefore acts as
\begin{equation}\label{hol:fh:eq:bm}
 \partial_{\mathrm{BM}}\Bigl(\sum_{n\ge n_0}c_np^n\Bigr)=-\sum_{n\ge n_0}nc_np^{n+1}
 =-p\,\vartheta_p\Bigl(\sum_{n\ge n_0}c_np^n\Bigr),
\end{equation}
first on monomials and then by strong additivity. It is not $D_z$: it moves
$p$.

\begin{lemma}[Triangular intrinsic jets; source~17]\label{hol:fh:lem:bmtriangular}
On that field $\partial_{\mathrm{BM}}^j=(-1)^jp^j\vartheta_p(\vartheta_p+1)\cdots
(\vartheta_p+j-1)$, so intrinsic and Euler jets of order at most $r$ are related
by an invertible triangular matrix over $\R(p)$.
\end{lemma}

\begin{proof}
$\vartheta_p(p^jg)=p^j(\vartheta_p+j)g$; induct from
$\partial_{\mathrm{BM}}=-p\vartheta_p$.
\end{proof}

\begin{theorem}[Rational-coefficient surreal differential freedom;
source~17]\label{hol:fh:thm:surreal}
Let $\mathfrak x_A=\sum_{n\in A}\omega^{-\mathsf b_n}\in\R((\omega^{-1}))
\subseteq\No$. For a finite-private-tail family of infinite sets, all
$\partial_{\mathrm{BM}}^j\mathfrak x_A$ are jointly algebraically independent
over $\R(\omega)$; for arbitrary finite families the relation ideal is that of
Theorem~\ref{hol:fh:thm:relations} with the correction $\mathsf c_0(\omega^{-1})$
and $\partial_{\mathrm{BM}}$ in place of $\vartheta_p$. In particular
$\dtrdeg_{\R(\omega)}\R((\omega^{-1}))=2^{\aleph_0}$ for $\partial_{\mathrm{BM}}$,
and the analogous statement holds over $\C(\omega)$.
\end{theorem}

\begin{proof}
Theorem~\ref{hol:fh:thm:numeric} and the classification with $k=\R$, then
Lemma~\ref{hol:fh:lem:bmtriangular}; each nonzero $\mathfrak x_A$ is a positive
infinitesimal; the binary-prefix family and the cardinality of the Laurent
field give the transcendence degree.
\end{proof}

Let $\mathcal B^{\mathrm{bd}}_\Gamma(k)$ be the ring of $f\in k((t^\Gamma))$ with
support bounded above, and $\mathcal F^{\mathrm{bd}}_\Gamma(k)$ its fraction
field; it contains $k[p,p^{-1}]$ for $p=t^\mu$.

\begin{lemma}[Bounded-support descent; source~17, after
\texttt{bst:thm:descent}]\label{hol:fh:lem:bounded}
Let $\mu$ be an order unit. If $g_1,\ldots,g_N\in k((p))$ are linearly
independent over $k(p)$, they are linearly independent over
$\mathcal F^{\mathrm{bd}}_\Gamma(k)$ in $k((t^\Gamma))$.
\end{lemma}

\begin{proof}
From $\sum_jb_jg_j=0$ with $b_j\in\mathcal B^{\mathrm{bd}}_\Gamma(k)$, project to
a coset $\gamma+\Z\mu$: the support of $b_j$ there is bounded below and above,
hence finite because $\mu$ is an order unit, so the projection is
$t^\gamma P_{j,\gamma}(p)$ with a Laurent polynomial $P_{j,\gamma}$.
Multiplication by $g_j\in k((p))$ preserves cosets, so
$\sum_jP_{j,\gamma}(p)g_j=0$ and every $P_{j,\gamma}=0$.
\end{proof}

\begin{corollary}[All intrinsic jets over the bounded-support field;
source~17]\label{hol:fh:cor:boundedjets}
For $p=\omega^{-1}$ and $\mathcal F^{\mathrm{bd}}=\mathcal F^{\mathrm{bd}}_\R(\R)
\subseteq\R((t^\R))\subseteq\No$, the jets $\partial_{\mathrm{BM}}^j\mathfrak x_A$
of a finite-private-tail family are algebraically independent over
$\mathcal F^{\mathrm{bd}}$, also with complex coefficients; for an arbitrary
finite family the ideal of each finite jet block is the extension of its
ideal over $\R(p)$.
\end{corollary}

\begin{proof}
Jet monomials lie in $\R((p))$ and are linearly independent over $\R(p)$ by
Theorem~\ref{hol:fh:thm:surreal}; apply Lemma~\ref{hol:fh:lem:bounded}, after
the affine change of Theorem~\ref{hol:fh:thm:relations} in the dependent case.
\end{proof}

\begin{remark}[An absolute base-field boundary; source~17]
\label{hol:fh:rem:absolute}
The witnesses $\mathfrak x_A$ and their jets lie in $\R((p))$, so they cannot be
independent over a base containing them, for instance $\R((p))$ or $\No$; the
external-variable scalar theorem does not contradict this. After a
noncofinal enlargement of the value group all integer supports can become
bounded, and the bounded-support base then absorbs every witness.
\end{remark}

\subsection{An exact domain in the full class}\label{hol:fh:sub:domain}

With $t=\omega^{-1}$ and valuations of surcomplex normal forms in that
convention, put
$\mathcal O_{\langle1\rangle}=\{0\}\cup\{x\in\No[i]^\times:v(x)\ge-m$ for some
$m\in\N\}$, equivalently the $x$ with $\abs x<\omega^m$ for some $m$.

\begin{theorem}[Exact common domain; source~17]\label{hol:fh:thm:domain}
For rapid heights and infinite $A$, the family $(\omega^{-\mathsf b_n}x^n)_{n\in A}$
is strongly summable exactly when $x\in\mathcal O_{\langle1\rangle}$, and every
termwise $D_z$-derivative has the same domain; the domain depends neither on
$A$ nor on the heights.
\end{theorem}

\begin{proof}
Write $x=c\,t^\gamma(1+w)$ with $v(w)>0$. The $n$th term has leading value
$\beta_n=\mathsf b_n+n\gamma$. If $\gamma\ge-m$, the increments of
$\mathsf b_n$ tend to infinity, so $\beta_n$ increases eventually, and the
supports lie in $\beta_n+M(\supp w)$, so Lemma~\ref{hol:lem:certificate}\,(b)
gives strong summability, with $(1+w)^n$ a finite expansion. If $\gamma<-m$
for every $m$, then for $n'>n$, $\beta_{n'}-\beta_n=(\mathsf b_{n'}-\mathsf b_n)
+(n'-n)\gamma<0$, an ordinary integer plus a negative infinite element, so the
leading exponents decrease strictly and Lemma~\ref{hol:lem:nonincreasing}
excludes summability. A $j$th derivative shifts $\beta_n$ by $-j\gamma$ and
multiplies by nonzero integers.
\end{proof}

\begin{proposition}[The domain is a valuation ring; source~17]
\label{hol:fh:prop:domainring}
$\mathcal O_{\langle1\rangle}$ is a valuation subring of $\No[i]$, with units the
elements whose valuations are bounded in absolute value by an integer and
maximal ideal $\{0\}$ together with the elements of positive infinite
valuation.
\end{proposition}

\begin{proof}
Integer lower bounds are preserved by sums and add under products; for $x\ne0$
one of $x,x^{-1}$ lies in the ring, both exactly when $v(x)$ is bounded on
both sides.
\end{proof}

\begin{remark}[The same shape as the domain of $\mathcal T_p$; merge]
\label{hol:fh:rem:domains}
For $p=\omega^{-1}$ the strong domain of $\mathcal T_p$ in its field is
$\mathcal O_{\Delta_\mu}$ (Theorem~\ref{hol:td:thm:domain}), the valuation ring of
the coarsening by the convex subgroup generated by the scale;
$\mathcal O_{\langle1\rangle}$ is the same construction in the whole class
$\No[i]$. Both are proper, as they must be
(Proposition~\ref{hol:ot:prop:fullclass}): $\omega^\omega$ is outside both.
\end{remark}

\begin{proposition}[Faithfulness on the domain; source~17]
\label{hol:fh:prop:faithful}
Every nonzero finite polynomial expression in $z$ and the $D_z$-jets of a
finite-private-tail family $(\mathcal G_A)$ with $p=\omega^{-1}$, with
coefficients in $\No[i]$, is nonzero at some $x\in\mathcal O_{\langle1\rangle}$,
which can be taken of arbitrarily large positive valuation; so the formal
independence over $\No[i](z)$ excludes identities on the common domain.
\end{proposition}

\begin{proof}
The expression $g=\sum g_nz^n$ is nonzero by
Theorem~\ref{hol:fh:thm:allscalars}, in the set-sized field of its coefficients
and $p$. The jets of $\mathcal G_A$ have coefficients of nonnegative value, so
$v(g_n)\ge L$ for one $L$. With $m$ least such that $g_m\ne0$ and a positive
$\rho>v(g_m)-L$, the term of degree $m$ is the unique least one at $x=t^\rho$,
since $v(g_n)+n\rho\ge L+n\rho>v(g_m)+m\rho$ for $n>m$; evaluation and
regrouping are legitimate by Theorem~\ref{hol:fh:thm:domain}.
\end{proof}

\begin{remark}[Three noninterchangeable operations; source~17]
\label{hol:fh:rem:three}
Strong summation on $\mathcal O_{\langle1\rangle}$, cofinal convergence in a
fixed order-unit workspace and convergence in the fine topology of the full
class are different notions, and the third is not inferred from the first
two. The derivative in Proposition~\ref{hol:fh:prop:faithful} is the termwise
$D_z$, not $\partial_{\mathrm{BM}}$ after a substitution $z=z(\omega)$. For
instance $\omega^a$ with real $a$ lies in $\mathcal O_{\langle1\rangle}$, and
$\omega^\omega$ does not: the series are entire on $\C((t^\Q))$ and not on
$\No[i]$.
\end{remark}

\subsection{Omnific codes and floor profiles}\label{hol:fh:sub:omnific}

Let $\Omega$ be a positive infinite surreal, $M_\Omega=\omega^\Omega$, and
\[
 \mathfrak z_A=M_\Omega\mathfrak x_A=\sum_{n\in A}\omega^{\Omega-\mathsf b_n}.
\]
All exponents are positive and the constant term is zero, so
$\mathfrak z_A\in\Oz$, a positive purely infinite omnific integer for infinite
$A$ \cite{repoodg}.

\begin{lemma}[The coding scale is transcendental; source~17]
\label{hol:fh:lem:scale}
$M_\Omega$ is transcendental over $\R((t^\R))$, and the support cosets of
$M_\Omega^jh$ and $M_\Omega^{j'}h'$, $h,h'\in\R((t^\R))$, are disjoint for
$j\ne j'$.
\end{lemma}

\begin{proof}
$j\Omega-r=j'\Omega-r'$ with real $r,r'$ would put $(j-j')\Omega$ in $\R$.
\end{proof}

\begin{theorem}[Independent omnific codes; source~17]\label{hol:fh:thm:codes}
For a finite-private-tail family of infinite sets, the $\mathfrak z_A$ are
algebraically independent over $\mathcal F^{\mathrm{bd}}(M_\Omega)$, hence over
$\R(\omega,M_\Omega)$, and the binary-prefix family gives $2^{\aleph_0}$ of
them. For an arbitrary finite family the algebraic relation ideal over
$\mathcal F^{\mathrm{bd}}(M_\Omega)$ is generated by $\boldsymbol\ell_j\cdot
(Y-M_\Omega\mathsf c_0(p))$, and its transcendence degree is $d_{\mathrm{pat}}$.
\end{theorem}

\begin{proof}
Corollary~\ref{hol:fh:cor:boundedjets} at order zero gives independence of the
pattern numbers over $\mathcal F^{\mathrm{bd}}$ inside $\R((t^\R))$; adjoining
$M_\Omega$ preserves it by Lemma~\ref{hol:fh:lem:scale}, comparing powers of
$M_\Omega$; multiplication by $M_\Omega$ is an invertible substitution. Multiply
\eqref{hol:fh:eq:patterns} by $M_\Omega$ for the general case.
\end{proof}

\begin{remark}[What the codes do not give; source~17]\label{hol:fh:rem:codes}
This is an \emph{algebraic} theorem: intrinsic differential independence of
the $\mathfrak z_A$ would require tracking the derivatives of $M_\Omega$. No
irreducibility, primality or divisibility statement is made, and multiplying
by $M_\Omega$ is not the forbidden evaluation $\mathcal G_A(M_\Omega)$ outside
$\mathcal O_{\langle1\rangle}$.
\end{remark}

For factorial heights and an ordinary real $a\ge0$,
Theorem~\ref{hol:fh:thm:domain} gives
$\mathcal G_A(\omega^{-1},\omega^a)=\sum_{n\in A}\omega^{an-n!}$; put
$\mathsf I_A(a)=\sum_{n\in A,\,(n-1)!\le a}\omega^{an-n!}$.

\begin{theorem}[Factorial floor profile; source~17]\label{hol:fh:thm:floor}
For infinite $A\subseteq\N_{\ge2}$ and real $a\ge0$,
$\lfloor\mathcal G_A(\omega^{-1},\omega^a)\rfloor_{\Oz}=\mathsf I_A(a)$, and
$0<\mathcal G_A(\omega^{-1},\omega^a)-\mathsf I_A(a)<1$; the cutoffs are the
$(n-1)!$ with $n\in A$.
\end{theorem}

\begin{proof}
$an-n!\ge0$ exactly when $(n-1)!\le a$, for finitely many $n$; at most one
exponent is zero, so $\mathsf I_A(a)$ has an integer constant term $0$ or $1$
and lies in $\Oz$. The rest has negative real exponents and coefficients one,
a positive infinitesimal. A positive nonzero omnific integer is at least $1$.
This is the case of \texttt{odg:thm:floor}~\cite{repoodg} in which the
real part is $0$ or $1$ and the infinitesimal part is positive.
\end{proof}

For $A=\N_{\ge2}$ the floors at $a=0,1,2,3,6$ are $0$, $1$, $\omega^2+1$,
$\omega^4+\omega^3$ and $\omega^{12}+\omega^{10}+1$.

\subsection{The role of the growth hypothesis}\label{hol:fh:sub:limits}

Rapid ratios isolate a maximal-degree target against every bounded-degree
competitor, and make the points $(n,\mathsf b_n)$ eventually avoid every
algebraic curve; $\mathsf b_n=n!$, $(n!)^2$ and $2^{n^2}$ qualify.

\begin{proposition}[A polynomial-height obstruction; source~17]
\label{hol:fh:prop:heat}
If $\mathsf b_n=\mathsf h(n)$ for an integer-valued polynomial $\mathsf h$ on
the chosen indices, then $\vartheta_p\mathcal G_A=\mathsf h(\vartheta)\mathcal
G_A$ for every $A$; for $\mathsf b_n=n^2$, $\vartheta_p\mathcal G_A=
\vartheta^2\mathcal G_A$.
\end{proposition}

\begin{proof}
Both sides multiply the coefficient of $p^{\mathsf b_n}z^n$ by $\mathsf h(n)$.
\end{proof}

This concerns \emph{mixed} derivations; it gives no pure $D_z$ equation over a
Hahn coefficient field. With $\mathsf b_n=n^2$ and $p=t^\mu$, $\mathcal
G_{\N_{\ge2}}$ is, up to its first two terms, the series $G_\mu$ of
Section~\ref{hol:ud:sub:surreal}, which is strongly entire at an order unit
and whose differential algebraicity is not decided
(Question~\ref{hol:fh:q:heights}).

For a proposed relation of top degree $\bar d$ and $p$-degree $\bar h$ the
factorial threshold $N\ge\max(3,4\bar d)$, $N!/2>\bar h$ is explicit, and
polarization reduces nonvanishing to a polynomial test at far-out points. This
is a finite certificate given effective access to enough private indices, not
a decision procedure: whether a computable set is infinite, or a pattern
recurs, may be undecidable. Intrinsic differentiation of
$\mathcal G_A(\omega^{-1},z(\omega))$ mixes coefficient and chain-rule terms, so
formal $D_z$ independence with frozen scalars cannot be substituted into it; no
single surreal number is transcendental over all of $\No$. The proofs use
finite polynomials, countable sequences and $(2^{\aleph_0})^{\aleph_0}=
2^{\aleph_0}$; no large-cardinal hypothesis, CH or GCH is needed.

\begin{remark}[Source~17's universal question; merge]
\label{hol:fh:rem:universal}
Source~17's Question~12.1 asks whether, in a characteristic-zero order-unit
Hahn field, every strongly entire solution of an algebraic ordinary
differential equation over $\KK(z)$ must be a polynomial. At its pin that was
Question~\ref{hol:q:nonlinear} as then posed; it is answered negatively by
$\mathcal T_p$ (Theorem~\ref{hol:td:main:boundary}), with least order three
(Theorem~\ref{hol:ot:main:second}), and it is not printed as an open question.
What remains is the classification of Question~\ref{hol:q:nonlinear} and its
companions.
\end{remark}

\subsection{Predecessors, repository comparison and priority}
\label{hol:fh:sub:predecessors}

\paragraph{Literature.}
Differential transcendence of a single ordinary power series with unbounded
ratios of successive nonzero exponents is classical: Ostrowski's 1920 paper,
Section~5 just before Theorem~11, attributes the criterion to
Gr\"onwall~\cite{ostrowski}. Source~17 read Ostrowski in the English
translation of Hansen, Stone and Wolfson; the Gr\"onwall original was not
inspected. So differential transcendence of an individual factorial-gap series
over $\C(p)$ is not new; the unified proof is given because it also controls
joint mixed jets and the whole relation ideal. Source~17 imports the
normal-form identification and the Berarducci--Mantova derivation~\cite{bm,bgexp}
and the normal-form definition of $\Oz$~\cite{lmfactor}; it consulted the
Wikipedia article on surreal numbers for orientation only, and located general
lacunary-series literature not used as a premise. No exhaustive MathSciNet,
zbMATH, thesis, historical-language or citation-network search was made. For
this merge the bibliographic data of the Ostrowski translation and of the
L'Innocente--Mantova preprint were checked; their content was not.

\paragraph{Repository comparison at the pin.}
Source~17 read, at \texttt{bcac55a}, the root README, recursive tree and
catalogue, this report's directory, catalogue entry and source~12's audit, and
the READMEs of the tail-span and bounded-support reports; README claims were
recorded as claims, not verified proofs, and this was not a line-by-line audit
of the 57 catalogued reports, their histories or the Lean declarations. The
tail-span report already has an exact cofinite-span relation theory and
numerical $\partial_{\mathrm{BM}}$ jets for different witnesses (its
\texttt{tail:thm:surreal}, \texttt{tail:eq:BMtriangular} and
\texttt{tail:rem:indexfamily}~\cite{repotail}); source~17's mechanism is
support, not coefficients, and its numerical base differs. The
bounded-support report already has factorial-support algebraic independence,
the witnesses $\eta_A=\sum_nc_A(n)\omega^{-n!}$, algebraically independent over
the bounded-support field, and the descent theorem, and states that it proves no differential
independence~\cite{repobst}; source~17 continues it differentially. These
overlaps are recorded, not hidden.

\paragraph{Priority.}
Not claimed new: single-series gap differential transcendence; normal forms,
the Berarducci--Mantova derivation and $\Oz$; lacunarity, polarization, row
rank, almost-disjoint families and the cardinality of a continuum family;
cofinite-span invariants as an idea; factorial-support algebraic independence
and bounded-support descent; existence of differentially transcendental
surreal numbers in general. The proposed contributions are the complete
ordinary and mixed relation ideals for finite-alphabet rapid-height families
with exact jet degrees, pure $D_z$ independence over the full scalar Hahn
field for strongly entire series with coefficients in $\Q(p)$, arbitrary
finite Hermite data in a continuum family, $\partial_{\mathrm{BM}}$-jet
independence over the bounded-support field, the exact common domain in
$\No[i]$, and the omnific codes and floors; priority is provisional, and a
correct theorem can be a rediscovery.

\subsection{What the single-scale part does not claim}
\label{hol:fh:sub:nonclaims}

Every limitation stated by source~17, in its manuscript, README, audit files or
verification records, is kept.
\begin{enumerate}[label=\textup{(F\arabic*)}]
\item Differential transcendence of one factorial-gap series is classical
(Gr\"onwall, via Ostrowski 1920) and not new; the original was not inspected.
\item Normal forms, the Berarducci--Mantova derivation, $\Oz$, lacunarity,
polarization, almost-disjoint families, cofinite spans and bounded-support
descent are not claimed new; nor is the existence of differentially
transcendental surreal numbers.
\item The universal order-unit problem is not solved by source~17 (it is
answered in this report by source~13; Remark~\ref{hol:fh:rem:universal}).
\item The codes carry no primality, irreducibility or divisibility statement,
and their independence is algebraic, not intrinsic differential.
\item The mixed theorem has base $k(p,z)$ with its two named derivations only.
\item The pure $z$ theorem freezes $E$; the witnesses are not transcendental over
$k((p))$, $\R((t^\R))$ or $\No$.
\item Only finite alphabets are classified; $\mathbf v_n=(1,n)$ shows the
obstruction.
\item The continuum family is not a differential transcendence basis and is not
uniformly computable.
\item The Hermite data are finite.
\item The domain theorem uses strong summation, not the fine topology;
$\partial_{\mathrm{BM}}$ is not $D_z$ after a substitution.
\item Bounded-support descent needs the integer scale to be cofinal.
\item Rapid ratios are sufficient, not necessary; polynomial heights fail the
mixed theorem.
\item The finite witness is not a decision procedure.
\item No large cardinals, CH or GCH are used.
\item No independent referee report and no Lean formalization; the 8{,}608
checks are regression checks, and their number is not a count of theorems.
\item Byte-for-byte reproducibility of the PDF is not claimed.
\item The repository comparison covered three reports' documentation, not every
report or Lean declaration; the literature search was targeted.
\item The floor theorem is proved only for ordinary real $a\ge0$.
\item \emph{[merge]} Remarks~\ref{hol:fh:rem:placement},
\ref{hol:fh:rem:domains} and~\ref{hol:fh:rem:universal} and the comparison with
Proposition~\ref{hol:cf:prop:scalar} are the merge's; no historical priority is
claimed for them.
\end{enumerate}

%======================================================================
% Polynomial-composition part: source 15. Labels use hol:pc:.
%======================================================================
\section{Polynomial-composition rigidity: profiles and the weighted theorem}
\label{hol:pc:sec:setting}

This section and the next contain the material of source~15, delivered in
batch~31 with sources~16--19 and merged last. The dilation theory of
Sections~\ref{hol:sec:qnecessity}--\ref{hol:sec:mixed} and the mixed criterion
of Theorem~\ref{hol:ud:main:mixed} concern arguments $\lambda z$ of degree one,
and there an order unit decides everything. Source~15 replaces the dominant
argument by a polynomial $P(z)$ of degree $d\ge2$. Then no hypothesis on
$\Gamma$ is needed at all: a product of $\bar s$ nonzero linear differential
transforms of $f(P(z))$ cannot be balanced by any polynomial in differential
transforms of $f(P_\nu(z))$ with $\deg P_\nu<d$ whose monomials have
\emph{composition weight} below $\bar sd$, unless $f$ is a polynomial
(Theorem~\ref{hol:pc:main:composition}). The proof keeps the Gauss data of
Section~\ref{hol:nl:sub:scaled}, reads how a polynomial argument rescales them,
and compares scales through an integer convexity inequality. An explicit
degree bound retains all exact indicial resonances; the linear case with
ordinary integer data has a complete omnific solution module; and every affine
hypersurface over $\Q$ is the entire-solution locus of one such equation,
which makes omnific solvability undecidable inside this rigid class.

\subsection{Source, conventions and the merge's choices}
\label{hol:pc:sub:conventions}

\begin{convention}[Setting of the polynomial-composition part]
\label{hol:pc:conv:field}
In Sections~\ref{hol:pc:sec:setting}--\ref{hol:pc:sec:arithmetic}, $k$ is an
arbitrary field of characteristic zero with the trivial valuation, $\Gamma$ is
as in Convention~\ref{hol:conv:group}, $\KK=k((t^\Gamma))$, and $v$, $\lc$,
$\res$, strong summability, $\EK$, $D_z$ and the Gauss data
$\gamma_\delta(f)$, $\operatorname{Act}_\delta(f)$, $\phi_\delta(f)$ of a nonzero
$f\in\EK$ are as in Conventions~\ref{hol:nl:conv:field}
and~\ref{hol:ot:conv:field}, \eqref{hol:nl:eq:gaussdata} and
\eqref{hol:ot:eq:active}; we put $\gamma_\delta(0)=+\infty$ and
$\vartheta=zD_z$. For nonzero $b\in\KK[z]$, $\lc_z(b)\in\KK^\times$ is its
leading coefficient in $z$; $\lc$ remains the leading Hahn coefficient. ``For large
$\delta$'' quantifies over a final segment of $\Gamma$ and presupposes no
countable cofinal sequence. A nonzero linear differential operator is written
$L=\sum_{r=0}^{R_L}b_r(z)D_z^r$ with $b_r\in\KK[z]$, the derivative acting on
what stands to its right. $D_z$ kills $\KK$ and is not the Berarducci--Mantova
derivation.
\end{convention}

For $k=\C$ one has $\KK=\K$ and $\EK=\E$. As for the parts before it, the
generality in $k$ is recorded for this part.

\paragraph{The source.}
Source~15 is \emph{Polynomial-Composition Rigidity at Surreal Scales: Weighted
nonlinear equations, finite degree bounds, and omnific Diophantine
universality} (24 pages, US letter, 23 September 2026), delivered as the
archive \texttt{Surreal\_Polynomial\_Composition\_Rigidity} (batch~31,
manuscript~01). It consulted the repository at
\begin{center}\small\texttt{bcac55ae6fd3f2b354e568ba1d2e94d496a2cc59}\end{center}
through a connector (a container clone failed), reading the root README, the
catalogue, this report's README and source~10's proof audit, an omnific
Diophantine audit, and the READMEs of the quotient and lattice reports; its
shipped file \path{15-polynomial-composition-SOURCE_AUDIT.md} lists them. At
that pin this report was the five-source text (sources~08--12). Source~15's
description of the project, linear rigidity, dilations, first-order nonlinear
rigidity and ``a higher-order theorem in fields whose value group has no order
unit'', was accurate there; the nonlinear picture has since been completed by
Theorems~\ref{hol:td:main:boundary}, \ref{hol:ot:main:second}
and~\ref{hol:nr:main:third}, which it does not use. Its statement that ``a
keyword search for Mahler material returned no matches'' was not correct at its
pin: the single-dilation report already discussed Mahler difference algebra
and cited Faverjon and Roques on linear Mahler equations with Hahn-series
solutions~\cite{faverjonroques,reposdhs}, and the computable-surreals report
cited the same paper; this report mentioned only the Skolem--Mahler--Lech
theorem. The comparison is corrected in Section~\ref{hol:pc:sub:predecessors}.

\paragraph{Renamed symbols.}
The following symbols of source~15 would collide with symbols of the earlier
parts. They are renamed here, in the new material only. The most important
row is the first: source~15's two indicial polynomials carry hol's names in
\emph{reverse}.
{\small
\begin{longtable}{@{}>{\raggedright\arraybackslash}p{0.27\textwidth}
>{\raggedright\arraybackslash}p{0.2\textwidth}
>{\raggedright\arraybackslash}p{0.45\textwidth}@{}}
\toprule
Source~15 & Here & Reason\\
\midrule
\endhead
exact indicial $\chi_L$; residue indicial $\bar\chi_L$ & $I_L$; $\chi_L$ &
For the linear differential polynomial $P=L[Y]$ they are exactly $I_P$ and
$\chi_P$ of \eqref{hol:nl:eq:chiI}: \emph{source~15's $\chi_L$ is this report's
$I_L$, not its $\chi_L$}.\\
shift $h_L$; $\alpha_L$ & $w_L$; $\beta_L$ & The corner weight $w_P$ and value
$\beta_P$ of \eqref{hol:nl:eq:dw} for $P=L[Y]$.\\
radius $\rho$; height $M_f(\rho)$; $I_f(\rho)$; $\operatorname{in}_\rho(f)$ & scale
$\delta$; $-\gamma_\delta(f)$; $\operatorname{Act}_\delta(f)$; $\phi_\delta(f)$ &
The Gauss data; source~15's height is the \emph{negative} of the Gauss value,
so its inequalities are reversed here.\\
reference radius $R$; $\mathcal T_\rho$; $\sigma$ (Proposition~2.4) &
$\delta_{\mathrm{ref}}$; $\mathcal E^{(\delta)}$; $\delta'$ & $R$ is an order
bound, $\mathcal T_p$ the theta descent, $\sigma$ a dilation.\\
$\kappa=v(\lc_z P)$ & $v(\lc_z P)$ & $\kappa_P$ is the degree cutoff
\eqref{hol:nl:eq:kappa}.\\
leading coefficient $\operatorname{lc}(P)$ of a polynomial & $\lc_z(P)$ &
$\lc$ is the leading Hahn coefficient of Section~\ref{hol:sub:conventions}.\\
$A(z)$; $Q(z,Y)$; $q$ lower arguments; $s$ top factors; $T_\nu$ &
$c_*(z)$; $\mathsf G(z,Y)$; $\bar m$; $\bar s$; $L'_\nu$ & $A[f]$, $Q_s$, $q$
(a dilation), $s$ (an order) and $T$ (an indeterminate) have other
meanings.\\
composition weight $w(\boldsymbol a)$; $H$, $b$; $H_{\rm top}$, $U_{\boldsymbol
a}$; $h_{\boldsymbol a}$, $b_{\boldsymbol a}$ & $\operatorname{cw}(\mathbf a)$;
$w_{\mathrm{top}}$, $\beta_{\mathrm{top}}$; $\gamma_{\mathrm{top}}$,
$\gamma_{\mathbf a}$; $w_{\mathbf a}$, $\beta_{\mathbf a}$ & $w_P$ is a corner
weight and $\wt$ an Euler weight; $H\in\Gamma$.\\
resonance set $\mathcal R$; bound $B$ & $\operatorname{Rsn}$;
$B_{\mathrm{cmp}}$ & $\mathcal R$ is a set of right sides in
Section~\ref{hol:nl:sec:first}; $B$ a threshold.\\
$D=\Z$ or $\Z[i]$; $R_{D,\Gamma}$; $R$ & $\mathfrak o$; $\mathfrak O_{\mathfrak
o,\Gamma}$; $\mathfrak R$ & $D_z$ is the derivative.\\
$E_z=zD$; projectors $S_j$; $L_{d,N}$ & $\vartheta$; $\mathrm{Pr}_j$;
$\mathrm{Ann}_{d,N}$ & $E_\eta$ is the formal exponential, $\mathsf S$ a shift.\\
$q=\max\abs{\boldsymbol a}$, $T=\max\sum ja_j$ (encoder) & $\deg p$,
$\operatorname{wd}(p)$ & As above.\\
theta $F$; even series $H$; char-$p$ series $H$; rank-infinity $G$ & $\Theta_p$;
$f_{\mathrm{ev}}$; $U$ of Example~\ref{hol:nl:ex:charp}; $G$ (kept) &
\eqref{hol:eq:theta}; $F$ is \eqref{hol:eq:explicitF}.\\
``Mahler equation'' & kept, with its meaning & Here it means equations in
$f(z^{b^j})$, argument substitution; not the exponent dilation $S_d$ of the
autonomous-dilation report~\cite{repoadr}, and not the scalar dilation of the
single-dilation report.\\
Theorems~5.1, 6.1, 7.1, 8.2, 9.1 & Theorem~\ref{hol:pc:main:composition} and
Theorems~\ref{hol:pc:thm:bound}, \ref{hol:pc:thm:matrix},
\ref{hol:pc:thm:linarith}, \ref{hol:pc:thm:universal} & Theorems~A--N of this
report are different theorems.\\
\bottomrule
\end{longtable}}
The shipped audit file \path{15-polynomial-composition-SOURCE_AUDIT.md} keeps
the source's notation.

\paragraph{Words with two meanings.}
A \emph{pullback} is $f\mapsto f\circ P$ for a polynomial $P$; a
\emph{dilation} is the degree-one case $P=\lambda z$. The \emph{composition
weight} of a monomial is Definition~\eqref{hol:pc:eq:weight}; it is neither the
corner weight $w_P$ nor the Euler weight $\wt$. An \emph{exact} resonance is an
ordinary integer root of $I_L$, a \emph{residue} resonance one of $\chi_L$;
Example~\ref{hol:nl:ex:resonance} shows they differ.

\paragraph{Results printed once.}
Source~15's Definition~2.1 and Lemma~2.2 are Definition~\ref{hol:def:entire} and
Lemma~\ref{hol:lem:certificate}\,(a); its Theorem~2.3, the exact entire-support
criterion, is Proposition~\ref{hol:cf:prop:criterion}, proved there by the same
two-scale argument; its Proposition~2.4 is Proposition~\ref{hol:prop:closure}
together with closure under polynomial composition, which is contained in
\texttt{ent:prop:operations} of the entire-functions report~\cite{repoentire}
(entire composition) and is also proved directly below; its Lemma~3.1 is
Corollary~\ref{hol:nl:cor:gauss} with the derivative estimate of
Lemma~\ref{hol:pc:lem:operatorupper} for $L=D_z^r$; its Proposition~3.2 is the
residue argument of Theorem~\ref{hol:th:thm:sampling}, over any $k$ (compare
Proposition~\ref{hol:prop:faithful}); the first half of its Lemma~3.3 is
Lemma~\ref{hol:nl:lem:active}; its Remark~3.6, that the Gauss value is not
$v(f(t^{-\delta}))$, is the remark after \eqref{hol:nl:eq:gaussdata}; its
Remark~4.4, $L=zD_z-N+\varepsilon$, is Example~\ref{hol:nl:ex:resonance}; the
conformal embedding in its Section~8.1 is \eqref{hol:eq:embedding}; the partial
theta series, the characteristic-$p$ series, the rank-infinity series and the
failure on the full Conway field in its Section~11 are \eqref{hol:eq:theta},
Example~\ref{hol:nl:ex:charp}, the series $F$ of Theorem~\ref{hol:thm:explicit} (source~15
uses $\bigoplus_{n\ge1}\R e_n$, this report $\Gamma_\infty$ of
\eqref{hol:eq:gammainfty}), Section~\ref{hol:sub:notallsurcomplex}
and Proposition~\ref{hol:ot:prop:fullclass}; and its remark that uncountable
cofinality forces $\EK=\KK[z]$ is the remark in Section~\ref{hol:nr:sub:examples}.

\paragraph{Where each numbered result went.}
Proposition~3.5 (pullback profile), Lemma~3.4 (integer convexity) and the
second half of Lemma~3.3 are Proposition~\ref{hol:pc:prop:pullback},
Lemma~\ref{hol:pc:lem:convex} and Lemma~\ref{hol:pc:lem:superlinear};
Lemmas~4.1, 4.3 and Proposition~4.2 are Lemmas~\ref{hol:pc:lem:operatorupper},
\ref{hol:pc:lem:polyleading} and Proposition~\ref{hol:pc:prop:operatorexact};
Theorem~5.1 and Corollaries~5.2, 5.3 are
Theorem~\ref{hol:pc:main:composition} (with its proof in
Section~\ref{hol:pc:sub:weighted}) and Corollaries~\ref{hol:pc:cor:linear},
\ref{hol:pc:cor:newton}; Theorem~6.1, Corollary~6.2 and Proposition~6.3 are
Theorem~\ref{hol:pc:thm:bound}, Corollary~\ref{hol:pc:cor:linearbound} and
Proposition~\ref{hol:pc:prop:extension}; Theorem~7.1 is
Theorem~\ref{hol:pc:thm:matrix}; Proposition~8.1 and Theorem~8.2 are
Proposition~\ref{hol:pc:prop:ompoly} and Theorem~\ref{hol:pc:thm:linarith};
Theorem~9.1, Remark~9.2 and Corollaries~9.3, 9.4 are
Theorem~\ref{hol:pc:thm:universal}, Remark~\ref{hol:pc:rem:universal} and
Corollaries~\ref{hol:pc:cor:undecidable}, \ref{hol:pc:cor:ten}; the examples
of Sections~10 and~11 are in Sections~\ref{hol:pc:sub:examples}
and~\ref{hol:pc:sub:boundaries}; Questions~12.1--12.8 are
Questions~\ref{hol:pc:q:critical}--\ref{hol:pc:q:formal}; Section~13 and the
appendices are merged into Sections~\ref{hol:pc:sub:predecessors},
\ref{hol:pc:sub:nonclaims} and~\ref{hol:sec:verification}.

\paragraph{The merge's choices.}
The review of source~15 for this merge found its proofs correct; the worked
equations were rechecked with SymPy. The merge adds three points, each marked.
First, Remark~\ref{hol:pc:rem:dichotomy} states the relation to
Theorems~\ref{hol:main:q} and~\ref{hol:ud:main:mixed}, to
Theorem~\ref{hol:thm:mixed} and to Corollary~\ref{hol:cor:detect}: with a
strictly dominant argument of degree at least two the order-unit dichotomy
disappears. Second, the Gauss data are those of this report, with source~15's
signs reversed, and its two indicial polynomials are identified with $I_P$ and
$\chi_P$. Third, the comparison with the Mahler and omnific literature of the
collection is completed (Section~\ref{hol:pc:sub:predecessors}). Source~15's
theorems, proofs, examples and non-claims are otherwise printed as delivered,
in the renamed notation.

\subsection{Pullback profiles}\label{hol:pc:sub:profiles}

For $\delta\in\Gamma$ let $\mathcal E^{(\delta)}$ be the set of
$\sum_nb_nz^n\in\KK[[z]]$ with $v(b_n)-n\delta\to+\infty$ cofinally. After
$z=t^{-\delta}X$ its elements have coefficients whose values tend cofinally to
infinity; it is closed under products (finite convolutions, with the two
minima normalized to zero), sums and $D_z$ (the multiplier $n$ has value zero),
and a family $(H_n)$ in it whose minimal weighted coefficient values tend
cofinally to infinity has a coefficientwise sum in it, which may be regrouped.
By Proposition~\ref{hol:cf:prop:criterion}, $\EK=\bigcap_\delta\mathcal
E^{(\delta)}$.

\begin{proposition}[Polynomial composition; source~15]
\label{hol:pc:prop:composition}
For $f=\sum a_nz^n\in\EK$ and $P\in\KK[z]$, the series $f\circ P=\sum_na_nP^n$
is defined coefficientwise by strongly summable families and lies in $\EK$;
$D_z(f\circ P)$ may be computed termwise. $P$ need not vanish at $0$.
\end{proposition}

\begin{proof}
Fix $\delta$ and put $\delta'=\max_j(j\delta-v([z^j]P))$. Every coefficient of
$a_nP^n$, weighted at $\delta$, has value at least $v(a_n)-n\delta'$, and these
tend cofinally to infinity by Proposition~\ref{hol:cf:prop:criterion}; so
$\sum a_nP^n$ exists in $\mathcal E^{(\delta)}$, with coefficients independent of
$\delta$. It lies in every $\mathcal E^{(\delta)}$. The same bounds justify
termwise differentiation.
\end{proof}

Ordinary formal substitution into a power series needs $P(0)=0$; strong
entireness supplies the convergence that replaces it.

\begin{lemma}[Superlinear escape; source~15]\label{hol:pc:lem:superlinear}
If $f\in\EK$ is not a polynomial, then for every integer $C$, every positive
integer $d$ and every fixed $\kappa\in\Gamma$,
$\gamma_{d\delta-\kappa}(f)+C\delta\to-\infty$ as $\delta\to+\infty$, that is,
it is eventually below every element of $\Gamma$.
\end{lemma}

\begin{proof}
Choose $m$ with $md>C$ and $a_m\ne0$. Then
$\gamma_{d\delta-\kappa}(f)+C\delta\le v(a_m)+m\kappa-(md-C)\delta$, and positive
integer multiples of an unbounded variable are cofinal, whether or not $\Gamma$
has an order unit.
\end{proof}

\begin{lemma}[Integer convexity; source~15]\label{hol:pc:lem:convex}
Let $d>e>0$ be integers and $\sigma,\tau,\delta_{\mathrm{ref}}\in\Gamma$ with
$d\tau\le e\sigma+(d-e)\delta_{\mathrm{ref}}$. Then for every nonzero $f\in\EK$
\begin{equation}\label{hol:pc:eq:convex}
 d\,\gamma_\tau(f)\ge e\,\gamma_\sigma(f)+(d-e)\,\gamma_{\delta_{\mathrm{ref}}}(f).
\end{equation}
\end{lemma}

\begin{proof}
For $a_n\ne0$, $n\ge0$ gives $d(v(a_n)-n\tau)\ge e(v(a_n)-n\sigma)+(d-e)(v(a_n)-n
\delta_{\mathrm{ref}})$, which is at least the right side of
\eqref{hol:pc:eq:convex}; take the minimum. No division in $\Gamma$ occurs.
\end{proof}

\begin{proposition}[Exact pullback profile; source~15]
\label{hol:pc:prop:pullback}
Let $P=cz^d+\cdots$ have degree $d\ge1$. For all large $\delta$, depending on
$P$ only, and every nonzero $f\in\EK$,
\begin{equation}\label{hol:pc:eq:pullback}
 \gamma_\delta(f\circ P)=\gamma_{d\delta-v(c)}(f),\qquad
 \operatorname{Act}_\delta(f\circ P)=d\operatorname{Act}_{d\delta-v(c)}(f),\qquad
 \phi_\delta(f\circ P)(X)=\phi_{d\delta-v(c)}(f)(\bar cX^d),
\end{equation}
with $\bar c=\res(t^{-v(c)}c)$.
\end{proposition}

\begin{proof}
For large $\delta$ the top term uniquely attains the Gauss value of $P$:
$v(c)-d\delta<v([z^j]P)-j\delta$ for $j<d$. With $\sigma=d\delta-v(c)$, the
polynomial $t^\sigma P(t^{-\delta}X)$ has integral coefficients and residue
$\bar cX^d$, and Proposition~\ref{hol:pc:prop:composition} with reduction gives
$\res\bigl(t^{-\gamma_\sigma(f)}f(P(t^{-\delta}X))\bigr)=\phi_\sigma(f)(\bar
cX^d)$, a nonzero polynomial with support $d\operatorname{Act}_\sigma(f)$.
\end{proof}

\subsection{Operators and their two indicial polynomials}
\label{hol:pc:sub:operators}

For $L=\sum_rb_r(z)D_z^r\ne0$ put
\begin{equation}\label{hol:pc:eq:indicial}
 w_L=\max_{b_r\ne0}(\deg b_r-r),\qquad
 I_L(T)=\sum_{\deg b_r-r=w_L}\lc_z(b_r)\fall Tr,\qquad
 \beta_L=\min_{\deg b_r-r=w_L}v(\lc_z b_r),
\end{equation}
and $\chi_L(T)=\sum\res(t^{-\beta_L}\lc_z b_r)\fall Tr$ over the $r$ with
$\deg b_r-r=w_L$ and $v(\lc_z b_r)=\beta_L$. For the linear differential
polynomial $L[Y]$ these are its corner weight $w_P$, full corner polynomial
$I_P$, corner value $\beta_P$ and residue corner polynomial $\chi_P$
(\eqref{hol:nl:eq:dw}, \eqref{hol:nl:eq:chiI}). Both $I_L$ and $\chi_L$ are
nonzero, since distinct falling factorials are independent, and $\chi_L$ has
finitely many roots in $\N$.

\begin{lemma}[Uniform operator bound; source~15]
\label{hol:pc:lem:operatorupper}
For all large $\delta$, depending on $L$ only, and every $u\in\EK$,
$\gamma_\delta(Lu)\ge\gamma_\delta(u)-w_L\delta+\beta_L$.
\end{lemma}

\begin{proof}
The term $c_{rj}z^jD_z^ru$, $c_{rj}=[z^j]b_r\ne0$, has Gauss value at least
$\gamma_\delta(u)-(j-r)\delta+v(c_{rj})$. For $j-r<w_L$ this exceeds the bound
for large $\delta$; at $j-r=w_L$, $j=\deg b_r$ and $v(c_{rj})\ge\beta_L$.
\end{proof}

\begin{proposition}[Exact eventual operator profile; source~15]
\label{hol:pc:prop:operatorexact}
If $u\in\EK$ is not a polynomial, then for all large $\delta$
\[
 \gamma_\delta(Lu)=\gamma_\delta(u)-w_L\delta+\beta_L,\qquad
 \phi_\delta(Lu)=\sum_{n\in\operatorname{Act}_\delta(u)}[X^n]\phi_\delta(u)\,
 \chi_L(n)X^{n+w_L};
\]
in particular $Lu\ne0$.
\end{proposition}

\begin{proof}
Take $\delta$ beyond Lemma~\ref{hol:pc:lem:operatorupper} and so large that
every active index exceeds $R_L$ and every root of $\chi_L$ in $\N$
(Lemma~\ref{hol:nl:lem:active}). Normalizing at the proposed value, the terms of
smaller shift and the top-shift terms with coefficient value above $\beta_L$
have positive value, and at input degree $n$ the surviving terms have the
residue multiplier $\chi_L(n)\ne0$; inactive coefficients contribute nothing.
The exponents $n+w_L$ are distinct and nonnegative.
\end{proof}

This is the corner identity of Section~\ref{hol:nl:sub:corner} for the linear
polynomial $L[Y]$, applied at every exterior scale.

\begin{lemma}[Exact polynomial leading term; source~15]
\label{hol:pc:lem:polyleading}
For a polynomial $u$ of degree $n$ with leading coefficient $a$,
$\deg Lu\le n+w_L$; if $I_L(n)\ne0$ equality holds with leading coefficient
$aI_L(n)$, and if $I_L(n)=0$ the degree is smaller or $Lu=0$.
\end{lemma}

\begin{proof}
The coefficient at degree $n+w_L$ is the top-shift contribution of $az^n$,
namely $aI_L(n)$; everything else has smaller degree. If $n+w_L<0$ every term
vanishes, and the falling factorials give $I_L(n)=0$.
\end{proof}

A residue resonance $\chi_L(n)=0$ need not be an exact one
(Example~\ref{hol:nl:ex:resonance}): residue resonance controls the exterior
proof, exact resonance the degree bounds.

\subsection{The weighted composition theorem}\label{hol:pc:sub:weighted}

Fix $P$ of degree $d\ge2$, nonzero operators $L_1,\ldots,L_{\bar s}$ ($\bar
s\ge1$), a nonzero $c_*\in\KK[z]$, and for $1\le\nu\le\bar m$ ($\bar m\ge0$)
nonconstant $P_\nu$ of degree $e_\nu<d$ and nonzero operators $L'_\nu$. Let
$\mathsf G(z,Y_1,\ldots,Y_{\bar m})=\sum_{\mathbf a\in\mathcal A}c_{\mathbf
a}(z)Y^{\mathbf a}$ in collected form, with nonzero $c_{\mathbf a}\in\KK[z]$
($\mathcal A=\varnothing$ if $\mathsf G=0$), and put
\begin{equation}\label{hol:pc:eq:weight}
 \operatorname{cw}(\mathbf a)=\sum_{\nu=1}^{\bar m}a_\nu e_\nu .
\end{equation}
The equation is
\begin{equation}\label{hol:pc:eq:main}
 c_*(z)\prod_{\ell=1}^{\bar s}L_\ell\bigl(f(P(z))\bigr)
 =\mathsf G\bigl(z,L'_1(f(P_1(z))),\ldots,L'_{\bar m}(f(P_{\bar m}(z)))\bigr).
\end{equation}

\begin{proof}[Proof of Theorem~\ref{hol:pc:main:composition}]
Suppose $f$ is not a polynomial. Put $\sigma=d\delta-v(\lc_z P)$ and
$\sigma_\nu=e_\nu\delta-v(\lc_z P_\nu)$, and take $\delta$ large for the finitely
many profiles below. By Proposition~\ref{hol:pc:prop:pullback} and
Lemma~\ref{hol:nl:lem:active}, $f\circ P$ is not a polynomial, so
Proposition~\ref{hol:pc:prop:operatorexact} applies to each $L_\ell(f\circ P)$,
and by Corollary~\ref{hol:nl:cor:gauss} the left side has the exact Gauss value
\begin{equation}\label{hol:pc:eq:top}
 \gamma_{\mathrm{top}}(\delta)=\bar s\,\gamma_\sigma(f)-w_{\mathrm{top}}\delta
 +\beta_{\mathrm{top}},\qquad
 w_{\mathrm{top}}=\deg c_*+\sum_\ell w_{L_\ell},\quad
 \beta_{\mathrm{top}}=v(\lc_z c_*)+\sum_\ell\beta_{L_\ell}.
\end{equation}
Choose one $\delta_{\mathrm{ref}}$ with $(d-e_\nu)\delta_{\mathrm{ref}}\ge e_\nu
v(\lc_z P)-d\,v(\lc_z P_\nu)$ for every $\nu$; then $d\sigma_\nu\le e_\nu\sigma+
(d-e_\nu)\delta_{\mathrm{ref}}$ identically in $\delta$, and
Lemma~\ref{hol:pc:lem:convex} gives $d\gamma_{\sigma_\nu}(f)\ge e_\nu
\gamma_\sigma(f)+(d-e_\nu)\gamma_{\delta_{\mathrm{ref}}}(f)$. With
$w_{\mathbf a}=\deg c_{\mathbf a}+\sum_\nu a_\nu w_{L'_\nu}$ and
$\beta_{\mathbf a}=v(\lc_z c_{\mathbf a})+\sum_\nu a_\nu\beta_{L'_\nu}$, the monomial
$\mathbf a$ has, by Lemma~\ref{hol:pc:lem:operatorupper} and
Proposition~\ref{hol:pc:prop:pullback}, Gauss value at least
$\gamma_{\mathbf a}(\delta)=\sum_\nu a_\nu\gamma_{\sigma_\nu}(f)-w_{\mathbf
a}\delta+\beta_{\mathbf a}$, also when a factor vanishes. Hence, with
$\abs{\mathbf a}=\sum a_\nu$,
\begin{equation}\label{hol:pc:eq:dominance}
 d\bigl(\gamma_{\mathbf a}-\gamma_{\mathrm{top}}\bigr)
 \ge-\bigl(\bar sd-\operatorname{cw}(\mathbf a)\bigr)\gamma_\sigma(f)
 +d(w_{\mathrm{top}}-w_{\mathbf a})\delta
 +\bigl(d\abs{\mathbf a}-\operatorname{cw}(\mathbf a)\bigr)
 \gamma_{\delta_{\mathrm{ref}}}(f)+d(\beta_{\mathbf a}-\beta_{\mathrm{top}}).
\end{equation}
The coefficient $\bar sd-\operatorname{cw}(\mathbf a)$ is a positive integer, so
by Lemma~\ref{hol:pc:lem:superlinear} the right side is eventually above every
element of $\Gamma$. On a common final segment, then, every monomial of the right
side has strictly larger Gauss value than the left side, which is impossible
for an identity; if $\mathsf G=0$ the nonzero left side already contradicts it.
\end{proof}

\begin{corollary}[Mixed differential--Mahler equations; source~15]
\label{hol:pc:cor:linear}
Let $P_0,\ldots,P_m\in\KK[z]$ be nonconstant with $d_m=\deg P_m>\max_{j<m}\deg
P_j$ and $d_m\ge2$, let $L_0,\ldots,L_m$ be linear differential operators over
$\KK[z]$ with $L_m\ne0$, and $g\in\KK[z]$. Every strongly entire solution of
\begin{equation}\label{hol:pc:eq:linear}
 \sum_{j=0}^mL_j\bigl(f(P_j(z))\bigr)=g(z)
\end{equation}
is a polynomial; so is every solution with rational coefficients, after a
common denominator is cleared.
\end{corollary}

\begin{proof}
Move the lower terms to the right: $\bar s=1$, $P=P_m$, $c_*=1$, a linear
$\mathsf G$ whose monomials have weights $\deg P_j<d_m$, and a forcing term of
weight zero. Clearing denominators multiplies the identity by a polynomial and
moves no coefficient through a derivative.
\end{proof}

With $P_j=z^{b^j}$, $b\ge2$, every strongly entire solution of a linear Mahler
equation, or of one whose pullbacks are multiplied by finite-order
differential operators, is a polynomial. The pure complex-analytic
antecedent is Bell, Coons and Rowland's~\cite[Lemma~6]{bcr}; the arbitrary-rank
mixed statement is source~15's.

\begin{remark}[Degree at least two removes the order-unit dichotomy; merge]
\label{hol:pc:rem:dichotomy}
The lower arguments in Corollary~\ref{hol:pc:cor:linear} may be dilations
$\lambda z$, of degree one. So every linear equation in $D_z$, finitely many
dilations $\sigma_\lambda$, and at least one polynomial argument of degree
$d\ge2$ strictly above the degrees of all other arguments has only polynomial
strongly entire solutions, \emph{for every $\Gamma$}, with or without an order
unit and whatever the valuations of the dilation parameters. By contrast,
Theorem~\ref{hol:main:q}, Theorem~\ref{hol:ud:main:mixed} and
Corollary~\ref{hol:cor:detect} concern equations all of whose arguments have
degree one; there a nonpolynomial solution exists exactly when some
(relative) dilation valuation is an order unit, and the partial theta series
is the witness. A dominant argument of degree at least two therefore detects
nothing about order units. Theorem~\ref{hol:thm:mixed} and
Corollary~\ref{hol:ud:cor:unitrigidity} are the degree-one mixed statements. The
hypothesis of a \emph{strictly} dominant degree is used:
Section~\ref{hol:pc:sub:boundaries} shows that equal degrees of the top
arguments, including degree one, allow nonpolynomial solutions. Nor does
Theorem~\ref{hol:pc:main:composition} say anything about algebraic differential
equations without a pullback, which are Theorems~\ref{hol:nl:main:first},
\ref{hol:td:main:boundary}, \ref{hol:ot:main:second}
and~\ref{hol:nr:main:third}.
\end{remark}

\begin{corollary}[A necessary Newton-weight condition; source~15]
\label{hol:pc:cor:newton}
If $f\in\EK$ is not a polynomial, every relation of the form
\eqref{hol:pc:eq:main} that it satisfies has a monomial of weight at least
$\bar sd$.
\end{corollary}

This restricts possible relations; it is not an algebraic or differential
independence statement, and relations with equal weights, or with cancellation
among several top-weight expressions, are not classified.

\subsection{Degree bounds and finite coefficient loci}\label{hol:pc:sub:degree}

Under the hypotheses of Theorem~\ref{hol:pc:main:composition} let
$\operatorname{Rsn}=\{N\in\N:I_{L_\ell}(dN)=0\text{ for some }\ell\}$, a finite
set, and
\begin{equation}\label{hol:pc:eq:B}
 B_{\mathrm{cmp}}=\max\Bigl(\{0\}\cup\operatorname{Rsn}\cup\Bigl\{\Bigl\lfloor
 \frac{w_{\mathbf a}-w_{\mathrm{top}}}{\bar sd-\operatorname{cw}(\mathbf a)}
 \Bigr\rfloor:\mathbf a\in\mathcal A\Bigr\}\Bigr)\in\N ,
\end{equation}
a ratio of ordinary integers, with no division in $\Gamma$.

\begin{theorem}[Resonant degree bound; source~15]\label{hol:pc:thm:bound}
Every nonzero strongly entire solution of \eqref{hol:pc:eq:main} has degree at
most $B_{\mathrm{cmp}}$; the zero series is tested directly, and the bound need not
be attained.
\end{theorem}

\begin{proof}
The solution is a polynomial. If its degree $N$ exceeds $B_{\mathrm{cmp}}$, then
$N\notin\operatorname{Rsn}$, and by Lemma~\ref{hol:pc:lem:polyleading} the left
side has degree $\bar sdN+w_{\mathrm{top}}$ with leading coefficient
$\lc_z(c_*)a^{\bar s}\lc_z(P)^{\bar sN}\prod_\ell I_{L_\ell}(dN)\ne0$, while the
monomial $\mathbf a$ has degree at most $\operatorname{cw}(\mathbf a)N+w_{\mathbf
a}$, which the definition of $B_{\mathrm{cmp}}$ makes smaller.
\end{proof}

\begin{corollary}[Linear bound; source~15]\label{hol:pc:cor:linearbound}
For \eqref{hol:pc:eq:linear}, with $w_j=w_{L_j}$, $d=d_m$ and
$\operatorname{Rsn}_m=\{N\in\N:I_{L_m}(dN)=0\}$, every nonzero entire solution
has degree at most
\[
 B_{\mathrm{lin}}=\max\Bigl(\{0\}\cup\operatorname{Rsn}_m\cup\Bigl\{\Bigl\lfloor
 \frac{w_j-w_m}{d-\deg P_j}\Bigr\rfloor:j<m,\ L_j\ne0\Bigr\}\cup\Bigl\{\Bigl\lfloor
 \frac{\deg g-w_m}{d}\Bigr\rfloor:g\ne0\Bigr\}\Bigr).
\]
\end{corollary}

Substituting $f=u_0+u_1z+\cdots+u_Bz^B$ with $B=B_{\mathrm{cmp}}$ and comparing
coefficients gives finitely many polynomial equations in the $u_j$ over the
field generated by the data, whose solutions in $\KK$ are exactly the entire
solutions; linear equations give a linear system. This is the composition
analogue of Theorem~\ref{hol:nl:thm:degrees} and
Corollary~\ref{hol:nl:cor:locus}. It is not a decision procedure: one must
represent the data and decide equalities, find the integer roots in
$\operatorname{Rsn}$, and solve the finite system. For rational data the first
two are effective; for arbitrary named surreal coefficients no algorithm is
asserted, and over the omnific integers even integer data can be undecidable
(Corollary~\ref{hol:pc:cor:undecidable}).

\begin{proposition}[Extension-stable coefficient description; source~15]
\label{hol:pc:prop:extension}
For an injective field map $\KK\hookrightarrow\KK'$ between Hahn fields of the
standing type, with the finite data transported, the same bound and the same
finite coefficient equations describe the entire solutions over $\KK'$; no
entire solution of higher degree appears.
\end{proposition}

\begin{proof}
Degrees, shifts and the vanishing of $I_{L_\ell}(dN)$ are preserved and
reflected, so $B_{\mathrm{cmp}}$ is unchanged; apply
Theorems~\ref{hol:pc:main:composition} and~\ref{hol:pc:thm:bound} over
$\KK'$. Polynomials are entire.
\end{proof}

This concerns the solutions of a fixed equation, which are polynomials; an
arbitrary infinite series can lose entireness under extension
(Section~\ref{hol:pc:sub:boundaries}).

\subsection{Vector equations}\label{hol:pc:sub:matrix}

For $F=(f_1,\ldots,f_r)\in\EK^r$ put $\gamma_\delta(F)=\min_i\gamma_\delta(f_i)$;
Lemmas~\ref{hol:pc:lem:convex} and~\ref{hol:pc:lem:superlinear} hold for it,
minimizing over components and indices.

\begin{theorem}[Matrix polynomial-composition rigidity; source~15]
\label{hol:pc:thm:matrix}
Let $P_0,\ldots,P_m$ be nonconstant, $d_m=\deg P_m>\max_{j<m}\deg P_j=:\max
d_j$, $d_m\ge2$, $A_0,\ldots,A_m\in\operatorname{Mat}_r(\KK[z])$ with $\det
A_m\ne0$, and $b\in\KK[z]^r$. Every $F\in\EK^r$ with
$\sum_{j=0}^mA_j(z)F(P_j(z))=b(z)$ is a polynomial vector. With $a=\det A_m$,
$C_j=\operatorname{adj}(A_m)A_j$, $c=\operatorname{adj}(A_m)b$ and entry
degrees, the largest component degree of a nonzero solution is at most
\[
 B_{\mathrm{mat}}=\max\Bigl(\{0\}\cup\Bigl\{\Bigl\lfloor\frac{\deg C_j-\deg a}
 {d_m-d_j}\Bigr\rfloor:C_j\ne0\Bigr\}\cup\Bigl\{\Bigl\lfloor\frac{\deg c-\deg a}
 {d_m}\Bigr\rfloor:c\ne0\Bigr\}\Bigr).
\]
\end{theorem}

\begin{proof}
Multiplying by the adjugate gives $aF(P_m)=c-\sum_{j<m}C_jF(P_j)$. If a component
is nonpolynomial, the left side has Gauss value
$\gamma_{d_m\delta-v(\lc_z P_m)}(F)-\deg(a)\delta+v(\lc_z a)$ for large $\delta$
(componentwise, by Proposition~\ref{hol:pc:prop:pullback}), while the right
terms are bounded below by $\gamma_{d_j\delta-v(\lc_z P_j)}(F)$ plus a fixed
integer multiple of $\delta$ and a fixed element of $\Gamma$; the argument of
Theorem~\ref{hol:pc:main:composition} with weights $d_j<d_m$ makes every right
term strictly larger in value. For the degree bound compare $\deg a+d_mN$ with
$\deg c$ and $\deg C_j+d_jN$.
\end{proof}

Nonsingularity cannot be dropped: with $A_m=\operatorname{diag}(1,0)$ and all
other terms zero, every $(0,f)$ is a solution. Matrices of differential
operators are not treated (Question~\ref{hol:pc:q:matrix}).

\section{Omnific coefficients, universality and boundaries of composition
rigidity}\label{hol:pc:sec:arithmetic}

\subsection{Omnific coefficients and linear solution modules}
\label{hol:pc:sub:omnific}

For $(\mathfrak o,k)=(\Z,\R)$ or $(\Z[i],\C)$ put $\mathfrak O_{\mathfrak
o,\Gamma}=\mathfrak o\oplus\Pi_{k,\Gamma}$, with
$\Pi_{k,\Gamma}=\{\sum_{\gamma<0}a_\gamma t^\gamma\}$ (well-ordered supports) and
the constant-term retraction $\operatorname{ct}:\mathfrak O_{\mathfrak
o,\Gamma}\to\mathfrak o$, a ring homomorphism because $\Pi_{k,\Gamma}$ is closed
under sums, under products and under multiplication by $\mathfrak o$, all
supports staying negative. Through \eqref{hol:eq:embedding}, allowing all of $\No$ as exponents
(each support a set) gives $\Oz=\Z\oplus\Pi$ and $\Oz[i]$, with the same
retraction \cite{repoodg,lmfactor}; these are not new.

\begin{proposition}[An elementary coefficient restriction; source~15]
\label{hol:pc:prop:ompoly}
A strongly entire series with all coefficients in $\mathfrak O_{\mathfrak
o,\Gamma}$ is a polynomial, with no equation.
\end{proposition}

\begin{proof}
Nonzero coefficients have value at most $0$; infinitely many would violate
Proposition~\ref{hol:cf:prop:criterion} at $\gamma=\beta=0$.
\end{proof}

So for omnific coefficients the analytic content of
Theorem~\ref{hol:pc:main:composition} is automatic; what an equation adds is
its degree bound and the structure of its coefficient locus.

\begin{theorem}[Linear omnific solution modules; source~15]
\label{hol:pc:thm:linarith}
Consider \eqref{hol:pc:eq:linear} with all arguments, operator coefficients and
forcing in $\mathfrak o[z]$, $\mathfrak o=\Z$ or $\Z[i]$, and let $\mathfrak R$
be $\mathfrak O_{\mathfrak o,\Gamma}$, or $\Oz$, or $\Oz[i]$, accordingly. An
effective procedure on the data either certifies that there is no polynomial
solution in $\mathfrak R[z]$ or gives $f_0,h_1,\ldots,h_r\in\mathfrak o[z]$
such that the solutions are exactly $f_0+\sum_iu_ih_i$, $u_i\in\mathfrak R$,
with unique parameters. Existence over $\mathfrak R$ is equivalent to
existence over $\mathfrak o$, and for a fixed Hahn workspace these are also
exactly its strongly entire solutions with coefficients in $\mathfrak R$.
\end{theorem}

\begin{proof}
The degree argument of Lemma~\ref{hol:pc:lem:polyleading} and
Corollary~\ref{hol:pc:cor:linearbound} is algebraic and applies to polynomials
over $\No$ and $\No[i]$; the integer roots of $I_{L_m}(d_mT)$ are computable.
Substituting a polynomial of degree $B_{\mathrm{lin}}$ gives $Mu=b$ over
$\mathfrak o$. Take a Smith normal form $UMV=S$ over the principal ideal domain
$\mathfrak o$, and put $u=Vy$, $b'=Ub$. The equations are $s_iy_i=b'_i$ and
zero rows. A solution over $\mathfrak R$ gives one over $\mathfrak o$ by $\operatorname{ct}$,
and conversely; so divisibility and zero-row tests decide existence. Since
$\mathfrak R$ is a domain, the homogeneous diagonal equations force $y_i=0$ for
$s_i\ne0$, the other $y_i$ are free, and $V$ gives the basis. The last
assertion is Proposition~\ref{hol:pc:prop:ompoly}.
\end{proof}

For $\mathfrak o=\Z$ and $\mathfrak R=\Oz$ the linear algebra is
\texttt{odg:thm:smith} and \texttt{odg:thm:linear} of the omnific Diophantine
report~\cite{repoodg}, applied to the finite system produced by the degree
bound; the Gaussian case is the same argument over $\Z[i]$. The proof uses
neither that $\Oz$ is a principal ideal domain nor any factorization theory.

\subsection{Universality and undecidability}\label{hol:pc:sub:universal}

Fix $N\ge0$. For $0\le j\le N$ put
\begin{equation}\label{hol:pc:eq:projector}
 \mathrm{Pr}_j=N!\prod_{0\le k\le N,\,k\ne j}\frac{\vartheta-k}{j-k}
 \in\Z[z]\langle D_z\rangle ;
\end{equation}
the denominator $(-1)^{N-j}j!(N-j)!$ divides $N!$, and
$\vartheta^i=\sum_jS(i,j)z^jD_z^j$ with Stirling numbers $S(i,j)$ of the second
kind. Each $\mathrm{Pr}_j$ has shift
zero and $\mathrm{Pr}_jf=N!u_jz^j$ for $f=\sum_{j\le N}u_jz^j$: interpolation in
the Euler eigenvalue.

\begin{theorem}[Universal hypersurface encoding; source~15]
\label{hol:pc:thm:universal}
For every nonzero $p=\sum_{\mathbf a}p_{\mathbf a}U^{\mathbf a}\in\Z[U_0,\ldots,
U_N]$ put $\operatorname{wd}(p)=\max\{\sum_jja_j:p_{\mathbf a}\ne0\}$,
$d=\deg p+\operatorname{wd}(p)+2$ and $\mathrm{Ann}_{d,N}=\prod_{k=0}^N
(\vartheta-dk)$. The equation
\begin{equation}\label{hol:pc:eq:encoding}
 \mathrm{Ann}_{d,N}\bigl(f(z^d)\bigr)=\sum_{\mathbf a}p_{\mathbf a}
 (N!)^{\deg p-\abs{\mathbf a}}z^{\operatorname{wd}(p)-\sum_jja_j}\prod_{j=0}^N
 (\mathrm{Pr}_jf(z))^{a_j}
\end{equation}
has integer data and satisfies the hypothesis of
Theorem~\ref{hol:pc:main:composition}; over every Hahn field of the standing
type its entire solutions are exactly $f=\sum_{j\le N}u_jz^j$ with
$p(u_0,\ldots,u_N)=0$, and its bound \eqref{hol:pc:eq:B} is exactly $N$.
\end{theorem}

\begin{proof}
Here $\bar s=1$ and every lower argument is $z$, so
$\operatorname{cw}(\mathbf a)=\abs{\mathbf a}\le\deg p<d$, and the solutions are
polynomials. $\mathrm{Ann}_{d,N}$ has shift $0$ and $I(T)=\prod_k(T-dk)$, so
$\operatorname{Rsn}=\{0,\ldots,N\}$; the monomial shifts are at most
$\operatorname{wd}(p)$ and $d-\abs{\mathbf a}\ge\operatorname{wd}(p)+2$, so every
fraction in \eqref{hol:pc:eq:B} has floor $0$ and $B_{\mathrm{cmp}}=N$. For $f$ of
degree at most $N$ the left side vanishes term by term and the right side is
$(N!)^{\deg p}z^{\operatorname{wd}(p)}p(u_0,\ldots,u_N)$, zero exactly when
$p(u)=0$ in characteristic zero.
\end{proof}

The margin $d=\deg p+\operatorname{wd}(p)+2$, rather than merely
$d>\max(\deg p,\operatorname{wd}(p))$, is what makes the bound exactly $N$ in
every case.

\begin{remark}[What universality says; source~15]\label{hol:pc:rem:universal}
After clearing denominators every affine hypersurface over $\Q$ arises. A finite
degree certificate therefore implies neither finitely many solutions, nor a
smooth solution set, nor a simple classification: the coefficient loci carry
the full complexity of affine algebraic geometry.
\end{remark}

\begin{corollary}[Omnific undecidability inside the rigid class; source~15]
\label{hol:pc:cor:undecidable}
No algorithm decides, from the integer data of an equation of
Theorem~\ref{hol:pc:main:composition}, whether it has a polynomial solution
with coefficients in $\Oz$; nor, in any fixed real Hahn workspace, whether it
has a strongly entire solution with coefficients in $\mathfrak O_{\Z,\Gamma}$.
\end{corollary}

\begin{proof}
By Theorem~\ref{hol:pc:thm:universal} and the retraction $\operatorname{ct}$, the equation
for $p$ has such a solution exactly when $p$ has a zero in $\Z^{N+1}$; the
Davis--Putnam--Robinson--Matiyasevich theorem~\cite{matiyasevich} makes that
undecidable.
\end{proof}

Undecidability of polynomial equations over $\Oz$ itself is
\texttt{odg:cor:H10} of the omnific Diophantine report~\cite{repoodg}; what is
new is the encoding into a rigid functional equation with a prescribed degree
bound.

\begin{corollary}[A fixed degree bound; source~15]\label{hol:pc:cor:ten}
Using Sun's undecidability theorem in eleven integer unknowns~\cite{sun}, the
undecidability persists for an effectively given subclass in which every
entire solution has degree at most ten, the top operator has order eleven and
every lower operator order at most ten.
\end{corollary}

\begin{proof}
Take $N=10$ in Theorem~\ref{hol:pc:thm:universal}; $d$ varies with $p$, but the
resonances and orders do not.
\end{proof}

This applies known Diophantine undecidability; it claims no optimal variable
count. The linear ordinary-data subclass is decidable
(Theorem~\ref{hol:pc:thm:linarith}), while the nonlinear subclass is already
undecidable at a fixed degree bound.

\subsection{Worked equations}\label{hol:pc:sub:examples}

\paragraph{Every degree attains the linear bound (source~15).}
For $d\ge2$ and $N\ge0$, $f(z^d)=z^{(d-1)N}f(z)$ has $B_{\mathrm{lin}}=N$;
comparing largest and smallest degrees gives $\deg f=\operatorname{ord}_zf=N$,
so the entire solutions are exactly $cz^N$.

\paragraph{A free omnific parameter (source~15).}
$f(z^2)-f(z)=z^2-z$ has bound one and entire solutions $z+c$, $c\in\KK$; its
omnific solutions are $z+c$, $c\in\Oz$, the simplest module of
Theorem~\ref{hol:pc:thm:linarith}.

\paragraph{Resonances that cannot be discarded (source~15).}
$(\vartheta-dN)f(z^d)=0$ has the solutions $cz^N$; a bound ignoring $I_L$ would
leave only constants. For
\[
 (zD_z-6)f(z^2+z)+(z^2D_z+1)f(z)=-3z^5-3z^4-2z^3-6z^2-8z-5
\]
the top resonance is $N=3$ while the nonresonant estimate is two; for
$f=a_0+\cdots+a_3z^3$ the left side is $-3a_3z^5+(-2a_2-3a_3)z^4+(-4a_2-2a_3)z^3
+(-3a_1-3a_2)z^2-4a_1z-5a_0$, and the unique solution is $f=z^3+2z+1$.

\paragraph{Nonlinear weighted equations (source~15).}
In $f(z^2)^2=f(z)^3+(D_zf(z))^2+z^4-z^3-z^2-3z-1$ the top weight is four, the
right monomials have weights three and two with shifts $0$ and $-2$, and the
forcing has weight $0$ and degree four, so $B_{\mathrm{cmp}}=1$; $f=az+b$ gives
$a^2=1$, $a^3=1$, $2ab=3a^2b-1$, hence the unique solution $f=z+1$. Likewise
$f(z^3)=f(z)^2+z^6-z^4-z^3-3z^2-2z$ has bound two and unique solution
$z^2+z+1$. The total degree three of the right side of the first equation
exceeds the two factors on the left; the weight, not the degree, is what
matters. These solutions were rechecked with SymPy for this merge.

\paragraph{A Diophantine encoder (source~15).}
For $N=1$ and $p=U_0^2+U_1^2-m$, $m\in\Z$: $\deg p=\operatorname{wd}(p)=2$,
$d=6$, $\mathrm{Pr}_0=1-\vartheta$, $\mathrm{Pr}_1=\vartheta$, and the equation is
$\vartheta(\vartheta-6)f(z^6)=z^2((1-\vartheta)f)^2+(\vartheta f)^2-mz^2$. Every
entire solution is $a+bz$ with $a^2+b^2=m$. For $m=-1$ there is no surreal
solution; for $m=1$ there are surreal and surcomplex families, and the real
omnific solutions are exactly $\pm1$ and $\pm z$, since $a^2+b^2=1$ forces $a,b$
finite, hence ordinary integers (an elementary observation not needed for
universality).

\subsection{Boundaries}\label{hol:pc:sub:boundaries}

\paragraph{Degree-one arguments are different (source~15).}
Over $\Gamma=\R$ with $p=t^\gamma$, $\gamma>0$, the partial theta series
$\Theta_p$ of \eqref{hol:eq:theta} is strongly entire and satisfies
$\Theta_p(z)=1+z\Theta_p(pz)$: both arguments have degree one, and
Corollary~\ref{hol:pc:cor:linear} does not apply. This is the report's own
witness (Theorem~\ref{hol:thm:thetadomain}), used here only to delimit the
theorem. Equal degrees also allow cancellation at higher degree: for a
nonpolynomial strongly entire even series $f_{\mathrm{ev}}$, for instance
$\sum_nt^{n^2}z^{2n}$ over $\R((t^\R))$, one has
$f_{\mathrm{ev}}(z^2)-f_{\mathrm{ev}}(-z^2)=0$.

\paragraph{Characteristic zero (source~15).}
With the prime $p$ and the series $U$ of Example~\ref{hol:nl:ex:charp},
also $D_z(U(z^{p+1}))=0$: infinitely
many indices share a residue, and Proposition~\ref{hol:pc:prop:operatorexact}
fails.

\paragraph{Rank, cofinality and enlargement (source~15).}
Over $\bigoplus_{n\ge1}\R e_n$ ($0<e_1\ll e_2\ll\cdots$), with no order unit,
$G=\sum_nt^{e_n}z^n$ is strongly entire (as $F$ in
Theorem~\ref{hol:thm:explicit}), so the theorem has
nonpolynomial inputs to exclude at infinite rank; with uncountable
cofinality every entire series is a polynomial and only the degree bounds and
encodings remain informative. Adjoining $h>ne_j$ for all $n,j$ destroys the
entireness of $G$ (Section~\ref{hol:sub:notallsurcomplex}), which does not
contradict Proposition~\ref{hol:pc:prop:extension}, and no nonpolynomial series
is strongly evaluable on all of $\No$ (Proposition~\ref{hol:ot:prop:fullclass}).

\begin{example}[A critical-weight equation; source~15]\label{hol:pc:ex:critical}
For $d\ge2$ the formal solutions over a field of characteristic zero of
$f(z^d)=f(z)^d$, a weight-critical case ($\operatorname{cw}=\bar sd=d$), are
exactly $0$ and $cz^N$ with $N\ge0$ and $c^{d-1}=1$. If $cz^N$ is the lowest
term, degree $dN$ gives $c=c^d$; writing $f=cz^NU$ with $U(0)=1$, $U(z^d)=U(z)^d$,
and a first nonconstant term $bz^r$ of $U$ would give $0$ on the left at degree
$r$ and $db\ne0$ on the right. This elementary classification is the
function-variable twin of \texttt{adr:lem:monomial} of the autonomous-dilation
report~\cite{repoadr}, where the exponent dilation $S_d$ replaces $z\mapsto
z^d$.
\end{example}

So the strict inequality is not necessary for every equation. It is
sufficient only: the case $\operatorname{cw}=\bar sd$ may still be rigid, and
several top-weight terms may carry further obstructions; no nonpolynomial
characteristic-zero counterexample at the boundary is asserted. Algebraic
differential equations without pullbacks, meromorphic solutions and several
variables are separate questions.

\paragraph{The formalization order proposed by source~15.}
Restricted-series tails; finite initial polynomials; integer convexity;
polynomial pullback; the operator residue profile; weighted domination; the
polynomial degree bound; Euler projectors; arithmetic reduction. The infinite
support argument and the finite algebraic certificate should stay separate,
and no Python check replaces the former. For linear equations with rational
or Gaussian rational data source~15 describes a certificate: the arguments and
operator coefficients, the degrees, shifts and $I_{L_j}$, the full integer-root
set of $I_{L_m}(d_mT)$ and $B_{\mathrm{lin}}$, and the coefficient matrix with a
row-reduction certificate, or, for integer or Gaussian-integer solutions, a
unimodular Smith certificate, since a rational nullspace basis need not be a
lattice basis; for non-effective data this is a certificate format, not an
algorithm.

\subsection{Predecessors, repository comparison and priority}
\label{hol:pc:sub:predecessors}

\paragraph{Literature.}
Bell, Coons and Rowland prove that a complex entire Mahler function is a
polynomial~\cite[Lemma~6]{bcr}; source~15 inspected the page with that lemma,
and does not claim the complex-analytic statement. Nishioka's
monograph~\cite{nishioka} is the standard reference for Mahler's method
(publisher metadata and scope only). The Hahn and omnific background is
L'Innocente--Mantova~\cite{lmfactor}; the undecidability inputs are
Matiyasevich~\cite{matiyasevich} (metadata and abstract) and Sun~\cite{sun}
(its abstract, eleven unknowns). Source~15 also looked at the abstract of Bell
and Coons's transcendence tests for Mahler functions and the Wikipedia
article on surreal numbers, for orientation only. Its searches were not
exhaustive. Two relevant references are missing from source~15 and are added
by this merge: Faverjon and Roques treat linear Mahler equations with
Hahn-series solutions algorithmically~\cite{faverjonroques}, as the
single-dilation report records~\cite{reposdhs}; and the autonomous-dilation
report~\cite{repoadr} records, through Nishioka--Nishioka, Mahler's 1983
criterion for convergent power-series solutions of $P(f(z),f(z^d))=0$ with
constant coefficients. Neither concerns strong entireness over an
arbitrary-rank Hahn field, and neither is used here; their content was not
rechecked by this merge.

\paragraph{Repository comparison at the pin.}
Source~15 read seven files at \texttt{bcac55a}, listed in its audit, not the
91-page text of this report as it then stood, and compared no report theorem by
theorem among the 57 then catalogued. Its keyword-search claim is corrected in
Section~\ref{hol:pc:sub:conventions}. Its statements about this report were
accurate at the pin. The universal quotient theorem of the quotient report and
the curve and lattice classifications of the omnific
reports~\cite{repoosq,repoodg} are not reproduced. The autonomous-dilation
report, which now discusses autonomous Mahler equations for the exponent
dilation, was not yet in the repository at the pin.

\paragraph{Priority.}
Not claimed: the complex entire-Mahler lemma, Mahler's method, normal forms and
the constant-term retraction, Smith normal forms, MRDP and Sun's theorem, the
support calculus and Gauss data of this report, and the elementary critical
classification. The proposed contributions are strict composition-weight
rigidity over arbitrary-rank Hahn fields, its exact-resonance degree bound and
extension-stable coefficient loci, the matrix version, the linear omnific
module theorem as a consequence, and the universal Euler-projector encoding
with its undecidability consequences; priority is provisional, and the word
``breakthrough'', which the source considered, is expressly disclaimed.

\subsection{What the polynomial-composition part does not claim}
\label{hol:pc:sub:nonclaims}

Every limitation stated by source~15, in its manuscript, audit file or
verification record, is kept.
\begin{enumerate}[label=\textup{(P\arabic*)}]
\item No independent referee report, no Lean formalization, no repository
build (a container clone failed), and no certified priority; the label
``breakthrough'' is disclaimed.
\item Bell--Coons--Rowland's Lemma~6, Nishioka, the MRDP theorem and Sun's
theorem are credited inputs; the variable count is not optimal, and the bound
ten is inherited.
\item Entireness is relative to one set-sized Hahn field; nothing is claimed
about unrestricted surreal analytic functions or entireness on $\No$.
\item $D_z$ is the function-variable derivative, not the Berarducci--Mantova
derivation.
\item The general higher-order algebraic differential rigidity problem is not
addressed (in this report it is Theorems~\ref{hol:td:main:boundary},
\ref{hol:ot:main:second} and~\ref{hol:nr:main:third}).
\item Polynomiality of entire series with omnific coefficients is elementary
(Proposition~\ref{hol:pc:prop:ompoly}).
\item The condition $\operatorname{cw}<\bar sd$ is sufficient, not necessary; the
critical case $\operatorname{cw}=\bar sd$ and several top products are open.
\item Corollary~\ref{hol:pc:cor:newton} is not an independence statement.
\item The finite coefficient locus is not a decision procedure, and no
algorithm is given for arbitrary named surreal coefficients.
\item The matrix theorem needs a nonsingular dominant matrix and does not treat
matrix differential operators.
\item Several variables and meromorphic solutions are not treated.
\item Proposition~\ref{hol:pc:prop:extension} concerns fixed equations only.
\item $\Oz$ is not claimed to be a principal ideal domain.
\item The classification of $f(z^d)=f(z)^d$ is elementary.
\item The 2{,}369 finite assertions do not cover infinite supports, all ranks or
universal statements, and do not certify novelty.
\item The repository comparison was limited to seven files, and the 57-report
catalogue was not compared theorem by theorem; its statement that the
repository contained no Mahler material was incorrect at the pin.
\item The universal quotient theorem and the curve and lattice results of the
omnific reports are not reproduced.
\item \emph{[merge]} Remark~\ref{hol:pc:rem:dichotomy}, the identification of
source~15's indicial polynomials with $I_P$ and $\chi_P$, and the added
Mahler references are the merge's; no historical priority is claimed for
them.
\end{enumerate}

%======================================================================
% Recurrence part: source 20. Labels use hol:rc:.
%======================================================================
\section{Hahn-valued recurrences: observables and Noetherian compression}\label{hol:rc:sec:setting}

This section and the next contain the material of source~20, delivered in
batch~33 after the fourteen principal theorems above had been merged. Its
subject is the \emph{ordinary} constant-coefficient recurrence
$y_{n+d}=a_{d-1}y_{n+d-1}+\cdots+a_0y_n$ whose coefficients and values lie in a
Hahn field, or in $\No[i]$; its index $n$ is always an ordinary natural number.
A single surreal term can carry infinitely many real coefficients, and
source~20 separates two kinds of information in such a sequence: the leading
scale, which is eventually affine on finitely many progressions, and a
prescribed lower Hahn or Conway coefficient, which is tame when the
characteristic roots are units and can be arbitrary otherwise.

Much of source~20's valuation theory was found independently of source~18 and
is printed once, in Sections~\ref{hol:ud:sec:setting}--\ref{hol:ud:sec:mixed}.
What is new here is the \emph{coefficient} side. All coefficient sequences of
a unit-root exponential polynomial lie in one finitely generated algebra of
sequences (Theorem~\ref{hol:rc:thm:algebra}); hence every set-sized family of
polynomial conditions on finitely many coefficients at a time reduces, along
the orbit, to a finite subfamily, and its hitting set consists of full
progressions modulo one period and a finite set
(Theorem~\ref{hol:rc:thm:compression}). In the opposite direction one fixed
Conway coefficient of a first-order omnific recurrence can reproduce any real
sequence (Theorem~\ref{hol:rc:thm:coding}). The recurrence statements over
$\No[i]$, $\Oz$ and $\Oz[i]$, an omnific recurrence whose degree decreases
forever, and determinantal profiles of matrix powers complete the part.
Source~20's Sections~9--10 on mixed operators re-derive special cases of
Theorem~\ref{hol:ud:main:mixed} and are recorded in
Section~\ref{hol:rc:sub:mixed}.

\subsection{Source, conventions and the merge's choices}
\label{hol:rc:sub:conventions}

\begin{convention}[Setting of the recurrence part]
\label{hol:rc:conv:field}
In Sections~\ref{hol:rc:sec:setting}--\ref{hol:rc:sec:surreal}, $k$ is an
arbitrary field of characteristic zero with the trivial valuation, $\Gamma$ is
as in Convention~\ref{hol:conv:group}, $\KK=k((t^\Gamma))$, and $v$, $\lc$,
$\res$, strong summability, $\EK$, $D_z$ and $\sigma_\lambda$ are as in
Convention~\ref{hol:ud:conv:field}: $D_z$ kills $\KK$ and is not the
Berarducci--Mantova derivation, and entireness is relative to this one
set-sized field. Sequences are indexed by ordinary $n\in\N$. A unit (an
element of valuation zero) is written $\lambda=\zeta(1+\tau)$ as in
\eqref{hol:ud:eq:unitdecomp}. An \emph{exponential polynomial} is
$\mathcal Y(n)=\sum_{j=1}^mR_j(n)\lambda_j^n$ with $R_j\in\KK[U]$ and
$\lambda_j\in\KK^\times$, as in Section~\ref{hol:ud:sub:sml}; it is
\emph{unit-root} if every $\lambda_j$ is a unit. For units
$\lambda_1,\ldots,\lambda_m$ with residues $\zeta_1,\ldots,\zeta_m$ put
$\langle\zeta\rangle=\langle\zeta_1,\ldots,\zeta_m\rangle\le k^\times$, a
finitely generated abelian group, whose torsion subgroup is therefore finite,
and let $\mathsf M_{\mathrm{tor}}$ be the exponent of that torsion subgroup
($\mathsf M_{\mathrm{tor}}=1$ when it is trivial).
\end{convention}

For $k=\C$ one has $\KK=\K$ and $\EK=\E$. The statements of
Section~\ref{hol:rc:sec:setting} hold verbatim for $\Gamma=0$, where $\KK=k$;
Convention~\ref{hol:conv:group} matters only where scales or entireness enter.

\paragraph{The source.}
Source~20 is \emph{Valuation Rigidity and Coefficient Universality in Surreal
Recurrences: A mixed differential--dilation theorem and the limits of
coefficientwise predictability} (22 pages, US letter, 23 September 2026),
delivered as the archive \texttt{surreal\_recurrences\_research} (batch~33,
manuscript~01; local number~20 in this directory). It consulted the repository
at
\begin{center}\small\texttt{efc5446229dbf61a97bdd5edae91dba20af9d131},\end{center}
where this report's \texttt{article.tex} is the blob
\texttt{9fd5bd460a19448b9078a4ef7f012e647f545e59}; its shipped file
\path{20-recurrences-SOURCE_AUDIT.md} lists the line ranges it read. At that
pin this report was the seven-source text (sources~08--14, 137 pages).
Sources~15--19 had been placed in this directory, with their audit and code
files, in \texttt{9d28e28}, an ancestor of the pin, but none of them had yet
been merged into the text; source~18 was merged later the same day, in
\texttt{2c4debb}. Source~20's description of this report was accurate at the
pin: the paragraph after Theorem~\ref{hol:thm:mixed} did leave the mixed
unit-valuation case unclassified when $\Gamma$ has an order unit, and pointed to
Question~\ref{hol:q:mixedunit}. That statement is \emph{stale} now:
source~18 answers Question~\ref{hol:q:mixedunit} positively for every
$\Gamma$ (Corollary~\ref{hol:ud:cor:unitrigidity}) and proves the
relative-scale criterion of Theorem~\ref{hol:ud:main:mixed}, which contains
source~20's Theorems~10.1, 10.3 and~10.4 (Section~\ref{hol:rc:sub:mixed}).
So source~20's description of its Section~10 as closing ``the repository gap'',
its Corollary~10.2 (``the case explicitly left open''), and its Remark~10.5 are
printed here as an independent second derivation, not as current
contributions of this report. Source~20 claims neither the pure differential
or dilation theorems, the escape mechanism, the order-unit definition, nor the
partial theta series; it credits them to this report.

\paragraph{Renamed symbols.}
The following symbols of source~20 would collide with symbols of the earlier
parts. They are renamed here, in the new material only.
{\small
\begin{longtable}{@{}>{\raggedright\arraybackslash}p{0.27\textwidth}
>{\raggedright\arraybackslash}p{0.2\textwidth}
>{\raggedright\arraybackslash}p{0.45\textwidth}@{}}
\toprule
Source~20 & Here & Reason\\
\midrule
\endhead
$K=k((t^\Gamma))$, $\mathcal E(K)$ & $\KK$, $\EK$ & $\K$ is $\C((t^\Gamma))$
and $\E$ its entire series in this report.\\
unit roots $q_i=c_i(1+\varepsilon_i)$; $u_n=\sum_iP_i(n)q_i^n$ &
$\lambda_j=\zeta_j(1+\tau_j)$; $\mathcal Y(n)=\sum_jR_j(n)\lambda_j^n$ & $q$ is
a dilation parameter; the unit-dilation notation of
Convention~\ref{hol:ud:conv:field} is kept.\\
envelope $W=A+M(S)$; $R_{i,\gamma}$ & $\Sigma=C+M(S)$; $r_{j,\gamma}$ & As in
Lemma~\ref{hol:ud:lem:envelope}; $R_j$ are the polynomial factors.\\
$C=\langle c_i\rangle$; period $M=\exp\operatorname{Tor}(C)$; class $r$ &
$\langle\zeta\rangle$; $\mathsf M_{\mathrm{tor}}$; $b$ & $\C$ is the complex
field; $\mathsf M$ is source~18's smaller period
(Lemma~\ref{hol:ud:lem:period}, Remark~\ref{hol:rc:rem:periods}); $M(S)$ is
the positive monoid; $r$ is a derivative order.\\
$\mathcal A=k[n,c_i^n]$; $R_0=k[T,X_1,\ldots,X_s]$ & $\mathfrak A$;
$k[U,X_1,\ldots,X_m]$ & $\mathcal A[f]$ is a differential expression of the
theta-descent part; $U$ is the polynomial variable of $k[U]$.\\
forward shift $E$ & $\mathrm{Sh}$ & As in Section~\ref{hol:ud:sub:sml}.\\
recurrence terms $u_n$ & $y_n$ & $u$ is the Laurent variable of
Section~\ref{hol:td:sub:laurent}; $y_n$ is otherwise local to the proof of
Theorem~\ref{hol:thm:unitorbit}.\\
vector $x\in K^d$, set $X$, $T\subseteq\Gamma$, $F\in k[X_1,\ldots,X_h]$ &
$\mathbf x\in\KK^{\bar d}$, $\mathcal X$, $\mathsf T$, $F\in
k[Z_1,\ldots,Z_h]$ & $x$ is an evaluation point, $d$ a degree, $T$ an
indeterminate.\\
purely infinite ideal $\mathcal J$ & $\Pi$ & $\Oz=\Z\oplus\Pi$ as in
Corollary~\ref{hol:td:cor:omnific}; $\jet\lambda j$ is a jet.\\
sequence $b$; $A_b=\sum b_m\omega^{\omega-m}$; set $S\subseteq\N$ &
$\mathbf b=(\mathbf b_m)$; $\mathfrak w_{\mathbf b}$; $A$ & $b$ is a residue
class; $S$ is a support alphabet.\\
Catalan root $\lambda$; $C(s)$, $C_k$ & $\lambda_-$; $\mathrm{Cat}(X)$,
$\mathrm{Cat}_k$ & $C_n$ are the Chebyshev polynomials of
\eqref{hol:td:eq:cheb}.\\
matrix $A$; $\delta_j(n)$ & $\mathbf A$; $\operatorname{mv}_j(n)$ & $\delta$
is a scale.\\
dilation group $G$; convex hull $H$ of $v(G)$ & $\mathfrak G$;
$\Delta_{\mathfrak G}$ & $G$ is the series of
Corollary~\ref{hol:nl:cor:theta}; convex subgroups are $\Delta$ here.\\
operator $\sum c_{\lambda,j,k}z^kD_z^j\sigma_\lambda$; $A_s(n)$, $B_j(n)$ &
$\sum\varkappa_{\ell,r,i}z^iD_z^r\sigma_{\lambda_\ell}$; $c_j(n)$ & As in
\eqref{hol:ud:eq:L} and \eqref{hol:ud:eq:cj}.\\
witness $\eta\notin H$; $\eta$ in examples & $\gamma$; $\epsilon$ &
$\eta=v(\tau)$ in Theorem~\ref{hol:thm:unitorbit}.\\
$\Theta_Q$; characteristic $p$ & $\Theta_p$; $\mathfrak p$ &
\eqref{hol:eq:theta}.\\
Theorems~4.1, 5.1, 5.5, 6.2, 7.1, 7.3, 10.1, 10.3, 10.4 & see below &
Theorems~A--O of this report are different theorems.\\
\bottomrule
\end{longtable}}
The shipped audit file \path{20-recurrences-SOURCE_AUDIT.md} keeps the
source's notation.

\paragraph{Words with two meanings.}
A \emph{unit} is a valuation-zero element (a ``valuation unit'' in source~20);
it is neither an \emph{order unit} (Definition~\ref{hol:def:orderunit}) nor a
unit of the ring $\Oz$, whose only units are $\pm1$. \emph{Coefficient} means a
Hahn coefficient $[t^\gamma]x$, a Conway coefficient $[\omega^\alpha]x$ (the
same number as $[t^{-\alpha}]x$ under \eqref{hol:eq:embedding}), a Taylor
coefficient $a_n$, or a recurrence coefficient $c_j(n)$; a
\emph{coefficient observable} is a sequence $n\mapsto[t^\gamma]y_n$ for one
fixed $\gamma$. \emph{Compression} here means \emph{Noetherian compression}, a
finite subfamily of coefficient conditions equivalent to the whole family
along an orbit (Theorem~\ref{hol:rc:thm:compression}); it is unrelated to the
bounded polynomial \emph{jet} compression of
Theorem~\ref{hol:th:thm:compression}, a statement of differential algebra.
Two \emph{periods} occur: source~18's $\mathsf M$ suffices for valuations,
while the hitting sets of Theorem~\ref{hol:rc:thm:compression} need
$\mathsf M_{\mathrm{tor}}$, which $\mathsf M$ divides
(Remark~\ref{hol:rc:rem:periods}); the tempting reading ``the valuation period
is also the period of every coefficient test'' is false. Finally
$\deg_\omega$ is the greatest exponent of a Conway normal form, neither the
degree of a polynomial nor a birthday.

\paragraph{Results printed once.}
Source~20 reproves several facts the report already contains. They are
printed once and cited. Its definitions of strong summability, of
$\mathcal E(K)$, of valuation units and of order units are
Definitions~\ref{hol:def:entire} and~\ref{hol:def:orderunit} and
Convention~\ref{hol:ud:conv:field}; its surreal translation
$t^\gamma\mapsto\omega^{-\gamma}$ is \eqref{hol:eq:embedding}, and its
localization of finitely many surreal data to a set-sized Hahn field is the
paragraph on omnific coefficients in Section~\ref{hol:ud:sub:surreal}. Its
Lemma~3.1 (support lemma) is Lemma~\ref{hol:lem:support}; its Lemma~3.2 and
Corollary~3.4 are Lemma~\ref{hol:cf:lem:exppoly} and its proof; its
Theorem~3.3, the classical Skolem--Mahler--Lech theorem, is imported there as
here. Its Theorem~4.1\,(a)--(b) is Lemma~\ref{hol:ud:lem:envelope}, proved by
the same finite binomial expansion and support count; its Remark~4.2 is the
remark after that lemma, and its Example~4.3,
$v\bigl((1+t^\epsilon)^n-1-nt^\epsilon\bigr)=2\epsilon$ for $n\ge2$ with
$y_0=y_1=0$, is the case $d=2$ of Example~\ref{hol:ud:ex:cancellation}. Its
Theorem~5.1 is Theorem~\ref{hol:ud:thm:unitvalues} with the admissible period
$\mathsf M_{\mathrm{tor}}$ in place of $\mathsf M$
(Remark~\ref{hol:rc:rem:periods}), and its Corollary~5.2 is the last clause of
that theorem together with its finitely many exceptional indices; its
Example~5.3 is Example~\ref{hol:ex:unitperiod}; its Lemma~5.4 is
Lemma~\ref{hol:ud:lem:lines}; its Theorem~5.5 is
Theorem~\ref{hol:ud:thm:affine}, which states the zero alternative for every
index of the class where source~20 states it beyond a threshold; its
Remark~5.6 is the third example of Remark~\ref{hol:rem:shortcut}. Its
Lemma~9.1 is Lemma~\ref{hol:ud:lem:operator} with
Corollary~\ref{hol:ud:cor:faithful}, and its Lemma~9.2 is
Lemma~\ref{hol:lem:escape} with Corollary~\ref{hol:cor:boundedcost}, with the
same closed exterior region $v(x)\le-B$; its Theorem~10.1 is
Theorem~\ref{hol:ud:thm:rigidity} for relative scale~$0$, and its
Corollary~10.2 is Corollary~\ref{hol:ud:cor:unitrigidity}. The entireness of
its theta witness is Theorem~\ref{hol:thm:thetadomain}, by the route of
Remark~\ref{hol:rem:thetasecond}, and its identities are
\eqref{hol:eq:thetafirst} and \eqref{hol:eq:thetahom}. Its torsion example
(Section~11.1) is the case $h=2$ of the paragraph on root-of-unity quotients in
Section~\ref{hol:ud:sub:examples} with Proposition~\ref{hol:prop:torsion}, and
its rank-two example (Section~11.2) is Example~\ref{hol:ud:ex:crossing} with
Example~\ref{hol:ex:ranktwo}.

\paragraph{Where each numbered result went.}
Theorem~4.1\,(c) is Theorem~\ref{hol:rc:thm:algebra}; the period paragraph
\eqref{hol:rc:eq:torperiod} of its Section~3 is Lemma~\ref{hol:rc:lem:torperiod};
Corollary~5.7 is Corollary~\ref{hol:rc:cor:split}; Definition~6.1,
Theorem~6.2, Corollary~6.3, Example~6.4 and Remark~6.5 are
Definition~\ref{hol:rc:def:conditions}, Theorem~\ref{hol:rc:thm:compression},
Corollary~\ref{hol:rc:cor:tails}, Example~\ref{hol:rc:ex:transient} and
Remark~\ref{hol:rc:rem:safeguards}; Theorem~7.1, Example~7.2, Theorem~7.3,
Corollary~7.4 and Remark~7.5 are Theorem~\ref{hol:rc:thm:surreal},
Example~\ref{hol:rc:ex:catalan}, Theorem~\ref{hol:rc:thm:coding},
Corollary~\ref{hol:rc:cor:nosml} and Remark~\ref{hol:rc:rem:quotients};
Corollary~8.1 is Corollary~\ref{hol:rc:cor:matrix}; Theorems~10.3 and~10.4
are Corollaries~\ref{hol:rc:cor:convex} and~\ref{hol:rc:cor:finitegroup};
Remark~10.5 is Remark~\ref{hol:rc:rem:stale}; Section~11 is
Section~\ref{hol:rc:sub:boundaries}, with its positive-characteristic formula
as Example~\ref{hol:rc:ex:charp}; Questions~12.1, 12.2, 12.5 and~12.11 are
recorded in the statuses of Questions~\ref{hol:ud:q:effective},
\ref{hol:ud:q:torsionblocks}, \ref{hol:ud:q:beyond}
and~\ref{hol:ud:q:formal}, and Questions~12.3, 12.4, 12.6--12.10 and~12.12 are
Questions~\ref{hol:rc:q:nonunit}--\ref{hol:rc:q:substitutions}, Question~12.7
re-scoped; the appendix and the audit file are merged into
Sections~\ref{hol:rc:sub:predecessors}, \ref{hol:rc:sub:nonclaims}
and~\ref{hol:sec:verification}. Theorem~\ref{hol:rc:main:observables} in the
introduction collects Theorems~\ref{hol:rc:thm:algebra},
\ref{hol:rc:thm:compression} and~\ref{hol:rc:thm:coding}.

\paragraph{The merge's choices.}
The review of source~20 for this merge found its proofs correct; its worked
examples were rechecked by hand and by its own suite. The merge adds five
points, each marked. First, Remark~\ref{hol:rc:rem:periods} compares the two
periods and shows that source~18's cannot replace source~20's for hitting
sets. Second, Section~\ref{hol:rc:sub:mixed} identifies source~20's
Theorems~10.3 and~10.4 as special cases of
Theorems~\ref{hol:ud:thm:rigidity} and~\ref{hol:ud:thm:dichotomy}, notes
through Example~\ref{hol:ud:ex:commonscale} that the inclusion is strict, and
shows in Remark~\ref{hol:rc:rem:finitegeneration} that finite generation cannot
be dropped from the convex-hull form of the criterion. Third,
Remark~\ref{hol:rc:rem:stream} credits the omnific-notations report, which
already contains the element of the coding theorem. Fourth,
Example~\ref{hol:rc:ex:charp} keeps source~20's exact positive-characteristic
valuation formula, of which Section~\ref{hol:ud:sub:examples} records a
subsequence. Fifth, the statuses of the related questions of source~18 are
updated and source~20's Question~12.7 is re-scoped
(Question~\ref{hol:rc:q:nonlinear}). Source~20's theorems, proofs, examples
and non-claims are otherwise printed as delivered, in the renamed notation.

\subsection{One finitely generated algebra of coefficient observables}
\label{hol:rc:sub:algebra}

The support envelope and the coefficient polynomials of
Lemma~\ref{hol:ud:lem:envelope} give more than a valuation theorem: they say
where every coefficient sequence lives.

\begin{theorem}[One finitely generated observable algebra; source~20]
\label{hol:rc:thm:algebra}
Let $\mathcal Y^{(1)},\ldots,\mathcal Y^{(\bar d)}$ be unit-root exponential
polynomials over $\KK$, let $\lambda_1,\ldots,\lambda_m$ list all their bases
and $\zeta_1,\ldots,\zeta_m$ the residues. There is one well-ordered
$\Sigma\subseteq\Gamma$ containing $\supp\mathcal Y^{(i)}(n)$ for every $i$
and every $n\ge0$. For every $i$ and every $\gamma\in\Gamma$ there are
$r^{(i)}_{j,\gamma}\in k[U]$ with
\begin{equation}\label{hol:rc:eq:coefficient}
 [t^\gamma]\,\mathcal Y^{(i)}(n)=\sum_{j=1}^m\zeta_j^n\,r^{(i)}_{j,\gamma}(n)
 \qquad(n\ge0),
\end{equation}
and so every coefficient observable $n\mapsto[t^\gamma]\mathcal Y^{(i)}(n)$
lies in the one finitely generated $k$-algebra
\begin{equation}\label{hol:rc:eq:algebra}
 \mathfrak A=k[\,n,\zeta_1^n,\ldots,\zeta_m^n\,]\subseteq k^{\N},
\end{equation}
the image of the evaluation homomorphism
$\operatorname{ev}:k[U,X_1,\ldots,X_m]\to k^{\N}$, $U\mapsto(n)_n$,
$X_j\mapsto(\zeta_j^n)_n$, sequences being multiplied pointwise.
\end{theorem}

\begin{proof}
Lemma~\ref{hol:ud:lem:envelope}, applied to each $\mathcal Y^{(i)}$ with its
own bases (the others entering with zero polynomial factor), gives a
well-ordered envelope $\Sigma^{(i)}$ and the identity
\eqref{hol:rc:eq:coefficient} for $\gamma\in\Sigma^{(i)}$; for
$\gamma\notin\Sigma^{(i)}$ take $r^{(i)}_{j,\gamma}=0$. The finite union
$\Sigma=\bigcup_i\Sigma^{(i)}$ is well ordered. The identity holds for
\emph{every} $n\ge0$, small $n$ included: the polynomials $\binom Ui$ vanish at
the integers $n<i$, and characteristic zero makes them polynomials over $k$.
Each summand $\zeta_j^nr^{(i)}_{j,\gamma}(n)$ is
$\operatorname{ev}(X_jr^{(i)}_{j,\gamma}(U))$. No sum over $n$ is asserted to
be strongly summable.
\end{proof}

\begin{remark}[Finitely generated, not finite-dimensional; source~20]
\label{hol:rc:rem:degrees}
The algebra $\mathfrak A$ has $m+1$ generators, but the observables do not lie
in a finite-dimensional subspace: already $[t^d](1+t)^n=\binom nd$ has degree
$d$ in $n$, so the degrees of the $r_{j,\gamma}$ are unbounded in $\gamma$, as
recorded after Lemma~\ref{hol:ud:lem:envelope}. The theorem gives finitely
many algebra generators for all coefficient sequences at once, not a finite
basis; and it is the finite generation, not a degree bound, that the next
subsection uses.
\end{remark}

For $q=1+t^\epsilon$ with $\epsilon>0$ the sequence
$y_n=q^n-1-nt^\epsilon=\sum_{i=2}^n\binom nit^{i\epsilon}$ vanishes for
$n\le1$ and has valuation $2\epsilon$ for $n\ge2$; replacing $q$ by its residue
$1$ would lose the whole sequence. This is source~20's Example~4.3, the case
$d=2$ of Example~\ref{hol:ud:ex:cancellation}: the observables keep every
principal-unit correction while each fixed coefficient is an ordinary
exponential polynomial over $k$.

\subsection{One period for all coefficient tests, and split recurrences}
\label{hol:rc:sub:period}

Source~20's Section~3 fixes the period
\begin{equation}\label{hol:rc:eq:torperiod}
 \mathsf M_{\mathrm{tor}}=\exp\operatorname{Tor}\langle\zeta\rangle,
\end{equation}
the exponent of the finite torsion subgroup of $\langle\zeta\rangle$, and
observes that it serves for all polynomial degrees and all finite products of
the residues at once.

\begin{lemma}[One period for all monomials in the residues; source~20]
\label{hol:rc:lem:torperiod}
Let $\zeta^{\alpha_1},\ldots,\zeta^{\alpha_s}$ be monomials
$\zeta^\alpha=\prod_j\zeta_j^{\alpha_j}$, $\alpha_i\in\Z^m$, let
$\varrho_1,\ldots,\varrho_s\in k[U]$ and $b\in\N$. The sequence
\[
 l\longmapsto\sum_{i=1}^s\varrho_i(\mathsf M_{\mathrm{tor}}l+b)\,
 \bigl(\zeta^{\alpha_i}\bigr)^{\mathsf M_{\mathrm{tor}}l+b}
\]
is either identically zero or has only finitely many zeros. The same
$\mathsf M_{\mathrm{tor}}$ serves for every such list, every choice of degrees
and every $b$; the finite zero sets need not have a common bound.
\end{lemma}

\begin{proof}
The bases in $l$ are the elements $(\zeta^{\alpha_i})^{\mathsf M_{\mathrm{tor}}}$;
combine equal ones. If
$(\zeta^{\alpha_i}/\zeta^{\alpha_{i'}})^{\mathsf M_{\mathrm{tor}}}$ is a root of
unity, then $\zeta^{\alpha_i}/\zeta^{\alpha_{i'}}\in\langle\zeta\rangle$ is a
torsion element, its order divides $\mathsf M_{\mathrm{tor}}$, and the two bases
are equal. So after combination the distinct bases have no root-of-unity
quotient, and Lemma~\ref{hol:cf:lem:exppoly} gives the alternative.
\end{proof}

\begin{remark}[Two periods; merge]\label{hol:rc:rem:periods}
Source~18's period $\mathsf M$ of Lemma~\ref{hol:ud:lem:period}, the least
common multiple of the orders of those quotients $\zeta_i/\zeta_j$ that are
roots of unity, divides $\mathsf M_{\mathrm{tor}}$, since each such quotient
lies in the torsion subgroup of $\langle\zeta\rangle$. Both are admissible in
Theorem~\ref{hol:ud:thm:unitvalues}, which with $\mathsf M_{\mathrm{tor}}$ is
source~20's Theorem~5.1. For hitting sets the smaller period does not suffice.
For $m=1$ and $\mathcal Y(n)=(-1)^n$ there is no quotient, so $\mathsf M=1$,
and indeed $v(\mathcal Y(n))=0$ for every $n$; but $\mathsf M_{\mathrm{tor}}=2$,
and the single condition $[t^0]x-1=0$ holds exactly for even $n$. Valuations
see only quotients of residues, because the least coefficient that is not an
identity is a combination with those residues as bases; coefficient tests see
arbitrary monomials in them, which is why
Theorem~\ref{hol:rc:thm:compression} is stated with
$\mathsf M_{\mathrm{tor}}$.
\end{remark}

\begin{corollary}[Recurrences with split characteristic polynomial;
source~20]\label{hol:rc:cor:split}
Let $(y_n)$ satisfy a constant-coefficient recurrence
$y_{n+d}=a_{d-1}y_{n+d-1}+\cdots+a_0y_n$ over $\KK$ whose characteristic
polynomial splits in $\KK$. There are $\mathsf M\ge1$ and $N\ge0$ such that on
each class $b$ modulo $\mathsf M$ either $y_n=0$ for all $n\ge N$ in the class,
or
\[
 v(y_n)=\beta_b+n\varpi_b\qquad(n\ge N,\ n\equiv b\pmod{\mathsf M}),
\]
with $\beta_b\in\Gamma$ and $\varpi_b$ the valuation of a nonzero
characteristic root. The leading coefficient on each nonzero class is
eventually an exponential polynomial over $k$ in $(n-b)/\mathsf M$, and
$\mathsf M$ may be chosen from the residues of the normalized nonzero roots
$t^{-v(\lambda)}\lambda$ alone. A zero characteristic root affects only a
finite initial segment.
\end{corollary}

\begin{proof}
The Jordan decomposition of the companion matrix over $\KK$ gives, for all $n$
at least the multiplicity of the root $0$, a representation
$y_n=\sum_jR_j(n)\lambda_j^n$ over the distinct nonzero roots, with
$\deg R_j$ below the multiplicity of $\lambda_j$; nilpotent blocks vanish
after finitely many steps. Apply Theorem~\ref{hol:ud:thm:affine} to the
exponential polynomial defined by this representation for all $n$, which
agrees with $y_n$ on that tail, and take the larger threshold. The statement on
leading coefficients is the one in Theorem~\ref{hol:ud:thm:unitvalues}, applied
to the winning group of roots.
\end{proof}

\subsection{Noetherian compression of infinitely many coefficient tests}
\label{hol:rc:sub:compression}

\begin{definition}[Coefficient-algebraic conditions; source~20]
\label{hol:rc:def:conditions}
A \emph{coefficient-algebraic condition} on
$\mathbf x=(x^{(1)},\ldots,x^{(\bar d)})\in\KK^{\bar d}$ is an equation
\[
 F\bigl([t^{\gamma_1}]x^{(i_1)},\ldots,[t^{\gamma_h}]x^{(i_h)}\bigr)=0,
 \qquad F\in k[Z_1,\ldots,Z_h],
\]
in finitely many coefficient coordinates. A \emph{coefficient-algebraic set}
$\mathcal X\subseteq\KK^{\bar d}$ is the common zero set of an arbitrary
set-sized family $(F_a)_{a\in I}$ of such conditions.
\end{definition}

The family may be infinite, but each member is an ordinary polynomial in
finitely many ordinary coefficients. For a prescribed $\mathsf T\subseteq\Gamma$
the support condition $\supp(x)\subseteq\mathsf T$ is the family
$[t^\gamma]x=0$, $\gamma\notin\mathsf T$. The condition that an arbitrary real
coefficient be an \emph{integer} is not of this form.

\begin{theorem}[Noetherian compression and one residual period; source~20]
\label{hol:rc:thm:compression}
Let $\mathbf y_n=(\mathcal Y^{(1)}(n),\ldots,\mathcal Y^{(\bar d)}(n))$ be a
finite vector of unit-root exponential polynomials, let
$\zeta_1,\ldots,\zeta_m$ be the residues of all their bases, and let
$\mathsf M_{\mathrm{tor}}$ be as in \eqref{hol:rc:eq:torperiod}. For every
coefficient-algebraic set $\mathcal X\subseteq\KK^{\bar d}$, defined by
$(F_a)_{a\in I}$:
\begin{enumerate}[label=\textup{(\alph*)}]
\item there is a finite $I_0\subseteq I$ such that, for \emph{every} $n\ge0$,
$\mathbf y_n\in\mathcal X$ if and only if $F_a$ vanishes at $\mathbf y_n$ for
all $a\in I_0$;
\item for each class $b$ modulo $\mathsf M_{\mathrm{tor}}$, either
$\mathbf y_{\mathsf M_{\mathrm{tor}}l+b}\in\mathcal X$ for every $l\ge0$, or
this holds for only finitely many $l$;
\item consequently $\{n\ge0:\mathbf y_n\in\mathcal X\}$ is a finite union of
full residue classes modulo $\mathsf M_{\mathrm{tor}}$ and a finite set.
\end{enumerate}
The same $\mathsf M_{\mathrm{tor}}$ serves for every $\mathcal X$; the finite
subfamily $I_0$ and the exceptional set depend on $\mathcal X$.
\end{theorem}

\begin{proof}
By Theorem~\ref{hol:rc:thm:algebra} every coordinate observable is
$\operatorname{ev}$ of a polynomial in $k[U,X_1,\ldots,X_m]$. For each $a\in I$
substitute such representatives for the finitely many coordinates in $F_a$;
this gives $\widehat F_a\in k[U,X_1,\ldots,X_m]$ whose value at
$(n,\zeta_1^n,\ldots,\zeta_m^n)$ is exactly $F_a$ evaluated at $\mathbf y_n$.
The ring $k[U,X_1,\ldots,X_m]$ is Noetherian (Hilbert's basis theorem), so the
ideal generated by all $\widehat F_a$ is generated by finitely many elements,
each a finite combination of finitely many $\widehat F_a$; collect those
indices into $I_0$. A point $(n,\zeta^n)$ at which the $\widehat F_a$,
$a\in I_0$, vanish kills the whole ideal, hence every $\widehat F_a$. This is
(a), for every $n$ and not only eventually.

Each $\widehat F_a(n,\zeta_1^n,\ldots,\zeta_m^n)$ is a finite sum of terms
$\varrho(n)(\zeta^\alpha)^n$ with $\alpha\in\N^m$. On a class modulo
$\mathsf M_{\mathrm{tor}}$ it is identically zero or has finitely many zeros by
Lemma~\ref{hol:rc:lem:torperiod}. If every condition vanishes identically on
the class, every orbit point of the class lies in $\mathcal X$; otherwise a
single condition that is not an identity there already leaves only finitely
many hits. This is (b),
and (c) follows by taking the finite union over classes.
\end{proof}

\begin{corollary}[Support and tail recurrence; source~20]
\label{hol:rc:cor:tails}
For a unit-root exponential polynomial $\mathcal Y$ and any
$\mathsf T\subseteq\Gamma$, the set $\{n:\supp\mathcal Y(n)\subseteq\mathsf
T\}$ is a finite union of full classes modulo $\mathsf M_{\mathrm{tor}}$ and a
finite set; if it meets a class infinitely often, it contains the whole
class. Along the orbit, the vanishing of all coefficients outside
$\mathsf T$ is equivalent to the vanishing of finitely many of them.
\end{corollary}

\begin{proof}
Apply Theorem~\ref{hol:rc:thm:compression} to the family $[t^\gamma]x=0$,
$\gamma\notin\mathsf T$; part~(a) chooses a finite subset of exactly these
equations.
\end{proof}

Under \eqref{hol:eq:embedding}, for a real-surreal unit-root sequence, the
absence of an infinitesimal tail ($\supp\subseteq\Gamma_{\le0}$), membership
in the purely infinite support space ($\supp\subseteq\Gamma_{<0}$) and the
vanishing of a prescribed truncation are instances of
Corollary~\ref{hol:rc:cor:tails}. Full membership in $\Oz$ also requires the
constant coefficient to be an ordinary integer, an arithmetic condition on an
exponential polynomial that is not a Zariski-closed coefficient condition and
is not disposed of here (Question~\ref{hol:rc:q:membership}).

\begin{example}[No common transient bound; source~20]
\label{hol:rc:ex:transient}
For $\mathcal Y(n)=(1+t)^n$ the single condition $[t^d]x=0$ holds exactly for
$n=0,\ldots,d-1$. The period is $\mathsf M_{\mathrm{tor}}=1$ for every $d$, but
the last exceptional index tends to infinity with $d$. A universal residual
period does not give a uniform transient length for all coefficient tests on
one orbit.
\end{example}

\begin{remark}[Why infinitely many conditions leave no hidden gap;
source~20]\label{hol:rc:rem:safeguards}
It would be invalid to take the union of infinitely many finite exceptional
sets and call it finite, and the proof does not do so. On each class either
every condition is an identity, or one condition that is not an identity
already leaves finitely many possible hits; separately, Noetherianity supplies
the finite subfamily. These are two independent safeguards against that
inference, as in Remark~\ref{hol:ud:rem:uniform} for valuations.
\end{remark}

%----------------------------------------------------------------------
\section{Surreal recurrences: leading scales, omnific coding, matrix powers
and mixed operators}\label{hol:rc:sec:surreal}

\subsection{Leading Conway exponents}\label{hol:rc:sub:leading}

For a nonzero $x$ in a workspace $K_\Gamma\subseteq\No[i]$ of
\eqref{hol:eq:embedding}, put $\deg_\omega(x)=-v(x)$: the greatest exponent of
its Conway normal form, the real and imaginary normal forms being combined
into one complex-coefficient normal form. For finitely many surreal or
surcomplex numbers such a workspace exists: the group generated by the union
of their supports is a set (Section~\ref{hol:ud:sub:surreal}). So every
statement below about a fixed finite recurrence takes place in one set-sized
Hahn field, and no proper-class product, sum or determinant is formed.

\begin{theorem}[Surreal and surcomplex recurrence profiles; source~20]
\label{hol:rc:thm:surreal}
Let $(y_n)$ satisfy a constant-coefficient recurrence of finite order over
$\No[i]$, with initial values in $\No[i]$. There are $\mathsf M\ge1$ and
$N\ge0$ such that for each class $b$ modulo $\mathsf M$ either $y_n=0$ for all
$n\ge N$ in the class, or
\begin{equation}\label{hol:rc:eq:degree}
 \deg_\omega(y_n)=n\,\mathsf d_b+\mathsf e_b\qquad(n\ge N,\
 n\equiv b\pmod{\mathsf M}),
\end{equation}
with $\mathsf d_b,\mathsf e_b\in\No$, $\mathsf d_b$ the degree of a nonzero
characteristic root. This applies in particular to every recurrence over $\Oz$
or $\Oz[i]$ with initial values in the same ring. If every nonzero
characteristic root lies in $\No$, $\mathsf M=2$ suffices; if every one is
positive, $\mathsf M=1$ suffices.
\end{theorem}

\begin{proof}
$\No[i]$ is algebraically closed, so the characteristic polynomial splits.
Its finitely many roots, the coefficients of a representation
$y_n=\sum_jR_j(n)\lambda_j^n$ on the tail after the nilpotent segment, and the
recurrence data lie in one set-sized Hahn field $\C((t^\Gamma))$ with
$\Gamma\le\No$, as above. Corollary~\ref{hol:rc:cor:split} applies there, and
negating the valuation gives \eqref{hol:rc:eq:degree}. For roots in $\No$ the
normalized leading coefficients lie in $\R^\times$, whose only nontrivial
torsion element is $-1$; so $\mathsf M_{\mathrm{tor}}$ of their residues
divides two, and it is an admissible period
(Remark~\ref{hol:rc:rem:periods}). Positive roots have positive leading
coefficients, which generate a torsion-free subgroup of $\R_{>0}$. Closure of
$\Oz$ and $\Oz[i]$ under ring operations keeps the terms in those rings.
\end{proof}

\begin{example}[An intercept larger than every ordinary time; source~20]
\label{hol:rc:ex:catalan}
Let $\lambda_-=\frac12\bigl(\omega-\sqrt{\omega^2-4}\bigr)$, the infinitesimal
root of $X^2-\omega X+1$. With $t=\omega^{-1}$,
\[
 \lambda_-=t\,\mathrm{Cat}(t^2),\qquad
 \mathrm{Cat}(X)=\sum_{k\ge0}\mathrm{Cat}_kX^k,\qquad
 \mathrm{Cat}_k=\frac1{k+1}\binom{2k}k,
\]
because $\mathrm{Cat}(X)=1+X\,\mathrm{Cat}(X)^2$. Put
$y_n=\omega^\omega\lambda_-^n$. Then
\begin{equation}\label{hol:rc:eq:catalan}
 y_{n+2}=\omega y_{n+1}-y_n,\qquad y_n\in\Pi,\qquad
 \deg_\omega(y_n)=\omega-n .
\end{equation}
For $n\ge1$, $\lambda_-^n$ has positive integer coefficients, supported at the
Conway exponents $-(n+2k)$, $k\ge0$, with $[t^{n+2k}]\lambda_-^n=
\frac n{n+2k}\binom{n+2k}k$; multiplication by $\omega^\omega$ moves them to
$\omega-n-2k>0$, a reverse well-ordered support. Also $y_0=\omega^\omega\in\Pi$,
so all initial values and all later terms are omnific, and the recurrence
follows from $\lambda_-^2=\omega\lambda_--1$. For every ordinary $n$, $y_n$ is
positive and infinite, and $y_{n+1}<y_n$ since $0<\lambda_-<1$. Its degree
decreases with slope $-1$ but never reaches zero at an ordinary time. This
agrees with Theorem~\ref{hol:rc:thm:surreal} ($\mathsf d=-1$,
$\mathsf e=\omega$) and contradicts any argument that silently treats
$\omega-n$ as though $n$ were cofinal in the surreal exponents.
\end{example}

\subsection{One fixed coefficient can encode any sequence}
\label{hol:rc:sub:coding}

For $x\in\No[i]$ and $\alpha\in\No$ let $[\omega^\alpha]x$ be the coefficient
of $\omega^\alpha$ in the normal form of $x$.

\begin{theorem}[A first-order omnific recurrence encodes any sequence;
source~20]\label{hol:rc:thm:coding}
For every real sequence $\mathbf b=(\mathbf b_m)_{m\ge0}$ there is
$\mathfrak w_{\mathbf b}\in\Pi$ such that
\[
 y_0=\mathfrak w_{\mathbf b},\qquad y_{n+1}=\omega y_n,\qquad y_n\in\Pi,
 \qquad[\omega^\omega]\,y_n=\mathbf b_n
\]
for every ordinary $n\ge0$. The same holds for complex sequences, with
$\mathfrak w_{\mathbf b},y_n\in\Pi+i\Pi\subseteq\Oz[i]$. For every
$A\subseteq\N$ one can therefore arrange
$\{n:[\omega^\omega]y_n\ne0\}=A$; when $A$ is nonempty and $\mathbf b$ takes
the values $0$ and $1$, every $y_n$ is positive.
\end{theorem}

\begin{proof}
Put
\begin{equation}\label{hol:rc:eq:code}
 \mathfrak w_{\mathbf b}=\sum_{m\ge0}\mathbf b_m\,\omega^{\omega-m},
\end{equation}
zero terms omitted. The exponents $\omega-m$ strictly decrease and are all
positive; every nonempty subset of them has a greatest element, the one of
least index. So the support is a reverse well-ordered set, of order type at
most $\omega$ in its decreasing enumeration, and \eqref{hol:rc:eq:code} is a
purely infinite omnific integer. Multiplication by the monomial $\omega^n$
translates exponents with no convolution:
$y_n=\omega^n\mathfrak w_{\mathbf b}=\sum_m\mathbf b_m\omega^{\omega+n-m}$. All
exponents stay positive, $n$ and $m$ being ordinary, and exactly the term
$m=n$ has exponent $\omega$. The complex case is the same coefficientwise, or
follows by separating real and imaginary parts. For the last assertion take
$\mathbf b$ the indicator of $A$; the leading coefficient is then $1$.
\end{proof}

\begin{corollary}[No coefficientwise Skolem--Mahler--Lech theorem for
omnific recurrences; source~20]\label{hol:rc:cor:nosml}
It is false that all fixed-Conway-coefficient zero sets, or nonzero sets, of
finite omnific recurrences are eventually periodic. Already the fixed
multiplier $\omega$, the fixed observed exponent $\omega$ and a varying purely
infinite initial value realize every subset of $\N$.
\end{corollary}

\begin{remark}[No conflict with the universal quotient; source~20]
\label{hol:rc:rem:quotients}
The map $x\mapsto[\omega^\omega]x$ is additive on $\Oz$ but not
multiplicative: it kills $\omega^{\omega/2}$ and sends its square to $1$. So
Theorem~\ref{hol:rc:thm:coding} does not contradict the theorem of the
set-sized-quotient report that every additive and multiplicative map from
$\Oz$ to a set-sized ring factors through the constant term
(\texttt{osq:thm:universal}~\cite{repoosq}); every $y_n$ above has constant
term zero. \emph{Related (batch~34).} The set-sized target in that theorem is
needed: the same report constructs unital ring homomorphisms
$\operatorname{sp}_a:\Oz\to\No[i]$, $a>0$, into the proper-class field $\No[i]$ with
$\operatorname{sp}_a(\omega^a)=-1$ (\texttt{osq:cr:thm:residually}), so they do
not factor through the constant term. Neither result bears on
Theorem~\ref{hol:rc:thm:coding}, whose observable is additive and not
multiplicative.
\end{remark}

\begin{remark}[The element is a known stream; merge]\label{hol:rc:rem:stream}
The number $\mathfrak w_{\mathbf b}$ is, with $a_n=\mathbf b_n$, the series
\texttt{onot:eq:guardbasic} of the omnific-notations report~\cite{repoonot},
which already observes that $\sum_na_n\omega^{\omega-n}=
\omega^\omega\sum_na_n(\omega^{-1})^n$ has positive exponents and a valid
support for every real sequence, and that an algorithm for each $a_n$ need
not decide whether all $a_n$ vanish. Source~20 does not cite it. The element
is therefore not new; what source~20 adds is its reading as a recurrence
observable, the orbit under multiplication by $\omega$ read at one fixed
exponent, and Corollary~\ref{hol:rc:cor:nosml}. The independent omnific codes
of Theorem~\ref{hol:fh:thm:codes} encode subsets of $\N$ differently, through
algebraic independence, not a coefficient observation.
\end{remark}

\paragraph{The exact distinction between the two sides (source~20).}
In Theorem~\ref{hol:rc:thm:compression} the roots are units, so the supports of
all iterates lie in one well-ordered envelope, and a fixed coefficient receives
only finitely many binomial contributions. In
Theorem~\ref{hol:rc:thm:coding} multiplication by $\omega$ has nonzero
valuation and translates the support, and the fixed observation reads a new
stored coefficient at each step. This movement of support, not nonlinear
arithmetic, permits universality. The leading exponent stays tame: if
$\mathbf b\ne0$ and $m_0$ is its first nonzero index, then
$\deg_\omega(y_n)=\omega+n-m_0$, as Theorem~\ref{hol:rc:thm:surreal}
requires. The exact leading-scale theorem and unrestricted lower coefficients
coexist in the simplest recurrence.

\subsection{Determinantal valuation profiles}\label{hol:rc:sub:matrix}

\begin{corollary}[Determinantal valuation profiles; source~20]
\label{hol:rc:cor:matrix}
Let $\mathbf A\in\operatorname{GL}_{\bar d}(\No[i])$, and for
$1\le j\le\bar d$ put
\[
 \operatorname{mv}_j(n)=\min\{v(\det\mathbf B):\mathbf B\text{ a }j\times j
 \text{ minor of }\mathbf A^n,\ \det\mathbf B\ne0\},\qquad
 \operatorname{mv}_0(n)=0,
\]
in a workspace containing the entries of $\mathbf A$. There is one finite
period such that every $\operatorname{mv}_j(n)$ and every difference
$\operatorname{mv}_j(n)-\operatorname{mv}_{j-1}(n)$ is eventually affine on each
residue class.
\end{corollary}

\begin{proof}
The $j\times j$ minors of $\mathbf A^n$ are the entries of
$(\bigwedge^j\mathbf A)^n$, and each entry satisfies the constant-coefficient
recurrence given by the characteristic polynomial of the finite matrix
$\bigwedge^j\mathbf A$. By Theorem~\ref{hol:rc:thm:surreal}, in its valuation
form, each is eventually zero or of eventually affine valuation on every class
of a finite period; refine the finitely many periods to one and take a common
threshold. Invertibility of $\mathbf A^n$ gives a nonzero minor of each size at
every $n$, so on each class some minor is active. The minimum of finitely many
eventually affine functions is eventually one of them
(Lemma~\ref{hol:ud:lem:lines}), and differences of eventually affine functions
are eventually affine.
\end{proof}

Over a valuation ring these minima are the sums of the first $j$ invariant-factor
valuations of a finite matrix, after a common monomial rescaling if
necessary; such a finite diagonal reduction is obtained directly by choosing
an entry of least valuation, clearing its row and column by divisibility in
the valuation ring, and repeating. The corollary gives exact eventual scale
profiles of these finite-dimensional invariants. It asserts neither a
real-valued Lyapunov limit nor eventual periodicity of all matrix
coefficients (source~20).

\subsection{Mixed operators: an independent derivation of part of
Theorem~\ref{hol:ud:main:mixed}}\label{hol:rc:sub:mixed}

Source~20's Sections~9--10 treat finite operators \eqref{hol:ud:eq:L} whose
multipliers lie in a torsion-free subgroup $\mathfrak G\le\KK^\times$. Distinct
elements of such a group have no root-of-unity quotient, so the hypothesis of
Theorems~\ref{hol:ud:thm:rigidity} and~\ref{hol:ud:thm:dichotomy} holds, and
source~18 needs less: faithfulness of the normal form needs only distinct
multipliers (Corollary~\ref{hol:ud:cor:faithful}), and rigidity only a
noncofinal relative scale. Source~20's reduction to a recurrence, its escape
lemma with the closed region, and its unit-valuation theorem are printed once
(Section~\ref{hol:rc:sub:conventions}); it found them independently of
source~18, whose text it had not seen. Its remaining two theorems are special
cases, printed here with its route.

\begin{corollary}[Mixed rigidity below a proper convex scale; source~20]
\label{hol:rc:cor:convex}
Let $\mathfrak G\le\KK^\times$ be torsion-free and $L$ a nonzero normal form
\eqref{hol:ud:eq:L} with distinct multipliers in $\mathfrak G$. If their
valuations lie in a proper convex subgroup $\Delta<\Gamma$, then $L$ has a
uniform exterior obstruction and degree bound as in
Theorem~\ref{hol:ud:thm:rigidity}; in particular every strongly entire
solution of $Lf=0$ is a polynomial.
\end{corollary}

\begin{proof}
The differences $v(\lambda_i)-v(\lambda_j)$ lie in $\Delta$, and a nonzero
element of a proper convex subgroup is not an order unit, since a convex
subgroup containing an order unit is $\Gamma$. So the relative scale is
noncofinal and Theorem~\ref{hol:ud:thm:rigidity} applies. Source~20's route
avoids the relative scale: by Theorem~\ref{hol:ud:thm:affine} the recurrence
coefficients have valuations $\beta_{j,b}+n\varpi_{j,b}$ with slopes in
$\Delta$; for any $\gamma>0$ outside $\Delta$, convexity gives
$n(\varpi_{0,b}-\varpi_{j,b})<\gamma$ for every $n\ge0$, whatever the sign of
the difference, so the costs are bounded by $C+\gamma$ with $C$ the largest
intercept difference, and Corollary~\ref{hol:cor:boundedcost} concludes. The
contradiction occurs in $\Gamma$ itself; no coarsened summability is used.
\end{proof}

The converse inclusion fails (merge): the multipliers of
Example~\ref{hol:ud:ex:commonscale} have the cofinal common valuation $\mu$,
which lies in no proper convex subgroup, yet
Theorem~\ref{hol:ud:thm:rigidity} excludes every nonpolynomial entire solution,
because only relative valuations matter. Source~18's theorem is strictly
stronger than source~20's Theorem~10.3.

\begin{corollary}[Exact finite-group criterion; source~20]
\label{hol:rc:cor:finitegroup}
Let $\mathfrak G\le\KK^\times$ be finitely generated and torsion-free, and let
$\Delta_{\mathfrak G}$ be the convex subgroup of $\Gamma$ generated by
$v(\mathfrak G)$. The following are equivalent.
\begin{enumerate}[label=\textup{(\roman*)}]
\item Some nonpolynomial $f\in\EK$ is annihilated by a nonzero finite operator
$\sum_{\lambda,r}p_{\lambda,r}(z)D_z^r\sigma_\lambda$ with multipliers
$\lambda\in\mathfrak G$ and $p_{\lambda,r}\in\KK[z]$.
\item $\Delta_{\mathfrak G}=\Gamma$.
\item Some $p\in\mathfrak G$ has positive order-unit valuation.
\end{enumerate}
When \textup{(iii)} holds, $\Theta_p$ of \eqref{hol:eq:theta} is a witness for
\textup{(i)}: it satisfies \eqref{hol:eq:thetahom}, with the multipliers
$1,p,p^2\in\mathfrak G$ and no derivative.
\end{corollary}

\begin{proof}
(i)$\Leftrightarrow$(iii): a finite list of distinct elements of $\mathfrak G$
has no root-of-unity quotient, so by Theorem~\ref{hol:ud:thm:dichotomy},
(i) holds if and only if $\abs{v(\lambda_i/\lambda_j)}$ is an order unit for
some $\lambda_i,\lambda_j\in\mathfrak G$; as $\lambda_i/\lambda_j\in\mathfrak G$
and $1\in\mathfrak G$, this is (iii), after inverting if the valuation is
negative. For the witness, $\Theta_p$ is a nonpolynomial element of $\EK$
(Lemma~\ref{hol:ud:lem:theta}), and the operator
$\sigma_1-(1+z)\sigma_p+pz\sigma_{p^2}$ of \eqref{hol:eq:thetahom} is a nonzero
normal form (Corollary~\ref{hol:ud:cor:faithful}).
(ii)$\Leftrightarrow$(iii), source~20's argument: for generators
$q_1,\ldots,q_s$ of $\mathfrak G$ the convex subgroup generated by
$v(\mathfrak G)$ is $\{\gamma:\abs\gamma\le n\eta\text{ for some }n\ge1\}$ with
$\eta=\max_i\abs{v(q_i)}$ (zero if every $v(q_i)=0$). If it is the nonzero
group $\Gamma$, then $\eta>0$ is an order unit, and a generator or its inverse
has valuation $\eta$; conversely an order-unit valuation forces
$\Delta_{\mathfrak G}=\Gamma$. Source~20 proves (i)$\Rightarrow$(ii) by
Corollary~\ref{hol:rc:cor:convex} and (iii)$\Rightarrow$(i) by the direct
estimate $h_{n+1}-h_n=n\beta+\delta>0$ for $n>A$, where $h_n=\binom n2\beta+n\delta$,
$\beta=v(p)$, $\delta=v(x)$ and $\abs\delta\le A\beta$, followed by the
support count of Remark~\ref{hol:rem:thetasecond}.
\end{proof}

When $\Gamma$ has an order unit, (ii) says that $v(\mathfrak G)$ is not
contained in the maximal proper convex subgroup $\Delta^{\mathrm{top}}$ of
Proposition~\ref{hol:ud:prop:top}; that proposition is the reformulation of
the negative case for an arbitrary finite list.

\begin{remark}[Finite generation cannot be dropped from (ii); merge]
\label{hol:rc:rem:finitegeneration}
The equivalence of (i) and (iii) uses no finite generation. The equivalence
with (ii) does. In $\Gamma_\infty$ of \eqref{hol:eq:gammainfty}, which has no
order unit, the torsion-free group $\mathfrak G$ generated by the monomials
$t^{\omega^j}$, $j\ge0$, has cofinal valuations, so $\Delta_{\mathfrak
G}=\Gamma_\infty$; but no element of $\Gamma_\infty$ is an order unit, and by
Theorem~\ref{hol:ud:thm:dichotomy}, or already by
Theorem~\ref{hol:cf:main:entire}\,(b), no nonpolynomial entire series satisfies
a nonzero mixed equation with multipliers in $\mathfrak G$.
\end{remark}

\begin{remark}[Which part source~20 counted as new; source~20]
\label{hol:rc:rem:stale}
Source~20 recorded that the partial-theta witness and the one-parameter
order-unit boundary were already in this report, and counted as its own
contribution the general exponential-polynomial valuation theorem and its use
for the mixed unit case and the finite multibase criterion; it also stressed
that it does not infer a mixed assertion by combining two pure ones. The
first sentence is accurate. The second was accurate against the text at its
pin and is stale against the current one: source~18, placed before that pin
and merged after it, proves the valuation theorem
(Theorems~\ref{hol:ud:thm:unitvalues} and~\ref{hol:ud:thm:affine}) and the
relative-scale criterion, of which source~20's Theorems~10.1, 10.3 and~10.4
are special cases. The two sources are independent derivations of the same
day; no priority between them is asserted.
\end{remark}

\subsection{Boundary examples}\label{hol:rc:sub:boundaries}

\paragraph{Torsion is a genuine obstruction (source~20).}
Over $\KK=k((t^\Q))$ let $p=t$, so that $\Theta_p$ is strongly entire, and put
$f(z)=\Theta_p(z^2)$. Then $f$ is nonpolynomial and strongly entire, and
$f(-z)-f(z)=0$; the operator $\sigma_{-1}-1$ is not zero on all formal series,
it merely kills every even series. So excluding only the formally zero
identity $\sigma_{-1}^2-1$ would not repair the theorems; a theory with torsion
must track residue-degree sectors and the operator's action on them. This is
the case $h=2$ of the paragraph on root-of-unity quotients in
Section~\ref{hol:ud:sub:examples} and of Proposition~\ref{hol:prop:torsion}.
Residual torsion, by contrast, is allowed throughout: $-(1+t)$ has infinite
order and is covered by Corollary~\ref{hol:ud:cor:unitrigidity}.

\paragraph{Two scales inside one rank-two workspace (source~20).}
Let $\Gamma=\Q e_0\oplus\Q e_1$ with $0<e_0\ll e_1$. Dilation by $t^{e_0}$,
together with any finite family of unit dilations generating a torsion-free
group, stays in the proper convex subgroup $\Q e_0$, and by
Corollary~\ref{hol:rc:cor:convex} supports no nonpolynomial entire solution of
a nonzero mixed equation; dilation by $t^{e_1}$ admits the partial theta
witness, $e_1$ being an order unit of the same group. So valuation rank alone
is not the classification parameter, as in Examples~\ref{hol:ex:ranktwo}
and~\ref{hol:ud:ex:crossing}.

\begin{example}[Characteristic zero cannot be omitted; source~20]
\label{hol:rc:ex:charp}
In $\mathbb F_{\mathfrak p}((t))$, for every $n\ge1$,
\[
 v\bigl((1+t)^n-1\bigr)=\mathfrak p^{v_{\mathfrak p}(n)} .
\]
Indeed, writing $n=\mathfrak p^am'$ with $\mathfrak p\nmid m'$,
$(1+t)^n=(1+t^{\mathfrak p^a})^{m'}$, whose first nonconstant term is
$m't^{\mathfrak p^a}\ne0$. These valuations have no eventually periodic finite
range although $1+t$ is a unit; the binomial-polynomial argument over $k$ and
the characteristic-zero zero-set theorem are both essential. The paragraph on
positive characteristic in Section~\ref{hol:ud:sub:examples} records the
subsequence $n=\mathfrak p^a$; the formula gives every $n$ (merge's
comparison). Positive-characteristic zero sets need a different
theory~\cite{derksen}.
\end{example}

\paragraph{Leading exponents are not birthdays or signs (source~20).}
$\deg_\omega$ records a leading normal-form exponent, not the birthday of the
surreal, and nothing here bears on Gonshor's product-birthday problem. An
eventually affine valuation does not give an eventually periodic sign:
ordinary real recurrence sequences have valuation zero whenever nonzero, while
their signs can oscillate with complex phases. No sign classification is used
or claimed.

\paragraph{Effectivity and the classical Skolem problem (source~20).}
The proofs give the finite exceptional sets existentially, and an arbitrary
coefficient exponential polynomial invokes the Skolem--Mahler--Lech theorem.
No algorithm locating all exceptional zeros is claimed; a general one would
solve the ordinary Skolem problem, by applying it to constant Hahn series.
Noetherian compression likewise gives the existence of a finite subfamily, not
an algorithm finding it from an arbitrary infinitely presented family. The
coding theorem is deliberately non-effective when $\mathbf b$ is
noncomputable: its initial value already contains the sequence. It gives no
finite effective representation of every subset of $\N$, and it does not show
that coefficient extraction is definable in the pure ring language of $\Oz$.

\subsection{Predecessors, repository comparison and priority}
\label{hol:rc:sub:predecessors}

\paragraph{Literature.}
Source~20 imports Conway normal forms and the Hahn-field viewpoint as
established foundations~\cite{bm}; the Hahn--Neumann support lemma, credited
historically to Neumann~\cite{neumann} (not audited as a historical paper) and
consulted in Bagayoko, Krapp, Kuhlmann, Panazzolo and Serra~\cite{bagayoko}
(version~2), for positive-support summability; and the classical
Skolem--Mahler--Lech theorem~\cite{lech,belldoc}, whose Bell account it read
and inspected visually. Fuchs and Heintze~\cite{fuchsheintze} obtain effective
valuation-growth bounds for recurrences over one-variable function fields
under nondegeneracy assumptions; source~20's exact eventual-affine statement
allows arbitrary Hahn supports and ordered value groups but sacrifices a
general effective exceptional-index bound, and does not improve their
constants. Krapp, Kuhlmann and Serra~\cite{kkslinrec} study generalized
recurrence relations indexed by exponent groups and their relation to
fraction fields: there the recurrence index is an exponent, here it is an
ordinary natural number and only the values lie in a Hahn field; the setups
are distinct and no theorem of theirs is imported. Derksen~\cite{derksen} is
context for positive characteristic. Formal summability is used in the
support-sensitive sense of the Hahn literature, not as numerical
convergence. The Wikipedia article on surreal numbers was read for
orientation only. For this merge the bibliographic data were checked; their
content was not rechecked, and no proof here rests on Fuchs--Heintze,
Krapp--Kuhlmann--Serra or Derksen.

\paragraph{Repository comparison at the pin.}
Source~20 read, at \texttt{efc5446}, the repository README and catalogue, this
report's \texttt{article.tex} in the line ranges listed in its audit
(1330--1500, the unit-orbit examples and differential recurrences; 2020--2150,
the sharp dilation criterion, partial theta and torsion; 2150--2228,
Lemma~\ref{hol:lem:bivariate}, Theorem~\ref{hol:thm:mixed} and the limitation
after it), and the README of the set-sized-quotient report~\cite{repoosq},
whose universal quotient theorem it uses only to explain the absence of a
contradiction (Remark~\ref{hol:rc:rem:quotients}), without auditing that
report. It did not read the placed audit file of source~18, nor the
omnific-notations report (Remark~\ref{hol:rc:rem:stream}). The comparison was
targeted, not a proof-by-proof audit of all reports, and the absence of a
matching theorem in it is not a proof of novelty. Its statements about this
report were accurate at the pin, as recorded in
Section~\ref{hol:rc:sub:conventions}.

\paragraph{Priority.}
Source~20 separates four groups. \emph{Imported foundations:} Hahn
multiplication and positive-support summability, Conway normal forms, real
closedness of $\No$, the Skolem--Mahler--Lech theorem and Hilbert's basis
theorem. \emph{Reproved repository mechanisms:} finite-step recurrence escape,
the noncofinal-scale obstruction and the partial-theta sharpness construction.
\emph{Candidate-original package:} the uniform coefficient algebra,
coefficient-constraint compression, arbitrary-rank exact valuation profiles,
their mixed unit and multibase application, and the juxtaposed omnific
coefficient-universality theorem. \emph{Not claimed:} independent priority
certification, formal verification, effective Skolem decision procedures,
arbitrary nonlinear mixed rigidity, and a classification of full omnific
membership sets for unrestricted recurrences. Against the current text, the
valuation profiles and the mixed application are source~18's, found
independently (Remark~\ref{hol:rc:rem:stale}), and the coding element is the
omnific-notations report's (Remark~\ref{hol:rc:rem:stream}); what remains
proposed by source~20 is the coefficient algebra, Noetherian compression with
its universal period, the recurrence statements over $\No[i]$, $\Oz$ and
$\Oz[i]$ with the decreasing omnific example, the coding reading with
Corollary~\ref{hol:rc:cor:nosml}, and the determinantal profiles.
``Candidate-original'' means developed and proved at manuscript level with no
matching result found in a targeted comparison; it does not establish
priority.

\subsection{What the recurrence part does not claim}
\label{hol:rc:sub:nonclaims}

Every limitation stated by source~20, in its manuscript, README, audit file or
verification record, is kept.
\begin{enumerate}[label=\textup{(V\arabic*)}]
\item Source~20 is an AI-assisted, unrefereed research draft with manuscript
arguments, not an independently certified breakthrough. No independent
referee report and no Lean formalization accompany it; no Lean build was run,
no declaration was added, and the repository's axiom audit does not cover
this part. Independent correctness and worldwide priority are not certified.
\item The Skolem--Mahler--Lech theorem, the Hahn--Neumann support lemma and
Hilbert's basis theorem are imported. No effective bound for exceptional
zeros and no algorithm locating them is given; such an algorithm would solve
the ordinary Skolem problem, which is not claimed solved.
\item Noetherian compression gives the existence of a finite subfamily only,
not an algorithm finding it from an infinitely presented family.
\item The degrees of the coefficient polynomials $r_{j,\gamma}$ are unbounded,
and a universal residual period gives no common transient bound for all
coefficient tests on one orbit (Example~\ref{hol:rc:ex:transient}).
\item No nonlinear mixed rigidity is asserted; source~20's nonlinear question
is re-scoped here (Question~\ref{hol:rc:q:nonlinear}).
\item No classification of full omnific membership sets: the integer
constant-term condition is not a coefficient-algebraic condition and is not
handled (Question~\ref{hol:rc:q:membership}).
\item The coding theorem is storage of a sequence in an infinite initial
value, not an algorithm that generates information; it is non-effective for
noncomputable sequences, gives no finite effective representation of subsets
of $\N$, and shows neither that coefficient extraction is definable in the pure
ring language of $\Oz$ nor any interpretation of second-order arithmetic.
\item Entireness is relative to one set-sized Hahn field; a series entire in
one workspace need not be entire after enlargement, and nothing concerns
simultaneous entireness on all of $\No[i]$. Statements over $\No[i]$ concern
each finite recurrence localized to a set-sized Hahn field.
\item $D_z$ fixes the coefficients and is not the Berarducci--Mantova
derivation; derivations of the coefficient field are not treated
(Question~\ref{hol:rc:q:derivations}).
\item Actual torsion among the dilation parameters and positive characteristic
are excluded, with counterexamples (Section~\ref{hol:rc:sub:boundaries}).
\item Nothing is claimed about real-valued Lyapunov limits or eventual
periodicity of the coefficients of matrix powers.
\item $\deg_\omega$ is neither a birthday nor a sign; Gonshor's
product-birthday problem is not addressed, and no sign classification is
claimed.
\item The effective constants of Fuchs and Heintze are not improved, and
their setting differs; the setup of Krapp, Kuhlmann and Serra is distinct.
\item Neumann's paper is credited historically, not audited; the
set-sized-quotient report was not audited, and its theorem is used only to
explain the absence of a contradiction.
\item The 4{,}605 finite checks concern finite identities and examples. They
prove neither the support lemma, the Skolem--Mahler--Lech theorem, the
infinite-support minimum argument nor the proper-class localization, and a
PDF build concerns document integrity, not truth. Their number measures no
proof coverage: $1{,}189$ of them compare elementary integer or lexicographic
facts that hold by construction (Section~\ref{hol:sec:verification}).
\item The repository comparison at \texttt{efc5446} was targeted, not
exhaustive; its description of this report is stale in the respects recorded
in Section~\ref{hol:rc:sub:conventions}, and its search was not a literature
priority audit.
\item \emph{[merge]} Remarks~\ref{hol:rc:rem:periods},
\ref{hol:rc:rem:stream} and~\ref{hol:rc:rem:finitegeneration}, the
identification of source~20's Theorems~10.3 and~10.4 as special cases of
Theorem~\ref{hol:ud:main:mixed}, the comparison in
Example~\ref{hol:rc:ex:charp}, and the re-scoping of
Question~\ref{hol:rc:q:nonlinear} are the merge's; no historical priority is
claimed for them.
\end{enumerate}

\section{Conclusion}

Strong entireness over a surreal Hahn workspace imposes a global scale
requirement that a finite linear differential equation cannot meet unless its
solution terminates. A dilation equation can meet that requirement, but
exactly when its valuation reaches every scale through integer multiples. The
decisive distinction is therefore not between Archimedean and non-Archimedean
fields, nor between rank one and higher rank; it is between \emph{a cofinal
dilation scale} and \emph{a scale trapped below another one}.

The obstruction is finite and robust. A linear recurrence forces a later
nonzero coefficient with a controlled valuation, and a large enough argument
turns an infinite chain of such coefficients into a nonincreasing sequence of
leading exponents. Descending supports and infinite reuse of one exponent are
two distinct ways in which the evaluation fails, and keeping them separate is
what makes one exclusion boundary closed and the other open. The mechanism
works with arbitrary Hahn coefficient tails, arbitrarily high valuation rank,
and sparse coefficient sequences, without replacing strong summability by
topological convergence and without any hidden completeness assumption.

The countable-rank example makes the separation concrete. The series
$\sum_{n\ge0}\omega^{-\omega^n}z^n$ is evaluable at every point of its full
workspace, yet no single nontorsion dilation in that workspace supplies a
finite linear equation for it. That gap between existence and finite linear
describability is the main mathematical contribution proposed here.

The nonlinear part pushes the same global requirement past linear
equations by a different mechanism. A finite algebraic relation between a
function and its first derivative is incompatible with nonpolynomial strong
entireness over any nontrivial characteristic-zero Hahn workspace: at one
large admissible scale the whole series leaves a finite initial polynomial,
whose top coefficient is controlled by a corner polynomial that is
automatically nonzero in order one. In higher order an exact cancellation of
that corner is a genuine obstruction to that method, and a fixed nonlinear
equation can then have polynomial solutions of unbounded degree.

The coefficient-field part removes the order restriction when the value
group has no order unit, by a third mechanism. A finite algebraic relation,
linearized along its own solution, confines every Taylor coefficient to one
finitely generated field, and such a field cannot supply independent
valuations at cofinally many Archimedean levels; a nonpolynomial entire
series needs them. So every nonpolynomial strongly entire series is then
differentially transcendental over the full field of constants, jointly with
its dilates, and every differentially algebraic quotient of entire series is
rational, while one explicit workspace still carries continuum many entire
series with jointly independent jets. With an order unit the mixed statement
fails.

The theta-descent part shows that the purely differential statement fails
too, so the order unit is the exact boundary. A bilateral theta series,
descended along $z=u+u^{-1}$ by the Chebyshev identity, becomes an ordinary
power series that is strongly entire exactly when its scale is an order unit,
and the classical elliptic identity gives it an explicit algebraic
differential equation of order three. Its residue corner vanishes, as it
must. The shift of the theta series by $u\mapsto p^2u$ then excludes every
equation of lower order, by three elementary difference-field arguments.

The order-three-threshold part shows that no other series does better. At the
infinitely many scales where a nonpolynomial entire series has a nonmonomial
initial polynomial, a fixed equation of order at most two reduces to one
homogeneous polynomial over the residue field, and logarithmic Euler
derivatives confine both endpoints of those initial polynomials to a finite
set, even when the leading-monomial test vanishes identically. So order three
is the least possible for every value group with an order unit. At order three
one further integer, the gap next to an endpoint, survives, and the theta
equation uses it with gap one. Which series attain order three is the open
part of Question~\ref{hol:q:nonlinear}.

The unit-dilation part closes the linear theory. Mixing derivatives with
dilations adds no new way to escape: the coefficients of a mixed recurrence
are exponential polynomials with Hahn coefficients, whose valuations are
eventually affine on finitely many progressions, and so for every value
group a nonpolynomial entire solution of some mixed linear equation exists
exactly when two multipliers are separated at the top Archimedean scale. In
particular a unit dilation, however close to a root of unity its residue,
adds nothing to derivatives.

The Newton-rigidity part adds a first result at order three. When the
exterior form of an equation has a nonzero first polar along the twisted
cubic of monomial jets, the two ends of a far initial polynomial demand next
exponents on opposite sides, and every cubic form vanishing on that curve has
such a polar; so order three with jet degree at most three is also rigid. In
degree four the test fails, the exact solutions of two boundary equations are
doubling binomials, and a sparse entire series shows that initial equations
valid at every scale need not lift to an identity.

The theta-hierarchy and single-scale parts measure what order three leaves
room for. Theta descents at scales with independent reciprocals have distinct
quadratic Gauss growth, so over $\R((t^\R))$ they are algebraically
independent, and polynomially weighted sums of them reach every large
differential order; yet every one is differentially algebraic. At the other
extreme a single scale already carries lacunary series with rational Taylor
coefficients that satisfy no algebraic differential equation at all, because a
maximal-degree band of heights cannot be cancelled.

The polynomial-composition part shows how much of the order-unit phenomenon
belongs to degree one. A polynomial argument of degree $d\ge2$ multiplies the
scale by $d$, so a superlinearly growing Gauss profile cannot be balanced by
arguments of lower degree unless their composition weight reaches the top
one; then every strongly entire solution is a polynomial, for every value
group, with a degree bound that keeps the exact indicial resonances. The
finite coefficient locus that remains can be as complicated as any affine
hypersurface over $\Q$, and its omnific points are undecidable in general.

The recurrence part reads the ordinary recurrences behind all these proofs one
coefficient at a time. When the characteristic roots are units, every iterate
lies in one well-ordered support envelope and each fixed coefficient is an
ordinary exponential polynomial over the residue field; all of them lie in one
finitely generated algebra, so any set-sized family of coefficient conditions
reduces along the orbit to finitely many, with hitting sets periodic modulo
one period up to a finite set. When a root moves the support, as $\omega$ does,
one fixed coefficient can store any real sequence, although the leading
exponent stays eventually affine.

\appendix

\section{Notation}\label{hol:app:notation}

\begin{center}
\begin{tabularx}{\textwidth}{@{}l X@{}}
\toprule
Symbol & Meaning\\
\midrule
$\Gamma$ & A fixed nonzero set-sized ordered abelian group; divisible only
where stated.\\
$\K$ & The complex Hahn field $\C((t^\Gamma))$.\\
$v(a),\lc(a),\res(a)$ & Least support exponent, its coefficient, and the
residue when $v(a)=0$.\\
$\E$, $\Dom(f)$ & Series strongly evaluable at every point of $\K$; the exact
set where a series is strongly evaluable.\\
$D_z$ & Formal derivative in $z$, zero on every element of $\K$.\\
$\sigma_q$ & Argument dilation $f(z)\mapsto f(qz)$.\\
$M(S)$ & The additive monoid generated by a well-ordered $S\subseteq\Gamma_{>0}$.\\
$\Delta_\beta$ & The smallest convex subgroup containing $\beta>0$,
\eqref{hol:eq:delta}.\\
$B$, $B^{+}$, $B_{\mathrm{rec}}$ & Exclusion thresholds in $\Gamma$;
$B_{\mathrm{rec}}$ divides by the jump length and needs $\Gamma$ divisible.\\
$\overline B^{\flat}_{\mathrm{rec}}$,
$\overline B_{\mathrm{rec}}$ & Their coarsened analogues in
$\Gamma/\Delta_{\abs{v(q)}}$.\\
$\mu$, $\beta$ & Usually a positive dilation valuation; an order unit only
when stated.\\
$\fall Xr$ & The ordinary falling factorial polynomial of degree $r$.\\
$\Theta_p$ & The partial theta series \eqref{hol:eq:theta}.\\
$\Gamma_\infty$ & The increasing-rank direct sum \eqref{hol:eq:gammainfty}.\\
\midrule
\multicolumn{2}{@{}l}{\emph{Nonlinear part,
Sections~\ref{hol:nl:sec:setting}--\ref{hol:nl:sec:consequences}}}\\
$k$, $\KK$, $\EK$ & Any characteristic-zero field; $k((t^\Gamma))$; its
strongly entire series (Convention~\ref{hol:nl:conv:field}).\\
$\Sfam_m$, $\Sfam_m^+$, $\rho$ & Series with strongly summable coefficient
family; those with coefficients of nonnegative valuation; coefficientwise
residue.\\
$\gamma_\delta$, $\phi_\delta$, $N_\delta$ & Gauss value, initial polynomial
and top active degree at the scale $t^{-\delta}$,
\eqref{hol:nl:eq:gaussdata}.\\
$d_P$, $w_P$, $\mathcal C_P$, $\beta_P$ & Corner data of a differential
polynomial, \eqref{hol:nl:eq:dw}.\\
$\chi_P$, $I_P$ & Residue and full corner polynomials,
\eqref{hol:nl:eq:chiI}.\\
$\kappa_P$ & The rational degree cutoff \eqref{hol:nl:eq:kappa}.\\
$\delta_0$ & A solution-dependent exclusion bound,
\eqref{hol:nl:eq:deltazero}.\\
$\vartheta_q$, $\wt$ & Positive-weight Euler operator and weight.\\
\midrule
\multicolumn{2}{@{}l}{\emph{Coefficient-field part,
Sections~\ref{hol:cf:sec:setting}--\ref{hol:cf:sec:families}}}\\
$\FE$ & The fraction field of $\EK$ in $\KK((z))$ (``meromorphic'').\\
$\jet\lambda jf$ & $(D_z^jf)(\lambda z)$, differentiation before dilation,
\eqref{hol:cf:eq:jet}.\\
$\rrk$, $\trdeg$, $\dtrdeg$ & Rational rank; transcendence degree;
differential transcendence degree.\\
$\cvr(f)$ & Coefficient-value rank \eqref{hol:cf:eq:cvr}.\\
$\Delta$, $v_\Delta$, $k_\Delta$, $\red_\Delta$ & A convex subgroup, the
coarsened valuation $v+\Delta$, its residue field and residue map.\\
$\Pi_{ij}$, $\theta$, $\xi_{ij}$, $\Xi$ & Evaluated partials, offset,
leading coefficients and linearization multiplier,
\eqref{hol:cf:eq:theta}--\eqref{hol:cf:eq:multiplier}.\\
$L$, $k_0\subseteq K_0$ & A coefficient subfield; a generic pair of
characteristic-zero fields.\\
$\mathcal F_\alpha$, $\code$, $R_\tau$ & The branch family
\eqref{hol:cf:eq:branches}, its binary code, the rank-one family.\\
\bottomrule
\end{tabularx}
\end{center}

\begin{center}
\begin{tabularx}{\textwidth}{@{}l X@{}}
\toprule
Symbol & Meaning\\
\midrule
\multicolumn{2}{@{}l}{\emph{Theta-descent part,
Sections~\ref{hol:td:sec:setting}--\ref{hol:td:sec:boundary}}}\\
$\mu$, $p$, $\Delta_\mu$ & A positive scale (an order unit only where
stated), $p=t^\mu$, the convex subgroup \eqref{hol:eq:delta}.\\
$C_n$, $\mathcal T_p$ & Monic Chebyshev polynomials \eqref{hol:td:eq:cheb};
the descended series \eqref{hol:td:eq:series}.\\
$\Psi_p$, $u$, $\check z$ & The bilateral theta series
\eqref{hol:td:eq:psi}; its variable; $u-u^{-1}$.\\
$\mathcal O_{\Delta_\mu}$, $x_m$, $\mathcal P_p$ & The strong domain of
$\mathcal T_p$; its zeros \eqref{hol:td:eq:roots}; $\prod_m(1-p^{2m})$.\\
$\mathcal Q_r$, $\mathfrak c$, $g_2$, $g_3$, $\Phi$ & Lambert series in the
nome $p^2$ and the constants \eqref{hol:td:eq:lambert}--\eqref{hol:td:eq:Phi}.\\
$\mathcal A[f]$, $\mathcal H[f]$, $\mathcal K[f]$, $\mathcal B[f]$ &
Differential expressions \eqref{hol:td:eq:A}--\eqref{hol:td:eq:KB}.\\
$\mathcal U_\mu$, $f^\natural$, $\vartheta_u$, $\mathsf S$ & The Laurent ring,
the substitution $z=u+u^{-1}$, $uD_u$, and $u\mapsto p^2u$.\\
$\psi_1,\psi_2,\psi_3$; $\mathcal L$, $\mathcal A_1$ & $\vartheta_u\Psi/\Psi$
and its two $\vartheta_u$-derivatives; $D_zf/f$ and $D_z\mathcal A[f]$.\\
\midrule
\multicolumn{2}{@{}l}{\emph{Order-three-threshold part,
Sections~\ref{hol:ot:sec:setting}--\ref{hol:ot:sec:theta}}}\\
$\Gamma_\Q$, $\KK_\Q$ & The divisible hull $\Gamma\otimes_\Z\Q$ and
$k((t^{\Gamma_\Q}))$.\\
$\operatorname{Act}_\delta$, $\underline N_\delta$ & Active indices at the
scale $\delta$ and the least of them \eqref{hol:ot:eq:active}; a
\emph{break} is a scale with $\underline N_\delta<N_\delta$.\\
$\vartheta_X$; $\mathsf u_\phi,\mathsf v_\phi,\mathsf w_\phi$ & $XD_X$;
$\vartheta_X\phi/\phi$ and its two $\vartheta_X$-derivatives.\\
$\mathfrak h_P$ & The exterior form of $P$ \eqref{hol:ot:eq:exterior}.\\
$\mathfrak q$, $i_0$, $\mathfrak r$, $\operatorname{Sp}(\mathfrak h)$ &
Dehomogenization, the power of $V$ removed, the remainder, the finite endpoint
spectrum (Lemma~\ref{hol:ot:lem:endpoints}).\\
$\Upsilon_{\mathfrak h}$, $\mathsf g^\pm$, $\delta_n$ & Endpoint--gap
polynomial \eqref{hol:ot:eq:upsilon}; upper and lower gaps; the break scale
$v(a_{n+1})-v(a_n)$.\\
\bottomrule
\end{tabularx}
\end{center}

\begin{center}
\begin{tabularx}{\textwidth}{@{}l X@{}}
\toprule
Symbol & Meaning\\
\midrule
\multicolumn{2}{@{}l}{\emph{Unit-dilation part,
Sections~\ref{hol:ud:sec:setting}--\ref{hol:ud:sec:mixed}}}\\
$\lambda_1,\ldots,\lambda_m$; $\zeta$, $\tau$ & Multipliers of a mixed
operator; residue and tail of a unit \eqref{hol:ud:eq:unitdecomp}.\\
$\mathcal Y(n)$, $R_j$; $\Sigma$, $r_{j,\gamma}$ & An exponential polynomial
$\sum_jR_j(n)\lambda_j^n$; its support envelope and coefficient polynomials
\eqref{hol:ud:eq:coefficientep}.\\
$\mathsf M$, $b$; $\beta_b$, $\varpi_b$ & Period and residue class; eventual
intercept and slope \eqref{hol:ud:eq:affine}.\\
$\mathrm{Sh}$ & The forward shift of sequences.\\
$\varkappa_{\ell,r,i}$, $d_L$, $s_0$, $c_j(n)$, $R_{j,\ell}$ & Coefficients of
\eqref{hol:ud:eq:L}, their $z$-degree bound, the least active shift, the
recurrence coefficients \eqref{hol:ud:eq:cj} and shift polynomials
\eqref{hol:ud:eq:Rjl}.\\
$\mathsf h$, $\bar e$, $\bar n$ & A polynomial right side, a bound on its
degree, an output degree.\\
$\operatorname{rs}(\lambda_1,\ldots,\lambda_m)$ & The relative scale
(Definition~\ref{hol:ud:def:rs}).\\
$\Delta^{\mathrm{top}}$ & The maximal proper convex subgroup when $\Gamma$ has
an order unit (Proposition~\ref{hol:ud:prop:top}).\\
$G_\mu$ & $\sum_nt^{n^2\mu}z^n=\Theta_{p^2}(pz)$, $p=t^\mu$; $G_1$ over
$\C((t^\Q))$ is the series $G$ of Corollary~\ref{hol:nl:cor:theta}.\\
\midrule
\multicolumn{2}{@{}l}{\emph{Newton-rigidity part,
Sections~\ref{hol:nr:sec:setting}--\ref{hol:nr:sec:boundary}}}\\
$\mathbf m_r(T)$; $\chi_{\mathfrak h}$ & The monomial jet curve
$(1,T,\ldots,T^r)$; $\mathfrak h(\mathbf m_r(T))$, so that
$\chi_{\mathfrak h_P}=\chi_P$.\\
$\mathfrak s$, $\mathfrak h_1$ & $Y_0Y_2-Y_1^2$ and the cofactor in
$\mathfrak h=\mathfrak s^e\mathfrak h_1$ (Remark~\ref{hol:nr:rem:conic}).\\
$\operatorname{Pol}_{\mathfrak h}$, $a_{\mathfrak h}$, $b_{\mathfrak h}$ & The
first polar \eqref{hol:nr:eq:polar} and its factorization
\eqref{hol:nr:eq:polarfactor}.\\
$\mathsf H_3$, $\mathsf Q_4$, $\mathsf B(m,n)$ & The quartic of order three,
the quadratic of order four and its pair kernel
\eqref{hol:nr:eq:H3}--\eqref{hol:nr:eq:pairkernel}.\\
$\mathsf A_1,\mathsf A_2$, $\mathsf n$, $\mathsf u_f$ & The auxiliary quadratics
and common-root parameter of Theorem~\ref{hol:nr:thm:H3};
$\vartheta f/f$.\\
$F_{\mathrm{sp}}$ & The sparse series \eqref{hol:nr:eq:sparse}.\\
$\mathbb L$; $\mathfrak M$, $\mathfrak M_1$, $\mathsf m$ & A field of external
constants; the fields and minimal polynomial of
Theorem~\ref{hol:nr:thm:systems}.\\
\bottomrule
\end{tabularx}
\end{center}

\begin{center}
\begin{tabularx}{\textwidth}{@{}l X@{}}
\toprule
Symbol & Meaning\\
\midrule
\multicolumn{2}{@{}l}{\emph{Theta-hierarchy part,
Sections~\ref{hol:th:sec:setting}--\ref{hol:th:sec:scope}}}\\
$\mathcal T_{t^\mu}$ & \eqref{hol:td:eq:series} with $p=t^\mu$, $\mu>0$ real.\\
$\operatorname{dord}(f)$ & The least $r$ with $f,\ldots,D_z^rf$ dependent
(Definition~\ref{hol:th:def:dord}).\\
$\mathfrak K_0\subseteq\mathfrak K_1$, $\partial$, $\mathrm d$, $\mathsf s$ &
Differential fields, their derivation, the universal derivation, an element with
$\partial\mathsf s=1$.\\
$\mathsf P_i$, $\mathsf c_{ij}$, $\mathsf W$, $\mathsf U_{ir}$, $\varepsilon_i$ &
Compression weights, grid coefficients, the wedge, jet parameters, formal
parameters (Theorem~\ref{hol:th:thm:compression}).\\
$\kappa_{\mathbf e}$, $\kappa'_{\mathbf e}$, $\delta^{(0)}$, $\delta^{\mathrm{cert}}$ &
Growth constants and thresholds of Proposition~\ref{hol:th:prop:threshold}.\\
$d_N$, $\mathfrak L_N$, $u_N$ & The compositum and its exact-order element
(Theorem~\ref{hol:th:thm:exactdimension}).\\
$\mathfrak V$; $\mathcal A_{\mathrm{DA}}$, $\mathcal D_{\mathrm{DA}}$ & A
$\Q$-independent set of positive reals; the differentially algebraic entire
series and their fraction field.\\
$\mathsf x$, $\gamma^{\mathcal U}$ & $-\vartheta_u\psi_1-2\mathcal Q_1(p^2)$; the
radius-one Gauss valuation on $\mathcal U_\mu$.\\
\midrule
\multicolumn{2}{@{}l}{\emph{Single-scale part,
Sections~\ref{hol:fh:sec:setting}--\ref{hol:fh:sec:realizations}}}\\
$\mathsf b_n$; $\mathsf N_A$, $\mathcal G_A$ & Rapid heights; the series of
Definition~\ref{hol:fh:def:rapid}.\\
$\vartheta_p$ & $p\,\mathrm d/\mathrm dp$.\\
$\mathbf v_n$, $A_{\mathbf v}$, $A_{\mathrm{fin}}$, $\mathsf c$, $\mathsf c_0$,
$\mathsf V$, $d_{\mathrm{pat}}$, $\boldsymbol\ell_j$ & Patterns, pattern classes,
finite corrections, pattern matrix and rank, kernel basis
(Section~\ref{hol:fh:sub:relations}).\\
$\bar d$, $\bar h$, $\mathbf T$, $\mathsf b_{\mathbf T}$; $\mathfrak P$,
$\mathsf W_{\mathfrak P}$ & Top degree, $p$-degree, target multiset, band start;
multihomogeneous part and its polarization.\\
$\mathfrak x_A$, $\partial_{\mathrm{BM}}$ & $\sum_{n\in A}\omega^{-\mathsf b_n}$;
the Berarducci--Mantova derivation.\\
$\mathcal B^{\mathrm{bd}}_\Gamma(k)$, $\mathcal F^{\mathrm{bd}}_\Gamma(k)$ &
Series with support bounded above; its fraction field.\\
$\mathcal O_{\langle1\rangle}$ & $\{0\}\cup\{x\in\No[i]^\times:v(x)\ge-m\}$.\\
$M_\Omega$, $\mathfrak z_A$, $\mathsf I_A(a)$ & $\omega^\Omega$; the omnific codes;
the floor profile.\\
\midrule
\multicolumn{2}{@{}l}{\emph{Polynomial-composition part,
Sections~\ref{hol:pc:sec:setting}--\ref{hol:pc:sec:arithmetic}}}\\
$\lc_z(b)$ & Leading coefficient in $z$ of $0\ne b\in\KK[z]$.\\
$\mathcal E^{(\delta)}$ & Series with $v(b_n)-n\delta\to+\infty$ cofinally.\\
$w_L$, $\beta_L$, $I_L$, $\chi_L$ & Shift, corner value, exact and residue
indicial polynomials of an operator \eqref{hol:pc:eq:indicial}.\\
$\delta_{\mathrm{ref}}$ & Reference scale of Lemma~\ref{hol:pc:lem:convex}.\\
$\bar s$, $\bar m$, $c_*$, $\mathsf G$, $L_\ell$, $L'_\nu$, $P_\nu$, $e_\nu$ &
Data of \eqref{hol:pc:eq:main}.\\
$\operatorname{cw}(\mathbf a)$ & Composition weight \eqref{hol:pc:eq:weight}.\\
$w_{\mathrm{top}}$, $\beta_{\mathrm{top}}$, $w_{\mathbf a}$, $\beta_{\mathbf a}$,
$\gamma_{\mathrm{top}}$, $\gamma_{\mathbf a}$ & Shifts, corner values and
Gauss bounds in the proof of Theorem~\ref{hol:pc:main:composition}.\\
$\operatorname{Rsn}$, $B_{\mathrm{cmp}}$, $B_{\mathrm{lin}}$, $B_{\mathrm{mat}}$ &
Resonances and degree bounds (Section~\ref{hol:pc:sub:degree},
Theorem~\ref{hol:pc:thm:matrix}).\\
$\mathfrak o$, $\mathfrak O_{\mathfrak o,\Gamma}$, $\Pi_{k,\Gamma}$,
$\operatorname{ct}$, $\mathfrak R$ & $\Z$ or $\Z[i]$; omnific coefficient ring;
its negative part; constant term; a coefficient ring of
Theorem~\ref{hol:pc:thm:linarith}.\\
$\mathrm{Pr}_j$, $\mathrm{Ann}_{d,N}$, $\operatorname{wd}(p)$ & Euler
projectors, annihilator and weighted degree of the encoder.\\
\bottomrule
\end{tabularx}
\end{center}

\begin{center}
\begin{tabularx}{\textwidth}{@{}l X@{}}
\toprule
Symbol & Meaning\\
\midrule
\multicolumn{2}{@{}l}{\emph{Recurrence part,
Sections~\ref{hol:rc:sec:setting}--\ref{hol:rc:sec:surreal}}}\\
$y_n$ & A recurrence sequence with values in $\KK$ or $\No[i]$.\\
$\langle\zeta\rangle$, $\mathsf M_{\mathrm{tor}}$ & The group generated by the
residues of unit bases; the exponent of its torsion subgroup
\eqref{hol:rc:eq:torperiod} (divisible by $\mathsf M$).\\
$\mathfrak A$, $\operatorname{ev}$ & The observable algebra $k[n,\zeta_1^n,\ldots,\zeta_m^n]$ of
\eqref{hol:rc:eq:algebra}; the evaluation map onto it from
$k[U,X_1,\ldots,X_m]$.\\
$\mathcal X$, $F_a$, $\mathsf T$ & A coefficient-algebraic set, its defining
conditions, a prescribed support set
(Definition~\ref{hol:rc:def:conditions}).\\
$\deg_\omega$; $\mathsf d_b$, $\mathsf e_b$ & $-v$, the greatest Conway
exponent; slope and intercept in \eqref{hol:rc:eq:degree}.\\
$\Pi$; $[\omega^\alpha]x$ & The purely infinite ideal, $\Oz=\Z\oplus\Pi$; a
Conway coefficient.\\
$\mathbf b$, $\mathfrak w_{\mathbf b}$ & A real sequence and its omnific code
\eqref{hol:rc:eq:code}.\\
$\lambda_-$, $\mathrm{Cat}_k$ & The infinitesimal root of $X^2-\omega X+1$;
Catalan numbers.\\
$\mathbf A$, $\operatorname{mv}_j(n)$ & An invertible matrix; the least
valuation of a nonzero $j\times j$ minor of $\mathbf A^n$.\\
$\mathfrak G$, $\Delta_{\mathfrak G}$ & A torsion-free group of dilation
parameters; the convex subgroup generated by $v(\mathfrak G)$.\\
\bottomrule
\end{tabularx}
\end{center}

\section{Proof-dependency and hypothesis audit}\label{hol:app:audit}

The logical dependencies are short. Differential rigidity uses the
integer-evaluation lemma, the coefficient conversion, and the escape chain
with its exclusion criterion. Dilation necessity uses the geometric-orbit
lemma or the unit-orbit theorem, the dilation conversion, and the same
criterion. Dilation sufficiency uses only the partial theta identity and the
support certificate. The mixed extension uses the bivariate orbit lemma. The
several-variable theorem uses countable polynomial avoidance, finite-block
regrouping and the one-variable theorem. The infinite-rank application uses a
direct strong-entireness proof and the completed classification. The
nonlinear part uses none of these linear arguments: first-order rigidity uses
the Hahn-family algebra (with the support calculus of
Section~\ref{hol:sec:support}), the large-active-degree lemma, the corner
identity and the order-one nondegeneracy of the corner; the higher-order
criterion replaces the last step by the hypothesis $\chi_P\ne0$; the Euler
theorem uses the same algebra in several variables with finite bounded-weight
sets. The coefficient-field part uses neither the escape chain nor the corner:
its all-order theorem uses the linearized coefficient recurrence, the
valuation-rank inequality, the principal convex bound and the exterior
obstruction (or the coarsened reduction), with the cofinal criterion of
Proposition~\ref{hol:cf:prop:criterion}; its mixed statements add the imported
Skolem--Mahler--Lech theorem. The theta-descent part uses none of these for
its existence direction: the exact domain uses the quadratic support
certificate and the nonincreasing-exponent obstruction; the equation of order
three uses the classical triple product and elliptic identity, proved as a
formal identity over $\Q$ and transferred by ring homomorphisms; the lower
bound uses the Laurent ring, the substitution homomorphism and three
difference-field obstructions. Its necessity direction is
Theorem~\ref{hol:cf:main:entire}\,(a). The order-three-threshold part uses
the all-scale criterion, a cofinal extension to the divisible hull, the
Hahn-family algebra and its residue homomorphism, local finiteness of active
indices and late breaks, the exterior reduction and the finite endpoint lemma;
its least-order statement adds the witness $\mathcal T_p$ and
Theorem~\ref{hol:cf:main:entire}\,(a), and its theta results the exterior form
of the cleared equation and a rational Riccati reduction. The unit-dilation
part uses the support calculus, the imported Skolem--Mahler--Lech theorem
through Lemma~\ref{hol:cf:lem:exppoly}, a common support envelope and period,
the least coefficient sequence that is not an identity, the eventual order of
finitely many affine functions, the escape chain with its exclusion
criterion, and, for existence, the partial theta identity and domain. The
Newton-rigidity part uses the machinery of the order-three-threshold part up
to the exterior form, then the first polar with its double factor, the
eventual side of integral roots through a $\Q$-linear functional and the
secant Jacobian; its exact classifications use only formal power series over a
field of characteristic zero. The theta-hierarchy part uses the coefficient
values of $\mathcal T_p$, Gauss multiplicativity, a unique dominant summand, the
K\"ahler criterion with a finite grid, the separant argument and, for the
continuum, a Hamel basis; its route to order three uses the Laurent ring and the
shift with a Gauss valuation. The single-scale part uses height separation,
polarization, Zariski density and linear algebra over $k(p,z)$, constant
extension, and, for its realizations, the positive support monoid, the
Berarducci--Mantova derivation and the normal-form definition of $\Oz$. The
polynomial-composition part uses the entire-support criterion, Gauss
multiplicativity and large active degrees, then the exact pullback and
operator profiles, integer convexity and superlinear escape; its degree bounds
use exact leading terms, its module theorem the Smith normal form and the
constant-term retraction, and its encoding Euler projectors. The recurrence
part uses the support envelope of the unit-dilation part and finite binomial
expansions for its observable algebra, Hilbert's basis theorem and a torsion
period with the imported Skolem--Mahler--Lech theorem for compression, the
algebraic closedness of $\No[i]$, a Jordan form and the periodic-affine
valuations for its surreal and matrix profiles, and only normal-form support
arithmetic for the coding theorem.

No unproved ring-theoretic conclusion from another report in this collection
is used. The real-valued coarsening of~\cite[Lemma~9.1]{repoentire} is cited
and not reproved, and no proof here depends on it. Higman's theorem is an
external standard ingredient; alternatively one may import the classical
Hahn--Neumann support lemma directly. The surreal interpretation imports the
normal-form field embedding, while every substantive function statement is
already proved in the abstract Hahn field. The coefficient-field part imports
one further external theorem without proof, the Skolem--Mahler--Lech theorem
in the form of Bell's Theorem~1.1~\cite{bell}, and only its statements with
several dilations use it. The theta-descent part imports the classical Jacobi
triple product~\cite{dlmftheta} and the constancy of an elliptic function
without poles; it proves the elliptic identity it needs. The unit-dilation
part imports the Skolem--Mahler--Lech theorem~\cite{lech,belldoc} for
exponential polynomials with at least two bases, and proves the finite-word
lemma that Lemma~\ref{hol:lem:support} imports from Higman
(Lemma~\ref{hol:ud:lem:words}). The Newton-rigidity part imports nothing
beyond the Jacobian criterion for algebraic independence in characteristic
zero, which it proves in one line. The theta-hierarchy part imports the theta
product and the Tate cubic identity, as source~13 did, and uses the axiom of
choice for a Hamel basis. The single-scale part imports the normalized
Berarducci--Mantova derivation~\cite{bm}, the normal-form identification and
definition of $\Oz$, and cites Gr\"onwall's single-series criterion without
using it. The polynomial-composition part imports the
Davis--Putnam--Robinson--Matiyasevich theorem~\cite{matiyasevich} and Sun's
eleven-unknown theorem~\cite{sun} for its two undecidability corollaries, and
cites the complex entire-Mahler lemma~\cite[Lemma~6]{bcr} without using it.
The recurrence part imports the Skolem--Mahler--Lech theorem~\cite{lech,belldoc},
Hilbert's basis theorem, the Hahn--Neumann support lemma~\cite{neumann,bagayoko}
and the algebraic closedness of $\No[i]$; it cites the set-sized-quotient
theorem~\cite{repoosq} only to explain the absence of a contradiction.

The potentially dangerous hypotheses, collected for review:
\begin{itemize}
\item Supports and the value group are sets; all operators are finite sums.
Power-series degrees and derivative orders are ordinary nonnegative integers;
dilation powers may be negative, as may the linearization offset $\theta$.
Formal Laurent series have only finitely many negative powers.
\item $\Gamma\ne0$, set-sized, ordered abelian. \emph{Divisibility is not
assumed} except in Corollary~\ref{hol:cor:periodic}, the refined threshold
\eqref{hol:eq:refinedthreshold}, Theorem~\ref{hol:thm:coarse}\,(b),
and Example~\ref{hol:ex:sparse}, each of which states it, and in
Lemma~\ref{hol:ot:lem:breaks}, which is applied, in the order-three-threshold
and Newton-rigidity parts, only after the cofinal extension of
Lemma~\ref{hol:ot:lem:extension} to the divisible hull.
\item Series are ordinary $\N$-indexed power series in one variable, except
in the several-variable Sections~\ref{hol:sec:multi}
and~\ref{hol:nl:sub:euler}, and in the formal Laurent fraction field
$\FE$. The main workspace rigidity equations have rational-function
coefficients in the variable, with arbitrary Hahn scalar coefficients.
The relative coefficient-field theorem
(Theorem~\ref{hol:cf:thm:coefficients}\,(a)) also allows coefficients in
$k_0((z))$; that generality is used in the joint-independence criterion.
\item The residue field is characteristic zero and trivially valued.
Algebraic closedness is not used. It is $\C$ in
Sections~\ref{hol:sec:intro}--\ref{hol:sec:certificates} and an arbitrary
characteristic-zero field in the nonlinear, coefficient-field, theta-descent,
order-three-threshold, unit-dilation, Newton-rigidity, single-scale,
polynomial-composition and recurrence parts (with $\R$ or $\C$ for their
omnific rings), and
$\R$ or $\C$ in the
theta-hierarchy part, which also fixes $\Gamma=\R$; the latter
generality is not transferred to the statements of the linear theorems
(Convention~\ref{hol:nl:conv:field}), although the unit-dilation part contains
Theorem~\ref{hol:main:ode}, the negative direction of
Theorem~\ref{hol:main:q} and Theorem~\ref{hol:thm:mixed} over every such field
(Remark~\ref{hol:ud:rem:contains}).
Characteristic zero is essential (Example~\ref{hol:nl:ex:charp}).
\item The single parameter $q$ in the main dilation classification is
nonzero and nontorsion. In the mixed-jet statements every $\lambda$ is
nonzero, and distinct parameters have non-root-of-unity quotients; an
individual parameter may be $1$; the same holds for the multipliers of the
unit-dilation part. Section~\ref{hol:sec:torsion} and
Section~\ref{hol:ud:sub:examples} explain the failure when this quotient
condition is dropped.
\item The derivative is $d/dz$, not a scalar surreal derivation.
\item Entireness is always relative to one fixed workspace, never to
$\No[i]$. Strong summability requires both well-ordering and local finiteness.
\item The conclusions of
Sections~\ref{hol:sec:intro}--\ref{hol:sec:certificates} concern
\emph{linear} equations. For every $\Gamma\ne0$, nonlinear differential
equations are treated only in order one, in higher order under
$\chi_P\ne0$, and for positive-weight Euler relations (the nonlinear part),
in order two (the order-three-threshold part), and in order three with total
jet degree at most three (the Newton-rigidity part); for explicit lacunary
series all orders are excluded at one order-unit scale (the single-scale part),
while for $\Gamma=\R$ minimal orders are unbounded (the theta-hierarchy part).
Linear mixed equations in
derivatives and finitely many dilations are treated for every $\Gamma$ (the
unit-dilation part), and so are equations in which one polynomial argument of
degree at least two strictly dominates the composition weights (the
polynomial-composition part). The recurrence part concerns ordinary
constant-coefficient recurrences indexed by $\N$, not equations for
functions, apart from its re-derived mixed statements.
Nonlinear differential transcendence in all orders, and algebraic
independence of derivatives at dilations with pairwise quotients that are not
roots of unity, are proved only when $\Gamma$ has no order unit, or for a
single series with infinite coefficient-value rank (the coefficient-field
part). No universal all-order rigidity theorem is claimed for groups with an
order unit, and none can hold: for them the theta-descent part gives a
nonpolynomial entire series with an algebraic differential equation of
order three, and three is then the least order. No general mixed-jet
independence theorem is claimed when
a quotient is a root of unity.
\end{itemize}

Two boundary behaviours deserve separate emphasis because they point in
opposite directions. The most delicate equality is in
Corollary~\ref{hol:cor:periodic}: weakly decreasing leading exponents already
obstruct strong summability, so the excluded region is closed, and
Proposition~\ref{hol:prop:expdomain} shows the bound is attained. The most
delicate strictness is in Theorem~\ref{hol:thm:coarse}: equal \emph{coarse}
values need not obstruct summability in $\Gamma$, so that excluded region must
be open, and Theorem~\ref{hol:thm:thetadomain} with
Section~\ref{hol:sub:coarseboundary} shows why.

Deliberate exclusions, recorded so that no reader supplies them by analogy:
finite-order $q$ is excluded from the main classification and has genuine
periodic-support counterexamples; linear mixed operators at unit valuation
are classified for every $\Gamma$ (Corollary~\ref{hol:ud:cor:unitrigidity}),
but multiplier sets with a root-of-unity quotient, prescribed operators at a
cofinal quotient and simultaneous systems are not. For groups with an order
unit, the universal nonlinear rigidity results cover orders one and two,
order three in total jet degree at most three, the nondegenerate higher-order
class and positive-weight Euler relations; the infinite-coefficient-value-rank
criterion gives additional all-order results for individual series, including
Theorem~\ref{hol:cf:thm:rankone}, while Theorem~\ref{hol:td:main:boundary}
shows that universal differential algebraic rigidity fails for them, with
least order three by Theorem~\ref{hol:ot:main:second}; which series have
minimal order three is not decided. The nonlinear part keeps the
non-claims of Section~\ref{hol:nl:sub:nonclaims}, the no-order-unit results
keep those of Section~\ref{hol:cf:sub:nonclaims}, the theta-descent part
keeps those of Section~\ref{hol:td:sub:nonclaims}, the
order-three-threshold part those of Section~\ref{hol:ot:sub:nonclaims}, the
unit-dilation part those of Section~\ref{hol:ud:sub:nonclaims}, the
Newton-rigidity part those of Section~\ref{hol:nr:sub:nonclaims}, the
theta-hierarchy part those of Section~\ref{hol:th:sub:nonclaims}, the
single-scale part those of Section~\ref{hol:fh:sub:nonclaims}, and the
polynomial-composition part those of Section~\ref{hol:pc:sub:nonclaims}
(critical weights, several top products and matrix differential pullbacks are
not treated), and the recurrence part those of
Section~\ref{hol:rc:sub:nonclaims} (no effective exceptional zeros, no
classification of omnific membership, and a coding theorem that stores rather
than computes); no
complete algorithm for arbitrary
surreal comparisons or cofinality decisions is claimed; the finite checks
prove no infinite theorem. The full classifications and the nonlinear,
coefficient-field, theta-descent, order-three-threshold, unit-dilation,
Newton-rigidity, theta-hierarchy, single-scale, polynomial-composition and
recurrence theorem packages remain without Lean verification or independent
refereeing.
The formalization ledger records the proved support and escape lemmas,
exponential and partial-theta domains, polynomial-orbit valuations,
constant-point avoidance, inward stability and torsion covariance,
including the partial-theta directions of the dilation classification.

\section{Reproducing the document and the finite checks}
\label{hol:app:reproduce}

The directory contains this standalone \LaTeX{} source, its compiled PDF, a
README, twenty source audit documents at the top level, twenty-two files under
\texttt{code/} (the thirteen verification programs of the sources merged here,
the build scripts of sources~10, 11, 13, 17, 19 and~20, and the
\texttt{Makefile}s of sources~12, 15 and~16) and thirty-two under
\texttt{data/} (the thirteen recorded outputs, the source build, layout and
visual reports, and the requirements files of sources~11, 13, 14, 15, 17
and~19), all belonging to the thirteen sources merged into this text. No source
document from the
inspected repository is redistributed, and no source manuscript is shipped. The bibliography is embedded: no BibTeX run, external
graphics, shell escape, or repository checkout is required, and a conventional
TeX Live or MiKTeX installation with the named packages suffices.

\begin{verbatim}
latexmk -pdf -interaction=nonstopmode -halt-on-error article.tex
python code/08-finite-recurrences-verify.py --output rerun-08.txt
python code/09-entire-hahn-holonomic-verification.py --output rerun-09.txt
python code/10-nonlinear-rigidity-verify.py --output rerun-10.json
pip install -r data/11-coarsening-differential-rigidity-requirements.txt
python code/11-coarsening-differential-rigidity-verify.py --output rerun-11.json
python code/12-coefficient-field-rigidity-verify.py --output rerun-12.json
pip install -r data/13-theta-descent-requirements.txt
python code/13-theta-descent-verify.py --precision 128 --output rerun-13.json
pip install -r data/14-order-three-threshold-requirements.txt
python code/14-order-three-threshold-verify.py
python code/18-unit-dilation-verify_finite.py --output rerun-18.json
pip install -r data/19-newton-rigidity-requirements.txt
python code/19-newton-rigidity-verify.py
python code/16-theta-hierarchy-verify.py --precision 128 --output rerun-16.json
pip install -r data/17-single-scale-requirements.txt
python code/17-single-scale-verify.py --output rerun-17.json
pip install -r data/15-polynomial-composition-requirements.txt
python code/15-polynomial-composition-verify.py --output rerun-15.json
python code/20-recurrences-verify.py > rerun-20.txt
\end{verbatim}

The \verb|--output| argument is required by the scripts of sources~08, 10,
12 and~13, so that a rerun need not overwrite the record distributed with the
report, and optional for those of sources~09 and~11; the scripts of
sources~10 and~12 refuse an existing output path (source~12's unless
\verb|--force| is given), source~13's overwrites one, and source~11's writes
\texttt{verification.json} in the current directory when no path is given.
Source~14's takes no arguments and always writes
\path{data/verification.json} next to its own \texttt{code/} directory: here a
new file beside the delivered record, overwritten by the next run, so it
should be run on a copy. Source~19's script behaves the same way and writes
the same file, \path{data/verification.json}; source~18's writes only the
path given by \verb|--output|, overwriting it, and otherwise only prints.
Source~16's writes \path{verification.json} in the current directory unless
\verb|--output| is given, and source~17's writes
\path{verification_results.json} beside itself in \texttt{code/}; both
overwrite, so pass new paths. Source~15's, like source~16's, writes
\path{verification.json} in the current directory unless \verb|--output| is
given, overwriting it. Source~20's takes no arguments and only prints its
record, so redirect it to a new file. All use exact integer and rational
arithmetic,
those of sources~11, 13, 14, 15, 17 and~19 through SymPy, and exit with a nonzero
status if a check fails. The build script
\path{code/10-nonlinear-rigidity-build.py} was written for source~10's own
manuscript: run here it would typeset this report, leave auxiliary files in
the directory and write \texttt{data/build\_report.json}. The shipped
\path{data/10-nonlinear-rigidity-build_report.json} and
\path{data/10-nonlinear-rigidity-layout_audit.json} describe the 23-page
source manuscript, not this report; use \texttt{latexmk} as above, on a copy.
Likewise \path{code/11-coarsening-differential-rigidity-build.py} looks for an
\texttt{article.tex} in \texttt{code/} and fails there, and
\path{code/12-coefficient-field-rigidity-Makefile} names the delivery paths
\texttt{article.tex} and \texttt{verify.py} and would rebuild
\texttt{article.pdf} in place; the shipped
\path{data/11-coarsening-differential-rigidity-build_audit.json} and
\path{data/12-coefficient-field-rigidity-build_report.json} describe the 20-
and 23-page source manuscripts, whose PDFs and sources, hashed in the first,
are not shipped. Source~13's \path{code/13-theta-descent-build.py} typesets
an \texttt{article.tex} in its own directory \texttt{code/}, where there is
none, and its \verb|--verify| option would call a \texttt{verify.py} there; run
here it creates \texttt{code/build/}, writes a failed first-pass log into it
and exits with an error. In its original delivery layout it writes only into
\texttt{build/}. The shipped \path{data/13-theta-descent-build_report.json}
describes the 21-page source manuscript, whose source and PDF, hashed there,
are not shipped. Likewise \path{data/14-order-three-threshold-build_report.json}
describes source~14's 26-page A4 manuscript, whose source and PDF, hashed
there, are not shipped; source~14 shipped no build script. Source~18 shipped
none either and no build record. Source~19's
\path{code/19-newton-rigidity-build.py} looks for an \texttt{article.tex}
in \texttt{code/} and fails there, and its \verb|--verify| option would call
a \texttt{code/verify.py} relative to \texttt{code/}, which does not exist;
\path{data/19-newton-rigidity-build_audit.json} describes the 24-page source
manuscript, whose source and PDF, hashed there, are not shipped. Source~16's
\path{code/16-theta-hierarchy-Makefile} and source~17's
\path{code/17-single-scale-build.py} were written for their delivery layouts
(Section~\ref{hol:sec:verification}); do not use them here, and the build and
layout records of both describe their 25- and 24-page manuscripts, which are
not shipped. The same holds for source~15's
\path{code/15-polynomial-composition-Makefile}, which would rebuild an
\texttt{article.pdf} in the current directory and whose \texttt{clean} target
deletes auxiliary files there, and for
\path{data/15-polynomial-composition-BUILD_REPORT.json}, which describes its
24-page manuscript. Source~20's \path{code/20-recurrences-build.sh} runs
\texttt{pdflatex} on an \texttt{article.tex} in its own directory and then
overwrites \texttt{verification.txt} there by redirection; here it stops at the
first \texttt{pdflatex} call and leaves \texttt{texput.log} in \texttt{code/}
(Section~\ref{hol:sec:verification}), so do not run it.

\clearpage
\begin{thebibliography}{99}
\addcontentsline{toc}{section}{References}
\raggedright

\bibitem{stanley}
R. P. Stanley,
\emph{Differentiably finite power series},
European Journal of Combinatorics \textbf{1} (1980), no.~2, 175--188.
Author-hosted copy:
\url{https://math.mit.edu/~rstan/pubs/pubfiles/45.pdf}.
Theorem~1.5 (printed page 176) gives the D-finite/P-recursive equivalence.
The finite coefficient conversion is reproved in this report for its stated
characteristic-zero Hahn coefficient field.

\bibitem{ramis}
J.-P. Ramis,
\emph{About the growth of entire functions solutions of linear algebraic
$q$-difference equations},
Annales de la Facult\'e des Sciences de Toulouse, Math\'ematiques,
S\'erie~6, \textbf{1} (1992), no.~1, 53--94.
\url{https://www.numdam.org/item/AFST_1992_6_1_1_53_0/}.

\bibitem{garoufalidis}
S. Garoufalidis,
\emph{The degree of a $q$-holonomic sequence is a quadratic quasi-polynomial},
Electronic Journal of Combinatorics \textbf{18} (2011), no.~2, Paper~P4.
Inspected manuscript: arXiv:1005.4580v4.
\url{https://arxiv.org/abs/1005.4580}.

\bibitem{divizio}
L. Di Vizio,
\emph{An ultrametric version of the Maillet--Malgrange theorem for nonlinear
$q$-difference equations},
Proceedings of the American Mathematical Society \textbf{136} (2008), no.~8,
2803--2814. Inspected manuscript: arXiv:0709.2464v2.
\url{https://arxiv.org/abs/0709.2464}.

\bibitem{andrewswarnaar}
G. E. Andrews and S. O. Warnaar,
\emph{The product of partial theta functions}.
Author-hosted manuscript; the definition on its first page is the series used
here, up to a sign convention.
\url{https://people.smp.uq.edu.au/OleWarnaar/pubs/Partial-thetas.pdf}.

\bibitem{higman}
G. Higman,
\emph{Ordering by divisibility in abstract algebras},
Proceedings of the London Mathematical Society (3) \textbf{2} (1952),
326--336.
\url{https://doi.org/10.1112/plms/s3-2.1.326}.

\bibitem{mueller}
J. M\"uller and A. Strohmaier,
\emph{The theory of Hahn meromorphic functions, a holomorphic Fredholm theorem
and its applications},
Analysis \& PDE \textbf{7} (2014), no.~3, 745--770.
\url{https://arxiv.org/abs/1205.0236}.
\url{https://doi.org/10.2140/apde.2014.7.745}.

\bibitem{bm}
A. Berarducci and V. Mantova,
\emph{Surreal numbers, derivations and transseries},
Journal of the European Mathematical Society \textbf{20} (2018), 339--390.
\url{https://arxiv.org/abs/1503.00315}.
\url{https://doi.org/10.4171/JEMS/769}.

\bibitem{bg}
O. Bournez and Q. Guilmant,
\emph{Surreal fields stable under exponential, logarithmic, derivative and
anti-derivative functions},
arXiv:2211.08396 (2022).
\url{https://arxiv.org/abs/2211.08396}.

\bibitem{kks}
E. Kaplan, L. S. Krapp and M. Serra,
\emph{Decomposing the automorphism group of the surreal numbers},
arXiv:2509.22374v3.
Used for the Hahn-field and surreal normal-form conventions, not for its
automorphism results.
\url{https://arxiv.org/html/2509.22374v3}.

\bibitem{repo}
\repo\ project,
\emph{Repository documentation and research-report catalogue},
snapshot \texttt{a3124af79f66b8b9c196d76b4cbc5ac3938907c4},
consulted 22 September 2026.
\url{https://github.com/VladimirReshetnikov/Surreal}.
The inspection used the root README and \texttt{docs/README.md} at this
commit; these are project documentation, not independent published
validation.

\bibitem{repoentire}
\repo\ project,
\emph{Cofinality, factorization and scalar extension for entire Hahn functions
at arbitrary rank},
\texttt{docs/surcomplex/entire-functions-at-arbitrary-rank/}, same pinned
snapshot.
\url{https://github.com/VladimirReshetnikov/Surreal/tree/a3124af79f66b8b9c196d76b4cbc5ac3938907c4/docs/surcomplex/entire-functions-at-arbitrary-rank}.
Its Lemma~9.1, labelled \texttt{ent:lem:coarsening}, is the real-valued
coarsening cited in Remark~\ref{hol:rem:coarsening}; its Theorem~9.3 is the
B\'ezout statement referred to in Section~\ref{hol:sub:position}. Source~11
inspected its README at \texttt{465a54b} for the polynomial intersection
obstruction for noncofinal extensions and for coarsened support methods.

\bibitem{repodifferential}
\repo\ project,
\emph{Differential algebra and differential equations over the surreal and
surcomplex numbers},
\texttt{docs/surcomplex/differential-equations/}, same pinned snapshot.
\url{https://github.com/VladimirReshetnikov/Surreal/tree/a3124af79f66b8b9c196d76b4cbc5ac3938907c4/docs/surcomplex/differential-equations}.
Cited to distinguish the scalar surreal derivation from $D_z$.

\bibitem{reporankone}
\repo\ project,
\emph{Rank-one non-Archimedean analytic geometry inside the surcomplex
numbers},
\texttt{docs/surcomplex/rank-one-berkovich/}, same pinned snapshot.
\url{https://github.com/VladimirReshetnikov/Surreal/tree/a3124af79f66b8b9c196d76b4cbc5ac3938907c4/docs/surcomplex/rank-one-berkovich}.
Cited for the pre-existing quadratic-exponent rank-one example.

\bibitem{repofoundations}
\repo\ project,
\emph{Foundations for Surreal and Surcomplex Mathematics},
\texttt{docs/foundations-and-computation/foundations/}.
Its Example \texttt{found:ex:internal} is the series $\sum t^{n^2}z^n$ over
$\C((t^\Q))$; cited by this merge, not by source~10.

\bibitem{repotail}
\repo\ project,
\emph{Tail-spans and differential transcendence in surreal and surcomplex Hahn
fields},
\texttt{docs/surreal/tail-spans-and-differential-transcendence/}, at
source~10's pin \texttt{048b72cf7cbfc8ab246e4f73788c10460cb3f6e0}.
\url{https://github.com/VladimirReshetnikov/Surreal/tree/048b72cf7cbfc8ab246e4f73788c10460cb3f6e0/docs/surreal/tail-spans-and-differential-transcendence}.
Its README was inspected by source~10 for function-class and derivative
boundaries, and again by sources~11 and~12 at their pin \texttt{465a54b},
for its continuum-sized families and its all-order results.

\bibitem{mantovamatusinski}
V. Mantova and M. Matusinski,
\emph{Surreal numbers with derivation, Hardy fields and transseries: a
survey}, arXiv:1608.03413, inspected version~2, Section~2 on normal forms
(Section~2.2 by source~11).
\url{https://arxiv.org/html/1608.03413v2}.
Published in \emph{Contemporary Mathematics} \textbf{697} (2017), 265--290;
source~13 consulted the publication records and the preprint for normal
forms, and source~14 the preprint for normal-form and Hahn-field background.

\bibitem{hayman}
W. K. Hayman,
\emph{The local growth of power series: a survey of the Wiman--Valiron
method},
Canadian Mathematical Bulletin \textbf{17} (1974), no.~3, 317--358.
\url{https://doi.org/10.4153/CMB-1974-064-0}.
Publisher metadata and abstract consulted by source~10.

\bibitem{cano}
J. Cano and P. Fortuny Ayuso,
\emph{Power series solutions of non-linear $q$-difference equations and the
Newton--Puiseux polygon},
Qualitative Theory of Dynamical Systems \textbf{21} (2022), article~123.
\url{https://doi.org/10.1007/s12346-022-00656-0}.
Introduction and Section~2 (initial and indicial polynomials) consulted by
source~10.

\bibitem{huluan}
P.-C. Hu and Y.-Z. Luan,
\emph{Non-Archimedean meromorphic solutions of functional equations},
arXiv:1311.5291v1 (2013).
\url{https://arxiv.org/abs/1311.5291}.
Field hypotheses and Theorems~1.1, 1.3 and~1.5 consulted by source~10;
introduction and field hypotheses by sources~11 and~12, as context only.

\bibitem{huyang}
P.-C. Hu and C.-C. Yang,
\emph{Meromorphic Functions over Non-Archimedean Fields},
Mathematics and Its Applications, Springer, Dordrecht, 2000.
\url{https://doi.org/10.1007/978-94-015-9415-8}.
Publisher metadata and description only; the book was not audited.

\bibitem{giansiracusa}
J. Giansiracusa and S. Mereta,
\emph{A general framework for tropical differential equations},
manuscripta mathematica \textbf{173} (2024), 1273--1304.
\url{https://doi.org/10.1007/s00229-023-01492-5}.

\bibitem{hurwitz}
A. Hurwitz,
\emph{Sur le d\'eveloppement des fonctions satisfaisant \`a une \'equation
diff\'erentielle alg\'ebrique},
Annales scientifiques de l'\'Ecole Normale Sup\'erieure (3) \textbf{6}
(1889), 327--332.
\url{https://doi.org/10.24033/asens.326}.
Bibliographic page and the attribution in~\cite{lku} consulted by source~11.

\bibitem{lku}
V. Laohakosol, K. Kongsakorn and P. Ubolsri,
\emph{Coefficients of differentially algebraic series},
Journal of the Australian Mathematical Society, Series~A \textbf{48} (1990),
402--412.
\url{https://doi.org/10.1017/S1446788700029955}.
Lemma~3 consulted by source~11.

\bibitem{kr}
C. Krattenthaler and T. Rivoal,
\emph{Arithmetic properties of the Taylor coefficients of differentially
algebraic power series}, arXiv:2502.09259v1 (2025).
\url{https://arxiv.org/abs/2502.09259}.
Abstract consulted by source~11, as context only.

\bibitem{ast}
R. Ait El Manssour, A.-L. Sattelberger and B. Teguia Tabuguia,
\emph{D-algebraic functions}, arXiv:2301.02512v3.
\url{https://arxiv.org/abs/2301.02512}.
Framework consulted by source~11 for context.

\bibitem{bell}
J. P. Bell,
\emph{A generalised Skolem--Mahler--Lech theorem for affine varieties},
arXiv:math/0501309v2 (2007).
\url{https://arxiv.org/abs/math/0501309}.
Its Theorem~1.1, the classical theorem over a field of characteristic zero,
is the form imported by source~12.

\bibitem{lech}
C. Lech,
\emph{A note on recurring series},
Arkiv f\"or Matematik \textbf{2} (1953), 417--421.
\url{https://doi.org/10.1007/BF02590997}.
Publisher metadata checked by source~12; the proof was not reconstructed.
Cited by source~18 as the primary source of the characteristic-zero theorem,
with its metadata verified, and by source~20 for the same theorem.

\bibitem{belldoc}
J. P. Bell,
\emph{The Skolem--Mahler--Lech theorem},
Documenta Mathematica, Extra Volume Mahler Selecta (2019), 173--178.
\url{https://ems.press/content/book-chapter-files/27398}.
Read by source~18, including the first two pages, for the precise
characteristic-zero statement and its exponential-polynomial formulation; a
different text from~\cite{bell}. Source~20 read and visually inspected the same
PDF.

\bibitem{neumann}
B. H. Neumann,
\emph{On ordered division rings},
Transactions of the American Mathematical Society \textbf{66} (1949),
202--252.
\url{https://doi.org/10.1090/S0002-9947-1949-0032593-5}.
Credited by source~18 for the support calculus, whose needed lemma it proves,
and by source~20 as the historical source of the support lemma, not audited
as a historical paper.

\bibitem{fuchsheintze}
C. Fuchs and S. Heintze,
\emph{On the growth of linear recurrences in function fields},
Bulletin of the Australian Mathematical Society \textbf{104} (2021), 11--20.
\url{https://doi.org/10.1017/S0004972720001094}.
Author version: \url{https://arxiv.org/abs/2006.11074}. Related work consulted
by source~18 through the author version and publisher metadata; not an input.
Cited by source~20 as related work whose effective constants it does not
improve.

\bibitem{gonshor}
H. Gonshor,
\emph{An Introduction to the Theory of Surreal Numbers},
London Mathematical Society Lecture Note Series~110, Cambridge University
Press, 1986.
\url{https://doi.org/10.1017/CBO9780511629143}.
Normal-form background for source~18; publisher metadata checked, the
monograph not reread.

\bibitem{derksen}
H. Derksen,
\emph{A Skolem--Mahler--Lech theorem in positive characteristic and finite
automata},
Inventiones Mathematicae \textbf{168} (2007), 175--224.
Author version: \url{https://arxiv.org/abs/math/0510583}.
Background for source~18's positive-characteristic direction, and for
source~20's; not an input to any proof.

\bibitem{pogudin2014}
G. Pogudin,
\emph{The primitive element theorem for differential fields with zero
derivation on the ground field}, arXiv:1409.3847.
\url{https://arxiv.org/abs/1409.3847}.
Consulted by source~16 as context for jet compression.

\bibitem{pogudin2018}
G. Pogudin,
\emph{A primitive element theorem for fields with commuting derivations and
automorphisms}, arXiv:1812.11375v3.
\url{https://arxiv.org/abs/1812.11375}.
Consulted by source~16 (HTML text) as context; not a substitute for
Theorem~\ref{hol:th:thm:compression}.

\bibitem{ostrowski}
A. Ostrowski,
\emph{\"Uber Dirichletsche Reihen und algebraische Differentialgleichungen},
Mathematische Zeitschrift \textbf{8} (1920), 241--298.
Read by source~17 in the English translation by E. C. Hansen, Y. Stone and
J. Wolfson, arXiv:2211.02088v1 (2022),
\url{https://arxiv.org/abs/2211.02088}; the passage before Theorem~11 in
Section~5 attributes the gap criterion to Gr\"onwall, whose paper was not
inspected.

\bibitem{lmfactor}
S. L'Innocente and V. Mantova,
\emph{A factorisation theory for generalised power series and omnific
integers}, arXiv:1710.07304v5 (2024).
\url{https://arxiv.org/abs/1710.07304}.
Used by source~17 for the normal-form definition of $\Oz$ only.

\bibitem{repobst}
\repo\ project,
\emph{Bounded-Support Hahn Arithmetic: cofinal descent, optimal support
thresholds and maximal transcendence},
\texttt{docs/surreal/transcendence-over-bounded-support/}. Its README was read
by source~17 at \texttt{bcac55a}; its \texttt{bst:thm:descent} contains the
descent of Lemma~\ref{hol:fh:lem:bounded}.

\bibitem{bgexp}
O. Bournez and Q. Guilmant,
\emph{Surreal fields stable under exponential and logarithmic functions},
arXiv:2201.08199v1 (2022).
\url{https://arxiv.org/abs/2201.08199}.
Its Corollary~2.7, the normal-form identification, consulted by source~12;
a different paper from~\cite{bg}.

\bibitem{repotate}
\repo\ project,
\emph{Hahn--Tate Uniformization at Arbitrary Valuation Rank},
\texttt{docs/surcomplex/hahn-tate-uniformization/}, at source~13's pin
\texttt{4cdeaec43ff1692ed2ad9f5bdf761a41872975a5}.
\url{https://github.com/VladimirReshetnikov/Surreal/tree/4cdeaec43ff1692ed2ad9f5bdf761a41872975a5/docs/surcomplex/hahn-tate-uniformization}.
Its theta series $\sum_{n\in\Z}(-1)^nq^{n(n-1)/2}u^n$, its theta domain
(\texttt{tate:thm:theta}) and its product and zeros
(\texttt{tate:prop:theta-product}) are cited in the theta-descent part;
source~13 inspected its README only.

\bibitem{repoodg}
\repo\ project,
\emph{Omnific Integers and Omnific--Diophantine Geometry},
\texttt{docs/surreal/omnific-diophantine-geometry/}. Its definition
\texttt{odg:def:rings} of $\Oz=\Z\oplus\Pi$ is cited by this merge for the
omnific integers of Corollary~\ref{hol:td:cor:omnific}.

\bibitem{bcr}
J. P. Bell, M. Coons and E. Rowland,
\emph{The rational--transcendental dichotomy of Mahler functions},
J. Integer Seq. \textbf{16} (2013), Article~13.2.10; arXiv:1210.2070.
\url{https://arxiv.org/abs/1210.2070}.
Its Lemma~6 is the complex entire-Mahler lemma credited by source~15.

\bibitem{nishioka}
K. Nishioka, \emph{Mahler Functions and Transcendence}, Lecture Notes in
Mathematics~1631, Springer, 1996. \url{https://doi.org/10.1007/BFb0093672}.
Cited by source~15 from publisher metadata only.

\bibitem{matiyasevich}
Yu. V. Matiyasevich, \emph{Diophantine representation of enumerable
predicates}, Math. USSR-Izv. \textbf{5} (1971), 1--28.
\url{https://doi.org/10.1070/IM1971v005n01ABEH001004}.

\bibitem{sun}
Z.-W. Sun, \emph{Further results on Hilbert's Tenth Problem}, Sci. China
Math. \textbf{64} (2021), 281--306; arXiv:1704.03504v7.
\url{https://arxiv.org/abs/1704.03504}.
Cited by source~15 only for undecidability in eleven integer unknowns.

\bibitem{faverjonroques}
C. Faverjon and J. Roques, \emph{Hahn series and Mahler equations: algorithmic
aspects}, J. London Math. Soc. (2) \textbf{110} (2024), e12945;
arXiv:2412.04928. \url{https://doi.org/10.1112/jlms.12945}.
Added by this merge from the single-dilation report's bibliography; not read
for this merge and not used.

\bibitem{reposdhs}
\repo\ project, \emph{A Single Dilation Recovers Hahn Support},
\texttt{docs/surcomplex/single-dilation-hahn-support/}. At \texttt{bcac55a} it
already discussed Mahler difference algebra and cited~\cite{faverjonroques}.

\bibitem{repoadr}
\repo\ project, \emph{Dilation Rigidity of Surreal Numbers and Omnific
Integers}, \texttt{docs/surcomplex/autonomous-dilation-relations/}. Its
\texttt{adr:lem:monomial} is the exponent-dilation twin of
Example~\ref{hol:pc:ex:critical}, and its README records Mahler's 1983
criterion through Nishioka--Nishioka.

\bibitem{repoosq}
\repo\ project, \emph{Set-Sized Quotients of the Omnific Integers},
\texttt{docs/surreal/set-sized-quotients-of-omnific-integers/}, and
\emph{Omnific Groups and Lattices},
\texttt{docs/surreal/omnific-groups-and-lattices/}.
Their READMEs were read by source~15 at \texttt{bcac55a}; nothing from them is
used. Source~20 read the first README at \texttt{efc5446} and cites its
universal quotient theorem, \texttt{osq:thm:universal}, only in
Remark~\ref{hol:rc:rem:quotients}, without auditing that report.

\bibitem{dlmftheta}
NIST Digital Library of Mathematical Functions, Chapter~20, \S20.5,
especially equations 20.5.3 and 20.5.9 (Jacobi's triple product), and
\S20.5(ii) (logarithmic derivatives).
\url{https://dlmf.nist.gov/20.5}. Accessed by source~13 on 23 September 2026.

\bibitem{dlmfelliptic}
NIST Digital Library of Mathematical Functions, Chapter~23, \S23.6
(relations to theta functions) and \S23.8 (trigonometric series and
products).
\url{https://dlmf.nist.gov/23.6}; \url{https://dlmf.nist.gov/23.8}. Accessed
by source~13 on 23 September 2026. Source~14 consulted \S23.3 (differential
equations of the Weierstrass function), \url{https://dlmf.nist.gov/23.3}, on
the same day.

\bibitem{dlmfpainleve}
NIST Digital Library of Mathematical Functions, Chapter~32, \S32.2
(differential equations; the definitions of Painlev\'e~I and~II).
\url{https://dlmf.nist.gov/32.2}. Accessed by source~14 on 23 September 2026.

\bibitem{woodbury}
\emph{Elliptic Curves and Modular Forms}, lecture notes hosted by
M.~Woodbury, 2010, \S16, printed page~32: the Tate equation and invariant
differential.
\url{https://www.mi.uni-koeln.de/~woodbury/research/ecnotes.pdf}.
Inspected by source~13 as a PDF image; the hosting attribution is not a claim
of authorship of every part.

\bibitem{ohyama}
Y. Ohyama,
\emph{Differential relations of theta functions},
Osaka Journal of Mathematics \textbf{32} (1995), 431--450.
\url{https://projecteuclid.org/journals/osaka-journal-of-mathematics/volume-32/issue-2/Differential-relations-of-theta-functions/ojm/1200786061.full}.
Publication record consulted by source~13; not used in place of a proof.

\bibitem{sebbar}
A. Sebbar,
\emph{Finite and infinite order differential relations for theta functions},
Milan Journal of Mathematics \textbf{84} (2016), 317--347.
\url{https://doi.org/10.1007/s00032-016-0261-6}.
Publisher abstract and bibliography consulted by source~13, and the
publisher abstract by source~14; full text not inspected by either. For the
merge of source~14 the bibliographic data were checked against the Crossref
record (volume~84, issue~2); the content was not verified.

\bibitem{bagayoko}
V. Bagayoko, L. S. Krapp, S. Kuhlmann, D. Panazzolo and M. Serra,
\emph{Automorphisms and derivations on algebras endowed with formal infinite
sums}, arXiv:2403.05827, version~2.
\url{https://arxiv.org/abs/2403.05827}.
Consulted by source~20 for positive-support summability.

\bibitem{kkslinrec}
L. S. Krapp, S. Kuhlmann and M. Serra,
\emph{Generalised power series determined by linear recurrence relations},
arXiv:2206.04126; published in Journal of Algebra (2025).
\url{https://arxiv.org/abs/2206.04126}.
Cited by source~20 to distinguish exponent-indexed recurrences from the
ordinary-indexed recurrences of the recurrence part; nothing is imported. A
different paper, and a different author list, from~\cite{kks}.

\bibitem{repoonot}
\repo\ project, \emph{Notations for Omnific Integers},
\texttt{docs/foundations-and-computation/omnific-notations/}. Its display
\texttt{onot:eq:guardbasic} is the element of Theorem~\ref{hol:rc:thm:coding};
cited by this merge, not by source~20.

\end{thebibliography}
\end{document}
